\xchapter{Advertisement}

THE following Tracts were published with the object of
contributing something towards the practical revival of doctrines,
which, although held by the great divines of our Church, at present have
become obsolete with the majority of her members, and are withdrawn from
public view even by the more learned and orthodox few who still adhere
to them. The Apostolic succession, the Holy Catholic Church, were
principles of action in the minds of our predecessors of the 17th
century; but, in proportion as the maintenance of the Church has been
secured bylaw, her ministers have been under the temptation of leaning
on an arm of flesh instead of her own divinely-provided discipline, a
temptation increased by political events and arrangements which need not
here be more than alluded to. A lamentable increase of sectarianism has
followed; being occasioned (in addition to other more obvious causes,)
first, by the cold aspect which the new Church doctrines have presented
to the religious sensibilities of the mind, next to their meagreness in
suggesting motives to restrain it from seeking out a more influential
discipline. Doubtless obedience to the law of the land, and the careful
maintenance of “decency and order,” (the topics in usage among us,)
are plain duties of the Gospel, and a reasonable ground for keeping in
communion with the Established Church; yet, if Providence has graciously
provided for our weakness more interesting and constraining motives, it
is a sin thanklessly to neglect them; just as it would be a mistake
to rest the duties of temperance or justice on the mere law of natural
religion, when they are mercifully sanctioned in the Gospel by the more
winning authority of our Saviour Christ. Experience has shown the
inefficacy of the mere injunctions of Church order, however scripturally
enforced, in restraining from schism the awakened and anxious sinner;
who goes to a dissenting preacher “because (as he expresses it) he
gets good from him:” and though he does not stand excused in God’s
sight for yielding to the temptation, surely the Ministers of the Church
are not blameless if, by keeping back the more gracious and consoling
truths provided for the little ones of Christ, they indirectly lead him
into it. Had he been taught as a child, that the Sacraments, not
preaching, are the sources of Divine Grace; that the Apostolical
ministry had a virtue in it which went out over the whole Church, when
sought by the prayer of faith; that fellowship with it was a gift and
privilege, as well as a duty, we could not have had so many wanderers
from our fold, nor so many cold hearts within it.

This instance may suggest many others of the
superior \emph{influence} of an apostolical over a mere secular method
of teaching. The awakened mind knows its wants, but cannot provide for
them; and in its hunger will feed upon ashes, if it cannot obtain the
pure milk of the word. Methodism and Popery are in different ways the
refuge of those whom the Church stints of the gifts of grace; they are
the foster-mothers of abandoned children. The neglect of the daily
service, the desecration of festivals, the Eucharist scantily
administered, insubordination permitted in all ranks of the Church,
orders and offices imperfectly developed, the want of Societies for
particular religious objects, and the like deficiencies, lead the
feverish mind, desirous of a vent to its feelings, and a stricter rule
of life, to the smaller religious Communities, to prayer and bible
meetings, and ill-advised institutions and societies, on the one
hand,—on the other, to the solemn and captivating services by which Popery gains its proselytes. Moreover, the multitude of men cannot
teach or guide themselves; and an injunction given them to depend on
their private judgment, cruel in itself, is doubly hurtful, as throwing
them on such teachers as speak daringly and promise largely, and not
only aid but supersede individual exertion.

These remarks may serve as a clue, for those who
care to pursue it, to the views which have led to the publication of the
following Tracts. The Church of Christ was intended to cope with human
nature in all its forms, and surely the gifts vouchsafed it are adequate
for that gracious purpose. There are zealous sons and servants of her
English branch, who see with sorrow that she is defrauded of her full
usefulness by particular theories and principles of the present age,
which interfere with the execution of one portion of her commission; and
while they consider that the revival of this portion of truth is
especially adapted to break up existing parties in the Church, and to
form instead a bond of union among all who love the Lord Jesus Christ in
sincerity, they believe that nothing but these neglected doctrines,
faithfully preached, will repress that extension of Popery, for which
the ever multiplying divisions of the religious world are too clearly
preparing the way.

OXFORD,

\emph{The Feast of All Saints}, 1834.

\xchapter{Thoughts On The Ministerial Commission, Respectfully Addressed to the Clergy}

Tract No. 1 (\emph{Ad Clerum})

I AM
but one of yourselves,—a Presbyter; and therefore I conceal my name,
lest I should take too much on myself by speaking in my own person. Yet
speak I must; for the times are very evil, yet no one speaks against
them.

Is not this so? Do not we ``look one upon
another,'' yet perform nothing? Do we not all confess the peril into
which the Church is come, yet sit still each in his own retirement, as
if mountains and seas cut off brother from brother? Therefore suffer me,
while I try to draw you forth from those pleasant retreats, which it has
been our blessedness hitherto to enjoy, to contemplate the condition and
prospects of our Holy Mother in a practical way; so that one and all may
unlearn that idle habit, which has grown upon us, of owning the state of
things to be bad, yet doing nothing to remedy it.

Consider a moment. Is it fair, is it dutiful, to
suffer our Bishops to stand the brunt of the battle without doing our
part to support them? Upon them comes ``the care of all the
Churches.'' This cannot be helped: indeed it is their glory. Not one of
us would wish in the least to deprive them of the duties, the toils, the
responsibilities of their high Office. And, black event as it would be
for the country, yet, (as far as they are concerned,) we could not wish
them a more blessed termination of their course, than the spoiling of
their goods, and martyrdom.

To them then we willingly and affectionately
relinquish their high privileges and honours; we encroach not upon the
rights of the SUCCESSORS
OF THE APOSTLES; we touch not their sword and crosier. Yet surely
we may be their shield-bearers in the battle without offence; and by our
voice and deeds be to them what Luke and Timothy were to St. Paul.

Now then let me come at once to the subject which
leads me to address you. Should the Government and Country so far forget
their GOD
as to cast off the Church, to deprive it of its temporal honours and
substance, \emph{on what} will you rest the claim of respect and
attention which you make upon your flocks? Hitherto you have been upheld
by your birth, your education, your wealth, your connexions; should
these secular advantages cease, on what must CHRIST'S
Ministers depend? Is not this a serious practical question? We know how
miserable is the state of religious bodies not supported by the State.
Look at the Dissenters on all sides of you, and you will see at once
that their Ministers, depending simply upon the people, become the
creatures of' the people. Are you content that this should be your
case? Alas! can a greater evil befal Christians, than for their teachers
to be guided by them, instead of guiding? How can we ``hold fast the
form of sound words,'' and ``keep that which is committed to our
trust,'' if our influence is to depend simply on our popularity? Is it
not our very office to \emph{oppose} the world? can we then allow
ourselves to \emph{court} it? to preach smooth things and prophesy
deceits? to make the way of life easy to the rich and indolent, and to
bribe the humbler classes by excitements and strong intoxicating
doctrine? Surely it must not be so;—and the question recurs, on what
are we to rest our authority, when the State deserts us?

CHRIST
has not left His Church without claim of its own upon the attention of
men. Surely not. Hard Master He cannot be, to bid us oppose the world,
yet give us no credentials for so doing. There are some who rest their
divine mission on their own unsupported assertion; others, who rest it
upon their popularity; others, on their success; and others, who rest it
upon their temporal distinctions. This last case has, perhaps, been too
much our own; I fear we have neglected the real ground on which our
authority is built,—our APOSTOLICAL
DESCENT.

We have been born, not of blood, nor of the will of
the flesh, nor of the will of man, but of GOD. The LORD JESUS CHRIST gave His SPIRIT to His Apostles; they in turn laid their
hands on those who should succeed them; and these again on others; and
so the sacred gift has been handed down to our present Bishops, who have
appointed us as their assistants, and in some sense representatives.

Now every one of us believes this. I know that some
will at first deny they do; still they do believe it. Only, it is not
sufficiently practically impressed on their minds. They \emph{do}
believe it; for it is the doctrine of the Ordination Service, which they
have recognised as truth in the most solemn season of their lives. In
order, then, not to prove, but to remind and impress, I entreat your
attention to the words used when you were made Ministers of CHRIST'S
Church.

The office of Deacon was thus committed to you: ``Take thou authority to execute the office of a Deacon in the Church
of GOD
committed unto thee: In the name,'' \&c.

And the priesthood thus:

``Receive the HOLY GHOST,
for the office and work of a Priest, in the Church of GOD, now committed unto thee by the
imposition of our hands. Whose sins thou dost forgive, they are
forgiven; and whose sins thou dost retain, they are retained. And be
thou a faithful dispenser of the Word of GOD, and of His Holy Sacraments: In the name,'' \&c.

These, I say, were words spoken to us, and received
by us, when we were brought nearer to GOD than at any other time of our lives. I know the grace of
ordination is contained in the laying on of hands, not in any form of
words;—yet in our own case, (as has ever been usual in the Church,)
words of blessing have accompanied the act. Thus we have confessed
before GOD
our belief, that through the Bishop who ordained us, we received the HOLY GHOST, the power to bind and to loose, to
administer the Sacraments, and to preach. Now \emph{how} is he able to
give these great gifts? \emph{Whence} is his right? Are these words
idle, (which would be taking GOD'S
name in vain,) or do they express merely a wish, (which surely is very
far below their meaning,) or do they not rather indicate that the
Speaker is conveying a gift? Surely they can mean nothing short of this.
But whence, I ask, his right to do so? Has he any right, except as
having received the power from those who consecrated him to be a Bishop?
He could not give what he had never received. It is plain then that he
but \emph{transmits}; and that the Christian Ministry is a \emph{succession}.
And if we trace back the power of ordination from hand to hand, of
course we shall come to the Apostles at last. We know we do, as a plain
historical fact: and therefore all we, who have been ordained Clergy, in
the very form of our ordination acknowledged the doctrine of the APOSTOLICAL
SUCCESSION.

And for the same reason, we must necessarily
consider none to be \emph{really} ordained who have not \emph{thus} been
ordained. For if ordination is a divine ordinance, it must be necessary;
and if it is not a divine ordinance, how dare we use it? Therefore all
who use it, all of \emph{us}, must consider it necessary. As well might
we pretend the Sacraments are not necessary to Salvation, while we make
use of the offices of the Liturgy; for when GOD
appoints means of grace, they are \emph{the} means.

I do not see how any one can escape from this plain
view of the subject, except, (as I have already hinted,) by declaring,
that the words do not mean all that they say. But only reflect what a
most unseemly time for random words is that, in which Ministers are set
apart for their office. Do we not adopt a Liturgy, \emph{in order to}
hinder inconsiderate idle language, and shall we, in the most sacred of
all services, write down, subscribe, and use again and again forms of
speech, which have not been weighed, and cannot be taken strictly?

Therefore, my dear Brethren, act up to your
professions. Let it not be said that you have neglected a gift; for if
you have the Spirit of the Apostles on you, surely this is a great gift.
``Stir up the gift of GOD
which is in you.'' Make much of it. Show your value of it. Keep it
before your minds as an honourable badge, far higher than that secular
respectability, or cultivation, or polish, or learning, or rank, which
gives you a hearing with the many. Tell \emph{them} of your gift. The
times will soon drive you to do this, if you mean to be still any thing.
But wait not for the times. Do not be compelled, by the world's
forsaking you, to recur as if unwillingly to the high source of your
authority. Speak out now, before you are forced, both as glorying in
your privilege, and to ensure your rightful honour from your people. A
notion has gone abroad, that they can take away your power. They think
they have given and can take it away. They think it lies in the Church
property, and they know that they have politically the power to
confiscate that property. They have been deluded into a notion that
present palpable usefulness, produceable results, acceptableness to your
flocks, that these and such like are the tests of your Divine
commission. Enlighten them in this matter. Exalt our Holy Fathers, the
Bishops, as the Representatives of the Apostles, and the Angels of the
Churches; and magnify your office, as being ordained by them to take
part in their Ministry.

But, if you will not adopt my view of the subject,
which I offer to you, not doubtingly, yet (I hope) respectfully, at all
events, CHOOSE
YOUR SIDE. To remain neuter much longer will be itself to take a
part. \emph{Choose} your side; since side you shortly must, with one or
other party, even though you do nothing. Fear to be of those, whose line
is decided for them by chance circumstances, and who may perchance find
themselves with the enemies of CHRIST,
while they think but to remove themselves from worldly politics. Such
abstinence is impossible in troublous times. HE
THAT IS NOT
WITH ME, IS AGAINST ME, AND HE THAT GATHERETH NOT WITH ME SCATTERETH
ABROAD.

[SIXTH EDITION.]

———————————————————————

These Tracts are continued in
Numbers, and sold at the price of 2d. for each sheet, or 7s. for 50
copies.

LONDON:
PRINTED FOR J. G. F. \& J. RIVINGTON,

ST. PAUL'S CHURCH YARD AND WATERLOO PLACE.

1840.

\xchapter{The Catholic Church}

Tract No. 2

No weapon that is formed against thee shall
prosper, and every tongue that shall rise against thee in judgment THOU SHALT CONDEMN.

IT
is sometimes said, that the Clergy should abstain from politics; and
that, if a Minister of CHRIST
is political, he is not a follower of him who said, ``My kingdom is not
of this world.'' Now there is a sense in which this is true, but, as it
is commonly taken, it is very false.

It is true that the mere affairs of this world
should not engage a Clergyman; but it is absurd to say that the affairs
of this world should not at all engage his attention. If so, this world
is not a preparation for another. Are we to speak when individuals sin,
and not when a nation, which is but a collection of individuals? Must we
speak to the poor, but not to the rich and powerful? In vain does St.
James warn us against having the faith of our LORD
JESUS CHRIST
with respect of persons. In vain does the Prophet declare to us the word
of the LORD,
that if the watchmen of Israel ``speak not to warn the wicked from his
way,'' ``his blood will be required at the watchman's hand.''

Complete our LORD'S declaration concerning the nature of His kingdom, and you
will see it is not at all inconsistent with the duty of our active and
zealous interference in matters of this world. ``If My kingdom were of
this world,'' He says, ``\emph{then would My servants fight}.''—Here
He has vouchsafed so to explain Himself, that there is no room for
misunderstanding His meaning. No one contends that His ministers ought
to use the weapons of a carnal warfare; but surely to protest, to warn,
to threaten, to excommunicate, are not such weapons. Let us not be
scared from a plain duty, by the mere force of a misapplied text. There
is an unexceptionable sense in which a clergyman may, nay, must be \emph{political}.
And above all, when the Nation interferes with the rights and
possessions of the Church, it can with even less grace complain of the
Church interfering with the Nation.

With this introduction let me call your attention
to what seems a most dangerous infringement on our rights, on the part
of the State. The Legislature has lately taken upon itself to remodel
the dioceses of Ireland; a proceeding which involves the appointment
of certain Bishops over certain Clergy, and of certain clergy under
certain Bishops, without the Church being consulted in the matter. I do
not say whether or not harm will follow from this particular act with
reference to Ireland; but consider whether it be not in itself an
interference with things spiritual.

Are we content to be accounted the mere creation of
the State, as schoolmasters and teachers may be, or soldiers, or
magistrates, or other public officers? Did the State make us? can it
unmake us? can it send out missionaries? can it arrange dioceses? Surely
all these are spiritual functions; and Laymen may as well set about
preaching, and consecrating the LORD'S
Supper, as assume these. I do not say the guilt is equal; but that, if
the latter is guilt, the former is. Would St. Paul, with his good will,
have suffered the Roman power to appoint Timothy, Bishop of Miletus, as
well as of Ephesus? Would Timothy at such a bidding have undertaken the
charge? Is not the notion of such an order, such an obedience, absurd?
Yet has it not been realized in what has lately happened? For in what is
the English state at present different from the Roman formerly? Neither
can be accounted members of the Church of CHRIST.
No one can say the British Legislature is in our communion, or that its
members are necessarily even Christians. What pretence then has it for
not merely advising, but superseding the Ecclesiastical power?

Bear with me, while I express my fear, that we do
not, as much as we ought, consider the force of that article of our
Belief, ``The One Catholic and Apostolic Church.'' This is a tenet so
important as to have been in the Creed from the beginning. It is
mentioned there as a \emph{fact}, and a fact \emph{to be believed}, and
therefore practical. Now what do we conceive is meant by it? As people
vaguely take it in the present day, it seems only an assertion that
there is a number of sincere Christians scattered through the world. But
is not this a truism? who doubts it? who can deny that there are people
in various places who are sincere believers? what comes of this? how is
it important? why should it be placed as an article of faith, after the
belief in the HOLY
GHOST?
Doubtless the only true and satisfactory meaning is that which our
Divines have ever taken, that there is on earth an existing Society,
Apostolic as founded by the Apostles, Catholic because it spreads its
branches in every place; i.e. the Church Visible with its Bishops, Priests, and Deacons. And this surely is a most important doctrine; for
what can be better news to the bulk of mankind than to be told that CHRIST
when He ascended, did not leave us orphans, but appointed
representatives of Himself to the end of time?

``The necessity of believing the Holy Catholic
Church,'' says Bishop Pearson in his Exposition of the Creed, ``appeareth first in this, that CHRIST
hath appointed it as the only way to eternal life ... CHRIST never appointed two ways to heaven, nor
did He build a Church to save some, and make another institution for
other men's salvation. There is none other name under heaven given
among men whereby we must be saved, but the name of JESUS; and that name is no otherwise given under heaven than in
the Church,'' ``This is the congregation of those persons here on
earth which shall hereafter meet in heaven … There is a necessity of
believing the Catholic Church, because except a man be of that he can be
of none. Whatsoever Church pretendeth to a new beginning, pretendeth at
the same time to a new Churchdom, and whatsoever is so new is none.''
This indeed is the unanimous opinion of our divines, that, as the
Sacraments, so Communion with the Church, is ``generally necessary to
salvation,'' in the case of those who can obtain it.

If then we express our belief in the existence of
One Church on earth from CHRIST'S
coming to the end of all things, if there is a promise it shall
continue, and if it is our duty to do our part in our generation towards
its continuance, how can we with a safe conscience countenance the
interference of the Nation in its concerns? Does not such interference
tend to destroy it? Would it not destroy it, if consistently followed
up? Now, may we sit still and keep silence, when efforts are making to
break up, or at least materially to weaken that Ecclesiastical Body
which we know is intended to last while the world endures, and the
safety of which is committed to our keeping in our day? How shall we
answer for it, if we transmit that Ordinance of GOD less entire than it came to us?

Now what am I calling on you to do? You cannot help
what has been done in Ireland; but you may protest against it. You may
as a duty protest against it in public and private; you may keep a
jealous watch on the proceedings of the Nation, lest a second act of the
same kind be attempted. You may keep it before you as a desirable object
that the Irish Church should at some future day meet in Synod and
protest herself against what has been done; and then proceed to
establish or rescind the State injunction, as may be thought expedient.

I know it is too much the fashion of the times to
think any earnestness for ecclesiastical rights unseasonable and absurd,
as if it were the feeling of those who live among books and not in the
world. But it is our \emph{duty} to live among books, especially to live
by ONE BOOK, and a very old one; and therein we
are enjoined to ``keep that good thing which is committed unto us,''
to ``neglect not our gift.'' And when men talk, as they sometimes do,
as if in opposing them we were standing on technical difficulties
instead of welcoming great and extensive benefits which would be the
result of their measures, I would ask them, (letting alone the question
of their beneficial nature, which \emph{is} a question,) whether this is
not being wise above that is written, whether it is not doing evil that
good may come. We cannot know the effects which will follow certain
alterations; but we can decide that the means by which it is proposed to
attain them are unprecedented and disrespectful to the Church. And when
men say, ``\emph{the day is past} for stickling about ecclesiastical
rights,'' let them see to it, lest they use substantially the same
arguments to maintain their position as those who say, ``The day is
past for being a Christian.''

Lastly, is it not plain that by showing a bold
front and defending the rights of the Church, we are taking the only
course which can make us respected? Yielding will not persuade our
enemies to desist from their efforts to destroy us root and branch. We
cannot hope by giving something to keep the rest. Of this surely we have
had of late years sufficient experience. But by resisting strenuously,
and contemplating and providing against the worst, we may actually
prevent the very evils we fear. To prepare for persecution may be the
way to avert it.

[FOURTH EDITION.]

———————————————————————

These Tracts are continued in
Numbers, and sold at the price of 2d. for each sheet, or 7s. for 50
copies.

LONDON: PRINTED FOR J. G. F. \& J. RIVINGTON,

ST. PAUL'S CHURCH YARD, AND WATERLOO PLACE.

1840.

\xchapter{Thoughts Respectfully Addressed to the Clergy On Alterations in the Liturgy}

Tract No. 3

ATTEMPTS
are making to get the Liturgy altered. My dear Brethren, I beseech you,
consider with me, whether you ought not to resist the alteration of even
one jot or tittle of it. Though you would in your own private judgments
wish to have this or that phrase or arrangement amended, is this a time
to concede one tittle?

Why do I say this? because, though most of you
would wish some immaterial points altered, yet not many of you agree in
those points, and not many of you agree what is and what is not
immaterial. If all your respective emendations are taken, the
alterations in the Services will be extensive; and though each will gain
something he wishes, he will lose more from those alterations which he
did not wish. Tell me, are the present imperfections (as they seem to
each) of such a nature, and so many, that their removal will compensate
for the recasting of much which each thinks to be no imperfection, or
rather an excellence?

There are persons who wish the Marriage Service
emended; there are others who would be indignant at the changes
proposed. There are some who wish the Consecration Prayer in the Holy
Sacrament to be what it was in King Edward's first book; there are
others who think this would be an approach to Popery. There are some who
wish the imprecatory Psalms omitted; there are others who would lament
this omission as savouring of the shallow and detestable liberalism of
the day. There are some who wish the Services shortened; there are
others who think we should have far more Services, and more frequent
attendance at public worship than we have.

How few would be pleased by \emph{any given}
alterations; and how many pained!

But once begin altering, and there will be no
reason or justice in stopping, till the criticisms of all parties are
satisfied. Thus, will not the Liturgy be in the evil case described in
the well-known story, of the picture subjected by the artist to the
observations of passers-by? And, even to speak at present of
comparatively immaterial alterations, I mean such as do not infringe
upon the doctrines of the Prayer Book, will not it even with these be a
changed book, and will not that new book be for certain an inconsistent
one, the alterations being made, not on principle, but upon chance
objections urged from various quarters?

But this is not all. A taste for criticism grows
upon the mind. When we begin to examine and take to pieces, our judgment
becomes perplexed, and our feelings unsettled. I do not know whether
others feel this to the same extent, but for myself, I confess there are
few parts of the Service that I could not disturb myself about, and feel
fastidious at, if I allowed my mind in this abuse of reason. First, e.g.
I might object to the opening sentences; ``they are not evangelical
enough; CHRIST is not mentioned in them; they
are principally from the Old Testament.'' Then I should criticise the
exhortation, as having too many words, and as antiquated in style. I
might find it hard to speak against the Confession; but ``the Absolution,''
it might be said, ``is not strong enough; it is a mere declaration, not
an announcement of pardon to those who have confessed.'' And so on.

Now I think this unsettling of the mind a frightful
thing; both to ourselves, and more so to our flocks. They have long
regarded the Prayer Book with reverence as the stay of their faith and
devotion. The weaker sort it will make sceptical; the better it will
offend and pain. Take, e.g. an alteration which some have offered in the
Creed, to omit or otherwise word the clause, ``He descended into \emph{hell}.''
Is it no comfort for mourners to be told that CHRIST
Himself has been in that unseen state, or Paradise, which is the alloted
place of sojourn for departed spirits? Is it not very easy to explain
the ambiguous word, is it any great harm if it is misunderstood, and is
it not very difficult to find any substitute for it in harmony with the
composition of the Creed? I suspect we should find the best men in the
number of those who would retain it as it is. On the other hand, will
not the unstable learn from us a habit of criticising what they should
never think of but as a divine voice supplied by the Church for their
need?

But as regards ourselves, the Clergy, what will be
the effect of this temper of innovation in us? We have the power to
bring about changes in the Liturgy; shall we not exert it? have we any security, if we once begin, that we shall ever end? Shall not we
pass from non-essentials to essentials? And then, on looking back after
the mischief is done, what excuse shall we be able to make for ourselves
for having encouraged such proceedings at first? Were there grievous
errors in the Prayer Book, something might be said for beginning, but
who can point out any? cannot we very well \emph{bear} things as they
are? does any part of it seriously disquiet us? no—we have before now
freely given our testimony to its accordance with Scripture.

But it may be said that ``we must conciliate an
outcry which is made; that some alteration is demanded.'' By whom? no one
can tell who cries, or who can be conciliated. Some of the laity, I
suppose. Now consider this carefully. Who are these lay persons? Are
they serious men, and are their consciences involuntarily hurt by the
things they wish altered? Are they not rather the men you meet in
company, worldly men, with little personal religion, of lax conversation
and lax professed principles, who sometimes perhaps come to Church, and
then are wearied and disgusted? Is it not so? You have been dining,
perhaps, with a wealthy neighbour, or fall in with this great Statesman,
or that noble Land-holder, who considers the Church two centuries behind
the world, and expresses to you wonder that its enlightened members do
nothing to improve it. And then you get ashamed, and are betrayed into
admissions which sober reason disapproves. You consider, too, that it is
a great pity so estimable or so influential a man should be disaffected
to the Church; and you go away with a vague notion that something must
be done to conciliate such persons. Is this to bear about you the solemn
office of a GUIDE
and TEACHER
in Israel, or to \emph{follow a lead}?

But consider what are the concessions which would
conciliate such men. Would immaterial alterations? Do you really think
they care one jot about the verbal or other changes which some
recommend, and others are disposed to grant? whether ``the unseen state''
is substituted for ``hell,'' ``condemnation'' for ``damnation,'' or the order
of Sunday Lessons is remodeled? No;—they dislike the \emph{doctrine}
of the Liturgy. These men of the world do not like the anathemas of the
Athanasian Creed, and other such peculiarities of our Services. But even
were the alterations, which would please them, small, are they the
persons whom it is of use, whom it is becoming to conciliate by
going out of our way?

I need not go on to speak against doctrinal
alterations, because most thinking men are sufficiently averse to them.
But, I earnestly beg you to consider whether we must not come to them if
we once begin. For by altering immaterials, we merely \emph{raise}
without \emph{gratifying} the desire of correcting; we excite the
craving, but withhold the food. And it should be observed, that the
changes called immaterial often contain in themselves the germ of some
principle, of which they are thus the introduction:—e.g. If we were to
leave out the imprecatory Psalms, we certainly countenance the notion of
the day, that love and love only is in the Gospel the character of ALMIGHTY
GOD and
the duty of regenerate man; whereas that Gospel, rightly understood,
shows His Infinite Holiness and Justice as well as His Infinite Love;
and it enjoins on men the duties of zeal towards Him, hatred of sin, and
separation from sinners, as well as that of kindness and charity.

To the above observations it may be answered, that
changes have formerly been made in the Services without leading to the
issue I am predicting now; and therefore they may be safely made again.
But, waving all other remarks in answer to this argument, is not this
enough, viz, that there \emph{is} peril? No one will deny that the rage
of the day is for concession. Have we not already granted (political)
points, without stopping the course of innovation? This is a fact. Now,
is it worth while even to risk fearful changes merely to gain petty
improvements, allowing those which are proposed to be such?

We know not what is to come upon us; but the writer
for one will try so to acquit himself now, that if any irremediable
calamity befalls the Church, he may not have to vex himself with the
recollections of silence on his part and indifference, when he might
have been up and alive. There was a time when he, as well as others,
might feel the wish, or rather the temptation, of steering a middle
course between parties; but if so, a more close attention to passing
events has cured his infirmity. In a day like this there are but two
sides, zeal and persecution, the Church and the world; and those who
attempt to occupy the ground between them, at best will lose their
labour, but probably will be drawn back to the latter. Be practical, I
respectfully urge you; do not attempt impossibilities; sail not as
if in pleasure boats upon a troubled sea. Not a word falls to the
ground, in a time like this. Speculations about ecclesiastical
improvements which might be innocent at other times, have a strength of
mischief now. They are realized before he who utters them understands
that he has committed himself.

Be prepared then for petitioning against any
alterations in the Prayer Book which may be proposed. And, should you
see that our Fathers the Bishops seem to countenance them, petition
still. Petition \emph{them}. They will thank you for such a proceeding. \emph{They
do not wish these alterations}; but how can they resist them without
the support of their Clergy? They consent to them, (if they do,) partly
from the notion that they are thus pleasing you. Undeceive them. They
will be rejoiced to hear that you are as unwilling to receive them as
they are. However, if after all there be persons determined to allow
some alterations, then let them quickly make up their minds \emph{how far}
they will go. They think it easier to draw the line elsewhere, than as
things now exist. Let them point out the limit of their concessions now;
and let them keep to it then; and, (if they can do this,) I will say
that, though they are not as wise as they might have been, they are at
least firm, and have at last come right.

THE BURIAL SERVICE

WE
hear many complaints about the Burial Service, as unsuitable for the use
for which it was intended. It expresses a hope, that the person
departed, over whom it is read, will be saved; and this is said to be
dangerous when expressed about all who are called Christians, as leading
the laity to low views of the spiritual attainments necessary for
salvation; and distressing the Clergy who have to read it.

Now I do not deny, I frankly own, it is sometimes
distressing to use the Service; but this it must ever be in the nature
of things; wherever you draw the line. Do you pretend you can
discriminate the wheat from the tares? of course not.

It is often distressing to use this Service,
because it is often distressing to think of the dead at all; not that
you are without hope, but because you have fear also.

\emph{How} many are there whom you know well enough
to dare to give any judgment about? Is a Clergyman only to express a
hope where \emph{he} has grounds for having it? Are not the feelings of
relatives to be considered? And may there not be a difference of
judgments? I may hope more, another less. If each is to use the precise
words which suit his own judgment, then we can have no words at all.

But it may be said, ``every thing of a \emph{personal}
nature may be left out from the Service.'' And do you really wish this?
Is this the way in which your flock will wish their lost friends to be
treated? a cold ``edification,'' but no affectionate valediction to the
departed? Why not pursue this course of (supposed) improvement, and
advocate the omission of the Service altogether.

Are we to have no kind and religious thoughts over
the good, lest we should include the bad?

But it will be said, that, at least we ought not to
read the Service over the flagrantly wicked; over those who are a
scandal to religion. But this is a very different position. I agree with
it entirely. Of course we should not do so, and truly the Church never
meant we should. She never wished we should profess our hope of the
salvation of habitual drunkards and swearers, open sinners, blasphemers,
and the like; not as daring to despair of their salvation, but thinking
it unseemly to honour their memory. Though the Church is not endowed
with a power of absolute judgment upon individuals, yet she is directed
to decide according to external indications, in order to hold up the \emph{rules}
of GOD'S
governance, and afford a type of it, and an assistance towards the
realizing it. As she denies to the scandalously wicked the LORD'S Supper, so does she deprive them of her
other privileges.

The Church, I say, does not bid us read the Service
over open sinners. Hear her own words introducing the Service. ``The
office ensuing is not to be used for any that die unbaptized, or
excommunicate, or have laid violent hands upon themselves.'' There is no
room to doubt \emph{whom} she meant to be excommunicated, open sinners.
Those therefore who are pained at the general use of the Service, should
rather strive to restore the practice of excommunication, than to alter
the words used in the Service. Surely, if we do not this, we are
clearly defrauding the religious, for the sake of keeping close to the
wicked.

Here we see the common course of things in the
world. We omit a duty. In consequence our services become inconsistent.
Instead of retracing our steps we alter the Service. What is this but,
as it were, to sin upon principle? While we keep to our principles, our
sins are inconsistencies; at length, sensitive of the absurdity which
inconsistency involves, we accommodate our professions to our practice.
This is ever the way of the world; but it should not be the way of the
church.

I will join heart and hand with any who will
struggle for a restoration of that ``godly discipline,'' the restoration
of which our Church publicly professes she considers desirable; but GOD forbid any one should so depart from her
spirit, as to mould her formularies to fit the case of deliberate
sinners! And is not this what we are plainly doing, if we alter the
Burial Service as proposed? we are recognizing the right of men to
receive Christian Burial, about whom we do not like to express a hope.
Why should they have Christian burial at all?

It will be said that the restoration of the
practice of Excommunication is impracticable; and that therefore the
other alternative must be taken, as the only one open to us. Of course
it is impossible, if no one attempts to restore it; but if all willed
it, how would it be impossible; and if no one stirs because he thinks no
one else will, he is arguing in a circle.

But, after all, what have we to do with
probabilities and prospects in matters of plain duty? Were a man the
only member of the Church who felt it a duty to return to the Ancient
Discipline, yet a duty is a duty, though he be alone. It is one of the
great sins of our times to look to consequences in matters of plain
duty. Is not this such a case? If not, prove that it is not; but do not
argue from consequences.

In the mean while I offer the following texts in
evidence of the duty.

Matth. xviii. 15-17. Rom. xvi. 17. 1 Cor. v. 7-13.
2 Thess. iii. 6, 14, 15. 2 Tim. iii. 5. Tit. iii. 10, 11. 2 John 10, 11.

THE PRINCIPLE OF UNITY

Testimony of St. Clement, the associate of St.
Paul, (Phil. iv. 3.) to the Apostolical Succession.

The
Apostles knew, through our LORD JESUS
CHRIST,
that strife would arise for the Episcopate. Wherefore having received an
accurate foreknowledge, they appointed the men I before mentioned, and
have given an orderly succession, that on their death other approved men
might receive in turn their office. Ep. i. 44.

Testimony of St. Ignatius, the friend of St. Peter,
to Episcopacy.

Your
celebrated Presbytery, worthy of GOD, is as closely knit to the Bishop, as the strings to a harp, and so by
means of your unanimity and concordant love JESUS CHRIST is
sung. Eph. 4.

There are who profess to
acknowledge a Bishop, but do every thing without him. Such men appear to
lack a clear conscience. Magn. 4.

He
for whom I am bound is my witness that I have not learned this doctrine
from mortal man. The Spirit proclaimed to me these words: ``Without the
Bishop do nothing.'' Phil. 7.

With these and other such strong passages in the
Apostolical Fathers, how can we permit ourselves in our present \emph{practical}
disregard of the Episcopal Authority? Are not we apt to obey only so far
as the law obliges us? Do we support the Bishop, and strive to move all
together with him as our bond of union and head; or is not our every-day
conduct as if, except with respect to certain periodical forms and
customs, we were each independent in his own parish?

[FIFTH EDITION.]

———————————————————————

These Tracts are continued in
Numbers, and sold at the price of 2d. for each sheet, or 7s. for 50
copies.

LONDON: PRINTED FOR J. G. F. \& J. RIVINGTON,

ST. PAUL'S CHURCH YARD, AND WATERLOO PLACE.

1840.

\xchapter{The Present Obligation of Primitive Practice}

Tract No. 6 (\emph{Ad Populum})

WHEN
we look around upon the present state of the Christian Church, and then
turning to ecclesiastical history acquaint ourselves with its primitive
form and condition, the difference between them so strongly acts upon
the imagination, that we are tempted to think, that to base our conduct
now on the principles acknowledged then, is but theoretical and idle. We
seem to perceive, as clear as day, that as a Primitive Church had its
own particular discipline and political character, so have we ours; and
that to attempt to revive what is past, is as absurd as to seek to raise
what is literally dead. Perhaps we even go on to maintain, that the
constitution of the Church, as well as its actual course of acting, is
different from what it was; that Episcopacy now is in no sense what it
used to be; that our Bishops are the same as the Primitive Bishops only
in name; and that the notion of an Apostolical Succession is ``a fond
thing.'' I do not wish to undervalue the temptation, which leads to this
view of Church matters; it is the temptation of sight to overcome faith,
and of course not a slight one.

But the following reflection on the history of the
Jewish Church may perhaps be considered to throw light on our present
duties.

1. Consider how exact are the injunctions of Moses
to his people. He ends them thus: ``These are the words of the covenant
which the LORD
commanded Moses to make with the children of Israel in the land of Moab,
beside the covenant which He made with them in Horeb … Keep therefore
the words of this covenant, and do them, that ye may prosper in all that
ye do … Neither with you only do I make this covenant and this oath;
but with him that standeth here this day before the LORD
our GOD,
and also with him that is not here with us this day.'' Deut. xxix.

2. Next, survey the history of the chosen people
for the several first centuries after taking possession of Canaan. The
exactness of Moses was unavailing. Can a greater contrast be conceived
than the commands and promises of the Pentateuch, and the history of the
Judges? ``Every man did that which was right in his own eyes.'' Judges
xvii. 6.

Samuel attempts a reformation on the basis of the
Mosaic Law; but the effort ultimately fails, as being apparently against
the stream of opinion and feeling then prevalent. The times do not allow
of it. Again, contrast the opulent and luxurious age of Solomon, though
the covenant was then openly acknowledged and outwardly accepted more
fully than at any other time, with the vision of simple piety and plain
straightforward obedience, which is the scope of the Mosaic Law. Lastly,
contemplate the state of the Jews after their return from the captivity;
when their external political relations were so new, the internal
principle of their government so secular, GOD'S arm apparently so far removed.
This state of things went on for centuries. Who would suppose that the
Jewish Law was binding in all its primitive strictness at the age when CHRIST
appeared? Who would not say that length of time had destroyed the
obligation of a projected system, which had as yet never been realized?

Consider, too, the impossible nature (so to say) of
some of its injunctions. An infidel historian somewhere asks scoffingly,
whether ``the ruinous law which required all the males of the chosen
people to go up to Jerusalem three times a year, was ever observed in
its strictness.'' The same question may be asked concerning the
observance of the Sabbatical year;—to which but a faint allusion, if
that, is made in the books of Scripture subsequent to the Pentateuch.

3. And now, with these thoughts before us, reflect
upon our SAVIOUR'S
conduct. He set about to fulfil the Law in its strictness, just as if He
had lived in the generation next to Moses. The practice of others, the
course of the world, was nothing to Him; He received and He obeyed. It
is not necessary to draw out the evidence of this in detail. Consider
merely His emphatic words in the beginning of Matt. xxiii. concerning
those, whom as individuals He was fearfully condemning. ``The Scribes
and Pharisees sit in Moses' seat: all therefore whatsoever they bid you
observe, that observe and do.''—Again, reflect upon the praise bestowed
upon Zacharias and his wife, that ``they were both righteous before GOD,
walking in all the commandments and ordinances of the LORD blameless.''—And upon the conduct
of the Apostles.

Surely these remarkable facts impress upon us the
necessity of going to the Apostles, and not to the teachers and oracles
of the present world, for the knowledge of our duty, as individuals and
as members of the Christian Church. It is no argument against a practice
being right, that it is neglected; rather we are warned against going
the broad way of the multitude of men.

Now is there any doubt in our minds, as to the
feelings of the Primitive Church regarding the doctrine of the
Apostolical Succession? Did not the Apostles observe, even in an age of
miracles, the ceremony of Imposition of Hands? And are not we bound, not
merely to acquiesce in, but zealously to maintain and inculcate the
discipline which they established?

The only objection which can be made to this view
of our duty, is, that the injunction to obey strictly is \emph{not}
precisely given to us, as it was in the instance of the Mosaic Law. But
is not the real state of the case merely this; that the Gospel appeals
rather to our love and faith, our divinely illuminated reason, and the
free principle of obedience, than to the mere letter of its injunctions?
And does not the conduct of the Jews just prove to us, that, \emph{though}
the commands of CHRIST were put before us ever so precisely,
yet there would not be found in any extended course of history a more
exact attention to them, than there is now; that the difficulty of
resisting the influence, which the world's actual proceedings exert upon
our imagination, would be just as great as we find it at present?

A SIN OF THE CHURCH

Remember from whence thou art fallen, and repent,
and do thy first works; or else I will come unto thee quickly, and will
remove thy candlestick out of his place, except thou repent.

———————

The following extract is from Bingham, Antiq. xv.
9.

In
the primitive ages, it was both the rule and practice of all in general,
both Clergy and Laity, to receive the Communion every LORD'S
day … As often as they met together for Divine Service on the LORD'S day, they were obliged to receive the Eucharist
under pain of Excommunication … And if we run over the whole history
of the three first ages, we shall find this to have been the Church's
constant practice ... We are assured farther, that in some places they
received the Communion every day.

Is there any one who will deny, that the Primitive
Church is the best expounder in this matter of our SAVIOUR'S will as conveyed through His
Apostles?

Can a learned Church, such as the English, plead
ignorance of His will thus ascertained?

Do we fulfil it?

Is not the regret and concern of pious and learned
writers among us, such as Bingham, at our neglect of it, upon record?

And is it not written, ``THAT SERVANT WHICH KNEW HIS LORD'S WILL, AND PREPARED
NOT HIMSELF, NEITHER DID ACCORDING TO HIS WILL, SHALL BE BEATEN WITH
MANY STRIPES?''

And putting aside this disobedience, can we wonder,
that faith and love wax cold, when we so seldom partake of the MEANS,
mercifully vouchsafed to us, of communion with our LORD and SAVIOUR?

OXFORD,

\emph{Oct}. 29, 1833.

[FIFTH
EDITION.]

———————————————————————

These Tracts are continued in
Numbers, and sold at the price of 2d. for each sheet, or 7s. for 50
copies.

LONDON: PRINTED FOR
J. G. F. \& J. RIVINGTON,

ST. PAUL'S CHURCH YARD, AND WATERLOO PLACE.

1840.

\xchapter{The Episcopal Church Apostolical}

Tract No. 7

THERE
are many persons at the present day, who, from not having turned their
minds to the subject, think they are Churchmen in the sense in which the
early Christians were, merely because they are Episcopalians. The extent
of their Churchmanship is, to consider that Episcopacy is the best form
of Ecclesiastical Polity; and, again, that it originated with the
Apostles. I am far from implying, that to go thus far is nothing; or is
not an evidence (for it is,) of a reverent and sober temper of mind;
still the view is defective. It is defective, because the expediency of
a system, though a very cogent, is not the highest line of argument that
may be taken in its defence: and because an opponent may deny the fact
of the Apostolicity of Episcopacy, and so involve its maintainer in an
argument. Doubtless the more clear and simple principle for a Churchman
to hold, is that of a \emph{Ministerial Succession}; which is undeniable
as a fact, while it is most reasonable as a doctrine, and sufficiently
countenanced in Scripture for its practical reception. Of this,
Episcopacy, i.e. \emph{Superintendence}, is but an accident; though, for
the sake of conciseness, it is often spoken of by us as synonymous with
it. It shall be the object of the following Tract to insist upon this
higher characteristic of our Church.

My position then is this;—that the Apostles
appointed successors to their ministerial office, and the latter in turn
appointed others, and so on to the present day;—and further, that the
Apostles and their Successors have in every age committed portions of
their power and authority to others, who thus become their delegates,
and in a measure their representatives, and are called Priests and
Deacons. The result is an Episcopal system, \emph{because} of the
practice of delegation; but we may conceive their keeping their powers
altogether to themselves, and in the same proportion in which this was
done would the Church polity cease to be Episcopalian. We may conceive
the Order of Apostolic Vicars (so to call it,) increased, till one
of them was placed in every village, and took the office of parish
Priest. I do not say such a measure would be justifiable or
pious;—doubtless it would be a departure from the rule of
antiquity—but it is conceivable; and it is useful to conceive it, in
order to form a clear notion of the Essence of the Church System, and
the defective state of those Christian Societies which are separate from
the Church Catholic. It is a common answer made to those who are called
High Churchmen, to say, that, ``if GOD had intended the \emph{form} of Church Government to be of
great consequence, He would have worded His will in this matter more
clearly in Scripture.'' Now enough has already been said to show the
irrelevancy of such a remark. We need not deny to the Church the
abstract right, (however we may question the propriety,) of altering its
own constitution. It is not merely \emph{because} Episcopacy is a \emph{better
or more scriptural form} than Presbyterianism, (true as this may be
in itself,) that Episcopalians are right, and Presbyterians are wrong;
but because the Presbyterian Ministers have assumed a power, which was
never intrusted to them. They have presumed to exercise the power of
ordination, and to perpetuate a succession of ministers, without having
received a commission to do so. This is the plain fact that condemns
them; and is a standing condemnation, from which they cannot escape,
except by artifices of argument, which will serve equally to protect the
self-authorized teacher of religion. If \emph{they} may ordain without
being sent to do so, \emph{others} may teach and preach without being
sent. They hold a middle position, which is untenable as destroying
itself; for if Christians can do without Bishops (i.e. Commissioned
Ordainers), they way do without Commissioned Ministers (i.e. the Priests
and Deacons). If an imposition of hands is necessary to convey one gift,
why should it not be to convey another?

1. As to the \emph{fact} of the Apostolical
Succession, i.e. that our present Bishops are the heirs and
representatives of the Apostles by successive transmission of the
prerogative of being so, this is too notorious to require proof. Every
link in the chain is known, from St. Peter to our present Metropolitans.
Here then I only ask, looking at this plain fact by itself, is there not
something of a divine providence in it? can we conceive that this
Succession has been preserved, all over the world, amid many
revolutions, through many centuries, \emph{for nothing}? Is it wise or
pious to despise or neglect a gift thus transmitted to us in matter
of fact, even if Scripture did not touch upon the subject?

2. Next, consider how \emph{natural} is the
doctrine of a succession. When an individual comes to me, claiming to
speak in the name of the Most High, it is natural to ask him for his
authority. If he replies, that we are all bound to instruct each other,
this reply is intelligible, but in the very form of it excludes the
notion of a ministerial order, i.e. a class of persons set apart \emph{from}
others for religious offices. If he appeals to some miraculous gift,
this too is intelligible, and only unsatisfactory when the alleged gift
is proved to be a fiction. No other answer can be given except a
reference to some person, who has given him license to exercise
ministerial functions; then follows the question, \emph{how} that
individual gained his authority to do so. In the case of the Catholic
Church, the person referred to, i.e. the Bishop, has received it from a
predecessor, and he from another, and so on, till we arrive at the
Apostles themselves, and thence our LORD
and. SAVIOUR.
It is superfluous to dwell on so plain a principle, which in matters of
this world we act upon daily.

3. Lastly, the \emph{argument from Scripture} is
surely quite clear to those who honestly wish direction for \emph{practice}.
CHRIST
promised He should be with His Apostles always, as ministers of His
religion, even unto the end of the world. In one sense the Apostles were
to be alive till He came again; but they all died at the natural time.
Does it not follow, that there are those now alive who represent them?
Now who were the most probable representatives of them in the generation
next their death? They surely, whom they had ordained to succeed them in
the ministerial work. If any persons could be said to have CHRIST'S
power and presence, and the gifts of ruling and ordaining, of teaching,
of binding and loosing, (and comparing together the various Scriptures
on the subject, all these seem included in His promise to be with the
Church always,) surely those, on whom the Apostles laid their hands,
were they. And so in the next age, if any were representatives of the
first representatives, they must be the next generation of Bishops, and
so on. Nor does it materially alter the argument, though we suppose the
blessing upon Ministerial Offices made, not to the Apostles, but to the
whole body of Disciples; i.e. the Church. For, even if it be the Church
that has the power of ordination committed to it, still it exercises
it through the Bishops as its organs; and the question recurs, \emph{how}
has the Presbytery in this or that country obtained the power? The
Church certainly has from the first committed it to the Bishops, and as
never resumed it; and the Bishops have no where committed it to the
Presbytery, who therefore cannot be in possession of it.

However, it is merely for argument sake that I make
this allowance, as to the meaning of the text in Matt. xxviii.; for our
LORD'S
promise of His presence ``unto the end of the world,'' was made to the
Apostles, \emph{by themselves}. At the same time, let it be observed
what force is added to the argument for the Apostolical Succession, by
the acknowledged existence in Scripture of the doctrine of a standing
Church, or permanent Body Corporate for spiritual purposes. For, if
Scripture has formed all Christians into one continuous community
through all ages, (which I do not here prove,) it is but according to
the same analogy, that the Ministerial Office should be vested in an
order, propagated from age to age, on the principle of Succession. And,
if we proceed to considerations of utility and expedience, it is plain
that, according to our notions, it is more necessary that a Minister
should be perpetuated by a fixed law, than that the community of
Christians should be, which can scarcely be considered to be vested with
any powers such as to require the visible authority which a Succession
supplies.

[SIXTH EDITION.]

———————————————————————

These Tracts are continued in
Numbers, and sold at the price of 2d. for each sheet, or 7s. for 50
copies.

LONDON: PRINTED FOR
J. G. F. \& J. RIVINGTON,

ST. PAUL'S CHURCH YARD, AND WATERLOO PLACE.

1840.

\xchapter{The Gospel a Law of Liberty}

Tract No. 8

IT
is a matter of surprise to some persons, that the ecclesiastical system
under which we find ourselves, is so faintly enjoined on us in
Scripture. One very sufficient explanation of the fact will be found in
considering that the Bible is not intended to teach us matters of \emph{discipline}
so much as matters of \emph{faith}; i.e. those doctrines, the reception
of which are necessary to salvation. But another reason may be
suggested, which is well worth our attentive consideration.

The Gospel is a Law of \emph{Liberty}. We are
treated as sons, not as servants; not subjected to a code of formal
commands, but addressed as those who \emph{love} GOD, and \emph{wish} to please Him. When a man gives orders to
those whom he thinks will mistake him, or are perverse, he speaks
pointedly and explicitly; but when he gives directions to friends, he
will trust much to their knowledge of his feelings and wishes, he leaves
much to their discretion, and tells them not so much what he would have
done in detail, as what are the objects he would have accomplished. Now
this is the way CHRIST
has spoken to us under the New Covenant; and apparently with this
reason, to \emph{try} us, whether or not we really love Him as our LORD and SAVIOUR.

Accordingly, there is no part, perhaps, of the
ecclesiastical system, which is not faintly traced in Scripture, and no
part which is much more than faintly traced. The question which a
reverend and affectionate faith will ask, is, ``what is \emph{most likely}
to please CHRIST?''
And this is just the question that obtains an answer in Scripture; which
contains just so much as \emph{intimations} of what is most likely to
please Him. Of course different minds will differ as to the degree of
clearness with which this or that practice is enjoined, yet I think no
one will consider the state of the case, as I have put it, exaggerated
on the whole.

Many duties are intimated to us by example, not by
precept—many are implied merely—others can only be inferred from a
comparison of passages—and others perhaps are contained only in the
Jewish Law. I will mention some specimens to assist the reflection of
the reader.

The early Christians were remarkable for keeping to
the Apostles' \emph{fellowship}. Who are \emph{more likely} to stand in
the Apostles' place since their death, than that line of Bishops which
they themselves began? for that the Apostles \emph{were} in some sense
or other to remain on earth to the end of all things, is plain from the
text, ``Lo, I am with you alway,'' \&c.

St. Paul set Timothy over the Church at Ephesus,
and Titus over the Churches of Crete; i.e. as Bishops; therefore it is \emph{safer}
to have Bishops now, it is more likely to be pleasing to Him who has
loved us, and bids us in turn love Him with the heart, not with formal
service.

Our LORD
committed the Administration of the LORD'S SUPPER
\emph{to His Apostles}: ``Do THIS in remembrance of Me;''—therefore the Church has ever
continued it in the hands of their Successors, and the delegates of
these.

From CHRIST'S
words, ``Suffer the little children,'' \&c. and from His blessing them,
we infer His desire that children should be brought near to Him in
baptism;—as we do also from St. Paul's conduct on several occasions.
Acts xvi. 15, 33. 1 Cor. i. 16.

So also we continue the practice of Confirmation,
from a desire to keep as near the Apostles' rule as possible.

Again, what little is there of express command in
the New Testament for our meeting together in \emph{public} worship, in
large congregations! Yet we see what the custom of the Apostolic Church
was from the book of Acts, 1 Cor. \&c. and we follow it.

In like manner, the words in Genesis ii. and the
practice of the Apostles in the Acts, are quite warrant enough for the
Sanctification of the LORD'S
Day, even though the fourth Commandment be not binding on us.

For the same reason we continue the Patriarchal and
Jewish rule of paying tithe to the Church. Some portion of our goods is
evidently due to GOD;—and
the ancient Divine Command is a \emph{direction} to us, which the
law of the land has made obligatory, in a case where reason and
conscience have no means of determining.

These may be taken as illustrations of a general
principle. And at this day it is most needful to keep it in view, since
a cold spirit has crept into the Church of demanding rigid demonstration
for every religious practice and observance. It is the fashion now to
speak of those who maintain the ancient rules of the ecclesiastical
system, not as zealous servants of CHRIST,
not as wise and practical expounders of His will, but as \emph{inconclusive
reasoners}, and \emph{fanciful theorists}, merely because, instead of
standing still and arguing, they have a heart to obey. Are there not
numbers in this day, who think themselves enlightened believers, yet who
are but acting the part of the husbandman's son in the Gospel, who said,
``I go, Sir,''—AND WENT NOT.

CHURCH REFORM

SURELY,
before the blessing of a Millenium were vouchsafed to us, if it be to
come, the whole Christian world has much to confess in its several
branches. Rome has to confess her Papal corruptions, and her cruelty
towards those who refuse to accept them. The Christian communities of
Holland, Scotland, and other countries, their neglect of the Apostolical
Order of Ministers. The Greek Church has to confess its saint-worship,
its formal fasts, and its want of zeal. The Churches of Asia their
heresy. All parts of Christendom have much to confess and reform. We
have our sins as well as the rest. Oh that \emph{we} would take the lead
in the renovation of the Church Catholic on Scripture principles!

Our greatest sin perhaps is the disuse of a ``godly
discipline.'' Let the reader consider,

1. The command.

``Put away
from yourselves the wicked person.'' ``A man that is a heretic, after the
first and second admonition, reject.'' ``Mark them which cause divisions
and offences, … and avoid them.''

2. The example, viz. in the Primitive Church.

``The
Persons or Objects of Ecclesiastical Censure were all such delinquents,
as fell into great and scandalous crimes after baptism, whether men or
women, priests or people, rich or poor, princes or subjects.'' Bing.
Antiq. xvi. 3.

3. The warning.

``Whosoever
... shall break one of these least commandments, and shall teach men so,
he shall be called the least in the kingdom of heaven.''

[NEW EDITION.]

———————————————————————

These Tracts are continued in
Numbers, and sold at the price of 2d. for each sheet, or 7s. for 50
copies.

LONDON: PRINTED FOR
J. G. F. \& J. RIVINGTON,

ST. PAUL'S CHURCH YARD, AND WATERLOO PLACE.

1840.

\xchapter{Heads of a Weekday Lecture, Delivered to a Country Congregation in —shire}

Tract No. 10

BEFORE
we meet again, we shall have celebrated the least of St. Simon and St.
Jude, the Apostles. You will be at your daily work, and will not have
the opportunity to attend the Service in Church. For that reason, it may
be as well, you should lay up some good thoughts against that day; and
such, by GOD'S
blessing, I will now attempt to give you.

As you well know, there were twelve Apostles; St.
Simon and St. Jude were two of them. They preached the Gospel of CHRIST; and they were like CHRIST,
as far as sinful man may be accounted like the Blessed SON of GOD. They were like CHRIST in their deeds and in their
sufferings. The Gospel for the festival [John xv. 17.] shows us this.
They were like CHRIST
in their \emph{works}, because CHRIST was a witness of the FATHER, and they were witnesses of CHRIST.
CHRIST
came in the name of GOD
the FATHER
ALMIGHTY;
He ``came and spoke,'' and ``did works which none other man did.'' In like
manner, the Apostles were sent to bear witness of CHRIST, to declare His power, His great mercy, His sufferings on
the cross for the sins of all men, His willingness to save all who come
to Him.

But again, they were like CHRIST in their \emph{sufferings}. ``If the
world hate you,'' He says to them, ``you know that it hated Me, before it
hated you. If ye were of the world, the world would love his own; but
because ye are not of the world, but I have chosen you out of the
world, therefore the world hateth you. Remember the word that I said
unto you, The servant is not greater than his lord. If they have
persecuted Me, they will also persecute you; if they have kept My
saying, they will keep yours also.''

Thus, they were like CHRIST in \emph{office}. I do not speak of
their holiness, their faith, and all their other high excellencies,
which GOD
the HOLY
GHOST
gave them. I speak now, not of their personal graces, but of their \emph{office},
of preaching, of witnessing CHRIST,
of suffering for being His servants. Men ought to have listened to them,
and honoured them; some did: but the many, the world, did not,—they \emph{hated}
them; they hated them, for their office-sake; not because they were
Paul, and Peter, and Simon, and Jude, but because they bore witness to
the SON
of GOD,
and were chosen to be His Ministers.

Here is a useful lesson for us at this day. The
Apostles indeed are dead; yet it is quite as possible for men still to
hate their preaching and to persecute them, as when they were alive. For
in one sense they are still alive; I mean, they did not leave the world
without appointing persons to take their place: and these persons
represent them, and may be considered with reference to us, as if they
were the Apostles. When a man dies, his son takes his property and
represents him; that is, in a manner he still lives in the person of his
son. Well, this explains how the Apostles may be said to be still among
us; they did not indeed leave their sons to succeed them as Apostles,
but they left \emph{spiritual} sons; they did not leave this life,
without first solemnly laying their hands on the heads of certain of
their brethren, and these took their place, and represented them after
their death.

But it may be asked, are these spiritual sons of
the Apostles still alive? no:—all this took place many hundred years
ago. These sons and heirs of the Apostles died long since. But then they
in turn did not leave the world without committing their sacred office
to a fresh set of Ministers, and they in turn to another, and so on even
to this day. Thus the Apostles had, first, spiritual sons; then
spiritual grandsons; then great grandsons; and so on, from one age to
another, down to the present time.

Again, it may be asked, \emph{who} are at this time
the successors and spiritual descendants of the Apostles? I shall
surprise some people by the answer I shall give, though it is very
clear, and there is no doubt about it; THE
BISHOPS. They stand in the place of the Apostles, as far as the
office of ruling is concerned \footnote{As far as the office of \emph{ruling}, not as far
as the office of \emph{teaching} is concerned. The Apostles were both \emph{inspired
teachers} (Acts ii. 3, 4), and Bishops (John xx. 21-23). Their
successors are Bishops only, not inspired teachers; and rule according
to the Apostles' teaching,—not absolutely, as the Apostles may be said
to have done.};
and, whatever we ought to do, had we lived when the Apostles were alive,
the same ought we to do for the Bishops. He that despiseth them,
despiseth the Apostles. It is our duty to reverence them for their
office-sake; they are the shepherds of CHRIST'S flock. If we knew them well, we should
love them for the many excellent graces they possess, for their piety,
loving-kindness, and other virtues. But we do not know them; yet still,
for all this, we may honour them as the Ministers of CHRIST, without going so far as to consider
their \emph{private} worth; and we may keep to their ``fellowship'', [Acts
ii. 42.] as we should to that of the Apostles. I say, we may all thus
honour them even without knowing them in private, because of their high
office; for they have the marks of CHRIST'S
presence upon them, in that they \emph{witness} for CHRIST, and \emph{suffer} for Him, as
the Apostles did. I will explain to you how this is.

There is a temptation which comes on many men to
honour no one, except such as they themselves know, such as have done a
favour or kindness to them personally. Thus sometimes people speak
against those who are put over them in this world's matters, as the
King. They say, ``What is the King to me? he never did me any good.'' Now,
I answer, whether he did or not, is nothing to the purpose. We are
bound, \emph{for} CHRIST'S sake, to honour him, \emph{because} he
is King, though he lives far from us; and this all well-disposed,
right-minded people do. And so, in just the same way, though for
much higher reasons, we must honour the Bishop, because he \emph{is} the
Bishop;—for his \emph{office}-sake;—because he is CHRIST'S Minister, stands in the place of the
Apostles, is the Shepherd of our souls on earth, while CHRIST is away. This is FAITH,
to look at things not as seen, but as unseen; to be as sure that the
Bishop is CHRIST'S
appointed Representative, as if we actually saw him work miracles as St.
Peter and St. Paul did, as you may read in the book of the Acts of the
Apostles.

But you will say, how do we know this, since we do
not see it? I repeat, the Bishops are Apostles to us, from their \emph{witnessing}
CHRIST,
and \emph{suffering} for Him.

1. They witness our L\textsc{ord} in
their very \emph{name}, for He is the true Bishop of our souls, as St.
Peter says, and they are Bishops. They witness CHRIST
in their \emph{station};—there is but one LORD to save us, and there is but one Bishop in
each place. The meetingers have no head, they are all of them mixed
together in a confused way; but we of CHRIST'S Holy Church (blessed, be GOD!) have one Bishop over us, and our Bishop
is the Bishop of ——. Many of you have seen him lately, when he
confirmed in our Church. That very \emph{confirmation} is another
ordinance, in which the Bishop witnesses CHRIST.
Our LORD
and SAVIOUR
confirms us with the SPIRIT
in all goodness; the Bishop is His figure and likeness, when he lays his
hands on the heads of children. Then CHRIST, (as we trust) comes to them, to confirm in them the grace
of Baptism. Moreover, the Bishop \emph{rules} the whole Church here
below, as CHRIST,
the true and eternal Sovereign, rules it above; and here again the
Bishop is a figure or witness of our LORD.
And further, it is the Bishop who is commissioned to make us Clergymen GOD'S
Ministers. He is CHRIST'S
instrument; and he visibly chooses those whom CHRIST vouchsafes to choose invisibly, to serve
in the Word and Sacraments of the Church. And thus, in one sense, it is
from the Bishop that the \emph{news of redemption and the means of grace}
have come to all men; this again is a witnessing CHRIST. I, who speak to you concerning CHRIST,
was ordained to do so by the Bishop; lie speaks in me,—as CHRIST wrought in him, and as GOD
sent CHRIST.
Thus the whole plan of salvation hangs together.—CHRIST the True Mediator above; His servant, the Bishop, His
earthly likeness; mankind, the subjects of His teaching; GOD
the Author of Salvation.

2. But I must now mention the more painful part of
the subject, i.e. the \emph{sufferings} of the Bishops, which is the
second mark of their being our living Apostles. I may say, Bishops have
undergone this trial in every age. As the first Apostles were hated and
opposed by the world, so have they ever been. I do not say they have
been always opposed in the same way. In these latter times, they have
experienced the lesser sufferings of bearing slander, reproach, threats,
vexations, and thwartings in their efforts to do good. Time was, when
they were even persecuted, cruelly slain by fire and sword. That time,
(though GOD avert it!) may come again. But, whether or not Satan is
permitted so openly to rage, certainly some kinds of persecution are to
be expected in our day; nay, such have begun. It is not so very long
since the great men of the earth told them to \emph{prepare for persecution};
it is not so very long since the mad people answered the summons, and
furiously attacked them, and seemed bent on destroying them, in all
parts of the country.

Yes! the day may come, even in this generation,
when the Representatives of CHRIST
are spoiled of their sacred possessions, and degraded from their civil
dignities. The day may come, when each of us inferior Ministers—when I
myself, whom you know—may have to give up our Churches, and be among
you, in no better temporal circumstances than yourselves; with no larger
dwelling, no finer clothing, no other fare, with nothing different
beyond those gifts, which I trust we received from the All-gracious GOD when we were made Ministers; and those
again, which have been vouchsafed to us before and after that time, for
the due fulfilment of our Ministry. Then you will look at us, not as
gentlemen, as now; not as your superiors in worldly station; but still,
nay, more strikingly so than now, still as messengers from Him, who
seeth and worketh in secret, and who judgeth not by outward appearance.
Then you will honour us, with a purer honour than many men do now,
namely, as those (if I may say so) who are intrusted with the keys of
heaven and hell, as the heralds of mercy, as the denouncers of woe to
wicked men, as intrusted with the awful and mysterious privilege of
dispensing CHRIST'S
Body and Blood, as far greater than the most powerful and the wealthiest
of men in our unseen strength and our heavenly riches. This may all come
in our day; \emph{we} must do our duty; go straight forward, looking
neither to the right hand nor the left, ``in patience possessing our
souls,'' watching and praying, and so preparing for the evil day. And
after all, if GOD'S loving kindness spares both us and you
the trial, still it will have been useful to have steadily thought about
it beforehand, and to have prepared our hearts to meet it.

OXFORD,

\emph{Nov}. 4, 1833.

[FIFTH EDITION.]

———————————————————————

These Tracts are continued in
Numbers, and sold at the price of 2d. for each sheet, or 7s. for 50
copies.

LONDON: PRINTED FOR
J. G. F. \& J. RIVINGTON,

ST. PAUL'S CHURCH YARD AND WATERLOO PLACE.

1840.

\xchapter{The Visible Church (In Letters to a Friend.)}

Tract No. 11 (Ad Scholas)

\xsection{Letter I.}

You wish to have my opinion on the doctrine of ``the
Holy Catholic Church,'' as contained in Scripture, and taught in the
Creed. So I send you the following lines, which perhaps may serve
through God's blessing, to assist you in your search after the truth in
this matter, even though they do no more; indeed no remarks, however
just, can be much more than an assistance to you. You must search for
yourself, and GOD must teach you.

I think I partly enter into your present
perplexity. You argue, that true \emph{doctrine} is the important matter
for which we must contend, and a \emph{right state of the affections} is
the test of vital religion in the heart: and you ask, ``Why may I not be
satisfied if my Creed is correct, and my affections spiritual? Have I
not in that case enough to evidence a renewed mind, and to constitute a
basis of union with others like minded? The love of CHRIST
is surely the one and only requisite for Christian communion here, and
the joys of heaven hereafter.'' Again you say, that —— and —— are
constant in their prayers for the teaching of the HOLY SPIRIT;
so that if it be true, that every one who asketh receiveth, surely they
must receive, and are in a safe state.

Believe me, I do not think lightly of these
arguments. They are very subtle ones; powerfully influencing the
imagination, and difficult to answer. Still I believe them to be mere
fallacies. Let me try them in a parallel case. You know the preacher at
——, and have heard of his flagrantly immoral life; yet it is
notorious that he can and does speak in a moving way of the love of CHRIST,
\&c. It is very shocking to witness such a case, which (we will hope)
is rare; but it has its use. Do you not think him in peril, in spite of
his impressive and persuasive language? Why?—You will say, his life is
bad. True; it seems then that more is requisite for salvation than an
orthodox creed, and keen sensibility; viz. consistent conduct.—Very
well then, we have come to an additional test of truth faith, obedience
to GOD'S word, and plainly a scriptural test,
according to St. John's canon, ``He who \emph{doeth} righteousness is
righteous.'' Do not you see then your argument is already proved to be
unsound? It seems that true doctrine and warm feelings are not enough.
How am I to know what \emph{is} enough? you ask. I reply, \emph{by
searching Scripture}. It was your original fault that, instead of
inquiring what GOD has told you is necessary for being
a true Christian, you chose out of your own head to \emph{argue} on the
subject;—e.g. ``I can never believe that to be such and such is not
enough for salvation,'' \&c. Now this is \emph{worldly wisdom}.

Let us join issue then on this plain ground,
whether or not the doctrine of ``the Church,'' and the duty of obeying it,
be laid down \emph{in Scripture}. If so, it is no matter as regards our
practice, whether the doctrine is primary or secondary, whether the duty
is much or little insisted on. A Christian mind will aim at obeying the \emph{whole}
counsel and will of GOD;
on the other hand, to those who are tempted arbitrarily to classify and
select their duties, it is written, ``Whosoever shall break one of these
least commandments, and shall teach men so, he shall be called the least
in the kingdom of heaven.''

And here first, that you may clearly understand the
ground I am taking, pray observe that I am not attempting to controvert
any one of those high evangelical points, on which perhaps we do not
altogether agree with each other. Perhaps you attribute less efficacy to
the Sacrament of Baptism than I do; bring out into greater system and
prominence the history of an individual's warfare with his spiritual
enemies; fix more precisely and abruptly the date of his actual
conversion from darkness to light; and consider that Divine Grace acts
more arbitrarily against the corrupt human will, than I think is
revealed in Scripture. Still, in spite of this difference of opinion, I
see no reason why you should not accept heartily the Scripture doctrine
of ``the Church.'' And this is the point I wish to press, not asking you
at present to abandon your own opinions, but to \emph{add to them} a
practical belief in a tenet which the Creed teaches and Scripture has
consecrated. And this surely is quite possible. The excellent Mr.
——, of ——, who has lately left ——, was both a Calvinist, and
a strenuous High-Churchman.

You are in the practice of distinguishing between
the Visible and Invisible Church. Of course I have no wish to maintain,
that those who shall be saved hereafter are exactly the same company
that are under the means of grace here; still I must insist on it, that
Scripture makes the existence of a Visible Church a condition of the
existence of the Invisible. I mean, the \emph{Sacraments} are evidently
in the hands of the Church Visible; and these, we know, are generally
necessary to salvation, as the Catechism says. Thus it is an undeniable
fact, as true as that souls will be saved, that a Visible Church must
exist as a means towards that end. The Sacraments are in the hands of
the Clergy; this few will deny, or that their efficacy is independent of
the personal character of the administrator. What then shall be thought
of any attempts to weaken or exterminate that Community, or that
Ministry, which is an appointed condition of the salvation of the elect?
But every one, who makes or encourages a schism, \emph{must} weaken it.
Thus it is plain, schism must be wrong in itself, even if Scripture did
not in express terms forbid it, as it does.

But further than this; it is plain this Visible
Church is a standing body. Every one who is baptized, is baptized \emph{into}
an existing community. Our Service expresses this when it speaks of
baptized infants being \emph{incorporated} into GOD'S
holy Church. Thus the Visible Church is not a voluntary association of
the day, but a continuation of one which existed in the age before us,
and then again in the age before that; and so back till we come to the
age of the Apostles. In the same sense, in which Corporations of the
State's creating, are perpetual, is this which CHRIST
has founded. This is a matter of fact hitherto: and it necessarily will
be so always, for is not the notion absurd of an unbaptized person
baptizing others? which is the only way in which the Christian community
can have a new beginning.

Moreover Scripture directly insists upon the
doctrine of the Visible Church, as being of importance. E.g. St. Paul
say;—``There is \emph{one body}, and one SPIRIT,
even as ye are called in one hope of your calling; one LORD, one faith, one baptism, one GOD
and FATHER
of all.'' Ephes. iv. 5, 6. Thus, as far as the Apostle's words go, it is
as false and unchristian, (I do not mean in degree of guilt, but in its
intrinsic sinfulness), to make more bodies than one, as to have many
Lords, many Gods, many Creeds. Now, I wish to know, how it is possible
for any one to fall into this sin, if Dissenters are clear of it? What
is the sin, if separation from the Existing Church is not it?

I have shown that there is a divinely instituted
Visible Church, and that it has been one and the same by successive
incorporation of members from the beginning. Now I observe further, that
the word Church, as used in Scripture, ordinarily means this
actually existing visible body. The possible exception to this rule, out
of about 100 places in the New Testament, where the word occurs, are
four passages in the Epistle to the Ephesians; two in the Colossians;
and one in the Hebrews. (Eph. i. 22; iii. 10, 21; v. 23-32. Col. i. 18,
24. Heb. xii. 23.)—And in some of these exceptions the sense is at
most but doubtful. Further, our SAVIOUR uses the word twice, and in both times of the Visible
Church. They are remarkable passages, and may here be introduced, in
continuation of my argument.

Matt. xvi. 18. ``Upon this rock I will build My
Church, and the gates of hell shall not prevail against it.'' Now I am
certain, any unprejudiced mind, who knew nothing of controversy,
considering the Greek word [\emph{ekkl\={e}sia}] means simply an \emph{assembly},
would have no doubt at all that it meant in this passage a visible body.
What right have we to disturb the plain sense? why do we impose a
meaning, arising from some system of our own? And this view is
altogether confirmed by the other occasion of our LORD'S
using it, where it can \emph{only} denote the Visible Church. Matt.
xviii. 17. ``If he (thy brother) shall neglect to hear them, tell it unto
the Church; but if he neglect to hear the Church, let him be unto thee
as a heathen man and a publican.''

Observe then what we gain by these two
passages;—the grant of \emph{power} to the Church; and the promise of \emph{permanence}.
Now look at the fact. The body then begun has continued; and has always
claimed and exercised the power of a corporation or society. Consider
merely the article in the Creed, ``The Holy Catholic Church;'' which
embodies this notion. Do not Scripture and History illustrate each
other?

I end this first draught of my argument, with the
text in 1 Tim. iii. 15, in which St. Paul calls the Church ``the pillar
and ground of the truth,''—which can refer to nothing but a Visible
Body; else martyrs may be invisible, and preachers, and teachers, and
the whole order of the Ministry.

My paper is exhausted. If you allow me, I will send
you soon a second Letter; meanwhile I sum up what I have been proving
from Scripture thus; that ALMIGHTY
GOD
might have left Christianity as a sort of sacred literature, as
contained in the Bible, which each person was to take and use by
himself; just as we read the works of any human philosopher or
historian, from which we gain practical instruction, but the
knowledge of which does not bind us to be Newtonians, or Aristotelians,
\&c., or to go out of our line of life in consequence of it. This, I
say, He might have done; but, in matter of fact, He has ordained
otherwise. He has actually set up a Society, which exists even this day
all over the world, and which (as a general rule) Christians are bound
to join; so that to believe in CHRIST
is not a mere opinion or a secret conviction, but a social or even a
political principle, forcing one into what is often stigmatized as party
strife, and quite inconsistent with the supercilious mood of those
professed Christians of the day, who stand aloof, and designate their
indifference as philosophy.

\xsection{Letter II.}

I AM
sometimes struck with the inconsistency of those, who do not allow us to
express the gratitude due to the Church, while they do not hesitate to
declare their obligation to individuals who have benefited them. To a
vow that they owe their views of religion and their present hopes of
salvation to this or that distinguished preacher, appears to them as
harmless, as it may be in itself true and becoming; but if a person
ascribes his faith and knowledge to the Church, he is thought to forget
his peculiar and unspeakable debt to that SAVIOUR
who died for him. Surely, if our LORD
makes man His instrument of good to man, and if it is possible to be
grateful to man without forgetting the Source of all grace and power,
there is nothing wonderful in His having appointed a company of men as
the especial medium of His instruction and spiritual gifts, and in
consequence, of His having laid upon us the duty of gratitude to it. Now
this is all I wish to maintain, what is most clearly (as I think)
revealed in Scripture, that the blessings of redemption come to us
through the Visible Church; so that, as we betake ourselves to a
Dispensary for medicine, without attributing praise or intrinsic worth
to the building or the immediate managers of its stores, in something of
the like manner we are to come to that One Society, to which CHRIST
has entrusted the office of stewardship in the distribution of gifts, of
which He alone is the Author and real Dispenser.

In the letter I sent you the other day, I made some
general remarks on this doctrine; now let me continue the subject.

First, the Sacraments, which are the ordinary means
of grace, are clearly in possession of the Church. Baptism is an
incorporation into a body; and invests with spiritual blessings, because
it is the introduction into a body so invested. In 1 Cor. xii. we are
taught first, the SPIRIT'S
indwelling in the Visible Church or body; I do not say \emph{in every
member of it}, but generally \emph{in} it;—next, we are told that
the SPIRIT baptizes individuals \emph{into}
that body. Again, the LORD'S
Supper carries evidence of its social nature even in its name; it is not
a solitary individual act, it is a joint communion. Surely nothing is
more alien to Christianity than the spirit of Independence; the peculiar
Christian blessing, i.e. the presence of CHRIST, is upon \emph{two or three} gathered together, not on mere
individuals.

But this is not all. The Sacraments are committed,
not into the hand of the Church Visible assembled together, (though even
this would be no unimportant doctrine practically) but into certain
definite persons, who are selected from their brethren for that trust. I
will not here determine who these are in each successive age, but will
only point out how far this principle itself will carry us. The doctrine
is implied in the original institution of the LORD'S
Supper, where CHRIST
says to His Apostles, ``Do this.'' Further, take that remarkable passage
in Matt. xxiv. 45-51. Luke xii. 42-46, ``Who then is that faithful and
wise Steward, whom his Lord shall make ruler over His household, to give
them their portion of meat in due season? Blessed is that servant, whom
his Lord, \emph{when He cometh}, shall find so doing!'' \&c. Now I do
not inquire \emph{who} in every age are the stewards spoken of, (though
in my own mind I cannot doubt the line of Bishops is that Ministry, and
consider the concluding verses fearfully prophetic of the Papal misuse
of the gift;—by the bye, at least it shows this, that bad men may
nevertheless be the channels of grace to GOD'S
``household,'') I do not ask who are the stewards, but surely the words, \emph{when
He cometh}, imply that they are to continue till the end of the
world. This reference is abundantly confirmed by our LORD'S parting words to the eleven; in which, after giving them
the baptismal commission, He adds, ``Lo! I am with you \emph{always},
even unto the end of the world.'' If then He was with the Apostles in a
way in which He was not present with teachers who were strangers to
their ``fellowship,'' (Acts ii. 42.) which all will admit, so, in like
manner, it cannot be a matter of indifference in \emph{any} age, what
teachers and fellowship a Christian selects; there must be those
with whom CHRIST is present, who are His ``Stewards,'' and
whom it is our duty to obey.

As I have mentioned the question of faithfulness
and unfaithfulness in Ministers, I may refer to the passage in 1 Cor.
iv. where St. Paul, after speaking of himself and others as ``\emph{Stewards}
of the mysteries of GOD,''
and noticing that ``it is required of Stewards, that a man be found
faithful,'' adds, ``With me it is a very small thing that I should be
judged of you or of man's judgment … therefore \emph{judge nothing before
the time}.''

To proceed, consider the following passage: ``Obey
them that have rule over you, and submit yourselves.'' Heb. xiii. 17.
Again, I do not ask \emph{who} these are; but whether this is not a
duty, however it is to be fulfilled, which multitudes \emph{in no sense}
fulfil. Consider the number of people, professing and doubtless in a
manner really actuated by Christian principle, who yet wander about from
church to church or from church to meeting, as sheep without a shepherd,
or who choose a preacher merely because he pleases their taste, and
whose first movement towards any clergyman they meet, is to examine and
criticize his doctrine: what conceivable meaning do they put upon these
words of the Apostle? Does any one \emph{rule over} them? do they in any
way \emph{submit themselves}? Can these persons excuse their conduct,
except on the deplorably profane plea, (which yet I believe is in their
hearts at the bottom of their disobedience,) that it matters little to
keep CHRIST'S ``least commandments,'' so that
we embrace the peculiar doctrines of His Gospel?

Some time ago I drew up a sketch of the Scripture
proof of the doctrine of the Visible Church; which with your leave I
will here transcribe. You will observe, I am not arguing for this or
that form of Polity, or for the Apostolical Succession, but simply the
duties of order, union, ecclesiastical gifts, and ecclesiastical
obedience; I limit myself to these points, as being persuaded that, when
they are granted, the others will eventually follow.

I. That there was a Visible Church in the Apostles'
day.

\begin{quote}

1. General texts. Matt. xvi. 18; xviii. 17. 1
Tim. iii. 15. Acts passim, \&c.

2. Organization of the Church.

\begin{quote}

(1.) Diversity of ranks. 1 Cor. xii. Eph. iv.
4-12. Rom. xii. 4-8. 1 Pet. iv. 10, 11.

(2.) Governors. Matt. xxviii. 19. Mark xvi. 15,
16. John xx. 22, 23. Luke xxii. 19, 20. Gal. ii. 9, \&c.

(3.) Gifts. Luke xii. 42, 43. John xx. 22, 23.
Matt. xviii. 18.

(4.) Order. Acts viii. 5, 6, 12, 14, 15, 17; xi.
22, 23; xi. 2, 4; ix. 27; xv. 2, 4, 6, 25; xvi. 4; xviii. 22; xxi.
17-19. conf. Gal. i. 1, 12. 1 Cor. xiv. 40. 1 Thess. v. 14.

(5.) Ordination. Acts vi. 6. 1 Tim. iv. 14; v.
22. 2 Tim. i. 6. Tit. i. 5. Acts xiii. 3. conf. Gal. i. 1, 12.

(6.) Ecclesiastical obedience. 1 Thess. v. 12,
13. Heb. xiii. 17. 1 Tim. v. 17.

(7.) Rules and discipline. Matt. xxviii. 19.
Matt. xviii. 17. 1 Cor. v. 4-7. Gal. v. 12, \&c. 1 Cor. xvi. 1,
2. 1 Cor. xi. 2, 16, \&c.

(8.) Unity. Rom. xvi. 17. 1 Cor. i. 10; iii. 3;
xiv. 26. Col. ii. 5. 1 Thess. v. 14. 2 Thess. iii. 6.

\end{quote}

\end{quote}

II. That the Visible Church, thus instituted by the
Apostles, was intended to continue.

\begin{quote}

1. Why should it not? The \emph{onus probandi}
lies with those who deny this position. If the doctrines and precepts
already cited are obsolete at this day, why should not the following
texts? e.g. 1 Pet. ii. 13, or e.g. Matt. vii. 14. John iii. 3.

2. Is it likely so elaborate a system should be
framed, yet with no purpose of its continuing?

3. The objects to be obtained by it are as
necessary now as then. (1.) Preservation of the faith. (2.) Purity of
doctrine. (3.) Edification of Christians. (4.) Unity of operation. Vid.
Epistles. to Tim. \& Tit. passim.

4. If system were necessary in a time of
miracles, much more is it now.

5. 2 Tim. ii. 2. Matt. xxviii. 20, \&c.

\end{quote}

Take these remarks, as they are meant, as mere
suggestions for your private consideration.

[SIXTH EDITION.]

———————————————————————

These Tracts are continued in
Numbers, and sold at the price of 2d. for each sheet, or 7s. for 50
copies.

LONDON: PRINTED FOR
J. G. F. AND J. RIVINGTON,

ST. PAUL'S CHURCH YARD, AND WATERLOO PLACE.

1840.

\xchapter{On the Apostolical Succession in the English Church}

Tract No. 15

WHEN
Churchmen in England maintain the Apostolical Commission of their
Ministers, they are sometimes met with the objection, that they cannot
prove it without tracing their orders back to the Church of Rome; a
position, indeed, which in a certain sense is true. And hence it is
argued, that they are reduced to the dilemma, either of acknowledging
they had no right to separate from the Pope, or, on the other hand, of
giving up the Ministerial Succession altogether, and resting the claims
of their pastors on some other ground; in other words, that they are
inconsistent in reprobating Popery, while they draw a line between their
Ministers and those of Dissenting Communions.

It is intended, in the pages that follow, to reply
to this supposed difficulty; but first, a few words shall be said, by
way of preface, on the doctrine itself, which we Churchmen advocate.

The Christian Church is a body consisting of Clergy
and Laity; this is generally agreed upon, and may here be assumed. Now,
what we say is, that these two classes are distinguished from each
other, and united to each other, by the commandment of GOD Himself; that the Clergy have a commission from GOD
ALMIGHTY
through regular succession from the Apostles, to preach the Gospel,
administer the Sacraments, and guide the Church; and, again, that in
consequence the people are bound to hear them with attention, receive
the Sacrament from their hands, and pay them all dutiful obedience. I
shall not prove this at length, for it has been done by others, and
indeed the common sense and understanding of men, if left to themselves,
would be quite sufficient in this case. I do but lay before the reader
the following considerations.

1. We hold, with the Church in all ages, that, when
our LORD,
after His resurrection, breathed on His Apostles, and said, ``Receive ye
the HOLY
GHOST,—as
My FATHER
hath sent Me, so send I you;'' He gave them the power of sending \emph{others}
with a divine commission, who in like manner should have the power of
sending others, and so on even unto the end; and that our LORD promised His continual assistance to these
successors of the Apostles in this and all other respects, when He said,
``Lo, I am with you,'' (that is, with you, and those who shall represent
and succeed you,) ``alway, even unto the end of the world.''

And, if it is plain that the Apostles left
successors after them, it is equally plain that the Bishops are these
Successors. For it is only the Bishops who have ever been called by the
title of Successors; and there has been actually a perpetual succession
of these Bishops in the Church, who alone were always esteemed to have
the power of sending other Ministers to preach and administer the
Sacraments. So that the proof of the doctrine seems to lie in a very
small space.

2. But, perhaps it may be as well to look at it in
another point of view. I suppose no man of common sense thinks himself
entitled to set about teaching religion, administering Baptism, and the
LORD'S
Supper, and taking care of the souls of other people, unless he has \emph{in
some way} been called to undertake the office. Now, as religion is a
business between every man's own conscience and GOD ALMIGHTY,
no one can have any right to interfere in the religious concerns of
another with the authority of a teacher, unless he is able to show, that
it is GOD
that has in some way called and sent him to do so. It is true, that men
may as \emph{friends} encourage and instruct each other with consent of
both parties; but this is something very different from the office of a
Minister of religion, who is entitled and called to ``exhort, rebuke, and''
``rule,'' ``with all authority,'' as well as love and humility.

You may observe that our LORD Himself did not teach the Gospel, without
proving most plainly that His FATHER had sent Him. He and His Apostles proved their divine
commission by miracles. As miracles, however, have long ago come to an
end, there must be some \emph{other} way for a man to prove his right to
be a Minister of religion. And what other way can there possibly be,
except a regular call and ordination by those who have succeeded to the
Apostles?

3. Further, you will observe, that all sects think
it necessary that their Ministers should be ordained by other Ministers.
Now, if this be the case, then the validity of ordination, even with
\emph{them}, rests on a \emph{succession}; and is it not plain that they
ought to trace that succession to the Apostles? Else, why are they
ordained at all? And, any how, if \emph{their} Ministers have a
commission, who derive it from private men, much more do the Ministers
of our Church, who actually do derive it from the Apostles. Surely those
who dissent from the Church have \emph{invented} an ordinance, as they
themselves must allow; whereas Churchmen, whether rightly or wrongly,
still maintain \emph{their} succession not to be an invention, but to be
GOD'S
ordinance. If Dissenters say, that \emph{order} requires there should be
some such \emph{succession}, this is true, indeed, but still it is only
a testimony to the mercy of CHRIST,
in having, as Churchmen maintain, \emph{given us} such a succession. And
this is \emph{all} it shows; it does nothing for \emph{them}; for, their
succession, not professing to come from GOD, has no power to restrain any fanatic from setting up to
preach of his own will, and a people with itching ears choosing for
themselves a teacher. It does but witness to a need, without supplying
it.

4. I have now given some slight suggestions by way
of evidence for the doctrine of the Apostolical Succession, from
Scripture, the nature of the case, and tile conduct of Dissenters. Let
me add a word on the usage of the Primitive Church. We know that the
succession of Bishops, and ordination from them, was the invariable
doctrine and rule of the early Christians. Is it not utterly
inconceivable, that this rule should have prevailed from the first age,
everywhere, and without exception, had it not been given them by the
Apostles?

But here we are met by the objection, on which I
propose to make a few remarks, that, though it is true there was a
continual Succession of pastors and teachers in the early Church who had
a divine commission, yet that no Protestants can have it; that we gave
it up, when our communion ceased with Rome, in which Church it still
remains; or, at least, that no Protestant can plead it without
condemning the Reformation itself, for that our own predecessors then
revolted and separated from those spiritual pastors, who, according to
our principles, then had the commission of JESUS
CHRIST.

Our reply to this is a flat denial of the alleged
facts on which it rests. The English Church did not revolt from those
who in that day had authority by succession from the Apostles. On
the contrary, it is certain that the Bishops and Clergy in England and
Ireland remained the same as before the separation, and that it was
these, with the aid of the civil power, who delivered the Church of
those kingdoms from the yoke of Papal tyranny and usurpation, while at
the same time they gradually removed from the minds of the people
various superstitious opinions and practices which had grown up during
the middle ages, and which, though never formally received by the
judgment of the whole Church, were yet very prevalent. I do not say the
case might never arise, when it might become the duty of private
individuals to take upon themselves the office of protesting against and
abjuring the heresies of a corrupt Church. But such an extreme case it
is unpleasant and unhealthy to contemplate. All I say here is, that this
was not the state of things at the time of the Reformation. The Church
then by its proper rulers and officers reformed itself. There was no new
Church founded among us, but the rights and the true doctrines of the
Ancient existing Church were asserted and established.

In proof of this we need only look to the history
of the times. In the year 1534, the Bishops and Clergy of England
assembled in their respective convocations of Canterbury and York, and
signed a declaration that the Pope or Bishop of Rome had no more
jurisdiction in this country by the word of GOD,
than any other foreign Bishop; and they also agreed to those acts of the
civil government, which put an end to it among us \footnote{Vide Collier, Eccl. Hist. v. ii. p. 94.}.

The people of England, then, in casting off the
Pope, but obeyed and concurred in the acts of their own spiritual
Superiors, and committed no schism. Queen Mary, it is true, drove out
after many years the orthodox Bishops, and reduced our Church again
under the Bishop of Rome, but this submission was only exacted by force,
and in itself null and void; and, moreover, in matter of fact it lasted
but a little while, for on the succession of Queen Elizabeth, the true
Successors of the Apostles in the English Church were reinstated in
their ancient rights. So, I repeat, there was no revolt, in any part of
these transactions, against those who had a commission from God; for it
was the Bishops and Clergy themselves, who maintained the just rights of
their Church.

But, it seems, the Pope has ever said, that our
Bishops were bound by the laws of GOD
and the Church to obey \emph{him}; that they were subject to him; and
that they had no right to separate from him, and were guilty in doing
so, and that accordingly they have involved the people of England in
their guilt; and, at all events, that \emph{they} cannot complain of
their flock disobeying and deserting them, when they have revolted from
the Pope. Let us consider this point.

Now that there is not a word in \emph{Scripture}
about our duty to obey the Pope, is quite clear. The Papists indeed say,
that he is the Successor of St. Peter; and that therefore he is Head of
all Bishops, because St. Peter bore rule over the other Apostles. But
though the Bishops of Rome were often called the Successors of St. Peter
in the early Church, yet every other Bishop had the \emph{same} title.
And though it be true, that St. Peter was the \emph{foremost} of the
Apostles, that does not prove he had any \emph{dominion} over them. The
eldest brother in a family has certain privileges and a precedence, but
he has no power over the younger branches of it. And so Rome has ever
had what is called the \emph{primacy} of the Christian Churches; but it
has not therefore any right to interfere in their internal
administration; not more of a right, than an elder brother has to meddle
with a younger brother's household.

And this is plainly the state of matters between us
and Rome, \emph{in the judgment of the Ancient Church} also, to which
the Papists are fond of appealing, and by which we are quite ready to
stand or fall. In early times, as is well known, all Christians thought
substantially alike, and formed one great body all over the world,
called the Church Catholic, or Universal. This great body, consisting of
a vast number of separate Churches, with each of them its own Bishop at
its head, was divided into a number of portions called Patriarchates;
these again into others called Provinces, and these were made up of the
separate Dioceses or Bishoprics. We have among ourselves an instance of
this last division in the Provinces of Canterbury and York, which
constitute the English Church, each of them consisting of a number of
distinct Bishoprics or Churches. The head of a Province was called
Archbishop, as in the case of Canterbury and York; the Bishops of those
two sees being, we know, not only Bishops with Dioceses of their own,
but having, over and above this, the place of precedence among the
Bishops in the same Province. In like manner, the Bishop at the head
of a Patriarchate was called the Patriarch, and had the place of honour
and certain privileges over all other Bishops within his own
Patriarchate. Now, in the early Christian Church, there were four or
five Patriarchates; e.g. one in the East, the Head of which was the
Bishop of Antioch; one in Egypt, the Head of which was the Bishop of
Alexandria; and, again, one in the West, the Head of which was the
Bishop of Rome. These Patriarchs, I say, were the Primates or Head
Bishops of their respective Patriarchates; and they had an order of
precedence among themselves, Rome being the first of them all. Thus the
Bishop of Rome, being the first of the Patriarchs in dignity, might be
called the honorary Primate of all Christendom.

However, as time went on, the Bishop of Rome, not
satisfied with the honours which were readily conceded to him, attempted
to gain \emph{power} over the whole Church. He seems to have been
allowed the privilege of \emph{arbitrating} in case of appeal from other
Patriarchates. If, \emph{e.g}. Alexandria and Antioch had a dispute, he
was a proper referee; or if the Bishops of those Churches were at any
time unjustly deprived of their sees, he was a fit person to interfere
and defend them. But, I say, he became ambitious, and attempted to \emph{lord
it} over GOD'S
heritage. He interfered in the internal management of other
Patriarchates; he appointed Bishops to sees, and Clergy to parishes
which were contained within them, and imposed on them various religious
and ecclesiastical usages illegally. And in doing so, surely he became a
remarkable contrast to the Holy Apostle, who, though inspired, and an
universal Bishop, yet suffered not himself to control the proceedings
even of the Churches he founded; saying to the Corinthians, ``not for
that we have dominion over your faith, but are helpers of your joy; for
by faith ye stand.'' 2 Cor. i. 24. This impressive declaration, which
seems to be intended almost as a prophetic warning against the times of
which we speak, was neglected by the Pope, who, among other tyrannical
proceedings, took upon him the control of the Churches in Britain, and
forbade us to reform our doctrine and usages, which he had no right at
all to do. He had no pretence for so doing, because we were altogether
independent of him; the English and Irish Churches, though in the West,
being exterior to his Patriarchate. Here again, however, some
explanation is necessary.

You must know, then, that from the first there were
portions of the Christian world, which were not included in any
Patriarchate, but were governed by themselves. Such were the Churches of
Cyprus, and such were the British Churches. This need not here be
proved; even Papists have before now confessed it. Now it so happened,
in the beginning of the 5th century, the Patriarch of Antioch, who was
in the neighbourhood of Cyprus, attempted against the Cyprian Churches
what the Pope has since attempted against us; viz. took measures to
reduce them under his dominion. And, as a sign of his authority over
them, he claimed to consecrate their Bishops. Upon which the Great
Council of the whole Christian world assembled at Ephesus, A.D. 451, made the following decree, which you
will find is a defence of England and Ireland against the Papacy, as
well as of Cyprus against Antioch.

``An innovation upon the Rule of the Church and the
Canons of the Holy Fathers, such as to affect the general liberties of
Christendom, has been reported to us by our venerable brother Rheginus,
and his fellow Bishops of Cyprus, Zeno, and Evagrius. Wherefore, since
public disorders call for extraordinary remedies, as being more
perilous, and whereas it is against ancient usage, that the Bishop of
Antioch should ordain in Cyprus, as has been proved to us in this
Council both in words and writing, by most orthodox men, We therefore
decree, that the Prelates of the Cyprian Churches shall be suffered
without let or hindrance to consecrate Bishops by themselves; and
moreover, that the \emph{same rule shall be observed also in other dioceses
and provinces every where, so that no Bishop shall interfere in another
province, which has not from the very first been under himself and his
predecessors}; and further, that if any one has so encroached and
tyrannized, he must relinquish his claim, that the Canons of the Fathers
be not infringed, nor the Priesthood be made an occasion and pretence
for the pride of worldly power, nor the least portion of that freedom
unawares he lost to us, which our LORD
JESUS CHRIST,
who bought the world's freedom, vouchsafed to us, when He shed His own
blood. Wherefore it has seemed good to this Holy Ecumenical Council,
that \emph{the rights of every province should be preserved pure and
inviolate, which have always belonged to it, according to the usage
which has ever obtained}, each Metropolitan having full liberty
to take a copy of the acts for his own security. And should any rule be
adduced repugnant to this decree, it is hereby repealed.''

Here we have a remarkable parallel to the dispute
between Rome and us; and we see what was the decision of the General
Church upon it. It will be observed, the decree is past \emph{for all
provinces in all future times}, as well as for the immediate
exigency. Now this is a plain refutation of the Romanists on their own
principles. \emph{They} profess to hold the Canons of the Primitive
Church: the very line they take, is to declare the Church to be one and
the same in all ages. Here then they witness against themselves. The
Pope \emph{has} encroached on the rights of other churches, and violated
the Canon above cited. Herein is the difference between his relation to
us, and that of any civil Ruler, whose power was in its origin illegally
acquired. Doubtless we are bound to obey the Monarch under whom we are
born, even though his ancestor were an usurper. Time legitimises a
conquest. But this is not the case in spiritual matters. The Church goes
by \emph{fixed laws}; and this usurpation has all along been counter to
one of her acknowledged standing ordinances, founded on reasons of
universal application.

After the Canon above cited, it is almost
superfluous to refer to the celebrated rule of the First Nicene Council,
A.D.
325, which, in defending the rights of the Patriarchates, expresses the
same principle in all its simple force and majesty.

``\emph{Let the ancient usages prevail}, which are
received in Egypt, Libya, and Pentapolis, relative to the authority of
the Bishop of Alexandria; as they are observed in the case of the Bishop
of Rome. And so in Antioch too, and other provinces, let the
prerogatives of the Churches be preserved.''

On this head of the subject, I will but notice,
that, as the Council of Ephesus controlled the ambition of Antioch, so
in like manner did St. Austin rebuke Rome itself for an encroachment of
another kind on the liberties of the African Church.

Bingham says,

``When Pope Zosimus and Celestine took upon them to
receive Appellants from the African Churches, and absolve those whom
they had condemned, St. Austin and all the African Churches sharply
remonstrated against this, as an irregular practice, \emph{violating}
\emph{the laws of unity}, and the settled rules of ecclesiastical
commerce; which required, that no delinquent excommunicated in one
Church should be absolved in another, without giving satisfaction to his
own Church that censured him. And therefore, to put a stop to this
practice and check the exorbitant power which Roman Bishops assumed to
themselves, they first made a Law in the Council of Milevis, That no
African Clerk should appeal to any Church beyond sea, under pain of
being excluded from communion in all the African Churches. And then,
afterwards, meeting in a general Synod, they dispatched letters to the
Bishop of Rome, to remind him how contrary this practice was to the
Canons of Nice, which ordered, That all controversies should be ended in
the places where they arose, before a Council and the Metropolitan.'' \footnote{Bingh. Antiq. xvi. 1. § 14.}

Thus I have shown, that our Bishops, at the time of
the Reformation, did but vindicate their ancient rights; were but acting
as grateful, and therefore jealous champions of the honour of the old
Fathers, and the sanctity of their institutions. Our duty surely in such
matters lies in neither encroaching nor conceding to encroachment; in
taking our rights as we find them, and using them; or rather in
regarding them altogether as trusts, the responsibility of which we
cannot avoid. As the same Apostle says, ``Let every man abide in the same
calling, wherein he is called.'' And, if England and Ireland had a plea
for asserting their freedom under any circumstances, much more so, when
the corruptions imposed on them by Rome even made it a duty to do so.

I shall answer briefly one or two objections, and
so bring these remarks to an end.

1. First, it may be said, that Rome has withdrawn
our orders, and excommunicated us; therefore we cannot plead any longer
our Apostolical descent. Now I will not altogether deny, that a
Ministerial Body might become so plainly apostate, as to lose its
privilege of ordination. But, however this may be, it is a little too
hard to \emph{assume}, as such an objection does, the very point in
dispute. \emph{When} we are proved to be heretical in doctrine, then
will be the time to begin to consider, whether our heresy is of so
grievous a character as to invalidate our orders; but, \emph{till then},
we may fairly and fearlessly maintain, that our Bishops are still
invested with the power of ordination.

2. But it may be said on the other hand, that if we
do not admit ourselves to be heretic, we necessarily must accuse the
Romanists of being such; and that therefore, on our own ground, we have
really no valid orders, as having received them from an heretical
Church. But even if Rome be so considered now, at least she was not
heretical in the primitive ages; no one will say that she was then
Antichrist \footnote{The following is from the
Life of Bernard Gilpin, vid. Wordsworth's Ecclesiastical Biography, vol.
iv. p. 94. ``Mr. Gilpin would often say that the Churches of the
Protestants were not able to give any firme and solid reason of their
separation besides this, to wit, that the Pope is Antichrist … The
Church of Rome kept the rule of faith intire, until that rule was
changed and altered \emph{by the Council of Trent, and from that time it
seemed to him a matter of necessitie to come out of the Church of Rome},
that so that Church which is true and called out from thence might
follow the word of God ... But he did not these things violently, but by
degrees.''}. Nay, as to the
middle ages, we may say with the learned Dr. Field, ``that none of those
points of false doctrine and error which Romanists now maintain, and we
condemn, were the doctrines of the Church before the Reformation
constantly delivered or generally received by all them that were of it,
but doubtfully broached, and devised without all certain resolution, or
factiously defended by some certain only, who as a dangerous faction
adulterated the sincerity of the Christian verity, and brought the
Church into miserable bondage.'' \footnote{See Field on the Church,
Appendix to book iii. where he proves all this. See also Birkbeck's
Protestant's Evidence.}
Accordingly, acknowledging and deploring all the errors of the middle
ages, yet we need not fear to maintain, that after all they were but the
errors of individuals, though of large numbers of Christians; and we may
safely maintain, that they no more interfere with the validity of the
ordination received by our Bishops from those who lived before the
Reformation, than errors of faith and conduct in a priest interfere with
the grace of the Sacraments received at his hands.

3. It may be said, that we throw blame on Luther,
and others of the foreign Reformers, who did act without the authority
of their Bishops. But we reply, that it has been always agreeable to the
principles of the Church, that, if a Bishop taught and upheld what was
contrary to the orthodox faith, the Clergy and people were not bound to
submit, but were obliged to maintain the true religion; and if
excommunicated by such Bishops, they were never accounted to be cut off
from the Church. Luther and his associates upheld in the main the true
doctrine; and though it is not necessary to defend \emph{every} act of
fallible men like them, yet we are fully justified in maintaining, that
the conduct of those who defended the truth against the Romish party,
even in opposition to their spiritual rulers, was worthy of great
praise. At the same time it is impossible not to lament, that they did
not take the first opportunity to place themselves under orthodox
Bishops of the Apostolical Succession. Nothing, as far as we can judge,
was more likely to have preserved them from that great decline of
religion, which has taken place on the Continent.

[NEW EDITION.]

———————————————————————

These Tracts are continued in
Numbers, and sold at the price of 2d. for each sheet, or 7s. for 50
copies.

LONDON: PRINTED FOR
J. G. F. \& J. RIVINGTON,

ST. PAUL'S CHURCH YARD, AND WATERLOO PLACE.

1840.

\xchapter{On Arguing concerning the Apostolical Succession}

Tract No. 19

MEN
are sometimes disappointed with the proofs offered in behalf of some
important doctrines of our religion; such especially as the necessity of
Episcopal Ordination, in order to constitute a Minister of CHRIST. They consider these proofs to be not so
strong as they expected, or as they think desirable. Now such persons
should be asked, whether these arguments they speak of are in their
estimation weak as a guide to their own practice, or weak in controversy
with hardheaded and subtle disputants. Surely, as Bishop Butler has
convincingly shown, the faintest probabilities are strong enough to
determine our \emph{conduct} in a matter of duty. If there be but a
reasonable likelihood of our pleasing CHRIST
more by keeping than by not keeping to the fellowship of the Apostolic
Ministry, this of course ought to be enough to lead those, who think
themselves moved to undertake the Sacred Office, to seek for a licence
to do so from it.

It is necessary to keep this truth distinctly in
view, because of the great temptation, that exists among us, to put it
out of sight. I do not mean the temptation, which results from
pride,—hardness of heart,—a profane disregard of the details and
lesser commandments of the Divine Law,—and other such like bad
principles of our nature, which are in the way of our honestly
confessing it. Besides these, there is a still more subtle temptation to
slight it, which will bear insisting on here, arising from an
over-desire to convince others, or, in other words, a desire to
out-argue others, a fear of seeming inconclusive and confused in our own
notions and arguments. Nothing, certainly, is more natural, when we hold
a truth strongly, than to wish to persuade others to embrace it also.
Nay, without reference to persuasion, nothing is more natural than to be dissatisfied in all cases with our own convictions of a principle
or opinion, nay suspicious of it, till we are able to set it down
clearly in words. We know, that, in all matters of thought, to write
down our meaning is one important means of clearing our minds. Till we
do so, we often do not know what we really hold, and what we do not
hold. And a cautious and accurate reasoner, when he has succeeded in
bringing the truth of any subject home to his mind, next begins to look
round about the view he has adopted, to consider what others will say to
it, and try to make it unexceptionable. At least we are led thus to
fortify our opinion, when it is actually attacked; and if we find we
cannot recommend it to the judgment of the assailant, at any rate we
endeavour to make him feel that it is to be respected. It is painful to
be thought a weak reasoner, even though we are sure in our minds that we
are not such.

Now, observe how these feelings will affect us, as
regards such arguments as were alluded to above; viz. such as are open
to exception, though they are sufficiently strong to determine our
conduct. A friend, who differs from us, asks for our reasons for our own
view. We state them, and he sifts them. He observes, that our
conclusions do not necessarily follow from our premises. E.g. to take
the argument for the Apostolical Succession derived from the ordination
of St. Paul and St. Barnabas (Acts xiii. 2, 3), he will argue, that
their ordination \emph{might} have been an accidental rite, intended
merely to commission them for their Missionary journey, which followed
it, in Asia Minor; again, that St. Paul's direction to Timothy (1 Tim.
v. 22), to ``lay hands suddenly on no man,'' \emph{may} refer to
confirmation, not ordination.

We should reply (and most reasonably too), that, \emph{considering
the undeniable fact} that ordination has ever been thought necessary
in the Church for the Ministerial Commission, our interpretation is the
most probable one, and therefore the safest to act upon; on which our
friend will think awhile, then shake his head, and say, that ``at all
events this is an \emph{unsatisfactory} mode of reasoning, that it does
not convince \emph{him}, that he is desirous of clearer light,'' \&c.

Now what is the consequence of such a discussion as
this on ourselves? not to make us \emph{give up} the doctrine, but
to make us afraid of \emph{urging} it. We grow lukewarm about it; and,
with an appearance of judgment and caution (as the world will call it),
confess that ``to rest the claims of our Clergy on an Apostolical Descent
is an unsafe and inexpedient line of argument; that it will not convince
men, the evidence not being sufficient; that it is not a practical way
of acting to insist upon it,'' \&c.—whereas the utmost that need be
admitted is, that it is out of place to make it the subject of a
speculative dispute, and to argue about it on that abstract logical
platform which virtually excludes a reference to conduct and duty. And
indeed, it would be no unwise caution to bear about us, wherever we go,
that our first business, as Christians, is to address men as responsible
servants of CHRIST,
not as antagonists; and that it is but a secondary duty (though a duty)
to ``refute the gainsayers.''

And, as on the one hand it continually happens,
that those who are most skilled in debate are deficient in sound
practical piety, so on the other it may be profitable to us to reflect,
that doctrines, which we believe to be most true, and which are received
as such by the most profound and enlarged intellects, and which rest
upon the most irrefragable proofs, yet may be above \emph{our}
disputative powers, and can be treated by us only with reference to our
conduct. And in this way, as in others, is fulfilled the saying of the
Apostle, that ``the preaching of the Cross is to them that perish
foolishness; but unto us, who are saved, it is the power of GOD
… Where is the wise? where is the scribe? where is the disputer of
this world? hath not GOD
made foolish the wisdom of this world? ... The foolishness of GOD is wiser than men; and the weakness of GOD
is stronger than men.''

ON RELUCTANCE TO CONFESS THE

APOSTOLICAL SUCCESSION

IF
a Clergyman is quite convinced that the Apostolical Succession is lost,
then of course he is at liberty to turn his mind from the subject. But
if he is not quite sure of this, it surely is his duty seriously to
examine the question, and to make up his mind carefully and
deliberately. For if there be a chance of its being preserved to us,
there is a chance of his having had a momentous talent committed to him,
which he is burying in the earth.

It cannot be supposed that any serious man would
treat the subject scoffingly. If any one is tempted to do so, let him
remember the fearful words of the Apostle. ``Esau, a \emph{profane person},
who for one morsel of meat, sold his birthright.''

If any are afraid that to insist on their
commission will bring upon them ridicule, and diminish their usefulness,
let them ask themselves whether it be not cowardice to refuse to leave
the event to GOD.
It was the reproach of the men of Ephraim, that, though they were ``harnessed
and carried bows,'' they ``turned themselves back in the day of battle.''

And if any there be, who take upon them to contrast
one doctrine of the Gospel with another, and preach those only which
they consider the more essential, let them consider our SAVIOUR'S words, ``These things ought ye to have done, and not to
leave the other undone.''

OXFORD,

\emph{Dec}. 23, 1833.

[SIXTH EDITION.]

———————————————————————

These Tracts are continued in
Numbers, and sold at the price of 2d. for each sheet, or 7s. for 50
copies.

LONDON: PRINTED FOR
J. G. F. \& J. RIVINGTON,

ST. PAUL'S CHURCH YARD, AND WATERLOO PLACE.

1840.

\xchapter{The Visible Church Letters to a Friend No. III.}

Tract No. 20 (Ad Scholas)

YOU
have some misgivings, it seems, lest the doctrine I have been advocating
``should lead to Popery.'' I will not, by way of answer, say that the
question is not, whether it will \emph{lead to Popery}, but whether it
is \emph{in the Bible}; because it would bring the Bible and Popery into
one sentence, and seem to imply the possibility of a ``communion'' between
``light and darkness.'' No; it is the very enmity I feel against the
Papistical corruptions of the Gospel, which leads me to press upon you a
doctrine of Scripture, which we are sinfully surrendering, and the
Church of Rome has faithfully retained.

How comes it that a system, so unscriptural as the
Popish, makes converts? because it has in it an element of truth and
comfort amid its falsehoods. And the true way of opposing it, is, not to
give up to them that element, which GOD'S
providence has preserved to us also, thus basely surrendering ``the
inheritance of our Fathers,'' but to claim it as our own, and to make use
of it for the purposes for which GOD
has given it to us. I will explain what I mean.

Before CHRIST came, Divine Truth was, as it were, a pilgrim in the world.
The Jews excepted, men who had portions of the SPIRIT of GOD, knew not their privilege. The whole force
and current of the external world was against them, acting powerfully on
their imagination, and tempting them to set sight against faith, to
trust the many witnesses who prophesied falsehood (as if) in the name of
the LORD,
rather than the still small voice which spoke within them. Who can
undervalue the power of this fascination, who has had experience of the
world ever so little? Who can go at this day into mixed society, who can
engage in politics or other active business, and not find himself
gradually drifting off from the true Rock on which his faith is built,
till he begins in despair to fancy, that solitude is the only safe place
for the Christian, or, (with a baser judgment,) that strict obedience
will not be required at the last day of those who have been engaged in
active life? If such is now the power of the world's enchantments,
surely much greater was it before our SAVIOUR came.

Now what did He do for us, in order to meet this
evil? His merciful Providence chose means which might act as a
counter-influence on the imagination. The visible power of the world
enthralled men to a lie; He set up a Visible Church, to witness the
other way, to witness for Him, to be a matter of fact, as undeniable as
the shining of the sun, that there \emph{was} such a principle as
conscience in the world, as faith, as fear of GOD;
that there \emph{were} men who considered themselves bound to live
as His servants. The common answer which we hear made every day to
persons who engage in any novel undertaking, is, ``You will get no one to
join you; nothing can come of it; you are singular in your opinion; you
do not take practical views, but are smit with a fancy, with a dream of
former times,'' \&c. How cheering it is to a person so circumstanced,
to be able to point to others elsewhere, who actually hold the same
opinions as himself, and exert themselves for the same objects! Why?
because it is an appeal to a \emph{fact}, which no one can deny; it is
an evidence that the view which influences him is something external to
his own mind, and not a dream. What two persons see, cannot be an ideal
apparition. Men are governed by such facts much more than by
argumentative proof. These act upon the imagination. Let a person be
told ten times over that an opinion is true, the \emph{fact} of its
being said becomes an argument for the truth of it; \emph{i.e}. it is so
with most men. We see from time to time the operation of this principle
of our nature in political matters. Our American colonies revolt; France
feels the sympathy of the event, and is revolutionized. Again, in the
same colonies, the Episcopal Church flourishes; we Churchmen at home
hail it as an omen of the Church's permanence among ourselves. On the
other hand, what can be more dispiriting than to find a cause, which we
advocate, sinking in some other country or neighbourhood, though there
be no reason for concluding, that, \emph{because} it has fallen
elsewhere, therefore it will among ourselves. In order then to supply
this need of our minds, to satisfy the imagination, and so to help our
faith, for this among other reasons CHRIST
set up a visible Society, His Church, to be as a light upon a hill, to
all the ends of the earth, while time endures. It is a witness of the
unseen world; a pledge of it; and a prefiguration of what hereafter will
take place. It prefigures the ultimate separation of good and bad, holds
up the great laws of GOD'S
Moral Governance, and preaches the blessed truths of the Gospel. It
pledges to us the promises of the next world, for it is something (so to
say) in hand; CHRIST
has done one work as the earnest of another. And it witnesses the truth
to the whole world; awing sinners, while it enspirits the fainting
believer. And in all these ways it helps forward the world to come; and
further, as the keeper of the Sacraments, it is an essential means of
the realizing it at present in our fallen race. Nor is it much to the
purpose, as regards our duty towards it, what are the feelings and
spiritual state of the individuals who are its officers. True it is,
were the Church to teach heretical doctrine, it might become incumbent
on us (a miserable obligation!) to separate from it. But, while it
teaches substantially the Truth, we ought to look upon it as one whole,
one ordinance of GOD,
not as composed of individuals, but as a house of GOD'S building;—as an instrument in His hand,
to be used and reverenced for the sake of its Maker.

Now the Papists have retained it; and so they have
the advantage of possessing an instrument, which is, in the first
place, suited to the needs of human nature; and next, is a special gift
of CHRIST,
and so has a blessing with it. Accordingly we see that in its measure
success follows their zealous use of it. They act with great force upon
the imaginations of men. The vaunted antiquity, the universality, the
unanimity of their Church puts them above the varying fashions of the
world, and the religious novelties of the day. And truly when one
surveys the grandeur of their system, a sigh arises in the thoughtful
mind, to think that we should be separate from them; Cum talis esses,
utinam noster esses!—But, alas! AN UNION IS IMPOSSIBLE. Their communion is
infected with heterodoxy; we are bound to flee it, as a pestilence. They
have established a lie in the place of GOD'S truth and, by their claim of immutability in doctrine,
cannot undo the sin they have committed. They cannot repent. Popery must
be destroyed; it cannot be reformed.

Now then what is the Christian to do? Is he forced
back upon that cheerless atheism (for so it practically must be
considered) which prevailed in the world before CHRIST'S
coming, poorly alleviated, as it was, by the received polytheisms of the
heathen? Can we conceive a greater calamity to have occurred at the time
of our Reformation, one which the Enemy of man would have been more set
on effecting, than to have entangled the whole of the Church Catholic in
the guilt of heterodoxy, and so have forced every one who worshipped in
spirit and in truth, to flee out of doors into the bleak world, in order
to save his soul? I do not think that Satan could have desired any event
more eagerly, than such an alternative; viz. to have forced Christians,
either to remain in communion with error, or to join themselves in some
such spontaneous union among themselves, as is dissolved as easily as it
is formed. Blessed be GOD! his malice has been thwarted. I do
believe it to be one most conspicuous mark of GOD'S adorable Providence over us, as great as
if we saw a miracle, that Christians in England escaped in the evil day
from either extreme, neither corrupted doctrinally, nor secularized
ecclesiastically. Thus in every quarter of the world, from North America
to New South Wales, a Zoar has been provided for those who would fain
escape Sodom, yet dread to be without shelter. I hail it as an omen amid
our present perils, that our Church will not be destroyed. He hath been
mindful of us; He will bless us. He has wonderfully preserved our Church
as a true branch of the Church universal, yet withal preserved it free
from doctrinal error. It is Catholic and Apostolic, yet not Papistical.

With this reflection before us, does it not seem to
be utter ingratitude to an astonishing Providence of GOD'S mercy, to be neglectful, as many
Churchmen now are, of the gift; to attempt unions with those who have
separated from the Church, to break down the partition walls, and to
argue as if religion were altogether and only a matter of each man's
private concern, and that the State and Nation were not bound to
prefer the Apostolical Church to all self-originated forms of
Christianity? But this is a point beside my purpose. Take the matter
merely in the light of human expedience. Shall we be so far less wise in
our generation than the children of this world, as to relinquish the
support which the Truth receives from the influence of a visible Church
upon the imagination, from the energy of operation which a well
disciplined Body ensures? Shall we not foil the Papists, not with their
own weapons, but with weapons which are ours as well as theirs? or, on
the other hand, shall we with a melancholy infatuation give them up to
them? Depend upon it, to insist on the doctrine of the Visible Church is
not to favour the Papists, it is to do them the most serious injury. It
is to deprive them of their only strength. But if we neglect to do so,
what will be the consequence? Break down the Divine Authority of our
Apostolical Church, and you are plainly preparing the way for Popery in
our land. Human nature cannot remain without visible guides; it chooses
them for itself, if it is not provided for them. If the Aristocracy and
the Church fall, Popery steps in. Political events are beyond our power,
and perhaps out of our sphere; but ecclesiastical matters are in the
hands of all Churchmen.

OXFORD,

\emph{Dec}. 24, 1833.

[SIXTH EDITION.]

———————————————————————

These Tracts are continued in
Numbers, and sold at the price of 2d. for each sheet, or 7s. for 50
copies.

LONDON: PRINTED FOR
J. G. F. \& J. RIVINGTON,

ST. PAUL'S CHURCH YARD, AND WATERLOO PLACE.

1840.

\xchapter{Mortification of the Flesh a Scripture Duty}

Tract No. 21 (Ad Populum)

IF we take the example of the Holy men of
Scripture as our guide, certainly bodily privation and chastisement are
a very essential duty of all who wish to serve GOD, and prepare themselves for His presence.

1. First, we have the example of Moses. His
recorded Fasts were miraculous; still they were Fasts, and the ordinance
was recommended to the notice of all believers afterwards, by the honour
put upon it. ``I abode in the mount forty days and forty nights; I
neither did eat bread nor drink water.'' Again; ``I fell down before the LORD, as at the first, forty days and forty nights; I did neither
eat bread nor drink water, because of all your sins.'' Deut. ix. 9, 18.
Fasting is in the former instance subservient to divine contemplation,
in the latter to humiliation and intercession for sinners.

Elijah. ``He said unto him, What manner of man was
he which came up to meet you, and told you these words? And they
answered him, He was an hairy man, and girt with a girdle of leather
about his loins. And he said, It is Elijah the Tishbite.'' 2 Kings i. 7,
8. It is indeed needless to show the ascetic character of him who was in
fact the chief and type of those who ``wandered about in sheepskins and
goatskins,'' ``in deserts, and in mountains, and in dens and caves of the
earth.'' He too fasted by the power of GOD
for forty days and nights; ``He arose and did eat and drink, and went in
the strength of that meat forty days and forty nights, unto Horeb the
mount of GOD'' 1 Kings xix. 8.

Daniel. ``I set my face unto the LORD GOD, to seek by prayer and supplications, with
fasting, and sackcloth, and ashes; and I prayed unto the LORD
my GOD,
and made my confession.'' Dan. ix. 3, 4. It must be observed, that Daniel
was not bound by any vow, as Samson and Samuel. Moreover it would appear
the gift of prophecy was given him in reward for his self-chastisements,
as the following passage shows. ``In those days I Daniel was mourning
three full weeks; I ate no pleasant bread, neither came flesh nor wine
in my mouth; neither did I anoint myself at all, till three whole weeks
were fulfilled … And he said unto me, O Daniel, a man greatly beloved,
understand the words that I speak unto thee, and stand upright; for unto
thee am I now sent … Fear not, Daniel ; for \emph{from the first day}
that thou didst set thine heart to understand, and \emph{to chasten thyself}
before thy GOD, thy words were heard, and I am come for thy words.'' Dan. x.
2, 3, 11, 12. Vide also Luke ii. 37. Acts x. 30.

2. Now here it will be objected, perhaps, that
these instances are taken from the Old Testament, and belong to the Law
of Moses, which is not binding on Christians.

I answer:

(1.) That in the above passages Fasting is
connected with moral acts, humiliation, prayer, meditation, which are
equally binding on us as on the Jews. Man is now what he was then; and
if affliction of the flesh was good then, it is now.

(2.) In matter of fact, \emph{private} Fasting,
such as instanced in the passages above quoted, was no special duty of
the Mosaic Law. Public Fasting, indeed, was on one occasion enjoined by
Moses himself, and on others by subsequent rulers; but this was in part
a ceremonial act, not a moral discipline, and was doubtless abolished
with the other rites of the Law.

``Of Fasts,'' says Lewis, ``there was no more than one
appointed by the Law of Moses, called the Fast of Expiation … The
great day of Expiation was a most severe Fast, kept every year upon the
tenth day of the month Tizri, which answers to our September … This
solemnity was observed with fasting and abstinence, not only from all
meat and drink, but from all other pleasure whatsoever; insomuch that
they did not wash their faces, much less anoint their heads, nor wear
their shoes … nor, (if their Doctors say true,) read any portion of
the law which would give them delight. They refrained likewise not only
from pleasure, but from labour, nothing being to be done upon this day,
but confessing of sins and repentance.'' \footnote{Lewis, Hebrew Republic, iv. 15}

Nay, it may rather be said, that the Jewish Law, as
such, was rather opposed than otherwise to austerities. The Nazarites
and Rechabites, being exceptions to the rule, are evidence of it. Vide,
on the other hand, Deut. xii. Eccles. v. 18 \footnote{Vide Spencer de Regg. Hebræor. lib. 3. diss. 1. ii. 3. diss. 4. i. 5,
\&c.}.

Such then being the character of the Law in its
formal letter, it tells just the contrary way to that which superficial
reasoners might expect. For it is most remarkable, first, that the
greatest prophets under it, such as Elijah and Daniel, were without
express command singularly austere and self-afflicting men, in the midst
of a people, who from the first went lusting after ``the fish which they
eat in Egypt freely; the cucumbers, and the melons, and the leeks, and
the onions, and the garlick, and said, Who shall give us flesh to eat?''
Next there is something of a very startling and admonitory nature in the
\emph{miraculous} fasts of Moses and Elijah, under this same imperfect
dispensation. The miracle evidently was for some purpose; yet it did not
sanction, in any direct way, any injunction of the Law. Was it not an
admonition to the Israelites, that there was a more excellent way of
obedience than that which ALMIGHTY GOD
as yet thought fit to promulgate by solemn enactment? Is it not an
intimation serviceable for Christian practice, as much as Moses'
announcement of the destined ``Prophet like unto him'' is intended for the
comfort of Christian faith?

Surely the duty of bodily discipline might be
rested on the answer to this plain question, \emph{Why} did Daniel use
austerities not enjoined by the Law?

3. Now turn to the New Testament, and observe what
clear light is therein thrown upon the duty already recommended to us by
the Old Testament Saints.

First, there is the instance of St. John the
Baptist. ``John came neither eating nor drinking,'' Matt. xi. 18: and his
disciples fasted, Matt. ix. 14.

Our SAVIOUR
did \emph{not} statedly fast; but here also the exception proves the
rule. He who did not fast statedly was the only one born of woman who
was untainted by sinful flesh; which seems to imply, that all who are
natural descendants of guilty Adam ought to fast.

He bade His disciples to fast. Consider his implied
precept, which is an express command to those who obey the Law of
Liberty. ``When thou fastest, anoint thy head, and wash thy face, that
thou appear not unto men to fast.'' Matt. vi. 17, 18.

Consider, moreover, the \emph{general austere}
character of Christian obedience, as enjoined by our LORD;—a circumstance much to be insisted on
in an age like this, when what is really self-indulgence is thought to
be a mere moderate and innocent use of this world's goods. I will but
refer to a few, out of many texts, which I am persuaded are now
forgotten by numbers of educated and amiable men who are fond of
extolling what they call the mild, tolerant, enlightened spirit of the
Gospel. Matt. v. 29, 30. vii. 13, 14. x. 37-39. Mark ix. 43-50. x. 25.
Luke xiv. 12, 26-33.

And reflect, too, whether the spirit of texts such
as the following will not move every true member of the Church Militant.
``The ark, and Israel, and Judah abide in tents; and my lord Joab, and
the servants of my lord, are encamped in the open fields; shall I then
go into mine house, to eat and to drink? … as thou liveth, and as thy
soul liveth, I will not do this thing.'' 2 Sam. xi. 11.

Now take the example of the Apostles. St. Peter was
fasting, when he had the vision which sent him to Cornelius: Acts x. 10.
The prophets and teachers at Antioch were fasting, when the HOLY GHOST revealed to them His purpose about Saul
and Barnabas: Acts iii. 2, 3. Vide also Acts xiv. 23. 2 Cor. vi. 5. xi.
27.

Weigh well the following text, which I am persuaded
many men would deny to be St. Paul's writing, had not a gracious
Providence preserved to us the epistle containing it. ``I keep under my
body, and bring it into subjection; lest that by any means, when I
have preached to others, I myself should be a cast-away.'' 1 Cor. ix. 27.

4. Lastly, Consider the practice of the Primitive
Christians.

The following account of the early Christian Fasts,
is from Bingham, Antiq. lib. xxi.

THE
QUADRAGESIMAL
OR LENT
FAST.—``The
Quadragesimal Fast before Easter,'' says Sozomen, ``some observe six
weeks, as the Illyrian and Western Churches, and all Libya, Egypt, and
Palestine; others make it seven weeks, as the Constantinopolitans and
neighbouring nations as far as Phœnicia; others fast three only of
those six or seven weeks, by intervals; others the three weeks next
immediately before Easter.''

The manner of observing Lent among those that were
piously disposed to observe it, was to abstain from all food till
evening. For anciently a change of diet was not reckoned a fast: but it
consisted in perfect abstinence from all sustenance for the whole day
till evening.

THE
FASTS OF
THE FOUR
SEASONS.—The
next Anniversary fasting days were those which were called \emph{Jejunia
quatuor temporum}, the Fasts of the Four Seasons of the Year …
These were at first designed … to beg a blessing of God upon the
several seasons of the year, or to return thanks for the benefits
received in each of them, or to exercise and purify both body and soul
in a more particular manner, at the return of these certain terms of
stricter discipline and more extraordinary devotion. [These afterwards
became the Ember Fasts.]

MONTHLY
FASTS.—In
some places they had also Monthly Fasts throughout the year excepting in
the two months of July and August … because of the sickness of the
season.

WEEKLY
FASTS.—Besides
these they had their weekly Fasts or, Wednesday and Friday, called the
Stationary Days, and Half-Fasts or Fasts of the Fourth and Sixth Days of
the Week … These Fasts being of continual use every week throughout
the Year, except in the Fifty Days between Easter and Pentecost, were
not kept with that rigour and strictness which was observed in the time
of Lent … [but] ordinarily held no longer than 9 o'clock, i.e. 3 in
the afternoon.

OXFORD,

\emph{The Feast of the Circumcision}.

1834.

[SIXTH EDITION.]

———————————————————————

These Tracts are continued in
Numbers, and sold at the price of 2d. for each sheet, or 7s. for 50
copies.

LONDON: PRINTED FOR
J. G. F. \& J. RIVINGTON,

ST. PAUL'S CHURCH YARD AND WATERLOO PLACE.

1840.

\xchapter{The Reformed Church}

Tract No. 31 (Ad Clerum)

All the people shouted with a great shout, when
they praised the LORD,
because the foundation of the House of the LORD was laid. But many of the Priests and
Levites, the chief of the fathers, who were ancient men that had seen
the first House, when the foundation of this House was laid before their
eyes, wept with a loud voice.—EZRA
iii. 11, 12.

SOME
remarks may, perhaps, be profitably made on the following well known
lines in Herbert's Church Militant, in which the text above quoted is
applied to our own period:

\begin{quote}

The second Temple could not reach
the first,

And the late Reformation never durst

Compare with ancient times and purer years,

But in the Jews and us, deserveth tears.

Nay, it shall every year decrease and fade,

Till such a darkness shall the world invade

At CHRIST'S last
coming, as His first did find;

Yet must their proportions be assigned

To these diminishings, as is between

The spacious world and Jewry to be seen.

\end{quote}

Surely there is a close analogy between the state
of the Jews after the captivity, and our own; and, if so, a clear
understanding and acknowledgment of it will tend to teach us our own
place and suggest to us our prospects.

1. It is scarcely necessary to notice the general
correspondence between the fortunes of the two Churches. Both Jews and
Christians ``left their first love,'' mixed with the world, were brought
under the power of their enemies, went into captivity, and at length,
through GOD'S
mercy, were brought back again from Babylon. Ezra and Nehemiah are the
forerunners of our Hookers and Lauds; Sanballat and Geshem of the
disturbers of our Israel. Samaria has set up its rival temple among us.

2. The second Temple lacked the peculiar treasures
of the Temple of Solomon, the Prince of Peace; such as the Ark, the
visible glory of GOD,
the tables of the Covenant, Aaron's rod, the manna, the oracle. In like
manner the Christian Church was, in the beginning, set up in unity;
unity of doctrine, or \emph{truth}, unity of discipline, or \emph{Catholicism},
unity of heart, or \emph{charity}. In spite of the heresies which then
disturbed the repose of Christians, consider the evidences which present
themselves in ecclesiastical history of their firm endurance of
persecution, their tender regard for the members of CHRIST,
however widely removed by place and language, their self-denying
liberality in supplying their wants, the close correspondence of all
parts of the body Catholic, as though it were but one family, their
profound reverential spirit towards sacred things, the majesty of their
religious services, and the noble strictness of their life and
conversation. Here we see the ``Rod'' of the Priesthood, budding forth
with fresh life; the ``Manna'' of the Christian ordinances uncorrupted;
the ``Oracle'' of Tradition fresh from the breasts of the Apostles; the ``Law,''
written in its purity on ``the fleshly tables of the heart;'' the ``Shechinah,''
which a multitude of Martyrs, Saints, Confessors, and gifted Teachers,
poured throughout the Temple. But where is our unity now? our
ministrations of self-denying love? our prodigality of pious and
charitable works? our resolute resistance of evil? We are reformed; we
have come out of Babylon, and have rebuilt our Church; but it is Ichabod;
``the glory is departed from Israel.''

3. The Jewish polity was, on its restoration, so
secularized, that the vestiges of a Theocracy scarcely remained in the
eyes of any but attentive believers. That it really existed as before,
is plain from the prophetic gift possessed by Caiaphas, wicked man as he
was. Consider the anomaly of the political relation of the Jews towards
the Ptolemies and Seleucidæ, their alliance with Rome, their dispersion
over the Roman Empire, their disuse of certain of the Mosaic ordinances,
the cruelties and blasphemies of Antiochus, the reign of Herod, and his
virtual rebuilding of the Temple, a remarkable omen as regards
ourselves. Turn to the restored Christian Church, and reflect upon the
perplexed questions concerning the union of Church and State, to which
the politics of the last three centuries have given rise: the tyrannical
encroachments of the civil power at various eras; the profanations at
the time of the Great Rebellion; the deliberate impiety of the French
Revolution; and the present apparent breaking up of Ecclesiastical
Polity every where, the innumerable schisms, the mixture of men of
different creeds and sects, and the contempt poured upon any show of
Apostolical zeal.

4. Consider the following passages, from the
Prophets, after the Captivity, and see if they do not apply to present
times.

Hagg. i. 4-10. ``Is it time for you, O ye, to dwell
in your ceiled houses, and this house lie waste? Now, therefore, thus
saith the LORD
of Hosts, consider your ways. Ye have \emph{sown much}, and \emph{bring in
little}; ye eat, but ye have not enough; ye drink, but ye are not
filled with drink; ye clothe you, but there is none warm; and he that
earneth wages, \emph{earneth wages to put it into a bag with holes},''
\&c.

Mal. i. 6-13. ``A son honoureth his father, and a
servant his master; if then I be a Father, where is Mine honour? and if
I be a Master, where is My fear? … Ye say, The \emph{table of the Lord is
polluted}, and \emph{the fruit thereof, even His meat, contemptible}.
Ye say also, Behold what a \emph{weariness} is it, … and \emph{ye brought
that which was torn, and the lame, and the sick}; thus ye brought an
offering; should I accept this of your hands, saith the LORD?''

Mal. ii. 1-9. ``And now, O ye Priests, this
commandment is for you ... And ye shall know that I have sent this
commandment unto you, that My covenant might be with Levi, saith the LORD of Hosts. My covenant was with him of life
and peace, and I gave them to him, for the fear wherewith he feared Me,
and was afraid before My Name. The Law of Truth was in his mouth, and
iniquity was not found in his lips; he walked with Me in peace and
equity, and did turn many away from iniquity. For the Priest's lips
should keep knowledge, and they shall seek the Law at his mouth; for he
is the messenger of the LORD
of hosts. \emph{But ye are departed out of the way}; ye have caused many
to stumble at the Law: \emph{ye have corrupted the covenant of Levi},
saith the LORD of Hosts. \emph{Therefore have I also
made you contemptible and base before all the people}.'' Does not the
history of the times of Hoadly and such as he, and our present trials,
throw light upon the parallel?

Mal. iii. 8, 9. ``Will a man rob GOD? yet ye have robbed Me; but ye say,
Wherein have we robbed Thee? \emph{in tithes and offerings}. \emph{Ye are
cursed with a curse}; for ye have robbed Me, even this whole nation.''

5. It is remarkable that, while the reinstated
Jewish Church was so deficient in zeal, piety, and consistent obedience,
and was punished by failure and disorganization; yet it never fell into
those gross and flagrant offences, which were the opprobrium of its
earlier period. \emph{It was clear of the sin of idolatry}.

6. Moreover consider the \emph{parties}, unknown to
the era of the Theocracy, which divided the Church after the captivity;
the Pharisees, Sadducees, and the rest; the necessary consequence of a
relaxation of the original principle of national union. The case is the
same in this day; as if the Church were already dead, new forms of
organization, multiplied varieties of life and action, show themselves
within her.

7. Lastly. The following texts suggest hope to all
true Christians. (Hagg. ii. 5-9.) ``\emph{According to the word that I
covenanted with you, when ye came out of Egypt}, so MY SPIRIT REMAINETH
AMONG YOU: fear ye not.'' He will be with us even in this base and
grovelling age, as with St. Paul, St. Cyprian, and St. Athanasius.

\begin{quote}

``Thou wilt; for Thou art Israel's
GOD;

And thine
unwearied arm

\emph{Is ready yet} with Moses' rod,'' \&c.

\end{quote}

``The glory of this latter house SHALL BE GREATER THAN OF THE FORMER, saith the
LORD of
Hosts.''

Strange it now seems before the event, how the
Church should close both with glory and yet in unbelief; yet surely, as
in the history of Jerusalem, so now both predictions will be at once
fulfilled. (Mal. iv. 1, 2.) ``The day cometh that shall burn as an oven,
and \emph{all the proud}, yea, and \emph{all who do wickedly shall be
stubble}: but \emph{unto you that fear My name shall the Sun of
Righteousness arise with healing in His wings}.''

And let it be remembered, that when our Lord seems
at greatest distance from His Church, then He is even at the doors.
Doubtless, when the Angel appeared in the Temple to Zacharias, the news
of a miraculous interposition was as great a marvel to the world at
large as if it were now noised abroad of one of our own Ministers in the
course of his Christian Service.

OXFORD,

\emph{The Feast of St. Mark}

———————————————————————

These Tracts are continued in
Numbers, and sold at the price of 2d. for each sheet, or 7s. for 50
copies.

LONDON: PRINTED FOR
J. G. F. \& J. RIVINGTON,

ST. PAUL'S CHURCH YARD, AND WATERLOO PLACE.

1840.

\xchapter{Primitive Episcopacy}

Tract No. 33 (Ad Scholas)

IN
primitive times the first step towards evangelizing a heathen country
seems to have been to seize upon some principal city in it, commonly the
civil metropolis, as a centre of operation; to place a Pastor, \emph{i.e}.
(generally) a Bishop there; to surround him with a sufficient number of
associates and assistants; and then to wait, till, under the blessing of
Providence, this Missionary College was able to gather around it the
scattered children of grace from the evil world, and invest itself with
the shape and influence of an organized Church. The converts would, in
the first instance, be those in the immediate vicinity of the Missionary
or Bishop, whose diocese nevertheless would extend over the heathen
country on every side, either indefinitely, or to the utmost extent of
the civil province; his mission being without restriction to all to whom
Christian faith had never been preached. As he prospered in the increase
of his flock, and set out his clergy to greater and greater distances
from the city, so would the homestead (so to call it) of his Church
enlarge. Other towns would be brought under his government, openings
would occur for stations in isolated places; till at length ``the burden
becoming too heavy for him,'' he would appoint others to supply his place
in this or that part of the province. To these he would commit a greater
or lesser share of his spiritual power, as might be necessary: sometimes
he would make them fully his representatives, or ordain them Bishops; at
other times he would employ presbyters for his purpose. These
assistants, or (as they were called) Chorepiscopi, would naturally be
confined to their respective districts; and if Bishops, an approximation
would evidently be made to a division of the large original diocese into
a number of smaller ones connected with and subordinate to the Bishop of
the metropolitan city. Thus, from the very Missionary character of the
Primitive Church, there was a tendency in its polity to what was
afterwards called the Provincial and Patriarchal system.

It is not, indeed, to be supposed that this was the
only way in which the graduated order of sees (so to call it)
originated; but, at least, it is one way. And there is this advantage in
remarking it: we learn from it, that large dioceses are the
characteristics of a church in its infancy or weakness; whereas, the
more firmly Christianity was rooted in a country, and the more vigorous
its rulers, the more diligently were its sees multiplied throughout the
ecclesiastical territory. Thus, St. Basil, in the fourth century,
finding his exarchate defenceless in the neighbourhood of Mount Taurus,
created a number of' dioceses to meet the emergency. These subordinate
sees may be called suffragan to the Metropolitan Church, whether their
respective rulers were mere representatives of the Bishop who created
them, \emph{i.e}., Chorepiscopi; or, on the other hand, substantive
authorities, sovereign within their own limits, though bound by external
ties to each other and to their Metropolitan. The most perfect state of
a Christian country would be, where there was a sufficient number of
separate dioceses; the next to it, where there were Chorepiscopi, or
Suffragan Bishops in the modern sense of the word.

Few persons, who have not expressly examined the
subject, are aware of the minuteness of the dioceses into which many
parts of Christendom were divided in the first ages. Some Churches in
Italy were more like our rural deaneries than what we now consider
dioceses; being not above ten or twelve miles in extent, and their sees
not above five or six miles from each other. Even now (or, at least, in
Bingham's time) the kingdom of Naples contains 147 sees, of which twenty
are Archbishoprics. Asia Minor is 630 miles long, 210 broad; yet in this
country there were almost 400 dioceses. Palestine is in length 160
miles, in breadth 120; yet the number of known dioceses amounted to 48.
Again, in the province of Syria Secunda, the see of Larissa (\emph{e.g}.)
was about 14 miles from Apamea, Arethusa 16 from Epiphania. And so
again, turning to the West, though the dioceses were generally larger,
as partaking more of a Missionary character, yet we shall find in
Ireland at one time from 50 to 60 sees.

Such was the character of the Primitive Regimen,
where Christianity especially flourished in the zeal and number of its
professors. But, where the country was mountainous or desert, the
inhabitants scanty, or but partially Christian, it was considered
advisable to leave all to the management of one chief Pastor, who
appointed assistants to himself according to his discretion, as the
circumstances of the times required. The office of these Chorepiscopi,
or country Bishops, was to preside over the country clergy, inquire into
their behaviour, and report to their principal; also to provide fit
persons for the inferior ministrations of the Church. They had the power
of ordaining the lower ranks of clergy, such as the readers,
sub-deacons, and exorcists; they might ordain priests and deacons with
the leave of the city Bishop, and administer the rite of confirmation;
and were permitted to sit and vote in councils. Thus their office bore a
considerable resemblance to that of our Archdeacons; except, of course,
that they had the power of ordination; whereas the latter are only
presbyters. And, in matter of fact, by such presbyters (\emph{visitors},
as they were called) they were superseded in the course of the fourth
and following centuries, till at length the Pope caused the order to be
set aside almost altogether in the ninth.

Little use was made of Suffragans during the middle
ages; but, at the time of our Reformation, Archbishop Cranmer felt the
deficiency of the English Church in respect of Bishoprics, and projected
several measures to supply it. The most complete was that of increasing
the number of dioceses; availing himself of existing circumstances, he
advised the King to apply the Abbey lands to the founding of twenty
additional sees. Bishop Burnet gives some of the particulars of this
attempt in the following passage:

``On the 23d of May, in the session of Parliament, a
bill was brought in by Cromwell for giving the King power to erect new
bishoprics by his letters-patent \footnote{It is scarcely necessary to observe, that
parliament was then the lay synod of the Church of England.}.
It was read that day for the first, second, and third time; and sent
down to the Commons. The preamble of it was, ‘that it was known what
slothful and ungodly life had been led by those who were called
religious. But that these houses might be converted to better uses; that
GOD'S
word might be better set forth; children brought up in learning; clerks
nourished in the universities; and that old decayed servants might have
livings; poor people might have almshouses to maintain them; readers of
Greek, Hebrew, and Latin, might have good stipends; daily alms might be
administered, and allowance might be made for mending of the highways,
and exhibitions for ministers of the Church; for these ends, if the King
thought fit to have more bishoprics or cathedral churches erected out of
the rents of these houses, full power was given him to erect and found
them, and to make rules and statutes for them, and such translations of
sees, or divisions of them, as he thought fit.' In the same paper, there
is a list of the sees which he intended to found; of which what was done
afterwards came so far short, that I know nothing to which it can be so
reasonably imputed, as the declining of Cranmer's interest at court, who
had proposed the erecting the new cathedrals and sees, with other things
mentioned in the preamble of the statute, as a great mean of reforming
the Church.'' \footnote{Burnet, Hist. Reform. iii.} Some of the
pro posed additional dioceses are then enumerated; Essex, Hertford,
Bedfordshire and Buckinghamshire, Oxford and Berkshire, Northampton and
Huntingdon, Middlesex, Leicester and Rutland, Gloucestershire,
Lancashire, Suffolk, Stafford and Salop, Nottingham and Derby, Cornwall.
As to the means by which they were to be endowed, no opinion is here
expressed on its lawfulness, as the present sketch is confined to the
consideration of the spiritual part of the ecclesiastical system. It is
scarcely necessary to add, that Cranmer's views were partly realized, in
the subsequent creation of the dioceses of Chester, Bristol, Gloucester,
Oxford, and Peterborough.

The same prelate, whose episcopate has had so
important an influence upon the constitution of our Church ever since,
also projected with great wisdom, a system of suffragan bishops or
Chorepiscopi, which he was able to bring into effect, and which lasted
till the reign of King James. Twenty-six such bishops were appointed;
the bishop of the diocese having the power of presenting two persons to
the king, who might choose either of them, and present them to the
archbishop of the province for consecration. These suffragans exercised
such jurisdiction as their principal gave them, or as had formerly been
committed to suffragans; their authority lasting no longer than he
continued their commission to them. ``These were believed,'' says Burnet \footnote{Hist. Reform. ii.}, ``to be the same with the Chorepiscopi in the primitive church
which, as they were begun before the first council of Nice, so they
continued in the Western Church till the 9th century, and then a
decretal of Damascus being forged, that condemned them, they were put
down every where by degrees, and now revived in England. The suffragan
sees were as follows: Thetford, Ipswich, Colchester, Dover, Guilford,
Southampton, Taunton, Shaftsbury, Molton, Marlborough, Bedford,
Leicester, Gloucester, Shrewsbury, Bristol, Penrith, Bridgwater,
Nottingham, Gantham, Hull, Huntingdon, Cambridge, Pereth, Berwick, St.
Germain's, and the Isle of Wight.''

After the disuse of suffragans in the reign of
James I. there was a fresh project for establishing them on the
Restoration. Charles, in one of his declarations, promises to increase
the number of bishops, in accordance with Archbishop Ussher's plan for
episcopal government. However, his intention was not put into execution, doubtless owing to existing circumstances, which reasonably
interfered with it.

The following extract is made from Bingham, Antiqu.
ix. 8: ``One great objection against the present diocesan episcopacy, and
that which to many may look the most plausible, is drawn from the vast
extent and greatness of most of the northern dioceses of the world,
which makes it so extremely difficult for one man to discharge all the
offices of the episcopal function ... The Church of England has usually
followed the larger model, and had very great and extensive dioceses;
for at first she had but seven bishoprics in the whole nation, and those
commensurate in a manner to the seven Saxon kingdoms. Since that time
she has thought it a point of wisdom to contract her dioceses, and
multiply them into above 20; and if she should think fit to add 40 or
100 more, she would not be without precedent in the practice of the
Primitive Church … In Ireland there are not now above half the number
of dioceses that there were before, and consequently they must needs be
larger by uniting them together. In England, there are more in number
than formerly, some new ones being created out of old ones, and at
present the whole number augmented to three times as many as they were
for some ages after the first conversion. Besides that, we have another
way of contracting dioceses in effect here in England appointed by law,
which law was never yet repealed; which is by devolving part of the
bishop's care upon the Chorepiscopi, or suffragan bishops, as the law
calls them:—a method commonly practised in the ancient Church in such
large dioceses as those of St. Basil and Theodoret, one of which had no
less than fifty Chorepiscopi under him, if Nazianzen rightly informs us.
And it is a practice which was continued here all the reign of Queen
Elizabeth, and even to the end of King James; and is what may be revived
again, whenever any bishop thinks his diocese too large, or his burden
too great to be sustained by himself alone.''

To the above statements, may be subjoined the
present number of souls, and the area of square miles, in certain of our
dioceses, as given in a pamphlet lately published, which has come into
the writer's hands since the foregoing was put on paper. (Vide Plan
for a New Arrangement, \&c. by Lord Henley.)

Souls
Square

Miles

Chester
1,806,722
4140

London
1,676,725
1942

York
1,526,288
5300

Lincoln
920,011
5775

Lichfield
978,655
3344

By this table, it is not here intended to insinuate
the necessity of any immediate measure of multiplying the English sees
or appointing suffragans, (the expediency of which is to be determined
by a variety of considerations, which it were unprofitable here to
detail,) but to show that the \emph{genius} of our ecclesiastical system
tends towards such an increase, and that the only question to be
determined is one of time. These statements are also made with the view
of keeping up in the minds of churchmen a recollection of the injury
which the Irish branch of our Church has lately sustained in the
diminution of its sees.

OXFORD,

\emph{The Feast of St. Philip and St. James}

P.S.—Since this was written, a new arrangement of
Dioceses and Sees has been made by authority of a Royal Commission,
composed of members the greater part of whom were laymen, and without
confirmation of their acts on the part of the Church.

[FOURTH EDITION.]

———————————————————————

These Tracts are continued in
Numbers, and sold at the price of 2d. for each sheet, or 7s. for 50
copies.

LONDON: PRINTED FOR
J. G. F. \& J. RIVINGTON,

ST. PAUL'S CHURCH YARD, AND WATERLOO PLACE.

1840.

\xchapter{Rites and Customs of the Church}

Tract No. 34 (Ad Scholas)

[\emph{Ho men oun pistos, h\={o}s chr\={e}, kai
errh\={o}menos oude deitai logou kai aitias, huper \={o}n an
epitachth\={e}i, all}'\emph{arkeitai t\={e}i paradosei mon\={e}i}].
\emph{Chrysost. in} 1 \emph{Cor. Hom}. 26.

He who is duly strengthened in faith, does not go
so far as to require argument and reason for what is enjoined, but is
satisfied with the \emph{tradition} alone.

———————

THE
reader of ecclesiastical history is sometimes surprised at finding
observances and customs generally received in the Church at an early
date, which have not express warrant in the Apostolic writings; \emph{e.g}.
the use of the cross in baptism. The following pages will be directed to
the consideration of this circumstance; with a view of suggesting from
those writings themselves, that a minute ritual was contemporaneous with
them, that the Apostles recognise it as existing anti binding, that it
was founded on religious \emph{principles}, and tended to the
inculcation of religious truth. Not that any formal proof is attainable
or conceivable, considering the brevity and subjects of the inspired
documents; but such fair evidence of the fact, as may recommend it to
the belief of the earnest and single-minded Christian. It is abundantly
evident that the Epistles were not written to prescribe and enforce the
Ritual of religion; all then we can expect, if it existed in the days of
the Apostles, is an occasional allusion to it in their Epistles as
existing, and a plain acquiescence in it: and thus much we find.

Let us consider that remarkable passage, (1 Cor. xi.
2-16.) which, I am persuaded, most readers pass over as if they could
get little instruction from it. St. Paul is therein blaming the
Corinthians for not adhering to the \emph{custom} of the Church, which
prescribed that men should wear their hair short, and that women should
have their head covered during divine service; a custom apparently
most unimportant, if any one ever was, but in his view strictly binding
on Christians. He begins by implying that it is one out of many rules or
traditions ([\emph{paradoseis}]) which he had given them, and they were
bound to keep. He ends by refusing to argue with any one who obstinately
cavils at it and rejects it: ``If any man seem to be contentious, we
have no such custom, neither the Churches of GOD.''
Here then at once a view is opened to us which is quite sufficient to
remove the surprise we might otherwise feel at the multitude of rites,
which were in use in the Primitive Church, but about which the New
Testament is silent; and further, to command our obedience to such as
come down to us from the first ages, and are agreeable to Scripture.

In accordance with this conclusion, is the clear
and forcible command given by the Apostle, (2 Thess. ii. 15.) ``Brethren, stand fast, and hold the traditions, which ye have been
taught, whether \emph{by word}, or our epistle.''

To return. St. Paul goes on to give the reason of
the usage, for the satisfaction of the weak brethren at Corinth. It was,
he implied, a symbol or development (so to say) of the principle of the
subordination of the woman to the man, and a memorial of the history of
our creation; nay it was founded in ``\emph{nature},'' \emph{i.e}.
natural reason. And lastly, it had a practical object: the woman ought
to have her head covered ``\emph{because} of the angels.'' We need not
stop to inquire \emph{what} this reason was; but it was a reason of a
practical nature which the Corinthians understood, though we may not. If
it mean, as is probable, ``because she is in the sight of the heavenly
angels,'' (1 Tim. v. 21.) it gives a still greater importance to the
ceremonies of worship, as connecting them with the unseen world.

It would seem indeed as if the very multiplicity of
the details of the Church ritual made it plainly impossible for St. Paul
to write them all down, or to do more than \emph{remind} the Corinthians
of his way of conducting religious discipline when he was among them. ``Be ye \emph{followers} of me;'' he says,
``I praise you that \emph{ye
remember me} in all things.'' It is evident there are ten thousand
little points in the working of any large system, which a present
instructor alone can settle. Hence it is customary at present, when
a school is set up, or when any novel manufacture in trade, or
extraordinary machinery, is to be brought into use, to set it going by
sending a person fully skilled in its practical details. Such was St.
Paul as regards the system of Christian discipline and worship; and when
he could not go himself, he sent Timothy in his place. He says in the
4th chapter: ``I beseech you, be ye followers of me. \emph{For this cause}
have I sent unto you Timotheus, who shall bring you into remembrance \emph{of
my ways which be in} CHRIST,
\emph{as I teach every where in every Church}.'' Here there is a like
reference to an uniform system of discipline,—whether as to Christian
conduct, worship, or Church government.

Another important allusion appears to be contained
in the 22d verse of the chapter above commented on. ``What, have ye not
houses to eat and drink in? or despise ye the \emph{Church of} GOD?'' This is remarkable as being a solitary
allusion in Scripture to \emph{houses} of prayer under the Christian
System, which nevertheless we know from \emph{ecclesiastical history}
were used from the very first. Here then is a most solemn ordinance of
primitive Christianity, which barely escapes, if it escapes, omission in
Scripture.

A passing allusion is made in another passage of
the same Epistle, to the use of the word Amen at the conclusion of the
Eucharistical prayer, as it is preserved after it and all other prayers
to this day. Thus the ritual of the Apostles descended to minutiæ, and
these so invariable in their use, as to allow of an appeal to them.

In the original institution of the Eucharist, as
recorded in the Gospels, there is no mention of \emph{consecrating} the
cup; but in 1 Cor. x. 16, St. Paul calls it ``the cup \emph{of blessing
which we bless}.'' This incidental information, vouchsafed to us in
Scripture, should lead us to be very cautious how we put aside other
usages of the early Church concerning this Sacrament, which do not
happen to be \emph{clearly} mentioned in Scripture; as \emph{e.g}. the
solemn offering of the elements to GOD by way of pleading His mercy through CHRIST,
which seems to have been universal in the early Church

As regards the same Sacrament, let us consider the
use of the word [\emph{leitourgount\={o}n}], \emph{ministering} (Acts
xiii. 2.); a word which, dropt (so to say) by accident, and
interpreted, as is reasonable, by its use in the services of the Jewish
Law, (Luke i. 23; Heb. x. 11.) remarkably coincides with the [\emph{leitourgia}]
of the Primitive Church, according to which the offering of the Altar
was intercessory, as pleading CHRIST’S
merits before the throne of grace.

Again, in 1 Cor. xv. 29, we incidentally discover
the existence of persons who are styled ``the baptized for the dead.''
Perhaps it is impossible to determine what is meant by this phrase, on
which little light is thrown by early writers. However, any how it seems
to refer to a \emph{custom} of the Church, which was so usual as to
admit of an appeal to it, which St. Paul approved, yet which he did not
in the Epistle directly enforce, and but casually mentions.

In 1 Cor. i. 16, St. Paul happens to inform us that
he baptized the \emph{household} of Stephanas. It has pleased the HOLY
SPIRIT
to preserve to us this fact; by which is detected the existence of a
rule of discipline for which the express doctrinal parts of Scripture
afford but indirect warrant, viz. the custom of household baptism. (Vid.
also Acts xvi. 15, 33.) This accidental disclosure accurately
anticipates the after practice of the early Church, according to which
families, infants included, were baptized, and that on a weighty
doctrinal \emph{reason}; viz. that all men were born in sin and in the
wrath of GOD,
and needed to be individually translated into that kingdom of grace,
into which baptism is the initiation.

These instances, then, not to notice others of a
like or a different kind, are surely sufficient to reconcile us to the
complete ritual system which breaks upon us in the writings of the
Fathers. If any parts of it indeed are contrary to Scripture, that is of
course a decisive reason at once for believing them to be additions and
corruptions of the original ceremonial; but till this is shown, we are
bound to venerate what is certainly primitive, and probably is
apostolic.

It will be remarked, moreover, that many of the
religious observances of the early Church are expressly built upon words
of Scripture, and intended to be a visible memorial of them, after the
manner of St. Paul’s directions about the respective habits of men and
women, which was just now noticed. Metaphorical or mystical
descriptions were represented by a corresponding literal action. Our LORD Himself authorised this procedure when He took up the
metaphor of the prophets concerning the fountain opened for our
cleansing (Zech. xiii. 1.) and represented it in the visible rite of
baptism. Accordingly, from the frequent mention of \emph{oil} in
Scripture as the emblem of spiritual gifts, (Is. lxi. 1-3, \&c.) it
was actually used in the Primitive Church in the ceremony of admitting
catechumens, and in baptizing. And here again they had the precedent of
the Apostles, who applied it in effecting their miraculous cures. (Mark
vi. 13. James v. 14.) And so from the figurative mention in Scripture of
\emph{salt}, as the necessary preparation of every religious sacrifice,
it was in use in the Western Church, in the ceremony of admitting
converts into the rank of catechumens. So again from Phil. ii. 10, it
was customary to bow the head at the name of JESUS. It were endless to multiply instances of
a similar pious attention to the very words of Scripture, as their
custom of continual public prayer from such passages as Luke xviii. 7;
or of burying the bodies of martyrs under the altar, from Rev. vi. 9; or
of the white vestments of the officiating ministers, from Rev. iv. 4.

Two passages on the subject from the Fathers shall
now be laid before the reader, by way of further illustration, and first
from Tertullian:

``Though
this observance has not been determined by any Scripture, yet it is
established by custom, which doubtless is derived from tradition. For
how can an usage ever obtain, which has not first been given by
tradition? But you say, even though tradition can be produced, still a
written (Scripture) authority must be demanded. Let us examine, then,
how far it is true, that a tradition itself, unless written in
Scripture, is inadmissible. Now I will give up the point at once, if it
is not already determined by instances of other observances, which are
maintained without any Scripture proof, on the mere plea of tradition,
and the sanction of consequent custom. To begin with baptism. Before we
enter the water, we solemnly renounce the devil, his pomp, and his
angels, in church in the presence of the Bishop. Then we are plunged in
the water thrice, and answer certain questions over and above what the LORD has determined in the written gospel. After coming out of it, we
taste a mixture of milk and honey; and for a whole week from that day we
abstain from our daily bath. The sacrament of the Eucharist, though
given by the LORD to all and at supper time, yet is celebrated in our
meetings before day break, and only at the
hand of our presiding ministers … We sign our forehead with the cross
whenever we set out and walk, go in or out, dress, gird on our sandals,
bathe, eat, light our lamps, sit or lie down to rest, whatever we do. If
you demand a scripture rule for these and such like observances, we can
give you none; all we say to you is, that tradition directs, usage
sanctions, faith obeys. That reason justifies this tradition, usage, and
faith, you will soon yourself see, or will easily learn from others;
meanwhile you will do well to believe that there is a law to which
obedience is due. I add one instance from the old dispensation. It is so
usual among the Jewish females to veil their head, that they are even
known by it. I ask where the law is to be found; the Apostle’s
decision of course is not to the point. Now if I no where find a law, it
follows that tradition introduced the custom, which afterwards was
confirmed by the Apostle when he explained the reason of it. These
instances are enough to show that a tradition, even though not in
Scripture, still binds our conduct, if a continuous usage be preserved
as the witness of it.''—Tertullian, de Coron. § 3.

Upon this passage it may be observed, that
Tertullian, flourishing A.D.
200, is on the one hand a very early witness for the existence of the
general doctrine which it contains, while on the other he gives no
sanction to those later customs, which the Church of Rome upholds, but
which cannot be clearly traced to primitive times.

St. Basil, whose work on the HOLY SPIRIT, § 66, shall next be cited, flourished
in the middle of the fourth century, 150 years after Tertullian, and was
of a very different school; yet he will be found to be in exact
agreement with him on the subject before us, viz. that the ritual of the
Church was derived from the Apostles, and was based on religious
principles and doctrines. He adds a reason for its not being given us in
Scripture, which we may receive or reject as our judgment leads us, viz.
that the rites were memorials of doctrines not intended for publication
except among baptized Christians, whereas the Scriptures were open to
all men. This at least is clear, that the ritual could scarcely have
been given in detail in Scripture, without imparting to the Gospel the
character of a burdensome ceremonial, and withdrawing our attention from
its doctrines and precepts.

``Of
those articles of doctrine and preaching, which are in the custody of
the Church, some come to us in Scripture itself, some are conveyed to us
by a continuous tradition in mystical depositories. Both have equal
claims on our devotion, and are received by all, at least by all who are
in any respect
Churchmen. For, should we attempt to supersede the usages which are not
enjoined in Scripture as if unimportant, we should do most serious
injury to Evangelical truth: nay, reduce it to a bare name. To take an
obvious instance; which Apostle has taught us in Scripture to sign
believers with the cross? Where does Scripture teach us to turn to the
east in prayer? Which of the saints has left us recorded in Scripture
the words of invocation at the consecration of the bread of the
Eucharist, and of the cup of blessing? Thus we are not content with what
Apostle or Evangelist has left on record, but we add other rites before
and after it, as important to the celebration of the mystery, receiving
them from a teaching distinct from Scripture. Moreover, we bless the
water of baptism, and the oil for anointing, and also the candidate for
baptism himself … After the example of Moses, the Apostles and Fathers
who modelled the Churches, were accustomed to lodge their sacred
doctrine in mystic forms, as being secretly and silently conveyed …
This is the reason why there is a tradition of observances independent
of Scripture, lest doctrines, being exposed to the world, should be so
familiar as to be despised … We stand instead of kneeling at prayer on
the Sunday; but all of us do not know the reason of this … Again,
every time we kneel down and rise up, we show by our outward action that
sin has levelled us with the ground, and the loving mercy of our Creator
has recalled us to heaven.''

The conclusion to be drawn from all that has been
said in these pages is this:—That rites and ordinances, far from being
unmeaning, are in their nature capable of impressing our memories and
imaginations with the great revealed verities; far from being
superstitious, are expressly sanctioned in Scripture as to their
principle, and delivered to the Church in their form by tradition.
Further, that they varied in different countries, according to the
respective founder of the Church in each. Thus \emph{e.g}., St. John and
St. Philip are known to have adopted the Jewish rule for observing
Easter-day; while other Apostles celebrated it always on a Sunday.
Lastly, that, although the details of the early ritual varied in
importance, and corrupt additions were made in the middle ages, yet
that, as a whole, the Catholic ritual was a precious possession; and if
we, who have rid ourselves of those corruptions, have lost not only the
possession, but the sense of its value, it is a serious question whether
we are not like men who recover from some grievous illness with the loss
or injury of their sight or hearing;—whether we are not like the Jews
returned from captivity, who could never find the rod of Aaron or
the Ark of the Covenant, which, indeed, had ever been hid from the
world, but then was removed from the Temple itself.

OXFORD,

\emph{The Feast of St. Philip and St. James}

[NEW EDITION.]

———————————————————————

These Tracts are continued in
Numbers, and sold at the price of 2d. for each sheet, or 7s. for 50
copies.

LONDON: PRINTED FOR
J. G. F. \& J. RIVINGTON,

ST. PAUL’S CHURCH YARD, AND WATERLOO PLACE.

1840.

\xchapter{The Grounds of our Faith}

Tract No. 45 (Ad Clerum)

EVERY
system of theology has its dangers, its tendencies towards evil. Systems
short of the truth have this tendency inherent in themselves, and in
process of time discover it, and work out the anticipated evil, which is
but the legitimate though latent consequence of their principles. Thus,
we may consider the present state of Geneva the fair result on the long
run of the system of self-will which was established there in the
sixteenth century. But even the one true system of religion has its
dangers on all sides, from the weakness of its recipients, who pervert
it. Thus the Holy Catholic doctrines, in which the Church was set up,
were corrupted into Popery, not legitimately, or necessarily, but by
various external causes acting on human corruption, in the lapse of many
ages. St. Paul's command of obedience to rulers, was changed into the
tyrannical rule of one Bishop over all countries; his recommendation of
an unmarried life, for certain religious objects, was made a rule of
celibacy in the case of the clergy. Now let us ask, what are the bad
tendencies of Protestantism? for this is a question which nearly
concerns ourselves. We are nearly 300 years from its rise in this
country; have any evils yet shown themselves from it? It is not here
proposed to examine the question at large; but a hint on one part of the
subject may be made in answer to it.

At the Reformation, the authority of the Church was
discarded by the spirit then predominant among Protestants, and
Scripture was considered as the sole document both for ascertaining and
proving our faith. The question immediately arose, ``Is this or that
doctrine in Scripture?''—and in consequence, various intellectual
gifts, such as argumentative subtilty, critical acumen, knowledge of the
languages, rose into importance, and became the interpreters of
Christian truth. Exposition lay through controversy. Now the natural
effect of disputation is to make us shun all but the strongest proofs,
those which an adversary will find substantial impediments in his line
of reasoning; and, therefore, to generate a cautious discriminative turn
of thought, to fix in the mind a \emph{standard} of proof simulating
demonstration, and to make light of mere probabilities. This
intellectual habit, resulting from controversy, would also arise from
the peculiar exercises of thought necessary for the accurate scholar or
antiquarian. It followed, that in course of time, all the delicate
shades of truth and falsehood, the unobtrusive indications of GOD'S will, the low tones of the ``still small voice,'' in which
Scripture abounds, were rudely rejected; the crumbs from the rich man's
table, which Faith eagerly looks about for, were despised by the
proud-hearted intellectualist, who (as if it were a favour in him to
accept the Gospel,) would be content with nothing short of certainty,
and ridiculed as superstitious and illogical whatever did not approve
itself to his own cold, hard, and unimpassioned temper. For instance, if
the cases of Lydia, of the jailer, of Stephanas, were brought to show
our LORD'S \emph{wish} as to the baptism of
households, the actions of his apostles to \emph{interpret} his own
commands, it was answered; ``This is no satisfactory \emph{proof}; it is
not \emph{certain} that every one of those households was not himself a
believer; it is not \emph{certain} there were any children among them:''—though
surely, in as many as \emph{three} households, the probability is on the
side which the Church has taken, especially viewing the texts in
connexion with our SAVIOUR'S words, ``Suffer the little
children,'' \&c. Again, while the observance of the LORD'S day was grounded upon the \emph{practice}
of the apostles, it was somehow felt, that this proof was not \emph{strong
enough} to bind the mass of Protestants: and so the chief argument
now in use is one drawn from the Jewish law, viz. the direct Scripture \emph{command},
contained in the fourth commandment.

Our S\textsc{aviour} has noticed the
frame of mind here alluded to, in Mark viii. 11, 12, where His feelings
and judgment upon it are also told us:—``And the Pharisees came forth,
and began to question with Him, seeking of Him \emph{a sign from heaven},
tempting Him. And \emph{He sighed deeply in His spirit}, and saith, Why
doth this generation seek after a sign? Verily I say unto you, \emph{There
shall no sign be given unto this generation}. And He left them.''

We are warned against the same hard, intractable
temper in the book of Psalms:—``I will inform thee, and teach thee in
the way wherein thou shalt go: \emph{and I will guide thee with Mine eye}.
Be ye not like to horse and mule, \emph{which have no understanding};
whose mouth must be held \emph{with bit and bridle}, lest they fall upon
thee.'' Ps. xxxii. 9, 10. This stubborn spirit, which yields to nothing
but violence, is determined to \emph{feel} CHRIST'S
yoke ere it submits to it, will not see except in broad day-light, and
like the servant who hid his talent, is ever making excuses, murmuring,
doubting, grudging obedience, and stifling docile and open-hearted
faith, is the spirit of ultra-Protestantism, \emph{i.e}. that spirit, to
which the principles of Protestantism \emph{tend}, and which they have
in a great measure \emph{realized}. On this subject the reader may
consult Nos. 4, 8, and 19, of this series of Tracts.

Now to apply this to the doctrines, at present so
much under-valued, which it is the especial object of these Tracts to
enforce.

When a clergyman has spoken strongly in defence of
Episcopacy, a hearer will go away saying, that there is much very able
and forcible, much very eloquent and excellent, in what he has just
heard; but after all, \emph{there is very little about Episcopacy in
Scripture}. This is the point to which a shrewd, clear-headed
reasoner will resort,—``\emph{after all};'' we come round and round to
it; the doctrine advocated is plausible, useful, generally received
hitherto;—granted, \emph{but} Scripture says very little about it.

Now it cannot be for a moment allowed, that
Scripture \emph{contains} little on the subject of Church Government;
though it may readily be granted that it \emph{obtrudes} on the reader
little about it. The doctrine is in it, not on it; not on the surface.
This need not be proved here, since the subject has been variously
considered in former Numbers of this series. But it may be useful in a few words to show how the state of the argument and controversy
concerning Episcopacy, illustrates the above remarks, and how parallel
it is to the state in which other religious truths are found, which no
Churchman ventures to dispute.

1. Now, in the first place, let us suppose, \emph{for
the sake of argument}, that Episcopacy is in fact not at all
mentioned in Scripture: even then it would be our duty to receive it.
Why? because the first Christians received it. If we wish to get at the
truth, no matter how we get at it, \emph{if} we get at it. If it be a
fact, that the earliest Christian communities were universally episcopal,
it is a reason for our maintaining Episcopacy; and \emph{in proportion}
to our conviction, is it incumbent on us to maintain it.

Nor can it be fairly dismissed as a non-essential,
or ordinance indifferent and mutable, though formerly existing over
Christendom; for, \emph{who} made us judges of essentials and
non-essentials? \emph{how} do we determine them? In the Jewish law, the
slightest transgression of the commandment was followed by the penalty
of death; vide Lev. viii. 35; x. 6. Does not its universality imply a
necessary connexion with Christian doctrine? Consider how such
reasonings would carry us through life; how the business of the world
depends on punctuality in minutes; how ``great a matter'' a mere spark
dropped on gunpowder ``kindleth.''

But, it may be urged, that we Protestants believe
the \emph{Scriptures} to contain the whole rule of duty.—Certainly
not; they constitute a rule of \emph{faith}, not a rule of \emph{practice};
a rule of \emph{doctrine}, not a rule of \emph{conduct} or \emph{discipline}.
Where (\emph{e.g}.) are we told in Scripture, that gambling is wrong? or
again, suicide? Our Article is precise: ``Holy Scripture containeth all
things necessary to salvation, so that whatsoever is not read therein,
\&c. it is not to be required of any man, that it should be \emph{believed}
as an article of \emph{faith}.'' Again it says, that the Apocrypha is \emph{not}
to be applied ``to establish any \emph{doctrine},'' implying that this is
the use of the canonical books.

2. However, let us pass from this argument, which
is but founded on a \emph{supposition}, that Episcopacy is not enjoined
in Scripture. Suppose we maintain, as we may well maintain, that it is
enjoined in Scripture. An objector will say, that at all events it is but obscurely contained therein, and cannot be drawn out from it
without a great deal of delicate care and skill. Here comes in the
operation of that principle of \emph{faith} in opposition to \emph{criticism},
which was above explained; the principle of being content with a little
light, where we cannot obtain sunshine. If it is \emph{probably}
pleasing to CHRIST,
let us maintain it. Now take a parallel case: \emph{e.g}. the practice
of infant baptism; where is this \emph{enjoined} in Scripture? No where.
Why do we observe it? Because the primitive Church observed it, and
because the Apostles in Scripture \emph{appear} to have sanctioned it,
though this is not altogether certain \emph{from} Scripture. In a
difficult case we do as well as we can, and carefully \emph{study} what
is most agreeable to our LORD and SAVIOUR. This is how our Church expresses it in
the xxviith Article: ``The baptism of young children is in any wise to be
retained in the Church, as \emph{most agreeable} with the institution of
CHRIST.''
This is true wariness and Christian caution; very different from that
spurious caution which ultra-Protestantism exercises. Let a man only be
consistent, and apply the same judgment in the case of Episcopacy: let
him consider whether the duty of keeping to Bishops, be not ``\emph{most
agreeable} with the institution of CHRIST.''
If, indeed, he denies this altogether, these remarks do not apply; but
they are addressed to waverers, and falsely moderate men, who cannot
deny, that the evidence of Scripture is in favour of Churchmen, but say
it is not strong enough. They say, that if Almighty GOD
had intended an uniformity in Church Government among Christians, He
would have spoken more clearly.

Now if they carried on this line of argument
consistently, they would not baptize their children: happily they are
inconsistent. It would be more happy still, were they consistent on the
other side; and, as they baptize their children, because it is safer to
observe than to omit the sacrament, did they also keep to the Church, as
the safer side. The received practice, then, of infant baptism seems a
final answer to all who quarrel with the Scripture evidence for
Episcopacy.

3. But further still, infant baptism, like
Episcopacy, is but a case of \emph{discipline}. What shall we say, when
we consider that a case of \emph{doctrine}, necessary doctrine,
doctrine the very highest and most sacred, may be produced, where the
argument lies as little on the surface of Scripture,—where the proof,
though \emph{most conclusive}, is as indirect and circuitous as that for
Episcopacy; viz. the doctrine of the Trinity? Where is this solemn and
comfortable mystery formally stated in Scripture, as we find it in the
creeds? Why is it not? Let a man consider whether all the objections
which he urges against the Scripture argument for Episcopacy may not be
turned against his own belief in the Trinity. It is a happy thing for
themselves that men are inconsistent; yet it is miserable to advocate
and establish a \emph{principle}, which, not in their own case indeed,
but in the case of others who learn it of them, leads to Socinianism.
This being considered, can we any longer wonder at the awful fact, that
the descendants of Calvin, the first Presbyterians, are at the present
day in the number of those who have denied the LORD
who bought them?

OXFORD,

\emph{The Feast of St. Luke}

[FOURTH EDITION.]

———————————————————————

These Tracts are continued in
Numbers, and sold at the price of 2d. for each sheet, or 7s. for 50
copies.

LONDON: PRINTED FOR
J. G. F. \& J. RIVINGTON,

ST. PAUL'S CHURCH YARD, AND WATERLOO PLACE.

1840.

\xchapter{Advertisement}

IN completing the second volume of a publication, to which the
circumstances of the day have given rise, it may be right to allude to a
change which has taken place in them since the date of its commencement.
At that time, in consequence of long security, the attention of members
of our Church had been but partially engaged in ascertaining the grounds
of their adherence to it; but the imminent peril to which all that is
dear to them has since been exposed, has naturally turned their thoughts
that way, and obliged them to defend it on one or other of the
principles which are usually put forward on its behalf. Discussions have
thus been renewed in various quarters, on points which had long remained
undisturbed; and, though numbers continue undecided in opinion, or take
up a temporary position in some one of the hundred middle points which
may be assumed between the two main theories in which the question
issues, and others, again, have deliberately entrenched themselves in
the modern or ultra-protestant alternative, yet, on the whole, there has
been much hearty and intelligent adoption, and much respectful study, of
those more primitive views maintained by our great Divines. As the
altered state of public information and opinion has a necessary bearing
on the efforts of those who desire to excite attention to the subject,
(in which number the writers of these Tracts are to be included,) it
will not be inappropriate briefly to state in this place, what it is
conceived is the present position of the great body of Churchmen with
reference to it.

While we have cause to be thankful for the sounder
and more accurate language which is now very generally adopted among
well-judging men on ecclesiastical subjects, we must beware of
over-estimating what has been done, and so becoming sanguine in our
hopes of success, or slackening our exertions to secure it. Many more
persons, doubtless, have taken up a profession of the main doctrine in
question, that, namely, of the One Catholic and Apostolic Church, than
fully enter into it. This is to be expected, it being the peculiarity of
all religious teaching, that words are imparted before ideas. A child
learns his Creed or Catechism before he understands it; and in beginning
any deep subject we are all but children to the end of our lives. The
instinctive perception of a rightly instructed mind, the \emph{primâ facie}
force of the argument, or the authority of our celebrated writers, have
all had their due and extensive influence in furthering the reception of
the doctrine, when once it was openly maintained; to which must be added
the prospect of the loss of state protection, which made it necessary to
look out for other reasons for adherence to the Church besides that of
obedience to the civil magistrate. Nothing, which has spread quickly,
has been received thoroughly. Doubtless there are a number of
seriously-minded persons, who think they admit the doctrine in question
much more fully than they do, and who would be startled at seeing that
realized in particulars, which they confess in an abstract form. Many
there are who do not at all feel that it is capable of a practical
application: and, while they bring it forward on special occasions, in
formal expositions of faith, or in answer to a direct interrogatory, let
it slip from their minds almost entirely in their daily conduct or their
religious teaching, from the long and inveterate habit of thinking and
acting without it. We must not then at all be surprised at finding, that
to modify the principles and motives on which men act is not the work of
a day; nor at undergoing disappointments, at witnessing relapses,
misconceptions, sudden disgusts, and, on the other hand, abuses and
perversions of the true doctrine, in the case of those who have taken it
up with greater warmth than discernment.

And in the next place, it will be found, that much
more has been done in awakening Churchmen to the truth of the Apostolical Commission as a fact, and to the admission of it as a duty,
than to the enjoyment of it as a privilege. If asked what is the use of
adhering to the Church, they will commonly answer, that it is commanded,
that all acts of obedience meet with their reward from Almighty God, and
this in the number; but the notion of the Church as the storehouse and
direct channel of grace, as a Divine Ordinance, not merely to be
maintained for order's sake, or because schism is a sin, but to be
approached joyfully and expectantly as a definite instrument, or rather
the appointed means, of spiritual blessings,—as an Ordinance which
conveys secret strength and life to every one who shares in it, unless
there be some actual moral impediment in his own mind,—this is a
doctrine which as yet is but faintly understood among us. Nay, our
subtle Enemy has so contrived, that by affixing to this blessed truth
the stigma of Popery, numbers among us are effectually deterred from
profiting by a gracious provision, intended for the comfort of our
faith, but in their case wasted.

The particular deficiency here alluded to may also
be described by referring to another form under which it shows itself,
viz. the \emph{à priori} reluctance in those who believe the
Apostolical Commission, to appropriate to it the power of consecrating
the Lord's Supper; as if there were some antecedent improbability in GOD'S gifts being lodged in particular observances, and
distributed in a particular way; and as if the strong wish, or moral
worth, of the individual could create in the outward ceremony a virtue
which it had not received from above. Rationalistic, or (as they may be
more properly called) carnal notions concerning the Sacraments, and, on
the other hand, a superstitious apprehension of resting in them, and a
slowness to believe the possibility of GOD'S
having literally blessed ordinances with invisible power, have, alas!
infected a large mass of men in our communion. There are those whose ``word
will eat as doth a canker;'' and it is to be feared, that we have been
over-near certain celebrated Protestant teachers, Puritan or
Latitudinarian, and have suffered in consequence. Hence we have almost embraced the doctrine, that GOD
conveys grace only through the instrumentality of the mental energies,
that is, through faith, prayer, active spiritual contemplations, or
(what is commonly called) communion with GOD,
in contradiction to the primitive view, according to which the Church
and her Sacraments are the ordained and direct visible means of
conveying to the soul what is in itself supernatural and unseen. For
example, would not most men maintain, on the first view of the subject,
that to administer the LORD'S
Supper to infants, or to the dying and apparently insensible, however
consistently pious and believing in their past lives, must be, under all
circumstances, and in every conceivable case, a superstition? and yet
neither practice is without the sanction of primitive usage. And does
not this account for the prevailing indisposition to admit that Baptism
conveys regeneration? Indeed, this may even be set down as the essence
of Sectarian Doctrine, (however its mischief may be restrained or
compensated, in the case of individuals,) to consider faith, and not the
Sacraments, as the proper instrument of justification and other gospel
gifts; instead of holding, that the grace of CHRIST comes to us altogether from without, (as
from Him, so through externals of His ordaining,) faith being but the \emph{sine
quâ non}, the necessary condition on our parts for duly receiving
it.

It has been with the view of meeting this cardinal
deficiency (as it may be termed) in the religion of the day, that the
Tract on Baptism, contained in the latter part of this volume, has been
inserted; which is to be regarded, not as an inquiry into one single or
isolated doctrine, but as a delineation, and serious examination of a
modern system of theology, of extensive popularity and great
speciousness, in its elementary and characteristic principles.

O\textsc{xford},

\emph{The Feast of All Saints}, 1835

\xchapter{The Visible Church Letter IV.}

Tract No. 47 (Ad Clerum)

I AM
sorry my delay has been so considerable in answering your remarks on my
Letters on the Church. Indeed it has been ungrateful in me, for you have
given me an attention unusual with the multitude of religious persons;
who, instead of receiving the arguments of others in simplicity, and
candour, seem to have a certain number of types, or measures of
professing Christians, set up in their minds, to one or other of which
they consider every one they meet with belongs, and who, accordingly,
directly they hear an opinion advanced, begin to consider whether the
speaker be a No. 1, 2, or 3, and having rapidly determined this, treat
his views with consideration or disregard, as it may be. I am far from
saying our knowledge of a person's character and principles should not
influence our judgment of his arguments; certainly it should have great
weight. I consider the cry ``measures not men,'' to be one of the many
mistakes of the day. At the same time there is surely is contrary
extreme, the fault of fancying we can easily look through men, and \emph{understand}
what each individual is; an arbitrary classing of the whole Christian
family under but two or three \emph{countenances}, and mistaking one man's
doctrine for another's. You at least have not called me am Arminian, or
a high Churchman, or a Borderer, or one of this or that school, and so
dismissed me.

To pass from this subject. You tell me that in my
zeal in advocating the doctrine of the Church Catholic and Apostolic, I
``use expressions and make assumptions which imply that the
Dissenters are without the pale of salvation.'' So let me explain myself
on these points.

You say that my doctrine of the one Catholic Church
in effect excludes Dissenters, nay, Presbyterians, from salvation. Far
from it. Do not think of me as of one who makes theories for himself in
his closet, who governs himself by book-maxims, and who, as being
secluded from the world, has no temptation to let his sympathies for
individuals rise against his abstract positions, and can afford to be
hard-hearted, and to condemn by wholesale the multitudes in various
sects and parties whom he never saw. I have known those among
Presbyterians whose piety, resignation, cheerfulness, and affection,
under trying circumstances, have been such, as to make me say to myself,
on the thoughts of my own higher privileges, ``Woe unto thee Chorazin,
woe unto thee Bethsaida!'' Where little is given, little will be
required; and that return, though little, has its own peculiar
loveliness, as an acceptable sacrifice to Him who singled out for praise
the widow's two mites. Was not Israel apostate from the days of
Jeroboam; yet were there not even in the reign of Ahab, seven thousand
souls who were ``reserved,'' an elect remnant? Does any Churchman wish to
place the Presbyterians, where, as in Scotland, their form of
Christianity is in occupation, in a worse condition under the Gospel
than Ephraim held under the Law? Had not the ten tribes the schools of
the Prophets, and has not Scotland at least the word of GOD? Yet what would be thought of the Jew who had maintained that
Jeroboam and his kingdom were in no guilt? and shall we, from a false
charity, from a fear of condemning the elect seven thousand, scruple to
say that Presbyterianism has severed itself from our temple privileges,
and undervalue the line of Levi and the house of Aaron? Consider our SAVIOUR'S discourse with the woman of
Samaria. While by conversing with her He tacitly condemned the Jews'
conduct in refusing to hold intercourse with the Samaritans, yet He
plainly declared that ``salvation was of the Jews.'' ``Ye worship ye know
not what;'' He says, ``we know what we worship.'' Can we conceive His
making light of the differences between Jew and Samaritan?

Further, if to whom much is given, of him much will
be required, how is it safe for us to make light of our privileges, if
we have them? is not this to reject the birth-right? to hide our talent
under a napkin? When we say that GOD
has done more for us than for the Presbyterians, this indeed \emph{may}
be connected with feelings of spiritual pride; but it need not. We may,
by so saying, provoke ourselves to jealousy; for we dare not deny that,
in spite of our peculiar privileges of communion with CHRIST, yet even higher saints \emph{may} lie
hid (to our great shame) among those who have not themselves the
certainty of our especial approaches to His glorious majesty. Was not
Elijah sent to a widow of Sarepta? did not Elisha cure Naaman? and are
not these instances set forward by our LORD
Himself as warnings to us ``not to be high-minded, but to fear;'' and,
again, as a gracious consolation when we think of our less favoured
brethren? Where is the narrowness of view and feeling which you impute
to me? Why may I not speak out, in order at once to admonish myself, and
to attempt to reclaim to a more excellent way those who are at present
severed from the true Church.

And what has here been said of an established
Presbyterianism, is true (in its degree) of dissent, when it has become
hereditary, and embodied in institutions.

Further, it is surely parallel with the order of
Divine Providence that there should be a variety, a sort of graduated
scale, in His method of dispensing His favour in CHRIST. So far from its being a strange thing that Protestant
sects are not ``in CHRIST,''
in the same fulness that we are, it is more accordant to the scheme of
the world that they should be between us and heathenism. It would be
strange if there were but two states, one absolutely of favour, one of
disfavour. Take the world at large, one form of paganism is better than
another. The North American Indians are theists, and as such more
privileged thorn polytheists. Mahometanism is a better religion than
Hindooism. Judaism is better than Mahometanism. One may believe that
long established dissent affords to such as are born and bred in it a
sort of pretext, and is attended with a portion of blessing, (where
there is no means of knowing better,) which does not attach to those who
\emph{cause} divisions, found sects, or wantonly wander from the
Church to the Meeting House;—that what is called an orthodox sect has
a share of Divine favour, which is utterly withheld from heresy. I am
not speaking of the next world, where we shall all find ourselves as
individuals, and where there will be but two states, but of existing
bodies or societies. On the other hand, why should the corruptions of
Rome lead us to deny her Divine privileges, when even the idolatry of
Judah did not forfeit hers, annul her temple-sacrifice, or level her to
Israel?

I say all this, merely, for the purpose of
suggesting to those who are ``weak,'' some idea of possible modes in which
Eternal Wisdom may reconcile the exuberance of his mercy in CHRIST to the whole race of man, with the
placing of it in its fulness in a certain ordained society and ministry.
For myself I prefer to rely upon the simple word of truth, of which
Scripture is the depository, and since CHRIST has told me to preach the whole counsel of GOD,
to do so fearlessly and without doubting; not being careful to find ways
of smoothing strange appearances in His counsels, and of obviating
difficulties, being aware on the one hand that His thoughts are not our
thoughts, nor our ways His ways, and on the other, that He is ever
justified in His sayings, and overcomes when He is judged.

OXFORD,

\emph{The Feast of All Saints}

[NEW EDITION.]

———————————————————————

These Tracts are continued in
Numbers, and sold at the price of 2d. for each sheet, or 7s. for 50
copies.

LONDON: PRINTED FOR
J. G. F \& J. RIVINGTON,

ST. PAUL'S CHURCH YARD, AND WATERLOO PLACE.

1840.

\xchapter{Advertisement}

THE
present Volume will be found to persevere in the change of plan adopted
in the latter part of the Second, the substitution of Tracts of
considerable extent of subject for the short and incomplete papers with
which the publication commenced. The reason of this change is to be
found in the altered circumstances under which they now make their
appearance. When the series began, the prospects of Catholic Truth were
especially gloomy, from the circumstance that irreligious principles and
false doctrines, which had hitherto been avowed only in the closet or on
paper, had just been admitted into public measures on a large scale,
with a probability of that admission becoming a precedent for future. A
great proportion of the Irish Sees had been suppressed by the State
against the Church's wish, all parties who were concerned to resist the
measure, acquiescing either in utter apathy or in despair. Scarcely a
protesting voice was heard, and the attempt to remonstrate was treated
on all hands with coldness and disapprobation. A sense of the dreariness
of such a state of things naturally led to those anxious appeals and
abrupt sketches of doctrine with which the Tracts opened. They were
written with the hope of rousing members of our Church to comprehend her
alarming position, of helping them to realize the fact of the gradual
growth, allowance, and establishment of unsound principles in the
management of her internal concerns; and, having this object, they
spontaneously used the language of alarm and complaint. They were
written, as a man might give notice of a fire or inundation, to startle
all who heard him, with only so much of doctrine and argument as might
be necessary to account for their publication, or might answer more
obvious objections to the views therein advocated.

This peculiarity in their composition has
occasioned them to be censured as intemperate and violent. If this be
true in such sense that they discovered any personal feeling,
bitterness, wrath, want of candour, unkindness, or reviling, of course
nothing can be said in their defence. Or, if they contain an extravagant
doctrine, crudely imagined, confusedly or hastily expressed, and
unsanctioned by our standard Divines, then, too, they are entitled to
very little respect. But if the charge of intemperance simply means that
they contain strong expressions upon high and delicate matters, suddenly
introduced, unexplained, and therefore obscure and harsh, though not
intrinsically erroneous, then by intemperance is meant nothing else than
want of judgment. Want of judgment, however, is commonly imputed to
proceedings which tend to defeat their object, though allowable in
themselves, and based upon true principles; and if so, the style of the
Tracts in question is \emph{not} injudicious, for their object has
\emph{not} been defeated. Naked statements, which offend the accurate
and cautious, are necessary upon occasions to infuse seriousness into
the indifferent.

These are the reasons, whether satisfactory or not
in the judgment of others, for the style and manner of the earlier
Tracts. When, however, from the circumstances of the times or from other
causes, more interest seemed to be excited among Churchmen concerning
those doctrines which it was their object to enforce, discussion became
more seasonable than the simple statements of doctrine with which the
series began; and their character accordingly changed.

It would be unbecoming to go into this detail in
this place, were not a prejudice entertained against these Tracts by
many who know them only by a few detached sentences, complete indeed in
themselves, and on the whole not unfairly selected, but which, so
detached, will not be understood in their true sense and bearings by
readers unacquainted with the language of our old divinity. Dr. Pusey's
valuable Pamphlet in answer to one objector, is, with the kind consent
of the Author, appended to this Advertisement [As Tract 77—\textbf{NR}].

OXFORD,

\emph{The Feast of All Saints}, 1836

\xchapter{Catena Patrum No. I. Testimony of Writers in the later English Church to the Doctrine of the Apostolical Succession}

Tract No. 74 (Ad Populum)

The Baptism of John, whence was it? from heaven or
of men? And they reasoned among themselves, saying, If we shall say,
From heaven, He will say unto us, Why did ye not then believe him? But
if we shall say, Of men, we fear the people; for all hold John as a
prophet.

PERSONS
who object to our preaching distinctly and unhesitatingly the doctrine
of the Apostolical succession, must be asked to explain, why we may not
do what our Fathers in the Church have done before us, or whether they
too, as well as we, are mistaken, or injudicious theorists, or Papists,
in so doing? This question is here plainly put to them; and at the same
time the attention of inquirers, who have not made up their minds on the
subject, is invited to the answer, if any is forthcoming, from the
parties addressed.

The doctrine in dispute is this; that CHRIST
founded a visible Church as an ordinance for ever, and endowed it once
for all with spiritual privileges, and set His Apostles over it, as the
first in a line of ministers and rulers, like themselves except in their
miraculous gifts, and to be continued from them by successive
ordination; in consequence, that to adhere to this Church thus
distinguished, is among the ordinary duties of a Christian, and is the means of his appropriating the Gospel blessings with an evidence of
his doing so not attainable elsewhere.

The passages quoted below contain, it is presumed,
this doctrine; but they are not intended as more than tokens and
suggestions of the full testimony, contained in the works of their great
authors.

\textbf{List of Authors cited.}

1. Bilson
23. Wake

2. Hooker
24. Potter

3. Bancroft
25. Nelson

4. Andrews
26. Kettlewell

5. Hall
27. Hicks

6. Laud
28. Law

7. Bramhall
29. Johnson

8. Mede
30. Dodwell

9. Mason
31. Collier

10. Sanderson
32. Leslie

11. Hammond
33. Wilson

12. Taylor
34. Bingham

13. Heylin
35. Skelton

14. Allestrie
36. Samuel Johnson

15. Pearson
37. Horne

16. Fell
38. W. Jones

17. Bull
39. Horsley

18. Stillingfleet
40. Heber

19. Ken
41. Jebb

20. Beveridge
42. Van Mildert

21. Sharp
43. Mant

22. Scott

\textbf{BILSON,
BISHOP.—Perpetual Government of Christ's Church,}

ch. ix. p. 105 \footnote{As quoted by Dr. Spry in his Bampton Lectures,
p. 311.}

It will happily [haply] be granted the Apostles
had their prerogative and pre-eminence above others in the Church of CHRIST;
but that limited to their persons, and during for their lives, and,
therefore, no reason can be made for their superiority, to force the
like to be received and established in the Church of CHRIST for all ages and places; since their office and function
are long since ceased, and no like power reserved to their successors
after them. I do not deny but many things in the Apostles were personal,
\&c. ... yet, that all their gifts ended with their lives, and no
part of their charge and power remained to their after-comers, may
neither be confessed by us, nor affirmed by any, unless we mean wholly
to subvert the Church of CHRIST
... The Scriptures, once written, suffice all ages for instruction; the
miracles then done, are for ever a most evident confirmation of their
doctrine; the authority of their first calling liveth yet in their
succession; and time and travel, joined with GOD'S
graces, bring pastors at this present to perfection; yet the Apostles'
charge to teach, baptize, and administer the LORD'S Supper, to bind and loose sinners in
heaven and in earth, to impose hands for the ordaining of pastors and
elders, these parts of the Apostolic function and charge are not
decayed, and cannot be wanted in the Church of GOD.
There must either be no Church, or else these must remain; for without
these no Church can continue.

\textbf{HOOKER,
PRESBYTER AND D\textsc{octor}.—Ecclesiastical
Polity},

Book v. § 77

… In that they are CHRIST'S ambassadors and His labourers, who
should give them their commission, but He whose most inward affairs they
manage? Is not GOD
alone the FATHER
of spirits? Are not souls the purchase of JESUS CHRIST? What angel in heaven could have said to
man, as our Lord did unto Peter, ``Feed my
sheep,—preach—baptize—do this in remembrance of Me. Whose sins ye
retain, they are retained; and their offences in heaven pardoned, whose
faults you shall on earth forgive?'' What think we? Are these
terrestrial sounds, or else are they voices uttered out of the clouds
above? The power of the ministry of GOD, translateth out of darkness into glory; it raiseth man from
the earth, and bringeth GOD
Himself from heaven; by blessing visible elements it maketh them
invisible graces; it giveth daily the HOLY GHOST;
it hath to dispose of that flesh which was given for the life of the
world, and that blood which was poured out to redeem souls; when it
poureth malediction upon the heads of the wicked, they perish; when it
revoketh the same, they revive. O wretched blindness, if we admire not
so great power; more wretched if we consider it aright, and,
notwithstanding, imagine that any but GOD
can bestow it! To whom CHRIST
hath imparted power, both over that mystical body which is the society
of souls, and over that natural which is Himself, for the knitting of
both in one, (a work which antiquity doth call the making of CHRIST'S body,) the same power is in such not
amiss both termed a kind of mark or character, and acknowledged to be
indelible … ``Receive the HOLY
GHOST;
whose sins soever ye remit, they are remitted; whose sins ye retain,
they are retained.'' Whereas, therefore, the other Evangelists had set
down, that CHRIST
did before His suffering promise to give His Apostles the keys of the
kingdom of heaven, and being risen from the dead, promised moreover at
that time a miraculous power of the HOLY
GHOST,
St. John addeth, that He also invested them even then with the power of
the HOLY
GHOST
for castigation and relaxation of sin, wherein was fully accomplished
that which the promise of the keys did import. Seeing therefore, that
the same power is now given, why should the same form of words
expressing it be thought foolish?

\emph{Ibid}. § 68.

Now the privilege of the visible Church of GOD
(for of that we speak) is to he herein like the ark of Noah, that, for
any thing we know to the contrary, all without it are lost sheep; yet in
this was the ark of Noah privileged above the Church, that whereas none
of them which were in the one could perish, numbers in the other are
cast away, because to eternal life our profession is not enough.

\textbf{BANCROFT,
ARCHBISHOP.—Sermon
preached at Paul's Cross}

This hath ever been reckoned a most certain ground
and principle in religion, that that Church, which maintaineth without
error the faith of CHRIST,
which holdeth the true doctrine of the Gospel in matters necessary to
salvation, and preacheth the same, which retaineth the lawful use of
those Sacraments only which CHRIST
hath appointed, and which appointeth vice to be punished, and virtue to
be maintained, notwithstanding in some other respects, and in some
points, it have many blemishes, imperfections, nay divers and sundry
errors, is yet to be acknowledged for the Mother of the faithful, the
House of GOD,
the Ark of Noah, the pillar of Truth, and the spouse of CHRIST. From which Church whosoever doth
separate himself, he is to be reckoned a schismatic or an heretic …

There are many causes set down by the said ancient
Fathers, why so many false prophets do go out into the world; but I will
only touch four; whereof I find the contempt of Bishops especially to be
one; for unto them, as St. Jerome saith, ever since St. Mark's time,
the care of Church Government hath been committed; they had authority
over the rest of the ministry … ``that the seed of schism might be
taken away, \&c.''

Read the Scriptures, but with sobriety; if any man
presuming upon his knowledge, seek further than is meet for him, besides

that he knoweth nothing as he ought to know, he
shall cast himself into a labyrinth, and never find that he seeketh for.
GOD hath
bound Himself by His promise unto His Church of purpose, that men by her
good direction might in this point be relieved; to whose godly
determination in matters of question, her dutiful children ought to
submit themselves without any curious or wilful contradiction. I could
bring many authorities to this effect.

\textbf{ANDREWS,
BISHOP AND
DOCTOR.—Sermons
on Whitsunday},

No. 9. (Works, p. 695.)

The HOLY
GHOST
may be received more ways than one. He hath many spiramina; [\emph{polytrop\={o}s}],
``in many manners'' He comes; and \emph{multiformis gratia} He comes
with. He and they carry the name of their cause; and to receive them is
to receive the SPIRIT.
There is a \emph{gratum faciens}, the saving grace of the SPIRIT, for one to save himself by, received by each, without
respect to others; and there is a \emph{gratis data} (whatever become of
us) serving to save others by, without respect to ourselves. And there
is [\emph{charis diakonias}], the grace of a holy calling, for it is a
grace, to be a conduit of grace any way. All these; and all from one and
the same SPIRIT.

That was here conferred, (in John xx. 22.) was not
the saving grace of inward sanctimony; they were not ``breathed on''
to that end. The Church to this day gives this still in her ordinations,
but the saving grace the Church cannot give; none but GOD can give that. Nor the \emph{gratis data} it is not. That came
by the tongues, both the gift of speaking divine languages, and the gift
of [\emph{apophthengesthai}], speaking wisely, and to the purpose: and
(we know) none is either the holier or the learneder by his ordination.

Yet a grace it is. For the very office itself is a
grace; \emph{mihi data est hæc gratia}, saith the Apostle, in more
places than one; and speaks of his office and nothing else. The
Apostleship was a grace, yet no saving grace. Else, should Judas have
been saved. Clearly then, it is the grace of their calling (this)
whereby they were sacred and made persons public, and their acts
authentical, and they enabled to do somewhat about he remission of
sins, that is not (of like avail) done by others, though perhaps, more
learned and virtuous than they, in that they have not the like \emph{mitto
vos}, nor the same \emph{accipite} that these have.

\emph{Ibid}.—Sermon on Absolution. (Appendix, p.
90.)

The power of remitting sin is originally in GOD,
and in GOD
alone. And CHRIST
our SAVIOUR,
by means of the union of the Godhead and manhood into one person, by
virtue whereof ``The Son of man hath power to forgive sins upon
earth.''

This power being thus solely vested in GOD,
He might, without wrong to any, have retained and kept to Himself, and
without means of word or sacrament, and without ministers, either
apostles or others, have exercised immediately by Himself from Heaven.
But we should then have said of the remission of sins, saith St. Paul, ``Who shall go up to heaven for it, and fetch it thence? for which
cause,'' saith he, ``the righteousness of faith speaketh thus, Say not
so, \&c.''

Partly this, but there should be no such difficulty
to shake our faith, as once to imagine to fetch CHRIST from heaven for the remission of our
sins; and partly also, because CHRIST, to whom alone this commission was originally granted,
having ordained Himself a body, would work by bodily things, and having
taken the nature of a man upon Him, would honour the nature He had so
taken, for these causes; that which was His, and His alone, He
vouchsafed to impart, and out of His commission to grant a commission,
and thereby to associate them to Himself, (it is His own word by the
prophet) and to make them [\emph{syergous}], that is \emph{cooperatores},
workers together with Him (as the Apostle speaketh) to the work of
salvation, both of themselves and of others. From GOD
then it is derived; from GOD
and to men. *
*
*

Now if we ask, to what men? the text is plain. They
to whom CHRIST
said this \emph{Remiseritis}, were the Apostles.

In the Apostles, (that we may come nearer yet) we
find three capacities as we may term them, 1. As Christians in general.
2. As preachers, priests, or ministers, more special. 3. As those
twelve persons, whom in strict propriety of speech, we term the
Apostles.

Some things that CHRIST spake to them, He spake to them as representing the whole
company of Christians; as his \emph{Vigilate}.

Some things to them, not as Christians, but as
preachers or priests; as His \emph{Ite prædicate Evangelium} and his \emph{Hoc
facile}; which no man thinketh all Christians may do.

And some things to themselves personally: as that
He had appointed them witnesses of His miracles and resurrection, which
cannot be applied but to them and them in person. It remaineth we
inquire, in which of these three capacities CHRIST
imparted to them this commission.

Not to the Apostles properly; that is, this was no
personal privilege to be in them, and to die with them, that they should
only execute it for a time, and none ever after them. GOD forbid we should so think it. For, this
power being more than needful for the world, (as in the beginning it was
said,) it was not to be either personal, or for a time; then those
persons dying, and those times determining, them in the ages following
(as we now in this) that should light into this prison or captivity of
sin, how could they or we receive any benefit by it? Of nature, it is
said by the heathen philosopher, that it does neither \emph{abundare in
superfluis}, nor \emph{deficere in necessariis}. GOD
forbid, but we should ascribe as much to GOD at the least, that neither He would ordain power superfluous
or more than needed, or else, it being needful, would appropriate it
unto one age, and leave all other destitute of it; and not rather, as
all writers both new and old take it, continue it successively to the
world's end.

And as not proper to the Apostles' persons, so
neither common to all Christians in general, nor in the persons of all
Christians conveyed to them. Which thing the very circumstances of the
text do evict. For He sent them first, and after inspired them; and
after both these, gave them this commission. Now all Christians are not
so sent, nor all Christians inspired with the grace or gift of the SPIRIT, that they were here. Consequently, it
was not intended to the whole society of Christians. Yea, I add, that
forasmuch as these two, both these two, must go before it, \emph{Missio}
and \emph{Inspiratio}, that though GOD
inspire some laymen, if I may have leave so to term them, with very
special graces of knowledge to this end, yet inasmuch as they have not
the former, of sending, it agreeth not to them, neither may they
exercise it, until they be sent, that is, until they have their calling
thereunto.

It being then neither personal nor peculiar to them
as Apostles, nor again common to all as Christians, it must needs be
committed to them as ministers, priests, or preachers; and consequently
to these that in that office and function do succeed them, to whom this
commission is still continued. Neither are they, that are ordained or
instituted to that calling, ordained or instituted by any other words or
verse than this (John xx. 23). Yet not so, that absolutely without them,
GOD cannot bestow it, on whom or when Him
pleaseth; or that He is bound to this means only, and cannot work
without it. For, \emph{gratia Dei non alligatur mediis}, the grace of GOD
is not bound but free, and can work without means either of word or
sacrament; and as without means, so without ministers, how and when to
Him seemeth good. But speaking of that which is proper and ordinary, in
the course by Him established, this is an ecclesiastical act, committed
as the residue of the ministry of reconciliation to ecclesiastical
persons. And if at any time He vouchsafe it by others that are not such,
they be in that case \emph{Ministri necessitatis, non officii}, in case
of necessity ministers, but by office not so. *
* *

The remission of sins, as it is from GOD
only, so it is by the death and blood-shedding of CHRIST alone; but for the applying of this unto us, there are
divers means established.
* *
* In
the institution of Baptism there is a power to that end.
* * 2. Again, there is also another power for
the remission of sins, in the institution of the Holy Eucharist. *
* 3.
Besides, in the word itself there is a like power ordained. ``Now are
you clean,'' saith CHRIST, (no doubt from their sins) \emph{propter
Sermonem hunc}. And the very name giveth as much, that it is
entitled, ``The Word of Reconciliation.'' 4. Further there is to the
same effect a power in prayer, and that in the priests' prayers, ``Call for the Priests,'' saith the Apostle,
``and let them pray for
the sick person, and if he have committed sin, it shall be forgiven
him.'' All and every of them, are acts for the remission of
sins: and in all and every of these, is the minister required, and they
cannot be dispatched without him.

But the ceremonies and circumstances that here
(John xx. 23.) I find used, prevail with me to think, that there is
somewhat here imparted to them, that was not before. For it carrieth no
likelihood, that our SAVIOUR
bestowing on them nothing here, but that which before He had, would use
so much solemnity, so diverse and new circumstances, no new or diverse
grace being here communicated.
* *
* take
it to be a power distinct from the former, and (not to hold you long) to
be the accomplishment of the promise made (Matt. xvi. 19.) of the power
of the keys, which here in this place and in these words is fulfilled;
and have therein for me, the joint consent of the Fathers. Which being a
different power in itself, is that which we call the act or benefit of
absolution; in which, as in the rest, there is in the due time and place
of it, an use for the remission of sins.

\textbf{HALL,
BISHOP AND
CONFESSOR.—On
Episcopacy},

Pt. iii. p. 9.

And for you, my dearly beloved brethren at home,
for CHRIST'S
sake, for the Church's sake, for your souls' sake be exhorted to
hold to this holy institution of your blessed SAVIOUR and His unerring apostles, and bless GOD
for Episcopacy. Do but cast your eyes a little back, and see what noble
instruments of GOD'S
glory He hath been pleased to raise up in this very Church of ours out
of this sacred vocation; what famous servants of GOD, what strong champions of truth, and renowned antagonists of
Rome and her superstitions; what admirable preachers, what incomparable
writers, yea, what constant and undaunted martyrs and confessors,
\&c. ... Neither doubt I but that it will please GOD, out of the same rod of Aaron, still to
raise such blossoms and fruit, as shall win Him glory to all eternity.
So you are to honour those your reverend pastors, to hate all factious
withdrawings from that government, which comes the nearest of any Church
upon earth to the Apostolical … Let me therefore confidently shut up
all with that resolute word of that blessed Martyr and Saint,
Ignatius … ``Let all things be done to the honour of GOD,
give respect to your bishop as you would GOD should respect you. My soul for theirs
which obey their bishop, presbyters, and deacons; GOD grant that my portion may be the same with
theirs.'' And let my soul have the same share with that blessed Martyr
that said so. Amen.

\textbf{LAUD,
ARCHBISHOP
AND MARTYR.—Conference
with Fisher},

xvi. 29.

``I am with you always unto the end of the
world.'' Yes most certain it is,—present by His SPIRIT; or else in bodily presence He continued not with His
Apostles, but during His abode on earth. And this promise of His
spiritual presence was to their successors; else, why ``to the end of
the world ?'' The Apostles did not, could not, live so long. But then
to the successors the promise goes no farther, than ``I am with you
always,'' which reaches to continual assistance, but not to divine and
infallible.

``The Comforter the HOLY GHOST shall abide with you for ever.'' Most
true again; for the HOLY
GHOST
did abide with the Apostles, according to CHRIST'S promise thus made, and shall abide with their
successors for ever, to comfort and preserve them.

\emph{Ibid}.—xxv. 15.

CHRIST
promised the Keys to St. Peter. (Matt. xvi.) True; but so did He to all
the rest of the Apostles (Matt. xviii. John xx.) and to their successors
as much as to his … St. Augustine is plain, ``If this were said only
of St. Peter, then the Church hath no power to do it,'' which GOD forbid! The keys therefore were given to
St. Peter and the rest in a figure of the Church, to whose power and for
whose use they were given. But there's not one key in all that bunch,
that can let in St. Peter's successors to a ``more powerful
principality'' universal than the successors of the other Apostles had.

\textbf{BRAMHALL,
ARCHBISHOP
AND CONFESSOR.—Vindication
of the Church of England.—Discourse III}

I do also acknowledge that Episcopacy was
comprehended in the Apostolic office, \emph{tanquam trigonus in tetragono},
and the distinction was made by the Apostles, with the approbation of CHRIST; that the angels of the seven Churches
in the Revelation were seven Bishops; that it is the most silly
ridiculous thing in the world, to calumniate that for a Papal
innovation, which was established in the Church before there was a Pope
at Rome: which hath been received and approved in all ages since the
very cradle of Christianity, by all sorts of Christians, Europeans,
Africans, Asiatics, Indians, many of which never had any intercourse
with Rome, nor scarcely ever heard of the name of Rome. If \emph{semper,
ubique, et ab omnibus} be not a sufficient plea, I know not what is.

But because I esteem them Churches not completely
formed, do I, therefore, exclude them from all hopes of salvation? or
esteem them aliens and strangers from the commonwealth of Israel? or
account them formal schismatics? No such thing. First, I know there are
many learned persons among them who do passionately affect Episcopacy;
some of which have acknowledged it to myself, that their Church would
never be rightly settled, until it was new moulded. Baptism is a
sacrament, the door of Christianity, a matriculation into the Church of
Christ: yet the very desire of it in case of necessity, is sufficient to
excuse from the want of actual Baptism. And is not the desire of
Episcopacy sufficient to excuse from the actual want of Episcopacy, in
like case of necessity? or should I censure these as Schismatics?

Secondly, there are others, who though they do not
long so much for Episcopacy, yet they approve it, and want it only out
of invincible necessity. In some places the sovereign prince is of
another communion; the Episcopal chairs are filled with Roman Bishops.
If they should petition for Bishops of their own, it would not be
granted. In other places the magistrates have taken away Bishops:
whether out of policy, because they thought that regiment not so proper
for their republics, or because they were ashamed to take away the
revenues, and preserve the order, or out of a blind zeal; they have
given an account to GOD: they owe none to me. Should I condemn all these as
schismatics for want of Episcopacy, who want it out of invincible
necessity?

Thirdly, there are others who have neither the same
desires, nor the same esteem of Episcopacy, but condemn it as an Anti-christian
innovation, and a rag of Popery. I conceive this to be most gross schism
materially. It is ten times more schismatical to desert, nay, to take
away (so much as lies in them) the whole order of Bishops, than to
subtract obedience from one lawful Bishop. All that can be said to
mitigate this fault is, that they do it ignorantly, as they have been
mistaught and misinformed. And I hope that many of them are free from
obstinacy, and hold the truth implicitly in the preparation of their
minds, being ready to receive it, when GOD
shall reveal it to them. How far this may excuse (not the crime but)
their persons from formal schism, either \emph{a toto} or \emph{a tanto},
I determine not, but leave them to stand or fall to their own Master.

But though these Protestants were worthy of this
contumely, yet surely the Romanists are no fit persons to object it,
whose \emph{opinastrety} did hinder an uniform reformation of the
Western Church. Who did invest Presbyters with Episcopal jurisdiction,
and the power of ordaining and confirming, but the Court of Rome, by
their commissions and delegations, for avaricious ends? And could they
think that the world would believe, that necessity is not as strong and
effectual a dispensation, as their mercenary Bulls? It is not at all
material, whether Episcopacy and Priesthood be two distinct orders, or
distinct degrees of the same orders, the one subordinate to the other;
whether Episcopal ordination do introduce a new character, or extend the
old. For it is generally confessed by both parties, Protestants and
Roman Catholics, that the same power and authority is necessary to the
extension of a character, or grace given by ordination, which is
required to the Institution of a Sacrament that is not human but divine.
These avaricious practices of that Court (though it be not commonly
observed) were the first source of these present controversies
about Episcopacy, and Ecclesiastical discipline, which do now so much
disturb the peace of the Church.

\emph{Ibid}.—Vindication of Grotius.—Discourse
III.

Excuse me for telling the truth plainly; many who
have had their education among Sectaries and Non-conformists have
apostated to Rome, but few or no right Episcopal Divines. Hot water
freezeth the soonest.

He addeth, that ``Grotius himself assures him
(whom he hath reason to believe) that there were not a few such among
the prelatical men.'' How! not a few such as these, who have apostated
from the Church of England? For ingenuity's sake let him tell us where
Grotius saith any such thing. Grotius hath not one word to his purpose,
when it is duly examined. But this it is to confute books in less time
than wise or modest men would require to read them.

Hitherto, he is not able to show us any tolerable
reason of his warning. But he showeth us the occasion, p. 82, ``Those
that unchurch either all or most of the Protestant Churches, and
maintain the Roman Church and not theirs to be true, do call us to a
moderate jealousie of them.'' This is far enough from proving his bold
suggestion, that they have a design to introduce the Pope into England.
So though all he say were true; yet he can conclude nothing from thence
to make good his accusation or insinuation. I wish he would forbear
these imperfect enthymematical forms of argument, which serve only to
cover deceit, and set down both his propositions expressly. His
assumption is wanting, which should be this: but a considerable party of
Episcopal divines in England do unchurch all or most of the Protestant
Churches, and maintain the Roman Church to be a true Church, and these
to be no true Churches. I can assent to neither of his prepositions, nor
to any part of them, as true, \emph{sub modo}, as they are alleged by
him.

First, I cannot assent to his major proposition,
that all those who make an ordinary personal uninterrupted succession of
Pastors, to be of the integrity of a true Church, (which is the ground
of his exception) have, therefore, an intention, or can justly be
suspected thereupon to have any intention to introduce the Pope. The
Eastern, Southern, and Northern Churches are all of them for such a
personal succession, and yet all of them utter enemies to the Pope.
Secondly, I cannot assent to his minor proposition, that either all or
any considerable part of the Episcopal divines in England, do unchurch
either all or most part of the Protestant Churches. No man is hurt but
by himself. They unchurch none at all, but leave them to stand or fall
to their own Master. They do not unchurch the Swedish, Danish, Bohemian
Churches, and many other Churches in Polonia, Hungaria, and those parts
of the world, who have an ordinary uninterrupted succession of Pastors,
some by the names of Bishops, others under the name of Seniors, unto
this day. (I meddle not with the Socinians:) they unchurch not the
Lutheran Churches in Germany, who both assert Episcopacy in their
confessions, and have actual superintendents in their practice, and
would have Bishops, name and thing, if it were in their power. Let him
not mistake himself: those Churches which he is so tender of, though
they be better known to us by reason of their vicinity, are so far from
being ``all or most part of the Protestant Churches,'' that being all
put together, they amount not to so great a proportion as the Britannick
Churches alone. And if one secluded out of them all those who want an
ordinary succession without their own faults, out of invincible
ignorance or necessity, and all those who desire to have an ordinary
succession either explicitly or implicitly, they will be reduced to a
little flock indeed.

But let him set his heart at rest. I will remove
this scruple out of his mind, that he may sleep securely upon both ears.
Episcopal divines do not deny those Churches to be true Churches,
wherein salvation may be had. We advise them, as it is our duty, to be
circumspect for themselves, and not to put it to more question, whether
they have ordination or not, or desert the general practice of the
universal Church for nothing, when they may clear it if they please.
Their case is not the same with those who labour under invincible
necessity. What mine own sense is of it, I have declared many years
since to the world in print; and in the same way received thanks, and a
public acknowledgment of my moderation from a French divine. And yet
more particularly in my reply to the Bishop of Chalcedon, Pres. p. 144.
and cap. i. p. 164. Episcopal divines will readily subscribe to the
determination of the learned Bishop of Winchester, in his answer to the
second epistle of Molineus. ``Nevertheless, if our form (of Episcopacy)
be of divine right, it doth not follow from thence, that there is not
salvation without it, or that a Church cannot consist without it. He is
blind who does not see churches consisting without it: he is
hard-hearted who denyeth them salvation. We are none of those
hard-hearted persons, we put a great difference between these things.
There may be something absent in the exterior regiment, which is of
divine right, and yet salvation to be had.'' This mistake proceedeth
from not distinguishing between the true nature and essence of a Church,
which we do readily grant them, and the integrity or perfection of a
Church, which we cannot grant them, without swerving from the judgment
of the Catholic Church.

\textbf{MEDE,
PRESBYTER.—Sermon
on Urim and Thummim. Works},

Book 1. p. 186.

The Ministers of Christ must be \emph{Lux Mundi},
the light of the world—\emph{Vos estis Lux Mundi}—``Ye are the
Light of the World: Ye are the World's Urim,'' saith Christ unto His
Apostles. ``For the lips of the priest should preserve knowledge, and
they should learn the law at his mouth.'' This light of knowledge, this
teaching knowledge, is the Urim of every Levite; and therefore Christ,
when He inspired his Apostles with knowledge of heavenly mysteries, He
sent a new Urim from above, even fiery tongues of Urim from heaven. He
sent no fiery heads, but fiery tongues; for it is not sufficient for a
Levite to have his head full of Urim, unless his tongue be a candle to
show it to others. There came, indeed, no Thummim [integrity or
perfection] from heaven, as there came an Urim; for though the Apostles
were secured from errors, they were not freed from sin; and yet we who
are Levites must have such a Thummim as may be gotten upon earth;
for St. Paul bids Titus in all things to show himself an example of good
works, and this is a Thummim of Integrity. But, besides this Thummim,
the Ministers of the Gospel have received from GOD
more especially another Thummim, like unto that which was proper to the
High Priest; namely, the power of binding and loosing, which is, as it
were, a power of oracle, to declare unto the people the remission of
their sins, by the acceptance of CHRIST'S
sacrifice.

\textbf{MASON,
PRESBYTER.—Vindiciæ
Ecclesiæ Anglicanæ},

i. 2.

\emph{Anglican}. Our Ministry is agreeable to
Divine Scripture, and therefore holy. Nor do we doubt, that, when the
Chief Shepherd shall appear, they who turn many to righteousness shall
shine as the stars for ever and ever. However, what is your argument
against our ministry?

\emph{Romanist}. Can a man be a lawful minister
without a lawful call?

\emph{Anglican}. Of course not.

\emph{Romanist}. If so, I pray tell we how the
Anglican Chinch can defend her ministry. Surely I may address each of
you in Harding's words to Jewel: ``What say you, my master? You bear
yourself as though Bishop of Salisbury; but how will you substantiate
your call? What is your warrant for ministering in the Word and
Sacraments?'' \&c. \&c. ... I ask thee, Is your call inward or
outward?

\emph{Anglican}. Both.

\emph{Romanist}. An outward call, to be lawful,
must be either immediately from CHRIST'S
mouth, as the Apostles were called, or mediately through the Church.

\emph{Anglican}. Well; we are called by GOD
through the Church; for it is He who gives ``Pastors and Doctors for
the perfecting of the Saints.''

Romanist. They who are called by GOD through the Church, must derive
their warrant and power by lawful succession from CHRIST and the Apostles. If you maintain you
have proceeded from this origin, it is your business to prove it clearly
to us; to set forth and trace your genealogy …

\emph{Anglican}. The Ministers of the Anglican
Church derive their imposition of hands in a lawful way from lawful
Bishops, possessed of a lawful authority; and therefore their call is
ordinary [not extraordinary, by miracles].

\emph{Romanist}. But whence have these Bishops
derived their power?

\emph{Anglican}. From GOD, through the hands of Bishops before them,
\&c. \&c.

\textbf{SANDERSON,
BISHOP AND
CONFESSOR.—Divine
Right of Episcopacy}

My opinion is, that Episcopal Government is not to
be derived merely from Apostolical practice or institution, but that it
is originally founded in the Person and Office of the Messias, our
blessed LORD
JESUS CHRIST;
who, being sent by our heavenly Father to be the Great Apostle (Heb.
iii. 1), Bishop, and Pastor (1 Peter ii. 25) of His Church, and anointed
to that office immediately after His baptism by John, with power and the
HOLY GHOST (Acts x. 37, 38), descending then upon
Him in a bodily shape (Luke iii. 22), did afterwards, before His
ascension into Heaven, send and empower His holy Apostles, in like
manner as His Father had before sent Him (John xx. 21), to execute the
same Apostolical, Episcopal, and Pastoral Office, for the ordering and
governing of His Church, until His coming again; and so the same office
to continue in them and their successors unto the end of the world.
(Matt. xxviii. 18-20.) This I take to be so clear, from these and other
like texts of Scripture, that if they shall be diligently compared
together, both between themselves and with the following practice of the
Churches of CHRIST,
as well in the Apostles' times as in the purest and primitive times
nearest thereunto, there will be left little cause why any man should
doubt thereof.

\textbf{HAMMOND,
PRESBYTER,
DOCTOR, AND
CONFESSOR.—On
the Power of the Keys. Preface}

That the prime act of power enstated by CHRIST
on His Apostles, as for the governing of the Church, (and exercising or
banishing all devils out of it,) so for the effectual performing that
great act of charity to men's souls, reducing pertinacious
sinners to repentance, should be so, either wholly dilapidated, or
piteously deformed, as to continue in the Church, only under one of
these two notions, either of an empty piece of formality, or of an
engine of state and secular contrivance, (the true Christian use of
shaming sinners into reformation, being well nigh vanished out of
Christendom,) might by an alien, or an heathen, much more by the
pondering Christian, be conceived very strange and unreasonable, were it
not a little clear that we are fallen into those times, of which it was
foretold by two Apostles, that in ``these last days, there should come
scoffers, walking after their own lusts,'' \&c. ... I shall design
to infer no further conclusion, but only this, that they which live ill
in the profession of a most holy faith … but especially they that
discharge and banish out of the Church those means which might help to
make the generality of Christians better, have the spirit of Antichrist
working in them, even when they think themselves most zealously busied
in beating down his kingdom. What those means are which might most
effectually tend to the amending the lives of Christians, I shall need
no farther to interpose my judgment, than, 1st, by submitting it to CHRIST, who put the keys into the
Apostles' hands, on purpose as a means to exemplify the end of His
coming … 2nd, by minding myself and others what the Apostles say of
this power, that it was given them [\emph{pros oikodom\={e}n}], to
build up the Church of CHRIST,
\&C.

\emph{Chapter} 3. The only difficulty remaining on
the point, will be, who are the Apostles' successors in that power;
and when the question is asked of that power, I must be understood of
the power of governing the Church peculiarly, (of which the power of the
keys was and is a principal branch,) for it must again be remembered
that the Apostles are to be considered under a double notion, first, as
planters, then as governors, of the Church … Which distinction being
premised, the question will now more easily be satisfied, being proposed
in these terms; who were the Apostles' successors in that power, which
concerned the governing their Churches which they planted? and first, I
answer, that it being a matter of fact, or story, later than that the
Scripture can universally reach to, it cannot be fully satisfied or
answered from thence ... but will in the full latitude, through the
universal Church in these times, be made clear, from the recent
evidences that we have, viz. from the consent of the Greek and Latin
Fathers, who generally resolve that Bishops are those successors.

\textbf{TAYLOR,
BISHOP,
CONFESSOR,
AND DOCTOR.—On
Episcopacy. Introduction}

Antichrist must come at last, and the great
apostasy foretold must be, and this not without means proportionable to
the production of so great declensions of Christianity. ``When ye hear
of wars and rumours of wars, be not afraid,'' says our blessed SAVIOUR, ``the end is not yet.'' It is not
war that will do ``this great work of destruction;'' for then it might
have been done long ere now. What then will do it? We shall know when we
see it. In the mean time, when we shall find a new device, of which,
indeed, the platform was laid, in Aërius and the Acephali, brought to a
good possibility of completing a thing, that whosoever shall hear, his
ears shall tingle, ``an abomination of desolation standing where it
ought not,'' ``\emph{in sacris},'' in holy persons, and places, and
offices, it is too probable that this is preparatory for the Antichrist,
and grand apostasy.

For if Antichrist shall exalt himself above all
that is called GOD,
and in Scripture none but kings and priests are such, ``dii vocati, dii
facti,'' I think we have great reason to be suspicious, that he that
divests both of their power, (and they are, if the king be Christian, in
very near conjunction,) does the work of Antichrist for him; especially
if the men whom it most concerns will but call to mind, that if the
discipline or government which CHRIST
hath instituted is that kingdom by which He governs all Christendom, (so
themselves have taught us,) when they (to use their own expressions)
throw CHRIST out of His kingdom; and then
either they leave the Church without a head, or else put Antichrist in
substitution.

We all wish that our fears in this and all things
else may be vain, that what we fear may not come upon us; but yet that
the abolition of episcopacy is the forerunner, and preparatory to the
great Apostasy, I have these reasons to show, at least, the
probability. First, \&c. *
* *

\emph{Sections 2 and 3}. This government was by
immediate substitution delegated to the Apostles, by CHRIST Himself, ``in traditione clavium,
inspiratione Spiritûs, in missione in Pentecosto.'' … This power so
delegated, was not to expire with their persons; for when the great
Shepherd had reduced His wandering sheep into a fold, He would not leave
them without ``guides to govern'' them, so long as the wolf might
possibly prey upon them, and that is, till the last separation of the
sheep from the goats. And this CHRIST
intimates in that promise, ``Ero vobiscum (Apostolis) usque ad
consummationen seculi.'' ``Vobiscum;'' not with your persons, for
they died long ago: but ``vobiscum et vestri similibus,'' with
Apostles to the end of the world. And, therefore, that the Apostolate
might be successive and perpetual, CHRIST gave them a power of Ordination, that by
imposing hands on others, they might impart that power which they
received from CHRIST.

\textbf{HEYLIN,
PRESBYTER
AND CONFESSOR.—On
Episcopacy},

i. 6.

The Church, at his (St. John's) departure, he
left firmly grounded in all the points of faith and doctrine, taught by
CHRIST
our SAVIOUR,
as well settled in the outward government, the polity and administration
of the same, which had been framed by the Apostles, according to the
pattern and example of their LORD
and Master. For being that the Church was born of seed immortal, and
they themselves, though excellent and divine, yet still mortal men, it
did concern the Church, in a high degree, to be provided of a
perpetuity, or, if you will, the immortality of overseers, both for the
sowing of this seed, and for the ordering of the Church, or the field
itself. This, since they could not do in person, they were to do it by
successors, who by their office were to be the ordinary pastors of the
Church, and the Vicars of CHRIST. Now, if you ask the Fathers who they
were that were accounted in their times and ages the successors of the
Apostles, they will with one accord make answer, that the Bishops were.

\textbf{ALLESTRIE,
PRESBYTER.—Sermons,
No. 16.}

The separateness of the functions of the Clergy,
the incommunicableness of their offices to persons not separated for
them, is so express a doctrine both of the letter of the text, and of
the HOLY
GHOST,
that sure I need not to say more, though several heads of probation
offer themselves; as first the condition of the callings, which does
divide from the community and sets them up above it. And here I might
tell you of ``bearing rule,'' of ``thrones,'' of ``stars,'' and ``Angels,'' and other words of a high sense, and yet not go out of the
Scripture bounds, although the dignity did not die with the Scripture
age, or expire with the Apostles; the age as low as Photius words it
thus, [\emph{To apostolikon, k. t. l.}] ``That Apostolical and Divine
dignity, which the Chief Priests are acknowledged to be possessed of by
right of succession.'' Styles which I could derive yet lower, and they
are of a prouder sound than those the modest bumble ears of this our age
are so offended with. But these heights, it may be, would give umbrages;
although it is strange that men should envy them to those, who are only
exalted to them, that they may with the more advantage take them by the
hands and lift them up to heaven. Those nearesses to things above do but
more qualify them to draw near to GOD,
on your behalf, that these your Angels also may see the face of your FATHER which is in heaven, and those
stars are, therefore, set in CHRIST'S right hand, that they may shed a blessing [blessed?]
influence on you from thence …

The censures of the Church, the burden of the keys,
which (passing by the private use of them in voluntary penitences, and
discipline upon the sick,) as they signify public exclusion out of the
Church, for scandalous enormities, and readmission into it upon
repentance, have been sufficiently evinced to belong to the governors of
the Church. The exercise of them is so much their work, that St. Paul
calls them ``the weapons of their spiritual warfare, by which they do
cast down imaginations, and every high thing that exalteth itself
against the knowledge of GOD, and bring into captivity every
thought to the obedience of CHRIST,'' (2 Cor. x. 4, 5.) a blessed victory, even for the
conquered, and these the only weapons to achieve it with. If those
who sin scandalously, and will not hear the admonition of the Church,
were cast out of the Church, if not religion, reputation would restrain
them somewhat; not to be thought fit company for Christians, would
surely make them proud against their vices. Shame, the designed effect
of their censures, hath great pungencies; the fear of it does goad men
into actions of the greatest hazard, and the most unacceptable; such as
have nothing lovely in them, but are wholly distasteful ... Now, the
infliction of these censures is so much the work to which Church
governors are called by the HOLY GHOST,
that they are equally called by Him to it and to Himself; both are alike
bestowed upon them. ``Receive the HOLY GHOST;
whose sins ye retain, they are retained.'' (John xx. 22.) And in the
first derivations of this office, it was performed with severities, such
as this age, I doubt, will not believe; and when they had no temporal
sword to be auxiliary to these spiritual weapons.

\textbf{PEARSON,
BISHOP AND
DOCTOR.—On
the Creed, Article ix.}

[After considering the Church as one, by reason of
its one foundation, faith, ministry of sacrament, hope, and charity, he
continues,—]

Lastly, all the Churches of GOD are united into one by the unity of
discipline and government, by virtue whereof the same CHRIST ruleth in them all. For they have
all the same pastoral guides appointed, authorized, sanctified, and set
apart by the appointment of GOD, by the direction of the SPIRIT, to direct and lead the people of GOD
in the same way of eternal salvation: as, therefore, there is no Church
where there is no order, no ministry; so, where the same order and
ministry is, there is the same Church.
* *
*

The necessity of believing the Holy Catholick
Church appeareth first in this, that CHRIST hath appointed it as the only way unto eternal life. We
read at the first, ``The LORD
added to the Church daily such as should be saved;'' and what was then
daily done hath been done since continually. CHRIST never appointed two ways to heaven; nor
did He build a Church to save some, and make another institution
for other men's salvation. ``There is none other name under heaven
given among men, whereby we must be saved,'' but the name of JESUS;
and that is no otherwise given under heaven than in the Church. As none
were saved from the Deluge but such as were within the ark of Noah,
formed for their reception by the command of GOD; as none of the first-born of Egypt lived, but such as were
within those habitations whose door-posts were sprinkled with blood by
the appointment of GOD
for their preservation; as none of the inhabitants of Jericho could
escape the fire or sword, but such as were within the house of Rahab,
for whose protection a covenant was made; so none shall ever escape the
eternal wrath of GOD,
which belong not to the Church of GOD.
This is the congregation of those persons here on earth which shall
hereafter meet in heaven. These are the vessels of the tabernacle
carried up and down, at last to be translated into and fixed in the
Temple.

Next, it is necessary to believe the Church of CHRIST,
which is but one, that, being in it, we may take care never to cast
ourselves, or be ejected, out of it. There is a power within the Church
to cast those out which do belong to it: for if any neglect to hear the
Church, saith our SAVIOUR,
let him be unto thee as a heathen man and a publican. By great and
scandalous offences, by incorrigible misdemeanors, we may incur the
censure of the Church of GOD;
and while we are shut out by them, we stand excluded out of Heaven. For
our SAVIOUR
said to His Apostles, upon whom He built his Church, ``Whosesoever sins
ye remit, they are remitted unto them; and whosesoever sins ye retain,
they are retained.'' Again, a man may not only passively and
involuntarily be ejected, but also may, by an act of his own, cast out
or eject himself, not only by plain and complete apostasy, but by a
defection from the unity of truth, falling into some damnable heresy; or
by an active separation, deserting all which are in communion with the
Catholick Church, and falling into an irrevocable schism.

\textbf{FELL,
BISHOP AND
CONFESSOR.——On
Ephesians},

v. 9.

Apostles, … Prophets, … Evangelists, …
Pastors, … and Teachers ... For the three first, some part of
their function was temporary and extraordinary; in what was ordinary and
perpetual, Bishops succeeded.

\textbf{BULL,
BISHOP AND
DOCTOR.—Vindication
of the English Church, § 24.}

We proceed, in the next place, to the constant
visibility and succession of Pastors in our Church … And here I make
him this fair proposal: Let him, or any one of his party, produce any
one solid argument to demonstrate such a succession of Pastors in the
Church of Rome, and I will undertake by the very same argument to prove
a like succession in our Church. Indeed, … the Author of the Letter is
concerned, no less than we are, to acknowledge such a succession of
lawful pastors in our Church, till the time of the Reformation; and if
we cannot derive our succession since, it is a hard case. But our
records, faithfully kept and preserved, do evidence to all the world an
uninterrupted succession of Bishops in our Church, canonically ordained,
derived from such persons in whom a lawful power of ordination was
seated by the confession of the Papists themselves. For the story of the
Nag's head Ordination is so putrid a fable, so often and so clearly
refuted by the writers of our Church, that the more learned and
ingenuous Papists are now ashamed to make use of it.

\textbf{STILLINGFLEET,
BISHOP.—Unreasonableness
of Separation; Preface}

Unthinking people ... are carried away with mere
noise and pretences, and hope these will secure them most against the
fears of Popery, who talk with most passion, and with least
understanding, against it; whereas no persons do really give them
greater advantages than these do. For, where they meet only with
intemperate railings, and gross misunderstandings of the state of the
controversies between them and us (which commonly go together), the most
subtle priests let such alone to spend their rage and fury; and when the
heat is over, they will calmly endeavour to let them see how grossly
they have been deceived in some things, and so wilt more easily
make them believe, they are as much deceived in all the rest. And thus
the East and West may meet at last, and the most furious antagonists may
become some of the easiest converts. This I do really fear will be the
case of many thousands among us, who now pass for most zealous
Protestants; if ever, which GOD
forbid, that religion should come to be uppermost in England. It is,
therefore, of mighty consequence for preventing the return of Popery,
that men rightly understand what it is. For, when they are as much
afraid of an innocent ceremony as of real idolatry, and think they can
worship images and adore the Host on the same grounds that they may use
the sign of the cross, or kneel at the Communion, when they are brought
to see their mistake in one case, they will suspect themselves deceived
in the other also.

When they find undoubted practices of the Ancient
Church condemned as Popish and Antichristian by their teachers, they
must conclude Popery to be of much greater antiquity than really it is;
and when they can trace it so very near the Apostles' times, they will
soon believe it settled by the Apostles themselves. For it will be very
hard to persuade any considering men, that the Christian Church should
degenerate so soon, so unanimously, so universally, as it must do, if
Episcopal government, and the use of some significant ceremonies, were
any parts of that apostasy … Three ways, Bishop Sanderson observes,
our dissenting brethren, though not intentionally and purposely, yet
really and eventually, have been the great promoters of the Roman
interest among us; (1) by putting-to their helping hand to the pulling
down of Episcopacy … (2) by opposing the interest of Rome with more
violence than reason; (3) by frequently mistaking the question, but
especially through the necessity of some false principle or other,
which, having once imbibed, they think themselves bound to maintain,
whatever becomes of the common cause of our Reformation.

\textbf{KEN,
BISHOP AND
CONFESSOR.—Exposition
of the Church Catechism}

I believe, O blessed and adorable Mediator, that
the Church is a society of persons, founded by Thy love to sinners,
united into one body of which Thou art the head, initiated by baptism,
nourished by the Eucharist, governed by pastors commissioned by Thee,
and endowed with the power of the keys, professing the doctrine taught
by Thee and delivered to the saints, and devoted to praise and to love
Thee.

I believe, O holy JESUS, that Thy Church, is holy like Thee its Author; holy by the
original design of its institution, holy by baptismal dedication, holy
in all its administrations which tend to produce holiness; and though
there will be always a mixture of good and bad in it in this world, yet
that it has always many real saints in it; and, therefore, all love, all
glory, be to Thee …

Glory be to Thee, O LORD my GOD, who hast made me a member of the
particular Church of England, whose faith, government, and worship are
Holy, and Catholic, and Apostolic, and free from the extremes of
irreverence and superstition, and which I firmly believe to be a sound
part of Thy Church universal, and which teaches me charity to those who
dissent from me; and, therefore, all love, all glory be to Thee.

\textbf{BEVERIDGE,
BISHOP AND
DOCTOR.—Sermon
on Christ's Presence with His Ministers}

In the first place, I observe, how much we are all
bound to acknowledge the goodness, to praise, magnify, and adore the
name of the most high GOD,
in that we were born and bred, and still live in a Church, wherein the
Apostolical line hath, through all ages, been preserved entire, there
having been a constant succession of such Bishops in it, as were truly
and properly successors to the Apostles, by virtue of that Apostolical
imposition of hands, which, being begun by the Apostles, hath been
continued from one to another, ever since their time, down to ours. By
which means, the same SPIRIT
which was breathed by our LORD
into His Apostles is, together with their office, transmitted to their
lawful successors, the pastors and governors of our Church at this time;
and acts, moves, and assists at the administration of the several parts
of the Apostolical office in our days, as much as ever. From whence
it follows, that the means of grace which we now enjoy are in themselves
as powerful and effectual as they were in the Apostles' days, \&c.
…

And this, I verily believe, is the great reason why
the devil has such a great spite at our Church, still stirring up
adversaries of all sorts against it,—Papists on the one hand, and
Sectaries on the other, and all, if possible, to destroy it; even
because the SPIRIT which is ministered in it, is so
contrary to his nature, and so destructive of his kingdom, that he can
never expect to domineer and tyrannize over the people of the land, so
long as such a Church is settled among them, and they continue firm to
it …

As for schism, they certainly hazard their
salvation at a strange rate, who separate themselves from such a Church
as ours is, wherein the Apostolical succession, the root of all
Christian communion, hath been so entirely preserved, and the word and
sacraments are so effectually administered; and all to go into such
assemblies and meetings, as can have no pretence to the great promise in
my text. For it is manifest, that this promise was made only to the
Apostles and their successors to the end of the world. Whereas, in the
private meetings, where their teachers have no Apostolical or Episcopal
imposition of hands, they have no ground to succeed the Apostles, nor by
consequence any right to the SPIRIT which our LORD hath; without which, although they preach their hearts out, I
do not see what spiritual advantage can accrue to their hearers by it,
\&c. …

\textbf{SHARP,
ARCHBISHOP.—Sermons,
Vol. vii. Of the Church}

``Go,'' He says, ``and make disciples of all
nations, baptizing them,'' \&c. … This commission of our SAVIOUR we may properly style the Charter of
the Church; and mind, I pray, what is contained in it. Our SAVIOUR
here declares the extent of His Church, and of what persons He would
have it constituted. It was to extend throughout all the world, and to
be made up of all nations. He here declares by whom He would have it
built and constituted, viz., the Apostles. He here declares upon what
grounds He would have it constituted, or upon what conditions any person
was to be received into it, viz., their becoming the disciples of JESUS CHRIST,
and undertaking to observe all that He has commanded. He here likewise
declares the form or the method by which persons were to be admitted
into this Church, and that was by being baptized in the name of the FATHER, and of the SON, and of the HOLY
GHOST.
And, lastly, He here promises the perpetual presence of His HOLY SPIRIT, both to assist the apostles and their
successors in the building and governing this Church, and to actuate and
enliven all the members of it … Thus, I am sure, I have given the true
notion of the Church, which the Scripture always intends, when it speaks
of the Church as the Body of CHRIST,
when it speaks of the Church which CHRIST purchased with His blood, when it speaks of the Church into
which we are baptized, when it speaks of the Church to which all those
glorious promises are made of the forgiveness of sins, of the perpetual
presence and assistance of the HOLY SPIRIT,
of the gates of hell never prevailing against it, and of everlasting
salvation in the world to come; I say, that Church is always meant of
the whole company of Christians dispersed over all the world, that
profess the common faith, (though perhaps none of them without mixture
of error,) and enjoy the administration of the word and sacrament, under
their lawful pastors and governors: all these people, wherever they
live, or by what name soever they call themselves, make up together that
one Body of CHRIST which we call the Catholic
Church.

\textbf{SCOTT,
PRESBYTER.—Christian
Life, Part ii. ch. 7.}

Another thing wherein those particular Churches,
into which the Catholic Church is distributed, do communicate with each
other, is, in the essentials of Christian regiment and discipline: for
though the particular modes and circumstances of Christian government
and discipline are not determined by divine institution, but left for
the most part free to the prudent ordering and disposal of the governors
of particular Churches, yet there is a standing form of government and
discipline in the Church, instituted by our SAVIOR
Himself, which, as I shall show hereafter, is this; that there should be
an episcopacy, or order of men, authorized in a continual succession
from the Apostles, (who were authorized by Himself) to oversee and
govern all those particular Churches into which the Church Catholic
should be hereafter distributed; to ordain, \&c. \&c. And this
being the standing government and discipline of the Catholic Church, no
particular Church or community of Christians can refuse to communicate
in it, without dividing itself from the communion of the Church
Catholic; I say, ``refuse to communicate in it,'' because it is
possible for a Church to be without this government and discipline,
which yet doth neither refuse it, nor the communion of any other Church
for the sake of it. A Church may be debarred of it by unavoidable
necessities, in despite of its power and against its consent … Though
this instituted government is necessary to the perfection of a Church,
yet it doth not therefore follow, that it is necessary to the being of
it … But though a community of Christians may be a true part of the
Catholic Church, and in communion with it, though it hath no episcopacy;
yet it is a plain case, that if it rejects the episcopacy, and separates
from the communion of it, it thereby wholly divides itself from the
communion of the Catholic Church.

\textbf{WAKE,
ARCHBISHOP.—Exposition
of the Doctrine of the English Church. Art. 15.}

The imposition of hands in Holy Orders, being
accompanied with a blessing of the HOLY
SPIRIT,
may, perhaps, upon that account, be called a kind of particular
sacrament. Yet since that grace, which is thereby conferred, whatever it
be, is not common to all Christians, nor by consequence any part of that
federal blessing which our blessed SAVIOUR
has purchased for us, but only a separation of him who receives it to a
special employ, we think it ought not to be esteemed a common sacrament
of the whole Church, as Baptism and the LORD'S Supper are … We confess that no man ought to exercise
the ministerial office till he be first consecrated to it. We believe
that it is the Bishop's part only to ordain. We maintain the
distinction of the several orders in the Church and though we have none
of them below a deacon, because we do not read that the Apostles had
any, yet we acknowledge the rest to have been anciently received in
the Church, and shall not therefore raise any controversy about them.

WAKE.—Art.
25.

Professing in our Creed a Holy Catholic Church, we
profess to believe not only that there was a Church planted by our SAVIOUR
at the beginning, that has hitherto been preserved by Him, and ever
shall be to the end of the world; but do in consequence undoubtedly
believe too, that this universal Church is so secured by the promises of
CHRIST, that there shall always be retained so
much truth in it, the want of which would argue that there could be no
such Church.

\textbf{POTTER,
ARCHBISHOP.—On
Church Government. Chap. v.}

First, then, it must be shown, that the office and
character of all persons, who are admitted into holy orders, extends
over the whole world, and it is manifest, in the first place, that the
Apostles had a general commission to ``teach and baptize,'' and to
execute all other parts of their office in all nations. As the bishops
of the Church have been shown to succeed the Apostles in all the parts
of their office which are of standing and constant use in the Church; so
we might reasonably conclude, though we have no farther proof of it,
that the office and character of bishop, and consequently of inferior
ministers, extends over all the world, because those of the Apostles,
their predecessors, did so; since there is scarce any reason why the
Apostles' authority should be universal, which will not hold, at least
in some degree, for the same extent of authority in the bishops, as will
appear from some of the following considerations: —

There is but one Catholic Church, whereof all
particular Churches are members, and therefore, when any spiritual
privilege or character is conferred on any particular Church, it must be
understood to extend over the whole Catholic Church: thus by Baptism,
men are not only made members of the particular Church where they happen
to be baptized, but of the Catholic Church over the world
and therefore whoever has been lawfully baptized in one
Church, has a right to partake of the LORD'S
Supper, and other Church privileges, in all other Churches, where he
happens to come; whereas, if baptism only admitted men into some
particular Church, they must be re-baptized, before they can lawfully be
received to communion in any diocese, where they have not been baptized
already.

If it was not thus in holy orders, that they who
have received them in one place, retain them in others; no minister
could have authority to preach the Gospel or to administer the
sacraments, or to exercise any other part of his functions beyond the
particular district in which he was ordained; the consequence whereas is
manifestly this, that the gospel of CHRIST
must not be propagated, nor any churches erected, in countries where
they had not stood even since the Apostles' times. For since there can
be no ministers without ordination, as was before proved, so then they,
who have been ordained in one country, may lawfully exercise their
respective functions in others, where there are no ordained ministers
already settled, or else those countries must remain for ever without
ministers, and consequently without sacraments and other public offices
of religion.

\textbf{NELSON,
CONFESSOR.—Festivals
and Fasts}

The Church being a regular society founded by CHRIST,
distinct from and independent of all other worldly societies, must
naturally make us suppose that He instituted some Officers for the
government of it … [The] Powers peculiar to the superior Order being
necessary for the good government of the Church, it is plain in fact
they did not expire with the Apostles. But, as our SAVIOUR
``glorified not Himself to be made an High Priest,'' but had His
commission from GOD
the FATHER,
so after His resurrection, He invested the Apostles with the same
commission His FATHER
had given unto Him: ``As My FATHER
hath sent Me, even so send I you: and He breathed on them, and said unto
them, Receive ye the HOLY
GHOST.''
In which commission is plainly contained the authority of ordaining
others, and a power to transfer that commission upon others, and those
upon others to the end of the world. And to show that it was not merely personal to the Apostles, our SAVIOUR
promises to be with them and their successors in the execution of this
commission, ``even unto the end of the world.'' ... This commission
the Apostles and their successors exercised in all places, and even in
opposition to the Rulers that then were; so that the Church subsisted as
a distinct society from the state, for above three hundred years, when
the civil government was only concerned to suppress and destroy it.
Indeed when the Church received the benefit of incorporation and
protection from the state, she was content to suffer some limitation as
to the exercise of these powers, and thought herself sufficiently
recompensed by the advantages that accrued to her by the incorporation.

\textbf{KETTLEWELL,
PRESBYTER
AND CONFESSOR.—Practical Believer, ii. 6.}

\emph{Question}. There remains yet one instance of
the Communion of the Primitive Christians, mentioned by St. Luke, viz.
their ``continuing in the Apostles' fellowship.'' (Acts ii. 42.) I
pray you what is meant by that?

\emph{Answer}. Owning their authority and
continuing under their government. They were appointed by CHRIST, as His deputies, to govern His Church;
and, therefore, to adhere to them, as the delegates of CHRIST, is called living ``in their
fellowship.''

\emph{Q}. But how can we live in their fellowship,
and adhere to their government, now they are dead?

\emph{A}. By adhering to and owning the authority
of our own Bishops, who are their successors, and rule the Church in
their stead.

\textbf{HICKS,
BISHOP AND
CONFESSOR.—Treatise
on the Episcopal Order, § 2.}

Can you, Sir, when you consider that Bishops are
appointed to succeed the Apostles, and, like them, to stand in CHRIST'S place, and exercise their
kingly, priestly, and prophetical office over their flocks; can you,
when you consider this, think it novel, or improper, or uncouth, to call
them spiritual princes, and their dioceses principalities?—when they
have every thing in their office which can denominate a
prince? For what is a prince but the chief ruler of a society, that hath
authority over the rest to make laws for it, to challenge the obedience
of all the members, and all ranks of men in it, and power to coerce
them, if they will not obey? And now, Sir, I pray you attend to what
follows, and then tell me, if the office of a Bishop contains not every
thing that is in the definition of a chief or a prince. St. Ignatius,
who was St. John's disciple, writes of the Bishop in his Epistle,
\&c.

\textbf{L\textsc{aw}.—Second
letter to the Bishop of Bangor}

``The priests of the sons of Levi shall come near;
for them hath the LORD
thy GOD
chosen to minister unto Him, and to bless in the name of the Lord.''
(Deut. xxi. 5.) Now, my LORD,
this is what we mean by the authoritative administration of the
Christian clergy: whether they be by way of benediction, or of any other
kind. We take them to be persons whom GOD
has chosen to minister unto Him, and to bless in His name. We imagine
that our SAVIOUR
was a greater priest and mediator than Aaron, or any of GOD'S former ministers. We are assured that CHRIST
sent His Apostles, as His FATHER
had sent Him; and that, therefore, they were His true successors; and
since they did commission others to succeed them in their office by the
imposition of hands, as Moses commissioned Joshua to succeed him, the
clergy who have succeeded the Apostles, have as divine a call and
commission to their work, as those who were called by our SAVIOUR; and are as truly His successors as the Apostles
themselves were.

\emph{Ibid}.—Postscript

The third objection against this uninterrupted
succession is this: that it is a popish doctrine, and ``gives Papists
advantage over us.'' The objection proceeds thus: ``We must not assert the
necessity of this succession, because the Papists say it is only to be
found in them.'' I might add, because some mighty zealous Protestants say
so too.

But if this be good argumentation, we ought not to
tell the Jews, or Deists, \&c., that there is any necessity of
embracing Christianity, because the Papists say, Christians can
only be saved in their Church. Again, we ought not to insist upon a true
faith, because the Papists say, that a true faith is only in their
communion. So that there is just as much Popery in teaching this
doctrine, as in asserting the necessity of Christianity to a Jew, or the
necessity of a right faith to a Socinian, \&c.

\textbf{JOHNSON,
PRESBYTER.—Unbloody
Sacrifice, Part II. Chap. 3.}

The Eucharist is one, as offered by priests, who
are one by their commission. It is very evident that it was not only our
SAVIOUR'S
intention, but His most passionate desire, that, as all His Apostles
received their commission from Him, so they might execute it with such a
harmony and consent of mind, that there might not be the least jarring
between them; for thus He prays for them; ``Keep through Thine own name
those whom Thou hast given Me, that they may be one, as We are.'' And the
foundation of our SAVIOUR'S
wishes and expectations for so perfect an union between His Apostles was
this, as is expressed by Himself, ``I have given them the Words which
Thou gayest Me,'' that is, He had committed to them the same treasures of
Divine truth which the FATHER
had before committed to him, \&c. … After His resurrection, He
does, with great solemnity, tell them, ``As My FATHER sent Me, even so send I you;'' from which words it is
evident, that the commission of all the Apostles was one and the same;
that it was such a commission as CHRIST Himself in His human nature, had received from His FATHER;
and even they who were not of the same order with the Apostles, but only
inferior Presbyters under them, yet by deriving their authority from the
same fountainhead, and exercising it in conformity to the instructions
which they received from them, they still kept the ``unity of the SPIRIT in the bond of peace.'' … It was upon this account that
Ignatius, Cyprian, and others, represent the whole college of Bishops
throughout the whole world as one person, sitting in one chair,
attending one altar; and that, therefore, is the one Eucharist which is
celebrated by this one priesthood: and St. Clement of Rome allows
nothing to be offered without the inspection of the high priest; and,
therefore, when a new altar is erected, a new Bishop ordained in
opposition to the former, then there is just occasion to ask that
question, as St. Paul did, ``Is CHRIST
divided?'' When two several pastors assume to themselves the privilege of
offering and consecrating the Sacrament not only in two distinct places,
but in contradiction to each other, and by two several inconsistent
claims; then it is evident, that one of them acts by no commission, for
if the true Eucharist can be had in two opposite assemblies, then CHRIST'S flesh ceases to be one.

\textbf{DODWELL,
CONFESSOR.—Discourse
on the one Priesthood, one Altar. Ch. 12.}

I observe that the Hierophanta, in their mysteries
represented a Divine Person. The same, in all probability, were the
thoughts of the primitive Christians, concerning their Bishops. This I
take to be the true design of that description of the Majestatic
Presence in the Revelations, to represent the Divine Presence, and
assistance in the Church, in as lively a way as possible, according to
the ways of Mystical Representation received in those times ... St. John
being particularly to affect the Churches he writes to, those of the
Lydian or Proconsular Asia, with a very feeling sense of the Divine
presence among them, (which might add the greater authority to his
several exhortations respectively,) he represents our SAVIOUR in a human visible shape; and
that the rest of the scene might be suitable, (that is, sensible also as
well as Himself,) he personates the Angels by their visible Bishops,
that so CHRIST
might be apprehended as present with the Bishops, as GOD was supposed to be wherever these Seven
Spirits were, which were peculiarly deputed to represent the Majestatic
Presence. This I take to be the reason why he confines his number, not
that by any geographical distinction those seven cities were
incorporated into a body, more than others of that province, but that he
had particular regard to that number of those Angels of the presence.
Therefore he makes seven candlesticks, alluding, as I have said, to the
like number of those in the tabernacle, as emblems of those seven
Churches. Therefore seven stars, alluding to the number of the Planets
and the Angels who presided over them, as emblems of the Bishops of
those Churches ... Thus it appears plainly, that the Bishops are
here represented in a mystical way; and how particularly suitable it
was, in this way, to personate them by the name of Angels. They were,
indeed, to perform the same office under CHRIST,
as a \emph{visible} human person, which the Angels were under Him as the
Logos, in reference to the restitution of souls to their original
dignity …

But because even His human nature, though visible
in itself, is yet invisible to us, therefore another way was thought of
for copying out that heavenly [\emph{telet\={e}}], even in the
ordinary external visible government of the Church. And here the Bishop
was to personate CHRIST
Himself, as the High Priest had, formerly, represented the Logos. The
seven deacons were to represent the Seven Mystical Angels, as I am very
apt to think, they were designed from the very original. I cannot think
it casual that the number first pitched on was exactly seven. But, that
which more confirms me in this opinion is the real suitableness of the
office of the Deacon to the Bishops, as representing the Logos in a
visible way, with that of those Angels to the same Logos, as He was
invisible. The office of the Angels in general is thus described, by the
Author of the Hebrews, that they are [`` ministering spirits, sent out
for a diaconate.''] These are exactly the very terms by which the Church
would have expressed the office of these Deacons, if she had been to
have described the same office as vested in mortal men … They (the
Angels) were to stand before the presence of GOD,
in a posture of readiness to be sent on messages by Him; and so were the
Deacons to stand before the Bishop, to be sent by him on his messages.
They were the ``eyes of the LORD
which ran to and fro through the whole earth.'' So also the Deacons are,
in the language of the Ancient Church, called the \emph{Oculi Episcopi},
for the same reason ... Now we may not wonder why the Bishops are called
Angels, in the forementioned mystical immediate relation to our SAVIOUR Himself as the chief ``Bishop of our
souls;'' because, indeed, in regard of Him, they bear no higher office
than that of Deacons … Accordingly the Primitive Church were extremely
rigorous in insisting on this very number of their Deacons, in all
places, as I have elsewhere showed. The council of Neocæsarea
imposed it as a universal rule, how great soever the Church were to
which the Deacons were to serve; ... a canon, which, though it were at
first designed only for their own province of Cappadocia, was,
notwithstanding, afterwards extended, first, to the Eastern Empire ...
afterwards to the Western ... Therefore, even then it is much more
probable that this number was already received in more Churches than
otherwise.

And now the comparisons of the Bishops in Ignatius
cannot seem so strange, these things being considered, as they did to
Blondell, who had considered none of them. They are generally designed
to express the sacredness and excellency of the persons which the clergy
bore in these mystical performances. Nor is there any thing in them that
is really affected or strained, much less blasphemous, no, nor any
extravagant flights of fancy … lf he were to compare them with the
first invisible archetypes of unity, (as that is, indeed, his great
design in those epistles, in opposition to the schisms then rising,)
then it was very proper for him to take notice only of the two orders
which were then immediately concerned in the office of ministration, and
then to compare them with GOD the Father, and the Logos; because as this
unity consists in the unity of the head, and the Scripture tells us that
the head of every man is CHRIST,
so also the same Scripture tells us that the Head of CHRIST is GOD ... These things, therefore, being thus
solidly laid down by the first fathers, in their disputes against their
contemporary Heretics and Schismatics, all the inferences thence deduced
against them, will follow naturally and undeniably ... It will follow,
that disunion from the Bishop was a disunion from CHRIST and the FATHER, and from all the invisible heavenly
Priesthood, and sacrifice, and intercession. It will follow that
disunion from any one ordinary, must consequently be a disunion from the
whole Catholic Church, seeing it is impossible for any, to continue a
member of CHRIST'S
mystical body, who is disunited from the mystical head of it. It will
follow that visible disunion from the external sacraments of the Bishop,
is in the consequence a disunion from the Bishop, and from the
whole Catholic Church in communion with him, who ought to ratify each
other's censures under pain of schism if they do not.

\textbf{COLLIER,
BISHOP AND
CONFESSOR.—Moral
Essays, Part III.}

'Tis the bulk and serviceableness of business, and
the use it has in the world, which makes an employment honourable. And
can any thing compare with the Apostles in this particular? Were they
not to form and instruct the Church, and to govern the most noble
society upon earth? Were they not to publish the Mysteries of
Redemption, the offers of the New Covenant, and the glories of the other
world … Fire in the figure of tongues sat upon the heads of each of
them. This was an emblem of the gift of languages, and the miracle was
as bright as the flame … This was a glorious attestation, this must
needs make their Commission undisputed, and their character indelible.
Should a Prince be proclaimed from the sky, anointed out of the Ampoul,
and crowned by an Angel, his authority could not be more visible ... I
can't help saying, that, in my opinion, a Prince made but a lean figure
in comparison with an Apostle. What is the magnificence of palaces, the
richness of furniture, the quality of attendance, what is all this to
the pomp of miracles, and the grandeur of supernatural power? … A
Prince can bestow marks of distinction, and posts of honour and
authority; but he can't give the HOLY GHOST, he can't register his favourites among
the quality of heaven, nor entitle them to the bliss of eternity.
No,—these powers were Apostolic privileges, and the enclosure of the
Church. The prerogative royal cannot stretch thus far; these jewels are
not to be found in the imperial crown ... I need not tell you how much
they suffered through their progress, and how gloriously they went off
into the other world. But before their departure they took care to
perpetuate their authority, and provide governors for the Church. Thus
the jurisdiction was conveyed to Bishops and Priests; this succession
has continued without interruption for above sixteen hundred years.

\textbf{LESLIE,
PRESBYTER
AND CONFESSOR.—Case
of the Regale and Pontificate}

When any constitution of civil government dissolves
itself, another immediately succeeds; or, if a Monarchy be turned into a
Commonwealth, or a Commonwealth into a Monarchy: and consequently that
which was dissolved, is no more; but we cannot say that the Church is no
more. There is still a Church, though in servitude, and nothing succeeds
to it; if it were dissolved, there would be no Church, but nothing would
come in its room, unless you will say a privation, that is, the want of
a Church … The Church is a society spread over the earth; and,
therefore, cannot be dissolved in any one kingdom or state; nor can the
concessions of any national Church oblige the Church Catholic; no, nor
oblige that national Church herself, otherwise than according to the
rules of the Catholic Church; more than a Committee of the House of
Lords or Commons can oblige the whole House, or govern themselves by any
other rules than those which are prescribed by the House.

The Church is laid as low and fenceless as the sand
under their [Atheism, Deism, \&c.] storms, which had long since
overwhelmed the City of GOD,
(after the change of her governors) if the Almighty promise (Matt. xvi.
18; xxviii. 20.) had not interposed to preserve some embers alive in the
midst of these torrents. And they will be preserved till the time
appointed by GOD
shall come, when His breath shall put new life in them, to lick up that
sea that now covers, but cannot drown them ... This is the city, the
society, over which the temporal governments of the earth have assumed
the dominion; and have said, ``Let us break their bonds asunder, and cast
away their cords from us.'' …

And let not so weak a thought arise in your minds,
as if all this were only the self-seeking of the Clergy, out of pride to
advance themselves. Alas! it must have the quite contrary effect with
any of them who consider what a heavy charge they have undertaken, and
what account will be exacted from them, for their faithful discharge of
it! That the blood of all those souls who perish through their
negligence or default, will be required at their hands! That they
have to wrestle, not only with flesh and blood, but against
principalities and powers, against the rulers of the darkness of this
world, against wicked spirits that are set up in high places! And
whoever opposes these with that truth and freedom that is necessary,
instead of honour must expect reproach and persecution; of which it is
not the least that they cannot vindicate the honour of CHRIST'S commission without being thought to seek their own glory.
Yet that must not hinder; the successors of the Holy Apostles must be
content to pass, as they did, ``through evil report and good report, as
deceivers, and yet true.''

\textbf{WILSON,
BISHOP,
CONFESSOR
AND DOCTOR.—Private
Thoughts}

``He that entereth not by the door into the
sheepfold,'' \&c. A lawful entrance, upon motives which aim at the
glory of GOD
and the good of souls; an external call and mission, from the apostolic
authority of Bishops.

``A stranger will they not follow;'' that is, they
ought not to follow such as break Catholic Unity …

Whoever is associated in the Priesthood of CHRIST,
ought, in imitation of Him, to sacrifice himself for the advantage of
His Church and for all the designs of GOD
…

``Bishops and Priests,'' saith St. Ambrose, ``are
honourable on account of the sacrifice they offer.'' The power of the
keys and the exercise of that power, the due use of confirmation, and
previous to that of examination ... are matters of infinite and eternal
concern … (\emph{At the Lord's Supper}. \emph{Before the Service begins}.)
May it please Thee, O GOD, who hast called us to this ministry, to
make us worthy to offer unto Thee this sacrifice for our own sins and
for the sins of the people. Accept our service and our persons, through
our LORD
JESUS CHRIST,
who liveth and reigneth with Thee and the HOLY GHOST,
One GOD,
world without end.—O reject not this people for me and for my sins.
Amen.

(\emph{Upon placing the Elements upon the Altar}.)
Vouchsafe to receive these Thy creatures from the hands of us sinners, O
Thou self-sufficient GOD!

(\emph{Immediately after the Consecration}.) We
offer unto Thee, our King and our GOD,
this bread and this cup. We give Thee thanks for these and for all Thy
mercies, beseeching Thee to send down Thy HOLY SPIRIT
upon this sacrifice, that He may make this bread the Body of Thy CHRIST,
and this cup the blood of Thy CHRIST;
and that all we, who are partakers thereof, may thereby obtain remission
of our sins and all other benefits of His passion. And together with us,
remember, O GOD,
for good, the whole mystical body of Thy Son; that such as are yet alive
may finish their course with joy, and that we, with all such as are dead
in the LORD,
may rest in hope and rise in glory, for Thy Son's sake, whose death we
now commemorate. Amen. May I adore Thee, O GOD,
by offering to Thee the pure and unbloody sacrifice, which Thou hast
ordained by JESUS
CHRIST.
Amen.

Whenever Church discipline meets with
discountenance, impieties of all kinds are sure to get head and abound.
And impieties unpunished do always draw down judgments. The same JESUS CHRIST, who appointed baptism for the receiving
men into His Church and family, has appointed excommunication, to shut
such out as are judged unworthy to continue in it … If baptism be a
blessing, excommunication is a real punishment; there being the same
authority for excommunication as for baptism. And if men ridicule it,
they do it at the peril of their souls.

\textbf{BINGHAM,
PRESBYTER.—Sermons
on Absolution. No. 2.}

In the first place, the commission of power to
ministers to retain and remit other men's sins, in whatever sense we
take it, is a great engagement on them to lead holy and pure lives
themselves. For it looks like an absurdity in practice, and is too often
really thought so, that men should be qualified to forgive other men's
sins, who are loaded with guilt and impurity themselves. There is
nothing so natural and obvious to us as, Physician, heal thyself: and,
therefore, if it be not a real objection against their office, yet it is
an unanswerable one against their persons. If it do not destroy the
tenor of their commission in the nature of the thing, yet it certainly
diminishes their authority and reputation in the opinion of men when
every profligate sinner can retort upon them and say, ``Thou that
teachest another, teachest thou not thyself? Thou that preachest a man
should not steal, dost thou steal? Thou that sayest a man should not
commit adultery, dost thou commit adultery? Thou that makest thy boast
of the law, through breaking the law, dishonourest thou GOD?''
It must needs take off very much from the veneration of the Sacrament of
Baptism, to have a man pretend to wash away the sins of others, who is
himself polluted and profane; and equally diminishes the reverence which
is due to the tremendous mystery of the Eucharist, to have it ministered
with unholy hands. It cannot relish well with men, to hear an
unsanctified mouth giving blessing to others, who in effect is cursing
himself; praying that the blood of CHRIST
may preserve others to eternal life, while he himself is eating and
drinking his own damnation, not discerning the LORD'S body. But above all, such a man cannot with any tolerable
decency or freedom, discharge the office of punishing and correcting
others, who is himself more justly liable to rebuke and censure. With
what face can he debar others from Baptism or the Eucharist, who is
himself unqualified to receive either? or exclude others from the
Church, who is himself unworthy to enter into it? Nothing can be a
greater engagement upon Ministers to lead holy and pure lives, than the
consideration of the commission which CHRIST
has given them, to retain or remit other men's sins, whether in a
sacramental way, or a declaratory way, or a precatory way, or a judicial
way: because without purity, they can by no means answer the end of this
office, and the nature of their trust; but their mal-administration will
rise up in judgment against them and condemn them.

2. A second thing which this office of retaining
and remitting sins requires of Ministers, is great diligence in their
studies and labours, without which they can never be able sufficiently
to discharge it. The Church, indeed, has made some part of this work
tolerably easy, by a prudent provision of many proper general forms of
absolution: to which in her wisdom she may add proper forms of
excommunication and judicial absolution. But when this is done, there
still remains a great deal more belonging to the full discharge of
this office, for which the Church can make no particular provision: and,
therefore, that must be left to the industry and diligence of Ministers,
in their particular studies and labours. And this requires both a
diffused knowledge, and great application; to be able to understand the
nature of all GOD'S laws, and the bounds and distinctions betwixt every virtue
and vice; to be able to resolve all ordinary cases of conscience, and
answer such doubts and scruples as are apt to arise in men's minds; to
know the qualifications of particular men, and the nature and degrees
and sincerity of their repentance, in order to give them a satisfactory
answer to their demands, and grant or refuse them the several sorts of
absolution, as they shall think proper, upon an impartial view of their
state and condition. He that thinks all this may be done without any
great labour and study, and a diligent search of the Holy Scriptures,
the rule and record of GOD'S will, scorns neither to understand the
nature of his office, nor the needs of men; nor what it is to stand in
the place of CHRIST,
and judge for him between GOD and man. The Priest's lips should preserve knowledge: and a
man that considers the large extent of that knowledge, together with the
great variety of cases and persons to which he may have occasion to
apply it, would rather be tempted to cry out with the Apostle, ``Who is
sufficient for these things?'' And if this be not an argument to engage a
man to industry in the office of a spiritual physician, it is hard to
say what is so.

\textbf{SKELTON,
PRESBYTER.—Discourse
71.}

The next thing the Puritans took offence at, was
the Hierarchy of the Church. They looked on the Bishops as the
instruments of papal tyranny, and the corrupters of true religion
...They were, it seems, so ignorant, as not to know, that the Bishops,
of all men, had most reason to oppose the usurpation of the Bishop of
Rome, who had made himself the only Bishop, and reduced all the rest to
ciphers. Nor did they consider, whether it was in the power of man to
abolish, at his discretion, an order of the Church, instituted by GOD
Himself, merely because the men who filled this order had degenerated,
together with all the rest of the Church, into superstition and
luxury. Here again the scheme of our opposers was not to reform, but to
destroy; and what was equally bold, to begin a new ministry, with hardly
any other mission than such as a number of men, and sometimes one man
only, wholly unauthorized, for aught that others could perceive, should
assume. From men thus sending themselves, or sent by we know not whom,
we are to receive the sacraments … We must not forget, however, that
these new orders lay claim to scriptural institution and primitive
example. What, all of them? And without succession? Do we hear of any
man in Scripture who ordained himself, or who presumed to take the
ministry of GOD'S
word and sacraments upon him, without being sent either immediately or
successively by CHRIST?
Or, can an instance of this kind be assigned during the first fourteen
centuries of the Church?

So sacred a thing is the succession of ordination,
that the HOLY
GHOST,
who had already enabled Barnabas and Saul to preach the word, ordered
them to be ``separated for the work whereunto He had called them, by
fasting, prayer, and imposition of hands,''—that is, to be ordained:
the SPIRIT of GOD
hereby plainly showing, that He himself would not break the successive
order of mission established in the Church.

\textbf{SAMUEL
JOHNSON.—[\emph{en
philosophou schemati presbeu\={o}n ton theion logon}]. Sermon 7.}

With regard to the order and government of the
Primitive Church, we may doubtless follow their [the ancients']
authority with perfect security; they could not possibly be ignorant of
laws executed, and customs practised by themselves; nor would they, even
supposing them corrupt, serve any interest of their own, by handing down
false accounts to posterity. We are, therefore, to inquire from them the
different Orders established in the Ministry from the Apostolic ages;
the different employments of each, and their several ranks,
subordinations, and degrees of authority. From their writings, we are to
vindicate the establishment of our Church, and by the same writings are
those who differ from us in these particulars to defend their conduct.

Nor is this the only, though perhaps the chief use
of these writers; for, in matters of faith and points of doctrine,
those at least who lived in the ages nearest to the times of the
Apostles undoubtedly deserve to be consulted. The oral doctrines and
occasional explications of the Apostles would not be immediately
forgotten, in the Churches to which they had preached, and which had
attended to them with the diligence and reverence which their mission
and character demanded. Their solution of difficulties, and
determinations of doubtful questions, must have been treasured up in the
memory of their audiences, and transmitted for some time from father to
son. Every thing, at least, that was declared by the inspired teachers
to be necessary to salvation, must have been carefully recorded; and,
therefore, what we find no traces of in Scripture, or the early Fathers,
as most of the peculiar tenets of the Romish Church, must certainly be
concluded to be not necessary. Thus, by consulting first the Holy
Scriptures, and next the writers of the Primitive Church, we shall make
ourselves acquainted with the will of GOD;
thus shall we discover the good way, and find that rest for our souls,
which will amply recompense our studies and inquiries.

\textbf{HORNE,
BISHOP AND
DOCTOR.—Charge
of Primary Visitation of his Diocese}

The Constitution and use of the Church of CHRIST
is another subject, on which our principles, for some years past, have
been very unsettled, and our knowledge precarious and superficial.
Ignorance is dangerous here, because there are so many whose interest it
is to flatter us in it, and take advantage of it. The definition of the
Church, contained in our Articles, was purposely less definitive than it
might have been, to avoid giving further offence to those whom we rather
wished to reconcile; but it does not appear, that the Church hath gained
any thing by its moderation; it hath rather lost; because in virtue of
that moderation, it hath been pleaded against us, that Ecclesiastical
Unity may be dispensed with, and that all our differences in this matter
are only problematical and immaterial.

But salvation is a gift of grace; that is, it is a
free gift, to which we have no natural claim. It is not to be conceived
within ourselves, but to be received, in consequence of our
Christian calling, from GOD
Himself, through the means of His Ordinances. These can no man
administer to effect, but by GOD'S
own appointment; at first, by His immediate appointment, and afterwards,
by succession and derivation from thence to the end of the world.
Without this rule we are open to imposture, and can be sure of nothing;
we cannot be sure that our ministry is effective, and that our
Sacraments are realities. We are very sensible the spirit of division
will never admit this doctrine; yet the spirit of charity must never
part with it. Writers and teachers who make it a point to give no
offence, treat these things very tenderly; but he who, in certain cases,
gives men no offence, will for that reason give them no instruction. \emph{Light
itself is painful to weak eyes; but delightful to them when grown
stronger, and reconciled to it with use; and he who was instrumental in
bringing them to a more perfect state of vision, though less acceptable
at first, may yet, for his real kindness, be more cordially thanked
afterwards, than if he had made the case and safety of his own person
the measure of his duty}. It is by no means evident, that the Church
hath ever recommended itself the more by receding from any of its just
pretensions: generosity obliges and secures a friend; but an enemy
construes it into weakness, and then it never does any good. Yet the
adversaries of the Church of England have always been persuading her to
make the experiment, and have promised great things from it; with what
views, it cannot be difficult to discover. It was an unhappy
circumstance, and had very ill effects, when some pious men \footnote{Mr. Wesley, \&c.}, of more zeal than discretion, who set out on the work of
reforming this nation, opened an asylum for penitents, which took in
people of all persuasions, without exception of any. It came to be
inferred from hence, that souls might be saved as well without as with a
Church; perhaps better; and when men have once begun to neglect roles,
they go on to despise them, and know not where to stop, till all things
are bought into confusion …

The ancient Church is the standard by which all
modern ones are to be examined; and unless a man knows what the
Church was in centuries before the Reformation, he will see but darkly
into the troubled waters of later times, in which faction and party have
confounded things; and it hath become as much the interest of some, that
the Church of CHRIST
should be found every where, as it is the desire of others that it
should be found no where … If we would guard against popular mistakes
in the subject at large, it will be necessary to examine first, what the
Church was under the Old Testament; for there we find its original
establishment, its form, its authority, its ministry, its unity and
uniformity, its maintenance, its independence; which things being so
particularly laid down, no new establishment is to be found in the
Epistles or the Gospels of the New Testament, but the ancient
constitution is referred to, to show us, in certain cases, what ought to
be from what had been … From the Scripture we should proceed next to
observe, what the Church was in the first ages of the Gospel, before
worldly policy, miscalled moderation, had any influence upon the
opinions of Christians. There is an epistle of St. Clement, on Church
unity and Church authority, with which all students in divinity should
be acquainted. It will teach them what the Christian society then was,
and what it ought to be. Ignatius and Cyprian, both of them martyrs,
will give further instruction. The latter is so particular and copious,
that a code of discipline might nearly be formed upon his authority.
With these preparations, we shall be the better able to judge, of what
happened at the Reformation, when many things were right and many wrong;
when the Church of England, by the singular blessing of GOD, preserved its constitution and its
doctrines, while many of the reformed fell off by degrees, some into
disorder, some into dissolution. What remains with us we must defend and
preserve; trusting that the same GOD
who hath raised this Church, when trodden down to the dust, will never
forsake us till we forsake Him ...

But I must now hasten, in the last place, to a
subject of more quietness or less suspicion [than the subject of civil
government], in which wise men of all persuasions are more nearly of a
mind; I mean, the conduct of the Christian life. Modern times and new modes of education have given too great a latitude in the articles
of dress, and dissipation, and self-indulgence. Every thing is to be
avoided which tends to diminish that gravity and seriousness which GOD expects to find in all those who are flying from the wrath to
come. It was observed of old, that when inconsiderate people are
avoiding one extreme, they commonly fall into another, while reason and
discretion keep the middle way. When Protestants laid aside the
austerities of superstition, they began to see less harm in the
liberties taken by the world. The kind of life to which the first
Christians conformed, hath been considered as a sort of heroic piety,
which had more of suffering and mortification than are now required of
us; as if the way to heaven could be easier, while the number of our
temptations is probably increasing from the refinement of modern times,
which, instead of giving us more liberty, call upon us for a greater
degree of caution and reserve …

To us JESUS CHRIST
is the pattern of holiness, the great exemplar of perfection, of whom we
are first to learn, what no heathen ever professed, to be ``meek and
lowly in heart;'' and accordingly, one of the best books extant on the
Spiritual life, is entitled, ``The Imitation of JESUS CHRIST.''
Its language is barbarous, but its matter is divine and heavenly, and
hath administered instruction and consolation to thousands of devout
Christians. The way of true devotion must still be understood to be the
same humble, secret, unaffected, unaspiring practice of piety, as it
used to be of old. The Cross, which JESUS
CHRIST
carried for our salvation, is still the true emblem of our profession,
from our baptism to our departure out of this life, and is to be borne
by us in our minds, as a daily admonition to patient suffering and
self-denial.

To assist us in the great duties of prayer and
meditation, books of devotion have their use; but to us of the clergy,
the liturgy of our Church is the best companion; and the daily use of it
in our churches and families is required by the canons. It cannot be
denied, that from various reasons prevailing amongst us, we are much
fallen off of late years, from the practice of weekly prayers in our
churches. Wherever this has been neglected, we should exhort the
people to the revival of it, if circumstances will possibly permit; and
alarm them against a mistake, to which they are all exposed, from a
fanatical prejudice of baneful influence, namely, that they come to
Church only to hear preaching; and hence they are indifferent, even on a
Sunday, to the prayers of the Church, unless there is a sermon.

\textbf{JONES
OF NAYLAND,
PRESBYTER.—Lectures
on Hebrews iii.}

The Church, in its nature, always was what it is
now, a society comprehending the souls as well as the bodies of men;
and, therefore, consisting of two parts, the one spiritual, answering to
the soul, and the other outward, answering to the body. Hence, some have
written much upon a visible Church and an invisible, as if they were two
things; but they are more properly one, as the soul and body make a
single person.

In the twelfth chapter of the Epistle to the
Hebrews, the Apostle gives such a description of that society, into
which Christians are admitted, as will show us the nature of it. ``Ye are
come,'' says he, ``unto Mount Zion,'' \&c. … The terms here used give
us a true prospect of the Church … This is that Zion of the Holy One
of Israel, to which the forces of the Gentiles were to flow from all
parts of the world ... the city of the living GOD,
distinguished from the cities of the world, as Jerusalem was from the
cities of the heathens, who dedicated their cities not to the living GOD, but to the names of their dead idols …
This, being the city of the living GOD, must be an immortal society, for the living GOD
does not preside over dead citizens; He is not the GOD of the dead, but the GOD of the living, and all the members of this
society live unto Him … It is, therefore, called the Heavenly
Jerusalem, because it is of a heavenly nature; and it is called the
Jerusalem which is above, which is free, and is the mother of us all ...
Its spiritual nature is further declared, in that it is said to
comprehend an innumerable company of angels … In the communion of the
Church the spirits of just men made perfect are also included. It is a
society which admits only the spirits of the living, and as such cannot
exclude the spirits of the dead; and this confirms what we said above, that the Church is a spiritual community, comprehending the dead
as well as the living …

But it is now to be shown, secondly, that as the
Church of GOD
hath always been the same in its nature, it hath likewise preserved the
same form in its external economy: the wisdom of GOD having so ordained, that the Christian Church under the Gospel
should not depart from the model of the Church under the Law. For as the
congregation of Israel was divided into twelve tribes, under the twelve
Patriarchs, so is the Church of CHRIST founded on the twelve Apostles, who raised to themselves a
spiritual seed amongst all the nations of the world … There were then
three orders of priests in the Jewish Church; there was the high priest,
and the sons of Aaron and the Levites. In the Church of CHRIST, there was the order of the Apostles,
besides whom there were the seventy Disciples sent out after them; and,
last of all, the Deacons were ordained to serve under both in the lower
offices of the Church. The same form is still preserved in every regular
Church of the world, which derives its succession and authority from the
Church of the Apostles: after whom the Bishops succeeded by their
appointment, such as Timothy and Titus, in their respective churches.
This authority has been opposed to the Christian as it was in the Jewish
Church: Corah and his company rose up against Moses and Aaron for
usurping a lordly authority over the people; so, in the later ages of
the Christian Church, a levelling principle hath prevailed, which has
appeared in many different shapes …

The Church has also been remarkably conformable to
itself in its sufferings. There never was a time, so far as we can
learn, when the true Church of GOD,
with its doctrines and institutions, was not hated and opposed by the
world: either persecuted and oppressed by powerful tyrants, or traduced
and insulted by lying historians.

\textbf{BISHOP}[sic]\textbf{.—Sermon
on Matt. xvi. 18, 19.}

The keys of the kingdom of Heaven here promised to
St. Peter ... must be something quite distinct from that with which it
hath generally been confounded, the power of remission and retention of
sins, conferred by our LORD,
after His resurrection, upon the apostles in general, and
transmitted through them to the perpetual succession of the priesthood.
This is the discretionary power lodged in the priesthood, of dispensing
the sacraments, and of granting to the penitent and refusing to the
obdurate the benefit and comfort of absolution. The object of this power
is the individual upon whom it is exercised, according to the particular
circumstances of each man's case. It was exercised by the Apostles in
many striking instances. It is exercised now by every priest, when he
administers or withholds the sacraments of baptism and the LORD'S
Supper, or, upon just grounds, pronounces or refuses to pronounce upon
an individual the sentence of absolution.

\textbf{HEBER,
BISHOP.—Sermons
in England, No. xii.}

We must return then, after all, (in ordinary cases,
and where an immediate and supernatural commission from the HOLY
GHOST is
neither proved nor pretended,) to the appointment and ordination of
those among our fellow-creatures who exercise a legitimate authority in
the Church of CHRIST,
and who, as being appointed by GOD,
are placed in GOD'S
stead, and commissioned by Him to dispense those graces which are
necessary for the feeding of His flock, and to designate those labourers
who are henceforth to work in His harvest.

And having arrived at this point of the discussion,
even if that discussion were to proceed no further, and if the
Scriptures had given us no information as to the persons by whom this
authority was to be exercised, the validity of our ordinations would
still be sufficiently plain, and the danger of separation from, or
rebellion against, our Church would be sufficiently great and alarming,
inasmuch as, where no distinct religious officer was instituted by GOD, the appointment of such officers must necessarily have
devolved on the collective Christian Church, and on those supreme
magistrates who, in every Christian country, are the recognized organs
of the public will and wisdom … It happens, however, to be in our
power to show (if not an explicit direction of CHRIST
for the form of our Church government, and the manner of appointing our
spiritual guides), yet a precedent so clear and a pattern so
definite, as to leave little doubt of the intentions of our Divine
Master, or of the manner in which those intentions were fulfilled by His
immediate and inspired Disciples. Nor will the force of such precedent
and example on the practice of succeeding Christians be regarded as
trifling by those who consider that it is on such grounds as these that
the obligation rests of many observances which are allowed by all
parties to be essential; among which maybe classed the baptism of
infants, the observance of the LORD'S Day, and our participation in the LORD'S Supper.

But, without entering into the question of the
absolute necessity of this rule, and without judging those other
national Churches which have departed from it, it is evident that those
Churches are most wise and most fortunate, who have continued in the
path which CHRIST and His Apostles have trodden before;
and that religious insubordination is then most unreasonable and most
dangerous, when exerted against a form of polity which the majority of
our fellow-Christians, the wisdom of our civil governors, and the full
stream of precedent, from the time of the Apostles themselves, combine
to recommend to our reverence.

We find, accordingly, that our LORD, on His own departure from the world,
committed, in most solemn terms, the government of His Church to His
Apostles. We find these Apostles, in the exercise of the authority thus
received, appointing Elders in every city, as dispensers of the word and
the sacraments of religion; and we find them also appointing other
Ecclesiastical Officers, who were to have the oversight of these Elders
themselves; and who, in addition to the powers which they enjoyed in
common with them, had the privilege, which the others had not, of
admitting, by the imposition of hands, those whom they thought fit to
the ministerial office …

And it is not too much to say, that we may
challenge those who differ from us, to point out any single period at
which the Church has been destitute of such a body of officers, laying
claim to an authority derived by the imposition of hands from the
Apostles themselves; or any single instance of a Church without this
form of government, till the Church of Geneva, at first from necessity,
and afterwards from a mistaken exposition of Scripture, supplied
the place of a single Bishop by the rule of an oligarchical presbytery.

\textbf{JEBB,
BISHOP.—Pastoral
Instructions. Discourse i.}

``And lo, I am with you alway, even unto the end of
the world;''—a promise not occasional or temporary, like that of
miraculous powers; but conveying an assurance that CHRIST Himself, will, in spirit and in power, be continually
present with His Catholic and Apostolick Church; with the bishops of
that Church, who derive from the Apostles by uninterrupted succession,
and with those inferior, but essential orders of the Church, which are
constituted by the same authority, and dedicated to the same service.

\textbf{VAN
MILDERT,
BISHOP.—Bampton
Lectures. Sermon viii.}

The system, of which the Apostles had laid the
foundation, was to be carried on through succeeding generations; but
with a gradual diminution of that extraordinary aid, which the
circumstances of the case rendered no longer necessary ... Yet since the
object to be attained was not temporary, but to continue from age to
age, the mode, the form, and the instrument to be employed, were still
to be conformable to the primitive institution. Accordingly, the
Apostles ordained successors to themselves, and took measures for
perpetuating in the Church a standing ministry of diverse orders and
gradations. In so doing, they showed in what sense we are to interpret
our LORD'S
assurance, that He would ``be with them always, even unto the end of the
world.'' …

The evidences, from the best historical records, to
the simple fact that a visible Church of this description has actually
subsisted from the time of our LORD
and His Apostles to this moment, are too well known to require a detail.
Nor is there any defect of similar evidence, to show that, whatever
errors or corruptions may have occasionally found admittance into it,
the Church itself has proved a successful instrument in the hands of
Providence, both of transmitting the unadulterated Word of GOD from generation to generation, and also of
promulgating and maintaining all its great fundamental truths; nay,
perhaps, of preserving even the very name as well as substance of
Christianity, which, humanly speaking, would probably have been long
since extinct, had it not been nurtured and cherished by this its
appointed guardian and protector …

Let us take, for instance, those articles of faith
which have already been shown to be essential to the Christian Covenant:
—the doctrines of the Trinity, of our LORD'S Divinity and Incarnation, of His Atonement and
Intercession, of our sanctification by the HOLY SPIRIT, of the terms of acceptance, and the
ordinances of the Christian Sacraments and Priesthood. At what period of
the Church have these doctrines, or either of them, been by any public
act disowned or called in question! We are speaking now, it will be
recollected, of what in the language of Ecclesiastical history, is
emphatically called THE
CHURCH;
that, which has from age to age borne rule, upon the ground of its
pretensions to Apostolical succession. And to this our inquiry is
necessarily restricted …

Surely, here is something to arrest attention;
something to awaken reflection; something which they who sincerely
profess Christianity, and are tenacious of the inviolability of its
doctrines, must contemplate with sentiments of awe and veneration. For,
though a sceptic may contend that this species of evidence does not
amount to a direct and demonstrative proof of the truth of the
doctrines; yet if they be not true, how shall we account for their
having been so uninterruptedly transmitted to these latter times? How
they have withstood the assaults of continued opponents? opponents,
wanting neither talents nor inclination to effect their overthrow? If
these considerations be deemed insufficient, let the adversary point out
by what surer tokens we shall discover any Christian community, duly
answering the Apostle's description, that it is ``built upon the
foundation of the Apostles and Prophets, JESUS CHRIST Himself being the chief Corner-Stone?''

\textbf{MANT,
BISHOP.—Parochial
Sermons, xxvii.}

Nor had He in this appointment a view to those
times only, in which the appointment was made; but He designed that it
should be extended to all future ages; for so we must understand the
words which He pronounced immediately after giving His apostles their
authority to baptize: ``Lo, I am with you always, even unto the end of
the world.'' A promise this which cannot be supposed to have respect to
the persons of the Apostles alone, who in the common course of nature
were soon to be taken from the world, to the end of which the promise
itself was to extend.

In conformity with this meaning, the Apostles, who
were themselves holy men and full of the HOLY GHOST,
did send other persons; to whom again, they gave power and authority to
send others, through whom the office of ministers of the Gospel has been
handed down in regular and uninterrupted succession from the Apostles to
the present time.

OXFORD,

\emph{The Feast of St. Mark}

[NEW EDITION.]

———————————————————————

These Tracts are continued in
numbers, and sold at the price of 2d. for each sheet, or 7s. for 50
copies.

LONDON: PRINTED FOR
J. G. F. \& J. RIVINGTON,

ST. PAUL'S CHURCH YARD AND WATERLOO PLACE.

1839.

\xchapter{Introduction}

On the history of the Breviary
\footnote{The authorities used in this account are
Gavanti's Thesaurus Rituum, cum notis Merari; Zaccaria's
Bibliotheca Ritualis; and Mr. Palmer's Origines Liturgicæ.}

The word \emph{Breviarium} first occurs in the work
of an author of the eleventh century, and is used to denote a compendium
or systematic arrangement of the devotional offices of the Church. Till
that time they were contained in several independent volumes according
to the nature of each. Such, for instance, were the \emph{Psalteria}, \emph{Homilaria},
\emph{Hymnaria}, and the like, to be used in the service in due course.
But at this memorable era, and under the auspices of the Pontiff who
makes it memorable, Gregory VII., an order was drawn up, for the use of
the Roman Church, containing in one all these different collections,
introducing the separate members of each in its proper place, and
harmonizing them together by the use of rubrics. Indeed, some have been
led to conclude that in its first origin the word Breviary was
appropriated to a mere collection of rubrics, not to the offices
connected by them. But even taking it in its present sense, it will be
obvious to any one who inspects the Breviary how well it answers to its
name. Yet even thus digested, it occupies four thick volumes of
duodecimo size.

Gregory VII. did but restore and harmonize these
offices; which seem to have existed more or less the same in their
constituent parts, though not in order and system, from Apostolic times.
In their present shape they are appointed for seven distinct seasons in
the twenty-four hours, and consist of prayers, praises, and
thanksgivings of various forms; and, as regards both contents and hours,
are the continuation of a system of worship observed by the Apostles and
their converts. As to contents, the Breviary Services consist of the
Psalms; of Hymns, and Canticles; of Lessons and Texts from inspired and
ecclesiastical authors; of Antiphons, Verses and Responses, and
Sentences; and of Collects. And analogous to this seems to have been the
usage of the Corinthian Christians, whom St. Paul blames for refusing to
agree in some \emph{common} order of worship; when they came together, \emph{every
one of them} saying a Psalm, a doctrine, a tongue, a revelation, an
interpretation \footnote{Vid. Socr. Hist. 22. Vide
also Lyra Apostolica, xv.}. On the
other hand, the Catholic seasons of devotions are certainly derived from
Apostolic usage. The Jewish observance of the third, sixth, and ninth
hours for prayer, was continued by the inspired founders of the
Christian Church. What Daniel had practised, even when the decree was
signed forbidding it, ``kneeling on his knees three times a day, and
praying, and giving thanks unto his GOD,'' St. Peter and the other Apostles were solicitous in
preserving. It was when ``they were all with one accord in one
place,'' at ``the \emph{third} hour of the day,'' that the HOLY GHOST
came down upon them at Pentecost. It was at the \emph{sixth} hour, that
St. Peter, ``went up upon the house-top to pray,'' and saw the vision
revealing to him the admission of the Gentiles into the Church. And it
was at the \emph{ninth} hour that ``Peter and John went up together
unto the temple.'' being ``the hour of prayer.'' But though these
were the more remarkable seasons of devotion, there certainly were
others besides them, in that first age of the Church. After our SAVIOUR'S departure, the Apostles, we are
informed, ``all \emph{continued} with one accord in prayer and
supplication, with the women, and Mary the mother of Jesus, and with His
brethren:'' and with this accords the repeated exhortation to pray
together without ceasing, which occurs in St. Paul's Epistles. It will
be observed that he insists in one passage on prayer to the abridgment
of sleep \footnote{Vid. Socr. Hist. 22. Vide
also Lyra Apostolica, xv.}; and one recorded
passage of his life exemplifies his precept. ``And at midnight Paul and
Silas prayed, and sang praises unto GOD,
and the prisoners heard them.'' Surely it is more natural to suppose
that this act of worship came in course, according to their wont, and
was only not omitted because of their imprisonment, somewhat after
Daniel's pattern, than that they should have gone aside to bear this
sort of indirect testimony to the Gospel.

Such was the Apostolic worship as far as Scripture
happens to have preserved it; that it was as systematic, and as
apportioned to particular times of the day, as in the aftertimes of
peace and prosperity, is not to be supposed; yet it seems to have been,
under ordinary circumstances, as ample and extended, as then. If St.
Paul thought a prison and a prison's inmates no impediment to vocal
prayer, we may believe it was no common difficulty which ever kept him
from it.

In subsequent times the Hours of prayer were
gradually developed from the three, or (with midnight) the four seasons,
above enumerated, to seven, viz. by the addition of Prime (the
first-hour), Vespers, (the evening), and Compline (bed-time); according
to the words of the Psalm, ``Seven times a day do I praise Thee,
because of Thy righteous judgments.'' Other pious and instructive
reasons existed, or have since been perceived, for this number. It was a
memorial of the seven days of creation; it was an honour done to the
seven petitions given us by our LORD in His prayer; it was a mode of pleading
for the influence of that Spirit who is revealed to us as sevenfold; on
the other hand, it was a preservative against those seven evil spirits,
which, are apt to return to the exorcised soul more wicked than he
who has been driven out of it; and it was a fit remedy of those seven
successive falls, which the Scripture says happen to ``the just man''
daily.

And, as the particular number of their Services
admitted of various pious meanings, so did each in its turn suggest
separate events in our SAVIOUR'S
history. He was born, and He rose again at midnight. At Prime, (or 7 A.M.
according to our reckoning,) He was brought before Pilate. At the third,
(or 9 A.M.)
He was devoted to crucifixion by the Jews, and scourged. At the sixth,
(or noon,) He was crucified. At the ninth, (or 3 P.M.) He expired. At Vespers, He was taken down from the cross;
at which hour He had the day before eat the Passover, washed His
Apostles' feet, and consecrated the Eucharist. At Completorium, or
Compline, He endured the agony in the garden.

These separate Hours, however, require a more
distinct notice. The night Service was intended for the end of the
night, when it was still dark, but drawing towards day; and, considering
that the hour for rest was placed soon after sunset, it did not infringe
upon the time necessary for repose. Supposing the time of sleep to
extend from 8 or 9 P.M.
to 3 or 4 in the morning, the worshipper might then rise without
inconvenience to perform the service which was called variously by the
name of Nocturns, or Matins, as we still differently describe the hours
in which it took place, as night or morning. It consists, when full, of
three parts or Nocturns, each made up of Psalms and Lessons; and it
ended in a Service, supposed to be used shortly before sunrise, and
called Lauds or Praises. This termination of the Nocturn Service is
sometimes considered distinct from it, so as to make eight instead of
seven Hours in the day; as if in accordance with the text, ``Give a
portion to seven, and also to eight.'' Accordingly it is sometimes
called by the name of Matins, instead of the Nocturns; and sometimes
both together are so called.

This subdivision of the night service has the
effect of dividing the course of worship into two distinct parts, of
similar structure with each other; the three Nocturns, Lauds, and Prime,
corresponding respectively to the three day hours (of the 3d, 6th, and
9th) Vespers and Compline. Of these the three day hours are made up of
Psalms, Hymns, and Sentences. These are the simplest of the
Services, and differ very little from each other through the year. Lauds
answer to Vespers, the sun being about to rise or about to set in the
one or the other respectively. Each contains five Psalms, a Text, Hymn,
Evangelical Canticle, Collect, and Commemoration of Saints. These hours
are the most ornate of the Services, and are considered to answer to the
morning and evening sacrifices of the Jews.

Prime and Compline were introduced at the same time
(the fifth century), and are placed respectively at the beginning of day
and the beginning of night. In each there is a Confession, four Psalms,
a Hymn, Text, and Sentences.

The ecclesiastical day is considered to begin with
the evening or Vesper service; according to the Jewish reckoning, as
alluded to in the text, ``In the evening, and morning, and at noon-day,
will I pray, and that instantly.'' The ancient Vespers are regarded by
some to be the most solemn hour of the day. They were sometimes called
the \emph{Officium Lucernarum} \footnote{Vid. Socr. Hist. 22. Vide
also Lyra Apostolica, xv.}.
Prayers were in some places offered while the lamps were lighting; and
this rite was called \emph{lumen offerre} \footnote{This ceremony must not be
confused with the \emph{Lucernarium}, or prayers at lighting the lamps;
which took place before the evening.}. The Mozarabic service supplies an instance of this, in which the
Office ran as follows:

\begin{quote}

``Kyrie
eleyson, Christe eleyson, Kyrie eteyson. Pater noster, \&c. In
nomine Domini Jesu Christi, lumen cum pace. R. Amen. Hoc est lumen
oblatum. R. Deo gratias.''

\end{quote}

On Festivals, the appropriate Services, beginning
on the evening of the preceding day, are continued over the evening of
the day itself; so that there are in such cases two Vespers, called the
First and the Second, of which the First are the more solemn.

This is the stated succession of the sacred offices
through the day, but the observance of the precise hours has not been
generally insisted on at any time, but has varied with local usages or
individual convenience. Thus the Matin and Laud Services may be
celebrated on the preceding evening, as is done (for instance) in the
Sestine Chapel at Rome during Passion week, the celebrated \emph{Miserere}
being one of the Psalms in Lauds. Prime may be used just before or after
sunrise; the Third soon after and soon after, the Sixth; the Ninth, near
dinner; Vespers and Compline, after dinner. Or Prime, the Third, Sixth,
and Ninth may come together two or three hours after sunrise. Noon,
which in most ages has been the hour for the meal of the day, is made to
divide the Services; there is a rule, for instance, against Compline
coming before dinner.

Such is the present order and use of the Breviary
Services, as derived more or less directly from Apostolic practice.
Impressed with their antiquity, our Reformers did not venture to write a
Prayer-Book of their own, but availed themselves of what was ready to
their hands: in consequence, our Daily Service is a compound of portions
of this primitive ritual, Matins being made up of the Catholic Matins,
Lauds, and Prime, and Even-song of Vespers and Compline. The reason why
these changes were brought about will be seen in the following sketch of
the history of the Breviary from the time of Gregory VII.

The word has been already explained to mean
something between a directory and an harmony of offices; but it is to be
feared there was another, and not so satisfactory reason for the use of
it. It implied an abridgment or curtailment of Services, and so in
particular of the Scripture readings, whether Psalms, or Lessons, at
least in practice. Of course there is no reason why the Church might
not, in the use of her discretion, limit as well as select the portions
of the inspired volume, which were to be introduced into her devotions;
but there were serious reasons why she should not defraud her children
of ``their portion of meat in due season;'' and it would seem, as if
the eleventh or at least the twelfth century, a time fertile in other
false steps in religion, must be charged also, as far as concerns Rome
and its more intimate dependencies, with a partial removal of the light
of the written Word from the Sanctuary. Whatsoever benefit attended the
adjustment of the offices in other respects, so far as the reading of
Scripture was omitted, it as productive of evil, at least in prospect.
An impulse was given, however slight in itself, which was followed up in
the Centuries which succeeded, and in all those churches which either
then, or in the course of time, adopted the usage of Rome.

Even now that usage is not universally received in
the Latin Communion, and it was in no sense enjoined on the whole
Communion till after the Council of Trent; but from the influence of the
papal see and of the monastic orders, it seems to have affected other
countries from a much earlier date. This influence would naturally be
increased by the circumstance that the old Roman Breviary had long
before Gregory's time been received in various parts of Europe: in
England, since the time of Gregory the Great, who, after the pattern of
Leo, and Gelasius before him, had been a Reformer of it; in Basle, since
the ninth century; in France and Germany by means of Pepin and
Charlemagne; while Gregory VII. himself effected its reception in Spain.
Other Breviaries however still were in use, as they are at this day. The
Ambrosian Breviary used in the Church of Milan, derives its name from
the great St. Ambrose; and in the ninth century Charles the Bald, while
sanctioning the use of the Roman, speaks also of the usage of Jerusalem,
of Constantinople, of Gaul, of Italy, and Toledo.

In Gregory's Breviary there are no symptoms of a
neglect of Scripture. It contains the offices for festival-days,
Sundays, and week-days; Matins on festivals having nine Psalms and nine
Lessons, and on Sundays eighteen Psalms and nine Lessons, as at present.
The course of the Scripture Lessons was the same as it had been before
his time; as it is preserved in a manuscript of the thirteenth century.
It will be found to agree in great measure both with the order of the
present Breviary and with our own. From Advent to Christmas were read
portions of the prophet Isaiah; from the Octave of the Epiphany to
Septuagesima, St. Paul's Epistle to the Romans; from Septuagesima to the
third Sunday in Lent, the book of Genesis, the i., xii. and xxvii. on
the Sundays to which they are allotted in our own offices; on the fourth
in Lent to Wednesday in Passion week, Jeremiah; from Easter to the third
Sunday after, the Apocalypse; from the third to the fifth, St. James;
from the Octave of the Ascension to Pentecost, the Acts; after the
Octave of Trinity to the last Sunday in July, the books of Kings; in
August, Proverbs; in September, Job, Tobit, Judith, and Esther; in
October, Maccabees; and in November, Ezekiel, Daniel, and other
Prophets.

Well would it have been if this laudable usage,
received from the first ages, and confirmed by Pope Gregory VII., had
been observed, according to his design, in the Roman Church; but his own
successors were the first to depart from it. The example was set in the
Pope's chapel of curtailing the sacred Services, and by the end of the
twelfth century it had been followed in all the churches in Rome, except
that of St. John Lateran. The Fratres Minores (Minorists or Franciscans)
adopted the new usage, and their Breviaries were in consequence
remarkable for the title ``secundum consuetudinem Romanæ Curiæ,''
contrary to the usage of such countries as conformed to the Roman
Ritual, which were guided by the custom of the churches in the city.
Haymo, the chief of this order, had the sanction of Gregory X., in the
middle of the thirteenth century, to correct and complete a change,
which, as having begun in irregularity, was little likely to have fallen
of itself into an orderly system; and his arrangements, which were
conducted on the pattern of the Franciscan devotions, nearly correspond
with the Breviary, as it at present stands.

Haymo's edition, which was introduced into the
Roman Church by Nicholas III. A.D.
1278, is memorable for another and still more serious fault. Graver and
sounder matter being excluded, apocryphal legends of Saints were used to
stimulate and occupy the popular mind; and a way was made for the use of
those Invocations to the Virgin and other Saints, which heretofore were
unknown in public worship. The addresses to the Blessed Mary in the
Breviary, as it is at present constituted, are such as the following:
the Ave Mary, before commencing every office through the day and at the
end of Compline; at the end of Lauds and Vespers, an Antiphon invocatory
of the Virgin; the Officium B. Mariæ, on the Sabbath or Saturday, and
sundry other offices, containing Hymns and Antiphons in her honour.
These portions of the Breviary carry with them their own plain
condemnation, in the judgment of an English Christian; no commendation
of the general structure and matter of the Breviary itself will have any
tendency to reconcile him to them; and it has been the strong feeling
that this is really the case, that has led the writer of these pages
fearlessly and securely to admit the real excellences, and to dwell upon
the antiquity, of the Roman Ritual. He has felt that, since the
Romanists required an unqualified assent to the \emph{whole} of the
Breviary, and that there were passages which no Anglican ever could
admit, praise the true Catholic portion of it as much as he might, he
did not in the slightest degree approximate to a recommendation of
Romanism. But to return;—these Invocations and Services to the Blessed
Virgin have been above enumerated, with a view of observing that, on the
very face of them, they do not enter into the \emph{structure} of the
Breviary; they are really, as they are placed, additions, and might
easily have been added at some future period, as (e.g.) was the case
with our own Thanksgiving, or the Prayer for the Parliament. This remark
seems to apply to all the intrinsically-exceptionable Addresses in the
Breviary; for as to the Confession at Prime and Compline, in which is
introduced the name of the Blessed Virgin and other Saints, this
practice stands on a different ground. It is not a simple gratuitous
Invocation made to them, but it is an address to ALMIGHTY
GOD \emph{in
His heavenly court}, as surrounded by His Saints and Angels,
answering to St. Paul's charge to Timothy, ``before GOD and the LORD JESUS
CHRIST
and the elect Angels,'' and to Daniel and St. John's address to the
Angels who were sent to them. The same may even be said of the
Invocation ``Holy Mary and all Saints,'' \footnote{It is observable that the
words ``Holy Mary'' do not occur in the ancient Monastic Breviaries.
The Confession at Prime and Compline does not occur in the Paris
Breviary, 1735.} \&c. in the Prime Service, which Gavanti describes as being of
very great antiquity. These usages certainly \emph{now} do but sanction
and encourage that direct worship of die Blessed Virgin and the Saints,
which is the great practical offence of the Latin Church, and so are a
serious evil; but it is worth pointing out, that, as on the one hand
they have more claim to be considered an integral part of the service,
so on the other, more can be said towards their justification than for
those addresses which are now especially under our consideration.

This is what occurs to observe on the first sight
of these Invocations; but we are not left to draw a conjectural judgment
about them. Their history is actually known, and their recent
introduction into the Church Services is distinctly confessed by Roman
ritualists.

The Ave Mary, for instance, is made up of the Angel's
salutation, ``Hail, thou,'' \&c. Elizabeth's ``Blessed art thou among
women,'' \&c. and the words, ``Holy Mary, Mother of GOD, pray for us sinners, now and in the hour
of our death.'' The last clause ``now and,'' \&c. was confessedly added
by the Franciscans in the beginning of the sixteenth century; and the
words preceding it, ``Holy Mary,'' \&c. which Gavanti, after Baronius,
wishes to attribute to the Council of' Ephesus (A.D.
431), are acknowledged by the later critics, Grancolas and Merari, to
have had no place in any form of prayer till the year 1508. Even the
Scripture portion of the Ave Mary, which, as Merari observes, is an
Antiphon rather than Prayer, and which occurs as such in the lesser
office of the Blessed Virgin, and in St. Gregory's Sacramentary in the
Mass Service for the fourth Sunday in Advent, is not mentioned by any
devotional writer, nor by Councils, nor Fathers, up to the eleventh
century, though they do enjoin the universal and daily use of the Creed
and LORD'S
Prayer, which are in the present Breviary used with it. It first occurs
among forms of prayer prescribed for the people in the statutes of Otho,
Bishop of Paris, A.D.
1195, who was followed, after the interval of a hundred years, by the
regulations of Councils at Oxford and elsewhere. Another space of at
least fifty years intervenes before the introduction of rosaries and
crowns in honour of the Virgin. As to the Roman Breviary, it did not
contain any part of the Ave Mary, till the promulgation of it by Pope
Pius V., after the Tridentine Council, A.D.
1550.

The four Antiphons to the Blessed Virgin, used at
the termination of the offices, are known respectively by their first
words: the \emph{Alma Redemptoris}, the \emph{Ave Regina}, the \emph{Regina
cœli}, and the \emph{Salve Regina}. Gavanti and Merari plainly tell
us that they are not to be found in ancient authors. The \emph{Alma
Redemptoris} is the composition of Hermannus Contractus, who died A.D.
1054. The author of the \emph{Ave Regina} is unknown, as is that of the \emph{Regina
cœli}. The Salve Regina is to be attributed either to Hermannus, or
to Peter of Compostella. Gavanti would ascribe the last words ``O clemens,
O pia, O dulcis,'' \&c. to St. Bernard, but Merari corrects him, the
work in which they are contained being supposititious. These Antiphons
seem to have been used by the Franciscans after Compline from the
thirteenth century: but are found in no Breviary before A.D. 1520.

The Saturday or Sabbath office of the Blessed
Virgin was introduced, according to Baronius, by the monks of the
Western church, about A.D.
1056.

The Officium Parvum B. V. M. was instituted by the
celebrated Peter Damiani at the same date. It is said indeed to have
been the restoration of a practice three hundred years old, and observed
by John Damascene; which it may well have been: but there is nothing to
show the identity of the Service itself with the ancient one, and that
is the only point on which evidence would be important. Thirty years
after its introduction by Damiani, it was made part of the daily worship
by decree of Urban II.

The Breviary then, as it is now received, is pretty
nearly what the Services became \emph{in practice} in Rome, and among
the Franciscans by the middle of the thirteenth century; the two chief
points of difference between it and the ancient Catholic Devotions,
being on the one hand its diminished allowance of Scripture reading, on
the other its adoption of uncertain legends, and of Hymns and Prayers to
the Virgin. However, the more grievous of these changes were not finally
made in the Breviary itself, till the Pontificate of Pius V., after the
Tridentine Council; at which time also it was imposed in its new form
upon all the Churches in communion with Rome, except such as had used
some other Ritual for above two hundred years. Not even at the present
day, however, is this Roman novelty, as it may be called, in universal
reception; the Paris Breviary, as corrected by the Archbishop of that
city, A.D.
1735, differs from it considerably in detail, though still disfigured by
the Invocations.

Before concluding this account of the Roman
Breviary, it is necessary to notice one attempt which was made in the
first part of the sixteenth century to restore it to a more primitive
form. In the year 1536, Quignonius, Cardinal of Santa Crux, compiled a
Breviary under the sanction of Clement VII., and published it under his
successor, Paul III. This Ritual, the use of which was permitted, but
not formally enjoined by the Holy See, was extensively adopted for forty
years, when it was superseded by the Franciscan Breviary, as the
now authorized one may be called, in consequence of a Bull of Pius V.
The Cardinal's Breviary was drawn up on principles far more agreeable to
those on which the Reformation was conducted, and apparently with the
same mixture of right and wrong in the execution. With a desire of
promoting the knowledge of Scripture, it showed somewhat of a rude
dealing with received usages, and but a deficient sense of what is
improperly called the \emph{imaginative} part of religion. His object
was to adapt the Devotions of the Church for private reading, rather
than chanting in choir, and so to encourage something higher than that
almost theatrical style of worship, which, when reverence is away, will
prevail, alternately with a slovenly and hurried performance, in the
performance of Church Music. Accordingly he left out the Versicles,
Responses, and Texts, which, however suitable in Church, yet in private
took more time, as he says, to find out in the existing formularies than
to read when found. He speaks in his preface expressly of the ``perplexus
ordo,'' on which the offices were framed. But his great reform was as
regards the reading of Scripture. He complains that, whereas it was the
ancient rule that the Psalms should be read through weekly and the Bible
yearly, both practices had been omitted. The Ferial or week-day service
had been superseded by the service for feast-days, as being shorter: and
for that reason every day, even through Lent, was turned into a
festival. To obviate the temptation which led to this irregularity, he
made the Ferial service about the length of that of the old feast day;
and he found space in these contracted limits for the reading of the
Psalms, and the whole Bible, except part of the Apocalypse, in the week
and the year respectively, by omitting the popular legends of the Saints
which had been substituted for them. He observes, that these
compositions had been sometimes introduced without any public authority,
or sanction of the Popes, merely at the will of individuals. Those which
he retained, he selected from authors of weight, whether of the Greek or
Latin Church. Besides, he omitted the Officium Parvum B. M. V., on the
ground that there were sufficient services in her honour independently
of it. In all his reforms he professes to be returning to the practice
of antiquity; and he made use of the assistance of men versed ``in
Latin and Greek, in divinity, and the jus pontificium.''

This Breviary was published in Rome, A.D. 1536, under the sanction, as has
been said, of Paul III. However, it was not of a nature to please the
divines of an age which had been brought up in the practice of the
depraved Catholicism then prevalent; and its real faults, as they would
appear to be, even enabled them to oppose it with justice. The Doctors
of the Sorbonne proceeded to censure it as running counter in its
structure to antiquity and the Fathers; and though they seem at length
to have got over their objections to it, and various editions at Venice,
Antwerp, Lyons, and Paris, showed that it was not displeasing to numbers
in the Roman Communion, it was at length superseded by the Bull of Pius
V. establishing the Franciscan Breviary, which had more or less grown
into use in the course of the preceding three hundred years.

This account of Cardinal Quignonius's Breviary, and
the circumstances under which it was compiled, will remind the English
reader of the introductory remarks concerning the Service of the Church,
prefixed to our Ritual; which he may read more profitably than
heretofore, after the above illustrations of their meaning. For this
reason they shall be here cited

``There was never any thing by the wit of man so
well devised, or so sure established, which in continuance of time hath
not been corrupted; as, among other things, it may plainly appear by the
Common Prayers, in the Church, commonly called Divine Service. The first
original and ground whereof, if a man would search out by the Ancient
Fathers, he shall find that the same was not ordained but of a good
purpose, and for a great advancement of godliness. For they so ordered
the matter, that all the whole Bible, (or the greatest part thereof,)
should be read over once every year; intending thereby that the Clergy,
and especially such as were Ministers in the Congregation, should (by
often reading and meditating on GOD'S Word) be stirred up to godliness
themselves, and be more able to exhort others by wholesome doctrine, not
to confute them that were adversaries to the truth; arid further, that
the people (by daily hearing of Holy Scripture read in the Church,)
might continually profit more and more in the knowledge of GOD, and be the more inflamed with the love of
His true religion.

``But these many years past, this godly and decent
order of the ancient Fathers hath been so altered, broken, and
neglected, by planting in uncertain Stories and Legends, with multitude
of Responds, Verses, vain Repetitions, Commemorations, and Synodals;
that commonly when any book of the Bible was begun, after three or four
chapters were read out, all the rest were unread. And in this sort the
book of Isaiah was begun in Advent, and the book of Genesis in
Septuagesima; but they were only begun, and never read through. After
like sort were other books of Holy Scripture used. And furthermore,
not-withstanding that the ancient Fathers have divided the Psalms into
seven portions, whereof every one was called a Nocturn, now of late time
a few of them have been daily said, and the rest utterly omitted.
Moreover, the number and hardness of the rules called the Pie, and the
manifold changings of the service, was the cause, that, to turn the book
only was so hard and intricate a matter, that many times there was more
business to find out what should be read, than to read it when it was
found out.

``These inconveniences therefore considered, here is
set forth such an Order, whereby the same shall be redressed. And for a
readiness in this matter, here is drawn out a Calendar for that purpose,
which is plain and easy to be understood; wherein (so much as may be)
the reading of holy Scripture is so set forth, that all things shall be
done in order, without breaking one piece from another. For this cause
be cut off Anthems, Responds, Invitatories, and such like things as did
break the continual course of the reading of the Scripture.''

It remains but to enumerate the selections from the
Breviary which follow. First has been drawn out, an Analysis of the
Weekly Service, as well for Sunday as other days. This is followed by an
ordinary Sunday Service at length, as it runs when unaffected by the
occurrence of special feast or season, in order to ground the reader,
who chooses to pursue the subject, in the course of daily worship as a
whole. With the same object a Week-day Service has also been drawn out.
Two portions of extraordinary Services are then added, one from the
Service for the Transfiguration, the other for the Festival of St.
Lawrence, with a view of supplying specimens of a more elevated and
impressive character. Next follows a design for a Service for March
21st, the day on which Bishop Ken was taken from the Church below, and
another for a Service of thanksgiving and commemoration for the
anniversaries of the days of death of friends or relations. These have
been added, to suggest to individual Christians a means of carrying out
in private the principle and spirit of those inestimable forms of
devotion which are contained in our authorized Prayer-Book. The series
is closed with an abstract of the Services for every day in Advent,
fitting on to sections 2 and 3, which contain respectively the types of
the Sunday and Week-day Service. Except by means of some such extended
portion, it is impossible for the reader to understand the general
structure, and appreciate the harmony of the Breviary.

Lastly, the writer of these pages feels he shall
have to ask indulgence for such chance mistakes, in the detail of the
following Services, as are sure to occur when an intricate system is
drawn out and set in order, with no other knowledge of it than is
supplied by the necessarily insufficient directions of a Rubric.

\xsection{§ 1. Analysis of the Seven Daily Services of the Church Catholic, as preserved in the Breviary.}

EVERY
Service but Compline is commenced with privately saying the Lord's
Prayer, and the Ave Mary, to which the Creed is added before Matins and
Prime. In like manner, after Compline, all three are repeated. Every
other Service ends with the Lord's Prayer in private, unless another
Service immediately follows. Concerning the introduction of the Ave
Mary, vid. supra, p. 11. This use of the Lord's prayer in private
before the beginning of the Service seems to have led the compilers of
King Edward's First Book to open with the Lord's Prayer, only said
aloud, not in private; but a pious custom has brought in again the
private prayer, as before, though without prescribing any particular
form. The compilers of King Edward's Second Book prefixed to the
Lord's Prayer, the Sentences, and an Exhortation, Confession, and
Absolution of their own. After these follows, ``O Lord, open thou our
lips,'' \&c. which stands first in the Breviary Service.

\xsubsection{1. MATINS, or Night Service, (after One, A.M.)}

\emph{Introduction}.

\emph{Verse}. O Lord, open Thou my lips.

\emph{Resp}. And my mouth shall shew forth Thy praise. 

(Each person to sign his lips with the Cross.)

\emph{Verse}. O God, make speed to save me.

\emph{Resp}. O Lord, make haste to help me.

(Each person to sign himself from the forehead to the breast.)

Glory be to the Father, \&c.

As it was, \&c. Amen.

(\emph{Ordinarily added}) Hallelujah. (\emph{i.e}. Praise ye the
Lord.)

\emph{Psalm} 95.—``O come let us sing,''
\&c. with a verse called an Invitatory, ``Let us worship the Lord:
our Maker,'' divided into two parts, the whole being used before the
1st, 3rd, and 8th verse, and at the end, and again after the \emph{Gloria
Patri}, and the latter part after the 4th and 9th, and between the \emph{Gloria}
and the whole. This Invitatory varies with the season, but its general
character is always preserved; e.g. in Advent, ``O Come let us
worship: the Lord, the King to come;'' or ``the Lord is at hand, O
come let us worship;'' again in Lent, ``It is not to you lost labour
that ye haste to rise up early: for the Lord hath promised a crown to
those who wait for Him.'' At Pentecost, ``Hallelujah, the Spirit of
the Lord hath filled the round world: O come let us worship,
Hallelujah.''

A Hymn follows according to the day, and terminates
the Introduction; then follow Psalms and Lessons, in one or three
Nocturns, according as the Service is for Weekday or Sunday.

\emph{On
Sunday}, Eighteen Psalms with Nine Lessons; viz.

Psalms 1,
2, 3, 6, 7, 8, 9, \& 10, 11, 12, 13, 14, 15.

A passage from Scripture, in three parts—(\emph{according to the time of the Year}.)

Psalms 16, 17, 18.

A passage from some Father of the Church, in three parts.

Psalms 19, 20, 21.

A comment on some passage of the Gospel, in three parts.

On Weekdays, Twelve Psalms with Three Lessons, viz.

\emph{On
Monday}.

Psalms 27, 28, 29, 30, 31, 32, 33, 34, 35, 36, 37, 38.

A passage in three parts from Scripture or the Fathers.

\emph{On
Tuesday}.

Psalms 39, 40, 41, 42, 44, 45, 46, 47, 48, 49, 50, 52.

A passage in three parts, \&c. as on \emph{Monday}.

\emph{On
Wednesday}.

Psalms 53, 55, 56; 57, 58, 59, 60, 61, 62, 64, 66, 68.

A passage, \&c.

\emph{On
Thursday}.

Psalms 69, 70, 71, 72, 73, 74, 75, 76, 77, 78, 79, 80.

A passage, \&c.

\emph{On
Friday}.

Psalms 81, 82, 83, 84, 85, 86, 87, 88, 89, 94, 96, 97.

A passage, \&c.

On
Saturday.

Psalms 98, 99, 100, or 92. (\emph{according to the day}) 101, 102,

103, 104, 105, 106, 107, 108, 109.

A passage, \&c.

———————

Then on every day of the week follows the

Te Deum Laudamus.

This noble Hymn follows in this place with especial
propriety, on Sundays and other Festivals; viz. after the reading the
words of Prophets and Apostles, and the writings and histories of Saints
and Martyrs, all of whom are commemorated in it. On all days it
impressively winds up the Service which precedes.

———————

\xsubsection{LAUDS;—(appended to the Matins towards the first twilight.)}

\emph{Verse}. ``O
God, make speed, \&c.

\emph{Resp}. O Lord,
make haste,''

Glory be, \&c.

As it was, \&c. Amen. Hallelujah.

\emph{Then
five} Psalms, \emph{viz. on}

\emph{Sunday}.

Psalms 93, 100, 63 and 67.

The Song of the Three Children.

Psalms 148-150.

\emph{Monday}.

Psalms 51, 5, 63. Song of Isaiah. (Is. xii.) Psalm 148.

\emph{Tuesday}.

Psalms 51, 43, 63. Song of Hezekiah. (Is. xxxviii.) Psalm 148.

\emph{Wednesday}.

Psalms 51, 65, 63. Song of Hannah. (1 Sam. ii. Psalm 148.)

\emph{Thursday}.

Psalms 51, 90, 63. Song of Moses. (Ex. xv.) Psalm 148.

\emph{Friday}.

Psalms 51, 143, 63. Song of Habakkuk. (Hab. iii.) Psalm 148.

\emph{Saturday}.

Psalms 51, 92, 63. Song of Moses. (Deut. xxxii.) Psalm 148.

\emph{The Service ends on all days with a} Text (\emph{Capitulum}),
a Hymn and a Collect (\emph{Oratio}), \emph{varying with the day and season};
the Song of Zacharias (\emph{Benedictus}) \emph{being interposed between
the Hymn and Collect, and several stated Collects, invocations, and
Sentences following}.

\xsubsection{2. PRIME, on the Rising of the Sun.}

\emph{Verse}. ``O God, make speed, \&c.

\emph{Resp}. O Lord, make haste, \&c'' and the \emph{Gloria Patri},
\&c. \emph{as before}.

\emph{A} Hymn, \emph{the same every day in the year}.
\emph{Then four}Psalms (\emph{except Saturday when there are three},) \emph{viz}.

\begin{quote}

Psalm 54;—

\emph{Sunday}.—Psalm 118, or (\emph{sometimes}) 93.

\emph{Monday}.—Psalm 24.

\emph{Tuesday}. —Psalm 25.

\emph{Wednesday}.—Psalm 26.

\emph{Thursday}.—Psalm 23.

\emph{Friday}.—Psalm 22.

\emph{Saturday}.—Psalm is \emph{omitted}.

\end{quote}

Then, throughout the week,—

\begin{quote}

Psalm 119, v. 1-32. in two parts.

\end{quote}

Then on \emph{Saturday} only, follows the Psalm
Quicumque, commonly called the Athanasian Creed. It is a far truer view
of this venerable composition, to consider it a Psalm or Hymn of praise,
and of concurrence in GOD'S
appointments, as Psalm 118 or 139, or the Te Deum, than as a formal
Creed; and by using it weekly, its living character and spirit are
incorporated into the Christian's devotions, and its influence on the
heart, as far as may be, secured. The time, too, should be observed. The
dawn of the first day of the week.

\emph{The Service concludes with a} Text (\emph{Capitulum});
\emph{with the} Lord's Prayer, \emph{privately}: \emph{a} Confession \emph{of
Priest to People}, \emph{and in turn of People to Priest}, \emph{and a
corresponding} Absolution: Sentences; Collect, \emph{the Third}, \emph{for
grace in our own Morning Service}, a Lesson from the Book of Martyrs;
an Invocation \emph{of St. Mary and All Saints}; Sentences, \emph{with the}
Lord's Prayer, \emph{privately}; Collect, \emph{the Second}, \emph{at the
end of our Communion Service}; a Short Lesson; \emph{and} Sentences.

\xsubsection{3. The THIRD (Nine A.M.) 4. SIXTH (noon.) 5. NINTH (Three P.M.)}

``O God, make speed,'' \&c. \emph{as before}.

\emph{A} Hymn, \emph{the same throughout the year},
\emph{at the same hours respectively}: \emph{then}—

\begin{quote}

\emph{At the Third}.—Psalm 119, v. 33-80. in
three parts.

\emph{At the Sixth}.—Psalm 119, v. 81-128. in three parts.

\emph{At the Ninth}.—Psalm 119, v. 129-176. in three parts.

\end{quote}

Thus the whole of the 119th Psalm is gone through
every day of the year.

\emph{Then} a Text (\emph{Capitulum}) \emph{and}
Sentences, \emph{with the} Lord's Prayer, \emph{privately}, \emph{varying
with the time of the year}. \emph{Then the} Collect \emph{for the day or
week}.

\xsubsection{6. VESPERS. (Evening.)}

``O God, make speed,'' \&c. \emph{as before}; \emph{then
Five} Psalms, \emph{viz. on}

\begin{quote}

\emph{Sunday}.—Psalms 110, 111, 112, 113, 114,
and 115.

\emph{Monday}.—Psalms 116, \emph{in two parts}, 117, 120, 121.

\emph{Tuesday}.—Psalms 122, 123, 124, 125, 126. 

\emph{Wednesday}.—Psalms 127, 128, 129, 130, 131.

\emph{Thursday}.—Psalms 132, 133, 135, 136, 137.

\emph{Friday}—Psalms 138, 139, 140, 141, 142.

\emph{Saturday}.—Psalms 144, 145, 146, 147, \emph{in two parts}.

\end{quote}

\emph{Then on every day a} Text (\emph{Capitulum}),
\emph{a} Hymn, \emph{and the} Collect, \emph{all varying with the day and
season}; \emph{between the Hymn and Collect always is interposed the}
Magnificat, \emph{sometimes with} Sentences \emph{after it}. \emph{The
Service ends}, \emph{as} Lauds, \emph{with} Collects \emph{and}
Invocations.

\xsubsection{7. COMPLINE. (Bed-time.)}

This Service is almost invariable throughout the
year.

It \emph{begins with a} Blessing \emph{for the ensuing
night}; \emph{a} Short Lesson; \emph{the} Confession \emph{and}
Absolution \emph{as at Prime}; Sentences; \emph{then four} Psalms, \emph{viz}.

\begin{quote}

Psalms 4, 31, down to v. 6, 9, 134.

\end{quote}

\emph{A} Hymn: Text (\emph{Capitulum}); Sentences;
Song of Simeon (Luke ii); Sentences \emph{with the} Lord's Prayer \emph{and}
Creed, \emph{privately}; Collect \emph{for safety during the night}. \emph{The
Service ends with an} Antiphon \emph{in praise of the} Virgin, \emph{and}
Collect \emph{upon it}.

———————

To this Sketch of the Services must be added a few
words concerning the Antiphons and Benedictions which occur throughout
them, but have not been noticed in their places.

The Antiphons or Anthems are sentences preceding
and succeeding the Separate Psalms and Songs, and are ordinarily verses
taken from the particular compositions to which they are attached. They
seem to answer the purpose of calling attention to what is coming, of
interpreting it, or of pointing out the particular part of it which is
intended to bear on the Service of the day; in all respects answering
the place of what is called by musicians a key note. They are repeated
at the end, as if to fix the impression or the lesson intended.

Antiphons are introduced in other connexions, as
before Collects; such are those, for instance, in the Commemorations of
the Blessed Virgin and the Apostles in Vespers and Lauds. Such,
too, are the Antiphons to the Blessed Virgin placed at the end of
Compline, to which especial attention was above directed (p. 11), on the
ground of their objectionable nature. They shall be here given, in order
to show clearly, as a simple inspection of them will suffice to do, the
utter contrariety between the Roman system, as actually existing, and
our own; which, however similar in certain respects, are in others so at
variance, as to make any attempt to reconcile them together in their
present state perfectly nugatory. Till Rome moves towards us, it is
quite impossible that we should move towards Rome; however closely we
may approximate to her in particular doctrines, principles, or views. In
reading the following, it should be recollected, indeed, that Antiphons
are not, strictly speaking, Prayers, but Sentences applied to the
particular purpose of meditation, thanksgiving, \&c.; yet the
following, taken together, are quite beyond the power of any defence
which might thence be available for less explicit compositions.

\emph{From Advent to the Purification}.

Alma Redemptoris Mater quæ pervia cœli

Porta manes, et stella maris, succurre cadenti,

Surgere qui curat, populo; to quæ genuisti,

Naturâ mirante, tuum sanctum Genitorem,

Virgo prius ac posterius, Gabrielis ab ore,

Sumens illud Ave, peccatorum miserere.
Kindly Mother of the Redeemer, who art ever of
heaven

The open gate, and the star of the sea, aid a fallen people,

Which is trying to rise again; thou who didst give birth,

While Nature marvelled how, to thy Holy Creator,

Virgin both before and after, from Gabriel's mouth

Accepting the All hail, be merciful towards sinners.

\emph{From the Purification to Good Friday}.

Ave, Regina cœlorum!

Ave, Domina Angelorum!

Salve radix, salve porta!

Ex quâ mundo lux est orta.

Gaude, Virgo gloriosa,

Super omnes speciosa;

Vale, o valde decora,

Et pro nobis Christum exora.
Hail, O Queen of the heavens

Hail, Lady of Angels

Hail, the root! hail the gate!

Whence to the world light is risen.

Rejoice, O glorious Virgin,

Beautiful above all;

Farewell, O thou most comely,

And prevail on Christ for us by thy prayer.

\emph{From Easier to the First Week complete
after Pentecost}.

Regina cœli, lætare,

Alleluia.

Quia quem meruisti portare, 

Alleluia.

Resurrexit, sicut dixit,

Alleluia.

Ora pro nobis Deum.

Alleluia.
Rejoice, O queen of heaven,

Hallelujah.

For He, whom for thy obedience' sake thou didst bear,

Hallelujah.

Is risen, as he said,

Hallelujah.

Pray thou God for us.

Hallelujah.

\emph{From Trinity tide to Advent}.

Salve Regina, mater misericordiæ, vita, dulcedo et spes
nostra, salve. Ad te clamamus exules, filii Hevæ. Ad te
suspiramus, gementes et flentes in hac lachrymarum valle. Eja
ergo advocata nostra, illos tuos misericordes oculos ad nos
converte, et Jesum, benedictum fructum ventris tui, nobis post
hoc exilium ostende. O clemens, O pia, O dulcis Virgo Maria.
Hail, O Queen, the mother of mercy, our life, sweetness, and
hope, hail. To Thee we exiles cry out, the sons of Eve. To Thee
we sigh, groaning and weeping in this valley of tears. Come
then, O our Patroness, turn thou on us those merciful eyes of
thine, and show to us, after this exile, Jesus the blessed fruit
of thy womb. O gracious, O pitiful, O sweet Virgin Mary.

These Antiphons have already been shown to be of
comparatively modern origin; which, indeed, is sufficiently clear from
their composition, independently of the question of doctrine. The
Absolutions and Benedictions on the other hand seem, from their
doctrinal character, to come from high antiquity. Wheatly remarks, that
the precise \emph{indicative} Absolution, such as it occurs in our
Visitation for the Sick, though altogether justifiable and edifying, did
not come into use till the twelfth century; that is, about the time of
the above innovations in commemorating the Blessed Virgin. Now the
Absolutions and Benedictions in the Breviary happen, on the contrary, to
be of a remarkably simple character; they are uniformly in the shape of \emph{petitions}
to ALMIGHTY GOD,
and they include the minister using them, being worded in the first, not
the second person. Again, in the \emph{quasi} Absolution, after the
stated Confession at Prime and Compline, it is to be noticed, that the
People absolve the Priest, before, and in the same words in which the
Priest absolves the People, as if vindicating to the body of Christians
that sacramental power, (whatever may be its degree,) which might
have seemed inconsistent with the special stress laid by Romanism on
Sacerdotal gifts. An Absolution occurs in each Nocturn between the
Psalms and Lessons: a short Benediction is pronounced before the reading
of each of the latter, being first asked for by the Reader.

\xsection{§ 2. Service for Sunday, June 21, 1801.}

Fourth Sunday after Pentecost

\xsubsection{1. MATIN SERVICE}

O LORD,
open Thou my lips.

\emph{And my mouth shall show forth
Thy praise}.

O God, make speed to save me.

\emph{O Lord make haste to help me}.

Glory be, \&c. \emph{Amen}.

Praise ye the Lord. (Hallelujah.)

(a) \emph{Invitatory}
Let us worship the
Lord: our Maker.

Let us worship, \&c.

\emph{Psalm} 95.
O come, let us
sing unto the Lord: let us heartily rejoice

in the strength of our salvation.

Let us come before His presence with
thanksgiving: and

shew ourselves glad in him with psalms.

Let us worship the Lord our Maker.

For the Lord is a great God: and a great King
above

all gods.

In His hand are all the corners of the earth:
and the

strength of the hills is his also.

Our Maker.

The sea is his, and he made it: and his hands
prepared

the dry land.

O come, let us worship and fall down: and
kneel before

the Lord our Maker.

For he is the Lord our God: and we are the
people of

his pasture, and the sheep of his hand.

Let us worship the Lord our Maker.

Today if ye will hear his voice, harden not
your hearts:

as in the provocation, and as in the day of temptation in the

wilderness.

When your fathers tempted me: proved me, and
saw

my works. 

Our Maker.

Forty years long was I grieved with this
generation, and

said: It is a people that do err in their hearts; for they have

not known my ways.

Unto whom I aware in my wrath: that they
should not

enter into my rest.

Let us worship the Lord our Maker.

Glory be, \&c.

As it was in the beginning, \&c. \emph{Amen}.

Our Maker.

Let us worship the Lord our Maker.

(b) \emph{Hymn}

[Nocte surgentes.]

Let us arise, and watch by night,

And meditate always;

And chant, as in our Maker's sight,

United hymns of praise.

So, singing with the Saints in bliss,

With them we may attain

Life everlasting after this,

And heaven for earthly pain.

Grant it to us, O Father, Son,

And Spirit, God of grace,

To whom all worship shall be done

In every time and place. \emph{Amen}.

———————

N\textsc{octurn}
1.

(c) \emph{Antiphon}
Serve ye the Lord.

\emph{Psalm} 1 (1)
B\textsc{lessed}
is the man that hath not walked in the counsel

of the ungodly, nor stood in the way of sinners: and hath

not sat in the seat of the scornful;

But his delight is in the law of the Lord:
and in his law

will he exercise himself day and night.

And he shall be like a tree planted by the
water-side:

that will bring forth his fruit in due season.

His leaf also shall not wither: and look,
whatever he

doeth, it shall prosper.

As for the ungodly, it is not so with them:
but they are

like the chaff, which the wind scattereth away from the face

of the earth. 

Therefore the ungodly shall not be able to
stand in the

judgment: neither the sinners in the congregation of the

righteous.

But the Lord knoweth the way of the
righteous: and the

way of the ungodly shall perish. Glory be, \&c.

\emph{Psalm} 2 (2)
Why do the heathen so
furiously rage together: and why

do the people imagine a vain thing?

The kings of the earth stand up, and the
rulers take coun-

sel together: against the Lord, and against his Anointed.

Let us break their bonds asunder: and cast
away their

cords from us.

He that dwelleth in heaven shall laugh them
to scorn:

the Lord shall have them in derision.

Then shall he speak unto them in his wrath:
and vex

them in his sore displeasure.

Yet have I set my King: upon my holy hill of
Sion.

I will preach the law, whereof the Lord hath
said unto

me: Thou art my Son, this day have I begotten thee.

Desire of me, and I shall give thee the
heathen for thine

inheritance: and the utmost parts of the earth for thy

possession.

Thou shalt bruise them with a rod of iron:
and break

them in pieces like a potter's vessel.

Be wise now therefore, O ye kings: be
learned, ye that

are judges of the earth.

Serve the Lord in fear: and rejoice unto him
with

reverence.

Kiss the Son, lest he be angry, and so ye
perish from the

right way: if his wrath be kindled, (yea, but a little) blessed

are all they that put their trust in him. Glory be, \&c.

\emph{Psalm} 3 (3)
Lord, how are they increased
that trouble me: many are

they that rise against me.

Many one there be that say of my soul: There
is no help

for him in his God.

But thou, O Lord, art my defender: thou art
my worship,

and the lifter up of my head.

I did call upon the Lord with my voice and he
heard

me out of his holy hill.

I laid me down and slept, and rose up again:
for the Lord

sustained me.

I will not be afraid for ten thousands of the
people: that

have set themselves against me round about.

Up, Lord, and help me, O my God: for thou
smitest all

mine enemies upon the check-bone; thou hast broken the

teeth of the ungodly.

Salvation belongeth unto the Lord: and thy
blessing is

upon thy people. Glory be, \&c. \emph{Psalm} 6 (4)
O Lord, rebuke me not in thine
indignation; neither

chasten me in thy displeasure.

Have mercy upon me, O Lord, for I am weak: O
Lord

heal me, for my bones are vexed.

My soul also is sore troubled: but, Lord, how
long wilt

thou punish me?

Turn thee, O Lord, and deliver my soul: O
save me for

thy mercies' sake.

For in death no man remembereth thee: and who
will

give thee thanks in the pit?

I am weary of my groaning; every night wash I
my bed:

and water my couch with my tears.

My beauty is gone for very trouble: and worn
away

because of all mine enemies.

Away from me, all ye that work vanity; for
the Lord hath

heard the voice of my weeping.

The Lord hath heard my petition: the Lord
will receive

my prayer.

All mine enemies shall be confounded, and
sore vexed:

they shall be turned back, and put to shame suddenly.

Glory be, \&c.

(c) \emph{Antiphon}
Serve the Lord in fear, and
rejoice unto Him with

reverence.

(d) \emph{Antiphon}
God is a righteous judge.

\emph{Psalm} 7 (5)
O Lord my God, in Thee have I
put my trust: save me

from all them that persecute me, and deliver me.

Lest he devour my soul hike a lion, and tear
it in pieces:

while there is none to help.

O Lord my God, if I have done any such thing:
or if there

be any wickedness in my hands;

If I have rewarded evil unto him that dealt
friendly with

me: yea, I have delivered him that without any cause is mine

enemy:

Then let mine enemy persecute my soul, and
take me:

yea, let him tread my life down upon the earth, and lay mine

honour in the dust.

Stand up, O Lord, in thy wrath, and lift up
thyself, because

of the indignation of mine enemies: arise up for me in

the judgment that thou hast commanded.

And so shall the congregation of the people
come about

thee: for their sakes therefore lift up thyself again.

The Lord shall judge the people; give
sentence with me, O

Lord: according to my righteousness, and according to the

innocency that is in me.

O let the wickedness of the ungodly come to
an end: but

guide thou the just.

For the righteous God trieth the very hearts
and reins. 

My help cometh of God: who preserveth them
that are true

of heart.

God is a righteous judge, strong and patient;
and God is

provoked every day.

If a man will not turn, he will whet his
sword: he hath

bent his bow, and made it ready.

He hath prepared for him the instruments of
death: he

ordaineth his arrows against the persecutors.

Behold, he travaileth with mischief: he hath
conceived

sorrow, and brought forth ungodliness.

He hath graven and digged up a pit: and is
fallen him-

self into the destruction that he made for other.

For his travail shall come upon his own head:
and his

wickedness shall fall on his own pate.

I will give thanks unto the Lord, according
to his righte-

ousness: and I will praise the name of the Lord most High.

Glory be, \&c.

\emph{Psalm} 8 (6)
O Lord our Governor, how
excellent is thy Name in all the

world: thou that hast set thy glory above the heavens!

Out of the mouth of very babes and sucklings
hast thou

ordained strength, because of thine enemies: that thou

mightest still the enemy, and the avenger.

For I will consider thy heavens, even the
works of thy

fingers: the moon and the stars, which thou hast ordained.

What is man, that thou art mindful of him:
and the son of

man, that thou visitest him?

Thou madest him lower than the angels: to
crown

with glory and worship.

Thou makest him to have dominion of the works
of thy

hands: and thou hast put all things in subjection under his

feet;

All sheep and oxen: yea, and the beasts of
the field;

The fowls of the air, and the fishes of the
sea: and a what-

soever walketh through the paths of the seas.

O Lord our Governor: how excellent is thy
Name in all

the world! Glory be, \&c.

\emph{Psalms} 9 \& 10 (7)
I will give thanks unto tree,
O Lord, with my whole heart:

I will speak of all thy marvellous works.

I will be glad and rejoice in thee: yea, my
songs will I

make of thy name, O thou most Highest.

While mine enemies are driven back; they
shall fall and

perish in thy presence.

For thou hast maintained my right and my
cause: thou art

set in the throne that judgest right.

Thou hast rebuked the heathen, and destroyed
the un-

godly: thou hast put out their name for ever and ever. 

O thou enemy, destructions are come to a
perpetual end:

even as the cities which thou hast destroyed: their memo-

rial is perished with them.

But the Lord shall endure for ever: he hath
also prepared

his seat for judgment.

For he shall judge the world in
righteousness: and mi-

nister true judgment unto the people.

The Lord also will be a defence for the
oppressed: even

a refuge in due time of trouble.

And they that know thy Name will put their
trust in

thee: for thou, Lord, hast never failed them that seek thee.

O praise the Lord which dwelleth in Sion:
shew the

people of his doings.

For when he maketh inquisition for blood, he
remem-

bereth them: and forgetteth not the complaint of the poor.

Have mercy upon me, O Lord; consider the
trouble

which I suffer of them that hate me: thou that liftest me

up from the gates of death.

That I may shew all thy praises within the
ports of the

daughter of Sion: I will rejoice in thy salvation.

The heathen are sunk down in the pit that
they made: in

the same net which they hid privily, is their foot taken.

The Lord is known to execute judgment: the
ungodly is

trapped in the work of his own hands.

The wicked shall be turned into hell: and all
the people

that forget God.

For the poor shall not alway be forgotten:
the patient

abiding of the meek shall not perish for ever.

Up, Lord, and let not man have the upper
hand: let the

heathen be judged in thy sight.

Put them in fear, O Lord: that the heathen
may know

themselves to be but men.

Why standest thou so far off, O Lord: and
hidest thy face

in the needful time of trouble?

The ungodly for his own lust doth persecute
the poor:

let them be taken in the crafty wiliness that they have

imagined.

For the ungodly hath made boast of his own
heart's desire:

and speaketh good of the covetous, whom God abhorreth.

The ungodly is so proud, that he careth not
for God:

neither is God in all his thoughts.

His ways are always grievous: thy judgments
are far above

out of his sight, and therefore defieth he all his enemies.

For he hath said in his heart, Tush, I shall
never be cast

down: there shall no harm happen unto me.

His mouth is full of cursing, deceit, and
fraud: under

his tongue is ungodliness and vanity. 

He sitteth lurking in the thievish corners of
the streets:

and privily in his lurking dens doth he murder the innocent:

his eyes are set against the poor.

For he lieth waiting secretly, even as a lion
lurketh he in

his den: that he may ravish the poor.

He doth ravish the poor: when he getteth him
into his

net.

He falleth down, and humbleth himself: that
the congre-

gation of the poor may fall into the hands of his captains.

He hath said in his heart, Tush, God hath
forgotten: he

hideth away his face, and he will never see it.

Arise, O Lord God, and lift up thine hand:
forget not

the poor.

Wherefore should the wicked blaspheme God:
while he

doth say in his heart, Tush, thou God carest not for it.

Surely thou hast seen it: for thou beholdest
ungodliness

and wrong.

That thou mayest take the matter into thy
hand: the

poor committeth himself unto thee; for thou art the helper

of the friendless.

Break thou the power of the ungodly and
malicious: take

away his ungodliness, and thou shalt find none.

The Lord is King for ever and ever: and the
heathen are

perished out of the land.

Lord, thou hast heard the desire of the poor:
thou pre-

parest their heart, and thine ear hearkeneth thereto.

To help the fatherless and poor unto their
right: that the

man of the earth be no more exalted against them. Glory

be, \&c.

\emph{Psalm} 11 (8)
In the Lord put I my trust:
how say ye then to my soul,

that she should flee as a bird unto the hill?

For, lo, the ungodly bend their bow, and make
ready their

arrows within the quiver: that they may privily shoot at

them which are true of heart.

For the foundations will be cast down: and
what hath the

righteous done

The Lord is in his holy temple: the Lord's
seat is in

heaven.

His eyes consider the poor: and his eyelids
try the chil-

dren of men.

The Lord alloweth the righteous: but the
ungodly, and

him that delighteth in wickedness, doth his soul abhor.

Upon the ungodly he shall rain snares, fire
and brimstone,

storm and tempest: this shall be their portion to drink.

For the righteous Lord loveth righteousness:
his counte-

nance will behold the thing that is just. Glory be, \&c. (d) \emph{Antiphon}
God is a righteous judge,
strong and patient: shall God

be angry every day?

(e) \emph{Antiphon}

Thou shalt keep them, O Lord.

\emph{Psalm} 12 (9)
Help me, Lord, for there is
not one godly man left: for

the faithful are minished from among the children of men.

They talk of vanity every one with his
neighbour: they

do but flatter with their lips, and dissemble in their double

heart.

The Lord shall root out all deceitful lips:
and the tongue

that speaketh proud things.

Which have said, with our tongue will we
prevail: we are

they that ought to speak; who is Lord over us?

Now for this comfortless troubles' sake of
the needy: and

because of the deep sighing of the poor;

I will up, saith the Lord: and will help
every one from

him that swelleth against him, and will set him at rest.

The words of the Lord are pure words: even as
the

silver which from the earth is tried, and purified seven times

in the fire.

Thou shalt keep them, O Lord: thou shalt
preserve him

from this generation for ever.

The ungodly walk on every side: when they are
exalted,

the children of men are put to rebuke. Glory be, \&c.

\emph{Psalm} 13 (10)
How long wilt thou forget me,
O Lord, for ever: how

long wilt thou hide thy face from me?

How long shall I seek counsel in my soul, and
be so

vexed in my heart: how long shall mine enemies triumph

over me?

Consider and hear me, O Lord my God: lighten
mine

eyes that I sleep not in death;

Lest mine enemy say, I have prevailed against
him: for

if I be cast down, they that trouble me will rejoice at it.

But my trust is in thy mercy: and my heart is
joyful in

thy salvation.

I will sing of the Lord, because he hath
dealt so lovingly

with me: yea, I will praise this Name of the Lord most

Highest. Glory be, \&c.

\emph{Psalm} 14 (11)
The fool hath said in his
heart: There is no God.

They are corrupt, and become abominable in
their doings:

there is none that doeth good, no, not one.

The Lord looked down from heaven upon the
children of

men: to see if there were any that would understand, and

seek after God.

But they are all gone out of the way, they
are altogether be-

come abominable: there is none that doeth good, no not one. 

Their throat is an open sepulchre; with their
tongues

have they deceived: the poison of asps is under their lips.

Their mouth is full of cursing and
bitterness: their feet

are swift to shed blood.

Destruction and unhappiness is in their ways,
and the

way of peace have they not known: there is no fear of God

before their eyes.

Have they no knowledge, that they are all
such workers

of mischief: eating up my people as it were bread, and call

not upon the Lord?

There were they brought in great fear, even
where no fear

was: for God is in the generation of the righteous.

As for you, ye have made a mock at the
counsel of the

poor: because he putteth his trust in the Lord.

Who shall give salvation unto Israel out of
Sion? When

the Lord turneth the captivity of his people: then shall

Jacob rejoice, and Israel shall be glad. Glory be, \&c.

\emph{Psalm} 15 (12)

Lord, who shall dwell in thy tabernacle: or
who shall

rest upon thy holy hill?

Even he that leadeth an uncorrupt life: and
doeth the

thing which is right, and speaketh the truth from his heart.

He that hath used no deceit in his tongue,
nor done evil

to his neighbour: and hath not slandered his neighbour.

He that setteth not by himself, but is lowly
in his own

eyes: and maketh much of them that fear the Lord.

He that sweareth unto his neighbour, and
disappointeth

him not: though it were to his own hindrance.

He that hath not given his money upon usury:
nor taken

reward against the innocent.

Whoso doeth these things shall never fall.
Glory

be, \&c.

(e) \emph{Antiphon}

Thou shalt keep them, O Lord: thou shalt
preserve him.

(f) \emph{Verse and Response}

I have thought upon thy Name, O Lord, in the
night

season.

\emph{And
have kept thy law}.

\emph{The Lord's Prayer}.

(\emph{Privately till the last two petitions}.)

Our Father, \&c.

And
lead us not into temptation,

\emph{But deliver us from evil}.

\emph{Absolution}

O Lord Jesu Christ, hear the prayers of thy
servants,

and have mercy upon us, who with the Father and the Holy

Ghost, livest and reignest world without end. \emph{Amen}. [Jube, Domine,

benedicere.]

\emph{Benediction} 1.

\emph{Lesson} 1.

1 Sam. xvii. 1-7.

\emph{Reader}.—Sir, pray for a blessing.

\emph{Minister}.—The Father everlasting,
bless us with a per-

petual blessing. \emph{Amen}.

Now the Philistines gathered together their
armies to

battle, and were gathered together at Shochoh, which \emph{belong- eth} to Judah, and pitched between Shochoh and Azekah,

in Ephes-dammim.

And Saul and the men of Israel were gathered
together

and pitched by the valley of Elah, and set the battle in array

against the Philistines.

And the Philistines stood on a mountain on
the one side,

and Israel stood on a mountain on the other side: and \emph{there was} a valley between them.

And there went out a champion out of the camp
of the

Philistines, named Goliath, of Gath, whose height \emph{was} six

cubits and a span.

And \emph{he had} an helmet of brass upon his
head, and he \emph{was}

armed with a coat of mail; and the weight of the coat \emph{was}

five thousand shekels of brass.

And \emph{he had} greaves of brass upon his
legs, and a target

of brass between his shoulders.

And the staff of his spear \emph{was} like a
weaver's beam; and

his spear's head \emph{weighed} six hundred shekels of iron;
and

one bearing a shield went before him.

But Thou, O Lord, have mercy upon us.

\emph{Thanks be to God}.

\emph{Response} 1.

\emph{Prepare your hearts unto the Lord, and
serve Him only; and he will deliver you
out of the hand of the Philis- tines}.

If ye do return unto the Lord with all your
hearts, then

put away the strange gods from among you.

\emph{And he will deliver you out of the hands
of the Philis- tines}.

\emph{Benediction} 2.

\emph{Reader}—Sir, pray for a blessing.

\emph{Minister}.—The only-begotten Son of
God, vouchsafe to

bless and succour us. \emph{Amen}.

\emph{Lesson} 2.

1 Sam. xvii. 8-11.

And he stood and cried unto the armies of
Israel, and said

unto them, Why are ye come out to set \emph{your} battle in
array?

\emph{am} I not a Philistine, and ye servants to Saul? choose
you a

man for you, and let him come down to me.

If he be able to fight with me, and to kill
me, then will 

we be your servants: but if I prevail against him, and kill

him, then shall ye be our servants, and serve us.

And the Philistine said, I defy the armies of
Israel this

day: give me a man that we may fight together.

When Saul and all Israel heard those words of
the

Philistine, they were dismayed, and greatly afraid.

But Thou, O Lord, have
mercy upon us.

\emph{Thanks be to God}.

\emph{Response} 2.

\emph{God is the avenger of all men. He hath
sent His angel,}

\emph{and taken me from my
Father's sheep; and hath}

\emph{anointed me with the
oil of His mercy}.

The Lord hath delivered me out of the mouth
of the lion,

and out of the paw of the bear.

\emph{And hath anointed me with the oil of His
mercy}.

\emph{The Benediction} 3.

\emph{Reader}—Sir, pray for a blessing.\emph{}

\emph{Minister}.—The grace of the Holy
Ghost illuminate our

minds and hearts. \emph{Amen}.

\emph{Lesson} 3.

1 Sam. xvii. 12-16.

Now David \emph{was} the son of that
Ephrathite of Beth-lehem-

judah, whose name was Jesse; and he had eight sons: and

the man went among men \emph{for} an old man in the days of

Saul.

And the three eldest sons of Jesse went, \emph{and}
followed

Saul to the battle: and the names of his three sons that

went to the battle \emph{were} Eliab the first-born; and next
unto

him, Abinadab; and the third, Shammah.

And David was the youngest; and the three
eldest fol-

lowed Saul.

But David went and returned from Saul, to
feed his

father's sheep at Bethlehem.

And the Philistine drew near, morning and
evening, and

presented himself forty days.

But thou, O Lord, have mercy upon us.

\emph{Thanks be to God}.

\emph{Response} 3.

\emph{The Lord, who delivered me out of the
mouth of the lion,}

\emph{and out of the paw of
the bear, He shall deliver me}

\emph{from the hands of mine
enemies}.

God hath sent forth His mercy and truth; my
soul is

among lions.

\emph{He shall deliver me from the hands of mine
enemies}.

Glory be, \&c.

\emph{He shall deliver me from the hands of mine
enemies}.

NOCTURN
II.

(g) \emph{Antiphon}

My goods.

\emph{Psalm} 16 (13)

Preserve me, O God: for in thee I have put my
trust.

O my soul, thou hast said unto the Lord: thou
art my

God, my goods are nothing unto thee.

All my delight is upon the saints, that are
in the earth:

and upon such as excel in virtue.

But they that run after another god: shall
have great

trouble.

Their drink-offerings of blood will I not
offer: neither

make mention of their names within my lips.

The Lord himself is the portion of mine
inheritance, and

of my cup: thou shalt maintain my lot.

The lot is fallen unto me in a fair ground:
yea, I have a

goodly heritage.

I will thank the Lord for giving me warning:
my reins

also chasten me in the night-season.

I have set God always before me: for he is on
my right

hand, therefore I shall not fall.

Wherefore my heart was glad, and my glory
rejoiced: my

flesh also shall rest in hope.

For why? thou shalt not leave my soul in
hell: neither

shalt thou suffer thy Holy One to see corruption.

Thou shalt show me the paths of life: in thy
presence is

the fulness of joy: and at thy right hand there is pleasure

for evermore. Glory be, \&c.

(g) \emph{Antiphon}

My goods are nothing unto Thee; in Thee have
I put

my trust, preserve me, O God.

(h) \emph{Antiphon}

Because of the words of Thy lips.

\emph{Psalm} 17 (14)

Hear the right, O Lord, consider my
complaint: and

hearken unto my prayer, that goeth not out of feigned lips.

Let my sentence come forth from thy presence:
and let

thine eyes look upon the thing that is equal.

Thou hast proved and visited mine heart in
the night

season; thou hast tried me, and shalt find no wickedness

in me: for I am utterly purposed that my mouth shall not

offend.

Because of men's works that are done
against the words

of my lips: I have kept me from the ways of the destroyer.

O hold thou up my goings in thy paths: that
my footsteps

slip not.

I have called upon thee, O God, for thou
shalt hear me:

incline thine ear to me, and harken unto my words. 

Show thy marvellous loving-kindness, thou
that art the

Saviour of them which put their trust in thee: from such as

resist thy right hand.

Keep me as the apple of an eye: hide me under
the

shadow of thy wings.

From the ungodly that trouble me: mine
enemies com-

pass me round about to take away my soul.

They are inclosed in their own fat: and their
mouth

speaketh proud things.

They lie waiting in our way on every side:
turning their

eyes down to the ground;

Like as a lion that is greedy of his prey:
and as it were

a lion's whelp lurking in secret places.

Up, Lord, disappoint him, and cast him down:
deliver

my soul from the ungodly, which is a sword of thine;

From the men of thy hand, O Lord, from the
men, I

say, and from the evil world: which have their portion

in this life, whose bellies thou fillest with thy hid trea-

sure.

They have children at their desire: and leave
the rest

of their substance for their babes.

But as for me, I will behold thy presence in
righteous-

ness: and when I awake up after thy likeness, I shall be

satisfied with it. Glory be, \&c.

(h) \emph{Antiphon}

Because of the words of Thy lips, I have kept
me from

the ways of the destroyer.

(i) \emph{Antiphon}

I will love Thee.

\emph{Psalm} 18 (15)
… O Lord, my strength; the Lord is my stony
rock, and

my defence: my Saviour, my God, and my might, in whom

I will trust: thy buckler, the horn also of my salvation, and

my refuge.

I will call upon the Lord, which is worthy to
be praised:

so shall I be safe from mine enemies.

The sorrows of death compassed me: and the
overflow-

ings of ungodliness made me afraid.

The pains of hell came about me: the snares
of death

overtook me.

In my trouble I will call upon the Lord: and
complain

unto my God.

So shall he hear my voice out of his holy
temple: and my

complaint shall come before him, it shall enter even into his

ears.

The earth trembled and quaked: the very
foundations also

of the hills shook, and were removed, because he was wroth. 

There went a smoke out in his presence: and a
consuming

fire out of his mouth, so that coals were kindled at it.

He bowed the heavens also, and came down: and
it was

dark under his feet.

He rode upon the cherubims, and did fly: he
came flying

upon the wings of the wind.

He made darkness his secret place: his
pavilion round

about him with dark water, and thick clouds to cover him.

At the brightness of his presence his clouds
removed:

hailstones, and coals of fire.

The Lord also thundered out of heaven, and
the Highest

gave his thunder: hail-stones, and coals of fire.

He sent out his arrows, and scattered them:
he cast forth

lightnings and destroyed them.

The springs of waters were seen, and the
foundations of

the round world were discovered, at thy chiding, O Lord:

at the blasting of the breath of thy displeasure.

He shall send down from on high to fetch me:
and shall

take me out of many waters.

He shall deliver me from my strongest enemy,
and from

them which hate me: for they are too mighty for me.

They prevented me in the day of my trouble:
but the

Lord was my upholder.

He brought me forth also into a place of
liberty: he

brought me forth, even because he had a favour unto me.

The Lord shall reward me after my righteous
dealing: ac-

cording to the cleanness of my hands shall he recompense me.

Because I have kept the ways of the Lord: and
have not

forsaken my God, as the wicked doeth.

For I have an eye unto all his laws: and will
not cast out

his commandments from me.

I was also uncorrupt before him: and eschewed
mine own

wickedness.

Therefore shall the Lord reward me after my
righteous

dealing: and according unto the cleanness of my hands in

his eye-sight.

With the holy thou shalt be holy: and with a
perfect man

thou shalt be perfect.

With the clean thou shalt be clean: and with
the forward

thou shalt learn frowardness.

For thou shalt save the people that are in
adversity: and

shalt bring down the high looks of the proud.

Thou also shalt light my candle: the Lord my
God shall

make my darkness to be light.

For in thee I shall discomfit an host of men:
and with

the help of my God I shall leap over the wall. 

The way of God is an undefiled way: the word
of the

Lord also is tried in the fire; he is the defender of all them

that put their trust in him.

For who is God, but the Lord: or who hath any
strength,

except our God?

It is God that girdeth me with strength of
war; and

maketh my way perfect.

He maketh my feet like harts' feet: and
setteth me up on

high.

He teacheth mine hands to fight: and mine
arms shall

break even a bow of steel.

Thou hast given me the defence of thy
salvation: thy

right hand also shall hold me up, and thy loving correction

shall make me great.

Thou shalt make room enough under me for to
go: that

my footsteps shall not slide.

I will follow upon mine enemies, and overtake
them:

neither will I turn again till I have destroyed them.

I will smite them, that they shall not be
able to stand:

but fall under my feet.

Thou hast girded me with strength unto the
battle: thou

shalt throw down mine enemies under me.

Thou hast made mine enemies also to turn
their backs

upon me: and I shall destroy them that hate me.

They shall cry, but there shall be none to
help them:

yea, even unto the Lord shall they cry, but he shall not

hear them.

I will beat them as small as the dust before
the wind: I

will cast them out as the clay in the streets.

Thou shalt deliver me from the strivings of
the people:

and thou shalt make me the head of the heathen.

A people whom I have not known: shall serve
me.

As soon as they hear of me, they shall obey
me: but the

strange children shall dissemble with me.

The strange children shall fail: and be
afraid out of their

prisons.

The Lord liveth, and blessed be my strong
helper: and

praised be the God of my salvation.

Even the God that seeth that I be avenged:
and subdueth

the people unto me.

It is he that delivereth me from my cruel
enemies, and

setteth me up above mine adversaries: thou shalt rid me

from the wicked man.

For this cause will I give thanks unto thee,
O Lord,

among the Gentiles: and sing praises unto thy Name.

Great prosperity giveth he unto his king: and
sheweth 

loving-kindness unto David his Anointed, and unto his seed

for evermore. Glory be, \&c.

(i) \emph{Antiphon}

I will love Thee, O Lord, my strength.

(j) \emph{Verse and Response}

Thou also, O Lord, shalt light my candle.

\emph{The Lord my God shall make my darkness to
be light}.

\emph{The Lord's Prayer}

(\emph{Privately}.)

Our Father, \&c.

And lead us not into temptation.

\emph{But deliver us from
evil}.

\emph{Absolution} 2.

His pity and mercy succour us, who with the
Father and

the Holy Ghost liveth and reigneth world without end. \emph{Amen}.

\emph{Benediction} 4.

\emph{Reader}.—Sir, pray for a blessing.

\emph{Minister}.—God the Father Almighty,
be favourable and

gracious to us. \emph{Amen}.

\emph{Lesson} 4.

(\emph{Sermon of St. Augustin, Bishop}.)

The children of Israel presented themselves
against their

adversaries forty days. Which signifieth the present world

having four seasons and four parts of the earth, and in which

the Christian people ceaseth not to war against Goliath and

his host, that is, with the devil and his angels. Yet they

could not conquer, except that Christ, the true David, had

descended with his staff, even the mystery of His Cross.

For before the coming of Christ, dearly beloved brethren,

the devil was at liberty; but when Christ came, He did to-

wards him what the Gospel speaketh of, ``No man can enter

into a strong man's house and spoil his goods, unless first

he bind the strong man.'' Christ therefore came, and bound

the devil.

Thou
then, O Lord, have mercy upon us.

\emph{Thanks be to God}.

\emph{Response} 4.

\emph{Saul hath slain his
thousands, and David his ten thou- sands;
for the hand of the Lord was with him, and he slew
the Philistines, and took away the reproach from
Israel}.

Is not this he of whom they sang in the
dance, saying,

Saul hath slain his thousands, and David his ten thousands.

\emph{For the hand of the Lord was with him, and
he slew the Philistines, and took away the reproach from Israel}.

\emph{Benediction} 5.

\emph{Reader}.—Sir, pray for a blessing.

\emph{Minister}.—Christ grant to us the
joys of endless life. \emph{Amen}. \emph{Lesson} 5.

(\emph{Sermon continued}.)

But some one will say, If the devil be bound,
how hath he

still such dominion? It is true, dearly beloved, that he hath

much dominion, but it is over the lukewarm and careless,

and those who fear not the Lord in truth. For he is bound,

as a dog that is chained, able to bite none but such as are

led by a fatal recklessness to close with him. Ye know,

my brethren, the foolishness of him whom a chained dog

biteth. Only beware thou of closing with him in the

likings and lusts of this world, and he will not dare to come

to thee. He can bark, he can vex; he can in no wise bite

except those that be willing. Not his violence, but his

blandishments hurt us; he doth not extort, he winneth our

consent.

But thou, O Lord, have mercy upon us.

\emph{Thanks be to God}.

\emph{Response} 5.

\emph{Ye mountains of Gilboa,
let there be no dew, neither let}

\emph{there
be rain upon you; for there the mighty men}

\emph{of
Israel are fallen}.

All ye mountains which are in his border, the
Lord shall

visit you, and ye shall pass from Gilboa.

\emph{For there the mighty men of Israel are
fallen}.

\emph{Benediction} 6.

\emph{Reader}.—Sir, pray for a blessing.

\emph{Minister}.—God kindle the fire of his
love in our hearts.

\emph{Amen}.

\emph{Lesson} 6.

(\emph{Sermon continued}.)

Therefore David came, and found the Jewish
people fight-

ing against the devil; and when there was no one to under-

take the single combat, he, bearing the figure of Christ, pro-

ceeded to the battle, with a staff in his hand, and went out

against Goliath. And then was shadowed out in him, what

was fulfilled in the Lord Jesus Christ. For Christ, the true

David, came, who, being to fight against the spiritual Goliath,

that is, the, devil, bare his own Cross. Ye see, my brethren,

where David struck the Philistine, in the forehead, which

had not been signed by the sign of the Cross. For as the staff

was the type of the Cross, so also the stone wherewith he

was struck, was a type of the Lord Christ.

Thou
therefore, O Lord, have mercy upon us.

\emph{Thanks be to God}.

\emph{Response} 6.

\emph{I took thee from the
sheepcote, saith the Lord, to be Shepherd over my people, and I was with thee whither- soever
thou wentest, and have cut off all thine enemies out of
thy sight.}

And I have made thee a great name, like unto
the name 

of the great men that are in the earth, and have caused thee

to rest from all thine enemies.

\emph{And I was with thee whithersoever thou
wentest, and have cut off all thine enemies out of Thy sight}.

Glory be to the Father, \&c.

\emph{And I was with thee whithersoever thou
wentest, and have cut off all thine enemies out of Thy sight}.

NOCTURN
III.

(k) \emph{Antiphon}

There is no speech.

\emph{Psalm} 19 (16)

The heavens declare the glory of God: and the
firmament

sheweth his handy work.

One day telleth another: and one night
certifieth another.

There is neither speech nor language: but
their voices

are heard among them.

Their sound is gone out into all lands: and
their words

into the ends of the world.

In them hath he set a tabernacle for the sun:
which

cometh forth as a bridegroom out of his chamber, and re-

joiceth as a giant to run his course.

It goeth from the uttermost part of the
heaven, and run-

eth about unto the end of it again: and there is nothing

hid from the heat thereof.

The law of the Lord is an undefiled law,
converting the

soul: the testimony of the Lord is sure, and giveth wisdom

unto the simple.

The statutes of the Lord are right, and
rejoice the heart,

the commandment of the Lord is pure, and giveth light

unto the eyes.

The fear of the Lord is clean, and endureth
for ever; the

judgments of the Lord are true, and righteous altogether.

More to be desired are they than gold, yea,
than much

fine gold: sweeter also than honey, and the honeycomb.

Moreover by them is thy servant taught: and
in keeping

of them there is great reward.

Who can tell how oft he offendeth: O cleanse
thou me

from my secret faults!

Keep thy servant also from presumptuous sins,
lest they

get the dominion over me: so shall I be undefiled, and in-

nocent from the great offence.

Let the words of my mouth, and the meditation
of my

heart: be always acceptable in thy sight.

O Lord: my strength, and my redeemer. Glory
be, \&c. (k) \emph{Antiphon}

There is neither speech nor language, but
their voices are

heard among them.

(l) \emph{Antiphon}

The Lord hear thee.

\emph{Psalm} 20 (17)
… in the day of trouble: the name of the God of

Jacob defend thee.

Send thee help from the sanctuary: and
strengthen thee

out of Sion.

Remember all thy offerings: and accept thy
burnt

sacrifice.

Grant thee thy heart's desire: and fulfil
all thy mind.

We will rejoice in thy salvation, and triumph
in the Name

of the Lord our God: the Lord perform all thy petitions.

Now know I that the Lord helpeth his
Anointed, and will

hear him from his holy heaven: even with the wholesome

strength of his right hand.

Some put their trust in chariots, and some in
horses: but

we will remember the name of the Lord our God.

They are brought down, and fallen: but we are
risen, and

stand upright.

Save, Lord, and hear us, O king of heaven:
when we

call upon thee. Glory be, \&c.

(l) \emph{Antiphon}

The Lord hear thee in the day of trouble.

(m) \emph{Antiphon}

The king.

\emph{Psalm} 21 (18)
… shall rejoice in thy strength, O Lord:
exceeding

glad shall he be of thy salvation.

Thou hast given him his heart's desire: and
hast not

denied him the request of his lips.

For thou shalt prevent him with the blessings
of goodness:

and shalt set a crown of pure gold upon his head.

He asked life of thee, and thou gavest him a
long life;

even for ever and ever.

His honour is great in thy salvation: glory
and great

worship shalt thou lay upon him.

For thou shalt give him everlasting felicity:
and make

him glad with the joy of thy countenance.

And why? because the king putteth his trust
in the Lord:

and in the mercy of the Most Highest he shall not miscarry.

All thine enemies shall feel thy hand; thy
right hand

shall find out them that hate thee.

Thou shalt make them like a fiery oven in
time of thy

wrath: the Lord shall destroy them in his displeasure, and

the fire shall consume them.

Their fruit shalt thou root out of the earth:
and their

seed from among the children of men. 

For they intended mischief against thee: and
imagined

such a device as they are not able to perform.

Therefore shalt thou put them to flight: and
the strings of

thy bow shalt thou make ready against the face of them.

Be thou exalted, Lord, in thine own strength:
so will we

sing, and praise thy power. Glory be, \&c.

(m) \emph{Antiphon}

The King shall rejoice in Thy strength, O
Lord.

(n) \emph{Verse and Response}

Be Thou exalted, Lord, in thine own strength.

\emph{So will we sing and
praise Thy power}.

\emph{The Lord's Prayer}\\
(\emph{Privately}.)

Our Father, \&c.

And
lead us not into temptation,

\emph{But deliver us from evil}.

\emph{Absolution} 3.

The Almighty and merciful Lord absolve us
from the

bonds of our sins. \emph{Amen}.

\emph{Benediction} 7.

\emph{Reader}.—Sir, pray for a blessing.

\emph{Minister}.—The reading of the Gospel
be to us salvation

and a defence. \emph{Amen}.

\emph{Lesson} 7.

(Luke v.)

At that time, as the people pressed upon
Jesus to hear the

word of God, he stood by the lake of Gennesareth.

And saw two ships standing by the lake; but
the fisher

-men were gone out of them, and were washing their nets.

And he entered into one of the ships, which
was Simon's,

and prayed him that he would thrust out a little from the

land. And he sat down and taught the people out of the ship.

Now when he had left speaking, he said unto
Simon,

Launch out into the deep, and let down your nets for a

draught.

And Simon answering, said unto him, Master,
we have

toiled all the night, and have taken nothing: nevertheless,

at thy word I will let down the net.

And when they had this done, they enclosed a
great mul-

titude of fishes: and their net brake.

And they beckoned unto \emph{their} partners
which were in the

other ship, that they should come and help them. And they

came, and filled both the ships, so that they began to sink.

When Simon Peter saw \emph{it}, he fell down
at Jesus' knees,

saying, Depart from me; for I am a sinful man, O Lord, \&c.

\emph{Homily of St. Ambrose, Bishop}

When the Lord wrought His various kinds of
healing

among many, the multitude could not be restrained, either

by time or place, from their eagerness to be healed. The

evening fell on them; they followed. The water met them:

they pressed on. Therefore he entered into the ship of

Peter. This is that ship, which according to St. Matthew

ever tosseth, according to St. Luke is filled with fishes, as an

emblem of the Church labouring in its beginning, overflow-

ing in its latter time; for by fishes are meant those who es-

cape the surge of this life. Christ still sleeps in the one amid

His disciples, and still directs them in the other; for He

sleeps to the fearful, but he watches over the perfect.

Thou then, O Lord, have mercy upon us.

\emph{Thanks be to God}.

\emph{Response} 7.

\emph{My transgressions are more in number than
the sands of}

\emph{the sea, and my sins
have multiplied; and I am not wor-}

\emph{thy to look up to the
height of heaven for the multitude of}

\emph{mine iniquities, for I
have provoked Thine anger; and}

\emph{done evil in Thy
sight}.

For I acknowledge my faults and my sin is
ever before

thee, for against Thee only have I sinned.

\emph{And done evil in Thy sight}.

\emph{Benediction} 8.

\emph{Reader}.—Sir, pray for a blessing.

\emph{Minister}.—The help of God abide with
us for ever. \emph{Amen}.

\emph{Lesson} 8.

(\emph{Homily continued}.)

No distress befalls this ship, in which
wisdom is pilot,

misbelief is unknown, faith fills the sails. For how can she

be in distress, being governed by him who is the very stay

of the Church? Therefore distress is only there where there

is scanty faith; where love is perfect, there is security.

Though others were told to let down their nets, to Peter

alone He said, Launch out into the deep, that is, the depth

of contemplation. For what is so deep as to see the depth

of His riches, to know the Son of God, to confess His Divine

Generation? Which, though man's understanding cannot

compass fully by searching out, the full assurance of faith

embraceth.

But thou, O Lord, have mercy upon us.

\emph{Thanks be to God}.

\emph{Response} 8.

(\emph{used on the Sundays after Trinity}.)

\emph{The two Seraphims cried
one to the other, Holy, holy,}

\emph{holy,
Lord God of hosts: all the earth is full of}

\emph{His
glory}.

There are three that bear record in heaven,
the Father,

the Word, and the Holy Ghost, and these Three are One.

\emph{Holy, holy, holy, Lord God of Hosts, all
the earth is full of His glory}.

Glory be, \&c.

\emph{All the earth is full of His glory}.

\emph{Benediction} 9.

\emph{Reader}.—Sir, pray for a blessing.

\emph{Minister}.—The King of Angels lead us
on to the society

of the inhabitants of heaven. \emph{Amen}. \emph{Lesson} 9.

(\emph{Homily continued}.)

For though I may not know how He was born,
yet I may

not be ignorant that He was born. The mode of His gene-

ration I know not, the Author of it I confess. We were not

witnesses of it, but we are witnesses of its revelation. If we

believe not God, whom shall we believe? For all belief

comes from sight or hearing. Sight is often deceived,

hearing is of faith.

\emph{Te Deum}.

We praise thee, O God: we acknowledge thee to
be the

Lord.

All the earth doth worship thee: the Father
everlasting.

To thee all Angels cry aloud: the Heavens,
and all the

Powers therein.

To thee Cherubim and Seraphim continually do
cry,

Holy, Holy, Holy Lord God of Sabaoth;

Heaven and earth are full of the Majesty: of
thy Glory.

The glorious company of the Apostles: praise
thee.

The goodly fellowship of the Prophets: praise
thee.

The noble army of Martyrs: praise thee.

The holy Church throughout all the world:
doth acknow-

ledge thee;

The Father: of an infinite Majesty;

Thine honourable, true: and only Son;

Also the Holy Ghost: The Comforter.

Thou art the King of Glory: O Christ.

Thou art the everlasting Son: of the Father.

When thou tookest upon thee to deliver man:
thou didst

not abhor the Virgin's womb.

When thou hadst overcome the sharpness of
death: thou

didst open the Kingdom of Heaven to all believers.

Thou sittest at the right hand of God: in the
Glory of the

Father.

We believe that thou shalt come: to be our
Judge.

We therefore pray thee, help thy servants:
whom thou

hast redeemed with thy precious blood.

Make them to be numbered with thy Saints: in
glory

everlasting.

O Lord, save thy people: and bless thine
heritage.

Govern them: and lift them up for ever.

Day by day: we magnify thee;

And we worship thy Name: ever world without
end.

Vouchsafe, O Lord: to keep us this day
without sin.

O Lord, have mercy upon us: have mercy upon
us.

O Lord, let thy mercy lighten upon us: as our
trust is in

thee.

O Lord, in thee have I trusted: let me never
be con-

founded.

\xsubsection{LAUDS SERVICE}

O God, make speed to save me.

\emph{O Lord, make haste to
help me}.

Glory
be, \&c. Amen.

Hallelujah.

(o) \emph{Antiphon}

Praise ye the Lord.

\emph{Psalm} 93 (1)

The Lord is King, and hath put on glorious
apparel:

the Lord hath put on his apparel, and girded himself

with strength

He hath made the round world so sure: that it
cannot be

moved.

Ever since the world began, hath thy seat
been prepared:

thou art from everlasting.

The floods are risen, O Lord, the floods have
lift up their

voice: the floods lift up their waves.

The waves of the sea are mighty, and rage
horribly: but

yet the Lord, who dwelleth on high, is mightier.

Thy testimonies, O Lord, are very sure:
holiness be-

cometh thine house for ever. Glory be, \&c.

(o) \emph{Ant}.

(p) \emph{Ant}.

\emph{Psalm} 100 (2)

O be joyful in the Lord, all ye lands: serve
the Lord with

gladness, and come before his presence with a song.

Be ye sure that the Lord he is God: it is he
that hath

made us, and not we ourselves; we are his people, and the

sheep of his pasture.

O go your way into his gates with
thanksgiving, and into

his courts with praise: be thankful unto him, and speak

good of his name.

For the Lord is gracious, his mercy is
everlasting: and his

truth endureth from generation to generation. Glory be, \&c.

(p) \emph{Ant}.

(q) \emph{Ant}.

\emph{Psalm} 63

\emph{and} 67 (3)

O God, thou art my God: early will I
seek thee.

My soul thirsteth for thee: my flesh also longeth after

thee: in a barren and dry land, where no water is.

Thus have I looked for thee in holiness: that I might

behold thy power and glory.

For thy loving-kindness is better than the life itself: my

lips shall praise thee.

As long as I live, will I magnify thee on this manner:

and lift up my hands in thy name.

My soul shall be satisfied, even as it were with marrow

and fatness: when my mouth praiseth thee with joyful lips. 

Have I not remembered thee in my bed, and thought upon

thee when I was waking?

Because thou hast been my helper: therefore under the

shadow of thy wings will I rejoice.

My soul hangeth upon thee: thy right hand hath up-

holden me.

These also that seek the hurt of my soul: they shall go

under the earth.

Let them fall upon the edge of the sword: that they may

be a portion for foxes.

But the King shall rejoice in God; all they also that

swear by him, shall be commended: for the mouth of them

that speak lies shall be stopped.

God be merciful unto us, and bless us: and shew us the

light of his countenance, and be merciful unto us;

That thy way may be known upon earth: thy saving health

among all nations.

Let the people praise thee, O God: yea, let all the people

praise thee.

O let the nations rejoice and be glad: for thou shalt judge

the folk righteously, and govern the nations upon earth.

Let the people praise thee, O God; let all the people praise

thee.

Then shall the earth bring forth her increase; and God,

even our own God, shall give us his blessing.

God shall bless us: and all the ends of the world shall fear

him. Glory be, \&c.

(q) \emph{Antiphon}

Praise ye the Lord. Praise ye the Lord.

(r) \emph{Antiphon}

The three children.

\emph{Song of the
Three Children}(4)

O all ye works of the Lord, bless ye the Lord: praise

him and magnify him for ever.

O ye angels of the Lord, bless ye the Lord: praise him,

and magnify him for ever.

O ye Heavens, bless ye the Lord: praise him, and
mag-

nify him for ever.

O ye Waters that be above the firmament, bless ye the

Lord: praise him, and magnify him for ever.

O all ye Powers of the Lord, bless ye the Lord: praise

him, and magnify him for ever.

O ye Sun and Moon, bless ye the Lord: praise him, and

magnify him for ever.

O ye Stars of Heaven, bless ye the Lord: praise him, and

magnify him for ever.

O ye Showers and Dew, bless ye the Lord: praise ye him,

and magnify him for ever. 

O ye Winds of God, bless ye the Lord: praise
him, and

magnify him for ever.

O ye Fire and Heat, bless ye the Lord: praise
him, and

magnify him for ever.

O ye Winter and Summer, bless ye the Lord:
praise him,

and magnify him for ever.

O ye Dews and Frosts, bless ye the Lord:
praise him, and

magnify him for ever.

O ye Frost and Cold, bless ye the Lord:
praise him, and

magnify him for ever.

O ye Ice and Snow, bless ye the Lord: praise
him, and

magnify him for ever.

O ye Nights and Days, bless ye the Lord:
praise him,

and magnify him for ever.

O ye Light and Darkness, bless ye the Lord:
praise him,

and magnify him for ever.

O ye Lightnings and Clouds, bless ye the
Lord: praise

him, and magnify him for ever.

O let the Earth bless the Lord: yea, let it
praise him, and

magnify him for ever.

O ye Mountains and Hills, bless ye the Lord:
praise him,

and magnify him for ever.

O all ye Green Things upon the Earth, bless
ye the Lord:

praise him, and magnify him for ever.

O ye Wells, bless ye the Lord: praise him,
and magnify

him for ever.

O ye Seas and Floods, bless ye the Lord:
praise him, and

magnify him for ever.

O ye Whales, and all that move in the waters,
bless ye the

Lord: praise him, and magnify him for ever.

O all ye Fowls of the air, bless ye the Lord
: praise him,

and magnify him for ever.

O all ye Beasts and Cattle, bless ye the
Lord: praise him,

and magnify him for ever.

O ye children of Men, bless ye the Lord:
praise him, and

magnify him for ever.

O let Israel bless the Lord: praise him, and
magnify him

for ever.

O ye Priests of the Lord, bless ye the Lord:
praise him,

and magnify him for ever.

O ye Servants of the Lord, bless ye the Lord:
praise him,

and magnify him for ever.

O ye Spirits and Souls of the Righteous,
bless ye the

Lord: praise him, and magnify him for ever.

O ye holy and humble Men of heart, bless ye
the Lord:

praise him, and magnify him for ever. 

O Ananias, Azarias, and Misael, bless ye the
Lord: praise

him, and magnify him for ever.

(r) \emph{Antiphon}

The three children, by command of the King,
were thrown

into the furnace, not fearing the flame of the fire, but saying,

Blessed be God.

(s) \emph{Antiphon}

Praise ye the Lord.

\emph{Psalm} 148, 149,

\emph{and} 150 (5)

O praise the Lord of heaven: praise him in
the height.

Praise him, all ye angels of his: praise him,
all his host.

Praise him, sun and moon; praise him, all ye
stars and

light.

Praise him, all ye heavens: and ye waters
that are above

the heavens.

Let them praise the Name of the Lord: for he
spake the

word, and they were made; he commanded, and they were

created.

He hath made them fast for ever and ever: he
hath given

them a law, which shall not be broken.

Praise the Lord upon earth: ye dragons, and
all deeps:

Fire and hail, snow and vapours: wind and
storm fulfil-

ling his word.

Mountains and all hills: fruitful trees and
all cedars;

Beasts and all cattle: worms and feathered
fowl;

Kings of the earth and all people: princes
and all judges

of the world;

Young men and maidens, old men and children,
praise the

Name of the Lord: for his Name only is excellent, and his

praise above heaven and earth.

He shall exalt the horn of his people; all
his saints shall

praise him: even the children of Israel, even the people

that serveth him.

O sing unto the Lord a new song: let the
congregation of

saints praise him.

Let Israel rejoice in him that made him: and
let the

children of Sion be joyful in their King.

Let them praise his Name in the dance: let
them sing

praises unto him with tabret and harp.

For the Lord hath pleasure in his people: and
helpeth the

meek-hearted.

Let the saints be joyful with glory: let them
rejoice in

their beds.

Let the praises of God be in their mouth: and
a two-

edged sword in their hands;

To be avenged of the heathen: and to rebuke
the people;

To bind their kings in chains: and their
nobles with links

of iron. 

That they may be avenged of them, as it is
written: Such

honour have all his saints.

O Praise God in his holiness: praise him in
the firma-

ment of his power.

Praise him in his noble acts: praise him
according to his

excellent greatness.

Praise him in the sound of the trumpet:
praise him upon

the lute and harp.

Praise him in the cymbals and dances: praise
him upon

the strings and pipe.

Praise him upon the well-tuned cymbals:
praise him

upon the loud cymbals.

Let every thing that hath breath: praise the
Lord.

(s) \emph{Antiphon}

Praise ye the Lord. Praise ye the Lord.
Praise ye the Lord.

(t) \emph{Text}

(\emph{Capitulum}.)

\emph{Rev}. vii. 12.

\emph{Minister}.—Blessing, and glory, and
wisdom, and thanks-

giving, and honour, and power, and might, be unto our

God for ever and ever. \emph{Amen}.

(u) \emph{Hymn}

[Ecce jam noctis.]
\emph{}

\emph{Thanks be to God}.

Paler have grown the shades of night,

And nearer draws the day,

Checkering the sky with streaks of light,

Since we began to pray.

To pray for mercy when we sin,

For cleansing and release,

For ghostly safety, and within

For everlasting peace.

Grant this to us, O Father, Son,

And Spirit, God of grace,

To whom all worship shall be done

In every time and place. \emph{Amen}.

(v) \emph{Verse and Response}

The Lord is king, and hath put on glorious
apparel.

\emph{The Lord hath put on
his apparel, and girded himself}

\emph{with
strength}.

(w) \emph{Ant}.

Jesus entered into the ship.

\emph{Benedictus}.

Blessed be the Lord God of Israel: for he
hath visited,

and redeemed his people:

And hath raised up a mighty salvation for us:
in the

house of his servant David.

As he spake by the mouth of his holy
prophets: which

have been since the world began; 

That we should be saved from our enemies: and
from the

hand of all that hate us;

To perform the mercy promised to our
forefathers: and to

remember his holy covenant.

To perform the oath which he sware to our
forefather

Abraham: that he would give us;

That we being delivered out of the hand of
our enemies:

might serve him without fear,

In holiness and righteousness before him: all
the days of

our life.

And thou, Child, shalt be called the prophet
of the High-

est: for thou shalt go before the face of the Lord to prepare

his ways;

To give knowledge of salvation unto his
people: for the

remission of their sins,

Through the tender mercy of our God: whereby
the Day-

spring from on high hath visited us;

To give light to them that sit in darkness,
and in the

shadow of death: and to guide our feet into the way of

peace.

(w) \emph{Antiphon}

Jesus entered into the ship, and when he was
set down

he taught the multitudes.

The
Lord be with you.

\emph{And
with thy spirit}.

(x) \emph{Collect}

(\emph{for the day and week}.)

Let us pray.

Grant to us, Lord, we beseech Thee, that the
course of

this world may be so peaceably ordered by Thy governance,

that Thy church may joyfully serve Thee in all godly quiet-

ness through the Lord. \emph{Amen}.

\emph{Antiphon}\\
(\emph{the following Commemo- rations are read at Vespers also}.)

Holy Mary, succour the wretched, help the
weak-hearted,

comfort the mourners, pray for the people, interpose for the

Clergy, intercede for the devoted females; let all feel thy

assistance who observe thy holy commemoration.

Pray for us, Holy Mother of God,

\emph{That we may be made
worthy of the promise of Christ}.

\emph{Collect}

Let us pray.

Grant, O Lord God, we beseech Thee, that we
Thy ser-

vants may ever prosper in perpetual health of body and

mind, and by the glorious intercession of the Blessed Mary,

Ever-Virgin, may be delivered from present sadness, and

enjoy eternal bliss.

\emph{Antiphon}

Glorious rulers of the earth, as they loved
each other in

this life, so in death they were not divided. 

Their sound is gone out into all lands.

\emph{And their words unto the ends of the world}.

Let us pray.

\emph{Collect}

O God, who when thy Apostle Peter walked on
the waves

savedst him from drowning with Thy right hand, and de-

livered his fellow Apostle Paul the third day from shipwreck

on the open sea, favourably hear us, and grant that, by the

deserts of both of them, we may obtain everlasting glory.

\emph{Antiphon}

Give peace in our time, O Lord, because there
is none

other that fighteth for us but only Thou, O God.

Peace be within thy walls.

\emph{And plenteousness within Thy palaces}.

Let us pray.

\emph{Collect}

O God, from whom all holy desires, all good
counsels, and

all just works do proceed, give unto Thy servants that peace

which the world cannot give; that both our hearts may be

set to obey Thy commandments, and also that by Thee we

being defended from the fear of our enemies, may pass our

time in rest and quietness, through Jesus Christ our Lord,

who liveth and reigneth with Thee in the unity of the Holy

Ghost, one God, world without end. \emph{Amen}.

The Lord be with you.

\emph{And
with thy spirit}.

Let us bless the Lord.

(y)

\emph{Thanks
be to God}.

(z)

May the souls of the faithful, through
God's mercy, rest in

peace. \emph{Amen}.

Continue

\xsection{§ 2. Service for Sunday, June 21, 1801 (continued).}

File 1 

\xsubsection{2. PRIME SERVICE}

O God, make speed to save me.

\emph{O Lord, make haste
to help me}.

Glory be, \&c. \emph{Amen}.
Hallelujah.

(aa) \emph{Hymn}\\
{}[Jam lucis orto.]

The star of morn to night succeeds,

We therefore meekly pray,

May God in all our words and deeds

Keep us from harm this day.

May he in love restrain us still

From tones of strife and words of ill,

And wrap around and close our eyes

To earth's absorbing vanities.

May wrath and thoughts that gender shame

Ne'er in our breasts abide,

And painful abstinences tame

Of wanton flesh the pride:

So when the weary day is o'er

And night and stillness come once more,

Blameless and clean from spot of earth

We may repeat, with reverent mirth,

Praise to the Father, as is meet,

Praise to the only Son,

Praise to the Holy Paraclete,

While endless ages run. \emph{Amen}.

(bb) \emph{Antiphon}
Praise ye the Lord.

\emph{Psalm} 53. (1)
Save me, O God, for thy
Name's sake: and avenge me in

thy strength.

Hear my prayer, O God: and hearken unto the
words of

my mouth.

For strangers are risen up against me: and
tyrants, which

have not God before their eyes, seek after my soul.

Behold, God is my helper: the Lord is with
them that

uphold my soul.

He shall reward evil unto mine enemies:
destroy thou

them in thy truth.

An offering of a free heart will I give thee,
and praise thy

Name, O Lord: because it is so comfortable. 

For he hath delivered me out of all my
trouble: and mine

eye hath seen his desire upon mine enemies. Glory be, \&c.

\emph{Psalm} 118. (2)
O give thanks unto the Lord,
for he is gracious: because

his mercy endureth for ever.

Let Israel now confess, that he is gracious:
and that his

mercy endureth for ever.

Let the house of Aaron now confess: that his
mercy en-

dureth for ever.

Yea, let them now that fear the Lord confess:
that his

mercy endureth for ever.

I called upon the Lord in trouble: and the
Lord heard

me at large.

The Lord is on my side: I will not fear what
man doeth

unto me.

The Lord taketh my part with them that help
me: there-

fore shall I see my desire upon mine enemies.

It is better to trust in the Lord: than to
put any confi-

dence in man.

It is better to trust in the Lord: than to
put any confi-

dence in princes.

All nations compassed me round about: but in
the Name

of the Lord will I destroy them.

They kept me in on every side, they kept me
in, I say,

on every side: but in the Name of the Lord will I destroy

them.

They came about me like bees, and are extinct
even as

the fire among the thorns: for in the name of the Lord I

will destroy them.

Thou hast thrust sore at me, that I might
fall: but the

Lord was my help.

The Lord is my strength, and my song: and is
become

my salvation.

The voice of joy and health is in the
dwellings of the

righteous: the right hand of the Lord bringeth mighty

things to pass.

The right hand of the Lord hath the
pre-eminence: the

right hand of the Lord bringeth mighty things to pass.

I shall not die, but live: and declare the
words of the

Lord.

The Lord hath chastened and corrected me: but
he hath

not given me over unto death.

Open me the gates of righteousness: that I
may go into

them, and give thanks unto the Lord.

This is the gate of the Lord: the righteous
shall enter

into it. 

I will thank thee, for thou hast heard me:
and art become

my salvation.

The same stone which the builders refused: is
become the

head-stone in the corner.

This is the Lord's doing: and it is
marvellous in our eyes.

This is the day which the Lord hath made: we
will re-

joice and be glad in it.

Help me now, O Lord: O Lord, send us now
prosperity.

Blessed be he that cometh in the name of the
Lord: we

have wished you good luck, ye that are of the house of the

Lord.

God is the Lord who hath showed us light:
bind the sa-

crifice with cords; yea, even unto the horns of the altar.

Thou art my God, and I will thank thee: thou
art my

God, and I will praise thee.

O give thanks unto the Lord, for he is
gracious: and his

mercy endureth for ever. Glory be, \&c.

\emph{Psalm} 119 (3)

1-16.
Blessed are those that are
undefiled in the way: and walk

in the law of the Lord.

Blessed are they that keep his testimonies:
and seek him

with their whole heart.

For they who do no wickedness: walk in his
ways.

Thou hast charged: that we shall diligently
keep thy com-

mandments.

O that my ways were made so direct: that I
might keep

thy statutes!

So shall I not be confounded: while I have
respect unto

all thy commandments.

I will thank thee with an unfeigned heart:
when I shall

have learned the judgments of thy righteousness.

I will keep thy ceremonies: O forsake me not
utterly.

Wherewithal shall a young man cleanse his
way: even by

ruling himself after thy word.

With my whole heart have I sought thee: O let
me not go

wrong out of thy commandments.

Thy words have I hid within my heart: that I
should not

sin against thee.

Blessed art thou, O Lord: O teach me thy statutes.

With my lips have I been telling: of all the
judgments of

thy mouth.

I have had as great delight in the way of thy
testimonies:

as in all manner of riches.

I will talk of thy commandments: and have
respect unto

thy ways.

My delight shall be in thy statutes: and I
will not forget

thy word. Glory be, \&c. \emph{Psalm} 119. (4)

17-32.

O do well unto thy servant: that I may live
and keep thy

word.

Open thou mine eyes: that I may see the
wondrous things

of thy law.

I am a stranger upon earth: O hide not thy
command-

ments from me.

My soul breaketh out for the very fervent
desire: that it

hath alway unto thy judgment.

Thou hast rebuked the proud: and cursed are
they that

do err from thy commandments.

O turn from me shame and rebuke: for I have
kept thy

testimonies.

Princes also did sit and speak against me:
but thy servant

is occupied in thy statutes.

For thy testimonies are my delight: and my
counsellors.

My soul cleaveth to the dust: O quicken thou
me, accord-

ing to thy word.

I have acknowledged my ways, and thou
heardest me: O

teach me thy statutes.

Make me to understand the way of thy
commandments:

and so shall I talk of thy wondrous works.

My soul melteth away for very heaviness:
comfort thou

me according unto thy word.

Take from me the way of lying: and cause thou
me to

make much of thy law.

I have chosen the way of truth: and thy
judgments have I

laid before me.

I have stuck unto thy testimonies: O Lord,
confound me

not.

I will run the way of thy commandments: when
thou hast

set my heart at liberty. Glory be, \&c.

\emph{The Creed of St. Athanasius}

Whosoever will be saved: before all things it
is necessary

that he hold the Catholic Faith.

Which Faith except every one do keep whole
and unde-

filed: without doubt he shall perish everlastingly.

And the Catholic Faith is this: That we
worship one God

in Trinity, and Trinity in Unity.

Neither confounding the Persons: nor dividing
the sub-

stance.

For there is one Person of the Father,
another of the Son:

and another of the Holy Ghost.

But the Godhead of the Father, of the Son,
and of the

Holy Ghost, is all one: the Glory equal, the Majesty co-

eternal.

Such as the Father is, such is the Son: and
such as the Holy

Ghost. 

The Father uncreate, the Son uncreate: and
the Holy

Ghost uncreate.

The Father incomprehensible, the Son
incomprehensible:

and the Holy Ghost incomprehensible.

The Father eternal, the Son eternal: and the
Holy Ghost

eternal.

And yet they are not three eternals: but one
eternal.

And also there are not three
incomprehensibles, nor three

uncreated: but one uncreated, and one incomprehensible.

So likewise the Father is Almighty, the Son
Almighty:

and the Holy Ghost Almighty.

And yet they are not three Almighties: but
one Almighty.

So the Father is God, the Son is God: and the
Holy

Ghost is God.

And yet they are not three Gods: but one God.

So likewise the Father is Lord, the Son Lord:
and the

Holy Ghost Lord.

And yet not three Lords: but one Lord.

For like as we are compelled by the Christian
verity: to

acknowledge every person by himself to be God and Lord;

So we are forbidden by the Catholic Religion:
to say,

There be three Gods, or three Lords.

The Father is made of none: neither created,
nor be-

gotten.

The Son is of the Father alone: not made, nor
created, but

begotten.

The Holy Ghost is of the Father and of the
Son: neither

made, nor created, nor begotten, but proceeding.

So there is one Father, not three Fathers;
one Son, not

three Sons; one Holy Ghost, not three Holy Ghosts.

And in this Trinity none is afore, or after
other: none is

greater, or less than another.

But the whole three Persons are co-eternal
together: and

co-equal.

So that in all things, as is aforesaid: the
Unity in Trinity,

and the Trinity in Unity, is to be worshipped.

He therefore that will be saved: must thus
think of the

Trinity.

Furthermore, it is necessary to everlasting
salvation: that

he also believe rightly the Incarnation of our Lord Jesus

Christ.

For the right Faith is, that we believe and
confess: that our

Lord Jesus Christ, the Son of God, is God and Man;

God, of the Substance of the Father, begotten
before the

worlds: and Man, of the Substance of his Mother, born in the

world. 

Perfect God, and perfect Man: of a reasonable
soul, and

human flesh subsisting:

Equal to the Father, as touching his Godhead:
and infe-

rior to the Father, as touching his Manhood.

Who although he be God and man: yet he is not
two, but

one Christ.

One; not by conversion of the Godhead into
flesh: but by

taking of the Manhood into God;

One altogether; not by confusion of
substance: but by

unity of Person.

For as the reasonable soul and flesh is one
man: so God

and Man is one Christ.

Who suffered for our salvation: descended
into hell, rose

again the third day from the dead.

He ascended into heaven; he sitteth on the
right hand of

the Father, God Almighty: from whence he shall come to

judge the quick and the dead.

At whose coming all men shall rise again with
their

bodies: and shall give account of their own works.

And they that have done good shall go into
life everlast-

ing: and they that have done evil, into everlasting fire.

This is the Catholic Faith: which except a
man believe

faithfully, he cannot be saved. Glory be, \&c.

(bb) \emph{Antiphon}

Praise ye the Lord, Hallelujah, Praise ye the
Lord.

(cc) \emph{Text}.

\emph{Capitulum}.

1 \emph{Tim}. i. 17.

\emph{Minister}.—To the King, eternal,
immortal, invisible, the

only wise God, be honour and glory for ever and ever. \emph{Amen}.

\emph{Thanks
be to God}.

(dd) \emph{Short Response}

\emph{O Christ, the Son of the living God, have mercy upon us.}

\emph{O Christ, \&c}.

(ee)
Thou who sittest at the right hand of the Father,

\emph{Have mercy upon us}.

Glory be to the Father, \&c.

\emph{O Christ, the Son of the living God have mercy upon us}.

O Christ, arise, help us.

\emph{And deliver us for thy Name's sake}.

(ff)

Lord have mercy.

\emph{Christ have mercy}.

Lord have mercy.

\emph{The Lord's Prayer}

(\emph{privately}.)
Our Father, \&c.

And lead us not into temptation,

\emph{But deliver us from
evil}.

\emph{The Creed}

(\emph{privately}.)

I believe, \&c.

The resurrection of the body.

\emph{And the life
everlasting}. \emph{Amen}. (\emph{In days of the double office, such as Christmas Day the Apostles' days, \&c. \&c. from ff to hh is omitted}.)
I have cried unto Three, O Lord.

\emph{And early shall my prayer come before Thee}.

O let my mouth be filled with Thy praise.

\emph{That I may sing of Thy glory and honour
all the day long}.

Lord, turn Thy face from my sins.

\emph{And blot out all my iniquities}.

Make me a clean heart, O God.

\emph{And renew a right spirit within me}.

Cast me not away from Thy presence.

\emph{And take not thy Holy Spirit from me}.

O give me the comfort of Thy help again.

\emph{And stablish me with Thy free Spirit}.

(gg)

(\emph{Here is some- times an insertion, vide} § 3.)
Our help is in the Name of the Lord.

\emph{Who hath made heaven and earth}.

\emph{Confession},

(\emph{This confession is not found in the Par- ris Breviary, 1735. Observations are made upon it above, pp. 11 and 21. It is repeated at Com- plin; that is at the end of the day, as Prime is the begin- ning}.)

\emph{Minister}.—I confess before God
Almighty, before the

Blessed Mary, Ever-Virgin, the blessed Michael Archangel,

the blessed John Baptist, the Holy Apostles Peter and Paul,

before all Saints, and you, my Brethren, that I have sinned

too much in thought, word, and deed. It is my fault, my

fault, my grievous fault. Therefore, I beseech the blessed

Mary, Ever-Virgin, the blessed Michael Archangel, the

blessed John Baptist, the holy Apostles Peter and Paul,

all Saints, and you, my Brethren, to pray the Lord our

God for me.

\emph{People}.—God Almighty have mercy upon
thee, absolve

thee from thy sins, and bring thee safe to life everlasting

\emph{Amen}.

\emph{People}.—I confess before God
Almighty, before the

Blessed Mary, Ever-Virgin, the blessed Michael Archangel,

the blessed John Baptist, the Holy Apostles Peter and Paul,

before all Saints, and thee, my Father, that I have sinned

too much in thought, word, and deed. It is my fault, my

fault, my grievous fault. Therefore, I beseech the blessed

Mary, Ever-Virgin, the blessed Michael Archangel, the

blessed John Baptist, the holy Apostles Peter and Paul,

all Saints, and thee, my Father, to pray the Lord our

God for me.

\emph{}

\emph{Minister}.—God Almighty have mercy
upon you, absolve

you from your sins, and bring you safe to everlasting life.

\emph{Amen}.

The Almighty and Merciful Lord grant to us
pardon,

absolution, and remission of all our sins. \emph{Amen}.

Vouchsafe, O Lord, to keep us.

\emph{This day without sin}.

O Lord, have mercy upon us.

\emph{Have mercy upon us}.

O Lord, let Thy mercy be showed upon us.

\emph{As we do put our trust
in Thee}.

Lord, hear our prayer.

\emph{And let our cry come
unto Thee}.

(hh)

The Lord be with you.

\emph{And with thy spirit}.

\emph{Collect}

Let us
pray.

O Lord God Almighty, who hast brought us to
the begin-

ning of this day, defend us in the same with Thy mighty

power, that we fall into no sin today, but all our words,

thoughts, and works may be ordered according to Thy

righteousness, through our Lord Jesus Christ, Thy Son,

who liveth and reigneth with Thee in the unity of the

Holy Ghost, world without end. \emph{Amen}.

The
Lord be with you.

\emph{And with thy spirit}.

Let us
bless the Lord.

\emph{Thanks be to God}.

\emph{Here the Martyr- ology is read}.

*
* *
*
*

Right
dear in the sight of the Lord.

\emph{Is the death of his saints}.

\emph{Collect}\\
(\emph{The words “Holy Mary” are not in the Monastic Bre- viaries}.)

\emph{Repeated thrice}.

Holy Mary and all the saints intercede for us
to the

Lord, that we may be worthy of His help and salvation, who

liveth and reigneth world without end. \emph{Amen}.

O God,
make speed to save me.

\emph{O Lord, make haste to help m}e.

Glory be, \&c.

\emph{As
it was, \&c}.

Lord,
have mercy.

\emph{Christ, have mercy}.

Lord,
have mercy.

\emph{The Lord's Prayer}

(\emph{privately}.)

Our
Father, \&c.

And lead us not into temptation,

\emph{But deliver us from
evil}.

O Lord, show thy Servants Thy work, and their
children

Thy glory.

\emph{And
the Glorious Majesty of the Lord our God be upon}

\emph{us: prosper Thou the' work of our hands upon us,}

\emph{O prosper Thou our handy work}.

Glory
be, \&c.

\emph{As it was}, \&c.

The
Lord be with you.

\emph{And with thy spirit}. \emph{Collect}

Let us pray.

O Lord God, King of heaven and earth,
vouchsafe this

day to direct and sanctify, to rule and govern our hearts

and bodies, our thoughts, words, and deeds in thy law, and

in the works of Thy commandments, that through Thy

most mighty protection, both here and ever, we may ob-

tain safety and freedom, O Saviour of the world, who livest

and reignest for ever. \emph{Amen}.

\emph{Benediction}

\emph{Reader}.—Sir,
pray for a blessing.

\emph{Minister}.—The
Lord Almighty order our lives and doings

in His peace. \emph{Amen}.

(ii) \emph{Short Lesson}

2 \emph{Thess}. iii. 5.

\emph{Reader}.—The Lord direct our hearts
and bodies in the

love of God, and the patient waiting for Christ.

Thou, then, O Lord, have mercy upon us.

\emph{Thanks be to God}.

\emph{Benediction}

Our help is in the name of the Lord.

\emph{Who hath made heaven
and earth}.

Bless ye.

\emph{Reader}.—May God.

\emph{Minister}.—The Lord bless us and
defend us from all evil,

and bring us to everlasting life; and may the souls of the

faithful, through God's mercy, rest in peace. \emph{Amen}.

\xsubsection{3. THIRD HOUR SERVICE}

O God, make speed to save me.

\emph{O Lord, make haste to save me}.

Glory be, \&c. Amen.
Hallelujah.

(jj) \emph{Hymn}

[Nunc Sancte

obis.]

Come, Holy Ghost, who ever One

Art with the Father and the Son;

Come, Holy Ghost, our souls possess,

With Thy full flood of Holiness.

Let mouth, and heart, and flesh combine

To herald forth our Creed divine,

And love so wrap our mortal frame,

Others may catch the living flame.

This grace on Thy redeem'd confer,

Father, Co-equal Son,

And Holy Ghost, the Comforter,

Eternal Three in One. \emph{Amen}.

(kk) \emph{Antiphon}

\emph{Psalm} 119 (1)

33-48.

Praise ye the Lord.

Teach me, O Lord, the way of thy statutes:
and I shall

keep it unto the end.

Give me understanding, and I shall keep thy
law: yea,

I shall keep it with my whole heart.

Make me to go in the path of thy
commandments: for

therein is my desire.

Incline my heart unto thy testimonies: and
not to

covetousness.

O turn away mine eyes lest they behold
vanity: and

quicken thou me in thy way.

O stablish thy word in thy servant: that I
may fear thee.

Take away the rebuke that I am afraid of: for
thy judg-

ments are good.

Behold, my delight is in thy commandments: O
quicken

me in thy righteousness.

Let thy loving mercy come also unto me, O
Lord: even

thy salvation, according unto thy word.

So shall I make answer unto my blasphemers:
for my

trust is in thy word. 

O take not the word of thy truth utterly out
of my month:

for my hope is in thy judgments.

So shall I alway keep thy law: yea, for ever
and ever.

And I will walk at liberty: for I seek thy
commandments.

I will speak of thy testimonies also, even
before kings:

and will not be ashamed.

And my delight shall be in thy commandments:
which I

have loved.

My hands also will I lift up unto thy
commandments,

which I have loved: and my study shall be in thy statutes.

Glory, \&c.

\emph{Psalm} 119. (2)

49-64.

O think upon thy servant, as concerning thy
word: where-

in thou hast caused me to put my trust.

The same is my comfort in my trouble: for thy
word hath

quickened me.

The proud have had me exceedingly in
derision: yet have

I not shrinked from thy law.

For I remembered thine everlasting judgments,
O Lord:

and received comfort.

I am horribly afraid: for the ungodly that
forsake thy

law.

Thy statutes have been my songs: in the house
of my

pilgrimage.

I have thought upon thy Name, O Lord, in the
night-

season: and have kept thy law.

This I had: because I kept thy commandments.

Thou art my portion, O Lord: I have promised
to keep

thy law.

I made my humble petition in thy presence
with my whole

heart: O be merciful unto me, according to thy word.

I called mine own ways to remembrance: and
turned my

feet unto thy testimonies.

I made haste, and prolonged not the time: to
keep thy

commandments.

The congregations of the ungodly have robbed
me: but I

have not forgotten thy law.

At midnight I will rise to give thanks unto
thee: because

of thy righteous judgments.

I am a companion of all them that fear thee:
and keep

thy commandments.

The earth, O Lord, is full of thy mercy: O
teach me thy

statutes. Glory be, \&c.

\emph{Psalm} 119. (3)

65-80.

O Lord, thou hast dealt graciously with thy
servant: ac-

cording unto thy word.

O learn me true understanding and knowledge:
for I have

believed thy commandments. 

Before I was troubled, I went wrong: but now
have I

kept thy word.

Thou art good and gracious: O teach me thy
statutes.

The proud have imagined a lie against me: but
I will

keep thy commandments with my whole heart.

Their heart is as fat as brawn: but my
delight hath been

in thy law.

It is good for me that I have been in
trouble: that I may

learn thy statutes.

The law of thy mouth is dearer unto me: than
thousands

of gold and silver.

Thy hands have made me and fashioned me: O
give me

understanding, that I may learn thy commandments.

They that fear thee will be glad when they
see me: be-

cause I have put my trust in thy word.

I know, O Lord, that thy judgments are right:
and that

thou of very faithfulness hast caused me to be troubled.

O let thy merciful kindness be my comfort:
according to

thy word unto thy servant.

O let thy loving mercies come unto me, that I
may live:

for thy law is my delight.

Let the proud be confounded, for they go
wickedly about

to destroy me: but I will be occupied in thy commandments.

Let such as fear thee, and have known thy
testimonies:

be turned unto me.

O let my heart be sound in thy statutes: that
I be not

ashamed. Glory be, \&c.

(kk) \emph{Antiphon}

Praise ye the Lord. Praise ye the Lord.
Praise ye the Lord.

(ll) \emph{Text}

(\emph{Capitulum}.)

1 John iv. 16.

\emph{Minister}.—God is love, and he that
dwelleth in love,

dwelleth in God, and God in him.

\emph{Thanks be to God}.

(mm) \emph{Short Response}.

\emph{Incline my heart, O
God: to thy Testimonies.}

\emph{Incline, \&c}.

Turn away mine eyes, lest they behold vanity,
and quicken

Thou me in Thy way.

\emph{To Thy testimonies}.

Glory be to the Father, \&c.

\emph{Incline my heart, O
God: to Thy testimonies}.

I said, Lord, have mercy upon me.

\emph{Heal my soul, for I
have sinned against Thee}.

Lord
have mercy.

\emph{Christ have mercy}.

Lord
have mercy.

\emph{The Lord's Prayer}

(\emph{privately}.)

Our Father, \&c.

And lead us not into
temptation,

But deliver us from evil.

Turn us again, Lord God of hosts.

\emph{Show us the light of
Thy countenance and we shall be}

\emph{whole}.

O Christ, arise, help us.

\emph{And deliver us for thy
name's sake}.

Lord, hear my prayer.

\emph{And let my cry come
unto Thee}.

(nn) \emph{Collect}

(\emph{for the day and week}.)

The
Lord be with you.

\emph{And with thy spirit}.

Let us pray.

Grant to us, Lord, we beseech Thee, that the
course of

this world may be so peaceably ordered by Thy providence,

that Thy Church may joyfully serve Thee in all godly

quietness, through the Lord.

The
Lord be with you.

\emph{And with thy spirit}.

Let us
bless the Lord.

\emph{Thanks be to God}.

May the souls of the faithful through God's
mercy rest in

peace. \emph{Amen}.

\xsection{4. SIXTH HOUR SERVICE}

O God, make speed to save me.

\emph{O Lord, make haste to help me}.

Glory be, \&c.
Amen. Hallelujah.

(oo) \emph{Hymn}

[Rector potens,

verax Deus.]

O God, the Lord of place and time,

Who orderest all things prudently,

Brightening with beams the opening prime,

And burning in the mid-day sky.

Quench Thou the fires of hate and strife,

The wasting fever of the heart;

From perils guard our feeble life,

And to our souls Thy peace impart.

This grace on Thy redeemed confer,

Father, Co-equal Son,

And Holy Ghost, the Comforter,

Eternal Three in One.—\emph{Amen}.

(pp) \emph{Antiphon}

\emph{Psalm} 119. (1)

81-96.

Praise ye the Lord.

My soul hath longed for thy salvation: and I
have a good

hope, because of thy word.

Mine eyes long sore for thy word: saying, O
when wilt

thou comfort me?

For I am become like a bottle in the smoke:
yet do I not

forget thy statutes.

How many are the days of thy servant: when
wilt thou

be avenged of them that persecute me

The proud have digged pits for me: which are
not after

thy law.

All thy commandments are true: they persecute
me

falsely; O be thou my help.

They had almost made an end of me upon earth:
but I

forsook not thy commandments.

O quicken me after thy loving-kindness: and
so shall I

keep the testimonies of thy mouth.

O Lord, thy word: endureth for ever in
heaven.

Thy truth also remaineth from one generation
to another:

thou hast laid the foundation of the earth, and it abideth. 

They continue this day according to thine
ordinance: for

all things serve thee.

If my delight had not been in thy law: I
should have

perished in my trouble.

I will never forget thy commandments: for
with them thou

hast quickened me.

I am thine, O save me: for I have sought thy
command-

ments.

The ungodly laid wait for me, to destroy me:
but I will

consider thy testimonies.

I see that all things come to an end: but thy
command-

ment is exceeding broad. Glory be, \&c.

\emph{Psalm} 119. (2)

97-112.
Lord, what love have I unto
thy law: all the day long is

my study in it!

Thou, through thy commandments, hast made me
wiser

than mine enemies: for they are ever with me.

I have more understanding than my teachers:
for thy

testimonies are my study.

I am wiser than the aged: because I keep thy
command-

ments.

I have refrained my feet from every evil way:
that I may

keep thy word.

I have not shrunk from thy judgments: for
thou teachest

me.

O how sweet are thy words unto my throat:
yea, sweeter

than honey unto my mouth!

Through thy commandments I get understanding:
there-

fore I hate all evil ways.

Thy word is a lantern unto my feet: and a
light unto my

paths.

I have sworn, and am stedfastly purposed: to
keep thy

righteous judgments.

I am troubled above measure: quicken me, O
Lord, ac-

cording to thy word.

Let the free-will offerings of my mouth
please thee, O

Lord: and teach me thy judgments.

My soul is always in my hand: yet do I not
forget thy law.

The ungodly have laid a snare for me: but yet
I swerved

not from thy commandments.

Thy testimonies have I claimed as mine
heritage for ever:

and why? they are the very joy of my heart.

I have applied my heart to fulfil thy
statutes always: even

unto the end. Glory be, \&c.

\emph{Psalm} 119 (3)

113-128.
I hate them that imagine evil
things: but thy law do I love.

Thou art my defence and shield: and my trust
is in thy

word. 

Away from me, ye wicked: I will keep the
command-

ments of my God.

O establish me according to thy word, that I
may live:

and let me not be disappointed of my hope.

Hold thou me up, and I shall be safe: yea, my
delight

shall be ever in thy statutes.

Thou hast trodden down all them that depart
from thy

statutes: for they imagine but deceit.

Thou puttest away all the ungodly of the
earth like dross:

therefore I love thy testimonies.

My flesh trembleth for fear of thee: and I am
afraid of

thy judgments.

I deal with the thing that is lawful and
right: O give me

not over unto mine oppressors.

Make thou thy servant to delight in that
which is good:

that the proud do me no wrong.

Mine eyes are wasted away with looking for
thy health:

and for the word of thy righteousness.

O deal with thy servant according unto thy
loving-mercy:

and teach me thy statutes.

I am thy servant; O grant me understanding;
that I may

know thy testimonies.

It is time for thee, Lord, to lay to thine
hand: for they

have destroyed thy law.

For I love thy commandments: above gold and
precious

stone.

Therefore hold I straight all thy
commandments: and all

false ways I utterly abhor. Glory be, \&c.

(pp) \emph{Antiphon}
Praise ye the Lord.
Hallelujah. Praise ye the Lord.

(qq) \emph{Text}.

(\emph{Capitulum}.)

Gal. vi. 2.
\emph{Minister}.—Bear ye one
another's burdens, and so fulfil

the law of Christ.

\emph{Thanks be to God}.

(rr) \emph{Short Response}
\emph{O
Lord, Thy word endureth: for ever in heaven. O Lord, \&c.}

Thy truth also remaineth from one generation
to another.

\emph{For ever in heaven}.

Glory be to the Father, \&c.

\emph{O Lord, thy word
endureth: for ever in heaven}.

The Lord is my shepherd, therefore can I lack
nothing.

\emph{He shall feed me
in a green pasture}.

Lord have mercy.

\emph{Christ have mercy}.

Lord have mercy.

\emph{The Lord's Prayer}

(\emph{privately}.)

Our Father, \&c.

And lead us not into temptation.

\emph{But deliver us
from evil}. 

Turn us again, O Lord God of hosts.

\emph{Show the light of
Thy countenance, and we shall be whole}.

O Christ arise, help us.

\emph{And deliver us for
Thy name-sake}.

O Lord, hear my prayer.

\emph{And let my crying
come unto Thee}.

(ss) \emph{Collect}

(\emph{for the day, \&c}.)

The Lord be with you.

\emph{And with thy spirit}.

Let us pray.

Grant to us, we beseech Thee, that the course
of this

world may be so peaceably ordered by thy governance, that

Thy Church may joyfully serve Thee in all godly quietness,

through the Lord.

The Lord be with you.

\emph{And with thy spirit}.

Let us bless the Lord.

\emph{Thanks be to God}.

May the souls of the faithful, through
God's mercy, rest

in peace. \emph{Amen}.

\xsubsection{5. NINTH HOUR SERVICE}

O God, make speed to save me.

\emph{O Lord, make haste to help me}.

Glory be, \&c.
Amen. Hallelujah.

(tt) \emph{Hymn}

[Rerum Deus

tenax vigor.]

O God, unchangeable and true,

Of all the life and power,

Dispensing light and silence through

Every successive hour.

Lord, brighten our declining day,

That it may never wane,

Till death, when all things round decay,

Brings back the morn again.

This grace on thy redeemed confer,

Father, Co-equal Son,

And holy Ghost, the Comforter,

Eternal Three in One. \emph{Amen}.

(uu) \emph{Antiphon}

\emph{Psalm} 119. (1)

119-144.
Praise ye the Lord.

Thy testimonies are wonderful: therefore doth
my soul

keep them.

When thy word goeth forth: it giveth light
and under-

standing unto the simple.

I opened my mouth, and drew in my breath: for
my de-

light was in thy commandments.

O look thou upon me, and be merciful unto me;
as thou

usest to do unto those that love thy Name.

Order my steps in thy word: and so shall no
wickedness

have dominion over me.

O deliver me from the wrongful dealings of
men; and so

shall I keep thy commandments.

Shew the light of thy countenance upon thy
servant: and

teach me thy statutes.

Mine eyes gush out with water: because men
keep not

thy law.

Righteous art thou, O Lord: and true is thy
judgment!

The testimonies that thou hast commanded: are
exceed-

ing righteous and true.

My zeal hath even consumed me, because mine
enemies

have forgotten thy words. 

Thy word is tried to the uttermost: and thy
servant

loveth it.

I am small and of no reputation: yet do I not
forget thy

commandments.

Thy righteousness is an everlasting
righteousness: and thy

law is the truth.

Trouble and heaviness have taken hold upon
me: yet is

my delight in thy commandments.

The righteousness of thy testimonies is
everlasting: O

grant me understanding, and I shall live. Glory be, \&c.

\emph{Psalm} 119. (2)

145-160.
I call with my whole heart:
hear me, O Lord, I will keep

thy statutes.

Yea, even unto thee do I call: help me, and I
shall keep

thy testimonies.

Early in the morning do I cry unto thee: for
in thy word

is my trust.

Mine eyes prevent the night-watches: that I
might be

occupied in thy words.

Hear my voice, O Lord, according unto thy
loving-kind-

ness: quicken me according as thou art wont.

They draw nigh that of malice persecute me:
and are far

from thy law.

Be thou nigh at hand, O Lord, for all thy
commandments

are true.

As concerning thy testimonies I have known
long since:

that thou hast grounded them for ever.

O consider mine adversity, and deliver me:
for I do not

forget thy law.

Avenge thou my cause, and deliver me: quicken
me ac-

cording to thy word.

Health is far from the ungodly: for they
regard not thy

statutes.

Great is thy mercy, O Lord: quicken me, as
thou art wont.

Many there are that trouble me, and persecute
me: yet

do I not swerve from thy testimonies.

It grieveth me when I see the transgressors:
because they

keep not thy law.

Consider, O Lord, how I love thy
commandments: O

quicken me according to thy loving-kindness.

Thy word is true from everlasting: all the
judgments of

thy righteousness endure for evermore. Glory be, \&c.

\emph{Psalm} 119. (3)

161-176.
Princes have persecuted me
without a cause: but my heart

standeth in awe of thy word.

I am as glad of thy word: as one that findeth
great spoils.

As for lies, I hate and abhor them: but thy
law do I

love. 

Seven times a day do I praise thee: because
of thy righte-

ous judgments.

Great is the peace that they have who love
thy law: and

they are not offended at it.

Lord, I have looked for thy saving health:
and done after

thy commandments.

My soul hath kept thy testimonies: and loved
them ex-

ceedingly.

I have kept thy commandments and testimonies:
for all

my ways are before thee.

Let my complaint come before thee, O Lord:
give me un-

derstanding according to thy word.

Let my supplication come before thee: deliver
me accord-

ing to thy word.

My lips shall speak of thy praise: when thou
hast taught

me thy statutes.

Yea, my tongue shall sing of thy word: for
all thy com-

mandments are righteous.

Let thine hand help me: for I have chosen thy
command-

ments.

I have longed for thy saving health, O Lord:
and in thy

law is my delight.

O let my soul live, and it shall praise thee:
and thy judg-

ments shall help me.

I have gone astray like a sheep that is lost:
O seek thy

servant; for I do not forget thy commandments. Glory

be, \&c.

(uu) \emph{Antiphon}
Praise ye the Lord.
Hallelujah. Praise ye the Lord.

(vv) \emph{Test}

(\emph{Capitulum}.)

1 \emph{Cor}. iv. 20.

\emph{Minister}.—For ye are bought with a
price. Therefore

glorify God in your body.

\emph{Thanks be to God}.

(ww) \emph{Short Response}
\emph{I call with my whole heart:
hear me, O Lord. I call, \&c}.

\emph{}I will keep Thy statutes.

\emph{} \emph{Hear me, O
Lord}.

\emph{}Glory be to the Father, \&c.

\emph{} \emph{I call with
my whole heart: hear me, O Lord}.

\emph{}Cleanse Thou me, O Lord, from my
secret faults.

\emph{} \emph{Keep back
Thy servant also from presumptuous sins}.

\emph{} \emph{} \emph{} \emph{}Lord have mercy.

\emph{} \emph{} \emph{} \emph{} \emph{Christ have mercy}.

\emph{} \emph{} \emph{} \emph{}Lord have mercy.

\emph{The Lord's Prayer}

(\emph{privately}.)
\emph{} \emph{} \emph{} \emph{} \emph{} \emph{}Our Father, \&c.

\emph{} \emph{}And
lead us not into temptation.

\emph{} \emph{} \emph{But deliver us from evil}. 

\emph{}Turn us again, Lord God of hosts.

\emph{} \emph{Show us the
light of Thy countenance and we shall be}\\
\emph{} \emph{} \emph{whole}.

\emph{}O Christ, arise, help us.

\emph{} \emph{And deliver
us for thy name's sake}.

\emph{}Lord, hear my prayer.

\emph{} \emph{And let my
cry come unto thee}.

(xx) \emph{Collect}

(\emph{for the day, \&c}.)

\emph{} \emph{}The
Lord be with you.

\emph{} \emph{} \emph{And with thy spirit}.

\emph{} \emph{} \emph{} \emph{}Let us pray.

\emph{}Grant to us, Lord, we beseech Thee,
that the course of this

world may be so peaceably ordered by Thy providence, that

Thy Church may joyfully serve Thee in all godly quiet-

ness, through the Lord.

\emph{} \emph{}The
Lord be with you.

\emph{} \emph{} \emph{} \emph{And with thy spirit}.

\emph{} \emph{}Let us
bless the Lord.

\emph{} \emph{} \emph{} \emph{Thanks be to God}.

\emph{}May the souls of the faithful, through
God's mercy, rest in

peace. \emph{Amen}.

\xsubsection{6. VESPER SERVICE}

O God, make speed to save me.

\emph{O Lord, make haste to help me}.

\emph{} \emph{}Glory
be, \&c. Amen. Hallelujah.

(yy) \emph{Antiphon} (1)

\emph{Psalm} 110. (1)

\emph{}The Lord said

… unto my Lord: Sit thou on my right hand, until I

make thine enemies thy footstool.

\emph{}The Lord shall send the rod of thy
power out of Sion: be

thou ruler, even in the midst among thine enemies.

\emph{}In the day of thy power shall the
people offer thee free-will

offerings with an holy worship: the dew of thy birth is of the

womb of the morning.

\emph{}The Lord aware, and will not repent:
Thou art a Priest for

ever after the order of Melchisedech.

\emph{}The Lord upon thy right hand: shall
wound even kings in

the day of his wrath.

\emph{}He shall judge among the heathen; he
shall fill the places

with the dead bodies: and smite in sunder the heads over

divers countries.

\emph{}He shall drink of the brook in the
way: therefore shall he

lift up his head. Glory be, \&c.

(yy) \emph{Antiphon} (1)
\emph{}The Lord said unto my
Lord: Sit Thou on my right hand.

(zz) \emph{Antiphon}. (2)

\emph{Psalm} 111. (2)

\emph{}All His commandments.

\emph{}I will give thanks unto the Lord with
my whole heart:

secretly among the faithful, and in the congregation.

\emph{}The works of the Lord are great:
sought out of all them

that have pleasure therein.

\emph{}His work is worthy to be praised, and
had in honour: and

his righteousness endureth for ever.

\emph{}The merciful and gracious Lord hath so
done his marvel-

lous works: that they ought to be had in remembrance.

\emph{}He hath given meat unto them that fear
him: he shall

ever be mindful of his covenant.

\emph{}He hath showed his people the power of
his works; that

he may give them the heritage of the heathen.

\emph{}The works of his hands are verity and
judgment: all his

commandments are true.

\emph{}They stand fast for ever and ever:
amid are done in truth

and equity. 

\emph{}He sent redemption unto his people: he
hath commanded

his covenant for ever; holy and reverend is his Name.

\emph{}The fear of the Lord is the beginning
of wisdom: a good

understanding have all they that do thereafter; the praise

of it endureth for ever. Glory be, \&c.

(zz) \emph{Antiphon} (2)
\emph{}All his commandments
are true: they stand fast for ever

and ever.

(aaa) \emph{Antiphon} (3)

\emph{Psalm} 112. (3)

\emph{}He hath great delight.

\emph{}Blessed is the man that feareth the
Lord; he hath great

delight in his commandments.

\emph{}His seed shall be mighty upon earth:
the generation of

the faithful shall be blessed.

\emph{}Riches and plenteousness shall be in
his house: and his

righteousness endureth for ever.

\emph{}Unto the godly there ariseth up light
in the darkness:

he is merciful, loving, and righteous.

\emph{}A good man is merciful, and lendeth:
and will guide his

words with discretion.

\emph{}For he shall never be moved: and the
righteous shall be

had in everlasting remembrance.

\emph{}He will not be afraid of any evil
tidings: for his heart

standeth fast, and believeth in the Lord.

\emph{}His heart is established, and will not
shrink: until he

see his desire upon his enemies.

\emph{}He hath dispersed abroad, and given to
the poor: and

his righteousness remaineth for ever; his horn shall be

exalted with honour.

\emph{}The ungodly shall see it, and it shall
grieve him: he

shall gnash with his teeth, and consume away; the desire

of the ungodly shall perish. Glory be, \&c.

(aaa) \emph{Antiphon} (3)
\emph{}He hath great delight
in his commandments.

(bbb) \emph{Antiphon} (4)

\emph{Psalm} 113. (4)
\emph{}The name of the Lord.

\emph{}Praise the Lord, ye servants: O praise
the name of the

Lord.

\emph{}Blessed be the Name of the Lord: from
this time forth

for evermore.

\emph{}The Lord's name is praised: from the
rising up of the

sun, unto the going down of the same.

\emph{}The Lord is high above all heathen:
and his glory above

the heavens.

\emph{}Who is like unto the Lord our God,
that hath his dwell-

ing so high: and yet humbleth himself to behold the things

that are in heaven and earth?

\emph{}He taketh up the simple out of the
dust: and lifteth the

poor out of the mire. 

\emph{}That he may set him with the princes:
even with the

princes of his people.

\emph{}He maketh the barren woman to keep
house: and to be

a joyful mother of children. Glory be, \&c.

(bbb) \emph{Antiphon} (4)
\emph{}Blessed be the name of
the Lord for evermore.

(ccc) \emph{Antiphon} (5)

\emph{Psalm} 114 \emph{and}

115. (5)

\emph{}We who live.

\emph{}When Israel came out of Egypt: and the
house of Jacob

from among the strange people,

\emph{}Judah was his sanctuary: and Israel
his dominion.

\emph{}The sea saw that, and fled: Jordan was
driven back.

\emph{}The mountains skipped like rams: and
the little hills

like young sheep.

\emph{}What aileth thee, O thou sea, that
thou fleddest: and

thou Jordan, that thou wast driven back?

\emph{}Ye mountains, that ye skipped like
rams: and ye little

hills like young sheep?

\emph{}Tremble, thou earth, at the presence
of the Lord: at the

presence of the God of Jacob;

\emph{}Who turned the hard rock into a
standing water: and

the flint-stone into a springing well.

\emph{}Not unto us, O Lord, not unto us, but
unto thy Name give

the praise: for thy loving mercy, and for thy truth's sake.

\emph{}Wherefore shall the heathen say: Where
is now their

God?

\emph{}As for our God, he is in heaven: he
hath done whatso-

ever pleased him.

\emph{}Their idols are silver and gold: even
the work of men's

hands.

\emph{}They have mouths and speak not: eyes
have they, and

see not.

\emph{}They have ears and hear not: noses
have they and smell

not.

\emph{}They have hands, and handle not; feet
have they, and

walk not: neither speak they through their throat.

\emph{}They that make them are like unto
them: and so are all

such as put their trust in them.

\emph{}But thou, house of Israel, trust thou
in the Lord: he

is their succour and defence.

\emph{}Ye house of Aaron, put your trust in
the Lord: he is

their helper and defender.

\emph{}Ye that fear the Lord, put your trust
in the Lord: he is

their helper and defender.

\emph{}The Lord hath been mindful of us, and
he shall bless us:

even he shall bless the house of Israel, he shall bless the

the house of Aaron. 

\emph{}He shall bless them that fear the
Lord: both small and

great.

\emph{}The Lord shall increase you more and
more: you and

your children.

\emph{}Ye are the blessed of the Lord: who
made heaven and

earth.

\emph{}All the whole heavens are the
Lord's: the earth hath he

given to the children of men.

\emph{}The dead praise not thee, O Lord:
neither all they that

go down into silence.

\emph{}But we will praise the Lord: from this
time forth for

evermore. Praise the Lord. Glory be, \&c.

(ccc) \emph{Antiphon} (5)
\emph{}We who live will praise
the Lord.

(ddd) \emph{Text}

(\emph{Capitulum}.)

2 \emph{Cor}. i. 3, 4.

\emph{Minister}.—Blessed be God, even the
Father of our

Lord Jesus Christ, the Father of mercies, and the God of

all comfort, who comforteth us in all our tribulation.

\emph{} \emph{Thanks be
to God}.

(eee) \emph{Hymn}

[Lucis creator

optime.]

Father of Lights, by whom each day,

\emph{}Is kindled out of night;

Who, when the heavens were made, didst lay

\emph{}Their rudiments in light:

Thou, who didst bind and blend in one,

\emph{}The glistening morn and evening pale,

Hear Thou our plaint, when light is gone

\emph{}And lawlessness and strife prevail.

Hear, lest the whelming weight of crime

\emph{}Wreck us with life in view

Lest thoughts and schemes of sense and time

\emph{}Earn us a sinner's due.

So may we knock at heaven's door,

\emph{}And strive the prize of life to win

Continually and evermore

\emph{}Guarded without and pure within.

This grace on Thy redeemed confer,

\emph{}Father, Co-equal Son,

And Holy Ghost the Comforter,

\emph{}Eternal Three in One. \emph{Amen}.

(fff) \emph{Verse and Response}

Let my prayer be set forth

\emph{In thy sight as the incense}.

(ggg)

\emph{Ant. Magnificat}
Master, we have toiled all the night.

\emph{}My soul doth magnify the Lord: and my
spirit hath

rejoiced in God my Saviour. 

\emph{}For he hath regarded: the lowliness of
his handmaiden.

\emph{}For, behold, from henceforth: all
generations shall call

me blessed.

\emph{}For he that is mighty hath magnified
me: and holy is

his Name.

\emph{}And his mercy is on them that fear
him: throughout all

generations.

\emph{}He hath shewed strength with his arm:
he hath scattered

the proud in the imagination of their hearts.

\emph{}He hath put down the mighty from their
seat: and hath

exalted the humble and meek.

\emph{}He hath filled the hungry with good
things: and the rich

he hath sent empty away.

\emph{}He remembering his mercy, hath holpen
his servant

Israel: as he promised to our forefathers, Abraham and

his seed for ever. Glory be, \&c.

(ggg) \emph{Antiphon}
\emph{}Master, we have toiled
all the night and have taken

nothing; nevertheless, at Thy word I will let down the net.

(hhh) \emph{Collect}

(\emph{for the day, \&c}.)

\emph{}The Lord be with you.

\emph{} \emph{And with
thy spirit}.

\emph{} \emph{} \emph{} \emph{} \emph{} \emph{} \emph{} \emph{}Let us pray.

\emph{}Grant to us, Lord, we beseech Thee,
that the course of

this world may be so peaceably ordered by Thy providence,

that Thy Church may joyfully serve Thee in all godly quiet-

ness through the Lord. \emph{Amen}.

\emph{The following Com- memorations are read at Lauds also}.

\emph{Antiphon}

\emph{}Holy Mary, succour the wretched, help
the weak-hearted,

comfort the mourners, pray for the people, interpose for the

clergy, intercede for the devoted women: let all feel thy

help, whoso celebrate thy holy commemoration.

\emph{}Pray for us, Holy Mother of God.

\emph{} \emph{That we may
be made worthy of the promises of Christ}.

\emph{} \emph{} \emph{} \emph{} \emph{} \emph{} \emph{} \emph{}Let us pray.

\emph{Collect}
\emph{}Grant, O Lord God, we
beseech Thee, that we thy ser-

vants may enjoy perpetual health of mind and body, and,

by the glorious intercession of the Ever-Virgin Mary, may

be delivered from present sorrow, and so have the fruition

of everlasting joy.

\emph{Antiphon}

\emph{}Peter the Apostle, and Paul the
teacher of the Gentiles,

they have taught us Thy law, O Lord. 

\emph{}Thou shalt make them princes in all
lands.

\emph{} \emph{They shall
remember Thy name from one generation to}\\
\emph{} \emph{} \emph{another}.

\emph{} \emph{} \emph{} \emph{} \emph{} \emph{} \emph{} \emph{}Let us pray.

\emph{}O God, who when thy Apostle Peter
walked on the waves,

savedst him from drowning with Thy right hand, and thrice

delivered his fellow Apostle Paul from shipwreck on the open

sea, favourably hear us, and grant, that by the deserts of both

of them we may obtain everlasting glory.

\emph{Antiphon}

\emph{}Give peace, in our time, O Lord,
because there is none

other that fighteth for us, but only Thou, O God.

\emph{}Peace be within Thy walls.

\emph{And plenteousness within Thy palaces}.

\emph{} \emph{} \emph{} \emph{} \emph{} \emph{} \emph{} \emph{}Let us pray.

\emph{Collect}
\emph{}O God, from whom all
holy desires, all good counsels, and

all just works do proceed, give unto Thy servants that peace

which the world cannot give; that both our hearts may be

set to obey Thy commandments, and also that by Thee we

being defended from the fear of our enemies, may pass our

time in rest and quietness, through Jesus Christ our Lord,

who liveth and reigneth with Thee in the unity of the Holy

Ghost, one God, world without end. \emph{Amen}.

\emph{} \emph{} \emph{} \emph{} \emph{} \emph{}The Lord be with you.

\emph{} \emph{} \emph{} \emph{} \emph{} \emph{} \emph{And with thy spirit}.

\emph{} \emph{} \emph{} \emph{} \emph{} \emph{}Let us bless the Lord.

\emph{} \emph{} \emph{} \emph{} \emph{} \emph{} \emph{Thanks be to God}.

\emph{}May the souls of the faithful, through
God's mercy, rest

in peace. \emph{Amen}.

\xsubsection{7. COMPLINE SERVICE}

\emph{Reader}.—Sir, pray for
a blessing.

\emph{Benediction}
\emph{Minister}.—The Lord
Almighty grant us a quiet night and

an end of toils. \emph{Amen}.

\emph{Short Lesson}

1 \emph{Pet}. v. 8.
\emph{}Brethren, be sober, be
vigilant; for your adversary the

devil, as a roaring lion, walketh about, seeking whom he may

devour; whom resist stedfast in the faith.

\emph{}But Thou, O Lord, have mercy upon us.

\emph{} \emph{} \emph{} \emph{} \emph{} \emph{} \emph{} \emph{Thanks be to God}.

\emph{}Our help is in the name of the Lord,

\emph{Who hath made heaven and earth}.

\emph{The Lord's Prayer}

(\emph{privately}.)
\emph{} \emph{} \emph{} \emph{} \emph{} \emph{} \emph{} \emph{}Our Father, \&c.

(\emph{vide supra, Prime Service}.)
\emph{Minister}.—I confess
before God Almighty, before the

blessed Mary, Ever-Virgin, the Blessed Michael Archangel,

the blessed John Baptist, the Holy Apostles Peter and Paul,

before all Saints, and you, my brethren, that I have sinned

too much in thought, word, and deed. It is my fault,

my fault, my grievous fault. Therefore I beseech the

blessed Mary, Ever-Virgin, the blessed Michael Archangel,

the blessed John Baptist, the Holy Apostles Peter and Paul,

all Saints, and you, my brethren, to pray the Lord our God

for me.

\emph{People}.—\emph{Almighty God, pity thee,
absolve thee from thy sins, and bring thee safe to life everlasting}. \emph{Amen}.

\emph{People}.—I confess before God
Almighty, before the blessed

Mary, Ever-Virgin, the blessed Michael Archangel, the

blessed John Baptist, the Holy Apostles Peter and Paul,

before all Saints, and thee, my father, that I have sinned too

much in thought, word, and deed. It is my fault, my fault,

my grievous fault. Therefore I beseech the blessed Mary,

Ever-Virgin, the blessed Michael Archangel, the blessed

John Baptist, the Holy Apostles Peter and Paul, all Saints,

and thee, my father, to pray the Lord our God for me.

\emph{Minister}.—Almighty God pity you,
absolve you from your

sins, and bring you safe to everlasting life. \emph{Amen}.

\emph{}The Almighty and merciful Lord grant
to us pardon, ab-

solution, and remission of all our sins. \emph{Amen}. \emph{}Turn thou us, O Lord
our Saviour.

\emph{} \emph{And let
thine anger cease from us}.

\emph{}O God, make speed to save me.

\emph{} \emph{O Lord,
make haste to help me}.

\emph{} \emph{} \emph{} \emph{} \emph{} \emph{} \emph{}Glory be, \&c.

\emph{} \emph{} \emph{} \emph{} \emph{} \emph{} \emph{}As it was, \&c.

\emph{Antiphon}

\emph{Psalm} 4.
\emph{}Have mercy upon me.

\emph{}Hear me, when I call, O God of my
righteousness: thou

hast set me at liberty when I was in trouble: have mercy

upon me, and hearken unto my prayer.

\emph{}O ye sons of men, how long will ye
blaspheme mine honour:

and have such pleasure in vanity, and seek after leasing?

\emph{}Know this also, that the Lord hath
chosen to himself the

man that is godly: when I call upon the Lord, he will hear me.

\emph{}Stand in awe, and sin not: commune
with your own heart,

and in your chamber, and be still.

\emph{}Offer the sacrifice of righteousness:
and put your trust in

the Lord.

\emph{}There be many that say: who will show
us any good?

\emph{}Lord, lift thou up: the light of thy
countenance upon us.

\emph{}Thou hast put gladness in my heart:
since the time that

their corn, and wine, and oil increased.

\emph{}I will lay me down in peace, and take
my rest: for it is

thou, Lord, only, that makest me dwell in safety. Glory

be, \&c.

\emph{Psalm} 31.
\emph{}In thee, O Lord, have I
put my trust: let me never be put

to confusion: deliver me in thy righteousness.

\emph{}Bow down thine ear to me: make haste
to deliver me.

\emph{}And be thou my strong rock, and house
of defence: that

thou rnayest save me.

\emph{}For thou art my strong rock and my
castle: be thou also

my guide, and lead me for thy Name's sake.

\emph{}Draw me out of the net that they have
laid privily for me:

for thou art my strength.

\emph{}Into thy hands I commend my spirit:
for thou hast re-

deemed me, O Lord, thou God of truth. Glory be, \&c.

\emph{Psalm} 91.
\emph{}Whoso dwelleth under
the defence of the Most High: shall

abide under the shadow of the Almighty.

\emph{}I will say unto the Lord, Thou art my
hope, and my

strong hold: my God, in him will I trust.

\emph{}For he shall deliver thee from the
snare of the hunter: and

from the noisome pestilence.

\emph{}He shall defend thee under his wings,
and thou shalt be

safe under his feathers: his faithfulness and truth shall be thy

shield and buckler. 

\emph{}Thou shalt not be afraid for any
terror by night: nor for

the arrow that flieth by day.

\emph{}For the pestilence that walketh in
darkness: nor for the

sickness that destroyeth in the noon-day.

\emph{}A thousand shall fall beside thee, and
ten thousand at thy

right hand: but it shall not come nigh thee.

\emph{}Yea, with thine eyes shalt thou
behold: and see the reward

of the ungodly.

\emph{}For thou, Lord, art my hope: thou hast
set thine house of

defence very high.

\emph{}There shall no evil happen unto thee:
neither shall any

plague come nigh thy dwelling.

\emph{}For he shalt give his angels charge
over thee: to keep thee

in all thy ways.

\emph{}They shall bear thee in their hands:
that thou hurt not

thy foot against a stone.

\emph{}Thou shalt go upon the lion and adder:
the young lion

and the dragon shalt thou tread under thy feet.

\emph{}Because he hath set his love upon me,
therefore will I

deliver him: I will set him up, because he hath known my

Name.

\emph{}He shalt call upon me, and I will hear
him: yea, I am

with him in trouble: I will deliver him, and bring him to

honour.

\emph{}With long life will I satisfy him: and
show him my sal-

vation. Glory be, \&c.

\emph{Psalm} 134.
\emph{}Behold now, praise the
Lord: all ye servants of the Lord;

\emph{}Ye that by night stand in the house of
the Lord: even in

the courts of the house of our God,

\emph{}Lift up your hands in the sanctuary:
and praise the Lord.

\emph{}The Lord, that made heaven and earth:
give thee blessing

out of Sion. Glory be, \&c.

\emph{Antiphon}
\emph{}Have mercy upon me, O
Lord, and hearken unto my

prayer.

\emph{Hymn}\\
{}[Te lucia ante

terminum.]

Now that the day-light dies away

\emph{}Ere we lie down and sleep,

Thee, Maker of the world, we pray

\emph{}To own us and to keep.

Let dreams
depart and visions fly,

\emph{}The offspring of the night;

Keep us, like shrines, beneath Thine eye,

\emph{}Pure in our foe's despite.

This grace
on Thy redeemed confer,

\emph{}Father, Co-equal Son,

And Holy Ghost the Comforter,

\emph{}Eternal Three in One. \emph{Amen}.

\emph{Text Jer}. xiv, 9.

\emph{Minister}.—Thou, O Lord, art in the
midst of us, and we

are called by thy name; leave us not, O Lord, our God.

\emph{Thanks be to God}.

\emph{Short Response}
Into thy hands, O Lord, I commend my spirit.

\emph{Into thy hands, \&c}.

For thou hast redeemed me, O Lord God of Truth.

\emph{I commend my spirit}.

Glory be, \&c.

\emph{Into thine hands, \&c}.

Keep us, O Lord, as the apple of an eye.

\emph{Hide us under the shadow of Thy wings}.

\emph{Antiphon}
\emph{}Save us, O Lord.

\emph{Nunc Dimittis}
\emph{}Lord, now lettest thou
thy servant depart in peace: ac-

cording to thy word.

\emph{}For mine eyes have seen: thy
salvation.

\emph{}Which thou hast prepared: before the
face of all people;

\emph{}To be a light to lighten the Gentiles:
and to be the glory

of thy people Israel. Glory be, \&c.

\emph{Antiphon}
\emph{}Save us, O Lord, when
we are awake, guard us when we

are asleep, that we may be with Christ when awake, and

may rest in peace.

\emph{} \emph{} \emph{} \emph{} \emph{}Lord, have mercy.

\emph{} \emph{} \emph{} \emph{} \emph{} \emph{Christ, hare mercy}.

\emph{} \emph{} \emph{} \emph{} \emph{}Lord, have mercy.

\emph{The Lord's Prayer}

(\emph{privately}.)
\emph{} \emph{} \emph{} \emph{}Our Father, \&c.

And lead us not into temptation.

\emph{But deliver us from evil}.

\emph{The Creed}

(\emph{privately}.)
\emph{} \emph{} \emph{} \emph{}I believe, \&c.

The resurrection of the body,

\emph{And the life everlasting}. \emph{Amen}.

Blessed art Thou, O Lord God of our fathers.

\emph{Greatly to be praised, and glorious for
ever}.

Let us bless the Father and the Son, with the Holy Ghost.

\emph{Let us praise and exalt Him for ever}.

Blessed art Thou, O Lord, in the firmament of heaven.

\emph{And greatly to be praised, and glorious,
and highly exalted}\\
\emph{} \emph{for
ever}. 

\emph{}The Almighty and merciful Lord bless
and protect us.

\emph{Amen}.

Vouchsafe, O Lord, to keep us,

\emph{This night without sin}.

Have mercy upon us, O Lord,

\emph{Have mercy upon us}.

O Lord, let Thy mercy be showed upon us.

\emph{As we do put our trust in Thee}.

Lord, hear my prayer.

\emph{And let my cry come unto Thee}.

\emph{} \emph{} \emph{} \emph{} \emph{} \emph{}The Lord be with you,

\emph{} \emph{} \emph{} \emph{} \emph{} \emph{} \emph{And with thy spirit}.

\emph{} \emph{} \emph{} \emph{} \emph{} \emph{} \emph{} \emph{}Let us pray.

\emph{Collect}
\emph{}Look down, we beseech
Thee, O Lord, upon this dwell-

ing, and drive far from it all snares of the enemy! let Thy

holy angels dwell in it, and preserve us in peace, and let

Thy blessing ever be upon us, through the Lord. \emph{Amen}.

\emph{} \emph{} \emph{} \emph{} \emph{} \emph{}The Lord be with you,

\emph{} \emph{} \emph{} \emph{} \emph{} \emph{} \emph{And with thy spirit}.

\emph{} \emph{} \emph{} \emph{} \emph{} \emph{}Let us bless the Lord.

\emph{} \emph{} \emph{} \emph{} \emph{} \emph{} \emph{Thanks be given to God}.

\emph{Benediction}
\emph{}The Almighty and
merciful Lord, the Father, and the

Son, and the Holy Ghost, bless and preserve us. \emph{Amen}.

———————

(\emph{This is one of the four Anti- phons animad- verted on}, pp.

11 and 23.)
\emph{}Hail, Queen! mother of
mercy, our life, sweetness, and

hope, hail! We exiles cry to thee, the children of Eve.

To thee we sigh with groans and weeping in this valley of

tears. Come then, O our advocate, turn on us thy merciful

eyes, and shew to us, after this banishment, Jesus, the blessed

fruit of thy womb. O kind, O pitiful, O sweet Virgin Mary.

\emph{}Pray for us, holy mother of God,

\emph{} \emph{That we may
be made worthy of the promises of Christ}.

\emph{Collect}
\emph{}Almighty everlasting
God, who hast prepared the body

and soul of Mary the glorious Virgin Mother, to be a fit

dwelling for Thy Son, by the co-operation of the Holy

Ghost, grant that we, who are gladdened by this comme-

moration of her, may by her pitiful intercession be deli-

vered from present ills, and eternal death, through the

same Jesus Christ our Lord. \emph{Amen}.

\emph{}The help of God remain with us always.
\emph{Amen}.

\xsection{§ 3. Week-day Service.}

For First
Monday in Advent

O Lord, open Thou, \&c.

And my mouth, \&c.

God, make speed, \&c.

O Lord, make haste, \&c.

Glory be, \&c. Amen. Hallelujah.

(a) \emph{Invitatory}

O come let us worship the Lord, the King approaching.

O come let us sing, \&c.

Glory, \&c.

(b) \emph{Hymn}
Verbum supernum prodiens, \&c.

\emph{Antiphon Psalm} 27. (1)

\emph{Psalm} 28. (2)

\emph{Antiphon}

The Lord is the strength.

The Lord is my light, \&c. Glory, \&c.

Unto Thee will I cry, \&c. Glory, \&c.

The Lord is the strength of my life.

\emph{Antiphon Psalm} 29. (3) 

\emph{Psalm} 30. (4)

\emph{Antiphon}

Worship the Lord.

Bring unto the Lord, \&c. Glory, \&c.

I will magnify Thee, \&c. Glory, \&c.

Worship the Lord with holy worship.

\emph{Antiphon Psalm} 31. (5)

\emph{Psalm} 32. (6)

\emph{Antiphon}

Deliver me.

In the, O God, have I put, \&c. Glory, \&c.

Blessed is he, \&c. Glory, \&c.

Deliver me in Thy righteousness.

\emph{Antiphon Psalm} 33. (7)

\emph{Psalm} 34. (8)

\emph{Antiphon}

It becometh well the just.

Rejoice in the Lord, \&c. Glory, \&c.

I will alway give thanks, \&c. Glory, \&c.

It becometh well the just to be thankful.

\emph{Antiphon Psalm} 35. (9)

\emph{Psalm} 36. (10) 

\emph{Antiphon}

Fight thou against them.

Plead thou my cause, \&c. Glory, \&c.

My heart showeth me, \&c. Glory, \&c.

Fight Thou against them, that fight against me. 

\emph{Antiphon}\\
\emph{Psalm} 37. (11)

\emph{Psalm} 38. (12)

\emph{Antiphon}

Commit thy way.

Fret not thyself, \&c. Glory, \&c.

Put me not to rebuke, \&c. Glory. \&c.

Commit thy way unto the Lord.

(f.) \emph{V}.

Out of Sion hath God appeared.

\emph{Reader}.—Sir, be pleased to bless us.

\emph{The Lord's Prayer}

(\emph{privately}.)

Our Father, \&c.

And lead us not. \&c.

\emph{But deliver us,
\&c}.

\emph{Absolution} 1.

O Lord Jesu Christ, hear the prayers of Thy
servants,

and have mercy upon us, who with the Father, \&c.

\emph{Benediction} 1.

\emph{Lesson} 1.

\emph{Isai}. i. 16-18.

\emph{Reader}.—Sir, be pleased to bless us.

The Father everlasting, \&c.

Wash you, wake you clean; put away the evil
of your

doings from before Mine eyes; cease to do evil;

Learn to do well; seek judgment, relieve the
oppressed;

judge the fatherless; plead for the widow:

Come now and let us reason together, saith
the Lord;

though your sins be as scarlet, they shall be as white as

snow; though they be red like crimson, they shall be as

wool.

But Thou, O Lord, have mercy upon us.

\emph{Thanks be to God}.

\emph{Response} 1.

\emph{Cherish the word, O Virgin Mary, which is
conveyed to thee from the Lord through the Angel; thou shalt conceive in thy womb, and bring forth both God and man: that thou mayest be blessed among all women}.

Thou shalt bring forth a son, and yet abide a
virgin:

thou shalt be with child, and be a mother, yet know not a

man.

\emph{That thou mayest be blessed among women}.

\emph{Benediction} 2.

\emph{Lesson} 2.

\emph{Isai}. i. 19-23.

\emph{Reader}.—Sir, be pleased to bless us.

The only begotten Son of God, \&c.

If ye be willing and obedient, ye shall eat
the good of the

land: but if ye refuse and rebel, ye shall be devoured with

the sword, for the mouth of the Lord hath spoken it.

How is the faithful city become an harlot? It
was full of

judgment; righteousness lodged in it, but now murderers. 

Thy silver is become dross, thy wine mixed
with water;

thy princes are rebellious and companions of thieves: every

one loveth gifts, and followeth after rewards: they judge not

the fatherless, neither doth the cause of the widow come

unto them.

But thou, O Lord, have mercy upon us.

\emph{Thanks be to God}.

\emph{Response} 2.

\emph{Let the heavens rejoice, and let the earth
be glad, let the hills be joyful together, for our Lord shall come, and He shall have pity upon the poor}.

In His time shall the righteous flourish; yes, and

abundance of peace.

\emph{And he shall have
pity upon the poor}.

\emph{Benediction} 3.

\emph{Lesson} 3.

\emph{Isai}. i. 24-28.

\emph{Reader}.—Sir, be pleased to bless us.

The grace of the Holy Ghost, \&c.

Therefore saith the Lord, the Lord of hosts,
the mighty

One of Israel, Ah, I will ease Me of Mine adversaries, and

avenge me of Mine enemies.

And I will turn My hand upon thee, and purely
purge

away thy dross, and take away all thy tin;

And I will restore thy judges as at the
first, and thy

counsellors as at the beginning: afterward thou shalt be

called the city of righteousness, the faithful city.

Zion shall be redeemed with judgment, and her
courts

with righteousness.

And the destruction of the transgressors and
of the sin-

ners shall be together, and they that forsake the Lord shall

be consumed.

But thou, O Lord, have mercy upon us.

\emph{Thanks be to God}.

\emph{Response} 3.

\emph{Strangers shall not pass through Jerusalem
any more; for in that day the mountains shall drop down new wine, and the hills shall flow with milk, saith the Lord}.

God shall come from Libanus, and the holy one
from the

mount of cedars.

\emph{For in that day, the mountains shall drop
down new wine, and the hills shall flow with milk, saith the Lord}.

Glory be to the Father, \&c.

\emph{For in that day
the mountains, \&c}.

\xsubsection{LAUDS}

O God, make speed, \&c.

\emph{O Lord, make haste, \&c.}

Glory be, \&c. Amen. Hallelujah.

\emph{Antiphon}

\emph{Psalm} 51. (1)

\emph{Antiphon}
Have mercy upon me.

... O God, after
Thy great goodness, \&c. Glory, \&c.

Have mercy upon me, O God.

\emph{Antiphon}

\emph{Psalm} 5. (2)

\emph{Antiphon}

Consider, O Lord.

Ponder my word, \&c. Glory, \&c.

Consider, O Lord, my meditation.

\emph{Antiphon}

\emph{Psalms} 63 \& 67. (3)

\emph{Antiphon}

O God, Thou art my God.

… Early will I
seek Thee, \&c. Glory, \&c.

O God, Thou art my God, early will I seek thee.

\emph{Antiphon}

\emph{Canticle}. Isa. xii. (4)

\emph{Antiphon}

Thine anger is turned away.

O Lord, I will praise Thee, \&c. Glory, \&c.

Thine anger is turned away, and thou comfortedst me.

\emph{Antiphon}

\emph{Ps}. 148-151. (5)

\emph{Antiphon}

O praise the Lord.

… of heaven,
praise him in the height, \&c. Glory, \&c.

O praise the Lord of heaven.

\emph{Text}

(\emph{Capitulum}.)

(t) \emph{Isai}. ii. 3.

Come ye and let us go up to the mountain of
the Lord, to

the house of the God of Jacob; and He will teach us of His

ways, and we will walk in His paths; for out of Zion

shall go forth the Law, and the Word of the Lord from

Jerusalem.

\emph{Thanks
be to God}.

(u) \emph{Hymn}
En clara vox redarguit, \&c.

(v) \emph{Verse and Response}

The voice of one crying in the wilderness,
Prepare ye

the way of the Lord.

\emph{Make His paths straight}.

(w) \emph{Antiphon The Benedictus}

The Angel of the Lord made annunciation to Mary.

Blessed be the Lord God, \&c.

\emph{Antiphon}

The Angel of the Lord made annunciation to
Mary, and

she conceived of the Holy Ghost. Hallelujah. 

(\emph{This portion of the Service down to the Collect, is pro- per to the Week- days of Advent, the Four Seasons, and Vigils, which are Fasts, except Christ- mas Eve; it is re- peated at Vespers after the Magnifi- cat with Psalm 51 for Psalm 130. The Lord's Prayer is all aloud}.)

Lord have mercy.

\emph{Christ
have mercy}.

Lord have mercy.

Our Father which art in heaven, hallowed be
Thy name,

Thy kingdom come, Thy will be done in earth, as it is in

heaven, Give us this day our daily bread, And forgive us our

trespasses, as we forgive them that trespass against us.

And lead us not into
temptation,

\emph{But deliver us from evil}.

I said, Lord, be merciful unto me.

\emph{Heal my soul, for I have sinned against
thee}.

Turn Thee again, O Lord, at the last.

\emph{And be gracious unto Thy servants}.

Let Thy merciful kindness, O Lord, be upon us.

\emph{Like as we do put our trust in Thee}.

Let Thy priests be clothed with righteousness.

\emph{And let Thy saints sing with joyfulness}.

O Lord, save the king.

\emph{And hear us what time we call upon Thee}.

O save Thy people, and give Thy blessing unto Thine

inheritance.

\emph{Feed them and set them up for ever}.

O think upon Thy congregation.

\emph{Whom Thou hast purchased and redeemed of
old}.

Peace be within Thy walls.

\emph{And plenteousness within Thy palaces}.

Let us pray for the faithful who are departed.

\emph{Grant them, O Lord, eternal repose, and
may perpetual light shine upon them}.

May they rest in peace.

\emph{Amen}.

For our brethren who are away.

\emph{My God, save Thy servants that put their
trust in Thee}.

For the afflicted and for captives.

\emph{Deliver Israel, O Lord,
out of all his troubles}.

Send thou help from the sanctuary.

\emph{And strengthen them out
of Sion}.

Lord, hear my prayer.

\emph{And let my cry come unto Thee}.

\emph{Psalm} 130.

Out of the deep, \&c. Glory, \&c.

Turn us again, Lord God of hosts.

\emph{Show the light of Thy countenance and we
shall be whole}.

Arise, O Christ, help us.

\emph{And deliver us for Thy name's sake}.

Lord, hear my prayer.

\emph{And let my cry come unto Thee}.

The Lord be with you.

\emph{And with thy spirit}.

Let us pray.

(x) \emph{Collect}

(\emph{Oratio}.)

(\emph{for the day and week}.)

Raise up, we pray Thee, O Lord, Thy power,
and come,

that we, being meet, may be snatched from the perils of our

sins by thy succour, and be saved by Thy deliverance, who

livest, \&c.

The Lord be with you.

\emph{And with thy spirit}.

Let us bless the Lord.

\emph{Thanks be to God}.

May the souls of the faithful rest in peace.

\emph{Amen}.

———————

\xsubsection{PRIME}

\begin{quote}

\begin{quote}

(\emph{Here are
marked merely the differences from the 4th Sunday in Pentecost. Some of them, as in Matins and
Lauds, are but peculiar to the season}.)

\end{quote}

\end{quote}

\emph{Antiphon}

\emph{Psalm} 54. (1)

\emph{Psalm} 24. (2)

\emph{Psalm} 119. (3)

1-32.

16-32. (4)
In that day the mountains shall drop down new
wine.

Save me, O God, \&c. Glory, \&c. (\emph{as
on Sunday}.)

The earth is the Lord's,
\&c. Glory, \&c.

Blessed are those, \&c. (\emph{as on Sunday}.)

O do well, \&c. (\emph{as
on Sunday}).

\emph{Antiphon}

In that day the mountains shall drop down new
wine, and

the hills shall flow with milk. Hallelujah.

(cc) \emph{Text}

(\emph{Capitulum}.)

\emph{Zech}. viii. 19, 20.

\emph{Minister}.—Love the Truth and Peace,
saith the Lord of

hosts.

\emph{Thanks be to God}.

(dd) \emph{Short Response}\\
(\emph{the same except instead of Thou that sittest, \&c}.) (ee)

O Christ, \&c.

Thou that shall come into the world—

Lord have mercy, \&c.

The Lord's Prayer

The Creed

Sentences

Our Father, \&c.

I believe, \&c.

I have cried, \&c. (\emph{down
to})

And stablish me, \&c.

(\emph{These additional Sentences are insert- ed, whenever there are the additional Sentences \& Psalms as above set down in Lauds and Ves- pers}.)
Deliver me, O Lord, from the evil man.

\emph{And preserve me from the wicked man}.

Deliver me from mine enemies, O God.

\emph{Defend me from them that rise up against
me}. 

O deliver me from the wicked doers.

\emph{And save me from the blood-thirsty men}.

So will I always sing praise unto Thy name.

\emph{That I may daily perform my vows}.

Hear us, O God of our salvation.

\emph{Thou that art the hope of all the ends of
the earth, and of them that remain in the broad sea}.

O God, make speed to save me.

\emph{O Lord, make haste to help me}.

Holy God, holy and strong, holy and eternal.

\emph{Have mercy upon us}.

Bless the Lord, O my soul.

\emph{And all that is within me, bless His holy
Name}.

Bless the Lord, O my soul.

\emph{And forget not all His benefits}.

Who forgiveth all thy sin.

\emph{And healeth all thine infirmities}.

Who saveth thy life from destruction.

\emph{And crowneth thee with mercy and
loving-kindness}.

Who satisfieth thy mouth with good things.

\emph{Making thee young and lusty as an eagle}.

Our help is, \&c.

\emph{The Confessions, Answers, and Pray- er of Absolution.}

I confess, \&c. God Almighty, \&c. I
confess, \&c. God

Almighty, \&c. The Almighty, \&c. vouchsafe, \&c.

\emph{Collect}

O Lord, God Almighty, \&c.

(\emph{And so on to the end, as on the 4th
Sunday in Pentecost, with the substitution of the following}.)

(ii) \emph{Short Lesson}

\emph{Isai}. xxxiii. 2.

O Lord, be gracious unto us. We have waited
for Thee.

Be Thou our arm every morning; our salvation also in the

time of trouble.

———————

\xsubsection{THIRD HOUR}

\emph{The
following are the only variations}.

(kk) \emph{Antiphon} (\emph{be- fore and after the Psalms}.)

Rejoice greatly, daughter of Zion: shout, O
daughter of

Jerusalem.
Hallelujah.

(ll) \emph{Text}

(\emph{Capitulum}.)

\emph{Jer}. xxxiii. 5.

Behold, the days come, saith the Lord, that I
will raise

unto David a righteous Branch, and a King shall reign and

prosper, and shall execute judgment and justice in the earth.

(mm) \emph{Short Response}

\emph{Come thou and save us: O Lord God of
hosts.}

\emph{Come Thou, \&c}.

Show the light of thy countenance and we
shall be whole.

\emph{O Lord God of Hosts}.

Glory be to the Father, \&c.

\emph{Come Thou and save us,
O Lord God of hosts}.

The heathen shall fear Thy name, O Lord.

\emph{And all the kings of
the earth Thy majesty}.

nn) \emph{Collect}

(\emph{for the week}.)
(as in Lauds.)

———————

\xsubsection{SIXTH HOUR}

(\emph{The
variations are})

(pp) \emph{Antiphon} (\emph{be- fore and after the Psalms}.)

Behold, the Lord shall come, and all His
Saints with Him,

and the light of that day shall be great. Hallelujah.

(qq) \emph{Text}

(\emph{Capitulum}.)

\emph{Jer}. xxiii. 6.

In His days Judah shall be saved, and Israel
shall dwell

safely; and this is the name whereby He shall be called,

The Lord our righteousness.

(rr) \emph{Short Response}

\emph{Show us Thy mercy: O Lord}.

\emph{Show
us, \&c}.

And grant us Thy Salvation.

\emph{O
Lord}.

Glory be to the Father, \&c.

\emph{Show us Thy mercy: O Lord}.

Remember me, O Lord, according to the favour
that Thou

bearest unto Thy people.

\emph{O
visit me with Thy Salvation}.

(ss) \emph{Collect for the week}
(\emph{as before}.)

———————

\xsubsection{NINTH HOUR}

(\emph{The
variations are})

(uu) \emph{Antiphon} (\emph{be- fore and after the Psalms}.)

Behold, the great Prophet shall come, and He
shall re-

build Jerusalem. Hallelujah.

(vv) \emph{Text}

(\emph{Capitulum}.)

\emph{Isai}. xiv. 1.

His time is near to come, and his days shall
not be pro-

longed: for the Lord will have mercy on Jacob, and will yet

choose Israel.

(ww) \emph{Short Response}

\emph{The
Lord shall arise upon thee: O Jerusalem}.

\emph{The Lord, \&c}.

And His glory shall be
seen upon thee.

\emph{O
Jerusalem}.

Glory be to the Father,
\&c.

\emph{The
Lord shall arise upon thee: O Jerusalem}. 

Turn us again, O Lord God of Hosts.

\emph{Show the light of thy countenance, and we shall be whole}.

(xx) \emph{Collect for the week}.
(\emph{as before}.)

———————

\xsubsection{VESPERS}

\emph{Antiphon}

\emph{Psalm} 116 (part 1) (l)

\emph{Antiphon}

O God, make speed, \&c.

\emph{O
Lord, make haste, \&c}.

He hath inclined.

I am well pleased, \&c.
Glory, \&c.

He hath inclined His ear unto me.

\emph{Antiphon}

\emph{Psalm} 116 (part 2)

(2)

\emph{Antiphon}

I believed

… and therefore
will I speak, \&c. Glory, \&c.

I believed, and therefore will I speak.

\emph{Antiphon}

\emph{Psalm} 117. (3)

\emph{Antiphon}

O Praise the Lord.

… all ye heathen,
\&c. Glory, \&c.

O Praise the Lord, all ye heathen.

\emph{Antiphon}

\emph{Psalm} 120. (4)

\emph{Antiphon}

I called upon the Lord.

When I was in trouble, \&c. Glory, \&c.

I called upon the Lord, and He heard me.

\emph{Antiphon}

\emph{Psalm} 121. (5)

\emph{Antiphon}

From whence cometh.

I will lift up, \&c. Glory, \&c.

From whence cometh my help.

(ddd) \emph{Text Capitulum}

\emph{Gen}. xlix. 10.

The sceptre shall not depart from Judah, nor
a lawgiver

from between his feet, until Shiloh come: and unto Him shall

the gathering of the people be.

\emph{Hymn}

Creator alme
siderum, \&c.

(fff) \emph{Verse and Response}

Drop down, ye heavens, from above, and let
the skies pour

down Righteousness.

\emph{Let the earth open, and let them bring
forth Salvation}.

(ggg) \emph{Antiphon}

Lift up thine eyes, O Jerusalem.

\emph{The Magnificat}

My soul doth magnify, \&c.

\emph{Antiphon}

Lift up thine eyes, O Jerusalem, and behold
the greatness

of thy king. Behold, the Saviour comes to loose thee from

thy chain.

(\emph{Then follow the same Sentences with Psalm} 51, \emph{instead
of Psalm} 130, \emph{which are inserted in Lauds}.)

(hhh) \emph{The Collect for the week}
(\emph{as before}.)

\emph{}The Lord be with you.

\emph{And with thy spirit}.

\emph{}Let us bless the Lord.

\emph{Thanks be to God}.

\emph{}May the souls of the faithful rest in
peace.

\emph{Amen}.

———————

\xsubsection{COMPLINE}

\begin{quote}

\begin{quote}

\emph{This office is invariable throughout the
year, except in the Anthems to the Virgin at the end}.

\emph{Instead of the Salve Regina as after
Pentecost, in Advent is used the}Alma Redemptoris.

\end{quote}

\end{quote}

\xsection{§ 4. Part of the Service for August 6th.}

\emph{The Feast
of the Transfiguration}

\xsubsection{MATINS}

O Lord, open thou my lips.

\emph{And my mouth shall
show forth thy praise}.

O God, make speed to save me.

\emph{O Lord, make haste
to help me}.

Glory be, \&c.

As it was, \&c. Amen. Hallelujah.

(a) \emph{Invitatory}
Let us worship Christ Most
High, the King of Glory.

Let us worship,
\&c.

\emph{Psalm} 95.
O come, let us sing unto the
Lord: let us heartily rejoice

in the strength of our salvation.

Let us come before his presence with
thanksgiving: and

show ourselves glad in him with psalms.

Let us worship Christ, \&c.

For the Lord is a great God: and a great King
above all

gods.

In his hand are all the corners of the earth:
and the

strength of the hills is his also.

The King of Glory.

The sea is his, and he made it: and his hands
prepared the

dry land.

O come, let us worship and fall down: and
kneel before

the Lord our Maker;

For he is the Lord our God: and we are the
people of his

pasture, and the sheep of his hand.

Let us worship Christ, \&c.

Today if ye will hear his voice, harden not
your hearts:

as in the provocation, and as in the day of temptation in the

wilderness.

When your fathers tempted me: proved me, and
saw my

works.

The King of Glory. 

Forty years long was I grieved with this
generation, and

said: It is a people that do err in their hearts: for they have

not known my ways.

Unto whom I sware in my wrath: that they
should not

enter into my rest.

Let us worship Christ, \&c.

Glory be, \&c.

The King of Glory.

Let us worship Christ
Most High, the King of Glory.

(b) \emph{Hymn}

[Quicunque

Christum quæ-

ritis.]

O ye who seek the Lord,

Lift up your eyes on high,

For there He doth the sign accord

Of His bright majesty.

We see a wondrous sight

That shall outlive all time,

Older than depth and starry height,

Limitless and sublime.

‘Tis He for Israel's fold

And heathen tribes decreed,

The King to Abraham pledged of old

And his eternal seed.

Prophets foretold His birth,

And witness'd when He came;

The Father speaks to all the earth,

To hear, and fear His name.

To Jesus, who displays

To babes His beaming face,

Be, with the Father, endless praise,

And with the Spirit of grace. \emph{Amen}.

——————

NOCTURN
I.

(c) \emph{Antiphon}
Thou madest Him a little lower
than the Angels, to crown

Him with glory and worship. Thou makest Him to have

dominion of the works of Thy hands.

\emph{Psalm} 8. (1)
O Lord our Governor, how
excellent is thy Name in all

the world thou hast set thy glory above the heavens! 

Out of the mouth of very babes and sucklings
hast thou

ordained strength, because of thine enemies: that thou

mightest still the enemy and the avenger.

For I will consider thy heavens, even the
works of thy

fingers: the moon and the stars which thou hast ordained.

What is man, that thou art mindful of him:
and the son

of man, that thou visitest him?

Thou madest him lower than the angels: to
crown him

with glory and worship.

Thou makest him to have dominion of the works
of thy

hands: and thou hast put all things in subjection under his

feet;

All sheep and oxen: yea, and the beasts of
the field;

The fowls of the air, and the fishes of the
sea: and what-

soever walketh through the paths of the seas.

O Lord our Governor: how excellent is thy
Name in all

the world! Glory be, \&c.

(c) \emph{Antiphon}
Thou madest Him a little lower
than the angels, to crown

Him with glory and worship. Thou makest Him to have

dominion of the works of Thy hands.

(d) \emph{Antiphon}

The Lord discovereth the thick bushes: in His
temple

doth every man speak of His honour.

\emph{Psalm} 29. (2)
Bring unto the Lord, O ye
mighty, bring your rams unto

the Lord: ascribe unto the Lord worship and strength.

Give the Lord the honour due unto his Name:
worship the

Lord with holy worship.

It is the Lord that commandeth the waters: it
is the glo-

rious God that maketh the thunder.

It is the Lord that ruleth the seas; the
voice of the Lord

is mighty in operation: the voice of the Lord is a glorious

voice.

The voice of the Lord breaketh the cedar
trees: yea, the

Lord breaketh the cedars of Libanus.

He maketh them also to skip like a calf:
Libanus also, and

Sirion, like a young unicorn.

The voice of the Lord divideth the flames of
fire; the voice

of the Lord shaketh the wilderness: yea, the Lord shaketh

the wilderness of Cades.

The voice of the Lord maketh the hinds to
bring forth

young, and discovereth the thick bushes: in his temple doth

every man speak of his honour.

The Lord sitteth above the water-flood: and
the Lord re-

maineth a King for ever.

The Lord shall give strength unto his people:
the Lord

shall give his people the blessing of peace. Glory be, \&c.
(d) \emph{Antiphon}

The Lord discovereth the thick bushes: in His
temple doth

every man speak of His honour.

(e) \emph{Antiphon}

Thou art fairer than the children of men:
full of grace are

Thy lips.

\emph{Psalm} 45. (3)

My heart is inditing of a good matter: I
speak of the

things which I have made unto the King.

My tongue is the pen: of a ready writer.

Thou art fairer than the children of men:
full of grace are

thy lips, because God hath blessed thee for ever.

Gird thee with thy sword upon thy thigh, O
thou most

Mighty: according to thy worship and renown.

Good luck have thou with thine honour: ride
on, because

of the word of truth, of meekness, and righteousness: and thy

right hand shall teach thee terrible things.

Thy arrows are very sharp, and the people
shall be sub-

dued unto thee: even in the midst among the King's

enemies.

Thy seat, O God, endureth for ever: the
sceptre of thy

kingdom is a right sceptre.

Thou hast loved righteousness, and hated
iniquity: where-

fore God, even thy God, hath anointed thee with the oil of

gladness above thy fellows.

All thy garments smell of myrrh, aloes, and
cassia: out of

the ivory palaces, whereby they have made thee glad.

Kings' daughters were among thy honourable
women:

upon thy right hand did stand the queen in a vesture of gold,

wrought about with divers colours.

Hearken, O daughter, and consider, incline
thine ear: for-

get also thine own people, and thy father's house.

So shall the King have pleasure in thy
beauty: for he is

thy Lord God, and worship thou him.

And the daughter of Tyre shall be there with
a gift: like as

the rich also among the people shall make their supplication

before thee.

The King's daughter is all glorious within:
her clothing

is of wrought gold.

She shall be brought unto the King in raiment
of needle-

work: the virgins that be her fellows shall bear her company,

and shall be brought unto thee.

With joy and gladness shall they be brought:
and shall

enter into the King's palace.

Instead of thy fathers, thou shalt have
children: whom

thou mayest make princes in all lands.

I remember thy name from one generation to
another, 

therefore shall the people give thanks unto thee, world

without end. Glory be, \&c.

(e) \emph{Antiphon}

Thou art fairer than the children of men:
full of grace

are Thy lips.

(1) \emph{Verse and Response}

Thou hast appeared glorious in the sight of the Lord.

\emph{Therefore the Lord hath clothed thee in
comely apparel}.

\emph{The Lord's Prayer}

(\emph{privately})

Our Father, \&c.

And lead us not into temptation.

\emph{But deliver us from
evil}.

\emph{Absolution} 1.

O Lord Jesus Christ, hear the prayers of thy
servants,

and have mercy upon us, who with the Father and the Holy

Spirit livest and reignest, world without end. \emph{Amen}.

\emph{Benediction} 1.

\emph{Reader}.—Sir, be pleased to bless us.

\emph{Minister}.—The Father everlasting,
bless us with an

eternal blessing. \emph{Amen}.

\emph{Lesson} 1.

2 \emph{Pet}. i. 10-14.

The rather, brethren, give diligence to make
your calling

and election sure; for if ye do these things, ye shall never

fall.

For so an entrance shall be ministered unto
you abun-

dantly, into the everlasting kingdom of our Lord and

Saviour Jesus Christ.

Wherefore I will not be negligent to put you
always in

remembrance of these things, though ye know them, and

be established in the present truth.

Yea, I think it meet, as long as I am in this
tabernacle,

to stir you up, by putting you in remembrance.

Knowing that shortly I must put off this my
tabernacle,

even as our Lord Jesus Christ hath showed me.

But thou, O Lord, have mercy upon us.

\emph{Thanks be to God}.

\emph{Response} 1.

\emph{Arise, shine, for thy light is come, and
the glory of the Lord is risen upon thee}.

And Gentiles shall come to thy light, and
kings to the

brightness of thy rising.

\emph{And the glory of the Lord is risen upon
them}.

\emph{Benediction} 2.

\emph{Reader}.—Sir, be pleased to bless us.

\emph{Minister}.—The only-begotten Son of
God, vouchsafe to

bless and help us. \emph{Amen}. \emph{Lesson} 2.

2 \emph{Pet}. i. 15-17.

Moreover, I wilt endeavour, that ye may be
able, after my

decease, to have these things always in remembrance.

For we have not followed cunningly devised
fables, when

we made known unto you the power and coming of our

Lord Jesus Christ, but were eye-witnesses of his majesty.

For he received from God the Father honour
and glory,

when there came such a voice to him from the excellent

glory, This is my beloved Son, in whom I am well pleased.

But
Thou, O Lord, have mercy upon us.

\emph{Thanks be to God}.

\emph{Response} 2.

\emph{In the bright cloud the
Holy Ghost was seen, the Father's}

\emph{voice
was heard: This is My beloved Son, in whom I am}

\emph{well
pleased; hear ye Him}.

A cloud overshadowed them, and the Father's
voice was

heard in thunder.

\emph{This is My beloved Son, in whom I am well
pleased; hear ye Him}.

\emph{Benediction} 3.

\emph{Reader}.—Sir, be pleased to bless us.

\emph{Minister}.—The grace of the Holy
Ghost enlighten our

thoughts and hearts. \emph{Amen}.

\emph{Lesson} 3.

2 \emph{Pet}. i. 18-20.

And this voice which came from heaven we
heard, when

we were with him in the holy mount.

We have also a more sure word of prophecy
whereunto

ye do well that ye take heed, as unto a light that shineth

in a dark place, until the day dawn, and the day-star arise

in your hearts:

Knowing this first, that no prophecy of the
scripture is

of any private interpretation.

For the prophecy came not in old time by the
will of man:

but holy men of God spake \emph{as they were} moved by the Holy

Ghost.

But
Thou, O Lord, have mercy upon us.

\emph{Thanks be to God}.

\emph{Response} 3.

\emph{Behold what manner of
love the Father hath bestowed upon}

\emph{us;
that we should be called the sons of God}.

We know that when He shall appear, we shall
be like

Him, for we shall see Him as He is.

\emph{That we should be
called the sons of God}.

Glory be to the Father, \&c.

\emph{That we should be
called the sons of God}.

NOCTURN
II.

(g) \emph{Antiphon}

Thou art of more honour and might than the
hills of the

robbers: the proud are robbed.

[Illuminans tu mi-

rabiliter à monti-

bus æternis, \&c.]

\emph{Psalm} 76. (4)

In Jewry is God known: his name is great in
Israel.

At Salem is his tabernacle: and his dwelling
in Sion.

There brake he the arrows of the bow: the
shield, the

sword, and the battle.

Thou art of more honour and might: than the
hills of the

robbers.

The proud are robbed, they have slept their
sleep: and

all the men, whose hands were mighty, have found nothing.

At thy rebuke, O God of Jacob: both the
chariot and horse

are fallen.

Thou, even thou art to be feared: and who may
stand in

thy sight when thou art angry?

Thou didst cause thy judgment to be heard
from heaven:

the earth trembled, and was still.

When God arose to judgment: and to help all
the meek

upon earth.

The fierceness of man shall turn to thy
praise: and the

fierceness of them shalt thou refrain.

Promise unto the Lord your God, and keep it,
all ye that

are around about him: bring presents unto him that ought

to be feared.

He shall refrain the spirit of princes: and
is wonderful

among the kings of the earth. Glory be, \&c.

(g) \emph{Antiphon}

Thou art of more honour and might than the
hills of the

robbers: the proud are robbed.

(h) \emph{Antiphon}

\emph{Psalm} 84. (5)

One day in thy courts is better than a
thousand.

O how amiable are thy dwellings: thou Lord of
hosts!

My soul hath a desire and longing to enter
into the

courts of the Lord: my heart and my flesh rejoice in the

living God.

Yea, the sparrow hath found her an house, and
the swal-

low a nest, where she may lay her young, even thy altars,

O Lord of hosts, my King and my God.

Blessed are they that dwell in thy house:
they will be

alway praising thee.

Blessed is the man whose strength is in thee:
in whose

heart are thy ways. 

Who going through the vale of misery use it
for a well:

and the pools are filled with water.

They will go from strength to strength: and
unto the

God of gods appeareth every one of them in Sion.

O Lord God of hosts, hear my prayer: hearken,
O God

of Jacob.

Behold, O God our defender: and look upon the
face of

thine Anointed.

For one day in thy courts: is better than a
thousand.

I had rather be a door-keeper in the house of
my God:

than to dwell in the tents of ungodliness. Glory be, \&c.

(h) \emph{Antiphon}

One day in thy courts is better than a
thousand.

(i) \emph{Antiphon}

Very excellent things are spoken of thee,
thou city of

God.

\emph{Psalm} 87. (6)

Her foundations are upon the holy hills: the
Lord loveth

the gates of Sion more than all the dwellings of Jacob.

Very excellent things are spoken of thee,
thou city of

God.

I will think upon Rahab and Babylon: with
them that

know me.

Behold ye the Philistines also: and they of
Tyre, with the

Morians: lo, there was he born.

And of Sion it shall be reported, that he was
born in her:

and the Most High shall stablish her.

The Lord shall rehearse it: when he writeth
up the people,

that he was born there.

The singers also and trumpeters shall he
rehearse: all

my fresh springs shall be in thee. Glory be, \&c.

(i) \emph{Antiphon}

Very excellent things are spoken of thee,
thou city of God.

(j) \emph{Verse and Response}

Thou crownest Him with glory and worship.

\emph{And makest him to have dominion of the
works of thy hands}.

\emph{The Lord's Prayer}

(\emph{privately})

Our Father, \&c.

And lead us not into temptation,

\emph{But deliver us from
evil}.

\emph{Absolution} 2.

His pity and mercy succour us, who with the
Father and

the Holy Ghost, liveth and reigneth, world without end.

\emph{Amen}.

\emph{Benediction} 4.

\emph{Reader}.—Sir, be pleased to bless us.

\emph{Minister}.—God, the Father Almighty,
be favourable and

gracious unto us. \emph{Amen}. \emph{Lesson} 4.

(\emph{Sermon of St. Leo, Pope}.)

The Lord revealed His glory before certain
chosen wit-

nesses, and brightens that bodily form which He had in

common with others with such splendour, that His face

was like to the sun's blaze, and His raiment all one with

the snow's whiteness. In which Transfiguration this was

the chief design, to remove from the hearts of the disciples

the scandal of the cross, that their faith might be proof

against the lowliness of His voluntary passion, by the

revelation of the excellence of His hidden dignity. And it

was no less a providence, that hereby the hope of the Holy

Church has a sure stay, by knowing how high a change is

in store for the whole body of Christ, so that the honour

first shown in the Head, might be shared in anticipation

by the members.

Thou
then, O Lord, have mercy upon us.

\emph{Thanks be to God}.

\emph{Response} 4.

\emph{They shall be satisfied with the
plenteousness of Thy house:}

\emph{and Thou shalt give
them drink of Thy pleasures, as out of}

\emph{the river}.

For with Thee is the well of life, and in Thy
light shall

we see light.

\emph{And Thou shalt give them drink of Thy
pleasures, as out of the river}.

\emph{Benediction} 5.

\emph{Reader}.—Sir, be pleased to bless us.

\emph{Minister}.—Christ grant to us the
joys of endless life. \emph{Amen}.

\emph{Lesson} 5.

(\emph{Sermon continued}.)

To confirm the Apostles and instruct them in
all know-

ledge, that miracle contained a further lesson. For Moses

and Elias, that is, the Law and the Prophets, appeared con-

versing with the Lord: that by the presence of five men

might be most fully accomplished what is written, ``By two

or three witnesses, every word shall be established.'' What

can be more stable, more firmly fixed, than that word, which

was heralded by the trumpet, both of the Old and the New

Testament, by the instruments of the ancient message in

unison with the Evangelical teaching? For the pages of

each covenant witness either to other, and the brightness

of the present glory does but disclose Him manifest and

clear, whom foregoing wonders had promised under the

veil of Mysteries.

But thou then, O Lord,
have mercy upon us.

\emph{Thanks be to God}.

\emph{Response} 5.

\emph{Master, it is good to be here; and let us
make here} 

\emph{three tabernacles, one
for Thee, one for Moses, and}

\emph{one for Elias}.

For he knew not what he said.

\emph{And let us make here three tabernacles,
one for Thee, one}

\emph{for Moses, and one for
Elias}.

\emph{Benediction} 6.

\emph{Reader}.—Sir, be pleased to bless us.\emph{}

\emph{Minister}.—God kindle the fire of His
love in our hearts.

\emph{Lesson} 6.

(\emph{Sermon continued}.)

Therefore stirred by these disclosures in
outward tokens,

the Apostle Peter, in scorn of things of the world, and in

disgust of what was earthly, was carried away by a sort of

transport into the longing after things eternal; and filled

with joy at all that vision, he desired to dwell with Jesus

there, where he was enjoying the manifestation of His glory.

Wherefore he said, ``Lord, it is good to be here; and, if

Thou wilt, let us build here three tabernacles, one for Thee,

one for Moses, and one for Elias.'' The Lord, however,

answered not to this proposal, intimating, not that it was

presumptuous, but that it was unbecoming to make it; see-

ing the world could not be saved but by Christ's death,

and that in the Lord's pattern faith should find its calling,

not to doubt of the promises of future blessedness, but

withal to understand that amid the trials of this life we must

ask for patience rather than for glory.

But
Thou, O Lord, have mercy upon us.

\emph{Thanks be to God}.

\emph{Response} 6.

\emph{If the ministration of death, written and
engraven in}

\emph{stones, was glorious,
so that the children of Israel}

\emph{could not stedfastly
behold the face of Moses for the}

\emph{glory of his
countenance, which glory was to be done}

\emph{away; much more shall
the ministration of the Spirit,}

\emph{which abideth, be
glorious}.

For Christ was counted worthy of more glory
than Moses,

inasmuch as he who hath builded the house hath more

honour than the house.

\emph{Much more shall the ministration of the
Spirit, which abideth, be glorious}.

Glory be to the Father, \&c.

\emph{Much more shall the ministration of the
Spirit, which abideth, be glorious}.

NOCTURN
III.

(k) \emph{Antiphon}

Tabor and Hermon shall rejoice in Thy Name:
Thou hast

a mighty arm.

\emph{Psalm} 89. (7)

My song shall be alway of the loving-kindness
of the

Lord: with my mouth will I ever be showing thy truth from

one generation to another.

For I have said, Mercy shall be set up for
ever: thy truth

shalt thou establish in the heavens.

I have made a covenant with my chosen: I have
sworn

unto David my servant;

Thy seed will I stablish for ever: and set up
thy throne

from one generation to another.

O Lord, the very heavens shall praise thy
wondrous

works: and thy truth, in the congregation of the saints.

For who is he among the clouds: that shall be
compared

unto the Lord?

And what is he among the gods: that shall be
like unto

the Lord?

God is very greatly to be feared in the
council of the

saints: and to be had in reverence of all them that are

round about him.

O Lord God of hosts, who is like unto thee:
thy truth,

most mighty Lord, is on every side.

Thou rulest the raging of the sea: thou
stillest the waves

thereof when they arise.

Thou hast subdued Egypt, and destroyed it:
thou hast

scattered thine enemies abroad with thy mighty arm.

The heavens are thine, the earth also is
thine: thou hast

laid the foundation of the round world, and all that

therein is.

Thou hast made the north and the south: Tabor
and

Hermon shall rejoice in thy Name.

Thou hast a mighty arm: strong is thy hand,
and high is

thy right hand.

Righteousness and equity are the habitation
of thy seat:

mercy and truth shall go before thy face.

Blessed is the people, O Lord, that can
rejoice in thee:

they shall walk in the light of thy countenance.

Their delight shall be daily in thy Name: and
in thy

righteousness shall they make their boast.

For thou art the glory of their strength: and
in thy

loving kindness thou shalt lift up our horns. 

For the Lord is our defence: the Holy One of
Israel is

our King.

Thou spakest sometime in visions unto thy
saints, and

saidst: I have laid help upon one that is mighty; I have

exalted one chosen out of the people.

I have found David my servant: with my holy
oil have

I anointed him.

My hand shall hold him fast: and my arm shall

strengthen him.

The enemy shall not be able to do him
violence: the son

of wickedness shall not hurt him.

I will smite down his foes before his face:
and plague

them that hate him.

My truth also and my mercy shall be with him:
and in

my Name shall his horn be exalted.

I will set his dominion also in the sea: and
his right hand

in the floods.

He shall call me, Thou art my Father: my God,
and my

strong salvation.

And I will make him my first-born: higher
than the kings

of the earth.

My mercy will I keep for him for evermore:
and my

covenant shall stand fast with him.

His seed also will I make to endure for ever:
and his

throne as the days of heaven.

But if his children forsake my law: and walk
not in my

judgments;

If they break my statutes, and keep not my
command-

ments: I will visit their offences with the rod, and their

sin with scourges.

Nevertheless my loving-kindness will I not
utterly take

from him: nor suffer my truth to fail.

My covenant will I not break, nor alter the
thing that is

gone out of my lips: I have sworn once by my holiness,

that I will not fail David.

His seed shall endure for ever: and his seat
is like as the

sun before me.

He shall stand fast for evermore as the moon:
and as the

faithful witness in heaven.

But thou hast abhorred and forsaken thine
Anointed:

and art displeased at him.

Thou hast broken the covenant of thy servant:
and cast

his crown to the ground.

Thou hast overthrown all his hedges: and
broken down

his strong holds. 

All they that go by spoil him: and he is
become a re-

proach to his neighbours.

Thou hast set up the right hand of his
enemies: and

made all his adversaries to rejoice.

Thou hast taken away the edge of his sword:
and givest

him not victory in the battle.

Thou hast put out his glory: and cast his
throne down to

the ground.

The days of his youth hast thou shortened:
and covered

him with dishonour.

Lord, how long wilt thou hide thyself, for
ever: and shall

thy wrath burn like fire

O remember how short my time is: wherefore
hast thou

made all men for nought?

What man is he that liveth, and shall not see
death: and

shall he deliver his soul from the hand of hell?

Lord, where are thy old loving-kindnesses:
which thou

swarest unto David in thy truth?

Remember, Lord, the rebuke that thy servants
have: and

how I do bear in my bosom the rebukes of many people?

Wherewith thine enemies have blasphemed thee,
and

slandered the footsteps of thine anointed: Praised be the

Lord for evermore.
Amen, and Amen.

(k) \emph{Antiphon}

Tabor and Hermon shall rejoice in Thy Name;
Thou hast

a mighty arm.

(l) \emph{Antiphon}

There is sprung up a light for the righteous,
and joyful

gladness for such as are true-hearted.

\emph{Psalm} 97. (8)

The Lord is King, the earth may be glad
thereof: yea, the

multitudes of the isles may be glad thereof.

Clouds and darkness are round about him:
righteousness

and judgment are the habitation of his seat.

There shall go a fire before him: and burn up
his enemies

on every side.

His lightnings gave shine unto the world: the
earth saw

it, and was afraid.

The hills melted like wax at the presence of
the Lord: at

the presence of the Lord of the whole earth.

The heavens have declared his righteousness:
and all the

people have seen his glory.

Confounded be all they that worship carved
images, and

that delight in vain gods: worship him, all ye gods.

Sion heard of it, and rejoiced: and the
daughters of Judah

were glad, because of thy judgments, O Lord.

For thou, Lord, art higher than all that are
in the earth:

thou art exalted far above all gods. 

O ye that love the Lord, see that ye hate the
thing which

is evil: the Lord preserveth the souls of his saints; he shall

deliver them from the hand of the ungodly.

There is sprung up a light for the righteous:
and joyful

gladness for such as are true-hearted.

Rejoice in the Lord, ye righteous: and give
thanks for a

remembrance of his holiness.

(l) \emph{Antiphon}

There is sprung up a light for the righteous:
and joyful

gladness for such as are true-hearted.

(m) \emph{Antiphon}

\emph{Psalm} 104. (9)

Praise the Lord, O my soul, O my God.

...
thou art become exceeding glorious: thou art clothed

with majesty and honour.

Thou deckest thyself with light as it were
with a garment:

and spreadest out the heavens like a curtain.

Who layeth the beams of his chambers in the waters: and

maketh the clouds his chariot, and walketh upon the wings

of the wind.

He maketh his angels spirits: and his
ministers a flaming

fire.

He laid the foundations of the earth: that it
never should

move at any time.

Thou coveredst it with the deep like as with
a garment: the

waters stand in the hills.

At thy rebuke they flee: at the voice of thy
thunder they

are afraid.

They go up as high as the hills, and down to
the valleys

beneath: even unto the place which thou hast appointed

for them.

Thou hast set them their bounds, which they
shall not

pass: neither turn again to cover the earth.

He sendeth the springs into the rivers: which
run among

the hills.

All beasts of the field drink thereof: and
the wild asses

quench their thirst.

Beside them shall the fowls of the air have
their habitation:

and sing among the branches.

He watereth the hills from above: the earth
is filled with

the fruit of thy works.

He bringeth forth grass for the cattle: and
green herb for

the service of men;

That he may bring food out of the earth, and
wine that

maketh glad the heart of man: and oil to make him a cheer-

ful countenance, and bread to strengthen man's heart.

The trees of the Lord also are full of sap:
even the cedars

of Libanus which he hath planted: 

Wherein the birds make their nests: and the
fir-trees are

a dwelling for the stork.

The high hills are a refuge for the wild
goats: and so are

the stony rocks for the conies.

He appointed the moon for certain seasons:
and the sun

knoweth his going down.

Thou makest darkness that it may be night:
wherein all

the beasts of the forest do move.

The lions roaring after their prey: do seek
their meat

from God.

The sun ariseth, and they get them away
together: and

lay them down in their dens.

Man goeth forth to his work, and to his
labour: until the

evening.

O Lord, how manifold are thy works; in wisdom
hast

thou made them all; the earth is full of thy riches.

So is the great and wide sea also: wherein
are things

creeping innumerable, both small and great beasts.

There go the ships, and there is that
Leviathan: whom

thou hast made to take his pastime therein.

These wait all upon thee: that thou mayest
give them

meat in due season.

When thou givest it them they gather it: and
when thou

openest thy hand they are filled with good.

When thou hidest thy face they are troubled:
when thou

takest away their breath they die, and are turned again to

their dust.

When thou lettest thy breath go forth they
shall be made:

and thou shalt renew the face of the earth.

The glorious Majesty of the Lord shall endure
for ever:

the Lord shall rejoice in his works.

The earth shall tremble at the look of him:
if he do but

touch the hills, they shall smoke.

I will sing unto the Lord as long as I live:
I will praise

my God while I have my being.

And so shall my words please him: my joy
shall be in

the Lord.

As for sinners, they shall be consumed out of
the earth:

and the ungodly shall come to an end: praise thou the

Lord, O my soul, praise the Lord.

\emph{Antiphon}

Praise the Lord, O my soul: O Lord my God.

(n) \emph{Verse and Response}

His honour is great in thy salvation.

\emph{Glory and great worship
shalt Thou lay upon him}. \emph{The Lord's Prayer}

(\emph{privately})

Our Father, \&c.

And lead us not into temptation.

\emph{But
deliver us from evil}.

\emph{Absolution} 3.

The Almighty and merciful Lord absolve us
from the

bonds of our sins. \emph{Amen}.

\emph{Benediction} 7.

\emph{Reader}.—Sir, be pleased to bless us.

\emph{Minister}.—The reading of the Gospel
be to us salvation

and defence. \emph{Amen}.

\emph{Lesson} 7.

\emph{Matt}. xvii. 1-5.

At that time Jesus taketh Peter, James, and
John his

brother, and bringeth them up into an high mountain apart,

And was transfigured before them: and his
face did shine

as the sun, and his raiment was white as the light.

And behold, there appeared unto them Moses
and Elias

talking with him.

Then answered Peter, and said unto Jesus,
Lord, it is

good for us to be here: if thou wilt, let us make here three

tabernacles; one for thee, one for Moses, and one for Elias.

While he yet spake, behold a bright cloud
overshadowed

them: and behold a voice out of the cloud, which said, This

is my beloved Son, in whom I am well pleased: hear ye him.

\emph{Homily of St. John Chrysostom}

Whereas the Lord spoke much concerning
dangers, His

own death, and the death and slaughter of His disciples,

and enjoined on them many severe and difficult things, and

all this, too, in the present life, and soon to come, while the

gain was but in hope and expectation, (for instance, that they

should save their life, if they lost it, that He will come in
the

glory of his Father and dispense rewards,) therefore, in order

to certify them even by sight, and to show them what that

glory was, with which he was to come, He shows it them as

far as they could understand it in this present life, and un-

veils it, lest they should grieve at the thought of their own

or their Lord's death, and chiefly Peter.

But thou, O Lord, have mercy on us.

\emph{Thanks be to God}.

\emph{Response} 7.

\emph{God hath called us with
a holy calling, according to His}

\emph{grace,
which is now made manifest by the glorious}

\emph{appearing
of our Saviour Jesus Christ}.

Who hath abolished death, and hath brought
life and im-

mortality to light.

\emph{By the glorious appearing of our Savior
Jesus Christ}. 

\emph{Benediction} 8.

\emph{Reader}.—Sir, be pleased to bless us.

\emph{Minister}.—The help of God abide with
us for ever. \emph{Amen}.

\emph{Lesson} 8.

(\emph{Homily continued}.)

And see what He does, after discoursing of
his kingdom

and of hell. For in that He said, ``He who finds his life,
shall

lose it, and whoever will lose it for my sake, the same shall

find it;'' and in that He said, ``He shall render to
every one

according to his works,'' He has pointed out both His king-

dom and hell. When then He had discoursed concerning

both, of His kingdom He granted the sight to human eyes,

but not of hell; since, needful as that might have been for the

uninstructed and unready, yet upright and clearsighted men,

as the Apostles, needed but to be confirmed by the better

part. This part indeed it was far more fitting He should

mention, yet He did not altogether pass over the other,

placing at times the terribleness of hell as if before the eyes,

as in his description of Lazarus, and of him who demanded

back the hundred pence.

But thou, O Lord, have mercy upon us.

\emph{Thanks be to God}.

\emph{Response}.

\emph{God, who commanded the
light to shine out of darkness,}

\emph{hath
shined in our hearts: to give the light of the know-}

\emph{ledge
of the glory of God in the face of Jesus Christ}.

Unto the godly there ariseth up light in the
darkness; he

is merciful, loving, and righteous.

\emph{To give the light of the knowledge of the
glory of God in the face of Jesus Christ}.

Glory be to the Father, \&c.

\emph{To give the light of the knowledge of the
glory of God in the face of Jesus Christ}.

\emph{Benediction} 9.

\emph{Reader}.—Sir, be pleased to bless us.

\emph{Minister}.—The King of Angels bring
us through to the

society of the inhabitants of heaven. \emph{Amen}.

\emph{Lesson} 9.

(\emph{Homily continued}.)

Think upon the greatness of mind in St.
Matthew, who has

not concealed the names of those who were preferred over

the rest: which St. John also shows often, when he notes

down the special praises of Peter so accurately and carefully.

In this fellow ship of Apostles there was no place for envy

or for vain-glory. Therefore He took apart with Him the

chief ones of the Apostles. Why took He those only? For

this reason, because they were superior to the rest. But why

did He not do so at once, but after six days? lest His other

disciples, or others generally should be troubled; for which

reason neither did He name those whom He was alone to

take with Him. 

Thou, then, O Lord, have
mercy upon us.

\emph{Thanks
be to God}.

\emph{Te Deum Vide} p. 47.

We praise thee, O God: we acknowledge thee to
be the

Lord.

All the earth doth worship thee: the Father
everlasting.

To thee all Angels cry aloud: the Heavens,
and all the

Powers therein, \&c.

———————

\xsubsection{LAUDS}

O God,
make speed to save me.

\emph{O Lord, make haste to help me}.

Glory be, \&c. As it
was, \&c. Amen. Hallelujah.

(o) \emph{Antiphon}

Jesus taketh Peter, James, and John his
brother, and

bringeth them up into an high mountain apart, and was

transfigured before them.

\emph{Psalm} 93. (1)

\emph{Vide} p. 38.

The Lord is King, and hath put on glorious
apparel: the

Lord hath put on his apparel, and girded himself with

strength.

He hath made the round world so sure: that it
cannot be

moved.

Ever since the world began, hath thy seat
been prepared:

thou art from everlasting, \&c.

Glory be, \&c.

(o) \emph{Antiphon}

Jesus taketh Peter, James, and John his
brother, and

bringeth them up into an high mountain apart, and was

transfigured before them.

(p) \emph{Antiphon}

His face did shine as the sun, and His
raiment was white

as snow. Hallelujah.

\emph{Psalm} 100, (2)

\emph{Vide} p. 48.

O be joyful, \&c.

Glory be, \&c.

(p) \emph{Antiphon}

His face did shine as the sun, and His
raiment was white

as snow. Hallelujah.

(q) \emph{Antiphon}

And behold there appeared unto them Moses and
Elias,

speaking with Jesus.

\emph{Psalm} 63 \emph{and} 67

\emph{Vide} p 48. (3)

O God, thou art my God, \&c.

God be merciful unto us, \&c.

Glory be, \&c.

(q) \emph{Antiphon}

And behold there appeared unto them Moses and
Elias,

speaking with Jesus. 

(r) \emph{Antiphon}

And Peter answered and said to Jesus, Lord,
it is good

to be here.

\emph{Song of the Three}

\emph{Children, vide}\\
p. 50. (4)

O all ye works of the Lord, bless ye the
Lord: praise him,

and magnify him for ever.

O ye Angels of the Lord, bless ye the Lord:
praise him,

and magnify him for ever.

O ye Heavens, bless ye the Lord: praise him,
and mag-

nify him for ever, \&c.

(r) \emph{Antiphon}

And Peter answered and said to Jesus, Lord,
it is good

to be here.

(s) \emph{Antiphon}

While he yet spake, behold, a bright cloud
overshadowed

them.

\emph{Psalm} 148, 149

\emph{and}150. \emph{Vide}\\
p. 52. (5)

O praise ye the Lord of heaven, \&c.

O sing unto the Lord, \&c.

O praise God, \&c.

(s) \emph{Antiphon}

While he yet spake, behold, a bright cloud
overshadowed

them.

(t) \emph{Text Phil}. iii. 20, 21.

\emph{Minister}.—We look for the Saviour,
our Lord Jesus Christ,

who shall change our vile body, that it may be like unto

His glorious body.

\emph{Thanks
be to God}.

(u) \emph{Hymn}

[Lux alma Jesu.]

Light of the anxious
heart,

Jesu,
Thy suppliants cheer;

Bid thou the gloom of
guilt depart,

And
shed Thy sweetness here.

Happy the man, whose
breast

Thou
mak'st Thy residence;

From God's right hand a
radiant guest,

Unseen
by fleshly sense.

Brightness of God above!

Unfathomable grace!

Vouchsafe a present fount
of love,

To
cleanse Thy chosen place.

To Thee, whom children
see,

The
Father ever blest,

The Holy Spirit, One and
Three,

Be
endless praise addrest. \emph{Amen}.

(v) \emph{Verse and Response}

A crown of pure gold is on His forehead.

\emph{With the sign of holiness, glory, and
honour}.

(w) \emph{Ant}.

And lo, a voice from the cloud, saying, This
is My beloved

Son, in whom I am well pleased. Hear ye Him. Hallelujah. \emph{Benedictus}

Blessed be the Lord God of Israel, \&c.

Glory be, \&c.

(w) \emph{Antiphon}

And lo, a voice from the clouds, saying, This
is My Be-

loved Son, in whom I am well pleased, Hear ye Him—

Hallelujah.

The Lord be with you.

\emph{And
with thy spirit}.

(x) \emph{Collect}

Let us
pray.

O God, who in the glorious transfiguration of
Thine Only

begotten hast sealed the treasure of the faith by the witness

of the ancient fathers, and by the voice coming down in a

cloud of light hast wonderfully shadowed forth the perfect

adoption of Thy sons, mercifully grant, that we may be

made fellow heirs and partners in the glory of our King,

through the same our Lord. \emph{Amen}.

\&c.
\&c. \&c.

(\emph{So on to the end of Lauds, in the Service for Sunday}.)

\xsection{§ 5. Part of the Service for August 10th.}

The Feast of St. Laurence, Deacon and Martyr

[In order to understand parts of the following
Service, it may be necessary for the reader to have some knowledge of
St. Laurence's history; which may fitly be conveyed in the following
translation from St. Ambrose's Offices, as found in the British
Magazine, for January, 1834.

``We must not omit mention of the blessed
Laurence, who, on the sight of Sextus, his Bishop, going to martyrdom,
began to weep, not so much at his passion, as his own orphanhood. So he
called out to him, 'Whither goest thou, O my father, without thy son?
Whither can a Priest be hurrying without his Deacon? Never as yet didst
thou offer sacrifice without an attendant. How have I displeased thee?
Hast thou found me a degenerate son? Peter let Stephen suffer before
him. Thou, too, O my Father, show thine own graces in my person, offer
up to God him whom thou hast begotten, nor seize the crown of martyrdom
without a noble company to answer your good thoughts concerning them.'

``The Prelate answered, 'Nay, son, I leave thee
not, neither forsake thee; a fiercer combat is in store for thee. We, as
the old, are allotted the lighter skirmish, but youth must bear off a
more glorious triumph over tyranny. Thou wilt soon be called upon; cease
thy tears; in three days thou shalt follow me. Ill would it seem for me
who hold the third rank in the sacred ministry, to press into the first.
I leave to thee the legacy of my own constancy.'''

In consequence, three days after, Laurence was
arrested, and, after other tortures, broiled to death on a gridiron. St.
Ambrose adds, ``when he was stretched upon the scorching gridiron, he
did but say, The meat is done; turn it over, and eat it.''' This
happened A.D. 258. Other particulars of his Martyrdom
will be found in the Service itself. It may be added, by way of
explaining an allusion in the above account, that ``Priest,'' in the
language of antiquity, means one who has the power of consecrating the
Eucharist, \emph{i.e}. Bishop and Presbyter, [\emph{Leitourgos}], or as
our Service seems sometimes to express it, \emph{Minister}; and that
this Deacon was the usual attendant on the Minister in the celebration.]

\xsubsection{FIRST VESPERS}

O God, make speed to save me.

\emph{O Lord, make haste
to help me}.

Glory be to the Father, \&c.

As it was in the beginning, \&c. Amen.
Hallelujah

(yy) \emph{Antiphon}
Laurence entered on his
Martyrdom, and confessed the

name of the Lord Jesus Christ.

\emph{Ps}. 110. p. 76. (1)
The Lord said, \&c.

(yy) \emph{Antiphon}
Laurence entered on his
Martyrdom, and confessed the

name of the Lord Jesus Christ.

(zz) \emph{Antiphon}

Laurence hath wrought a good work, who by the
sign of

the Cross gave sight to the blind.

\emph{Ps}. 111. p. 76. (2)
I will give thanks, \&c.

(zz) \emph{Antiphon}
Laurence hath wrought a good
work, who by the sign of

the Cross gave sight to the blind.

(aaa) \emph{Antiphon}

My soul hangeth upon Thee, because my flesh
is burned

in the fire for Thee, O my God.

\emph{Ps}. 112. p. 77. (3)
Blessed is the man, \&c.

(aaa) \emph{Antiphon}
My soul hangeth upon Thee,
because my flesh is burned

in the fire for Thee, O my God.

(bbb) \emph{Antiphon}

The Lord hath sent His Angel, and hath
delivered me

from the midst of the fire, and I am not tormented.

\emph{Ps}. 113. p. 77. (4)
Praise the Lord, ye servants,
\&c.

(bbb) \emph{Antiphon}
The Lord hath sent his Angel,
and hath delivered me

from the midst of the fire, and I am not tormented.

(ccc) \emph{Antiphon}

Blessed Laurence prayed, saying, I give Thee
thanks, O

Lord, because I have been found worthy to enter Thy gates.

\emph{Psalm} 117. (5)
O praise the Lord, all ye
heathen: praise him, all ye

nations.

For his merciful kindness is ever more and
more towards

us: and the truth of the Lord endureth for ever. Praise

the Lord.

(ccc) \emph{Antiphon}
Blessed Laurence prayed,
saying, I give Thee thanks, O

Lord, because I have been found worthy to enter Thy gates. 

(ddd) \emph{Text}\\
2 \emph{Cor}. ix. 6.

\emph{Minister}.—Brethren, he which soweth
sparingly shall

reap sparingly; and he which soweth bountifully shall reap

bountifully.

\emph{Thanks be to God}.

(eee) \emph{Hymn}

[Invicte martyr

unicum. \emph{It is re- markable that this Hymn, which is the only one of those here trans- ated which sa- vours of Roman- ism, is the only one, except one other, which is not known to be ancient. The rest are either Ambro- sian or Gregorian, except one, which is by Prudentius}.

(Quicunque Chri-

stum.) \emph{Of the Hymn} Lux Alma

Jesu, \emph{the Trans- lator cannot dis- cover the Author}.]

Martyr of Christ, thy fight is won!

Following the Father's only Son,

O'er thy fall'n foes thou triumphest,

In heavenly courts a risen guest.

Use thou for us thy gift of prayer

To cleanse thy brethren's sin,

To sweeten earth's infectious air,

And gain us peace within.

For ever broken is the chain

That bound thy
body's hallowed fane;

As God hath
given thee, break the tie

Which links our
hearts to vanity.

To God the
Father, God the Son,

And
God the Paraclete,

Be praise, while
circling ages run

Beneath
the Eternal's feet. \emph{Amen}.

(fff) \emph{Verse and Response}

Thou hast crowned him with glory and worship.

\emph{And makest Him to
have dominion of the works of Thy hands}.

(ggg) \emph{Antiphon}

On the hot bars I denied Thee not, my God:
and, when

brought to the fire, I confessed Thee, O Christ. Thou hast

proved my heart, and visited me in the night: thou hast

tried me by fire, and hast found no wickedness in me.

\emph{Magnificat}
My
soul doth magnify the Lord, \&c.

\emph{Antiphon}
On the hot bars I denied thee
not, \&c.

The Lord be with you.

\emph{And with thy spirit}.

Let us pray.

\emph{Collect}
Almighty God, who gavest to
blessed Laurence to over-

come the fire of his torture, grant to us, we beseech Thee,

to extinguish the flames of our vices, through the Lord.

\emph{Amen}.

(\emph{The Service ends as above in Vespers},
p. 80.)

\xsubsection{MATIN SERVICE}

O Lord, open Thou my lips:

\emph{And my mouth shall show
forth Thy praise}.

O God, make speed to save me.

\emph{O Lord, make haste to
help me}.

Glory be to the Father, \&c.

As it was in the
beginning, \&c. Amen. Hallelujah.

(a) \emph{Invitatory}
Blessed Laurence, the Martyr
of Christ, is crowned, and

triumphs in heaven: come, let us worship.

Blessed Laurence, the Martyr of Christ,
\&c.

\emph{Psalm} 95
O come, let us sing unto the
Lord: let us heartily rejoice

in the strength of our salvation.

Let us come before his presence with
thanksgiving: and

show ourselves glad in him with psalms.

Blessed Laurence, the
Martyr of Christ, is crowned,

and
triumphs in heaven: come, let us worship.

For the Lord is a great God: and a great King
above all

gods.

In his hand are all the corners of the earth:
and the

strength of the hills is his also.

Come,
let us worship.

The sea is his, and he made it: and his hands
prepared

the dry land.

O come, let us worship and fall down: and
kneel before

the Lord our Maker.

For he is the Lord our God: and we are the
people of

his pasture, and the sheep of his hand.

Blessed Laurence, the
Martyr of Christ, is crowned,

and
triumphs in heaven: come, let us worship.

Today if ye will hear his voice, harden not
your hearts:

as in the provocation, and as in the day of temptation in

the wilderness;

When your fathers tempted me: proved me, and
saw my

works. 

Come,
let us worship.

Forty years long was I grieved with this
generation, and

said: It is a people that do err in their hearts, for they have

not known my ways;

Unto whom I sware in my wrath: that they
should not

enter into my rest.

Blessed Laurence, the
Martyr of Christ, is crowned,

and
triumphs in heaven: come, let us worship.

Glory be to the Father, \&c.

As it was in the beginning, \&c. \emph{Amen}.

Come, let us worship.

Blessed Laurence, the Martyr of Christ, is
crowned, and

triumphs in heaven: come,
let us worship.

(b) \emph{Hymn}

[Deus tuorum

militum.]

O God, of Thy soldiers

The Portion and Crown,

Spare sinners, who hymn

The praise of the Blest;

Earth's bitter joys,

Its lures and its frown,

He weighed them and scorned,

And so is at rest.

Thy Martyr he ran

All valiantly o'er

A highway of blood

For the prize Thou hast given.

We kneel at Thy feet,

And meekly implore,

Our pardon may wait

On his triumph in heaven.

Honour and praise

To the Father and Son

And the Spirit be done

Now and always. \emph{Amen}.

NOCTURN
I.

(c) \emph{Antiphon}
Whither speedest thou without
thy son, O my father?

Whither, holy Priest, hurriest thou without attendant?

\emph{Psalm} 1. (1)
Blessed is the man, that hath
not walked in the counsel of

the ungodly, nor stood in the way of sinners: and hath not

sat in the seat of the scornful.

But his delight is in the law of the Lord:
and in his law

will he exercise himself day and night.

And he shall be like a tree planted by the
water-side:

that will bring forth his fruit in due season.

His leaf also shall not wither: and look,
whatsoever he

doeth, it shall prosper.

As for the ungodly, it is not so with them:
but they are

like the chaff, which the wind scattereth away from the

face of the earth.

Therefore the ungodly shall not be able to
stand in the

judgment: neither the sinners in the congregation of the

righteous.

But the Lord knoweth the way of the
righteous: and the

way of the ungodly shall perish. Glory be, \&c.

(c) \emph{Antiphon}
Whither speedest thou, without
thy son, O my father?

Whither, holy Priest, hurriest thou without attendant?

(d) \emph{Antiphon}

Forsake me not, O holy father, for I have
just laid out

those treasures which thou gavest me in trust.

\emph{Psalm} 2. (2)
Why do the heathen so
furiously rage together: and why

do the people imagine a vain thing?

The kings of the earth stand up, and the
rulers take

counsel together: against the Lord, and against his

Anointed.

Let us break their bonds asunder: and cast
away their

cords from us.

He that dwelleth in heaven shall laugh them
to scorn:

the Lord shall have them in derision.

Then shall he speak unto them in his wrath:
and vex

them in his sore displeasure.

Yet have I set my King: upon my holy hill of
Sion.

I will preach the law, whereof the Lord hath
said unto

me: Thou art my Son, this day have I begotten thee.

Desire of me, and I shall give thee the
heathen for thine

inheritance: and the utmost parts of the earth for thy

possession. 

Thou shalt bruise them with a rod of iron:
and break

them in pieces like a potter's vessel.

Be wise now therefore, O ye kings; be
learned, ye that

are judges of the earth.

Serve the Lord in fear: and rejoice unto him
with re-

verence.

Kiss the Son, lest he be angry, and so ye
perish from the

right way: if his wrath be kindled, (yea, but a little)

blessed are all they that put their trust in him. Glory be,

\&c.

(d) \emph{Antiphon}
Forsake me not, O holy father,
for I have now laid out

those treasures which thou gavest me in trust.

(e) \emph{Antiphon}

I desert thee not, O my son, neither do I
forsake thee;

but a higher conflict for the faith of Christ is in store for

thee.

\emph{Psalm} 3. (3)
Lord, how are they increased
that trouble me: many are

they that rise against me.

Many one there be that say of my soul: There
is no help

for him in his God.

But thou, O Lord, art my defender: thou art
my worship,

and the lifter up of my head.

I did call upon the Lord with my voice; and
he heard me

out of his holy hill.

I laid me down and slept, and rose up again:
for the Lord

sustained me.

I will not be afraid for ten thousands of the
people: that

have set themselves against me round about.

Up, Lord, and help me, O my God: for thou
smitest all

mine enemies upon the cheek-bone; thou hast broken the

teeth of the ungodly.

Salvation belongeth unto the Lord: and thy
blessing is

upon thy people. Glory be, \&c.

(e) \emph{Antiphon}
I desert thee not, O my son,
neither do I forsake thee;

but a higher conflict for the faith of Christ is in store for

thee.

(f) \emph{Verse and Response}

Thou hast crowned him with glory and worship.

\emph{And makest him to have dominion of the works of Thy hands}.

\emph{The Lord's Prayer}

(\emph{privately}.)

Our Father, \&c.

And lead us not in to temptation,

\emph{But deliver us from evil}.

\emph{Absolution} 1.
O Lord Jesus Christ, hearken unto the prayers of
thy 

servants, and have mercy upon us, who livest and reignest

with the Father and Holy Ghost, world without end. \emph{Amen}.

\emph{Benediction} 1.

\emph{Reader}.—Sir, be pleased to bless us.

\emph{Minister}.—The Father everlasting
bless us with a per-

petual blessing. \emph{Amen}.

\emph{Lesson} 1.

\emph{Ecclus}. li. 1-5.
I will thank thee, O Lord and
King, and praise thee, O

God my Saviour, I do give praise unto Thy name;

For thou art my defender and helper, and hast
preserved

my body from destruction, and from the snare of the slan-

derous tongue, and from lips that forge lies, and hast been

mine helper against mine adversaries;

And hast delivered me, according to the
multitude of Thy

mercies and greatness of Thy name, from the teeth of them

that were ready to devour me, and out of the hands of such

as sought after my life, and from the manifold afflictions

which I had;

From the choking of fire on every side, and
from the

midst of the fire which I kindled not.

From the depth of the belly of hell, from an
unclean

tongue, and from lying words, from an unjust king, and

from an unrighteous tongue.

But thou, O Lord, have mercy upon us.

\emph{Thanks be to God}.

\emph{Response} 1.

\emph{Laurence, the Deacon,
wrought a good work, who en-}\\
\emph{lightened
the blind by the sign of the Cross: and gave}\\
\emph{the
treasures of the Church to the poor}.

He hath dispersed, he hath given to the poor,
and his

righteousness remaineth for ever.

\emph{And he gave the treasures of the Church to
the poor}.

\emph{Benediction} 2.

\emph{Reader}—Sir, be pleased to bless us.\emph{}\\
\emph{Minister}.—The only begotten Son
of God vouchsafe to

bless us and help us. \emph{Amen}.

\emph{Lesson} 2.

\emph{Ecclus}. li. 6-9.
My soul shall praise the Lord
even unto death, my life

was near to the hell beneath.

They compassed me on every side, and there
was no man

to help me: I looked for the succour of men, but there was

none.

Then thought I upon thy mercy, O Lord, and
upon Thy

acts of old, how Thou deliverest such as wait for Thee,

and savest them out of the hands of the enemies. 

Then I lifted up my supplication from the
earth, and

prayed for deliverance from death.

But Thou, O Lord, have mercy upon us.

\emph{Response} 2.

\emph{Thanks be to God}.

\emph{My child, be not afraid, for I am with
thee, saith the}\\
\emph{Lord. When thou
passest through the fire, thou shalt}\\
\emph{not be burned,
neither shall the smell of fire pass upon}\\
\emph{thee}.

I will deliver thee from the bond of the
wicked, and

rescue thee from the hand of the mighty.

\emph{When thou passest through the fire, thou
shalt not be burned, neither shall the smell of fire pass upon thee}.

\emph{Benediction} 3.

\emph{Reader}.—Sir, be pleased to bless us.

\emph{Minister}.—The grace of the Holy
Spirit enlighten our

thoughts and hearts. \emph{Amen}.

\emph{Lesson} 3.

\emph{Ecclus}. li. 10-12.
I called upon the Lord, the
Father of my Lord, that He

would not leave me in the days of my trouble, and in the

time of the proud, when there was no help.

I will praise Thy name continually, and will
sing praise

with thanksgiving; and so my prayer was heard.

For thou savedst me from destruction, and
deliveredst me

from the evil time: therefore will I give thanks, and praise

Thee, and bless Thy name, O Lord.

But Thou, O Lord, have mercy upon us.

\emph{Response} 3.

\emph{Thanks be to God}.

\emph{They bound down his
limbs upon the bars; but while they ministered live coals, the deacon of Christ laughed them to
scorn. Blessed Laurence, Martyr of Christ, inter- cede
for us}.

My night has no darkness, but all things grow
clear in the

light.

\emph{Blessed Laurence, Martyr of Christ,
intercede for us}.

Glory be to the Father, \&c.

As it was in the beginning, \&c.

\emph{Blessed Laurence, Martyr of Christ,
intercede for us}.

NOCTURN
II.

(g) \emph{Antiphon}
Blessed Laurence prayed
saying, Lord Jesus Christ, God

from God, have mercy on Thy servant.

\emph{Psalm} 4. (4)
Hear me when I call, O God of
my righteousness: thou

hast set me at liberty when I was in trouble: have mercy

upon me, and hearken unto my prayer.

O ye sons of men, how long will ye blaspheme
mine

honour: and have such pleasure in vanity, and seek after

leasing?

Know this also, that the Lord hath chosen to
himself the

man that is godly: when I call upon the Lord, he will

hear me.

Stand in awe, and sin not: commune with your
own heart

and in your chamber, and be still.

Offer the sacrifice of righteousness: and put
your trust in

the Lord.

There be many that say: Who will show us any
good?

Lord, lift thou up: the light of thy
countenance upon us.

Thou hast put gladness in my heart: since the
time that

their corn, and wine, and oil increased.

I will lay me down in peace, and take my
rest: for it is

thou, Lord, only, that makest me dwell in safety. Glory

be, \&c.

(g) \emph{Antiphon}
Blessed Laurence prayed
saying, Lord Jesus Christ, God

from God, have mercy on thy servant.

(h) \emph{Antiphon}

Romanus said to blessed Laurence, I see
before thee a

young man of fair countenance, hasten to baptize me.

\emph{Psalm} 5. (5)
Ponder my words, O Lord:
consider my meditation.

O hearken thou unto the voice of my calling,
my King

and my God: for unto thee will I make my prayer.

My voice shalt thou hear betimes, O Lord:
early in the

morning will I direct my prayer unto thee, and will look

up.

For thou art the God that hast no pleasure in
wickedness:

neither shall any evil dwell with thee.

Such as be foolish shall not stand in thy
sight: for thou

hatest all them that work vanity.

Thou shalt destroy them that speak leasing:
the Lord will

abhor both the blood-thirsty and deceitful man.

But as for me, I will come into thine house,
even upon the

multitude of thy mercy: and in thy fear will I worship

toward thy holy temple. 

Lead me, O Lord, in thy righteousness,
because of mine

enemies: make thy way plain before my face.

For there is no faithfulness in his mouth:
their inward

parts are very wickedness.

Their throat is an open sepulchre: they
flatter with their

tongue.

Destroy thou them, O God: let them perish
through their

own imaginations: cast them out in the multitude of their

ungodliness: for they have rebelled against thee.

And let all them that put their trust in thee
rejoice: they

shall ever be giving of thanks, because thou defendest

them: they that love thy Name shall be joyful in thee.

For thou, Lord, wilt give thy blessing unto
the righteous:

and with thy favourable kindness wilt thou defend him as

with a shield. Glory be, \&c.

(h) \emph{Antiphon}
Romanus said to blessed
Laurence, I see before thee a

young man of fair countenance, hasten to baptize me.

(i) \emph{Antiphon}

Blessed Laurence answered, My night has no
darkness,

but all things grow clear in the light.

\emph{Psalm} 8. (6)
O Lord, our Governor, how
excellent is thy name in all

the world: thou that hast set thy glory above the heavens!

Out of the mouth of very babes and sucklings
hast thou

ordained strength, because of thine enemies: that thou

mightest still the enemy and the avenger.

For I will consider thy heavens, even the
works of thy

fingers: the moon and the stars which thou hast ordained.

What is man, that thou art mindful of him:
and the Son

of man that thou visitest him?

Thou madest him lower than the angels: to
crown him

with, glory and worship.

Thou makest him to have dominion of the works
of thy

hands: and thou hast put all things in subjection under

his feet;

All sheep and oxen: yea, and the beasts of
the field;

The fowls of the air, and the fishes of the
sea: and what-

soever walketh through the paths of the seas.

O Lord, our Governor: how excellent is thy
name in all

the world! Glory be, \&c.

(i) \emph{Antiphon}
Blessed Laurence answered, My
night has no darkness,

but all things grow clear in the light.

(j) \emph{Verse and Response}

Lord, thou hast set upon his head.

\emph{A crown of precious
stones}. \emph{The Lord's Prayer},

(\emph{privately}.)
Our Father, \&c.

And lead us not into
temptation,

\emph{But
deliver us from evil}.

\emph{Absolution} 2.
His pity and mercy help us,
who with the Father and the

Holy Spirit liveth and reigneth, world without end. \emph{Amen}.

\emph{Benediction} 4.

\emph{Reader}.—Sir, be pleased to bless us.

\emph{Minister}.—God the Father Almighty be
favourable and

gracious unto us. \emph{Amen}.

\emph{Lesson} 4.

(\emph{Sermon of St. Leo, Pope}.)
While the fury of the heathen
powers raged against the

most chosen members of Christ, those chiefly who were in

the order of the Ministry, Laurence the Deacon, having not

only the dispensation of the Sacraments, but also of the

Church's store, incited the impious persecutor, who promised

himself a double price in one man, the gain of the sacred

treasure and the ruin of him who surrendered it. Instigated

therefore by this twofold flame, avarice and hatred of the

truth, to rob him of his treasure and of Christ, he demands

of the spotless Sacristan, those stores of the Church of

which he was greedy. To whom he, most holy Deacon, by

way of showing where they really were laid up, presented

vast multitudes of Christian poor, on whose food and

clothing he had expended wealth which could not be taken

away, but was irrevocably his from the sanctity of his

using.

But thou, O Lord, have mercy upon us.

\emph{Thanks be to God}.

\emph{Response} 4.

\emph{Whither speedest thou
without thy son, O my father? Whither
hurriest thou, holy priest, without thy Dea- con?
Thou never yet hast offered sacrifice without attendant}.

What hath in me displeased thy fatherly love?
in what

have I come short? make trial of his fitness to whom

thou hast entrusted the dispensation of the Lord's blood.

\emph{Thou never yet hast offered sacrifice
without attendant}.

\emph{Reader}.—Sir, be pleased to bless us.

\emph{Benediction} 5.
\emph{Minister}.—Christ
grant us the joys of eternal life. \emph{Amen}.

\emph{Lesson} 5.

(\emph{Sermon, continued}.)
Therefore the disappointed
robber roars aloud, and kind-

ling into hatred of that religion, which had introduced such

an application of worldly goods, he attempts the plunder of

another treasure-house, not of gold or silver, to rob it of that

store which was of a more holy costliness. He bids Lau-

rence renounce Christ, and prepares against the stubborn 

courage of that Deacon's heart, dreadful tortures; and when

the first prove fruitless, he proceeds to fiercer. He tears and

shreds his limbs with continued scourging, next he gives

orders to broil them over the fire, so that, being stretched

upon the red hot bars, first on one side, then on the other,

the torment might be the greater, and the punishment more

protracted.

But thou, O Lord, have
mercy upon us.

\emph{Thanks be to God}.

\emph{Response} 5.

\emph{Forsake me not, O holy father, for I have
already laid}\\
\emph{out my stores. I
desert thee not, my son, neither do I}\\
\emph{forsake thee; but a
fiercer conflict for the faith of Christ}\\
\emph{is in store for
thee}.

We, as aged men, receive the onset of the
skirmish, thou

being young, will have to bear off a more glorious triumph

over the persecutor; the Deacon shall follow his Bishop on

the third day.

\emph{I desert thee not, my son, neither do I
forsake thee: but a fiercer conflict for the faith of Christ is in store for thee}.

But thou, O Lord, have mercy upon us.

\emph{Thanks be to God}.

\emph{Benediction} 6.

\emph{Reader}.—Sir, be pleased to bless us.

\emph{Minister}.—God kindle the fire of His
love in our hearts.

\emph{Amen}.

\emph{Lesson} 6.
Thou gainest nothing, thou
availest nothing, O savage

cruelty! The mortal frame is gradually released from thy

tortures. Laurence departs heavenward, and thy flames fail

thee. The love of Christ surpassed the flame, and the fire

which burned around him was duller than that which was

kindled within him. O Persecutor, thou hast spent thy rage

upon the Martyr; thou hast spent it, and added to his

palm, while adding to his pain. For what part of thy de-

vice has not redounded to the conqueror's glory, when even

the instruments of his suffering are converted into deco-

rations of his triumph? Let us then rejoice, dearly beloved,

with a spiritual joy, and glory in the Lord concerning the

most blessed end of this famous man. God is wonderful

in His Saints, in whom He hath ordained for us a sanction

and an example, and hath so shown forth His glory through

the whole world, that from the rising to the setting sun,

among the refulgent lights of the Diaconate, Rome became

as honoured in her Laurence, as Jerusalem in Stephen.

But thou, O Lord, have mercy upon us.

\emph{Thanks be to God}.
\emph{Response} 6.
\emph{Blessed Laurence cried out
and said, My God I worship, Him alone I serve, and
therefore I fear not your tortures}.

My night hath no darkness, but all things
grow clear in

the light.

\emph{And therefore I fear not your tortures}.

Glory be to the Father,
\&c.

\emph{And therefore I fear not your tortures}.

NOCTURN
III.

(k) \emph{Antiphon}
They bound down his limbs upon
the bars: but while

they laid underneath live coals, the Deacon of Christ laughs

them to scorn.

\emph{Psalm} 15 (7)
Lord, who shall dwell in thy
tabernacle: or who shall rest

upon thy holy hill?

Even he that leadeth an uncorrupt life: and
doeth the thing

which is right, and speaketh the truth from his heart.

He that hath used no deceit in his tongue,
nor done evil to

his neighbour: and hath not slandered his neighbour.

He that setteth not by himself, but is lowly
in his own

eyes: and maketh much of them that fear the Lord.

He that sweareth unto his neighbour, and
disappointeth

him not: though it were to his own hindrance.

He that hath not given his money upon usury:
nor taken

reward against the innocent.

Whoso doeth these things shall never fall.
Glory be, \&c.

(k) \emph{Antiphon}
They bound his limbs upon the
bars; but while they laid

underneath live coals, the Deacon of Christ laughs them

to scorn.

(l) \emph{Antiphon}

Thou hast tried me with fire, and hast found
no wicked-

ness in me.

\emph{Psalm} 17. (8)
Hear the right, O Lord,
consider my complaint: and

hearken unto my prayer, that goeth not out of feigned lips.

Let my sentence come forth from thy presence:
and let

thine eyes look upon the thing that is equal.

Thou hast proved and visited mine heart in
the night-

season; thou hast tried me, and shalt find no wickedness

in me: for I am utterly purposed that my mouth shall not

offend.

Because of men's works, that are done
against the words

of thy lips: I have kept me from the ways of the destroyer.

O hold thou up my goings in thy paths: that
my footsteps

slip not.

I have called upon thee, O God, for thou
shalt hear me:

incline thine ear to me, and hearken unto my words.

Show thy marvellous loving-kindness, thou
that art the

Saviour of them which put their trust in thee: from such as

resist thy right hand.

Keep me as the apple of an eye: hide me under
the shadow

of thy wings.

From the ungodly that trouble me: mine
enemies compass

me round about to take away my soul.

They are inclosed in their own fat: and their
mouth

speaketh proud things.

They lie waiting in our way on every side:
turning their

eyes down to the ground;

Like as a lion that is greedy of his prey:
and as it were a

lion's whelp lurking in secret places.

Up, Lord, disappoint him, and cast him down:
deliver

my soul from the ungodly, which is a sword of thine;

From the men of thy hand, O Lord, from the
men, I say,

and from the evil world: which have their portion in this

life, whose bellies thou fillest with thy hid treasure.

They have children at their desire: and leave
the rest of

their substance for their babes.

But as for me, I will behold thy presence in
righteous-

ness; and when I awake up after thy likeness, I shall be

satisfied with it. Glory be, \&c.

(l) \emph{Antiphon}
Thou hast tried me with fire,
and hast found no wicked-

ness in me.

(m) \emph{Antiphon}

When I was questioned, I confessed the Lord:
when I am

burned, I give thanks.

\emph{Psalm} 21. (9)
The King shall rejoice in thy
strength, O Lord: exceed-

ing glad shall he be of thy salvation.

Thou hast given him his heart's desire: and
hast not

denied him the request of his lips.

For thou shalt prevent him with the blessings
of good-

ness: and shalt set a crown of pure gold upon his head.

He asked life of thee, and thou gavest him a
long life:

even for ever and ever.

His honour is great in thy salvation: glory
and great

worship shalt thou lay upon him.

For thou shalt give him everlasting felicity:
and make

him glad with the joy of thy countenance. 

And why? because the king putteth his trust
in the Lord:

and in the mercy of the Most Highest he shall not miscarry.

All thine enemies shall feel thy hand: thy
right hand

shall find out them that hate thee.

Thou shalt make them like a fiery oven in
time of thy

wrath: the Lord shall destroy them in his displeasure, and

the fire shall consume them.

Their fruit shalt thou root out of the earth:
and their seed

from among the children or men.

For they intended mischief against thee; and
imagined

such a device as they are not able to perform.

Therefore shalt thou put them to flight: and
the strings

of thy bow shalt thou make ready against the face of them.

Be thou exalted, Lord, in thine own strength:
so will we

sing, and praise thy power. Glory be, \&c.

(m) \emph{Antiphon}

When I was questioned, I confessed the Lord:
when I

am burned, I give thanks.

(n) \emph{Verse and Response}

His honour is great in thy salvation.

\emph{Glory and great worship
shalt thou lay upon him}.

\emph{The Lord's Prayer}

(\emph{privately}.)

Our Father, \&c.

And lead us not into temptation,

\emph{But
deliver us from evil}.

\emph{Absolution} 3.

The Almighty and merciful Lord absolve us
from the

chain of our sins. \emph{Amen}.

\emph{Benediction} 7.

\emph{Reader}.—Sir, be pleased to bless us.

\emph{Minister}.—May the reading of the
Gospel be to us salva-

tion and a defence. \emph{Amen}.

\emph{Lesson} 7.

\emph{John} xii. 24, 25.

At that time Jesus said to His disciples,
Verily, verily, I

say unto you, except a corn of wheat fall into the ground

and die, it abideth alone; but if it die, it bringeth forth

much fruit. He that loveth his life, shall lose it, and he that

hateth his life in this world, shall keep it unto life eternal.

\emph{Homily of St. Augustine}

The Lord Jesus Himself was the corn of wheat
to be put

to death, and to be multiplied; to be put to death by the

unbelief of the Jews, to be multiplied by the belief of the

Gentiles. Therefore, exhorting us to trace the footsteps of

His passion, He says, ``He that loveth his life shall lose it.''

Which may be understood in two ways. He that loveth,

shall lose; that is, if thou lovest, thou shalt lose. If thou

wouldest possess life in Christ, fear not that death for

Christ which is necessary. Or otherwise; He that

loveth his life, shall lose it. Love it not lest thou really
lose

it; love it not here, lest thou lose it eternally.

But
thou, O Lord, have mercy upon us.

\emph{Thanks
be to God}.

\emph{Response}

\emph{On the hot bars I denied Thee not, my God;
and when}

\emph{brought to the fire, I
confessed the Lord Jesus Christ}.

Thou hast proved, O Lord, and visited mine
heart in the

night-season.

\emph{And when brought to the fire, I confessed
the Lord Jesus Christ}.

\emph{Benediction} 8.

\emph{Reader}.—Sir, be pleased to bless us.

\emph{Minister}.—May he whose festival we
keep, intercede for

us to the Lord. \emph{Amen}.

\emph{Lesson} 8.

(\emph{homily, continued}.)

The latter of these two seems rather to be
the sense of

the Gospel. For it goes on, ``And he that hateth his life

in this world, shall keep it unto life eternal.'' Therefore,

as is said before, ``He that loveth,'' that is, ``in this world,''

he surely shall lose it; but ``he who hateth,'' namely, ``in

this world,'' shall keep it unto life eternal, A great and

marvellous saying, how it should be that a man should love

his life to its destruction, and hate it to its preservation. If

thou hast loved it perversely, then thou really hatest it; if

thou hast hated rightly, then thou hast loved it. Blessed are

they who so hate it while really saving it, as not to lose it

while loving it.

But
Thou, O Lord, have mercy upon us.

\emph{Thanks
be to God}.

\emph{Response} 8.

\emph{O Hippolytus, if thou hast faith in the
Lord Jesus Christ,}

\emph{I will both show thee
treasures, and promise thee life ever-}

\emph{lasting}.

The blessed Laurence said to Hippolyus, If
thou hast

faith in the Lord Jesus Christ.

\emph{I will both show thee treasures, and
promise thee life ever- lasting}.

Glory be to the Father, \&c.

As it was, \&c.

\emph{I will both show thee treasures, and
promise thee life ever- lasting}.

\emph{Benediction} 9.

\emph{Reader}.—Sir, be pleased to bless us.

\emph{Minister}.—The King of Angels lead us
on to the fellow-

ship of the inhabitants of heaven. \emph{Amen}. \emph{Lesson} 9.

(\emph{Homily continued}.)

But beware of the desire of self-murder
stealing on thee,

as if from the precept of hating thy life in this world. For

hence certain evil-tempered and perverse men, and to them-

selves more cruel and wicked murderers, give themselves

to the flames, drown themselves in the water, break their

bones down precipices, and so perish. This is not from

Christ's teaching, who even answered to the devil, suggest-

ing to Him such a fall, ``Get thee behind me, Satan; it is

written, Thou shalt not tempt the Lord thy God.'' But to

Peter he said, signifying by what death he should glorify

God, ``When thou wast young, thou girdedst thyself and

walkedst whither thou wouldest; but when thou art old,

another shall gird thee, and shall carry thee whither thou

wouldest not.'' Where he sufficiently intimated that he who

follows Christ's footsteps must be put to death, not by

himself but by another.

\emph{Te Deum}

We praise thee, O God; we acknowledge thee to
be the

Lord.

All the earth doth worship thee: the Father
everlasting.

To thee all Angels cry aloud: the heavens and
all the

Powers therein, \&c.

\xsection{§ 6. Matin Service for March 21.}

Bishop Ken's day

[For Social or Private Devotion.]

O Lord, open Thou my lips.

\emph{And my mouth shall show forth Thy praise}.

O God, make speed to save me.

\emph{O Lord make haste to help me}.

Glory be, \&c.

As it was, \&c. \emph{Amen}.

Praise to Thee, O Lord, King of eternal glory.

(a) \emph{Invitatory with Psalm} 95.

\emph{Vide} pp. 17 \& 21.

O come, let us worship the Lord, the King of
Confessors.

O come, let us sing unto the Lord: let us
heartily re-

joice, \&c.

(b) \emph{Hymn}

In witness of his Lord,

In humble following
of his Saviour dear,

This is the man to wield the unearthly sword,

Warring unharmed with
sin and fear.

Who, Lord, uncalled by
Thee,

Dare touch Thy
Spouse, Thy very self below?

Or who dare count him summoned worthily,

Except Thine hand and
seal he show?

Where can Thy seal be
found,

But on the chosen
seed from age to age,

By Thine anointed heralds duly crowned,

As kings and priests,
Thy war to wage?

\emph{Or this}.

Lord, and what shall
this man do?

Ask'st thou,
Christian, for thy friend?

If his love for Christ be true,

Christ hath told thee
of his end:

This is he whom God approves,

This is he whom Jesus loves.

Ask not of Him more than
this,

Leave it to his
Saviour's breast,

Whether early called to bliss,

He in youth shall
find his rest,

Or armed at his station wait,

Till his Lord be at
the gate.

Gales from heaven, if so
He will,

Sweeter melodies can
wake,

On the lonely mountain rill,

Than the meeting
waters make,

Who hath the Father and the Son,

May be left, but not
alone.

———————

NOCTURN
I.

(c) \emph{Antiphon}
Blessed is the man whose
delight is in the law of the

Lord.

\emph{Psalm} 1. (1)
Blessed is the man that,
\&c.

(c) \emph{Antiphon}
Blessed is the man whose
delight is in the law of the Lord,

who hath not walked in the counsel of the ungodly, nor

stood in the way of sinners, and hath not sat in the seat of

the scornful.

(d) \emph{Antiphon}

Desire of me, and I will give thee the
heathen for thine

inheritance.

\emph{Psalm} 2. (2)
Why do the heathen, \&c.

(d) \emph{Antiphon}
Desire of me, and I will give
thee the heathen for thine

inheritance, and the utmost parts of the earth for thy pos-

session. Thou shalt bruise them with a rod of iron.

(e) \emph{Antiphon}

Thou, O Lord, art my worship.

\emph{Psalm} 3. (3)
Lord, how are they, \&c.

(e) \emph{Antiphon}
Thou, O Lord, art my worship,
and the lifter up of my

head. I did call upon the Lord with my voice, and He

heard me out of His holy hill.

(1) \emph{Verse and Response}

The Lord loved him and adorned him.

\emph{And clothed him in
a robe of glory}.

\emph{The Lord's Prayer}

(\emph{privately}.)

Our Father, \&c.

And lead us not into temptation.

\emph{But deliver us from evil}. \emph{Absolution} 1.
O Lord Jesus Christ,
mercifully hear the supplications of

Thy people, and grant us Thy peace all the days of our life,

who livest and reignest with the Father and the Holy Ghost,

world without end. \emph{Amen}.

\emph{Benediction} 1.
\emph{}\\
\emph{Reader}.—Sir, be pleased to bless
us.

\emph{Minister}.—The Lord bless us and keep
us. \emph{Amen}.

\emph{Lesson} 1.

1 \emph{Tim}. iii. 1-6.
This \emph{is} a true saying,
If a man desire the office of a Bishop,

he desireth a good work.

A Bishop then must be blameless, the husband
of one wife,

vigilant, sober, of good behaviour, given to hospitality, apt to

teach;

Not given to wine, no striker, not greedy of
filthy lucre;

but patient, not a brawler, not covetous:

One that ruleth well his own house, having
his children in

subjection with all gravity;

(For if a man know not how to rule his own
house, how

shall he take care of the Church of God?)

Not a novice, lest being lifted up with pride
he fall into

the condemnation of the devil.

But thou, O Lord, have mercy upon us.

\emph{Thanks be to God}.

\emph{Response} 1.

\emph{Well done, thou good and faithful servant,
thou hast been}\\
\emph{faithful over
a few things, I will make thee ruler over}\\
\emph{many things:
enter thou into the joy of thy Lord}.

Lord, Thou deliveredst unto me five talents,
behold, I have

gained beside them five talents more.

\emph{Enter thou into the joy of thy Lord}.

\emph{Benediction} 2.

\emph{Reader}.—Sir, be pleased to bless us.

\emph{Minister}.—The Lord make His face to
shine upon us,

and be gracious unto us. \emph{Amen}.

\emph{Lesson} 2.

\emph{Tit}. i. 7-11.
For a Bishop must be
blameless, as the steward of God;

not self-willed, not soon angry, not given to wine, no striker,

not given to filthy lucre;

But a lover of hospitality, a lover of good
men, sober, just,

holy, temperate;

Holding fast the faithful word as he hath
been taught, that

he may be able by sound doctrine both to exhort and to con-

vince the gainsayers.

For there are many unruly and vain talkers
and deceivers,

specially they of the circumcision. 

Whose mouth must be stopped, who subvert
whole

houses, teaching things which they ought not, for filthy

lucre's sake.

But Thou, O Lord, have mercy upon us.

\emph{Thanks be to God}.

\emph{Response} 2.

\emph{Let Thy Thummim
and Thy Urim be with Thy holy one,}

\emph{whom Thou didst prove at Massah, and with whom Thou}

\emph{didst strive at the waters of Meribah; they shall put}

\emph{incense before Thee, and whole burnt sacrifice upon}

\emph{Thine altar}.

Bless, Lord, his substance, and accept the
work of his

hands; smite through the loins of them that rise against him,

and of them that hate him, that they rise not again:

\emph{They shall put incense before Thee, and
whole burnt sacrifice upon Thine altar}.

\emph{Benediction} 3.

\emph{Reader}.—Sir, be pleased to bless us.

\emph{Minister}.—The Lord lift up His
countenance upon us, and

give us peace. \emph{Amen}.

\emph{Lesson} 3.

\emph{Tit}. ii. 1-8.
But speak thou the things
which become sound doc-

trine.

That the aged men be sober, grave, temperate,
sound in

faith, in charity, in patience.

The aged women likewise, that \emph{they be}
in behaviour as

becometh holiness, not false accusers, not given to much wine,

teachers of good things;

That they may teach the young women to be
sober, to love

their husbands, to love their children.

\emph{To be} discreet, chaste, keepers at
home, good, obedient to

their own husbands, that the word of God be not blasphemed.

Young men likewise exhort to be sober-minded.

In all things showing thyself pattern of good
works; in

doctrine \emph{showing} uncorruptness, gravity, sincerity,

Sound speech, that cannot be condemned; that
he that is

of the contrary part may be ashamed, having no evil thing

to say of you.

But Thou, O Lord, have mercy upon us.

\emph{Thanks be to God}.

\emph{Response} 3.

\emph{Let the saints be joyful with glory, let
them rejoice in their}\\
\emph{beds, let the
praises of God be in their mouths, and a two-}\\
\emph{edged sword in
their hands: to bind their kings in chains,}\\
\emph{and their
nobles with links of iron}. 

That they may be avenged of them, as it is
written, Such

honour have all His saints.

\emph{To bind their kings in chains, and their
nobles with links of iron}.

Glory be, \&c.

\emph{To bind their kings in chains, and their
nobles with links of iron}.

———————

NOCTURN
II.

(g) \emph{Antiphon}
Thou hast put gladness in my
heart.

\emph{Psalm} 4. (4)
Hear me when I call, \&c.

(g) \emph{Antiphon}
Thou hast put gladness in my
heart, since the time that

their corn, and wine, and oil increased.

(h) \emph{Antiphon}

Lead me, O Lord, in Thy righteousness.

\emph{Psalm} 5. (5)
Ponder my words, \&c.

(h) \emph{Antiphon}
Lead me, O Lord, in Thy
righteousness, because of mine

enemies; make Thy way plain before my face.

(i) \emph{Antiphon}

Out of the mouth of very babes and sucklings
hast Thou

ordained strength.

\emph{Psalm} 8. (6)
O Lord, our Governor, \&c.

(i) \emph{Antiphon}
Out of the mouth of very babes
and sucklings hast Thou

ordained strength, because of Thine enemies, that Thou

mightest still the enemy and the avenger.

(j) \emph{Verse and Response}

The Lord hath chosen Him as a priest unto
Himself.

To sacrifice to Him
the offering of praise.

\emph{The Lord's Prayer}

(\emph{privately}.)

Our Father, \&c.

And lead us not into temptation,

\emph{But deliver us from evil}.

\emph{Absolution} 2.

Grant, O Lord, we beseech Thee, that we, who
for our

evil deeds do worthily deserve to be punished, by the com-

fort of Thy grace may mercifully be relieved, who livest and

reignest with the Father and the Holy Ghost, ever one

God, world without end. \emph{Amen}.

\emph{Benediction} 4.

\emph{Reader}.—Sir, be pleased to bless us.

\emph{Minister}. —The love of God be upon
us now and for ever.

\emph{Lesson} 4.
Thomas Ken, the son of an
ancient family, was born at

Berkhamstead, in Hertfordshire, in the year of grace, 1637,

and educated at Winton and Oxford, on the foundation of

William of Wykeham, of famous memory, some time Bishop

of Winton. Admitted into holy orders, he commenced a

course of preaching at St. John's church, near Winton, where

there was no preacher, with so blessed an effect, that many

Anabaptists came over to the Church, and received baptism at

his hands. That he might have time for these active duties,

for study also, and for prayer, he restricted himself to but

one sleep, rising at one or two of the clock in the morning,

and sometimes sooner; which practice grew into a habit, and

continued with him almost till his last illness. After a while

he was made chaplain to the king's niece, the Princess of

Orange, and passed over to Holland; where, after gaining

her entire esteem for his most prudent behaviour and strict

piety, he fell under the displeasure of the Prince, for inter-

fering with one of his courtiers, who had seduced a young

English gentlewoman, and was eventually obliged to leave

the royal service.

But Thou, O Lord, have mercy upon us.

\emph{Thanks be to God}.

\emph{Response} 4.

\emph{Princes have
persecuted me without a cause; but my heart}

\emph{standeth in awe of Thy word}.

I will speak of Thy testimonies also, even
before kings,

and will not be ashamed.

\emph{But my heart standeth in awe of Thy word}.

\emph{Benediction} 5.

\emph{Reader}.—Sir, be pleased to bless us.

\emph{Minister}.—The grace of our Lord
Jesus Christ be upon us,

now and for ever. \emph{Amen}.

\emph{Lesson} 5.
Afterwards, when he was at
Winton, in his Prebendal

house, the king came thither with his court and mistress,

whom his harbingers ordered to be lodged where Ken

dwelt; but fearing God more than the face of the king,

he refused her admittance, and obliged her to seek another

lodging. And herein was seen how winning is holy severity,

even with those who suffer for it; for a vacancy soon occur-

ring in the see of Bath and Wells, the king himself, as his

own especial act, named Ken to fill it; and he was con-

secrated thereunto out St. Paul's day, in the year 1684.

Moreover, during the king's last illness, he was admitted

to his chamber, and gave close attendance at his bedside

for at least three whole days and nights, without any in-

termission, watching at proper intervals to suggest pious 

thoughts and prayers; in which time one of the King's

mistresses coming in, the Bishop prevailed with him to

have her removed, and induced him further to send for the

Queen, and ask her forgiveness for his long infidelity.

And though he was not able to do all he had wished for

his dying Sovereign, he did all that was allowed him.

But Thou, O Lord, have mercy upon us.

\emph{Thanks be to God}.

\emph{Response} 5.

\emph{Many shall be
purified, and made white, and tried: and}

\emph{many of them that sleep in the dust of the earth shall}

\emph{awake}.

And they that be wise shall shine as the
brightness of

the firmament, and they that turn many to righteousness

as the stars for ever and ever.

\emph{And many of them that sleep in the dust of
the earth shall awake}.

\emph{Benediction} 6.

\emph{Reader}.—Sir, be pleased to bless us.

\emph{Minister}.—The fellowship of the Holy
Ghost be with us

now and for ever. \emph{Amen}.

\emph{Lesson} 6.
When the King's brother
succeeded to the throne, Ken

showed his loyalty to him indeed, but in word was free

spoken. At length, when the King advanced in daring, the

Bishop refusing to submit the conduct of the Church to his

pleasure, and mindful of the rights of Christ's heritage, was

in consequence committed, together with six of his brethren,

to the Tower, on a charge of treason. Afterwards, when the

King and his family were dethroned for arbitrary doings, he

showed his true loyalty to him by refusing to acknowledge

the new dynasty, and lost his station in the state rather than

violate his allegiance. Being driven from his see by the civil

power, he died in obscurity in the year 1710. Thus he gave

to Cæsar the things that be Cæsar's, and to God the things

that be God's. He was as meek, gentle, and affectionate in

his bearing, as he was bold in the cause of the Gospel; and

he took his troubles cheerfully and lightly. He possessed, in

an especial way, that most excellent gift of charity. Once,

when four thousand pounds fell to his see, he gave great

part of it to the French Protestants then under persecution;

and, when he was deprived, all his means, after the sale of

goods at his palace and elsewhere, was not more than seven

hundred pounds. When state interests interfered with the

prosperity of the Church in Scotland, he said, he conceived

great hopes that God would have compassion on the English

branch of it, if she did but compassionate and support her

sister: and he gave testimony concerning his belief
shortly

before his death, saying that he died in the Holy Catholic and

Apostolic Faith, professed by the whole Church before the

disunion of east and west. Such was Ken, a burning and shining

light, bringing back primitive times.

But thou, O Lord, have mercy upon us.

\emph{Thanks be to God}.

\emph{Response} 6.

\emph{I said I have
laboured in vain, I have spent my strength}

\emph{for nought and in vain; yet surely my judgment is}

\emph{with the Lord, and my work with my God}.

He hath made my mouth like a sharp sword: in
the

shadow of His hand hath He hid me.

\emph{Surely my judgment is with the Lord, and
my work with}\\
\emph{my God}.

Glory be, \&c.

\emph{Surely my judgment is with the Lord, and
my work with}\\
\emph{my God}.

———————

NOCTURN
III.

(k) \emph{Antiphon}
Lord, he shall dwell in Thy
tabernacle.

\emph{Psalm} 15. (7)
Lord, who shall dwell, \&c.

(k) \emph{Antiphon}
Lord, he shall dwell in Thy
tabernacle: he hath led an

uncorrupt life, he hath done the thing that is right.

(l) \emph{Antiphon}

He asked life of Thee, and Thou gavest him a
long life.

\emph{Psalm} 21. (8)
The king shall rejoice,
\&c.

(l) \emph{Antiphon}
He asked life of Thee, and
Thou gavest him a long life:

glory and great worship shalt Thou lay upon him: Thou

shalt set a crown of pure gold upon his head.

(m) \emph{Antiphon}

He shall receive the blessing from the Lord.

\emph{Psalm} 24. (9)
The earth is the Lord's,
\&c.

(m) \emph{Antiphon}
He shall receive the blessing
from the Lord, and righte-

ousness from the God of his salvation; for this is the

generation of them that seek him.

(n) \emph{Verse and Response}

The key of David will I lay upon his
shoulder.

\emph{He shall open and
none shall shut, and He shall shut and none shall open}.

\emph{The Lord's Prayer}

(\emph{privately}.)

Our Father, \&c.

And lead us not into temptation.

\emph{But deliver us from evil}. \emph{Absolution} 3.
Grant to us, merciful Lord,
that, whereas we are sore let

and hindered in running the race that is set before us, Thy

bountiful grace and mercy may speedily help and deliver us,

who livest and reignest with the Father and the Holy Ghost,

ever one God, world without end. \emph{Amen}.

\emph{Benediction} 7.

\emph{Reader}.—Sir, be pleased to bless us.

\emph{Minister}.—The reading of the Gospel
be to us salvation

and a defence. \emph{Amen}.

\emph{Lesson} 7.

\emph{Luke} xxii. 25-30.
And he said unto them, The
Kings of the Gentiles exer-

cise lordship over them: and they that exercise authority

upon them are called benefactors.

But ye shall not be so: but he that is
greatest among you,

let him be as the younger: and he that is chief, as he that

doth serve.

For whether is greater, he that sitteth at
meat, or he that

serveth? \emph{is} not he that sitteth at meat? but I am among

you as he that serveth.

Ye are they which have continued with me in
my temp-

tations.

And I appoint unto you a kingdom, as my
Father hath

appointed unto me;

That ye may eat and drink at my table in my
kingdom,

and sit on thrones, judging the twelve tribes of Israel.

(\emph{From the works of Jeremy Taylor, Bishop}.)
The nature of honour is to be
a reward of virtue; and

by how much greater the reward is, by so much the greater

is the honour which is conferred. What honour shall it

then be, when God shall give unto those who served Him,

not only to tread upon the stars, to inhabit the palaces of

honour, to be lords of the world, but, transcending all that

is created, and finding nothing among his riches sufficient

to reward them, shall give then His own Infinite Essence, to

enjoy, as a recompence of their holiness, not for a day, but

to all eternity. O happy labour of the victorious, and glo-

rious combat of the just, against the vices and temptations

of the world, when victory deserves so inestimable a crown!

How great shall be that glory, when a just soul shall, in the

presence of an infinite number of angels, sit in the same

throne with Christ: and shall, by the just sentence of God,

be praised for a conqueror over the world, and the invisible

powers of hell! What can it desire more, than to be par-

taker of all those Divine goods, and even to accompany

Christ in the same throne? How cheerfully do they bear 

all afflictions for Christ, who with a lively faith and certain

hope apprehend so sublime honours.

Thou, then, O Lord,
have mercy upon us.

\emph{Thanks be to God}.

\emph{Response} 7.

\emph{Whosoever shall confess Me before men, him
shall the Son of}\\
\emph{man also
confess before the Angels of God}.

To him that overcometh will I grant to sit
with Me in

My throne.

\emph{Him shall the Son of man also confess
before the Angels of God}.

\emph{Benediction} 8.

\emph{Reader}.—Sir, be pleased to bless us.

\emph{Minister}.—Grace be with all them
that love our Lord

Jesus Christ in sincerity. \emph{Amen}.

\emph{Lesson} 8.
If the applause of men, and
the good opinion which they

have from others, be esteemed an honour, what shall be the

applause of heaven, and the good opinion not only of Saints

and Angels, but of God himself, whose judgment cannot

err? David took it for a great honour, that the daughter of

his king was judged a reward of his valour; God surpasses

this, and honours so much the service of his elect, that He

pays their merits with no less a reward than Himself. Be-

sides this, he who is most known, and is praised and cele-

brated for good and virtuous by the greatest multitude, is

esteemed the most glorious and honourable person: but all

this world is a solitude in respect of the citizens of heaven,

where innumerable angels approve and praise the virtuous

actions of the just; and they likewise are nothing: and all

creatures, men, and angels, are but a solitary wilderness, in

respect of the Creator. What man so glorious upon earth,

whose worth and valour hath been known to all? Those

who were born before him could not know him; but the

just in heaven shall be known by all, past and to come, and

by all the angels, and by the King of men and angels.

But thou, O Lord,
have mercy upon us.

\emph{Thanks be to God}.

\emph{Response} 8.

\emph{In the sight of the unwise they seemed to
die, and their de-}\\
\emph{parture is
taken for misery, but they are in peace}.

Though they be punished in the sight of men,
yet is their

hope full of immortality.

\emph{But they are in
peace}.

Glory be to the Father, \&c.

\emph{But they are in
peace}. 

\emph{Benediction} 9.

\emph{Reader}.—Sir, be pleased to bless us.

\emph{Minister}.—The Lord of peace Himself
give you peace

always by all means. \emph{Amen}.

\emph{Lesson} 9.
The honour of the just in
heaven depends not, like that

of the earth, upon accidents and reports, nor is exposed to

dangers, or measured by the discourse of others; but in itself

contains its own glory and dignity. The Romans erected

statues unto those whom they intended to honour; because,

being mortal, there should something remain after death, to

make their persons and services, which they had done to

the common weal, known to posterity; but in heaven there

is no need of this artifice, because those, which are there

honoured, are immortal, and shall have in themselves some

character engraved, as an evident and clear token of their

noble victories and achievements; what greater honour than

to be friends of God, sons, heirs, and kings in the realm of

heaven?

\emph{Te Deum}
We praise Thee, O God; we
acknowledge Thee to be

the Lord.

All the earth doth worship Thee: the Father
everlasting.

To Thee all angels cry aloud; the heavens,
and all the

Powers therein, \&c.

\xsection{§ 7. Service in Commemoration of the Dead in Christ.}

[It is perhaps scarcely necessary to observe, that
the complaints and supplications contained in the Psalms and Lessons
selected, are introduced in memory of those temporal miseries,
especially in sickness and dying, from which death opens an escape. On
the subject of the state of the Dead, \emph{vide}No. 5 of Mr. Dodsworth's recent Sermons.]

\xsubsection{FIRST VESPERS}

\emph{Antiphon}
Behold, O Lord, how that I am
Thy servant.

\emph{Psalm}\emph{}116. (l)

(\emph{Instead of Glory be, \&c}.)
I am well pleased, \&c.

With thee is the well of life.

\emph{And in Thy light
they shall see light}.

\emph{Antiphon}
Behold, O Lord, how that I am
Thy servant; Thou hast

broken my bonds in sunder.

\emph{Antiphon}

My soul hath long dwelt among.

\emph{Psalm}\emph{}120. (2)
When I was in trouble, \&c.

With Thee, \&c.

\emph{And in Thy, \&c}.

\emph{Antiphon}
My soul hath long dwelt among
them that are enemies

unto peace.

\emph{Antiphon}

The Lord shall preserve Thee from all evil.

\emph{Psalm}\emph{}121. (3)
I will lift up mine eyes,
\&c.

With Thee, \&c.

\emph{Antiphon}
The Lord shall preserve Thee
from all evil; yea, it is even

He that shall keep Thy soul.

\emph{Antiphon}

I look for the Lord.

\emph{Psalm}\emph{}130. (4)
Out of the deep, \&c.

With Thee, \&c.

\emph{Antiphon}
I look for the Lord, my soul
doth wait for Him.

\emph{Antiphon}

Though I walk in the midst of trouble.

\emph{Psalm} 138. (5)
I will give thanks unto Thee,
\&c.

With Thee, \&c.

\emph{Antiphon}
Though I walk in the midst of
trouble, yet shalt Thou

refresh me. 

(fff) \emph{Verse and Response}

I heard a voice from heaven, saying unto me,

\emph{Blessed are the dead that die in the Lord}.

(ggg) \emph{Ant}.
All that the Father giveth Me.

\emph{Magnificat}
My
soul doth magnify, \&c. Glory be, \&c.

\emph{Ant}.
All that the Father giveth Me,
shall come to Me; and him

that cometh to Me, I will in no wise cast out.

\emph{The Lord's Prayer}\emph{}\\
(\emph{privately}.)

Our Father, \&c.

And lead us not into temptation,

\emph{But deliver us from evil}.

\emph{Psalm}\emph{}146.
Praise the Lord, O my soul,
\&c.

With Thee is the well of life.

\emph{And in Thy light
shall they see light}.

In
the valley of the shadow of death.

\emph{They
shall fear no evil}.

Into
Thy hands we commend their spirit.

\emph{For Thou hast
redeemed them, O Lord, Thou God of truth}.\emph{}\\
Lord hear our prayer.

\emph{And let our cry come unto Thee}.

The Lord be
with you.

\emph{And with thy spirit}.

Let us pray.

O God, the Maker and Redeemer of all
believers, grant to

all Thy servants a merciful judgment at the last day, that they,

in the face of all creatures, may then be acknowledged as Thy

true children, through our Lord Jesus Christ. \emph{Amen}.

O Almighty God, who hast knit together
thine elect in one

communion and fellowship in the mystical body of Thy Son

Christ our Lord, grant us grace to follow Thy blessed Saints

in all virtuous and godly living, that we may come to those

unspeakable joys which Thou hast prepared for them that

unfeignedly love Thee, through Jesus Christ our Lord. \emph{Amen}.

Almighty God, with whom do live the
spirits of them that

depart hence in the Lord, and with whom the souls of the

faithful, after they are delivered from the burden of the flesh,

are in joy and felicity, we give Thee hearty thanks for that

it pleased Thee, as on this day, to deliver our dear brother

out of the miseries of this sinful world; beseeching Thee

that it may please Thee of Thy gracious goodness, shortly

to accomplish the number of Thine elect, and to hasten Thy

kingdom, that we, with all those that are departed in the true

faith of Thy holy name, may have our perfect consummation

and bliss, both in body and soul, in Thy eternal and ever-

lasting glory, through Jesus Christ our Lord. \emph{Amen}. 

O God, merciful and faithful, the aid
of all that trust in

Thee, keep safely under the shadow of thy wings our-

selves, our relations, and friends, and all believers, even

Thy whole Church, that we may enjoy Thy presence alway,

and increase in Thy Holy Spirit more and more, till we

come to Thine everlasting kingdom, through our Lord and

Saviour Jesus Christ. \emph{Amen}.

With Thee is the well
of life.

\emph{And in Thy light shall they see light}.

May their souls rest
in peace. \emph{Amen}.

———————

\xsubsection{MATINS}

O Lord, open Thou my lips.

\emph{And my mouth shall
show forth Thy praise}.

O God, make speed to save me.

\emph{O Lord, make haste
to help me}.

Glory be, \&c. Amen. Hallelujah.

(a) \emph{Invitatory}
Let us worship the Lord of
spirits: for all live unto Him.

\emph{Psalm} 95.
O come, let us sing, \&c.

With Thee is the well of life.

\emph{And in Thy light shall they see light}.

———————

NOCTURN
I.

\emph{Antiphon}
In Thy fear will I worship.

\emph{Psalm} 5.
Ponder my words, \&c.

\emph{Antiphon}
With Thee, \&c.

In Thy fear will I worship towards Thy Holy
Temple.

\emph{Antiphon}

The Lord hath heard my petition.

\emph{Psalm} 6.
O Lord, rebuke me not, \&c.

With Thee, \&c.

\emph{Antiphon}
The Lord hath heard my
petition, the Lord will receive

my prayer. 

\emph{Antiphon}

We will rejoice.

\emph{Psalm} 20.
The Lord hear thee in the day
of trouble, \&c.

With Thee, \&c.

\emph{Antiphon}
We will rejoice in Thy
salvation, and triumph in the

name of the Lord our God.

(f) \emph{Verse and Response}

In the valley of the shadow of death.

\emph{They shall fear no evil}.

\emph{The Lord's Prayer}

(\emph{privately}.)

Our Father, \&c.

\emph{Lesson} 1.

\emph{Job} ix. 11-20.
Lo, he goeth by me, and I see \emph{him}
not: he passeth on

also, but I perceive him not.

Behold, he taketh away, who can hinder him?
who will

say unto him, What doest thou?

\emph{If} God will not withdraw his anger,
the proud helpers do

stoop under him.

How much less shall I answer him, \emph{and}
choose out my

words \emph{to reason} with him?

Whom, though I were righteous, \emph{yet}
would I not answer,

\emph{but} I would make supplication to my judge.

If I had called, and he had answered me; \emph{yet}
would I not

believe that he had hearkened unto my voice.

For he breaketh me with a tempest, and
multiplieth my

wounds without cause.

He will not suffer me to take my breath, but
filleth me

with bitterness.

If I \emph{speak} of strength, lo, \emph{he is}
strong: and if of judgment,

who shall set me a time \emph{to plead}?

If I justify myself, mine own mouth shall
condemn me: \emph{if}

\emph{I say} I \emph{am} perfect, it shall also prove me
perverse.

\emph{Response} 1.

\emph{I know that my
Redeemer liveth, and that he shall stand}

\emph{on the latter day upon the earth; and in my flesh shall}

\emph{I see God}.

Whom I shall see for myself, and my eyes
shall behold,

and not another.

\emph{And in my flesh shall I see God}.

\emph{Lesson} 2.

\emph{Job} x. 1-7.

My soul is weary of my life; I will leave my
complaint

upon myself; I will speak in the bitterness of my soul.

I will say unto God, Do not condemn me; show
me

wherefore thou contendest with me.

\emph{Is it} good unto thee that thou
shouldest oppress, that thou

shouldest despise the work of thine hands, and shine upon

the counsel of the wicked? 

Hast thou eyes of flesh? or seest thou as man
seeth?

\emph{Are} thy days as the days of man? \emph{are}
thy years as man's

days?

That thou enquirest after mine iniquity, and
searchest

after my sin?

Thou knowest that I am not wicked: and \emph{there
is} none

that can deliver out of thine hand.

\emph{Response} 2.

\emph{I should utterly have fainted, but that I
believe verily to see the goodness of the Lord in the land of the living}.

O tarry thou the Lord's leisure: be strong
and He shall

comfort thine heart, and put thou thy trust in the Lord.

\emph{I believe verily to see the goodness of
the Lord in the land of the living}.

\emph{Lesson} 3.

\emph{Job} x. 14-22.

If I sin, then thou markest me, and thou wilt
not acquit

me from mine iniquity.

If I be wicked, woe unto me; and \emph{if} I
be righteous, \emph{yet}

will I not lift up my head. I \emph{am} full of confusion;
therefore

see thou mine affliction;

For it increaseth. Thou huntest me as a
fierce lion; and

again thou showest thyself marvellous upon me.

Thou renewest thy witnesses against me, and
increasest

thine indignation upon me; changes and war \emph{are} against
me!

Wherefore then hast thou brought me forth out
of the

womb? Oh that I had given up the ghost, and no eye had

seen me!

I should have been as though I had not been;
I should

have been carried from the womb to the grave.

\emph{Are} not my days few? cease \emph{then},
\emph{and} let me alone, that

I may take comfort a little.

Before I go \emph{whence} I shall not return,
\emph{even} to the land of

darkness and the shadow of death;

A land of darkness, as darkness \emph{itself};
and of the shadow

of death, without any order, and \emph{where} the light \emph{is}
as

darkness.

\emph{Response} 3.

\emph{The people that
walked in darkness have seen a great}

\emph{light; they that dwelt in the land of the shadow of}

\emph{death, upon them hath the light shined}.

Awake thou that sleepest, and arise from the
dead, and

Christ shall give thee light.

\emph{They that dwell in the land of the shadow
of death, upon them hath the light shined}. 

With thee is the well of life,

\emph{And in thy light
shall we see light}.

\emph{They that dwell in the land of the shadow
of death, upon them hath the light shined}.

———————

NOCTURN
II.

\emph{Antiphon}
He shall feed me.

\emph{Psalm} 23. (4)
The Lord is my shepherd,
\&c.

With Thee, \&c.

\emph{Antiphon}
He shall feed me in a green
pasture.

\emph{Antiphon}

Mine eyes are ever looking.

\emph{Psalm} 25. (5)
Unto Thee, O Lord, will I lift
up my soul.

With Thee, \&c.

\emph{Antiphon}
Mine eyes are ever looking
unto the Lord, for he shall

pluck my feet out of the net.

\emph{Antiphon}

In the time of trouble.

\emph{Psalm} 27. (6)
The Lord is my light, \&c.

With Thee, \&c.

\emph{Antiphon}
In the time of trouble He
shall hide me in His tabernacle.

(j) \emph{Verse and Response}

Let the Saints be joyful with glory.

\emph{Let them rejoice
in their beds}.

\emph{The Lord's Prayer}

(\emph{privately}.)

Our Father, \&c.

\emph{Lesson} 4.

\emph{Job} xiii. 20-28.
Only do not two \emph{things}
unto me: then will I not hide my-

self from thee.

Withdraw thine hand far from me: and let not
thy dread

make me afraid.

Then call thou, and I will answer: or let me
speak, and

answer thou me.

How many \emph{are} mine iniquities and sins?
make me to

know my transgression and my sin.

Wherefore hidest thou thy face, and holdest
me for thine

enemy?

Wilt thou break a leaf driven to and fro? and
wilt thou

pursue the dry stubble?

For thou writest bitter things against me,
and makest

me to possess the iniquities of my youth. 

Thou puttest my feet also in the stocks, and
lookest nar-

rowly unto all my paths; thou settest a print upon the heels

of my feet.

And he, as a rotten thing, consumeth, as a garment that

is moth-eaten.

\emph{Response} 4.

\emph{Comfort us again now after the time that thou hast}

\emph{plagued us, and for the years wherein we have}

\emph{suffered
adversity}.

Show Thy servants Thy work, and their children Thy

glory.

\emph{For the years wherein we have suffered adversity}.

\emph{Lesson} 5.

\emph{Job} xvi. 7-22.

But now he hath made me weary: thou hast made

desolate all my company.

And thou hast filled me with wrinkles, \emph{which} is a witness

\emph{against me}: and my leanness rising up in me, beareth wit-

ness to my face.

He teareth \emph{me} in his wrath, who hateth me: he gnasheth

upon me with his teeth; mine enemy sharpeneth his eyes

upon me.

They have gaped upon me with their mouth, they have

smitten me upon the cheek reproachfully, they have

gathered themselves together against me.

God hath delivered me to the ungodly, and turned me

over into the hands of the wicked.

I was at ease, but he hath broken me asunder: he hath

also taken \emph{me} by my neck, and shaken me to pieces, and

set me up for his mark.

His archers compass me round about, he cleaveth my

reins asunder, and doth not spare; he poureth out my gall

upon the ground.

He breaketh me with breach upon breach, he runneth

upon me like a giant.

I have sewed sackcloth upon my skin, and defiled my

horn in the dust.

My face is foul with weeping, and on my eye-lids \emph{is} the

shadow of death;

Not for \emph{any} injustice in mine hands: also my prayer \emph{is}

pure.

O earth, cover thou not my blood, and let my cry have

no place.

Also, now, behold, my witness \emph{is} in heaven, and my

record \emph{is} on high.

My friends scorn me: but \emph{mine} eye poureth out \emph{tears}

unto God. 

Oh that one might plead for a man with God, as a man

\emph{pleadeth} for his neighbour!

When a few years are come, then I shall go the way

\emph{whence} I shall not return.

\emph{Response} 5.

\emph{My God hath sent his Angel, and hath shut the lions'}

\emph{mouths that they have not hurt me}.

Forasmuch as before Him innocency was found in me;

\emph{He hath shut the lions' mouths that they have not hurt me}.

\emph{Lesson} 6.

Job xiv. 1-10.

Man \emph{that is} born of a woman \emph{is} of few days, and full of

trouble.

He cometh forth like a flower, and is cut down; he

fleeth also as a shadow, and continueth not.

And dost thou open thine eyes upon such an one, and

bringest me into judgment with thee.

Who can bring a clean thing out of an unclean? Not

one.

Seeing his days \emph{are} determined, the number of his months

\emph{are} with thee, thou hast appointed his bounds that he cannot

pass.

Turn from him that he may rest, till he shall accomplish,

as an hireling, his day.

For there is hope of a tree, if it be cut down, that it will

sprout again, and that the tender branch thereof will not

cease.

Though the root thereof wax old in the earth, and the

stock thereof die in the ground;

\emph{Yet} through the scent of water it will bud, and bring

forth boughs like a plant.

But man dieth, and wasteth away: yea, man giveth up

the ghost, and where \emph{is} he?

\emph{Response}

\emph{Thy sun shall no more go down, neither shall thy moon}

\emph{withdraw itself, for the Lord shall be thine everlasting}

\emph{light}.

And the days of thy mourning shall be ended.

\emph{For the Lord shall be thine everlasting light}.

With thee is the well of life, \&c.

\emph{For the Lord shall be thine everlasting light}.

———————

NOCTURN
III.

\emph{Antiphon}

He brought me out.

\emph{Psalm} 40. (7)

I waited patiently, \&c.—With Thee,
\&c.

\emph{Antiphon}

He brought me out of the horrible pit, out of
the mire

and clay.

\emph{Antiphon}

Heal my soul.

\emph{Psalm} 41. (8)

Blessed is he, \&c.—With Thee, \&c.

\emph{Antiphon}

Heal my soul, for I have sinned against Thee.

\emph{Antiphon}

My soul is athirst.

\emph{Psalm} 42. (9)

Like as the hart, \&c.—With Thee,
\&c.

\emph{Antiphon}

My soul is athirst for God, yea, even for the
living God;

when shall I come to appear before the presence of God?

(n) \emph{Verse and Response}

O deliver not the soul of Thy turtledove unto
the multi-

tude of the enemies.

\emph{And forget not the congregation of the
poor for ever}.

\emph{The Lord's Prayer}

(\emph{privately}.)
Our Father, \&c.

\emph{Lesson} 7.

\emph{Job} xvii. 11-16.

My days are passed, my purposes are broken
off, \emph{even}

the thoughts of my heart.

They change the night into day: the light \emph{is}
short be-

cause of darkness.

If I wait, the grave \emph{is} mine house: I
have made my bed

in the darkness.

I have said to corruption, Thou \emph{art} my
father: to the

worm, \emph{Thou art} my mother and my sister.

And where is now my hope? as for my hope, who
shall

see it?

They shall go down to the bars of the pit,
when \emph{our} rest

together \emph{is} in the dust.

\emph{Response} 7.

\emph{Thy dead men shall
live, together with my dead body shall they
arise}.

Awake and sing, ye that
dwell in dust, for thy dew is as

the dew of herbs, and the earth shall cast out the dead.

\emph{Together with my dead body shall they come}.

\emph{Lesson} 8.

\emph{Job} xix. 20-27.

My bone cleaveth to my skin and to my flesh,
and I am

escaped with the skin of my teeth.

Have pity upon me, have pity upon me, O ye my

friends; for God hath touched me.

Why do ye persecute me as God, and are not
satisfied

with my flesh? 

Oh that my words were now written, oh that
they were

printed in a book!

That they were graven with an iron pen and
lead, in the

rock for ever!

For I know \emph{that} my Redeemer liveth,
and \emph{that} he shall

stand at the latter \emph{day} upon the earth:

And \emph{though} after my skin \emph{worms}
destroy this \emph{body}, yet in

my flesh shall I see God:

Whom I shall see for myself, and mine eyes
shall behold,

and not another; \emph{though} my reins be consumed within me.

\emph{Response} 8.

\emph{If Thou, Lord, wilt be
extreme to mark what is done amiss,
O Lord, who may abide it? but there is mercy with
Thee},

Wash me throughly from my wickedness, and
cleanse

me from my sin.

\emph{But there is mercy with Thee}.

\emph{Lesson} 9.

\emph{Jonah} ii. 2-10.

I cried by reason of mine affliction unto the
Lord, and he

heard me; out of the belly of hell cried I, \emph{and} thou
heard-

est my voice.

For thou hadst cast me into the deep, in the
midst of the

seas; and the floods compassed me about: all thy billows

and thy waves passed over me.

Then I said, I am cast out of thy sight; yet
I will look

again toward thy holy temple.

The waters compassed me about, \emph{even} to
the soul; the

depth closed me round about, the weeds were wrapt about

my head.

I went down to the bottoms of the mountains:
the earth

with her bars \emph{was} about me for ever: yet hast thou
brought

up my life from corruption, O Lord my God.

When my soul fainted within me, I remembered
the Lord,

and my prayer came in unto thee, into thine holy temple.

They that observe lying vanities, forsake
their own mercy.

But I will sacrifice unto thee with the voice
of thanks-

giving; I will pay \emph{that} that I have vowed: salvation \emph{is}
of

the Lord.

And the Lord spake unto the fish, and it
vomited out

Jonah upon the dry land.

\emph{Response} 9.

\emph{I am the Resurrection
and the Life, he that believeth in me,
though he were dead, yet shall he live}.

And whosoever liveth and believeth in Me
shall never die.

\emph{Though he were dead,
yet shall he live}. 

With Thee is the well of life, \&c.

\emph{Though he were dead,
yet shall he live}.

———————

\xsubsection{LAUDS}

\emph{Antiphon}

Cast me not away.

\emph{Psalm} 51. (1)

Have mercy upon me, \&c. With Thee,
\&c.

\emph{Antiphon}

Cast me not away from thy presence, and take
not Thy

Holy Spirit from me.

\emph{Antiphon}

Have I not remembered Thee?

\emph{Psalm} 63. (2)

O God, Thou art my God, \&c. With Thee,
\&c.

\emph{Antiphon}

Have I not remembered Thee in say bed, and
thought

upon Thee when I was waking?

\emph{Antiphon}

I am like a green olive tree.

\emph{Psalm} 52. (3)

O God, thou art my God, \&c. With Thee,
\&c.

I am like a green olive tree in the house of God; my

trust is in the tender mercy of God for ever and ever.

\emph{Antiphon}

The living, the living, he shall praise Thee.

\emph{Song of Ezekias} (4)

\emph{Is}. xxxviii. 10-20.

I said, in the cutting off of my days I shall
go to the gates

of the grave: I am deprived of the residue of my years, \&c.

I said, I shall not see the Lord, even the
Lord in the

land of the living: I shall behold man no more with the

inhabitants of the world.

Mine age is departed, and is removed from me
as a shep-

herd's tent: I have cut off like a weaver my life: he will

cut me off with pining sickness: from day \emph{even} to night
wilt

thou make an end of me.

I reckoned till morning, \emph{that} as a
lion so will he break

all my bones: from day even to night wilt thou make an

end of me.

Like a crane \emph{or} a swallow, so did I
chatter: I did mourn

as a dove; mine eyes fail \emph{with looking} upward: O Lord, I

am oppressed; undertake for me.

What shall I say? he hath both spoken unto
me, and

himself hath done \emph{it}: I shall go softly all my years in
the

bitterness of my soul.

O Lord, by these \emph{things men} live, and
in all these \emph{things is} the life of my spirit: so wilt thou recover me, and make

me to live. 

Behold, for peace I had great bitterness: but
thou hast

in love to my soul \emph{delivered it} from the pit of
corruption:

for thou hast cast all my sins behind thy back.

For the grave cannot praise thee, death can \emph{not}
celebrate

thee: they that go down into the pit cannot hope for thy

truth.

The living, the living, he shall praise thee,
as I do this

day: the father to the children shall make known thy

truth.

The Lord \emph{was ready} to save me:
therefore we will sing

my songs to the stringed instruments all the days of our

life, in the house of the Lord. With Thee, \&c.

\emph{Antiphon}

The living, the living, he shall praise Thee,
as I do this

day.

\emph{Antiphon}

Let every thing that hath breath.

\emph{Ps}. 148, 149, \&

150. (5)

O praise the Lord of heaven.

With Thee, \&c.

\emph{Antiphon}

Let every thing that hath breath praise the
Lord.

(v) \emph{Verse and Response}

I heard a voice from heaven, saying unto me,

\emph{Blessed are they that
die in the Lord}.

(w) \emph{Antiphon}

I am the Resurrection and the life.

\emph{Benedictus}

Blessed be the Lord God of Israel, \&c

With Thee, \&c.

\emph{Antiphon}

I am the Resurrection and the life; he that
believeth in

Me, though he were dead, yet shall he live, and whosoever

liveth and believeth in Me, shall never die.

\emph{The Lord's Prayer}

(\emph{privately}.)

Our Father, \&c.

\emph{Psalm} 126.

When the Lord turned again, \&c.

\emph{The Service terminates with the Collects
and Sentences used at Vespers}, p. 147, 148.

With Thee is the well, \&c.

May their souls rest in peace. \emph{Amen}.

\xsection{§ 8. Service for the Sundays in Advent.}

[Sections 2 and 3 form the respective bases
to this and the following section, the letters and numbers marking the
substitutions: the Hymns at Matins, Lauds, and Vespers are the only
portion of the Services not given. Christmas Day is here supposed to
fall on a Sunday. The Service for St. Thomas’s necessarily affects the
Lauds and dependent Hours on the 21st: with this exception no notice is
taken of Festivals, which materially alter the Advent Services as
actually performed in the Latin Church. It is instructive to compare the
Sunday Services as they here stand, with ours, which are formed from
them.]

———————

\xsubsection{MATINS}

O Lord, open, \&c.

\emph{And our mouth,
\&c}.

O God, make speed, \&c.

\emph{O Lord, make
haste, \&c}.

Glory be, \&c. Amen. Hallelujah.

(a) \emph{Invitatory}
\emph{Sundays} 1 \& 2. O
come, let us worship: the Lord our

King is coming.

\emph{Sundays} 3 \& 4. O come, let us
worship: the Lord is

now at hand.

(b) \emph{Hymn}

Verbum supernum prodiens, \&c.

———————

NOCTURN
I.

(c) \emph{Antiphon}
Behold, the mighty King: shall
come with great power

to save the nations. Hallelujah.

(d ) \emph{Antiphon}
Strengthen ye: the weak hands,
and confirm the feeble

knees; say to them of fearful heart, Behold your God will

come and save you. Hallelujah.

(e) \emph{Antiphon}
Be strong: fear not, behold
your God will come with

vengeance, even God with a recompense: He will come

and save you. (f) \emph{Verse and Response}
Out of Sion hath God appeared

in perfect beauty.

SUNDAY
1.

\emph{Lesson} 1.

\emph{Isa}. i. 1-3.
The vision of Isaiah the son of
Amoz, which he saw con-

cerning Judah and Jerusalem in the days of Uzziah, Jotham,

Ahaz, and Hezekiah, kings of Judah.

Hear, O heavens, and give ear, O earth: for the Lord hath

spoken, I have nourished and brought up children, and

they have rebelled against me.

The ox knoweth his owner, and the ass his master’s crib:

but Israel doth not know, my people doth not consider.

\emph{Response} 1.
\emph{I look afar off, and behold, I see the power of God coming, and a cloud covering the whole earth: Go ye forth to meet him, and say: Tell us, whether Thou be He: who shall rule over the people
Israel}.

All ye children of the earth and sons of men, the rich and

the poor together.

\emph{Go ye forth to meet him, and say: Tell us, whether Thou be He: who shall rule over the people
Israel}.

Hear, O Thou Shepherd of Israel, Thou that leadest

Joseph like a sheep.

\emph{Tell us, whether Thou be He: who shall rule over the people
Israel}.

Lift up your heads, O ye gates, and be ye lift up, ye ever-

lasting doors, and the King of glory shall come in.

\emph{Who shall rule over the people
Israel}.

Glory be to the Father, \&c.

\emph{I look afar off, and behold, I see the power of God, \&c}.

\emph{Lesson} 2.

\emph{Isa}. i. 4-6.
Ah sinful nation, a people laden with iniquity, a seed of

evil doers, children that are corrupters: they have forsaken

the Lord, they have provoked the Holy One of Israel unto

anger, they are gone away backward.

Why should ye be stricken any more? ye will revolt more

and more: the whole head is sick, and the whole heart

faint.

From the sole of the foot even unto the head \emph{there is} no

soundness in it; \emph{but} wounds, and bruises, and putrifying

sores: they have not been closed, neither bound up, neither

mollified with ointment. \emph{Response} 2.
\emph{I saw in
the night visions, and behold, One like unto the Son of
man come with the clouds of heaven, and there was
given him dominion, and glory, and a kingdom: that
all people, nations, and languages should serve Him}.

His dominion is an everlasting dominion,
which shall

not pass away, and His kingdom that which shall not be

destroyed.

\emph{That all people,
nations, \&c}.

\emph{Lesson} 3.

\emph{Isa}. i. 7-9.
Your country \emph{is}
desolate, your cities are burned with fire:

your land, strangers devour it in your presence, and \emph{it is}

desolate, as overthrown by strangers.

And the daughter of Zion is left as a cottage
in a vine-

yard, as a lodge in a garden of cucumbers, as a besieged

city.

Except the Lord of hosts had left unto us a
very small

remnant, we should have been as Sodom, \emph{and} we should

have been like unto Gomorrah.

\emph{Response} 3.
\emph{The
angel Gabriel was sent to Mary, a virgin espoused to}\\
\emph{Joseph,
announcing to her the Word, and the virgin was}\\
\emph{troubled
at the light: Fear not, Mary, thou hast found}\\
\emph{favour
with the Lord, and behold thou shall conceive and}\\
\emph{bring
forth a son, and He shall be called the Son of the}\\
\emph{Highest}.

The Lord God shall give unto Him the Throne
of His

father David, and He shall reign over the house of Jacob

for ever.

\emph{Behold,
thou shalt conceive, \&c}.

Glory be to the Father,
\&c.

\emph{Behold,
thou shalt conceive, \&c}.

SUNDAY
2.

\emph{Lesson} 1.

\emph{Isa}. xi. 1-4.
And there shall come forth a
rod out of the stem of Jesse,

and a Branch shall grow out of his roots:

And the Spirit of the Lord shall rest upon
him, the spirit

of wisdom and understanding, the spirit of counsel and

might, the spirit of knowledge and of the fear of the Lord;

And shall make him of quick understanding in
the fear

of the Lord: and he shall not judge after the sight of his

eyes, neither reprove after the hearing of his ears:

But with righteousness shall he judge the
poor, and re-

prove with equity for the meek of the earth: and he shall

smite the earth with the rod of his mouth, and with the

breath of his lips shall he slay the wicked. \emph{Response}.
\emph{O
Jerusalem, thy salvation shall come quickly: wherefore}\\
\emph{art
thou consumed with grief? is there any counselor}\\
\emph{in
thee? because sorrow hath changed thee:}\\
\emph{I will save thee
and deliver thee: fear not}.

For I am the Lord thy God; the Holy One of
Israel, thy

Redeemer.

\emph{I will save thee and
deliver thee: fear not}.

\emph{Lesson} 2.

\emph{Is}. xi. 4-7.
And he shall smite the earth
with the rod of his mouth,

and with the breath of his lips shall he slay the wicked.

And righteousness shall be the girdle of his
loins, and

faithfulness the girdle of his reins.

The wolf also shall dwell with the lamb, and
the leopard

shall lie down with the kid; and the calf and the young lion

and the fatling together; and a little child shall lead them.

And the cow and the bear shalt feed; their
young ones

shall lie down together: and the lion shall eat straw like

the ox.

\emph{Response} 2.
\emph{The
Lord, my God, shall came, and all the saints with}\\
\emph{Thee,
and in that day there shall be great light: and}\\
\emph{it
shall be in that day that living waters shall go out}\\
\emph{from
Jerusalem, and the Lord shall be king for ever:}\\
\emph{over
all the earth}.

Behold, the Lord shall come with power, and
His kingdom

shall be in His hand, and dominion and sovereignty,

\emph{Over
all the earth}.

\emph{Lesson} 3.

\emph{Is}. xi. 8-10.
And the sucking child shall
play on the hole of the asp,

and the weaned child shall put his hand on the cockatrice’

den.

They shall not hurt nor destroy in all my
holy mountain:

for the earth shall be full of the knowledge of the Lord, as

the waters cover the sea.

And in that day there shall be a root of
Jesse, which shall

stand for an ensign of the people; to it shall the Gentiles

seek: and his rest shall be glorious.

\emph{Response} 3.
\emph{City of
Jerusalem, weep thou not, for the Lord hath}\\
\emph{sorrowed
for thee: and He shall take from thee all}\\
\emph{tribulation}.

Behold, the Lord shall come in strength, and
his arm

shall rule.

\emph{And He shall take,
\&c}.

Glory be to the Father, \&c.

\emph{And He shall take,
\&c}. (go on to (g))

SUNDAY
3.

\emph{Lesson} 1.

\emph{Is}. xxvi. 1-6.
In that day shall this song be
sung in the land of Judah;

We have a strong city; salvation will \emph{God} appoint \emph{for}
walls

and bulwarks.

Open ye the gates, that the righteous nation
which keepeth

the truth may enter in.

Thou wilt keep him in perfect peace, \emph{whose}
mind \emph{is} stayed

\emph{on thee}: because he trusteth in thee.

Trust ye in the Lord for ever: for in the
Lord Jehovah \emph{is}

everlasting strength.

For he bringeth down them that dwell on high:
the lofty city,

he layeth it low; he layeth it low, even to the ground: he

bringeth it \emph{even} to the dust.

The foot shall tread it down, \emph{even} the
feet of the poor, \emph{and}

the steps of the needy.

\emph{Response} 1.
\emph{Behold,
the Lord shall appear on a white cloud: and}\\
\emph{with
him the thousands of His saints; and he shall}\\
\emph{have
written on his garment and on his thigh, King}\\
\emph{of
kings, and Lord of lords}.

He shall appear in the end and shall not lie:
if He tarry,

wait for Him; for He shall surely come:

\emph{And with Him the
thousands, \&c}.

\emph{Lesson} 2.

\emph{Is}. xxvi. 7-10.

The way of the just \emph{is} uprightness;
thou, most upright,

dost weigh the path of the just.

Yes, in the way of thy judgments, O Lord,
have we

waited for thee: the desire of \emph{our} soul \emph{is} to thy
name, and

to the remembrance of thee.

With my soul have I desired thee in the
night; yea, with

my spirit within me will I seek thee early: for when thy

judgments \emph{are} in the earth, the inhabitants of the world

will learn righteousness.

Let favour be showed to the wicked, \emph{yet}
will he not learn

righteousness: in the land of uprightness will he deal un-

justly, and will not behold the majesty of the Lord.

\emph{Response} 2.
\emph{O
Bethlehem, city of the Most High God, out of thee}\\
\emph{shall
go forth the Ruler of Israel, whose goings forth}\\
\emph{have
been from everlasting, and He shall be great in}\\
\emph{the
midst of the whole earth: and there shall be peace}\\
\emph{in
our land, when He shall come}.

He shall speak peace to the Gentiles, and His
dominion

shall be from sea to sea.

\emph{And there shall be
peace, \&c}.

\emph{Lesson} 3.

\emph{Is}. xxvi. 11-14.
Lord, \emph{when} thy hand is
lifted up, they will not see: \emph{but}

they shall see, and be ashamed for \emph{their} envy at the
people;

yea, the fire of thine enemies shall devour them. 

Lord, thou wilt ordain peace for us: for thou
also hast

wrought all our works in us.

O Lord our God, \emph{other} lords beside
thee have had domi-

nion over us: \emph{but} by thee only will we make mention of
thy

name.

\emph{They} are dead, they shall not live; \emph{they
are} deceased, they

shall not rise: therefore hast thou visited and destroyed them,

and made all their memory to perish.

\emph{Response} 3.
\emph{He that
shall come will come, and will not tarry; and}\\
\emph{fear
shall no more be in our borders; for He is our}\\
\emph{Saviour}.

He will subdue our iniquities; and Thou wilt
cast all their

sins into the depths of the sea.

\emph{For
He is our Saviour}.

Glory be to the Father,
\&c.

\emph{For
He is our Saviour}. (go on (g))

SUNDAY
4.

\emph{Lesson} 1.

\emph{Is}. xxxv. 1-7.
The wilderness and the
solitary place shall be glad for

them; and the desert shall rejoice, and blossom as the rose.

It shall blossom abundantly, and rejoice even
with joy and

singing: the glory of Lebanon shall be given unto it, the ex-

cellency of Carmel and Sharon, they shall see the glory of

the Lord, \emph{and} the excellency of our God.

Strengthen ye the weak hands, and confirm the
feeble

knees.

Say to them \emph{that are} of a fearful
heart, Be strong, fear not:

behold, your God will come \emph{with} vengeance, \emph{even}
God \emph{with}

a recompence; He will come and save you.

Then the eyes of the blind shall be opened,
and the ears

of the deaf shall be unstopped.

Then shall the lame \emph{man} leap as an
hart, and the tongue of

the dumb sing: for in the wilderness shall waters break out,

and streams in the desert.

And the parched ground shall become a pool,
and the

thirsty land springs of water.

\emph{Response} 1.
\emph{Blow the
trumpet in Sion, summon the nations, proclaim}\\
\emph{to
all people and say: Behold God our Saviour shall}\\
\emph{come}.

Proclaim and make it heard; speak and cry
aloud.

\emph{Behold, God our Saviour
shall come}.

\emph{Lesson} 2.

\emph{Is}. xxxv. 7-10.
In the habitation of dragons,
where each lay, \emph{shall be} grass

with reeds and rushes. 

And an highway shall be there, and a way, and
it shall

be called The way of holiness; the unclean shall not pass

over it; but it \emph{shall be} for those: the wayfaring men,
though

fools, shall not err therein.

No lion shall be there, nor \emph{any}
ravenous beast shall go up

thereon; it shall not be found there: but the redeemed shall

walk \emph{there}.

And the ransomed of the Lord shall return,
and come to

Zion with songs and everlasting joy upon their heads: they

shall obtain joy and gladness, and sorrow and sighing shall

flee away.

\emph{Response} 2.
\emph{The
sceptre shall not depart from Judah, nor a lawgiver from
between his feet, until Shiloh come: and unto Him shall
the gathering of the people be}.

His eyes are redder than wine, and His teeth
whiter than

milk.

\emph{And unto Him shall,
\&c}.

\emph{Lesson} 3.

\emph{Is}. xli. 1-4.
Keep silence before me, O
islands; and let the people

renew \emph{their} strength: let them come near; then let them

speak: let us come near together to judgment.

Who raised up the righteous man from the
east, called him

to his foot, gave the nations before him, and made \emph{him}
rule

over kings? he gave \emph{them} as the dust to his sword, \emph{and}
as

driven stubble to his bow.

He pursued them, \emph{and} passed safely; \emph{even}
by the way \emph{that}

he had not gone with his feet.

Who hath wrought and done \emph{it}, calling
the generations

from the beginning? I the Lord, the first, and with the

last; I \emph{am} he.

\emph{Response} 3.
\emph{I must
decrease, but He must increase; he that cometh}\\
\emph{after
me, was in being before me: whose shoes’ latchet}\\
\emph{I
am not worthy to unloose}.

I baptized you with water: but He shall
baptize you with

the Holy Ghost.

\emph{Whose
shoes’ latchet, \&c}.

Glory be to the Father, \&c.

\emph{Whose shoes’ latchet, \&c}. (go on
(g))

———————

NOCTURN
II.

(g) \emph{Antiphon}
Rejoice greatly: O daughter of
Jerusalem; behold thy

king cometh unto thee; O Sion, be not afraid, for thy

salvation cometh quickly.

(h) \emph{Antiphon}
Our king: shall come, even
Christ, whom John proclaimed

as the Lamb which was to come.

(i) \emph{Antiphon}
Behold, I come: quickly, and
My reward is with Me, to

give to every man according as his work shall be.

(j) \emph{Verse and Response}
Send ye the
Lamb to the Ruler of the land.

\emph{From
Sela to the wilderness, unto the mount of the}

\emph{daughter of Sion}.

SUNDAY
1.

\emph{Lesson} 4.

(\emph{From St. Leo}.)
Our Saviour, in His account of
the coming of the

kingdom of God, and of the end of the world, addressed to

His Apostles, and in them to the whole Church, bids them

beware, lest, at any time, their hearts should be weighed

down with surfeiting and drunkenness, and cares of this life:

which warning, beloved brethren, we know to belong to us

specially, to whom the threatened day, though hidden, is cer-

tainly near.

\emph{Response} 4.
\emph{Hail
thou that art highly favoured; the Lord is with}\\
\emph{thee:
the Holy Ghost shall come upon thee, and the}\\
\emph{power
of the Highest shalt overshadow thee; therefore}\\
\emph{that
Holy Thing which shall be born of Thee shall be}\\
\emph{called
the Son of God}.

How shall this be, seeing I know not a man?
and the

angel answered and said to her,

\emph{The
Holy Ghost shall come upon thee, \&c}.

\emph{Lesson} 5.
For whose coming it is fitting
that all mankind should

make ready, lest it should surprise any by gluttony or

worldly cares. For daily experience proves, my beloved,

that the keenness of the mind is blunted by fulness of drink,

and the vigour of the heart clouded by excess of meat; so

that love of eating even does injury to our bodily health,

unless a rule of temperance counteracts the seduction and

refuses to indulgence what afterwards would be a burden.

\emph{Response} 5.
\emph{We look
for the Saviour, the Lord Jesus Christ; who}\\
\emph{shall
change this vile body that it may be fashioned like}\\
\emph{unto
His glorious body}.

Let us live soberly, righteously, and godly
in this present 

world, looking for the blessed hope and the glorious appear-

ance of the great God.

\emph{Who
shall change, \&c}.

\emph{Lesson} 6.
For though the flesh hath no
desires apart from the soul,

and derives its senses from whence it gains animation, yet

it is the part of the soul to deny some things to that body

which is made subject to her, and by an inward discretion

to hold back what is outward from mischief; that, being

often at leisure from bodily desires, she may devote herself

to divine wisdom in the palace of the mind, and while the

noise of earthly cares is altogether hushed, may rejoice in

holy meditations and everlasting pleasures.

\emph{Response} 6.
\emph{O my
Lord, send, I pray Thee, by the hand of him whom}\\
\emph{Thou
wilt send: behold the affliction of Thy people: as}\\
\emph{Thou
hast spoken, come: and deliver us}.

Hear, O Thou Shepherd of Israel, Thou that
leadeth

Joseph like a sheep, Thou that sittest between the cherubim.

\emph{As Thou host spoken,
come: and deliver us}.

Glory
be to the Father, \&c.

\emph{And deliver us}. (go on (k))

SUNDAY
2.

(\emph{From St. Jerome}.)
And there shall go forth a rod
out of the root of Jesse. As

far as the beginning of the vision, or burden of Babylon,

which Isaiah the son of Amoz saw, all this prophecy relates

to Christ. Which I would explain by portions, lest if it be

all set forth, and treated at once, it confuse the memory of

the reader. The Jews understand our Lord by the rod and

flower from the rod of Jesse, the rod signifying his power,

the flower his comeliness.

\emph{Response} 4.
\emph{Behold,
the Lord shall come, our defender, the Holy One}\\
\emph{of
Israel: with the crown of His kingdom upon Him}.

And he shall reign from sea to sea, and from
the river unto

the ends of the earth.

\emph{With the crown, \&c}.

\emph{Lesson} 5.
But we, by the rod from the
root of Jesse, understand the

holy Virgin Mary, who had no shrub belonging to her: of

whom we read before, Behold, a Virgin shall conceive, and

shall bear a son. And by the flower is meant the Lord and

Saviour, who says in the Song of songs, I am the rose of

Sharon, and the lily of the valleys.

\emph{Response} 5.
\emph{As a
mother comforteth her sons, so will 1 comfort you,}\\
\emph{saith
the Lord: and from Jerusalem, the city I have}\\
\emph{chosen,
shall succour come to you; and you shall see,}\\
\emph{and
your heart shall rejoice}. 

I will give to Sion salvation, and to
Jerusalem My glory.

\emph{And
you shall see, \&c}.

\emph{Lesson} 6.
Therefore upon this flower,
which shall suddenly spring

up from the stem and root of Jesse by the Virgin Mary, the

Spirit of the Lord shall rest; for in Him it hath seemed

good that all the fulness of the Godhead should dwell bodily;

not by portions, as in the case of other saints, but, according

to the Hebrew gospel of the Nazarenes, all the fountain of

the Holy Spirit shall descend upon Him. But the Lord is that

Spirit; and where the Spirit of the Lord is, there is liberty.

\emph{Response} 6.
\emph{O
Jerusalem, thou shall plant thy vine upon thy moun-}\\
\emph{tains.
Leap for joy, for the day of the Lord shall}\\
\emph{come.
Arise, O Sion, be turned unto the Lord thy}\\
\emph{God;
rejoice and be glad, O Jacob: for thy Saviour}\\
\emph{shall
come in the midst of the nations}.

Rejoice greatly, O daughter of Sion: shout, O
daughter

of Jerusalem.

\emph{For thy Saviour shall come in the midst of
the nations}.

Glory
be to the Father, \&c.

\emph{For thy Saviour shall come, \&c}.
(go on (k))

SUNDAY
3.

\emph{Lesson} 4.

(\emph{From St. Leo}.)
The time of the year, and our
own religious custom, leads

me to proclaim to you, dearly beloved, with the anxiety of a

pastor, the fast of the 10th month, in which for the comple-

tion of the gathering of all fruits, is poured out most suitably

to the Lord their giver, the tribute of abstaining from them;

for what can be more available than fasting? by observing

which we draw near to God, and resisting the devil, over-

come our pleasant vices.

\emph{Response} 4.
\emph{Weep
not, O Egypt, for thy Ruler shall come to thee,}\\
\emph{before
whose coming the depths shall be moved: to}\\
\emph{deliver
His people from the hand of power}.

Behold the Lord of hosts shall come, thy God
with great

power.

\emph{To deliver, \&c}.

\emph{Lesson} 5.
For fasting has ever been the
food of virtue. In a word,

from abstinence proceed chaste thoughts, moderate wishes,

sober purposes: and by voluntary inflictions, the desires of

the flesh are mortified, the graces of the spirit strengthened.

But since, not by fasting only, the health of our souls is

gained, let us make almsgiving the supplement of fasting.

Let us give to active virtue, what we diminish from indul- 

gence: Let the abstinence of the faster be the refreshment

of the poor.

\emph{Response} 5.
\emph{Her time
is well nigh come, and her days shall not be}\\
\emph{prolonged:
the Lord shall have pity upon Jacob, and}\\
\emph{Israel
shall be saved}.

Return, O virgin of Israel, return into thy
cities.

\emph{The Lord shall have
pity, \&c}.

\emph{Lesson} 6.
Let us give attention to the
defence of widows, the wel-

fare of orphans, the consolation of mourners, the reconcilia-

tion of enemies. Befriend the foreigner, assist the oppressed,

clothe the naked, nurse the sick; that whoso of us by

righteous labours shall offer to God, the author of all good,

such pious sacrifice, may receive in turn for the same the

reward of His heavenly kingdom. Let us keep fast on

Wednesdays and Fridays; let us keep Vigil on the Sabbath

with the blessed Apostle Peter, whose deeds of grace work-

ing together with our prayers, we may obtain what we seek

through our Lord Jesus Christ, who with the Father and the

Holy Ghost, liveth and reigneth, world without end. \emph{Amen}.

\emph{Response} 6.
\emph{The Lord
shall descend as the rain into a fleece of wool:}\\
\emph{righteousness
shall spring forth in His days, yea, and}\\
\emph{abundance
of peace}.

And all kings shall worship Him, all the
heathen shall

serve Him.

\emph{Righteousness
shall spring forth, \&c}.

Glory be to the Father,
\&c.

\emph{Righteousness
shall spring forth, \&c}. (go on (k))

SUNDAY
4.

\emph{Lesson} 4.

(\emph{From St. Leo}.)
If we understand truly and
wisely, dearly beloved, the

history of our creation, we shall find that man was therefore

formed after the likeness of God, that he might imitate the

Creator; and that the natural dignity of our race consists in

the image of divine kindness shining within us as in a

mirror. Towards which we are daily renewed by the grace

of the Saviour, so that what fell in the first Adam, is raised

again in the second.

\emph{Response} 4.
\emph{Unto us
a child is born, and He shall be called the mighty
God: He shall sit upon the throne of David His
Father and shall reign: and the dominion shall be upon
His shoulder}. 

And in Him shall all tribes of the earth be
blessed, and

all nations shall serve Him.

\emph{He shall sit, \&c}.

\emph{Lesson} 5.
The sole cause of our recovery
is the pitifulness of God,

whom we should not love, unless He first loved us, and

scattered the darkness of our ignorance with the light of

His truth; which the Lord declaring by holy Isaiah saith,

I will bring the blind by a way that they know not, and

paths of which they were ignorant will I make them tread;

I will make their darkness light, and their crooked straight.

This will I do to them, and not forsake them. And again,

I am found of them who sought me not, and I am manifested

to them that asked not for me.

\emph{Response} 5.
\emph{Behold the fulness of time
is now come, in which God sent}\\
\emph{forth His Son upon
the earth, born of a virgin, made under}\\
\emph{the Law: to redeem
those that are under the Law}.

For His great love wherewith He loved us, God
sent forth

His Son in the likeness of sinful flesh.

\emph{To redeem those, \&c}.

\emph{Lesson} 6.
How this was fulfilled the
Apostle John teaches, saying,

We know that the Son of God is come, and has given us an

understanding that we may know him that is true, and we

are in His true Son. And again, We love God, because He

first loved us. Therefore God, by loving us, restores in us

His own image, and that He may find in us the likeness of

His own goodness, He giveth whence we also may work

what He worketh, by lighting the lamps of our minds, and

kindling in us the flame of His own love, that we may love

not Him only, but all that He loves.

\emph{Response} 6.
\emph{O Virgin
of Israel, return to thy cities; how long wilt}\\
\emph{thou
be put aside and be in sorrow? thou shalt con-}\\
\emph{ceive
the Lord and Saviour, a new oblation in the earth:}\\
\emph{men
shall go forward to salvation}.

I have loved thee with an everlasting love;
therefore with

loving kindness have I drawn thee.

\emph{How long wilt thou be
put aside and be in sorrow? Thou}\\
\emph{shalt
conceive thy Lord and Saviour, a new oblation in}\\
\emph{the
earth: men shall go forward to salvation}.

Glory be to the Father, \&c.

\emph{Men shall go forward to
salvation, \&c}.

———————

NOCTURN
III.

(k) \emph{Antiphon}
The Angel Gabriel: spake unto
Mary, saying, Hail thou

that art highly favoured, the Lord is with thee. Blessed art

thou among women.

(l) \emph{Antiphon}
Mary said: What meaneth this
salutation ? for my soul

is troubled, and I shall bear my King, and yet remain a

virgin.

(m) \emph{Antiphon}
Before the coming: of the
great King, let the hearts of

men be cleansed, to walk worthily to meet Him, for behold

He will come, and will not tarry.

(n) \emph{Verse and Response}
The Lord will come forth of
his holy place.

He will come to save His people.

SUNDAY
1.

\emph{Lesson} 7.

\emph{Luke} xxi. 25-33.
\emph{At that time Jesus said to
his disciples}.

There shall be signs in the sun, and in the
moon, and in

the stars; and upon the earth distress of nations, with per-

plexity; the sea and the waves roaring;

Men’s hearts failing them for fear, and for
looking after

those things which are coming on the earth: for the powers

of heaven shall be shaken.

And then shall they see the Son of man coming
in a cloud

with power and great glory.

And when these things begin to come to pass,
then look

up, and lift up your heads; for your redemption draweth

nigh.

And he spake to them a parable: Behold the
fig-tree, and

all the trees;

When they now shoot forth, ye see and know of
your own

selves that summer is now nigh at hand.

So likewise ye, when ye see these things come
to pass,

know ye that the kingdom of God is nigh at hand.

Verily I say unto you, This generation shall
not pass

away, till all be fulfilled.

Heaven and earth shall pass away: but my
words shall

not pass away.

(\emph{Homily of Pope Gregory}.)
Our Lord and Redeemer,
desirous of finding us ready,

foretels what evils shall attend the ageing world, in order to

sober us from the love of it. He makes known how great

strokes shall precede its approaching end; that, if we will

not fear God in prosperity, at least when afflicted by His

strokes, we may dread His near judgment. \emph{Response} 7.

\emph{Behold a virgin shall conceive and bear a Son, saith}\\
\emph{the Lord, and His
name shall be called Wonderful, the}\\
\emph{mighty God}.

Upon the throne of David and over his kingdom
shall

He reign for ever.

\emph{And his name shall, \&c}.

\emph{Lesson} 8.
Shortly before the passage of
the holy Gospel, which my

brethren have been hearing, the Lord had said, Nation

shall rise against nation, and kingdom against kingdom:

and great earthquakes shall be in divers places, and pesti-

lences and famines. Then after some additions, He says

what you have heard: There shall be signs in the sun, and

in the moon, and in the stars; and upon the earth distress

of nations with perplexity; the sea and the waves roaring.

Of all these signs certainly some we see fulfilled already,

others we fear as soon coming.

\emph{Response} 8.
\emph{Hear the
word of the Lord, O ye nations, and declare it}\\
\emph{in
the ends of the earth; and say to the isles afar off,}\\
\emph{Our
Saviour shall come}.

Declare it, and make it heard; speak, and cry
out.

\emph{And say to the isles,
\&c}.

\emph{Lesson} 9.
For nation is rising against
nation, and distress from

them presses upon the countries, more in the events we see

than in the books we read. You know too how often we

hear from other parts of the world, of earthquakes over-

whelming cities; pestilences we suffer without respite.

We do not yet openly behold signs in the sun, moon, and

stars: but the alteration of the atmosphere betokens that

they are not far off.

\emph{Response} 9.
\emph{Behold
the days come, saith the Lord, that I will raise}\\
\emph{up
unto David a righteous Branch, and a King shall}\\
\emph{reign
and prosper, and shall execute judgment and}\\
\emph{justice
in the earth: and this is His name whereby He}\\
\emph{shall
be called: The Lord our Righteousness}.

In his days Judah shall be saved and Israel
shall dwell

safely.

\emph{And this is His name
whereby He shall be called: The}\\
\emph{Lord
our Righteousness}.

Glory be to the Father, \&c.

\emph{The Lord is our
Righteousness}.

(\emph{End of Matins}. \emph{The Te Deum is not sung in Advent}.)

SUNDAY
2.

\emph{Lesson} 7.

\emph{Matt}. xi. 2-9.

Now when John had heard in the prison the
works of

Christ, he sent two of his disciples.

And said unto Him, Art thou he that should
come, or do

we look for another?

Jesus answered and said unto them, Go and
show John

again those things which ye do hear and see:

The blind receive their sight, and the lame
walk, the

lepers are cleansed, and the deaf hear, the dead are raised

up, and the poor have the gospel preached to them.

And blessed is \emph{he}, whosoever shall not
be offended in me.

And as they departed, Jesus began to say unto
the mul-

titudes concerning John, What went ye out into the wil-

derness to see? A reed shaken with the wind?

But what went ye out for to see? A man
clothed in soft

raiment? behold, they that wear soft \emph{clothing} are in
king’s

houses.

But what went ye out for to see? A prophet?
yea, I say

unto you, and more than a prophet.

(\emph{Homily of Pope Gregory}.)

So many signs and so many miracles, were to
earls be-

holder not an offence but a wonder. Yet the mind of un-

believers took this serious offence at Him, that after all

His miracles He should be seen to die. Whence also Paul

saith, We preach Christ crucified, to the Jews a stumbling-

block, and to the Gentiles foolishness. For it seemed to

man to be foolish, that for man the Source of life should

die; and thus man took offence at Him for the very thing,

whence he ought to have felt himself the more His debtor.

For God is so much the more worthy of honour from men,

by how much He has for men undergone dishonour.

\emph{Response} 7.

\emph{The Lord shall come
from Samaria to the Eastern Gate,}

\emph{and
shall come to Bethlehem, walking upon the waters}

\emph{of
the redemption of Judah: then shall every man be}

\emph{saved,
for behold He cometh}.

And his throne shall be prepared in mercy,
and He shall

sit on it in truth.

\emph{Then
shall every man, \&c}.

\emph{Lesson} 8.

What means His saying, that He is blessed
whosoever is

not offended in Him, but to mark, in plain words, the abject

and degrading nature of His death? as if He openly said,

I do what is marvellous, yet I submit to suffer what is abject.

Since then I shortly follow thee in dying, men must specially

beware, lest, though venerating My signs, they despise My

death. \emph{Response} 8.

\emph{Make haste and tarry
not, O Lord: and deliver thy}

\emph{people}.

Come, O Lord, and do not tarry; absolve Thy
people

from their iniquities.

\emph{And
deliver Thy people}.

\emph{Lesson} 9.

But let us hear what he said to the crowd
concerning

this same John, when He had dismissed John’s disciples:

What went ye out into the wilderness for to see? a reed

shaken with the wind? which he means to be answered in

negative. For as the wind takes it, a reed moves this way

or that. And what is the reed, but the carnal mind? which,

according as it falls in with popularity or reproach, at once

inclines to the one side or the other.

\emph{Response} 9.

\emph{Behold the Lord shall
come down with glory and His}

\emph{power
with him: to visit His people in peace, and to}

\emph{establish
upon them everlasting life}.

Behold our Lord shall come with power.

\emph{To visit His people,
\&c}.

Glory be to the Father, \&c.

\emph{To visit His people,
\&c}. (\emph{End of Matins}.)

SUNDAY
3.

\emph{Lesson} 7.

\emph{John} i. 19-28.

And this is the record of John, when the Jews
sent priests

and Levites from Jerusalem to ask him, Who art thou

And He confessed, and denied not but
confessed, I am

not the Christ.

And they asked him, What then? art thou
Elias? And

he saith, I am not. Art thou that prophet? And he

answered, No.

Then said they unto him, Who art thou? that
we may

give an answer to them that sent us. What sayest thou of

thyself?

He said, I am the voice of one crying in the
wilderness,

Make straight the way of the Lord, as said the prophet

Esaias.

And they which were sent were of the
Pharisees.

And they asked him, and said unto him, Why
baptizest

thou then, if thou be not that Christ, nor Elias, neither

that prophet?

John answered them, saying, I baptize with
water: but

there standeth one among you, whom ye know not;

He it is, who coming after me is preferred
before me,

whose shoes latchet I am not worthy to unloose. 

These things were done in Bethabara beyond
Jordan,

where John was baptizing.

(\emph{Homily of Pope Gregory}.)

From the words of this lesson, dearest
brethren, we have

proof of John’s humility: who, whereas he had such grace

that he was taken for the Christ, yet preferred to remain

unmoveably in what he was, lest he should be carried

away by human opinion into fancies above what he was.

For he confessed and denied not, but confessed I am not

the Christ. But in saying, I am not, he denied what he was

not, not what he was: that by speaking truth he might

become a member of Him, whose name he did not fraudu-

lently appropriate. While then he is not ambitious of the

name, he becomes a member of Christ; because in aiming

at humbly confessing his own weakness, he was found

worthy of truly possessing His majesty.

\emph{Response} 7.

\emph{Come, O Lord, and do
not tarry; absolve Thy people}

\emph{from
their iniquities: and bring back the scattered}

\emph{ones
into their own land}.

Stir up, O Lord, Thy power, and come, and
save us.

\emph{And bring back, \&c}.

\emph{Lesson} 8.

But when we recollect our Saviour’s saying
in another

passage, the words of this passage raise a perplexing ques-

tion. For elsewhere, when the Lord was asked by His

disciples concerning the coming of Elias, He answered,

Elias is already come, and they have done unto him what-

soever they listed, and, if ye will k now, John himself is

Elias. Yet John, when he was asked, said, I am not Elias.

What means it, dearly beloved brethren, that what the

truth affirms, the prophet of the truth denies? The two

statements are very different; It is he, and, I am not he.

How then is he the prophet of the truth, if he agrees not

with the sayings of that truth itself?

\emph{Response} 8.

\emph{Behold, the root of
Jesse shall come for the salvation of}

\emph{the
people; unto Him shall the Gentiles seek: and His}

\emph{name
shall be glorious}.

And the Lord God shall give unto him the
throne of His

father David, and He shall reign over the house of Jacob

for ever.

\emph{And his name, \&c}.

\emph{Lesson} 9.

But if we inspect the matter minutely, what
seems in-

consistent, will be found consistent. For the angel said to

Zacharias concerning John, He shall go before Him in the

spirit and power of Elias. He then is said to be coming 

in the spirit and power of Elias, inasmuch as he went

before the first coming of the Lord, as Elias will go before

His second coming. As he will be the forerunner of the

Judge, so John was the forerunner of the Redeemer. John

then was in the spirit of Elias, not in his person. What

then the Lord declares of the spirit, John denies of the

person.

\emph{Response} 9.

\emph{The Lord shall teach us
His ways, and we will walk in}

\emph{His
paths: for out of Sion shall go forth the law,}

\emph{and
the word of the Lord from Jerusalem}.

Come ye, and let us go up into the mountain
of the Lord,

and to the house of the God of Jacob.

\emph{For out of Sion, \&c}.

Glory
be to the Father, \&c.

\emph{For out of Sion, \&c}.
(\emph{End of Matins}.)

SUNDAY
4.

\emph{Lesson} 7.

\emph{Luke} iii. 1-9.

Now in the fifteenth year of the reign of
Tiberius Cæsar,

Pontius Pilate being governor of Judea, and Herod being

tetrarch of Galilee, and his brother Philip tetrarch of Iturea

and of the regions of Trachonitis, and Lysanias the tetrarch

of Abilene,

Annas and Caiaphas being the high priests,
the word of

God came unto John the son of Zacharias in the wilderness.

And he came into all the country about
Jordan, preaching

the baptism of repentance for the remission of sins;

As it is written in the book of the words of
Esaias the

prophet, saving, The voice of one crying in the wilderness,

Prepare ye the way of the Lord, make His paths straight.

Every valley shall be filled, and every
mountain and hill

shall be brought low; and the crooked shall be made

straight, and the rough ways \emph{shall} be made smooth;

And all flesh shall see the salvation of God.

Then said he to the multitude that came forth
to be

baptized of him, O generation of vipers, who hath warned

you to flee from the wrath to come?

Bring forth therefore fruits worthy of
repentance, and

begin not to say within yourselves, We have Abraham to

\emph{our} father: for I say unto you, That God is able of these

stones to raise up children unto Abraham.

And now also the axe is laid unto the root of
the trees;

every tree therefore which bringeth not forth good fruit is

hewn shown, and cast into the fire. (\emph{Homily of Pope Gregory}.)

John said to the multitudes which went out to
be bap-

tized by him, O generation of vipers, who hath warned

you to flee from the wrath to come? For the wrath to

come is the punishment of uttermost vengeance; which

the sinner is not able then to flee, who has not now

recourse to the laments of penitence. And it is observa-

ble, that evil offspring, imitating the conduct of evil pa-

rents, are called generation of vipers: for in envying and

persecuting the good, repaying evil, injuring their neigh-

bours, in all these ways following the doings of their

carnal fathers, they are as though poisonous children of

poisonous parents.

\emph{Response} 7.

\emph{I have sworn, saith the
Lord, that I will no longer be}

\emph{wroth
upon the earth; for the mountains and the hills}

\emph{shall
receive My righteousness: and My covenant of}

\emph{peace
shall be with Jerusalem}.

My salvation is near to come, and My
righteousness

to be revealed.

\emph{And My covenant, \&c}.

\emph{Lesson} 8.

But since we have already sinned, since we
are involved

in the usage of bad habits, let him say what we must do,

in order to flee from the wrath to come. It follows; Bring

forth therefore fruits, worthy of repentance. In these words

it is observable that the friend of the bridegroom enjoins

not only fruits, but worthy fruits of repentance. It is one

thing to bring forth the fruit, another worthy fruits of re-

pentence. By way of explaining what is meant by worthy,

it must be borne in mind, that whoso has committed

nothing unlawful, he may lawfully use what is lawful; and

do his works of charity without necessarily relinquishing

the world.

\emph{Response} 8.

\emph{We will not depart from
Thee, O Lord, Thou wilt quicken}

\emph{us,
and we will call upon Thy name: show us Thy}

\emph{countenance,
and we shall be whole}.

Remember me, O Lord, according to the favour
that

Thou bearest unto Thy people: O visit me with Thy sal-

vation.

\emph{Show us Thy countenance},
\&c.

\emph{Lesson} 9.

But if any one has incurred the guilt of
fornication, or

perchance, what is still more heinous, adultery, be so much

more should he deprive himself of things lawful, as he

remembers he has committed what is unlawful. Nor should

the fruit of a good work be the same, where a man has 

sinned less, or sinned more, been betrayed into none, or

into some crimes, or into many. In saying then “worthy

fruits of penance,” the conscience of every one is suitably

addressed, that he may obtain by so much the more advan-

tage of good works by penance, as he has incurred heavier

loss by sinning.

\emph{Response} 9.

\emph{Behold ye, how great is
He, who comes to save the}

\emph{nations.
He is the king of righteousness; of whose}

\emph{generation
there is no end}.

The Forerunner is for us entered, made for
ever an high

priest after the order of Melchizadek.

\emph{Of
whose generation, \&c}.

Glory be to the Father,
\&c.

\emph{Of
whose generation, \&c}.

———————

\xsubsection{LAUDS}

(o) \emph{Antiphon}

(\emph{also at prime on the respec- tive Sundays and the rest of each week till Dec}.

17.)

S\textsc{unday} 1.

In that day the mountain shall drop down new
wine,

and the hills shall flow with milk. Hallelujah.

SUNDAY
2.

Behold, the Lord shall come in the clouds of
heaven with

great power. Hallelujah.

SUNDAY
3.

The Lord will come and will not tarry, and
will bring to

light the hidden things of darkness, and manifest himself

to all nations. Hallelujah.

SUNDAY
4.

Blow the trumpet in Sion, for the day of the
Lord is at

hand; behold, he shall come and save us. Hallelujah.

Hallelujah.

(p) \emph{Antiphon}

(\emph{also at Third till Dec}. 17.)

SUNDAY
1.

Rejoice greatly, daughter of Sion: shout, O
daughter of

Jerusalem. Hallelujah.

SUNDAY
2.

We have a strong city; salvation will God
appoint for

walls and bulwarks. Open ye the gates, for God is with

us. Hallelujah.

SUNDAY
3.

Rejoice, O Jerusalem, with great joy, for the
Saviour

shall come unto thee. Hallelujah.

SUNDAY
4.

Behold, the Desire of all nations shall come,
and the house

of the Lord shall be filled with glory. Hallelujah.

(q) \emph{Antiphon}

(\emph{also at sixth, till Dec}. 17.)

SUNDAY
1.

Behold, the Lord shall come and all his
saints with Him,

and the light in that day shall be great. Hallelujah.

SUNDAY
2.

Behold, the Lord shall appear, and will not
lie; though

He tarry, wait for Him, because He will surely come, He

will not tarry. Hallelujah.

SUNDAY
3.

I will give salvation to Sion, and my glory
to Jerusalem.

Hallelujah.

SUNDAY
4.

The crooked shall be made straight, and the
rough places

plain; come, Lord, and do not tarry. Hallelujah.

(r) \emph{Antiphon}

SUNDAY
1.

Ho, every one that thirsteth, come ye to the
waters;

seek ye the Lord while he may be found. Hallelujah.

SUNDAY
2.

The mountains and the hills shall break forth
before him

into singing; and all the trees of the field shall clap their

hands: for the Lord the king shall come to his everlasting

kingdom. Hallelujah. Hallelujah.

SUNDAY
3.

Every mountain and hill shall be made low,
and the

crooked shall be made straight, and the rough places plain:

Come, O Lord, and do not tarry. Hallelujah.

SUNDAY
4.

The Lord shall come: meet ye Him, and say,
Great is

His beginning, and of His kingdom shall be no end. The

mighty God our king, the Prince of peace. Hallelujah.

Hallelujah.

(s) \emph{Antiphon}

(\emph{also at Ninth, till Dec}. 17.)

SUNDAY
1.

Behold the great Prophet shall come, and He
shall re-

build Jerusalem. Hallelujah.

SUNDAY
2.

Behold, our Lord shall come with power, and
shall en-

lighten the eyes of His servants. Hallelujah.

SUNDAY
3.

Let us live righteously and godly, waiting
for the blessed

hope and coming of the Lord.

SUNDAY
4.

Thy Almighty word leaped down from heaven out
of

Thy royal throne. Hallelujah.

(t) \emph{Text}

(\emph{also at Third, and Vespers on the respective Sundays}.)

SUNDAY
1.

Brethren, it is high time to wake out of
sleep, for now is

our salvation nearer than when we believed.

SUNDAY
2.

Brethren, whatsoever things were written
aforetime, were

written for our learning, that we through patience and

comfort of the Scriptures might have hope.

SUNDAY
3.

Brethren, rejoice in the Lord always; and
again I say,

Rejoice. Let your moderation be known unto all men. The

Lord is at hand.

SUNDAY
4.

Brethren, let a man so account of us, as
ministers of

Christ, and stewards of the mysteries of God. Moreover, it

is required of stewards, that a man be found faithful.

(u) \emph{Hymn}

En clara vox redarguit,
\&c.

(v) \emph{Verse and Response}

The voice of one crying in the wilderness;
Prepare ye

the way of the Lord.

\emph{Make his paths straight}.

(w) (\emph{Antiphon of the Benedictus}.)

SUNDAY
1.

The holy Ghost shall come upon thee, O Mary;
fear

not, thou shalt conceive in thy womb the Son of God.

Hallelujah.

SUNDAY
2.

When John had heard in the prison the works
of Christ,

he sent two of his disciples, and said unto Him, Art Thou

He that should come, or do we look for another?

SUNDAY
3.

Upon the throne of David and upon his
kingdom, He

shall sit for ever. Hallelujah.

SUNDAY
4.

Hail thou that art highly favoured; the Lord
is with

thee; blessed art thou among women. Hallelujah.

(x) \emph{Collect}

(\emph{also at Third, Sixth, Ninth, and Vespers on the respective Sundays}.)

SUNDAY
1.

Raise up, we beseech Thee, O Lord, Thy power,
and

come, that we being found meet may be snatched from

the perils of our sins by Thy succour, and may be saved

by Thy deliverance, who livest and reignest with God

the Father, \&c.

SUNDAY
2.

Stir up, O Lord, our heart, to prepare the
ways of Thy

only begotten Son, that by His coming we may be made

meet to serve Thee with purified minds, who liveth with

Thee, \&c.

SUNDAY
3.

Give ear, we beseech Thee, O Lord, to our
prayers, and

lighten the darkness of our mind by the grace of Thy

visitation, who livest and reignest, \&c.

SUNDAY
4.

Raise up, we pray Thee, Thy power, and come
among

us, and with great might succour us, that by the aid of Thy

grace, what our sins do hinder, the bounty of Thy prop-

tiation may forward, who livest, \&c.

———————

\xsubsection{PRIME}

(bb) \emph{Antiphon}

Same as (o).

(ee) \emph{Verse in Short Response}\\
(ii) \emph{Short Lesson},

\emph{Isa}. xxxiii. 2.

Thou that art now coming into the world.

O Lord, be gracious unto us; we have waited
for Thee;

be Thou our arm every morning, our salvation also in the

time of trouble.

———————

\xsubsection{THIRD HOUR}

(kk) \emph{Antiphon}

The same as (p).

(ll) \emph{Text}

The same as (t).

(mm) \emph{Short Response}

\emph{Come Thou and save us: O Lord God of
hosts.}

\emph{Come Thou, \&c}.
(\emph{Repeated}.)

Show the light of Thy countenance, and we
shall be whole.

\emph{O Lord God of hosts}.

Glory be to the Father, and to the Son, and
to the Holy

Ghost, \&c.

\emph{Come Thou and save us,
O Lord God of hosts}.

The heathen shall fear Thy name, O Lord.

\emph{And all the kings of
the earth Thy majesty}.

(nn) \emph{Collect}

The same as (x).

———————

\xsubsection{SIXTH HOUR}

(pp) \emph{Antiphon}
The same as (q).

(qq) \emph{Text}

SUNDAY
1.

The night is far spent,
the day is at hand, let us therefore

cast away the works of darkness, and put upon us the

armour of light.

SUNDAY
2.

Now the God of patience
and consolation grant you to be

like minded one toward another, according to Jesus Christ,

that ye may with one mind and one mouth glorify God,

even the Father of our Lord Jesus Christ.

SUNDAY
3.

Be careful for nothing;
but in every thing by prayer and

supplication, with thanksgiving, let your requests be made

known unto God.

SUNDAY
4.

But with me it is a very
small thing that I should be

judged of you, or of man’s judgment; yea I judge not

myself.

(rr) \emph{Short Response}

\emph{Show us Thy mercy, O Lord.}

\emph{Show us, \&c}.

And grant us Thy salvation,

\emph{O Lord}. 

Glory be, \&c.

\emph{Show us Thy mercy, O Lord}.

Remember me, O Lord, according to the favour
that

thou bearest unto Thy people.

\emph{O visit me with
Thy salvation}.

(ss) \emph{Collect}
The same as (x.)

———————

\xsubsection{NINTH HOUR}

(uu) \emph{Antiphon}
The same as (s).

(vv) \emph{Text}

SUNDAY
1.

Let us walk honestly as
in the day, not in rioting and

drunkenness, not in chambering and wantonness, not in

strife and envying: but put ye on the Lord Jesus Christ.

SUNDAY
2.

Now the God of hope fill
you with all joy and peace in

believing, that ye may abound in hope, through the power

of the Holy Ghost.

SUNDAY
3.

And the peace of God
which passeth all understanding,

shall keep your hearts and minds in Christ Jesus our Lord.

SUNDAY
4.

Therefore judge nothing
before the time, until the Lord

come, who both will bring to light the hidden things of

darkness and will make manifest the counsels of the heart,

and then shall every man have praise of God.

(ww) \emph{Short Response}
\emph{The
Lord shall arise upon thee: O Jerusalem}.

\emph{The Lord, \&c}.

And His glory shall be seen upon thee.

\emph{O Jerusalem}.

Glory be to, \&c.

\emph{The Lord shall arise upon thee: O Jerusalem}.

Turn us again, Lord God of hosts.

\emph{Show the light of Thy countenance, and we
shall be whole}.

(xx) \emph{Collect}
The same as (x).

———————

\xsubsection{VESPERS}

(yy)

(zz)

(aaa) \emph{Antiphon the}\\
(bbb) \emph{same as}

(ccc)
(0)

(p)

(q) \emph{respectively}.

(r)

(s)

(ddd) \emph{Text}
The same as (t).

(eec) \emph{Hymn}
Creator alme siderum, \&c.

(fff) \emph{Verse and Response}
Drop down, ye heavens, from
above, and let the skies

pour down righteousness.

\emph{Let the earth open, and let them bring
forth salvation}.

(ggg) \emph{Antiphon of the Magnif}.

SUNDAY
1.

Fear not, Mary, for thou
hast found favour with God:

and behold, thou shalt conceive, and shalt bear a son.

Hallelujah.

SUNDAY
2.

Art Thou He that should
come? or look we for another?

Tell John the things which ye see. The blind receive their

sight, the dead are raised, to the poor the Gospel is preached.

Hallelujah.

SUNDAY
3.

Blessed art thou, O
Mary, which believedst in the Lord:

there shall be a performance in thee of those things which

were said to thee of the Lord.

SUNDAY
4.

(\emph{being Dec}. 18.)

O Lord, and Ruler of the
House of Israel, who appear-

edst unto Moses in the flame of a burning bush, and

gavest to him the Law in Sinai, come to redeem us with a

stretched out arm.

(hhh) \emph{Collect}
The same as (x).

\xsection{§ 9. Service for Week Days in Advent.}

\xsubsection{MATINS}

(a) \emph{Invitatory}

(\emph{the same as the respective Sundays}.)

O Lord, open, \&c. Glory be, \&c.

1st \emph{and} 2d \emph{week}. O come, let us
worship: the Lord our

king approaching.

3d \emph{and} 4th \emph{week}. O come, let us
worship: the Lord is

close at hand.

(b) \emph{Hymn}

(\emph{all four weeks}.)

Verbum supernum prodiens, \&c.

\emph{Monday}
\emph{Tuesday}
\emph{Wednesday}
\emph{Thursday}
\emph{Friday}
\emph{Saturday}

\emph{Antiphon}
The Lord is

the
That I

offend not
O that the

Lord would
Haste Thee:

O Lord God
Sing we

merrily:
For the

Lord hath

done:

(1) \emph{Psalm}
27
39
53
69
81
98

(2) \emph{Psalm}
28
40
55
70
82
99

\emph{Antiphon}
strength: of

my life.
in my

tongue.
deliver: His

people out

of captivity.

to deliver

me.
unto God

our strength.
marvellous

things.

\emph{Antiphon}

Worship the

Lord:

Heal: my

soul, for

For my soul

Be thou:

Thou only,

art the most

O be joyful:

(3) \emph{Psalm}
29
41
56
71
83
100

(4) \emph{Psalm}
30
42
57
72
84
101

\emph{Antiphon}
with holy

worship.
I have sin-

ned against

Thee.

trusteth

in Thee.
my strong

hold.
highest over

all the earth.
the Lord all

ye lands.

\emph{Antiphon}

Deliver me:

in

My heart is

inditing

Do ye judge

the thing

Think upon

the tribe:

Lord, Thou

art become

Let my

crying:

(5) \emph{Psalm}
31
44
58
73
85
102

(6) \emph{Psalm}
32
45
59
74
86
103

\emph{Antiphon}
Thy righte-

ousness.
of a good

matter.
is right: O

ye sons of

men.

of Thine in-

heritance.
gracious:

unto Thy

land.
come unto

Thee.

\emph{Antiphon}

It becometh

well:

A very pre-

sent help:

O be Thou:

our

Unto Thee:

do we

Her founda-

tions: are

Praise the

Lord.

(7) \emph{Psalm}
33
46
60
75
87
104

(8) \emph{Psalm}
34
47
61
76
88
105

\emph{Antiphon}
the just to

be thankful.

in trouble.
help in

trouble.
give thanks.
upon the

holy hills.
O my soul.

\emph{Antiphon}

Fight

against them

Great is the

Lord:

My soul

truly

Thou art the

God:

Praised be

the Lord

O visit me:

with

(9) \emph{Psalm}
35
48
62
77
89
106

(10) \emph{Psalm}
36
49
64
78
94
107

\emph{Antiphon}
that fight

against me.
and highly

to be

praised.

waiteth still

upon God.
that doest

wonders.
for ever-

more.
Thy salva-

tion.

\emph{Antiphon}

Commit thy

way:

The Lord,

even the

most

O praise:

our God

Be merciful:

Sing unto

the Lord:

I will give

great

thanks:

(11) \emph{Psalm}
37
50
66
79
96
108

(12) \emph{Psalm}
38
52
68
80
97
109

\emph{Antiphon}
unto the

Lord.
mighty God:

hath spoken.
ye people.
unto our

sins, O

Lord.
and praise

His name.
unto the

Lord with

my mouth.

(f) \emph{Verse and Response}

(\emph{as on Sundays}.)
Out of
Sion hath God appeared

\emph{In perfect beauty}.

\emph{The Lord's Prayer}
Our Father, \&c.

\emph{Absolution}

(\emph{These numbers answer to the numbers affixed to the Absolutions and Benedictions in} § 2.)
M.
T.
W. Th.
F. S.

1.——2.——3. 1.——2.——3.

1\emph{st Monday}
\emph{Benediction} 1.—1 \footnote{The Benedictions run as follows:

M.
T.
W.
Th.
F.
S.

\emph{Benediction}
1st.—
1.
4.
7.
1.
4.
7.

2d.—
2.
5.
8.
2.
5.
8.

3d.—
3.
6.
9.
3.
6.
9.}.

\emph{Lesson} 1. \emph{Isaiah} i. 16-18 \footnote{These lessons will
illustrate the complaint in the preface of our Prayer

Book, that ``when any book of the Bible was begun,'' in the Breviary, ``after

3 or 4 chapters were read out, all the rest were unread.'' No chapter
is finished.}.

\emph{Response} 1. \emph{Cherish the Word, O
Mary the Virgin, which}\\
\emph{is conveyed to thee
from the Lord by the angel; thou shalt}\\
\emph{conceive in thy
womb and bring forth God and man; that}\\
\emph{thou mayest be
called blessed among all women}.

Yea, thou shalt bring forth a son, yet abide
a virgin; thou

shalt be with child and be a mother, yet know not a man.

\emph{That
thou mayest be blessed among all women}.

\emph{Benediction} 2.—2.

\emph{Lesson} 2. \emph{Isaiah} i. 19-23.

\emph{Response} 2. \emph{Let the heavens
rejoice, and let the earth be}\\
\emph{glad, let the hills
be joyful together before the Lord, for}\\
\emph{He shall come, and
shall have pity on the poor}.

In His time shall the righteous flourish, yea
and abund-

ance of peace.

\emph{And
He shall have pity on the poor}.

\emph{Benediction} 3.—3.

\emph{Lesson} 3. \emph{Isaiah} i. 24-28.

\emph{Response} 3. \emph{Strangers shall not
pass through Jerusalem}\\
\emph{any more, for in
that day the mountains shall drop down}\\
\emph{new wine, and the
hills shall flow with milk, saith the}\\
\emph{Lord}.

God shall come from Lebanon, and the Holy One
from

the woody mountain.

\emph{For
in that day the mountains, \&c}. (\emph{End of Matins}.)

1\emph{st Tuesday}

\emph{Benediction} 1.—4.

\emph{Lesson} 1. \emph{Isaiah} ii. 1-3.

\emph{Response} 1. \emph{Ye mountains of Israel,
spread forth your}\\
\emph{branches, and bud
into flower, and bear fruit; it is full}\\
\emph{time that the day
of the Lord come}.

Drop down, ye heavens, from above, and let
the skies

pour down righteousness; let the earth open and bring

forth salvation.

\emph{It
is full time that the day of the Lord come}.

\emph{Benediction} 2.—5.

\emph{Lesson} 2. \emph{Isaiah} ii. 4-6.

\emph{Response} 2. \emph{Let the mountains break
forth with joy, and}\\
\emph{the hills with
righteousness, for the Lord, the light of the}\\
\emph{world, is come with
power}.

Out of Zion shall go forth the law, and the
word of the

Lord from Jerusalem.

\emph{For
the Lord, the light, \&c}.

\emph{Benediction} 3.—6.

\emph{Lesson} 3. \emph{Isaiah} ii. 7-9.

\emph{Response} 3. \emph{Behold, I come, the
Lord your God from}\\
\emph{Timon: to visit you
in peace}.

I will look upon you, and make you to
increase ; ye shall

be multiplied, and I will establish My covenant with you.

\emph{To
visit you in peace}. (\emph{End of Matins}.)

1\emph{st Wednesday}

\emph{Benediction} 1.—7 † \footnote{Where the number of this
Benediction is marked with a †, it runs thus:

``May He bless us who
liveth and reigneth world without end. \emph{Amen}.''}.

\emph{Lesson} 1. \emph{Isaiah} iii. 1-4.

\emph{Response} 1. \emph{Christ our King shall
come, whom John}\\
\emph{announced as the
Lamb that was to come}.

Kings shall keep peace before him, and the
Gentiles shall

supplicate Him.

\emph{Whom
John announced as the Lamb that was to come}.

\emph{}\\
\emph{Benediction} 2—8.

\emph{Lesson} 2. Isaiah iii. 5-7.

\emph{Response} 2. \emph{Ezekiel prophesied long
before, I saw a gate}\\
\emph{that was shut;
behold, the everlasting God went forth}\\
\emph{from it for the
salvation of the world: and again it was}\\
\emph{shut, figuring the
Virgin, who after the birth remained a}\\
\emph{virgin}. 

The gate which thou sawest, the Lord alone
shall pass

through it.

\emph{And
again it was, \&c}.

(\emph{As Response}9\emph{of Sunday}.)

\emph{Benediction} 3.—9.

\emph{Lesson} 3. \emph{Isaiah} iii. 8-11.

\emph{Response}. 3. \emph{Behold, the days come,
saith the Lord, that I}\\
\emph{will raise up unto
David a righteous branch, and a king}\\
\emph{shall reign and
prosper, and shall execute judgment and}\\
\emph{justice in the
earth; and this is His name whereby He shall}\\
\emph{be called, The Lord
our Righteousness}.

In His days Judah shall be saved, and Israel
shall dwell

safely.

\emph{And
this is His name, \&c}. (\emph{End of Matins}.)

1\emph{st Thursday}\\
(\emph{as on Monday}.)

\emph{Benediction} 1.—1.

\emph{Lesson} 1. \emph{Isaiah} iv. 1-3.

\emph{Response} 1. \emph{Cherish the Word,
\&c}.

(\emph{As Response}2\emph{of Sunday}.)

\emph{Benediction} 2.—2.

\emph{Lesson} 2. \emph{Isaiah} v. 1-4.

\emph{Response} 2. \emph{I saw in the night
visions, and behold, one like}\\
\emph{unto the Son
of man came with the clouds of heaven; and}\\
\emph{there was
given Him dominion, and glory, and a kingdom:}\\
\emph{that all
people, nations, and languages should serve Him}.

His dominion is an everlasting dominion,
which shall not

pass away, and His kingdom that which shall not be

destroyed.

\emph{That all people, \&c}.

(\emph{As Response} 3

\emph{of Sunday}.)

\emph{Benediction} 3.—3.

\emph{Lesson} 3. \emph{Isaiah} v. 5-7.\emph{}\\
\emph{Response} 3. \emph{The angel Gabriel
was sent to Mary, a virgin}\\
\emph{espoused to
Joseph, announcing to her the Word, and the}\\
\emph{virgin was
troubled at the light: Fear not, Mary, thou}\\
\emph{hast found
favour with the Lord: behold, thou shalt con-}\\
\emph{ceive and
bring forth a Son, and He shall be called the Son}\\
\emph{of the Highest}.

The Lord God shall give unto Him the throne
of His

father David, and He shall reign over the house of Jacob for

ever.

\emph{Behold, thou shalt conceive, \&c}. (\emph{End of Matins}.)

1\emph{st Friday}\\

(\emph{As Response} 4

\emph{of Sunday}.)

\emph{Benediction} 1.—4. 

\emph{Lesson} 1. \emph{Isaiah} vi. 1-3.

\emph{Response} 1. \emph{Hail thou that art
highly favoured, the Lord}\\
\emph{is with thee;
the Holy Ghost shall come upon thee, and}\\
\emph{the power of
the highest shall overshadow thee, therefore}\\
\emph{that holy
thing which shall be born of thee shall be called}\\
\emph{the Son of God}.

How shall this be, seeing I know not a man?
And the

Angel answered and said to her.

\emph{The Holy Ghost shall come upon thee, \&c}.

(\emph{As Response} 5

\emph{of Sunday}.)

\emph{Benediction} 2.—5.

\emph{Lesson} 2. \emph{Isaiah} vi. 4-7.\emph{}\\
\emph{Response} 2. \emph{We wait for the
Saviour, the Lord Jesus}\\
\emph{Christ, who
shall change our vile body, that it may be}\\
\emph{fashioned like
unto his glorious body}.

Let us live soberly, righteously, and godly
in this present

world, looking for that blessed hope, and the glorious appear-

ing of the great God.

\emph{Who shall change, \&c}.

(\emph{As Response} 6

\emph{of Sunday}.)

\emph{Benediction} 3.—6.

\emph{Lesson} 3. \emph{Isaiah} vi. 8-10.

\emph{Response} 3. \emph{O my Lord, send, I pray
Thee, by the hand of}\\
\emph{him whom thou
wilt send; behold the affliction of thy}\\
\emph{people: as
Thou hast spoken, come, and deliver us}.

Hear, O Thou Shepherd of Israel, Thou that
leadest

Joseph like a sheep, Thou that sittest between the

cherubim.

\emph{As Thou hast spoken, \&c}. (\emph{End of Matins}.)

l\emph{st Saturday}\\

(\emph{As Response} 7

\emph{of Sunday}.)

\emph{Benediction} 1.—7 †.

\emph{Lesson} 1. \emph{Isaiah} vii. 1-3.

\emph{Response} 1. \emph{Behold, a virgin shall
conceive, and shall bring}\\
\emph{forth a Son,
saith the Lord, and his name shall be called}\\
\emph{Wonderful, the
Mighty God}.

Upon the throne of David, and upon his
kingdom He shall

sit for ever.

\emph{And His name shall be called Wonderful, the Mighty God}.

(\emph{As Response} 8

\emph{of Sunday}.)

\emph{Benediction} 2.—8.

\emph{Lesson} 2. \emph{Isaiah} vii. 4-6.

\emph{Response} 2. \emph{Hear the Word of the
Lord, O ye nations, and}\\
\emph{declare it in
the ends of the earth; and say to the isles}\\
\emph{afar off, our
Saviour shall come}.

Declare it and make it heard; speak and cry
out.

\emph{And say to the isles, \&c}. 

(\emph{As Response} 9 \emph{of Sunday and Res- ponse} 3 \emph{of Wednes- day}.)

\emph{Benediction} 3.—9.

\emph{Lesson} 3. \emph{Isaiah} vii. 10-15.

\emph{Response} 3. \emph{Behold the day is come},
\&c. (\emph{End of Matins}.)

2\emph{d Monday}

\emph{Benediction} 1.—1.

\emph{Lesson} 1. \emph{Isaiah} xiii. 1-4.

\emph{Response} 1. \emph{Cherish the Word,
\&c}.

(\emph{The Response as}

\emph{on} 1\emph{st Monday}.)

\emph{Benediction} 2.—2.

\emph{Lesson} 2. \emph{Isaiah} xiii. 4-8.

\emph{Response} 2. \emph{Let the heavens
rejoice, \&c}.

\emph{Benediction} 3.—3.

\emph{Lesson} 3. \emph{Isaiah} xiii. 9-11.

\emph{Response} 3. \emph{Strangers shall not
pass, \&c}. (\emph{End}.)

2\emph{nd Tuesday}

\emph{Benediction} 1.—4.

\emph{Lesson} 1. \emph{Isaiah} xiv. 1-4.

\emph{Response} 1. \emph{Ye mountains of Israel,
\&c}.

(\emph{Responses as on}

1\emph{st Tuesday}.)

\emph{Benediction} 2.—5.

\emph{Lesson} 2. \emph{Isaiah} xiv. 3-6.

\emph{Response} 2. \emph{Let the mountains,
\&c}.

\emph{Benediction} 3.—6.

\emph{Lesson} 3. \emph{Isaiah} xiv. 12-15.

\emph{Response} 3. \emph{Behold, I come, \&c}.
(\emph{End}.)

2\emph{d Wednesday}

\emph{Benediction} 1.—7 †.

\emph{Lesson} 1. \emph{Isaiah} xvi. 1-4.

\emph{Response} 1. \emph{Christ our king, \&c}.

(\emph{As on} 1\emph{st}

\emph{Wednesday}.)

\emph{Benediction} 2.—8.

\emph{Lesson} 2. \emph{Isaiah} xvi. 4-6.

\emph{Response} 2. \emph{Ezekiel prophesied,
\&c}.

(\emph{As Response} 9

\emph{of Sunday}.)

\emph{Benediction} 3.—9.

\emph{Lesson} 3. \emph{Isaiah} xvi. 7, 8.

\emph{Response} 3. \emph{Behold, the Lord shall
come, descending in}\\
\emph{glory, and His
power with Him: to visit His people in}\\
\emph{peace, and
establish upon them everlasting life}.

Behold our Lord shall come with power.

\emph{To visit His people, \&c}. (\emph{End}.)

2\emph{d Thursday}\\
(\emph{As Response} 1

\emph{of Sunday}.)

\emph{Benediction} 1.—1.

\emph{Lesson} 1. \emph{Isaiah} xix. 1, 2 

\emph{Response} 1. \emph{O Jerusalem, thy
salvation shall come quickly;}

\emph{wherefore art thou
consumed with grief? is there any}

\emph{counsellor in thee?
because sorrow hath changed thee. I}

\emph{will save thee and
deliver thee, fear not}.

For I am the Lord thy God, the Holy one of
Israel, thy

Redeemer.

\emph{I
will save thee and deliver thee, fear not}.

(\emph{As Response} 2

\emph{of Sunday}.)

\emph{Benediction} 2.—2.

\emph{Lesson} 2. \emph{Isaiah} xix. 3-6.

\emph{Response} 2. \emph{The Lord my God shall
come and all the}

\emph{Saints with Thee, and
in that day there shall be great}

\emph{light; and it shall be
in that day that living waters shall}

\emph{go out from Jerusalem,
and the Lord shall be king for}

\emph{ever over all the earth}.

Behold, the Lord shall come with power, and
His king-

dom shall be in His hand, and dominion and sovereignty.

\emph{Over
all the earth}.

(\emph{As Response} 3

\emph{of Sunday}.)

\emph{Benediction} 3.—3.

\emph{Lesson} 3. \emph{Isaiah} xix. 11-13.

\emph{Response} 3. \emph{City of Jerusalem, weep
thou not, for the Lord}

\emph{has sorrowed for thee;
and He shall take from thee all}

\emph{tribulation}.

Behold, the Lord shall come in strength, and
His arm shall

rule.

\emph{And
He shall take, \&c}. (\emph{End}.)

2\emph{d Friday}\\

(\emph{As Response} 4

\emph{of Sunday}.)

\emph{Benediction} 1—4.

\emph{Lesson} 1. \emph{Isaiah} xxiv. 1-3.

\emph{Response} 1. \emph{Behold, the Lord shall
come, our defender, the}

\emph{Holy one of Israel,
having the crown of His kingdom upon}

\emph{Him}.

And He shall reign from sea to sea, and from
the river

even unto the ends of the earth.

\emph{Having
the crown of his kingdom upon him}.

(\emph{As Response} 5

\emph{of Sunday}.)

\emph{Benediction} 2—5.

\emph{Lesson} 2. \emph{Isaiah} xxiv. 4-6.

\emph{Response} 2. \emph{As a mother comforteth
her sons, so will I}

\emph{comfort you, saith the
Lord; and from Jerusalem, the}

\emph{city I have chosen,
shall succour come to you; and you}

\emph{shall see and your
heart shall rejoice}.

I will give to Zion salvation, and to
Jerusalem, My glory,

\emph{And
you shall see, \&c}. 

(\emph{As Response} 6

\emph{of Sunday}.)

\emph{Benediction} 3.—6.

\emph{Lesson} 3. \emph{Isaiah} xxiv. 7-15.

\emph{Response} 3. \emph{O Jerusalem, thou shalt
plant thy vine upon}

\emph{thy mountains; leap for
joy, for the day of the Lord shall}

\emph{come; arise, O Zion, be
turned unto the Lord thy God;}

\emph{rejoice and be glad, O
Jacob: for thy Saviour shall come}

\emph{in the midst of the
nations}.

Rejoice greatly, O daughter of Zion, shout, O
daughter of

Jerusalem.

\emph{For
thy Saviour, \&c}. (\emph{End}.)

2\emph{d Saturday}\\

(\emph{As Response} 7

\emph{of Sunday}.)

\emph{Benediction} 1.—7 †.

\emph{Lesson} 1. \emph{Isaiah} xxv. 1-4.

\emph{Response} 1. \emph{The Lord shall come
from Samaria to the}

\emph{eastern gate, and shall
come to Bethlehem, walking upon}

\emph{the waters of the
redemption of Judah; then shall every}

\emph{man be saved, for,
behold, He cometh}.

And His throne shall be prepared in mercy,
and He shall

sit on it in truth.

\emph{Then
shall every man be saved, for, behold, He cometh}.

(\emph{As Response} 8

\emph{of Sunday}.)

\emph{Benediction} 2.—8.

\emph{Lesson} 2. \emph{Isaiah} xxiv. 4-7.

\emph{Response} 2. \emph{Make haste, and tarry
not, O Lord: and}

\emph{deliver Thy people}.

Come, O Lord, and do not tarry: absolve thy
people from

their iniquities.

\emph{And
deliver Thy people}.

(\emph{As Response} 9

\emph{of Sunday and} 3

\emph{of Wednesday}.)

\emph{Benediction} 3.—9.

\emph{Lesson} 3. \emph{Isaiah} xxv. 8-12.

\emph{Response} 3. \emph{Behold, the Lord shall
come, \&c}. (\emph{End}.)

3\emph{d Monday}

(\emph{As Response} 1

\emph{of Sunday}.)

\emph{Benediction} 1.—1.

\emph{Lesson} 1. \emph{Isaiah} xxviii. 1-3.

\emph{Response} 1. \emph{Behold the Lord shall
appear, \&c}.

(\emph{As Response} 2

\emph{of Sunday}.)

\emph{Benediction} 2.—2.

\emph{Lesson} 2. \emph{Isaiah} xxviii. 4-7.

\emph{Response} 2. \emph{O Bethlehem, city,
\&c}.

(\emph{As Response} 3

\emph{of Sunday}.)

\emph{Benediction} 3.—3.

\emph{Lesson} 3. \emph{Isaiah} xxviii. 16-18.

\emph{Response} 3. \emph{He that shall come,
will come, \&c}. (\emph{End}.)

3\emph{d Tuesday}

(\emph{As Response} 4

\emph{of Sunday}.)

\emph{Benediction} 1.—4.

\emph{Lesson} 1. \emph{Isaiah} xxx. 18-20.

\emph{Response} 1. \emph{Weep not, O Egypt,
\&c}.

(\emph{As Response} 5

\emph{of Sunday}.)

\emph{Benediction} 2.—5.

\emph{Lesson} 2. \emph{Isaiah} xxx. 22-25.

\emph{Response} 2. \emph{Her time is well nigh,
\&c}.

(\emph{As Response} 6

\emph{of Sunday}.)

\emph{Benediction} 3.—6.

\emph{Lesson} 3. \emph{Isaiah} xxx. 26-28.

\emph{Response} 3. \emph{The Lord shall descend
as the rain, \&c}. (\emph{End}.)

3\emph{d Wednesday}

\emph{Benediction} 1.—7.

\emph{Lesson} 1. \emph{Luke} i. 26-38.

\emph{Homily of St. Ambrose}

(\emph{Lib}. 2 \emph{in Lucam}.)

(\emph{for Ember day}.)

Divine mysteries lie hid, nor is it easy, as
the prophet

says, for any one among men to know the purpose of God.

Yet, from the other deeds and precepts of the Lord and

Saviour, we are able to understand that it was a matter of

solicitous purpose that she should be chosen above others

to bear the Lord, who was already betrothed to a husband.

For why did she not conceive before she was
betrothed?

probably, lest it should be said that she had conceived in

adultery.

\emph{Response} 1. \emph{O Jerusalem, that
bringest good tidings, lift}

\emph{up thy voice, be not
afraid; say unto the cities of Judah,}

\emph{Behold your God! Behold
the Lord God will come}

\emph{whom we did wait for}.

O Sion, which bringest good tidings; get thee
up into the

high mountains, lift up thy voice with strength.

\emph{Say
unto the cities of Judah, Behold your God! Be-}

\emph{hold, the Lord God will come whom we did wait for}.

(\emph{Homily con- tinued}.)
\emph{}

\emph{Benediction} 2.—8.

\emph{Lesson} 2.

And the angel came in unto her. Recognize the
virgin in

her behaviour, in her modesty, in the annunciation, in the

mystery. To be frightened belongs to virgins; and to dread

the approach, to feel abashed at the address of man. Let

woman learn to imitate the settled resolve of modesty herein

displayed. Alone was she in the inner chamber, seen by no

man, found by the angel alone. Alone without companion, 

alone without witness, lest she should be insulted by any

unworthy voice, she is saluted by the angel.

\emph{Response} 2. \emph{There shall come a star
out of Jacob, and a}

\emph{sceptre shall arise out
of Israel, and shall smite through}

\emph{the princes of Moab:
and all the earth shall be His}

\emph{possession}.

All kings shall fall down before Him, all
nations shall do

Him service.

\emph{And
all the earth, \&c}.

(\emph{Homily con- tinued}.)

\emph{Benediction} 3.—9.

\emph{Lesson} 3. For the mystery of so high a
message was to

be divulged by the mouth, not of man, but of an Angel.

Today is it first heard; The Holy Ghost shall come upon

thee. It is heard, it is believed. Lastly she says, Behold the

handmaid of the Lord; be it unto me according to thy word.

Observe her humility, observe her devotion. She calls her-

self the handmaid, who is chosen to be the mother of the

Lord; nor is she elated by the sudden promise.

\emph{Response} 3. \emph{The Lord Almighty shall
soon come; and}

\emph{His name shall be
called Emmanuel}.

Righteousness shall arise in His days, and
abundance of

peace.

\emph{And
his name, \&c}. (\emph{End}.)

3\emph{d Thursday}
\emph{}

\emph{Benediction} 1.—1.

\emph{Lesson} 1. \emph{Isaiah} xxxiii. 1, 2.

\emph{Response} 1. \emph{The Lord shall go forth
and fight against}

\emph{the nations; and His
feet shall stand in that day upon}

\emph{the mount of Olives,
towards the east}.

And He shall be exalted above the hills, and
all nations

shall flow unto Him.

\emph{And
His feet, \&c}.

\emph{Benediction} 2.—2.

\emph{Lesson} 2. \emph{Isaiah} xxxiii. 3-6.

\emph{Response} 2. \emph{The forerunner is for
us entered, a Lamb}

\emph{without spot; made an
high priest for ever and ever,}

\emph{after the order of
Melchizedek}.

He is the king of righteousness, and of His
generation

there is no end.

\emph{Made
an high priest, \&c}.

\emph{Benediction} 3.—3.

\emph{Lesson} 3. \emph{Isaiah} xxxiii. 14-17.

\emph{Response} 3. \emph{The Gentiles shall see
Thy righteousness}, 

\emph{and all kings Thy
glory: and Thou shalt be called by}

\emph{a new name, which the
mouth of the Lord shall name}.

Thou shalt also be a crown of glory in the
hand of the

Lord, and a royal diadem in the hand of thy God.

\emph{And
Thou shalt be called, \&c}. (\emph{End}.)

3\emph{d Friday}

\emph{Homily of St. Ambrose, Lib}.

2 \emph{in Luc}.

(\emph{For Ember day}.)

\emph{Benediction} 1.—7.

\emph{Lesson} 1. \emph{Luke} i. 39-45.

Those who exact faith, are used to give
grounds for it.

Accordingly the angel when announcing what was hidden

to give grounds for faith, by the precedent of the elder and

barren woman, announced the conception of the Virgin, by

way of declaring that with God all is possible whatsoever

He will. Mary, on hearing this, not as incredulous at the

heavenly voice, nor as uncertain as to the message, nor as

doubtful of the precedent, but as joyful from her resolve,

religious by her calling, speedy for joy, proceeded into the

mountainous country. Whither indeed should she hasten,

who was filled with God, but to what is upward?

\emph{Response} 1. \emph{Send ye the Lamb to the
Ruler of the land,}

\emph{from Sela to the
wilderness, unto the mount of the}

\emph{Daughter of Sion}.

O Lord, show Thy mercy upon us, and grant us
Thy sal-

vation.

\emph{From
Sela to the wilderness, unto the mount of the}

\emph{daughter of Sion}.

(\emph{Homily con- tinued}.)

\emph{Benediction} 2.—8.\emph{}

\emph{Lesson} 2. Learn hence also, ye holy
women, what at-

tention ye should pay to your relations who are with child.

Mary, who before lived in solitude in the inner chambers,

was not deterred from going abroad by her virgin modesty,

nor from her eagerness by the mountainous journey, nor

from her dutifulness by its length. A virgin leaves her

home and goes into the mountains with haste, a virgin mind-

ful of her duty, unmindful of the hardships, called not by

her sex, but by affection. Learn hence, O virgins, not to go

about from house to house, not to loiter in public places, not

to talk freely in society. Mary, loitering at home, in haste

when abroad, abode for three months with her kinswoman.

\emph{Response} 2. \emph{Drop down ye heavens
from above, and let}

\emph{the skies pour down
righteousness; let the earth open}

\emph{and bring forth
salvation}. 

Send ye the Lamb to the Ruler of the land,
from Sela to

the wilderness, unto the mount of the daughter of Sion.

\emph{Let
the earth, \&c}.

\emph{Homily con- tinued}.

\emph{Benediction} 3.—9.\emph{}

\emph{Lesson} 3. Ye have been taught, O
virgins, the modesty

of Mary, observe the humility. She came as kinswoman to

her next of kin, as junior to her elder; nor came alone, but

first saluted her. For it beseems a virgin, the purer she is,

to be the more humble also. Let her be used to defer to her

seniors. Let her be a pattern of humility, who has taken

the profession of chastity. Here too is a rule of doctrine, as

well as a pattern of piety. For let it be considered that the

superior comes to the inferior that this inferior may be aided:

Mary to Elizabeth, Christ to the Baptist.

\emph{Response} 3. \emph{The fields of the
desert have brought forth the}

\emph{plant of the odour of
Israel, for behold, our God shall}

\emph{come with strength: and
His brightness with Him}.

Out of Sion the perfection of beauty our God
shall mani-

fest Himself.

\emph{And
His brightness with Him}. (\emph{End}.)

3\emph{d Saturday}

\emph{Homily of Pope Gregory in Evang}. 20.

(\emph{For Ember day}.)

\emph{Benediction} 1.—7.

\emph{Lesson} 1. \emph{Luke} iii. 1-3.

The time when the forerunner of our Redeemer
received

the office of preaching, is marked by mention of the

Roman Emperor, and the Kings of Judea. For, whereas he

came to proclaim Him, who was to redeem some from Judea

and many from the heathen, therefore, the date of his

preaching is marked by the ruler of the Gentiles and the

governor of the Jews. And whereas heathenism was to be

brought in, and Judea scattered for the guilt of its misbelief,

the very description of the earthly sovereign signifies this;

the prince of the Roman state being described as one, of

Judea as fourfold.

\emph{Response} 1. \emph{There shall come forth
a rod out of the stem}

\emph{of Jesse, and a branch
shall grow out of His roots; and}

\emph{righteousness shall be
the girdle of His loins, and faith-}

\emph{fulness the girdle of
His reins}.

And the Spirit of the Lord shall rest upon
him; the

spirit of wisdom and understanding, the spirit of counsel

and might.

\emph{And
righteousness shall be, \&c}.

\emph{Homily con- tinued}.

\emph{Benediction} 2.—8. 

\emph{Lesson} 2. And indeed, the voice our
Redeemer has de-

clared, Every kingdom divided against itself shall be

brought to desolation. It is plain, therefore, that Judea

had come to the end of its kingdom, being parted among

so many kings. Also, it is suitably noted what priests, as

well as what kings, were in power at the time: and, whereas

John the Baptist was proclaiming Him who was at once

king and priest, Luke the Evangelist marked the date of

that proclamation by mention of the kingdom and the

priesthood.

\emph{Response} 2. \emph{There shall be a root
of Jesse, which shall}

\emph{rise to judge the
Gentiles: on Him shall the Gentiles}

\emph{trust: and His name
shall be blessed for evermore}.

The kings shall shut their mouths at Him: the
Gentiles

shall implore Him,

\emph{And
His name, \&c}.

\emph{Homily con- tinued}.

\emph{Benediction} 3.—9.

\emph{Lesson} 3. And he came into all the
regions of Jordan,

proclaiming the baptism of repentance for the remission of

sins. It is plain to every reader, that John not only pro-

claimed the baptism of repentance, but also administered it

to some, yet he could not administer it for the remission of

sins. For the remission of sins is imparted to us only in

Christ's baptism. The words then are observable; \emph{pro- claiming} the baptism of repentance for the remission of
sins.

He proclaimed that cleansing Baptism which he could not

administer: that, as he forerun the Incarnate Word of the

Father by the word of proclamation, so he might forerun

also the baptism of repentance, in which sins are forgiven,

with his own baptism, by which sins could not be forgiven.

\emph{Response} 3. \emph{Come, O Lord, and do
not tarry; absolve Thy}

\emph{people from their sins;
and recover the scattered ones into}

\emph{their own land}.

O Lord raise up Thy power, and come to save
us.

\emph{And
recover, \&c}. (\emph{End}.)

4\emph{th Monday}\\

(\emph{As Response} 1

\emph{of Sunday}.)

\emph{Benediction} 1.—1.

\emph{Lesson} 1. \emph{Isaiah} xli. 8-10.

\emph{Response} 1. \emph{Blow the trumpet,
\&c}.

(\emph{As Response} 2

\emph{of Sunday}.)

\emph{Benediction} 2.—2.

\emph{Lesson} 2. \emph{Isaiah} xli. 11-13.

\emph{Response} 2. \emph{The sceptre shall not
depart, \&c}. 

(\emph{As Response} 3

\emph{of Sunday}.)

\emph{Benediction} 3.—3.

\emph{Lesson} 3. \emph{Isaiah} xli. 14-16.

\emph{Response} 3. \emph{I must decrease, \&c}.
(\emph{End}.)

4\emph{th Tuesday}\\

(\emph{As Response} 4

\emph{of Sunday}.)

\emph{Benediction} 1.—4.

\emph{Lesson} 1. \emph{Isaiah} xlii. 1-4.

\emph{Response} 1. \emph{Unto us as a child,
\&c}.

(\emph{As Response} 5

\emph{of Sunday}.)

\emph{Benediction} 2.—5.

\emph{Lesson} 2. \emph{Isaiah} xlii. 5-7.

\emph{Response} 2. \emph{Behold the fulness of
time, \&c}.

(\emph{As Response} 6

\emph{of Sunday}.)

\emph{Benediction} 3.—6.

\emph{Lesson} 3. \emph{Isaiah} xlii. 10-13.

\emph{Response} 3. \emph{O virgin of Israel,
\&c}. (\emph{End}.)

4\emph{th Wednesday}

(\emph{As Response} 7

\emph{of Sunday}.)

\emph{Benediction} 1.—7 †.

\emph{Lesson} 1. \emph{Isaiah} li. 1-3.

\emph{Response} 1. \emph{I have sworn, \&c}.

(\emph{As Response} 8

\emph{of Sunday}.)

\emph{Benediction} 2.—8.

\emph{Lesson} 2. \emph{Isaiah} li. 4-6.

\emph{Response} 2. \emph{And so will not we,
\&c}.

(\emph{As Response} 9

\emph{on Sunday}.)

\emph{Benediction} 3.—9.

\emph{Lesson} 3. \emph{Isaiah} li. 7, 8.

\emph{Response} 3. \emph{Consider how great,
\&c}. (\emph{End}.)

4\emph{th Thursday}

(\emph{As Response} 1

\emph{of Sunday and Response} 1 \emph{of Monday}.)

\emph{Benediction} 1.—1.

\emph{Lesson} 1. \emph{Isaiah} lxiv. 1-4.

\emph{Response} 1. \emph{Blow the trumpet,
\&c}.

(\emph{As Response} 2

\emph{of Sunday and Response} 2 \emph{of Monday}.)

\emph{Benediction} 2.—2.

\emph{Lesson} 2. \emph{Isaiah} lxiv. 5-7.

\emph{Response} 2. \emph{The sceptre shall not
depart, \&c}.

(\emph{As Response} 3

\emph{of Sunday and Response} 3 \emph{of Monday}.)

\emph{Benediction} 3.—3.

\emph{Lesson} 3. \emph{Isaiah} lxiv. 8-11.

\emph{Response} 3. \emph{I must decrease, \&c}.
(\emph{End}.)

4\emph{th Friday}\\

(\emph{As Response} 4

\emph{of Sunday and Response} 1 \emph{of Tuesday}.)

\emph{Benediction} 1.—4. 

\emph{Lesson} 1. \emph{Isaiah} lxvi. 5-8.

\emph{Response} 1. \emph{Unto us as a child,
\&c}.

(\emph{As Response} 5

\emph{of Sunday and Response} 2 \emph{of Tuesday}.)

\emph{Benediction} 2.—5.

\emph{Lesson} 2. \emph{Isaiah} lxvi. 9-12.

\emph{Response} 2. \emph{Behold the fulness of
time, \&c}.

(\emph{As Response} 6

\emph{of Sunday and Response} 3 \emph{of Tuesday}.)

\emph{Christmas Eve, \&c. \&c}.

\emph{Christmas Day, \&c. \&c}.

\emph{Benediction} 3.—6.

\emph{Lesson} 3. \emph{Isaiah} lxvi. 13-16.

\emph{Response} 3. \emph{O virgin of Israel,
\&c}. (\emph{End}.)

\xsubsection{LAUDS}

O God, make speed, \&c.

\emph{O Lord, make haste, \&c}.

Glory be, \&c. Amen. Hallelujah.

\emph{Antiphons and Psalms}

1\emph{st}, 2\emph{d}, \& 3\emph{d}

\emph{Monday}.
4\emph{th Monday} (the 19th.)

(o) Have mercy.

(o) Behold, the Lord will come, the King of
the whole

earth.

(1) \emph{Psalm} 51.

(o) Have mercy

on me, O God.
(o) Behold, the Lord will
come, the King of the whole

earth; blessed are they that are ready to meet Him.

(p) Consider.

(p) When the Son of man cometh.

(2) \emph{Psalm} 5.

(p) Consider my

meditation.
(p) When the Son of man
cometh, will He find faith on

the earth?

(q) O God, Thou

art my God.

(q) Behold, the fulness of time is now come.

(3) \emph{Psalm} 63 \& 67.

(q) O God, Thou

art my God,

early will I seek

Thee.
(q) Behold, the fulness of
time is now come, in the which

God sent forth his Son to the earth.

(r) Thine anger.

(r) With joy shall ye draw water.

(4) Song of Isaiah. (\emph{Isaiah} xii.)

(r) Thine anger

is turned away,

and Thou com-

fortedst me.
(r) With joy shall ye draw
water out of the wells of sal-

vation. 

(s) O praise the

Lord.

(s) The Lord shall go forth of His holy
place.

(5) \emph{Psalm} 148-150.

(s) O praise the

Lord of heaven.
(s) The Lord shall go forth of
His holy place, He shall

come to save his people. (\emph{go on} (t))

1\emph{st}, 2\emph{d} \& 3\emph{d}

\emph{Tuesday}

4\emph{th Tuesday}, (the 20th.)

(o) Cleanse me.

(o) Drop down, ye heavens.

(1) \emph{Psalm} 51.

(o) Cleanse me.
(o) Drop down, ye heavens,
from above, and let the skies

pour down righteousness; let the earth open and bring

forth salvation.

(p) The help.

(p) Send ye the Lamb.

(2) \emph{Psalm} 43.

(p) The help of

my countenance

and my God.
(p) Send ye the Lamb to the
Ruler of the land, from

Sela in the wilderness, unto the mount of the daughter of

Zion.

(q) Early.

(q) That Thy way may be known.

(3) \emph{Psalm} 63 \& 67.

(q) Early will I

seek Thee, my

God.
(q) That Thy way may be known
upon earth, Thy saving

health among all nations.

(r) Save us O

Lord.

(r) Reward, O Lord.

(4) Song of Hezekiah. (\emph{Isaiah} xxxviii.)

(r) Save us, O

Lord, all the

days of our life.
(r) Reward, O Lord, them that
wait for Thee, that Thy

prophets be found faithful.

(s) Praise him.

(s) The Law was given by Moses.

(5) \emph{Psalm} 148-150.

(s) Praise Him,

all ye angels of

His.
(s) The Law was given by
Moses, but grace and truth

came by Jesus Christ. (\emph{go on} (t))

1\emph{st}, 2\emph{d} \& 3\emph{d}

\emph{Wednesday}.

4\emph{th Wednesday}, (the 21st, St.
Thomas's.)

(o) Wash me

throughly.

(o) This is My commandment, that ye love one
another

as I have loved you.

(1) \emph{Psalm} 51.
(1) \emph{Psalm} 94.

(o) Wash me

throughly from

my wickedness.
(o) This is My commandment,
that ye love one another

as I have loved you. 

(p) Thou, O

Lord.

(p) Greater love hath no man than this, that
a man lay

down his life for his friends.

(2) \emph{Psalm} 65.

(2) \emph{Psalm} 100.

(p) Thou, O

Lord, art prais-

ed in Sion.

(p) Greater love hath no man
than this, that a man lay

down his life for his friends.

(q) My lips

shall praise

Thee.

(q) Ye are my friends, if ye do whatsoever I
command

you, saith the Lord.

(3) \emph{Ps}. 63 \& 67.

(3) \emph{Ps}. 63 \& 67.

(q) My lips

shall praise

Thee as long as

I live, my God.

(q) Ye are my friends, if ye
do whatsoever I command

you, saith the Lord.

(r) The Lord

shall judge.

(r) Blessed are the peace-makers, blessed are
the pure in

heart, for they shall see God.

(4) \emph{Song of Hannah} (1 \emph{Sam}.

ii.)

(4) \emph{Song of the Three
Children}. (\emph{Benedicite Omnia Opera}.)

(r) The Lord

shall judge the

ends of the

earth.

(r) Blessed are the
peace-makers, blessed are the pure in

heart, for they shall see God.

(s) Praise Him.

(s) In your patience possess ye your souls.

(5) \emph{Ps}. 148-150.

(5) \emph{Psalm} 148-150.

(s) Praise Him,

all ye heavens.

(s) In your patience possess
ye your souls.

1\emph{st}, 2\emph{d} \& 3\emph{d}

\emph{Thursday}.

4\emph{th Thursday} (the 22d.)

(o) Against

Thee only have

I sinned.

(o) The Lord Almighty shall come.

(1) \emph{Psalm} 51.

(o) Against

Thee only have

I sinned, have

mercy upon me,

O Lord.

(o) The Lord Almighty shall
come to save His people.

(p) Lord.

(p) Turn Thee, Lord.

(2) \emph{Psalm} 90.

(p) Lord, Thou

hast been our

refuge.

(p) Turn Thee, Lord, a little
while, and delay not to come

to Thy servants.

(q) Have I not

thought.

(q) The Lord who is to reign.

(3) \emph{Psalm} 63 \& 67.

(q) Have I not

thought upon

Thee when I

was waking.

(q) The Lord who is to reign
shall come from Sion;

Emmanuel is His mighty name.

(r) I will sing.

(r) He is my God.

(4) \emph{Song of Moses} (\emph{Exod}. xv.)

(r) I will sing

unto the Lord,

for he hath tri-

umphed glori-

ously.

(r) He is my God, and I will
prepare him an habitation;

my father's God, and I will exalt Him. 

(s) O praise God.

s) The Lord is our lawgiver.

(5) \emph{Psalm} 148-150.

(s) O praise God

in His holiness.

(s) The Lord is our lawgiver,
the Lord is our king: He

will come and save us.

1\emph{st}, 2\emph{nd}, \& 3\emph{rd}

\emph{Friday}.

4\emph{th Friday} (the 23d.)

(o) Stablish me.

(o) Stand still.

(1) \emph{Psalm} 51.

(o) Stablish me

with Thy free

Spirit.

(o) Stand still, and ye shall
see the salvation of the Lord.

(p) Hearken unto

me.

(p) To Thee, O Lord, have I lift up my soul.

(2) \emph{Psalm} 143.

(p) Hearken unto

me for Thy truth

and righteousness'

sake.

(p) To Thee, O Lord, have I
lift up my soul; come, O

Lord, and deliver me, for I flee unto thee.

(q) Shew us the

light.

(q) Come, O Lord, and do not tarry.

(3) \emph{Psalm} 63 and 67.

(q) Shew us the

light of Thy

countenance.

(q) Come, O Lord, and do not
tarry: absolve Thy people

Israel from their iniquities.

(r) O Lord.

(r) God shall come from Lebanon.

(4) \emph{Song of Habakhuk} (\emph{Hab}.
iii.)

(r) O Lord, I

have heard Thy

speech, and

was afraid.

(r) God shall come from
Lebanon, and His brightness

shall be as a stream.

(s) Praise Him.

(s) But I will look towards the Lord.

(5) \emph{Psalm} 148-150.

(s) Praise Him

in the cymbals

and dances.

(s) But I will look towards
the Lord, and will wait for

God my Saviour.

1\emph{st and} 2\emph{nd Saturday}.

3\emph{rd Saturday} (the 17th.)

(o) O be favour-

able.

(o) The prophets announced.

(1) \emph{Psalm} 51.

(o) O be favour-

able and gracious.

(o) The prophets announced
that the Saviour should be

born of the Virgin Mary.

(p) It is a good

thing.

(p) The Spirit of the Lord.

(2) \emph{Psalm} 92.

(p) It is a good

thing to give

thanks unto the

Lord.

(p) The Spirit of the Lord is
upon me; He hath sent

me to preach good tidings to the poor. 

(q) All the ends

of the world.

(q) for Sion's sake will I not hold my peace.

(3) \emph{Psalm} 63 and 67.

(q) All the ends

of the world

shall fear Him.

(q) For Sion's sake will I not
hold my peace, until the

righteousness thereof goeth forth as brightness.

(r) Ascribe ye

greatness.

(r) My speech shall drop as the rain.

(3) \emph{Song of Moses} (Deut. xxxii.)

(r) Ascribe ye

greatness unto

our God.

(r) My speech shall drop as
the rain, and our God shall

come down upon us as the dew.

(s) Praise him.

(s) Tell the people, and say.

(5) \emph{Psalm} 148-150.

(s) Praise Him

upon the well-

tuned cymbals.

(s) Tell the people, and say,
Behold God our Saviour

shall come.

(t) \emph{Text}

Isa. ii. 3.

(\emph{for every day}.)

Come ye, and let us go up to the mountain of
the Lord,

to the house of the God of Jacob; and He will teach us of His

ways, and we will walk in His paths; for out of Zion shall

go forth the law, and the Word of the Lord from Jerusalem.

(u) \emph{Hymn}

En clara vox redarguit.

(v) \emph{Verse and Response}

(\emph{for every day}.)

The voice of one crying in the wilderness,
Prepare ye

the way of the Lord.

\emph{Make his paths straight}.

(w) \emph{Antiphon of the Benedictus}\\
1\emph{st Monday}

The angel of the Lord made announcement to
Mary, and

she conceived of the Holy Ghost. Hallelujah.

1\emph{st Tuesday}

Before they came together,
Mary was found with child of

the Holy Ghost. Hallelujah.

1\emph{st Wednesday}

Out of Sion shall go forth the
law, and the Word of the

Lord from Jerusalem.

1\emph{st Thursday}

Blessed art thou among women,
and blessed is the fruit

of thy womb.

1\emph{st Friday}

Behold God and man of the
house of David shall come to

sit on the throne. Hallelujah.

1\emph{st Saturday}

O Sion, be not afraid: behold
thy God shall come.

Hallelujah.

2\emph{nd Monday}

The Almighty Lord shall come
from heaven, and in His

hand is honour and dominion.

2\emph{nd Tuesday}

The Lord shall arise upon
thee, O Jerusalem; and his

glory shall be seen in thee.

2\emph{nd Wednesday}

Behold I send My messenger,
who shall prepare My way

before Thee. 2\emph{nd Thursday}

Thou art He that should come,
O Lord, whom we look

for, to save Thy people.

2\emph{nd Friday}

Say ye, strengthen the
weak-hearted; behold our Lord

God shall come.

2\emph{nd Saturday}

The Lord shall lift up an
ensign among the nations, and

shall gather together the dispersed of Israel.

3\emph{rd Monday}

A rod shall go forth out of
the root of Jesse, and all the

earth shall be filled with the glory of the Lord; and all

flesh shall see the salvation of God.

3\emph{rd Tuesday}

Thou, Bethlehem, in the land
of Judah, art not the least

among the princes of Judah; for out of thee shall come

forth a governor, that shall rule my people Israel.

3\emph{rd Wednesday}

The Angel Gabriel was sent to
a Virgin, Mary, who was

espoused to Joseph.

3\emph{rd Thursday}

Watch ye in heart, for the
Lord our God is at hand.

3\emph{rd Friday}

As soon as the voice of thy
salutation sounded in mine

ears, the babe leaped in my womb for joy. Hallelujah.

3\emph{rd Saturday}

How shall this be, Angel of
God, seeing I know not a

man? Hear thou, Mary the Virgin, the Holy Ghost shall

come upon thee, and the power of the Highest shall over-

shadow thee.

4\emph{th Monday}

The Lord saith, Repent ye, for
the kingdom of heaven

is at hand.

4\emph{th Tuesday}

Awake, awake, put on strength,
O arm of the Lord.

4\emph{th Wednesday}

(the 21st.)

Be afraid, for on the fifth
day our Lord shall come to you.

4\emph{th Thursday}

I will place salvation in Sion,
and in Jerusalem my glory.

4\emph{th Friday}

Behold all things are
accomplished, which were said by

the Angel concerning Mary the Virgin.

\emph{Christmas Eve, \&c}.

(x) \emph{Collect}

(\emph{the same as in the four Sundays respectively}.)

(\emph{Except on the three Ember days, when it is as follows}:)

1\emph{st Week}. Raise up we
beseech Thee, \&c.

2\emph{nd Week}. Stir up, O Lord, \&c.

3\emph{rd Week}. Give ear, we beseech Thee,
\&c.

4\emph{th Week}. Raise up, we beseech Thee,
\&c.

(x) \emph{Collect for Ember Wednesday}

Grant to us, we beseech Thee,
Almighty God, that when

the festival of our redemption cometh, it may both furnish

us with aids for this present life, and impart to us the re-

wards of everlasting bliss, through the Lord, \&c.

\emph{Friday}

Raise up, we beseech Thee, O
Lord, Thy power, and

come, that they who rely upon Thy mercy may be speedily

delivered from all adversity, who livest and reignest, \&c.

\emph{Saturday}

O God, who seest that we are
afflicted by our own cor-

ruption, mercifully grant that we may be consoled by Thy

visitation, who livest, \&c.

\xsubsection{PRIME, THIRD, SIXTH, NINTH}

((o) (p) (q) (s))
\emph{The Service is the same as
on the respective Sundays, except that from the} 17\emph{th inclusive, the Antiphons in the four
Services are taken successively from the} 1\emph{st}, 2\emph{d}, 3\emph{d},
\emph{and} 5\emph{th Antiphons at Lauds, on the same day};

\emph{That in} Prime, \emph{there
is a Psalm according to the day of the week, instead of the Benedicite}, (vid. § 1. supra.) \emph{the
Athana- sian Creed is omitted, and the} Text (c) is Love the truth
and

peace, saith the Lord of Hosts, \emph{instead of} To the King,
\&c.

\emph{And in the Third, Sixth, and Ninth, the
collect}(nn, ss, xx)

\emph{is the same as at the Lauds} (x) \emph{of each day; and that
the Text for every day is as follows}:

\emph{At the Third}

(ll)

Behold the days come, saith the Lord, that I
wilt raise

unto David a righteous branch, and a king shall reign and

prosper, and shall execute judgment and justice in the earth.

\emph{At the Sixth}

(qq)
In his days Judah shall be
saved and Israel shall dwell

safely; and this is His name whereby He shall be called,

the Lord our righteousness.

\emph{At the Ninth}

(vv)

(\emph{and on the} 21\emph{st at Prime also, instead of the Lectio Brevis}.)
Her time is near to come, and
her days shall not be pro-

longed; for the Lord will have mercy on Jacob, and will

yet choose Israel.

\xsubsection{VESPERS}

\emph{Monday}
\emph{Tuesday}
\emph{Wednesday}
\emph{Thursday}
\emph{Friday}
\emph{Saturday}

(yy) \emph{Ant}.
He hath

inclined
We will

go into
They shall

not be

ashamed

when they

speak
And all
Even

before the

gods will
Blessed is

the

\emph{Psalm} (1)
116 (\emph{part})
122
127
132
138
144

(yy) \emph{Ant}.
His ear

unto me.
the house

of the Lord.
with their

enemies in

the gate.

his trouble.
I sing

praise unto

Thee.
Lord my

strength.

(zz) \emph{Ant}.
I believed:

and there-

fore
O thou

that

dwellest
Blessed are

all they
Behold

how good

and joyful

a thing

it is,
O Lord,

thou hast
Every day

will I

\emph{Psalm} (2)
116 (\emph{part})
123
128
133
139
145

(zz) \emph{Ant}.
will I speak.
in the

heavens

have mercy

upon us.

that fear

the Lord.
brethren,

to dwell

together

in unity.
searched

me out and

known me.
give thanks

unto Thee.

(aaa) \emph{Ant}.
Praise the

Lord
Our help

is in the
Many a

time have
Whatsoever

the Lord
Preserve

me from
While I

live will

\emph{Psalm} (3)
117
124
129
135
140
146

(aaa) \emph{Ant}.
all ye

heathen.
name of

the Lord.
they fought

against me

from my

youth up.

pleased,

that did

He.
the wicked

man.
I praise

the Lord.

(bbb) \emph{Ant}.
I called

upon the

Lord,
Do good,

O Lord,

unto
Out of the

deep have I
For His

mercy
Lord, I

call upon

Thee:
A joyful

and

pleasant

thing

\emph{Psalm} (4)
120
125
130
136
141
147 (\emph{part})

(bbb) \emph{Ant}.
and He

heard me.
those who

are good,

and turn our

hearts.

called unto

Thee, O

Lord.
endureth

for ever.
haste Thou

unto me.
it is to be

thankful.

(ccc) \emph{Ant}.
From

whence

cometh
Then

were we

like
O Israel,

trust
Sing us

one of
Thou art

my portion
Praise the

Lord

\emph{Psalm} (5)
121
126
131
137
142
147 (\emph{part})

(ccc) \emph{Ant}.
my help.
unto them

that dream.
in the

Lord.
the songs of

Sion.
in the land

of the

living.
O

Jerusalem.

(ddd) \emph{Text}

(\emph{On Saturdays the Text is that of Lauds and Vespers of the following Sundays}.)
The Sceptre shall not depart
from Judah, nor a lawgiver

from between his feet, until Shiloh come, and unto Him shall

the gathering of the people, \&c.

(eee) \emph{Hymn}

Creator alme siderum, \&c.

(fff) \emph{Verse and Response}
Drop down, ye heavens, from
above, and let the skies

pour down righteousness. \emph{Let the earth open and
bring forth salvation}.

(ggg)

1\emph{st Monday}
\emph{Antiphon of the Magnificat}.

Lift up thine eyes, O Jerusalem, and behold
the greatness

of thy King. Behold thy Saviour cometh to loose thee from

thy chains.

1\emph{st Tuesday}
Seek ye the Lord while He may
be found; call upon

Him while He is near. Hallelujah.

1\emph{st Wednesday}
One that is mightier than I
shall come after Me, whose

shoes' hatchet I am not worthy to unloose.

1\emph{st Thursday}
I will wait for the Lord my
Saviour, and will attend upon

Him while He is near. Hallelujah.

1\emph{st Friday}
Out of Egypt have I called My
Son. He shall come to

save His people.

1\emph{st Saturday}
Come, O Lord, to visit us in
peace, that we may rejoice

before Thee with a perfect heart.

2\emph{nd Monday}
Behold the king, the Lord of
the earth shall come; and

He shall take away the yoke of our captivity.

2\emph{nd Tuesday}
The voice of one crying in the
wilderness, Prepare ye the

way of the Lord, make the paths of our God straight.

2\emph{nd Wednesday}
O Sion, thou shalt be renewed,
and thou shalt see thy

Holy One, who is to come unto thee.

2\emph{nd Thursday}
He who shall come after me,
was in being before me:

whose sandals I am not worthy to unloose.

2\emph{nd Friday}
O sing unto the Lord a new
song, His praise is from the

ends of the earth.

2\emph{nd Saturday}
There was no God made before
Me, neither shall there

be after Me; for every knee shall bow to Me, and every

tongue confess to God.

3\emph{rd Monday}
All generations shall call me
blessed; because God hath

regarded the low estate of His handmaiden.

3\emph{rd Tuesday}
Awake, awake, stand up, O
Jerusalem, loose the chains

off thy neck, O captive daughter of Sion.

3\emph{rd Wednesday}
Behold the handmaid of the
Lord; be it unto me accord-

ing to thy word.

3\emph{rd Thursday}
Rejoice ye with Jerusalem, and
leap for joy in her all ye

that love her for ever.

3\emph{rd Friday}
This was the witness of John;
He that cometh after me,

was in being before me.

3\emph{rd Saturday}

(the 17th.)

[\emph{O Sapientia}.]

(\emph{This is the first of a series of Majores Antiphonæ, beginning on this day}. \emph{They are said whole before and after the Magnificat}.)

O eternal Wisdom, which proceedest from the
mouth of

the Most High, reaching from one end of creation unto

the other, mightily and harmoniously disposing all things,

come Thou teach us the way of understanding. 4\emph{th Monday} (19)
O Root of Jesse, who art
placed for a sign of the people,

before whom kings shall shut their mouths, whom the

Gentiles shall supplicate; come Thou to deliver us, do

not tarry.

(\emph{Vide} 4\emph{th Sunday for the
intervening Antiphon}.)

4\emph{th Tuesday} (20)
O Key of David and Sceptre of
the house of Israel, who

openest and none shutteth, who shuttest and none openeth,

come Thou, and bring forth the captive from the house of

bondage, who sitteth in darkness and in the shadow of

death.

4\emph{th Wed}. (21)
O rising Brightness of the
Everlasting Light and Sun

of Righteousness, come Thou and enlighten those who sit

in darkness and in the shadow of death.

4\emph{th Thurs}. (22)
O King and the Desire of all
nations, and chief Corner-

stone, who makest two to be one, come Thou and save man

whom Thou formedst from the clay.

4\emph{th Friday} (23)
O Emmanuel, our King and
Lawgiver, the gatherer of the

people and their Saviour, come Thou to save us, O

Lord our God.

\emph{Christmas Eve, \&c}.

(hhh) \emph{Collect}

(\emph{The same as on the respective Sundays: each Saturday be- longing to the Sunday follow- ing}.)

1\emph{st week}. Raise up, we beseech Thee, \&c.

2\emph{d week}. Stir up, O Lord, \&c.

3\emph{d week}. Give ear, we beseech Thee, \&c.

4\emph{th week}. Raise up, we beseech Thee, \&c.

OXFORD,

\emph{The Feast of St. John the Baptist}.

———————

ERRATUM.

In pages 35, 36, 41, 42, 45, 46, 63, 82, the words ``Jube,
Domine, benedicere''

are translated by mistake, ``Sir, pray for a blessing,'' instead of, ``Sir,
be pleased to bless us.''

[FOURTH EDITION.]

These Tracts are sold at the price of 2d. for each
sheet, (except No. 75) or 7s. for 50 copies.

LONDON: PRINTED FOR
J. G. F. \& J. RIVINGTON,

ST. PAUL'S CHURCH YARD, AND WATERLOO PLACE.

1840.

\xchapter{Catena Patrum No. II. Testimony of Writers in the later English Church to the Doctrine of Baptismal Regeneration}

Tract No. 76 (Ad Populum)

CONSIDERING
the confidence and zeal with which modern and unscriptural views on the
subject of Christian Baptism are put forth at the present time, it will
not be unseasonable to present the reader with some testimonies from the
writings of Anglican Divines in behalf of the doctrine of Baptismal
Regeneration. By this doctrine is meant, first, that the Sacrament of
Baptism is not a mere \emph{sign} or \emph{promise},
but actually a \emph{means} of grace, an \emph{instrument}, by which, when rightly received, the soul is admitted to the
benefits of CHRIST'S
Atonement, such as the forgiveness of sin, original and actual,
reconciliation to GOD,
a new nature, adoption, citizenship in CHRIST'S kingdom, and the inheritance of heaven,—in a word,
Regeneration. And next, Baptism is considered to be rightly received,
when there is no positive obstacle or hindrance to the reception in the
recipient, such as impenitence or unbelief would be in the case of an
adult; so that infants are necessarily right recipients of it, as not
being yet capable of actual sin. So much as these two positions is
certainly held by every one of the authors of the following passages,
though it is impossible to bring out their full meaning in such brief
extracts, however carefully selected.

There is a variety of questions connected with the
subject beyond the two positions above set down, on which the writers
under review differ more or less from each other, but not so as in the
slightest degree to interfere with their clear and deliberate
maintenance of these. Such, for instance, as the following:—Whether
grace be given in and through the water, or only contemporaneously with
it. Again, whether Baptism, strictly speaking, \emph{conveys} the
blessings annexed to it, or simply \emph{admits} into a state gifted
with those blessings, as being the initiatory rite of the covenant of
mercy. Or, again, whether or not Baptism, besides washing away past sin,
admits into a state in which, for sins henceforth committed, Repentance
stands in place of a Sacrament, so as to ensure forgiveness without
specific ordinance; or whether the Holy Eucharist is that ordinance; or
whether the full and explicit absolution of sin after Baptism is
altogether put off till the day of judgment. Or, again, there may be
difference of opinion as to the state of infants dying unbaptized. Or,
again, whether Regeneration is an instantaneous work completed in
Baptism, or admits of degrees and growth. Or, again, whether or not the
HOLY SPIRIT
can utterly desert a soul once inhabited by Him, except to quit it for
ever. Or, whether the change in the soul made by Baptism is indelible,
for good or for evil; or may be undone, as if it had never been. Or, how
far the enjoyment of the grace attached to it is suspended on the
condition of our doing our part in the covenant. All these are
questions, far from unimportant, but which do not at present come into
consideration; the one point, maintained in the following extracts,
being, that infants are by and at baptism unconditionally translated
from a state of wrath into a state of grace and acceptance for CHRIST'S
sake.

\textbf{List of Authors cited.}

1. Jewell
22. Scott

2. Hooker
23. Jenkin

3. Andrews
24. Sherlock

4. Donne
25. Wall

5. Field
26. Potter

6. Jackson
27. Nelson

7. Laud
28. Waterland

8. Bramhall
29. Kettlewell

9. Hammond
30. Hickes

10. Taylor
31. Johnson

11. Heylin
32. Leslie

12. Allestrie
33. Wilson

13. Barrow
34. Bingham

14. Thorndike
35. Skelton

15. Pearson
36. Horne

16. Bull
37. Jones

17. Comber
38. Heber

18. Ken
39. Jebb

19. Patrick
40. Van Mildert

20. Beveridge
41. Mant

21. Sharp

N. B. It would be easy to extend this list, were it
necessary: vid. Cosin's Devotions, Stanhope's Boyle Lectures,
\&c.

\textbf{JEWELL,
BISHOP.—Treatise
on Sacraments}

``They (the sacraments) are not bare signs; it
were blasphemy so to say. The grace of GOD
doth always work with His sacraments; but we are taught not to seek that
grace in the sign, but to assure ourselves by receiving the sign, that
it is given us by the thing signified. We are not washed from our sins
by the water, we are not fed to eternal life by the bread and wine, but
by the precious blood of our SAVIOUR
CHRIST,
that lieth hid in these sacraments.'' p. 263.

For this cause are infants baptized, because they
are born in sin, and cannot become spiritual but by this new birth of
the water and the Spirit. They are the heirs of the promise; the
covenant of GOD'S
favour is made unto them. GOD
said to Abraham, ``I will establish my covenant between me and thee,
and thy seed after thee.'' `` Therefore,'' saith the Apostle, ``If
the root be holy, so are the branches.'' And again, ``The unbelieving
husband is sanctified by the wife, and the unbelieving wife is
sanctified by the husband; else were your children unclean; but now are
they holy.'' When the disciples rebuked those that brought little
children to CHRIST,
that He might touch them, he said, ``Suffer the little children to come
unto me, and forbid them not, for of such is the kingdom of GOD.'' And again, ``Their angels
always behold the face of my Father which is in heaven,'' p. 265.

``The water wherein we are baptized doth not
cleanse the soul;'' but, ``the blood of JESUS CHRIST
His SON
doth cleanse us from all sin.'' Not the water, but the blood of CHRIST
reconcileth us unto GOD,
strengtheneth our conscience, and worketh our redemption. We must seek
salvation in CHRIST
alone, and not in any outward thing. Hereof saith Cyprian, ``Remissio peccatorum, sive per baptismum, sive per alia sacramenta
donetur, proprie Spiritus Sancti est. Verborum solemnitas,'' \&c. ``The remission of sins, whether it be given by baptism, or by any
other sacraments, does properly appertain to the HOLY
GHOST.
The solemnity of the words, and the invocation of GOD'S holy Name, and the outward signs
appointed to the ministry of the priest by the institution of the
Apostles, work the visible outward sacrament. But touching the substance
thereof, it is the HOLY
GHOST
that worketh it.'' St. Ambrose also saith, ``Vidisti fontem, vidisti
sacerdotem,'' \&c. ``Thou hast seen the water, thou hast seen the
priest, thou hast seen those things which thou mightest see with the
eyes of thy body, and with such sight as man hath: but those things
which work and do the deed of salvation, which no eye can see, thou hast
not seen.

``Such a change is made in the sacrament of
baptism. Through the power of GOD'S
working the water is turned into blood. They that be washed in it
receive the remission of sins; their robes are made clean in the blood
of the Lamb. The water itself is nothing; but by the working of GOD'S SPIRIT, the death and merits of our LORD
and SAVIOUR
CHRIST
are thereby assured unto us.

``A figure hereof was given at the Red Sea. The
children of Israel passed through in safety; but Pharaoh and his whole
army were drowned. Another figure hereof was given in the ark. The whole
world was drowned, but Noah and his family were saved alive. Even so in
the fountain of baptism, our spiritual Pharaoh, the devil, is choked:
his army, that is, our sins are drowned, and we are saved. The wicked of
the world are swallowed in concupiscence and vanities, and we abide safe
in the ark: GOD
hath chosen us to be a peculiar people to Himself; we walk not after the
flesh, but after the SPIRIT,
therefore we are in CHRIST
JESUS,
and there is now no condemnation unto us.

``Now touching the minister of this sacrament,
whether he be a good man or an evil man, godly or godless, an heretic or
a Catholic, an idolater or a true worshipper of GOD: the effect is

all one, the value or worthiness of the sacrament
dependeth not of man, but of GOD.
Man pronounceth the word, but GOD settleth our hearts with grace: man toucheth or washeth us
with water, but GOD
maketh us clean by the cross of CHRIST. It is not the minister, but CHRIST Himself which is the Lamb of GOD
that taketh away the sins of the world.'' p. 266.

\emph{Ibid}.—Reply to M. Harding's Censure, p.
249.

And forasmuch as these two sacraments, being both
of force like, these men [the Romanists] to advance their fantasies in
the one, by comparison so much abase the other: and specially for the
better opening of Chrysostom's mind, I think it good, briefly and by
the way, somewhat to touch what the old Catholic Fathers have written of
GOD'S invisible working in the Sacrament of
Baptism. Dionysius generally of all mysteries writeth thus: ``Angeli
Deum,'' \&c. ``The angels being creatures spiritual, so far forth
as it is lawful for them, behold GOD,
and his godly power. But we are led as we may, by sensible outward
tokens,'' (which he calleth images) ``unto the contemplation of
heavenly things.'' The Fathers, in the Council of Nice, say thus: ``Baptism must be considered, not with our bodily eyes, but with the
eyes of our mind. Thou seest the water; think thou of the power of GOD, that in the water is hidden. Think thou that the water is
full of heavenly fire, and of the sanctification of the HOLY
GHOST.''
Chrysostom, speaking likewise of baptism, saith thus: ``Ego non aspectu
judico ea, quæ videntur, sed mentis oculis,'' \&c. ``The things
that I see, I judge not by sight, but by the eyes of my mind. The
heathen, when he heareth the water of baptism, taketh it only for plain
water: but I see not simply, or barely, that I see; I see the cleansing
of the soul by the SPIRIT of GOD.'' So likewise saith Nazianzenus: ``Mysterium (baptismi) majus est, quam ea quæ videntur;''
``The
mystery of baptism is greater than it appeareth to the eye.'' So St.
Ambrose: ``Aliud est, quod visibiliter agitur: aliud quod invisibiliter
celebratur:'' ``In baptism there is one thing done visibly to the eye;
another thing is wrought invisibly to the mind.''

Again he saith, ``Believe not only the bodily eyes
(in this sacrament of baptism): the thing that is not seen, is better
seen: the thing that thou seest is corruptible; the thing that thou
seest not is for ever.'' To be short, in consideration of these
invisible effects, Tertullian saith, ``The HOLY
GHOST
cometh down and halloweth the water.'' St. Basil saith, ``The kingdom
of heaven is there set open.'' Chrysostom saith, ``GOD Himself in baptism, by his invisible power, holdeth thy
head.'' St. Ambrose saith, ``The water hath the grace of CHRIST; in it is the presence of the
Trinity.'' St. Bernard saith, ``Let us be washed in his blood.''

By the authorities of thus many ancient Fathers, it
is plain, that in the sacrament of baptism, by the sensible sign of
water, the invisible grace of GOD
is given unto us.''

\textbf{HOOKER,
PRESBYTER
AND DOCTOR.—On
Ecclesiastical Polity},

Book v. 60.

Unless as the SPIRIT is a necessary inward cause, so water were a necessary
outward mean, to our regeneration, what construction should we give unto
those words wherein we are said to be new born, and that [\emph{ex hudatos}],
even of water? Why are we taught, that with water GOD doth purify and cleanse His Church?
Wherefore do the Apostles of CHRIST term baptism a bath of regeneration? What purpose had they
in giving men advice to receive outward baptism, and in persuading them
it did avail to remission of sins? If outward baptism were a cause in
itself possessed of that power, either natural or supernatural, without
the present operation whereof no such effect could possibly grow, it
must then follow, that seeing effects do never prevent the necessary
causes out of which they spring, no man could ever receive grace before
baptism: which being apparently both known, and also confessed to be
otherwise, in many particulars, although in the rest we make not baptism
a cause of grace; yet tile grace which is given them with their baptism,
doth so far forth depend on the very outward sacrament, that GOD
will have it embraced, not only as a sign or token what we receive, but
also as an instrument or means whereby we receive grace, because
baptism is a sacrament which GOD hath instituted in His Church, to the end that they which
receive the same might thereby be incorporated into CHRIST; and so through His most precious
merit obtain, as well that saving grace of imputation which taketh away
all former guiltiness, as also that infused divine virtue of the HOLY
GHOST
which giveth to the powers of the soul their first disposition towards
future newness of life. There are that elevate too much the ordinary and
immediate means of life, relying wholly upon the bare conceit of that
eternal election, which notwithstanding includeth a subordination of
means, without which we are not actually brought to enjoy what GOD secretly did intend; and therefore, to build upon GOD'S
election, if we keep not ourselves to the ways which He hath appointed
for men to walk in, is but a self-deceiving vanity. When the Apostle saw
men called to the participation of JESUS
CHRIST,
after the Gospel of GOD
embraced, and the sacrament of life received, he feareth not then to put
them in the number of elect saints; he then accounteth them delivered
from death, and clean purged from all sin. Till then, notwithstanding
their preordination unto life, which none could know of, saving GOD, what were they, in the Apostle's own
account, but children of wrath, as well as others, plain aliens,
altogether without hope, strangers, utterly without GOD in this present world? So that by
sacraments, and other sensible tokens of grace, we may boldly gather,
that He whose mercy vouchsafeth now to bestow the means, hath also
sithence intended us that whereunto they lead. But let us never think it
safe to presume of our own last, and by bare conjectural collections of
his first intent and purpose, the means failing that should come
between. Predestination bringeth not to life without the grace of
eternal vocation, wherein our baptism is implied. For as we are not
naturally men without birth, so neither are we Christian men in the eye
of the Church of GOD
but by new birth; nor according to the manifest ordinary course of
divine dispensation new born, but by that baptism which both declareth
and maketh us Christians. In which respect, we justly hold it to be
the door of our actual entrance into GOD'S
house, the first apparent beginning of life, a seal perhaps to the grace
of election before received: but to our sanctification here, a step that
hath not any before it.

\emph{Ibid}. 64.

Were St. Augustine now living, there are which
would tell him for his better instruction, that to say of a child, it is
elect, and to say, it doth believe, are all one; for which cause, sith
no man is able precisely to affirm the one of any infant in particular,
it followeth that precisely and absolutely we ought not to say the
other. Which precise and absolute terms are needless in this case. We
speak of infants as the rule of piety alloweth both to speak and think.
They that can take to themselves in ordinary talk, a charitable kind of
liberty to name men of their own sort GOD'S
dear children, (notwithstanding the large reign of hypocrisy,) should
not, methinks, be so strict and rigorous against the Church, for
presuming as it doth of a Christian innocent. For when we know how CHRIST in general hath said, Of such is the kingdom of heaven,
which kingdom is the inheritance of GOD'S elect; and do withal behold how His providence hath called
them unto the first beginnings of eternal life, and presented them at
the well-spring of new birth, wherein original sin is purged, besides
which sin there is no hindrance of their salvation known to us, as
themselves will grant; hard it were that having so many fair inducements
whereupon to ground, we should not be thought to utter, at the least, a
truth as probable and allowable in terming any such particular infant an
elect babe, as in presuming the like of others whose safety,
nevertheless, we are not absolutely able to warrant.

\textbf{ANDREWS,
BISHOP AND
DOCTOR.—On
the Holy Ghost},

Serm. vii.

Now CHRIST
is baptized. And no sooner is He so, but He falls to His prayers, \emph{Indigentia
mater orationis}, (we say) want begets prayer: therefore, yet there
wants somewhat—a part and that a chief part of baptism is still
behind.

There goes more to baptism, if it be as it should
be, than \emph{baptismus fluminis}, yea (I may boldly say,) there
goes more to it, if it be as it should, than \emph{baptismus sanguinis}.
CHRIST
``came in water and blood, not in water only, but in water and
blood:'' that is not enough, except the ``SPIRIT
also bear witness.'' So \emph{baptismus flaminis} is to come too. There
is to be a Trinity beneath,—l. water, 2. blood, and 3. the SPIRIT, to answer to that above: but (the SPIRIT'S
baptism coming too) in the mouth of all three, all is made sure, all
established thoroughly. This is it, He prays for, as man.

For the baptism of blood that was due to every one
of us, (and each of us to have been baptized in his own blood, to have
had three such immersions?) that hath CHRIST
quit us of when he was asked by the prophet, ``How his robes came so
red?'' He says, ``He had been in the wine-press;'' but there He had
been, and that ``He bad trod alone, and not one of the people with
Him;'' none but He there; in that, spare us in that.

But the other two parts He sets down precisely to
Nicodemus (and in him, to us all,)—l. water, 2. and the HOLY GHOST …

St. Paul tells us (Col. ii.) that besides the
circumcision, that was the manufacture, there was another made without
hands. There is so in baptism, besides the hand seen, that casts on the
water, the virtue of the HOLY
GHOST is
there, working, without hands, what here was wrought.

And for this CHRIST prays; that then it might, might then, and might ever, be
joined to that of the water. Not in his baptism only, but in the
people's; and (as he afterwards enlarges His prayer) in all others
that ``should ever after believe in His name:'' that what in His
(here) was, in all theirs might be: what in this first, in all
following; what in CHRIST'S, in all Christians; heaven might
open, the HOLY
GHOST
come down, the FATHER
be pleased to say over the same word, \emph{toties quoties}, so oft as
any Christian man's child is brought to his baptism. CHRIST hath prayed, now,

See the force of His prayer. Before it heaven was
mured up, no dove to be seen, no voice to be heard, \emph{Altum silentium}.
But straight upon it (as if they had but waited the last word of His
prayer) all of them follow immediately.

Heaven opens, first. For, if when the lower heaven
was shut three years, Elias was able with his prayer to open it, (it is
our SAVIOUR
in the next chapter following,) and bring down rain; the prayer of CHRIST (who is more of might than many
such as Elias) shall it not be much more of force, to enter the Heaven
of heavens, the highest of them all, and to bring down thence the waters
``above the heavens,'' even the heavenly graces of the HOLY SPIRIT?

For, so when our SAVIOUR cried, (John vii.) ``If any thirst,'' \&c. This (saith
St. John) He spake of the SPIRIT.'' For the SPIRIT and His graces are the very supercelestial water; one drop
whereof, infused into the waters of Jordan, will give them an admirable
power to pierce even into the innermost parts of the soul: and to
baptize it, (that is) not only to take out the stains of it, and make it
clean; but further, give it a tincture, lustre, or gloss; for so is
baptism properly of [\emph{bapt\={o}}], taken from the dyer's fat,
and is a dyeing or giving a fresh colour, and not a bare washing only.

Always, the opening of heaven, opens unto us, that
no baptism without heaven open: and so, that baptism is \emph{de cœlo, non
ab hominibus}, from heaven, not of men. So it was here; so is it to
be holden for ever. 2. And from heaven; not \emph{clanculum} (as
Prometheus is said to get his fire), but [\emph{ane\={o}ichth\={e}nai}],
orderly, by a fair door set open, in the view of much people; for all
that were present saw the impression in the sky. Which door was not
mured up again; for we find it still open, (Apoc. iii.) and we find that
keys were made, and given of it, after this. 3. And all this, that there
might not only be a passage for these down, but for us up. For heaven
gate, \emph{ab hoc exemplo}, doth ever open at baptism in sign, he that
new cometh from the fount hath then right of entrance in thither. Then
(I say) when by baptism he is cleansed; for, before, \emph{Nihil inquinatum},
nothing defiled can enter there.

\textbf{DONNE,
PRESBYTER.—Serm.
xxxi.}

p. 309.

The water of Baptism, is the water that runs
through all the Fathers; all the Fathers that had occasion to dive or
dip in these waters (to say anything of them) make these first waters,
in the creation, the figure of baptism. Therefore Tertullian makes the
water, \emph{Primam sedem Spiritus Sancti}, the progress, and the
settled house, the voyage, and the harbour, the circumference, and the
centre of the HOLY GHOST.
And therefore St. Hierome calls these waters, \emph{Matrem mundi}, the
Mother of the world: and this in the figure of baptism. The waters
brought forth the whole world, were delivered of the whole world, as a
mother is delivered of a child; and this, \emph{in figura baptismi}, to
foreshew that the waters also should bring forth the Church; that the
Church of GOD should be born of the Sacrament of Baptism.
So says Damascen, and he establishes it with better authority than his
own. The divine Basil saith (saith he) ``The SPIRIT of GOD wrought upon the waters in the creation,
because he meant to do so after, in the regeneration of man. And
therefore, \emph{Pristinam sedem recognoscens conquiescit}, till the HOLY GHOST have moved upon our children in baptism,
let us not think all done that belongs to those children; and when the HOLY
GHOST
hath moved upon those waters, so in baptism, let us not doubt of His
power and effect upon all those children that die so. We know no means
how those waters could have produced a minnow, a shrimp, without the SPIRIT of GOD had moved upon them; and by this motion of
the SPIRIT
of GOD,
we know they produce whales, and leviathans. We know no ordinary means
of any saving grace for a child but baptism; neither are we to doubt of
the fulness of salvation, in them that have received it. And for
ourselves, \emph{mergimur et emergimus}, in baptism we are sunk under
water, and then raised above the water again; which was the manner of
baptizing in the Christian Church, by immersion and not by aspersion,
till of late times: \emph{Affectus et amores}, (says he,) our corrupt
affections, and our inordinate love of this world is that, that is to be
drowned in us; \emph{Amor securitatis}, a love of peace, and holy
assurance, and acquiescence in GOD'S ordinance, is that that lifts us above water.

Therefore that Father puts all upon the due
consideration of our baptism: and as St. Jerome says, Certainly he that
thinks upon the last Judgment advisedly, cannot sin thus; so he that
says with St. Augustine, Let me make every day to GOD,
this confession; \emph{Domine}, \&c. O LORD my GOD, O holy, holy, holy LORD
my GOD;
I consider that I was baptized in thy name, and what thou promised me,
and what I promised thee then, and can I sin this sin? can this sin
stand with those conditions, those stipulations which passed between us
then? The SPIRIT of GOD is motion, the SPIRIT of GOD is rest too; and in due consideration of
baptism, a true Christian is moved, and settled too; moved to a sense of
the breach of his conditions, settled in the sense of the mercy of his GOD,
in the merits of his CHRIST,
upon his godly sorrow. So these waters are the waters of baptism.

\textbf{FIELD,
PRESBYTER.—Of
the Church},

book i. chap. xii.

This was the fault of sundry in the Primitive
Church; and which was yet more to be condemned, many did therefore defer
and put off their baptism, that so whatsoever evil things they did in
the mean time, might in that laver of new birth be washed away, thereby
taking greater liberty to offend, for that they had so present means of
full remission, and perfect reconciliation; so making that which was
ordained against sin, and for the weakening and overthrow of it, to be
an encouragement thereunto, and to give life and strength unto it.

\textbf{JACKSON,
PRESBYTER
AND DOCTOR.—On
CHRIST'S
exercise of his everlasting Priesthood},

ch. i. (vol. iii. p. 271.)

It is no part of our Church's doctrine or
meaning, that the washing, or sprinkling infants' bodies with
consecrated water, should take away sins by its own immediate virtue. To
affirm thus much implies, as I conceive, a contradiction to that
apostolical doctrine. ``The like figure whereunto even Baptism doth
also now save us (not the putting away the filth of the flesh, but
the answer of a good conscience towards GOD)
by the resurrection of JESUS
CHRIST,
who is gone into heaven,'' \&c. 1 Pet. iii. 21. The meaning of our
Church intends no further than thus: That if this sacrament of Baptism
be duly administered, the blood, or bloody sacrifice of CHRIST, or (which is all one) the influence of His SPIRIT
doth always accompany, or is concurrent to this solemn act. But whether
this influence of His Spirit or virtual presence of His body and blood
be either immediately or only terminated to the soul and spirit of the
party baptized, or have some virtual influence upon the water of Baptism
as a mean to convey the Grace of Regeneration unto the soul of the party
baptized, whilst the water is poured upon him, is too nice and curious a
question, in this age, for sober Christians to debate or contend about.
It may suffice to believe that this sacramental pledge hath a virtual
presence of CHRIST'S Blood, or some real influence from His Body,
concomitant, though not consubstantiated to it, which is prefigured or
signified by the washing or sprinkling the body with water.

But it will be, or rather is objected, but only by
private or some saucy spirits, That if the doctrine of our Church were
true and sound, then all that be rightly baptized should be undoubtedly
saved, being once washed or cleansed from their sins. The objection were
of some force, if the Church of England did hold or maintain such
doctrine or tenets as they do which make or favour it; to wit, That the
sins of the elect only are remitted by Baptism, or by Sacrament of CHRIST'S Body and Blood; or, that sins once remitted cannot be
remitted afresh; or, that the party which is once pardoned for his sins,
before committed, cannot afterwards be condemned. The orthodoxal truth
is, That albeit the original sin of children truly baptized in the name
of CHRIST,
or the actual sins of young or elder men so baptized, and the sins of
their forefathers (so far as it concerns men of riper years to repent of
them both) be so truly remitted in Baptism, that neither young men nor
old may be baptized again; yet the stipulation of a good conscience,
wherein the internal Baptism (as St. Peter tells) doth consist, may and
ought, by the law of GOD
and of CHRIST'S
Church, to be reiterated.

And this stipulation of every Christian, male or
female, though baptized after they have passed their nonage for civil
contracts, ought to be resumed or reacknowledged as often as they intend
to receive the sacramental pledges of CHRIST'S
Body and Blood, either privately or in the public congregation. But for
all such as have been baptized in their infancy, the personal resumption
or ratification of that vow, which their fathers and mothers in GOD
did make for them at the sacred laver, is to be exacted of them \emph{ore
tenus}, in some public congregation, before they can be lawfully
admitted to be public communicants of CHRIST'S
Body and Blood.

\emph{Ibid}.—Ch. lv. (p. 298.)

If either the actual sins of all men, or the sins
of the elect in special, had been so remitted by CHRIST'S death, as some conceive they were,
that is, absolutely pardoned before they were committed, there had been
no end or use of CHRIST'S
Resurrection in respect of us; no need of Baptism: yet was Baptism, from
the hour of His resurrection, necessary unto all that did believe in His
death and resurrection. The urgent and indispensable necessity of
Baptism, especially in respect of actual believers, is not anywhere more
emphatically intimated than in St. Peter's answer to the Jews, whose
hearts were pierced with sorrow that they had been the causes of CHRIST'S death. They in this stound or sting of conscience
demand, ``Men and brethren, what shall we do? And Peter answered them,
Repent, and be baptized every one of you, in the name of Jesus CHRIST, for the remission of sins. And
they that gladly received the word were baptized the same day.'' Acts
ii. 37, 38, 41. These men had been deeply tainted with sin, not original
only, but with sins actual of the worst kind; guilty they were, in a
high degree, of the death of the SON
of GOD,
yet had they as well their actual as their original sins remitted by
Baptism. It is then unsound and imperfect doctrine, that sin original
only is taken away or remitted by Baptism; for whatsoever sins are
remitted or taken away by CHRIST'S
death, the same sins are in the same manner remitted and taken away by
Baptism into His death; actual sins are remitted, in such as are
guilty of actual sins when they are baptized, though only sin original
be actually remitted in those which are not guilty of actual sins, as in
infants. No man's sins are actually remitted before he be actually
guilty of them.

The question is, how either sin original is
remitted, or how any work of Satan is dissolved by Baptism; and this
question, in the general, is rightly resolved, by saying, ``They are
remitted by faith.'' But this general resolution sufficeth not, unless
we know the object of our faith in this particular. Now the particular
object of our faith, of that faith by which sins (whether by Baptism or
otherwise) are remitted, is not our general belief in CHRIST;
even our belief in CHRIST
dying for us in particular, will not suffice, unless it include our
belief of the everlasting virtue of His bloody sacrifice, and of His
everlasting priesthood for purifying and cleansing our souls. No sins be
truly remitted unless they be remitted by the office or exercise of His
priesthood and whilst so remitted they are not remitted by any other
sacrifice than by the sole virtue of His body and blood, which He ``once offered for all,'' for the sins of all. It is not the virtue or
efficacy of the consecrated water in which we were washed, but the
virtue of his blood which was once shed for us, and which, by Baptism,
is sprinkled upon us, or communicated unto us, which immediately
cleansed us from all our sins. From this everlasting virtue of this His
bloody sacrifice, faith, by the ministry of Baptism, is immediately
gotten in such as had it not before. And in such as have faith before
they be baptized, the guilt of actual sins is remitted by the exercise
or act of faith, as it apprehends the everlasting efficacy of this
sacrifice, and by the prayer of faith, and supplication unto our High
Priest. Faith, then, is as the mouth or appetite by which we receive
this food of life, and is a good sign of health; but it is the food
itself received, which must continue health and strengthen spiritual
life in us; and the food of life is no other than CHRIST'S body and blood; and it is our High Priest Himself which
must give us this food.

Baptism, saith St. Peter, (1 Pet. iii. 20.) doth
save us. What Baptism doth save us? not the putting away the filth
of the flesh, (yet this is the immediate effect of the water in
Baptism,) but the answer (or stipulation) of a good conscience towards GOD! But how doth this kind of Baptism, or this
concomitant of Baptism save us? The Apostle, in the same place, tells
us, ``by the resurrection of JESUS
CHRIST.''
``The answer or stipulation of a good conscience,'' includes an
illumination of our Spirits by the SPIRIT
of GOD;
a qualification by which we are made sons of light, being before the
sons of darkness. But, that by this qualification we become the sons of
light; that this qualification is, by Baptism, wrought in us; that by
this qualification, however wrought in us, we are saved from our sins;
all this is immediately from the ``virtue of CHRIST'S
Resurrection.'' That is, as you have heard before, He was consecrated
by the sufferings of death to be an everlasting Priest, and by His
resurrection from death, His body and blood became an everlasting
propitiation for sins, an inexhaustible fountain of grace, by which we
are purified from the dead works of sin.

\emph{Ibid}.—Of Christ's session at the right
hand of God. ch. xvii. p. 170.

[St. Paul] saith, ``that all that are baptized are
dead to sin;'' that is, first, they are ``dead unto it by solemn vow
or profession.'' Secondly, they are said to be ``dead unto sin, or sin
to be dead in them,'' inasmuch as they in Baptism receive an antidote
from GOD by which the rage and poison of it might
easily be assuaged or expelled, so they would not either receive that
grace or means which GOD
in Baptism exhibits unto them in vain, or use it amiss. So we may say
that any popular disease is quelled or taken away, after a sovereign
remedy be found against it, which never fails; so men will seek for it,
reasonably apply for it, and observe that diet which the physician, upon
the taking of it, prescribes unto them. Some in our times there be (and
more, I think, than have been in all the former) which deny all
baptismal grace. Others there be which grant some grace to be conferred
by Baptism, even unto infants; but yet these restrain it only to infants
elect. And this they take to be the meaning of our Church's Catechism,
wherein children are taught to believe [That as CHRIST, the second person in the Trinity, did redeem them and all
mankind; so the HOLY
GHOST
(the third person) doth sanctify them, and all the elect people of GOD.]

But can any man be persuaded that it was any part
of our Church's meaning, to teach children when they first make
profession of their faith, to believe, that they are of the number of
the elect; that is, of ``such as cannot finally perish?'' This were to
teach them their faith backwards, and to seek the kingdom of heaven not \emph{ascendendo},
by ascending, but \emph{descendendo}, by descending from it. For higher
than thus St. Paul himself, in his greatest perfection, could not
possibly reach; no, nor the blessed angels, which have kept their first
station almost these 6000 years. Yet certain it is, that our Church
would have every one, at the very first profession of his faith, to
believe that he is one of the elect people of GOD.

\textbf{LAUD,
ARCHBISHOP
AND MARTYR.—Conference
with Fisher, § 15}.

First, that Baptism is necessary to the salvation
of infants (in the ordinary way of the Church, without binding GOD to the use and means of that
sacrament, to which He hath bound us) is expressed in St. John iii. ``Except a man be born of water,'' \&c. So, no baptism, no
entrance. Nor can infants creep in any other ordinary way. And this is
the received opinion of all the ancient Church of CHRIST.

And, secondly, that infants ought to be baptised,
is first, plain by evident and direct consequence out of Scripture. For
if there be no salvation for infants in the ordinary way of the Church
but by Baptism, and this appear in Scripture, as it doth, then out of
all doubt, the consequence is most evident out of that Scripture, that
infants are to be baptized, that their salvation may be certain. For
they which cannot help themselves, must not be left only to
extraordinary helps, of which we have no assurance, and for which we
have no warrant at all in Scripture, while we, in the mean time, neglect
the ordinary way and means commanded by CHRIST. Secondly, it is very near an expression in Scripture
itself. For when St. Peter had ended that great Sermon of his, Acts ii.,
he applies two comforts unto them, verse 38, ``Amend your life,''
\&c. And then, v. 39, He infers, ``For the promise is made,''
\&c. The promise; what promise? What? why the promise of
sanctification by the HOLY
GHOST.
By what means; Why, by Baptism. For it is expressly, ``Be baptized, and
ye shall receive.'' And as expressly, ``This promise is made to you
and to your children.''

\textbf{BRAMHALL,
ARCHBISHOP
AND CONFESSOR.—Of
persons dying without Baptism},

p. 979.

The discourse which happened the other day, about
your little daughter, I had quite forgotten till you were pleased to
mention it again last night. If any thing did fall from me, which gave
offence to any there present, I am right sorrowful, but I hope there did
not; as, on the other side, if any occasion of offence had been given to
me, I should readily have sacrificed it to that reverend respect, which
is due to the place—your table,—anciently accounted a sacred thing,
and to the lord of it, yourself. This morning, lying musing in my bed,
it produced some trouble to me, to consider how passionately we are all
wedded to our own parties, and how apt we are all to censure the
opinions of others before we understand them, while, our want of charity
is a greater

error in ourselves, and more displeasing to ALMIGHTY
GOD,
than any of those supposed assertions which we condemn in others,
especially when they come to be rightly understood. And to show thus
particular breach is not so wide, nor the more moderate of either party
so disagreeing, as is imagined, I digested these sudden meditations,
drawn wholly, in a manner from the grounds of the Roman schools; and so
soon as I was risen, I committed them to writing.

First, there is a great difference to be made
between the sole want of Baptism upon invincible necessity, and the
contempt or wilful neglect of Baptism when it may be had. The latter we
acknowledge to be a damnable sin, and, without repentance and GOD'S extraordinary mercy, to exclude a man
from all hope of salvation. But yet if such a person, before his death,
shall repent and deplore his neglect of the means of grace, from his
heart, and desire, with all his soul, to be baptized, but is debarred
from it invincibly, we do not, we dare not pass sentence of
condemnation upon him; nor yet the Roman Catholics themselves. The
question then is, whether the want of Baptism, upon invincible
necessity, do evermore infallibly exclude from heaven?

Secondly, we distinguish between the visible sign,
and the invisible grace; between the exterior sacramental ablution, and
the grace of the sacrament, that is, interior regeneration. We believe
that whosoever hath the former, hath the latter also, so that he do not
put a bar against the efficacy of the sacrament by his infidelity or
hypocrisy, of which a child is not capable. And therefore our very
Liturgy doth teach, that a child baptized, dying before the commission
of actual sin, is undoubtedly saved.

Thirdly, we believe that without baptismal grace,
that is, regeneration, no man can enter into the kingdom of GOD. But whether GOD
hath so tied and bound himself to His ordinances and sacraments, that He
doth not or cannot confer the grace of the sacraments, extraordinarily,
where it seemeth good to His eyes, without the outward element: this is
the question between us.

\textbf{HAMMOND,
PRESBYTER,
CONFESSOR,
AND DOCTOR.— Sermon XV.—A New Creature}.

It is observable, that our state of nature and sin
is, in Scripture, expressed ordinarily by old age, the natural sinful
man; that is, all our natural affections that are born and grow up with
us, are called the old man; as if, since Adam's fall, we were decrepit
and feeble, and aged as soon as born, as a child begotten by a man in a
consumption never comes to the strength of a man, is always weak, and
crazy, and puling, hath all the imperfections and corporal infirmities
of age before he is out of his infancy. And, according to this ground,
the whole analogy of Scripture runs; all that is opposite to the old
decrepit state, to the dotage of nature, is new. The new covenant, Mark
i. 27. The language of believers, new tongues, Mark xvi. 17. A new
commandment, John xiii. 34. A new man, Eph. ii. 15. In sum, the state of
grace is expressed by [\emph{panta kaina}], all is become new, 2 Cor. v.
17. So that old and new, as it divides the Bible, the whole state
of things, the world; so it doth that to which all these serve, man;
every natural man, which hath nothing but nature in him, is an old man,
be he never so young, is full of years, even before he is able to tell
them. Adam was a perfect man when he was but a minute old, and all his
children are old even in the cradle, nay, even dead with old age, Eph.
ii. 5. And, then, consequently, every spiritual man, which hath somewhat
else in him than he receiveth from Adam, he that is born from above,
John iii. 3, [\emph{genn\={e}th\={e}i an\={o}then}], (for it
may be so rendered from the original, as well as born again, as our
English read it), he that is by GOD'S
SPIRIT
quickened from the old death, Eph. ii. 5, he is, contrary to the former,
a new man, a new creature; the old eagle hath cast his beak and is grown
young; the man, when old, hath entered the second time into his
mother's womb, and is born again; all the grey hairs and wrinkles fall
off from him, as the scales from blind Tobit's eyes, and he comes
forth a refined, glorious, beauteous new creature: you would wonder to
see the change. So that you find, in general, that the Scripture
presumes it, that there is a renovation, a casting away the old coat, a
youth and spring again in many men, from the old age and weak bedrid
state of nature. Now that you may conceive wherein it consists, how this
new man is brought forth in us, by whom it is conceived, and in what
womb it is carried, I will require no more of you, than to observe and
understand with me, what is meant by the ordinary phrase in our divines,
a new principle, or inward principle of life, and that you shall do
briefly thus. A man's body is naturally a sluggish, inactive,
motionless, heavy thing, not able to stir or move the least animal
motion, without a soul to enliven it; without that, it is but a carcase,
as you see at death, when the soul is separated from it, it returns to
be but a stock or lump of flesh; the soul bestows all life and motion on
it, and enables it to perform any work of nature. Again, the body and
soul together, considered in relation to somewhat above their power and
activity, are as impotent and as motionless as before the body without
the soul. Set a man to remove a mountain, and he will heave, perhaps, to
obey your command, but in event will do no more towards the displacing of it, than a stone in the street could do: but now, let an
omnipotent power be annexed to this man, let a supernatural spirit be
joined to this soul, and then will it be able to overcome the proudest,
stoutest difficulty in nature. You have heard, in the Primitive Church,
of a grain of faith removing mountains; and believe me, all miracles are
not yet outdated. The work of regeneration, the bestowing of a spiritual
life on one dead in trespasses and sins, the making of a carcase walk,
the natural old man to spring again, and move spiritually, is as great a
miracle as that …

For the third question, when this new principle
enters: first you are to know, that it comes into the heart in a
threefold condition; first, as an harbinger; secondly, as a private
secret guest; thirdly, as an inhabitant or housekeeper. As it is an
harbinger, so it comes to fit and prepare us for itself; trims up, and
sweeps, and sweetens the soul, that it may be readier to entertain him
when he comes to reside; and that he doth (as the ancient gladiators had
their \emph{arma prælusoria},) by skirmishing with our corruptions,
before he comes to give them a pitched battle; he brandishes a flaming
sword about our ears, and as by a flash of lightning, gives us a sense
of a dismal, hideous state; and so somewhat restrains us from excess and
fury; first, by a momentary remorse, then by a more lasting, yet not
purifying flame, the spirit of bondage. In sum, every check of
conscience, every sigh for sin, every fear of judgment, every desire of
grace, every motion or inclination toward spiritual good, be it ever so
short-winded, is \emph{praæludium spiritus}, a kind of John Baptist to
CHRIST, something that GOD sent before ``to prepare the ways of the LORD.''
And thus the SPIRIT
comes very often; in every affliction, every disease, (which is part of
GOD'S
discipline, to keep us in order,) in brief, at every sermon that works
upon us at the hearing: then, I say, the lightning flashes in our eyes;
we have a glimpse of His SPIRIT,
but cannot come to a full sight of it: and thus He appears to many, whom
He will never dwell with. Unhappy men, that cannot lay hold on Him, when
He comes so near them! and yet somewhat more happy than they that never
came within ken of Him; stop their ears when He spake to them even
at this distance. Every man in the Christian Church hath frequently, in
his life, a power to partake of GOD'S
ordinary preparing graces: and it is some degree of obedience, though no
work of regeneration, to make good use of them; and if he without the
inhabitance of the SPIRIT,
cannot make such use as he should, yet to make the best he can: and
thus, I say, [i.e. in a \emph{parallel} way] the SPIRIT appears to the unregenerate, almost
every day of our lives. 2ndly, when this SPIRIT comes a guest to lodge with us, then He
is said to enter; but till by actions and frequent obliging works, he
makes himself known to his neighbours, as long as he keeps his chamber,
till he declare himself to be there, as long he remains a private secret
guest, and that is called the introduction of the form, that makes a man
to be truly regenerate; when the seed is sown in his heart, when the
habit is infused, and that is done sometimes discernibly, sometimes not
discernibly, but seldom, as when Saul was called in the midst of his
madness, Acts ix., he was certainly able to tell a man the very minute
of his change, of his being made a new creature. Thus they which have
long lived in an enormous Antichristian course, do many times find
themselves strucken on a sudden, and are able to date their
regeneration, and tell you punctually how old they are in the SPIRIT.
Yet because there be many preparations to this SPIRIT, which are not this SPIRIT, many presumptions in our hearts false
grounded, many tremblings and jealousies in those that have it, great
affinity between faith natural and spiritual: seeing it is a SPIRIT that thus enters, and not as it
did light on the Disciples, in a bodily shape, it is not an easy matter
for any one to define the time of his conversion. Some may guess
somewhat nearer than others, as remembering a sensible change in
themselves; but, in a word, the surest discerning of it is in its
working, not at its entering. I may know that now I have the SPIRIT,
better than at what time I came to it. Undiscernibly GOD'S supernatural agency interposes
sometimes in the mother's womb, as in John Baptist springing in
Elizabeth at Mary's salutation, (Luke i. 41.) and perhaps in Jeremy, (Jer.
i. 5.) ``Before thou camest out of the womb, I sanctified thee;'' and
(in Isa. xlix. 5.) ``The LORD
that formed me from the womb to be his servant.'' But this
divine address attends most ordinarily till the time of our Baptism,
when the SPIRIT,
accompanying the outward sign, infuses itself into their hearts, and
there seats and plants itself, and grows up with the reasonable soul,
keeping even their most luxuriant years within bounds; and as they come
to an use of their reason, to a more and more multiplying this habit of
grace into holy spiritual acts of faith and obedience; from which it is
ordinarily said, that infants baptized have habitual faith, as they may
be also said to have habitual repentance, and habits of all other
graces, because they have the root and seed of those beauteous,
healthful flowers, which will actually flourish there when they come to
years. And this, I say, is so frequent to be performed at Baptism, that
ordinarily it is not wrought without that means, and in those means we
may expect it, as our Church doth in our Liturgy, where she presumes, at
every Baptism, that it hath pleased GOD to regenerate the infant by his HOLY SPIRIT. And this may prove a solemn piece of
comfort to some, who suspect their state more than they need, and think
it impossible that they should be in a regenerate condition, because
they have not as yet found any such notable change in themselves, as
they see and observe in others. These men may as well be jealous they
are not men, because they cannot remember when their soul came to them:
if they can find the effects of spiritual life in themselves, let them
call it what they will, a religious education, or a custom of
well-doing, or an unacquaintedness with sin; let them comfort themselves
in their estate, and be thankful to GOD who visited them thus betimes; let
it never trouble them that they were not once as bad as other men, but
rather acknowledge GOD'S
mercy, who had prevented such a change, and by uniting them to Him in
the cradle, hath educated and nursed them up in familiarity with the SPIRIT.

\textbf{TAYLOR,
BISHOP,
CONFESSOR,
AND DOCTOR.—Life
of Christ},

sect. 9.—On Baptism, part ii. 16.

Thirdly, in baptism we are born again; and this
infants need in the present circumstances, and for the same great reason
that men of age and reason do. For our natural birth, is either of itself insufficient, or is made so by the fall of Adam, and the
consequent evils, that nature alone, or our first birth, cannot bring us
to heaven, which is a supernatural end, that is, an end above all the
power of our nature as now it is. So that if nature cannot bring us to
heaven, grace must, or we can never get thither; if the first birth
cannot, a second must: but the second birth spoken of in Scripture is
baptism; ``a man must be born of water and the Spirit.'' And therefore
baptism is [\emph{loutron palingenesias}], ``the laver of a new
birth.'' Either then infants cannot go to heaven any way that we know
of, or they must be baptized. To say they are left to GOD,
is an excuse, and no answer; for when GOD hath opened the door, and calls that the ``entrance into heaven,'' we do not leave them to GOD, when we will not carry them to him in the
way which He hath described, and at the door which Himself hath opened:
we leave them indeed, but it is but helpless and destitute: and though GOD
is better than man, yet that is no warrant to us; what it will be to the
children, that we cannot warrant or conjecture. And if it be objected,
that to the new birth are required dispositions of our own, which are to
be wrought by and in them that have the use of reason; besides that,
this is wholly against the analogy of a new birth, in which the person
to be born is wholly a passive, and hath put into him the principle that
in time will produce its proper actions; it is certain that they that
can receive the new birth, are capable of it. The effect of it is a
possibility of being saved, and arriving to a supernatural felicity. If
infants can receive this effect, then also the new birth, without which
they cannot receive the effect. And if they can receive salvation, the
effect of the new birth, what hinders them but they may receive that,
that is in order to that effect, and ordained only for it, and which is
nothing of itself, but in its institution and relation, and which may be
received by the same capacity, in which one may be created, that is, a
passivity, or a capacity obediential?

Fourthly; concerning pardon of sins, which is one
great effect of baptism, it is certain that infants have not that
benefit, which men of sin and age may receive. He that hath a sickly
stomach, drinks wine, and it not only refreshes his spirits, but
cures his stomach: he that drinks wine, and hath not that disease,
receives good by his wine, though it does not minister to so many needs;
it refreshes, though it does not cure him: and when oil is poured upon a
man's head, it does not always heal a wound, but sometimes makes him a
cheerful countenance, sometimes it consigns him to be a king, or a
priest. So it is in baptism: it does not heal the wounds of actual sins,
because they have not committed them; but it takes off the evil of
original sin: whatsoever is imputed to us by Adam's prevarication, is
washed off by the death of the second Adam, into which we are baptized.

\textbf{HEYLIN,
PRESBYTER
AND CONFESSOR.—On
the Apostles' Creed},

Art. x. Chap. vi.

In which, (Article the 27th) lest any
should object, as Dr. Harding did against Bishop Jewell, that we make
baptism to be nothing but a sign of regeneration, and that we dare not
say, as the Catholic Church teacheth, according to the Holy Scriptures, ``That in and by baptism, sins are fully and truly remitted, and, put
away,'' we will reply with the said most reverent and learned prelate,
(a man who well understood the Church's meaning), That we confess, and
have ever taught, that in the Sacrament of Baptism, by the death and
blood of CHRIST, is given remission of all manner of sins; and that not in
half, or in part, or by way of imagination and fancy, but full, whole,
and perfect of all together; and that if any man affirm, that ``Baptism
giveth not full remission of sins,'' it is no part nor portion of our
doctrine. To the same effect also saith judicious Hooker, ``Baptism is
a Sacrament,'' \&c. [quoted above] … But because these were
private men, neither of which, for aught appears, had any hand in the
first setting out of the Book of Articles, (which was in the reign of
King Edward the Sixth,) though Bishop Jewell had in the second edition,
when they were reviewed and published in Queen Elizabeth's time; let
us consult the Book of Homilies, made and set out by those who composed
the Articles; and there we find, that by GOD'S mercy and the virtue of that sacrifice which our High
Priest and SAVIOUR
CHRIST JESUS,
the SON
of GOD,
once offered for us upon the cross, we do obtain GOD'S grace, and remission, as well of our original sin in
baptism, as of all actual sin committed by us after baptism, if we truly
repent and turn unfeignedly unto Him again. Which doctrine of the Church
of England, as it is consonant to the Word of GOD in Holy Scripture, so is it also most agreeable to the common
and received judgment of pure antiquity. For in the Scripture it is said
expressly by St. Peter, \&c. \&c. This also was the judgment of
the ancient writers, and that too long before the starting of the
Pelagian heresies, to which much is ascribed by some as to the advancing
of the efficacy and fruit of baptism, by succeeding Fathers. For thus
Tertullian; ``Now (saith he) do the waters daily preserve the people of
GOD, death being destroyed and overthrown by
the washing away of sins; for where the guilt is taken away, there is
the punishment remitted also.'' St. Cyprian thus; ``That the remission
of sins, whether given in baptism, or by any other of the sacraments, is
properly to be ascribed to the HOLY
GHOST.''
The African Fathers in full Council do affirm the same, and so doth
Origen also for the Alexandrian, of both which we shall speak anon in
the point of Pædobaptism. Thus Nyssen for the Eastern churches: ``Baptism (saith he) is the expiation of our sins, the remission of our
offences, the cause of our new birth and regeneration.'' Thus do the
Fathers in the Constantinopolitan Council profess their faith in one
baptism (or being only once baptized) ``for the remission of sins.''
And finally, that this was the doctrine of the Church in general, before
Augustine's time, who is conceived to be the first that did advance
the power and efficacy of baptism to so great a height, in opposition to
the Pelagian heresies, appears by a byword grown before his time into
frequent use; the people being used to say, when they observed a man to
be too much addicted to his lusts and pleasures, Let him alone to take
his pleasure, ``for as yet the man is not baptized.'' More of this we
shall see anon in that which follows. Nor is this only \emph{Primitive},
but good \emph{Protestant} doctrine, as is most clear and evident by
that of Zanchius, whom only I shall instance in, of the later writers. ``When the minister baptizeth, I believe that C\textsc{hrist}
with His own hand reacheth as it were from heaven, besprinkleth the
infant with His blood to the remission of sins, by the hand of that man
whom I see besprinkling him with the waters of baptism.'' So that I
cannot choose but marvel how it comes to pass, that it must now be
reckoned for a point of Popery, that the ``Sacraments are instrumental
causes of our justification,'' or of the ``remission of our sins,''
or that it is a point of learning, of which neither the Scriptures, nor
the reformed religion, have taught us anything. So easy a thing it is to
blast that with Popery, which any way doth contradict our own private
fancies.

\textbf{ALLESTRIE,
PRESBYTER.—Serm.
ii. p. 23}.

In our Israel by our covenant there is as much of
this required, for we were all initiated into our profession by washing,
``regenerated in a laver,'' and ``born again of water,'' becoming so
\emph{Tertullian's sanctitatis designati}, set aside for holiness,
consecrated to cleanness, and made the votaries of purity: how clean a
thing then must a Christian be who must be washed into the name? nor is
he thus washed only in the font, there was a more inestimable ``fountain opened for sin and uncleanness.'' (Apoc. xi. 5.)
``JESUS
CHRIST
hath washed us in his own blood;'' and Heb. ix. 14. ``The blood of CHRIST did purge our consciences from dead
works to serve the living GOD.'' How great is our necessity of being clean, when to provide
a means to make us so, GOD
opens his SON'S
side, and our laver is drawn out of the heart of CHRIST. Yet we have more effusions to
contribute to it. (1 Cor. vi. 11.) ``But ye are washed,'' \&c. and
we must ``be baptized with the HOLY
GHOST
and with fire.'' A laver of flame also to wash away our scurfe as well
as sallages, and beyond all these, some of us have been purged too with
the fiery trial, and molten in the furnace of affliction, to separate
our dross and purify us from alloy, that we may be clean and refined
too, may become Christians of the highest \emph{carrect}.

\textbf{BARROW,
PRESBYTER
AND DOCTOR.—Of
the Holy Ghost}.

Serm. xlv. vol. iii. p. 370.

The memorial therefore of that most gracious and
glorious dispensation, [of the HOLY
GHOST at
Pentecost, \&c.] the Christian Church wisely and piously hath
continually preserved, obliging us at this time peculiarly to bless GOD
for that incomparable and inestimable gift conferred then most visibly
upon the Church, and still really bestowed upon every particular member
duly incorporated thereinto.

I say bestowed upon every particular member of the
Church, for the evangelical covenant doth extend to every Christian; and
a principal ingredient thereof is the collation of this SPIRIT, which is the finger of GOD,
whereby (according to the Prophet Jeremy's description of that
covenant) ``GOD'S
law is put into their inward parts, and written in their hearts!''
inscribed (as St. Paul allusively speaketh) not with ink, but by the SPIRIT, \&c. not only as the Jewish law,
represented from without to the senses, but impressed within upon the
mind and affections; whence GOD'S
SPIRIT
is called the SPIRIT
of promise, the donation thereof being the peculiar promise of the
Gospel; and the end of our SAVIOUR'S
undertaking is by St. Paul declared, ``that we might receive the
promise of the SPIRIT
by faith;'' that is, by embracing Christianity, might partake thereof,
according to GOD'S
promise; and the apostolical ministry or exhibition of the Gospel is
styled ``the ministration of the Spirit,'' and tasting `` of the
heavenly gift, and participation of the HOLY
GHOST,''
is part of a Christian's character; and the inception of Christianity
is described by St. Paul, ``But we are bound to give thanks,'' \&c.
(2 Thess. ii. 13); and our SAVIOUR
instructed Nicodemus, that no man can enter into the kingdom of GOD
(that is, become a Christian, or subject of GOD'S spiritual kingdom,) without being regenerated by water,
and by the Spirit, that is, without baptism, and the spiritual grace
attending it, according as St. Peter doth in the words adjoining to our
text imply, that the reception of the HOLY
SPIRIT
is annexed to holy baptism: ``Repent (saith he) and be baptized every
one,'' \&c. ... ``for the promise (that great promise of the HOLY
GHOST)
is unto you,'' \&c. … that is, the HOLY SPIRIT
is promised to all, how far soever distant in place and time, whoever
shall be invited unto, and shall embrace the Christian profession. St.
John also maketh it to be a distinctive mark of those, in whom CHRIST abideth, and who dwell in CHRIST,
that is, of all true Christians, to have this SPIRIT; ``Hereby [saith he] we know that he abideth in us by the
SPIRIT,''
\&c. … and St. Paul denieth him to be a good Christian who is
destitute thereof. ``Now (saith he) if any man have not the SPIRIT,''\&c. … ``and know ye not, (saith he to the
Corinthians) that ye are the temple,'' \&c. … that is, Do ye not
understand this to be a common privilege of all Christians, such as ye
profess yourselves to be? And the conversion of men to Christianity he
thus expresseth, ``After the kindness and love of GOD
our SAVIOUR,''
\&c. (Tit. iii. 4.) ... And all pious dispositions qualifying us for
entrance into heaven and happiness (faith, charity, devotion, every
grace, every virtue) are represented to be the fruits of the HOLY SPIRIT. And the union of all Christians into
one body; the Catholic society of all truly faithful people, doth
according to St. Paul, result from this one SPIRIT, as a common soul animating and actuating them: ``For (saith
he) by one SPIRIT
are they all baptized,'' \&c. …

In fine, whatever some few persons, or some petty
sects (as the Pelagians of old, the Socinians now) may have deemed, it
hath been the doctrine constantly, and with very general consent
delivered in the Catholic Church, that to all persons by the holy
mystery of baptism duly initiated to Christianity, or admitted into the
communion of CHRIST'S body, the grace of GOD'S
HOLY SPIRIT,
certainly is bestowed, enabling them to perform the conditions of piety
and virtue then undertaken by them; enlightening their minds, rectifying
their wills, purifying their affections, directing and assisting them in
their practice; the which holy gift (if not abused, ill treated, driven
away, or quenched by their ill behaviour) will perpetually be continued,
improved, and increased to them; it is therefore by Tertullian (in his
prescriptions against heretics,) reckoned as part of that fundamental
rule which was grounded upon the general tradition and consent of
the Christian Church, that ``CHRIST hath sent the virtue of the HOLY
GHOST,
in his room, which doth act believers;'' to which that article doth
answer of the Apostolical Creed, in which we profess to believe the HOLY GHOST, meaning, I suppose, thereby, not only
the bare existence of the HOLY GHOST,
but also its gracious communication and energy.

\textbf{THORNDIKE,
PRESBYTER.—Book
iii. chap. viii}.

It is demanded in the second place, what is that
regeneration by the HOLY
GHOST,
and wherein it consists, whereof infants that are baptized can be
thought capable. For the wild conceits of those that imagine them to
have faith in CHRIST
(which without actual motion of the mind, is not), require miracles to
be wrought of course, by baptizing, that the effect thereof may come to
pass. And if the state of grace (which the habitual grace of GOD'S SPIRIT
either supposed or inferreth) is not to be attained but by the
resolution of embracing the covenant of grace, (as, by all the premises,
it is not otherwise attended), it will be every whit as hard to say what
is that habitual grace, that is said to be poured into the souls of
infants that are baptized, being nothing else but a facility in doing
what the Covenant of Grace requireth. But, if we conceive the
regeneration of Infants that are baptized to consist in the habitual
assistance of GOD'S
SPIRIT,
the effects whereof are to appear, in making them able to perform that
which their Christianity requires at their hands, so soon as they shall
understand themselves to be obliged by it; we give reason enough of the
effect of their baptism, whether they die or live, and yet become not
liable to any inconvenience. For supposing the assistance of GOD'S SPIRIT
assigned them by the promise of baptism, to take effect when their
bodily instruments enable the soul to act as Christianity requireth; if
the soul, by death, come to be discharged of them, can any thing be said
why original concupiscence, which is the law of the members, should
remain any more, to impeach the subjection of all faculties to the law
of GOD'S SPIRIT? Or will it be any thing strange, that
when they come to be taught Christianity, the same SPIRIT of GOD should be taught to sway them, to embrace
it of their own choice, and not only in compliance with the will of
their parents? Yet is this no more, than the regeneration of infants by
water and the HOLY
GHOST
importeth; that the SPIRIT
of GOD
should be habitually present, to make those reasons which GOD hath given to convince the world, that they
ought to be Christians, both discernible to the understanding, and
weighing down the choice; whereas, those that are converted from being
enemies to GOD,
(that is to say, at those years, when no man can be converted to GOD,
that is not His enemy before), though the SPIRIT of GOD
knock at their hearts without, striving to cast out the strong man that
is within doors, and to make a dwelling for itself in the heart, are
possessed by a contrary principle, till they yield GOD'S SPIRIT that entertainment which GOD
requireth.

\textbf{PEARSON,
BISHOP AND
DOCTOR.—Exposition
of the Creed}.

Article ix.

Being therefore we are that the preaching remission
of sins belongeth not only certainly, but in some sense peculiarly, to
the Church of CHRIST,
it will be next considerable how this remission is conferred upon any
person in the Church.

It is certain that forgiveness of sins was promised
to all who were baptized in the name of CHRIST; and it cannot be doubted but all persons who did perform
all things necessary to the receiving the ordinance of baptism, did also
receive the benefit of the ordinance, which is \emph{remission of sins}.
``John did baptize in the wilderness, and preach the baptism of
repentance for the remission of sins.'' (Mark i. 4.) And St. Peter made
this exhortation of his first sermon, ``Repent, and be baptized, every
one of you, in the name of JESUS
CHRIST,
for the remission of sins.'' (Acts ii. 38.) In vain doth doubting and
fluctuating Socinus endeavour to evacuate the evidence of this
Scripture; attributing the remission either to repentance without
consideration of baptism, or else to the public profession of faith made
in baptism; or if any thing must be attributed to baptism itself, it
must be nothing but a declaration of such remission. For how will these
shifts agree with that which Ananias said unto Paul, without any mention
either of repentance or confession, ``Arise and be baptized, and
wash away thy sins:'' (Acts xxii. 16.) and that which St. Paul, who was
so baptized, hath taught us concerning the Church, that CHRIST doth ``sanctify and cleanse it
with the washing of water.'' (Eph. v. 26.) It is therefore sufficiently
certain that baptism, as it was instituted by CHRIST after the pre-administration of St.
John, wheresoever it was received with all qualifications necessary in
the person accepting, and conferred with all things necessary to be
performed by the person administering, was most infallibly efficacious,
as to this particular, that is, to the remission of all sins committed
before the administration of this sacrament.

\textbf{BULL,
BISHOP AND
DOCTOR.—Sermon
vii}.

``And besides this,'' \&c. (2 Pet. i. 5.) As
if he had said, You have now, GOD
be thanked, escaped the pollutions of the world, and are truly, I hope,
converted to Christianity, and in baptism have been regenerated by the HOLY
GHOST
(that he means by their being \emph{made partakers of the divine nature}).
This indeed is a very great achievement, and an invaluable mercy of GOD, vouchsafed to you; yet I beseech
you, rest not here: but \emph{besides this, giving all diligence, add to
your faith, virtue}, \&c.

\textbf{COMBER,
PRESBYTER.—Part
iii. sect. iii}.

p. 201.

We must not presently turn our backs upon GOD
so soon as the holy rite [baptism] is finished, but complete the
solemnity by thanksgiving and prayer; and that we may do both, not only
with the spirit, but with understanding, the minister doth here teach us
what must be the subjects of our praises and petitions.

I. Our praises must look back upon the grace
already showed, and the benefits which are already given to this infant,
which are principally two 1. Internally it is regenerated; 2. Externally
it is grafted into CHRIST'S
Church, for which we must give hearty thanks to ALMIGHTY GOD.
To which we must add, II. Our prayers, which must look forward upon the
grace which will be needful to enable it to live answerable to this
estate into which it is admitted; and this we must beg of ALMIGHTY GOD also, or else the former blessings will be
altogether in vain. Now all this is plain, that no more would need to be
added, but only that some with Nicodemus are apt to say, ``How can
these things be?'' (John iii. 9.) Judging it impossible that so great a
matter as regeneration can be effected so soon, and by so mean an
instrument as they account it; whereas the effect is to be ascribed to
the divine power of the Author, not to the intrinsic efficacy of the
outward means: yet in regard we can never bless GOD
heartily for a mercy unless we believe He hath bestowed it, we must
labour to remove these scruples by a fuller account of this baptismal
regeneration, that we may not withhold the divine praises, by our
doubting and unbelief. The word regeneration is but twice (that I know
of) used in Scripture; first, Matt. xix. 28, ``Ye that have followed me
in the regeneration;'' where though (by altering the
point—``followed me, in the regeneration when the Son of Man,''
\&c.) it may signify the resurrection yet as we read, it signifies
the renewing of men by the Gospel and baptism. Secondly, (Tit. iii. 5.) ``He saved us by the laver of regeneration, and renewing of the HOLY GHOST,''
which is a paraphrase upon that of our SAVIOUR, (John iii.) ``Except a man be born of
water,'' \&c. And, because persons, come to age before their
conversion, are first taught and persuaded by the Word of GOD, the language of Holy Writ enlarges the
metaphor, and saith, Such as are begotten by the Word of GOD,
(1 Cor. iv. 15.) and then born again or regenerated in baptism. In like
manner speak the Fathers, who do constantly and unanimously affirm, that
we are regenerated in, or by baptism. So that we must next inquire
wherein this regeneration doth consist. And first, whereas both children
and those of riper years are by nature dead in sin, so that they lie
under the guilt and power thereof: our gracious FATHER
doth here in baptism seal a covenant with us, wherein He promiseth to
pardon us; and when this deadly load is removed, the soul receives as it
were a new life, and takes new hopes and courage, being restored to the
divine favour, and being set free from the sad expectations of
unavoidable condemnation for former sin, original in infants, and both
it and actual in those of riper ears. Before this covenant we were dead
in law, and by the pardon of our sins we are begotten again to a lively
hope; and herein stands the first particular of our regeneration,
viz. in the remission of sins, wherefore both Scripture and antiquity
teach us, that baptism is the means for remission of sin, and hence they
join pardon and regeneration commonly together, because this forgiveness
puts us into a new estate, and an excellent condition in comparison of
that which our natural birth had left us in.

Secondly; But further, by baptism we gain new
relations, and old things being done away all things become new, \&c.
… Thirdly; Our corrupt nature is changed in baptism, and there is a
renovation effected thereby, both as to the mortification of the old
affections, and the quickening of the new, by the HOLY
SPIRIT,
which is hereby given to all that put no bar or impediment unto it. This
was the Ancients' doctrine, who affirmed a real change to be wrought,
and believed the SPIRIT
to be therein bestowed as GOD
had promised, (Ezek. xxxvi. 25, 26.) ``That he would sprinkle clean
water,'' \&c. …And it is manifest that in the first ages of the
Church there were abundant of gifts and graces miraculously bestowed
upon Christians in their baptism, and no doubt if the catechumens of our
days who are at age, would prepare themselves as strictly by repentance,
fasting, and prayer, as they of old did, they should find incomparable
effects of this sacred laver, if not in as miraculous measures, yet to
as real purposes; that is, they should be truly regenerated, and their
hearts changed by the influence of the Divine SPIRIT.
But some may doubt whether infants be regenerate in this sense, because
they are not capable of giving any evidences of their receiving the SPIRIT, nor doth there any immediate effects of
their regeneration appear; hence the Palagians denied it, but they are
therefore condemned by the Milevitan Council, Can. ii. and confuted by
St. Aug. ad Bon. lib. iii. It is confessed they can show no visible
signs of spiritual life in the operations thereof, no more can they of
their having a rational soul, for some time, and yet we know they have
the power of reason within them; and since all infants are alike, either
all do here receive a principle of new life, or none receive it;
wherefore I see no reason why we may not believe as the ancients did,
that GOD'S
grace (which is dispensed according to the capacity of the suscipient) is here given to infants to heal their nature, and that He
bestowed on them such measures of His SPIRIT as they can receive; for the malignant
effects of the first Adam's sin are not larger than the free gift
obtained by the second Adam's righteousness. (Rom. v. 15, 18.) And if
it be asked how it comes to pass, then, that so many children do
afterwards fall off to all impurity? I answer, so do too many grown
persons also, and neither infants nor men are so regenerated in this
life, as absolutely to extinguish the concupiscence: for the flesh still
will lust against the Spirit: but thus GOD
gives the SPIRIT
also to lust against the flesh. (Gal. v.)

\textbf{KEN,
BISHOP AND
CONFESSOR.—Exposition
of Church Catechism},

p. 136.

Glory be to thee, O most indulgent Love, who in our
baptism dost give us the HOLY
SPIRIT
of Love, to be the principle of new life and of love in us, to infuse
into our souls a supernatural, habitual grace, and ability to obey and
love thee, for which all love, all glory be to thee.

Glory be to thee, O compassionate Love, who, when
we were conceived and ``born in sin,'' of sinful parents, when we
sprang from a root wholly corrupt, and were all ``children of wrath,''
hast in our baptism ``made us children'' of thy own heavenly FATHER by adoption and grace; when we were heirs of hell, hast
made us heirs of heaven, even joint heirs with thy own self, of thy own
glory; for which, with all the powers of my soul, I adore and love thee.

\textbf{PATRICK,
BISHOP.—On
Baptism},

p. 441.

The sum of all is, that hereby we are regenerated
and born again. It is the sacrament of the new birth, by which we are
put into a new state, and change all our relations: so that whereas
before we were only the children of Adam, we are now taken to be the
children of GOD;
such of whom He will have a fatherly care, and be indulgent and merciful
unto. We have now a relation likewise to CHRIST
as our Head, and to the HOLY
GHOST
as the giver of life and grace. Yea, herein He grants remission of sin,
and we are sanctified, and set apart to His uses. We being hereby given
to Him, and He accepting of us, do become His possession and proper
goods, and cannot without being guilty of the foulest robbery, sin
against GOD. We are made hereby the temples of the HOLY
GHOST,
the place where He, and nothing else is to inhabit; and being by this
consecrated to Him, He likewise then enters upon His possession, and we
are said thereby to receive the HOLY
GHOST;
so that if we run into sin, we defile His house, and commit the greatest
profaneness and impiety, and may be said very truly to do despite to the
SPIRIT
of GOD
whereby we were sanctified.

\textbf{BEVERIDGE,
BISHOP AND
DOCTOR.—On
admission into the Church by Baptism}.

Vol. i. Serm. xxxv. p. 304.

But what he means by being ``born of water and the
SPIRIT,''
is now made a question: I say now, for it was never made so till of late
years. For many ages together none doubted of it, but the whole
Christian world took it for granted, that our SAVIOUR,
by these words, meant only that except a man be baptized according to
His institution, he cannot enter into the kingdom of GOD; this being the most plain and obvious sense of the words,
forasmuch as there is no other way of being born again of water, as well
as of the SPIRIT
but only in the Sacrament of Baptism.

To understand what He means by being born again, we
must call to mind what He saith in another place, ``My kingdom is not
of this world;'' (John xviii. 36.) though it is in this world, it is
not of it; it is not a secular or earthly kingdom, but a kingdom purely
spiritual and heavenly: ``It is not meat and drink, but righteousness,
and peace, and joy in the HOLY
GHOST
(Rom. xiv. 17.) And therefore when a man is born into this world, he is
not thereby qualified for the kingdom of GOD,
nor hath any right or title to it, no more than as if he had not been
born at all; but before he enter into that, he must be born again, he
must undergo another kind of birth than he had before: he was before
born of the flesh, he must now be born of the SPIRIT;
otherwise he cannot be capable of entering into such a kingdom as is
altogether spiritual. Thus our LORD
Himself explains his own meaning in my text, by adding immediately in
the next words, ``That which is born of the flesh, is flesh,'' \&c.
… As if He had said, He that is born, as all men are at first, only of
the flesh, such a one is altogether carnal and sensual; and so can be
affected with nothing but the sensible objects of this world. But he
that is born of the SPIRIT
of GOD,
thereby becomes a spiritual creature, and so is capable of those
spiritual things of which the kingdom of GOD
consisteth, ``even of righteousness, and peace, and joy in the HOLY
GHOST.''
And he whose mind is changed, and turned from darkness to light, and
from the power of Satan unto GOD,
is truly said to be born again; because he is quickened with another
kind of life than he had before; and to be born of the SPIRIT of GOD, because it is by it that this new and
spiritual life is wrought in him. So that he is now born into another
world, even into the kingdom of GOD,
where he hath GOD
himself, of whom he is born, for his Father, and the kingdom of GOD
for his portion and inheritance. And therefore it is, that except a man
be thus born of the SPIRIT,
it is impossible he should enter into the kingdom of GOD, seeing he can enter into it no other way,
than by being born of the SPIRIT.

But that we may thus be born of the SPIRIT, we must be born also of water,
which our SAVIOUR
here puts in the first place. Not as if there was any such virtue in
water, whereby it could regenerate us, but because this is the rite or
ordinance appointed by CHRIST,
wherein to regenerate us by His HOLY SPIRIT;
our regeneration is wholly the act of the SPIRIT of CHRIST. But there must be something done on our
parts in order to it, and something that is instituted and ordained by CHRIST
Himself, which in the Old Testament was circumcision; in the New,
baptism, or washing with water; the easiest that could be invented, and
the most proper to signify His cleansing and regenerating us by His HOLY SPIRIT.
And seeing this is instituted by CHRIST Himself, as we cannot be born of water without the SPIRIT,
neither can we, in an ordinary way, be born of the SPIRIT without water, used or applied in
obedience and conformity to His institution. CHRIST hath joined them together, and it is not
in our power to part them: he that would be born of the SPIRIT,
most be born of water too …

As baptizing necessarily implies the use of water,
so our being made thereby disciples of CHRIST, as necessarily implies our partaking of His SPIRIT:
for all that are baptized, and so made the disciples of CHRIST, are thereby made the members of His
body; and are therefore said to be baptized into CHRIST, (Rom. vi. 5. Gal. iii. 27.) But they
who are in CHRIST,
members of His body, must needs partake of the SPIRIT that is in Him their Head. Neither doth
the SPIRIT
of CHRIST
only follow upon, but certainly accompanies the Sacrament of Baptism,
when duly administered according to His institution. For as St. Paul
saith, ``By one SPIRIT
we are all baptized into one body.'' (1 Cor. xii. 13.) So that in the
very act of baptism, the SPIRIT
unites us unto CHRIST,
and makes us members of His body; and if of His body, then of His Church
and kingdom, that being all His body. And therefore all who are rightly
baptized with water, being at the same time baptized also with the HOLY
GHOST,
and so born of water and the SPIRIT,
they are, \emph{ipso facto}, admitted into the kingdom of GOD,
established upon earth, and if it be not their own fault, will as
certainly attain to that which is in heaven.

\emph{Ibid}. p. 306 —This I would desire all here
present to take special notice of, that you may not be deceived by a
sort of people risen up among us, who being led, as they pretend, by the
light within them, are fallen into such horrid darkness, and damnable
heresies, that they have quite laid aside the Sacrament of Baptism, and
affirm, in flat contradiction to our SAVIOUR'S
words, that they may be saved without it. I pray GOD to open their eyes, that they may not go blindfold into
eternal damnation. And I advise you all, as you desire not to apostatize
from the Christian religion, and as you tender your eternal salvation,
take heed that you be never seduced by them, under any pretence
whatsoever; but rather, if you be acquainted with any of them, do what
you can to turn them from darkness to light, from the power of Satan
unto GOD again; that they may obtain
forgiveness of their sins, and inheritance among them who are sanctified
by faith in Him, who saith, ``Except a man be born of water,'' \&c.

Not only a man, in contradiction to a child, or a
woman, but as it is in the original, [\emph{ean m\={e} tis}], except
any one, any human creature whatsoever, man, woman, or child, ``except
he be born of water,'' \&c. ... So that our LORD
is so far from excluding children from baptism, that He plainly includes
them, speaking in such general terms, on purpose that we may know that
no sort of people, old or young, can ever be saved without it. And so He
doth too, where He commands, as was observed before, that ``All nations
should be made disciples by being baptized in the name of,'' \&c.
… For, under all nations, children must needs be comprehended, which
make a great, if not the greatest part of all nations. And although
these general expressions be sufficient to demonstrate the necessity of
Infant Baptism, yet foreseeing that ignorant and unlearned people would
be apt to wrest the Scriptures to their own destruction, He elsewhere
commands children particularly to be brought unto Him, saying, ``Suffer
the little children,'' \&c. (Mark x. 14.) But if the kingdom of GOD
consist of children, as well as other people, they must of necessity be
baptized, or born of water and the SPIRIT;
for otherwise, He Himself saith, ``They cannot enter into the
kingdom.''

Hence it is, that we find the Apostles baptizing
whole families, children, if any, as well as others: and the whole
Catholic Church, in all places and ages ever since, hath constantly
admitted the children of the believing parents into the Church, by
baptizing them according to the institution and command of our SAVIOUR; none ever making any question of it, but all Christians,
all the world over, taking it for granted that it ought to be done, till
of late years.

\textbf{SHARP,
ARCHBISHOP.—Vol.
v. Sermon v}.

p. 71.

There is the same relation between CHRIST and Christians, that there is
between the vine and the branches; the same necessity of communication
of vital influences from the root to the branch in the one as in
the other: which communication of influences is made by the HOLY SPIRIT of GOD, derived from CHRIST, and diffusing Himself into every
particular member of the whole body of Christians. Hence it is
Christians are so frequently called the Temples of the HOLY GHOST. ``Know ye not,'' saith St. Paul, ``that ye are,'' \&c.; and, again,
``Know ye not, that your bodies
are the members,'' \&c. which he explains presently after thus: ``Know ye not that your bodies are the Temples,'' \&c. And the same
St. Paul, in the eighth to the Romans, lays the foundation of our
relation to CHRIST, and our hopes of eternal life, in the
very thing, viz. the SPIRIT
of GOD
his dwelling in us; as may be there seen more at large.

This, then, being the privilege of all Christians,
that by their being consecrated to CHRIST, they have a right to the continual presence of the HOLY
GHOST in
their souls; or, if you will, GOD
hath so great a right and property in them, that He sends down His HOLY
SPIRIT
to take possession of them, in order to the securing and sealing them
for His own in the other world; we may easily, from hence, gather what
it is \emph{to grieve the Holy Spirit}, (which is the thing we are now
inquiring into,) viz. We then grieve Him, when being already Christians
in profession, we either will not vouchsafe Him a lodging in our hearts,
which He doth desire; and, in order to the obtaining it, makes frequent
applications to our souls by His holy motions; or, when we have already
given Him entertainment, we carry ourselves so unbecomingly towards Him,
as to tempt Him to forsake us. We then grieve the HOLY
SPIRIT,
when, having taken upon ourselves the covenant of Baptism, and thereby
consecrated and consigned ourselves to GOD,
we either refuse to admit the SPIRIT to take possession of us, or having admitted Him, do not
show that respect, nor observe that decency, nor express that kindness,
that is due to so worthy a guest: but by our rude, and unmannerly, and
ill-natured behaviour towards Him, put such affronts upon Him as highly
provoke Him to quit his habitation.

\textbf{SCOTT,
PRESBYTER.—Christian
Life, chap. ii. sect i}.

p. 354.

Second sort of the HOLY GHOST'S operations, viz. that which He
ordinarily doth, and always hath done, and will always continue to do;
for upon the cessation of these His miraculous operations, the HOLY
GHOST
did not wholly withdraw Himself from mankind, but He still continues
mediating with us under CHRIST,
in order to the reconciling our wills and affections to GOD, and subduing that inveterate malice and
enmity against Him, which our degenerate nature hath contracted. For it
is by this blessed SPIRIT
that CHRIST
hath promised to be with us to the end of the world. (Matt. xxviii. 20.)
And CHRIST
Himself hath assured us, that upon His ascension into heaven He would ``pray His Father, and He should give us another Comforter,'' meaning
this HOLY GHOST, ``that he might abide with us for
ever;'' (John xiv. 16.) and, accordingly, the HOLY GHOST is vitally united to the Church of CHRIST,
even as souls are united to their bodies. For as there is one Body, the
Church, so here is one SPIRIT,
\emph{i.e}. the HOLY
GHOST,
which animates that Body, (Eph. iv. 4.) and hence the unity of the
Church is in the foregoing verse called the unity of the SPIRIT: because as the soul, by diffusing itself through all the
parts of the body, unites them together, and keeps them from flying
abroad, and dispersing into atoms; so the HOLY SPIRIT,
by diffusing Himself throughout this mystical Body, joins and unites all
its parts together, and makes it one separate and individual
corporation. So that, when by Baptism we are once incorporated into this
body, we are entitled to, and do at least, \emph{de jure}, participate
of the vital influences of the HOLY
GHOST,
who is the soul of it; and accordingly, as Baptism joins us to that
body, of which this Divine SPIRIT
is the soul; so it also conveys that Divine SPIRIT to us. So that, as in natural bodies, those ligaments which
unite and tie the parts to one another, do also convey life and spirit
to them all; so also in this mystical body, those federal rites of
Baptism and the LORD'S
Supper, which are, as it were, its nerves and arteries, that join and
confederate its members to one another, are also the conveyance of that
spiritual life from the HOLY
GHOST,
which moves and actuates them all. And hence the ``washing of
regeneration,'' and ``the renewing of the HOLY
GHOST,''
the ``being born of water and of the HOLY GHOST,''
are put together as concurrent things; and in Acts ii. 38. Baptism
is affirmed to be necessary to our receiving the HOLY GHOST; and if by Baptism we receive the HOLY
GHOST,
that is a right and title to His grace and influence, then must the HOLY GHOST be still supposed vitally united to the
Church, whereof we are made members by our Baptism, and, like an
omnipresent soul, to be diffused all through it, and to move and actuate
every part of it by His heavenly grace and influence.

\textbf{JENKIN,
PRESBYTER.—On
Christian Religion, vol. ii}.

p. 427.

Baptism is very agreeable to the nature of the
Christian Religion, being a plain and easy rite, and having a natural
significancy of that purity of heart, which it is the design of the
Gospel to promote and establish in the world; and it is fitted to
represent to us the cleansing of our souls by the Blood of CHRIST, and the grace of purity and holiness, which is conveyed in
this sacrament, and the spirit of regeneration which is conferred by it.
Tit. iii. 5.

\textbf{SHERLOCK,
BISHOP.—Vol.
ii. Disc. vii}.

You see the power of Baptism, and the blessings
that are annexed to it, to which all are entitled who partake in the
Baptism of CHRIST:
for Himself He was neither born nor baptized but for our sakes; that the
blessings of both might descend on us, who, through faith, are heirs
together with Him of the promises of GOD.

By Baptism the gates of heaven are set open to us,
and the way paved for our return to our native country. By Baptism we
are declared to be such sons of GOD
in whom He will delight, and whom He will appoint to be heirs of His
kingdom. By Baptism we receive the promise of the SPIRIT, by which we cry, Abba, Father.

Are not these great privileges? And is not here
room for mighty expectations? And yet how unsuitable to these claims do
the circumstances of a Christian's life often appear? He is upon the
road to heaven, you say, and the gates stand open to receive him; but
how does he stumble and fall like other men, and sometimes lose his way,
and wander long, bewildered in night and darkness? Or, if he keeps the
road, how lazily does he travel, as if he were unwilling to come to
his journey's end, and afraid to see the country which he is going to
possess? The Christian only, of all men, pretends to supernatural power
and strength, and an intimate acquaintance with the SPIRIT
of GOD;
and yet how hardly does he escape the pollutions of the world, and how
often look back, with languishing eyes, upon the pleasures, riches, and
honours of this life? And though he boasts of more than human strength,
yet how does he sometimes sink below the character and dignity even of a
man? Ye sons of GOD,
for such ye are, how do ye die like the children of men, and how like is
your end to theirs?

And what must we say of these things? Is the
promise of GOD
become of none effect? Is Baptism sunk into mere outward ceremony, and
can no longer reach to the purifying the heart and mind? The fact must
not be disputed; it is too evident, at least in these our days, that the
lives of Christians do not answer to the manifold gifts and graces
bestowed on them.

\textbf{WALL,
PRESBYTER.—On
infant Baptism, part ii. chap. vi}.

I believe Calvin was the first that ever denied
this place (John iii. 5.) to mean Baptism. He gives another
interpretation, which he confesses to be new. This man did, indeed,
write many things in defence of Infant Baptism. But he has done ten
times more prejudice to that cause, by withdrawing (as far as in him
lay) the strength of this text of Scripture (which the ancient
Christians used as a chief ground of it) by that forced interpretation
of his, than he has done good to it by all his new hypotheses and
arguments. What place of Scripture is more fit to produce, for the
satisfaction of some plain and ordinary man, (who, perhaps, is not
capable of apprehending the force of the consequences by which it is
proved from other places,) that he ought to have his child baptized,
than this, (especially if it were translated in \emph{English}, as it
should be,) where our SAVIOUR says, that no \emph{person} shall come to heaven without
it? meaning, at least in GOD'S \emph{ordinary} way.

\textbf{POTTER,
ARCHBISHOP.—of
Church Government, chap. i}.

p. 14.

Whoever wilfully neglects to be made a member of
the Christian Church, does, by necessary consequence, deprive himself of
all the privileges which belong to it; just as in any civil corporation
they who are not members of it can plead no right to any of its
privileges. This has already been shown to be the sense of CHRIST, and the same is constantly affirmed by the Christian
writers of all ages. ``They who do not come into the Church [saith Irenæus]
do not partake of the SPIRIT,
but deprive themselves of life.'' For where the Church is, there is the
SPIRIT
of GOD.
And in St. Cyprian's opinion, he cannot have GOD for his Father who has not the Church for his Mother.

Hence the privileges of the Christian Church, such
as Remission of Sins, the Grace of the HOLY SPIRIT,
and Eternal Life, are commonly said to be annexed to Baptism, this being
the constant rite of initiation into the Church. Thus, in Ananias's
exhortation to St. Paul, ``Arise, and be baptized,'' \&c. … St.
Barnabas expressly affirms, that ``Baptism procures remission of
sins;'' and proves, from the Scriptures, that they who are baptized,
are received into GOD'S favour, whereas all the rest of mankind
lie under His displeasure. Peter thus exhorts his new converts ``Repent, and be baptized every one of you,'' \&c. … Our blessed
SAVIOUR joins Faith and Baptism together, as
necessary conditions of salvation: ``Except a man be born of water and
of the SPIRIT,''
\&c. … And in another place, ``He that believes, and is baptized,
shall be saved.'' From these, and like passages of Scripture, the
Primitive Church constantly inferred, that where the Gospel had been
sufficiently propounded, no man could be saved, without Baptism actually
obtained, or earnestly desired. Whence Tertullian calls it the ``happy
sacrament of water, whereby we are washed from the sins of our former
blindness, and delivered into eternal life.'' And Cyprian gives this
reason, why the Baptism of infants should not be delayed so long as the
eighth day after their birth, that (since it is said in the Gospel,
that ``the Son of Man came not to destroy men's souls, but to save
them'') it is our duty, as far as in us lies, to take care that no soul
shall be destroyed.

\textbf{NELSON,
CONFESSOR.—Festivals
and Fasts},

p. 115.

By this means [by Baptism] the children of
believers are entered into covenant with GOD under the Gospel, as they were under the Law by circumcision;
and that infants are capable of this federal relation, is plainly
declared by Moses; (Deut. xxix. 11.) and since they are the offspring of
Adam, and consequently obnoxious to death by his fall, how can they be
made partakers of that redemption which CHRIST
hath purchased for the children of GOD, if they do not enjoy the advantage of that method which is
alone appointed by CHRIST
for them to become members of GOD'S kingdom? For JESUS himself has assured us, ``Except one be born of water,''
\&c. … And therefore it was the constant custom of the Primitive
Church to administer Baptism to infants for the remission of sins. And
this practice was esteemed, by the best tradition, to be derived from
the Apostles themselves.

\textbf{WATERLAND,
PRESBYTER.—On
Regeneration. 2}.

The second is the case of infants. Their innocence
and incapacity are to them instead of repentance, which they do not
need, and of actual faith, which they cannot have. They are capable of
being savingly born of water and the SPIRIT,
and of being adopted into sonship with what depends thereupon; because,
though they bring no virtues with them, no positive righteousness, yet
they bring no obstacle, no impediment. They stipulate, they enter into
contract, by their sureties, upon a presumptive and interpretative
consent: they become consecrated in solemn form to FATHER, SON,
and HOLY
GHOST;
pardon, mercy, and other covenant privileges are made over to them; and
the HOLY
SPIRIT
translates them out of their state of nature (to which a curse belongs)
to a state of grace, favour, and blessing; this is their regeneration.

\textbf{KETTLEWELL,
PRESBYTER
AND CONFESSOR.—On
the Creed— Article, Forgiveness of Sins},

p. 685.

\emph{Ques}.—For whose sake doth ALMIGHTY GOD allow us all this benefit of forgiveness?

\emph{Ans}.—For JESUS CHRIST,
who, as you have seen, died for our sins, and gave His blood a ransom,
to purchase for us this pardon of them. ``He is set forth a
propitiation,'' \&c, (Rom. iii. 25.) And thus we shall receive all
this mercy for His sake, when, with the disposition before expressed, we
devoutly pray to GOD
for it in His name.

\emph{Ques}.—By the promises of the Gospel, I see
that this forgiveness is assured to all Christians upon the terms which
you have described. But is it in any signs and tokens outwardly
dispensed to them?

\emph{Ans}.—Yes; both in the Holy Sacraments and
in the sacerdotal Absolution. Which ways of ministering this
forgiveness, as well as the forgiveness itself, are noted in some
ancient Creeds: this article being thus professed in St. Cyprian's
Form at Baptism: ``I believe the Remission of Sins by the Church.''

\emph{Ques}.—Is this forgiveness dispensed to us
in the Sacrament of Baptism?

\emph{Ans}.—Yes; and that most amply, the water
of Baptism washing off the stain of all former sins. ``Be baptized, and
wash away thy sins,'' said Ananias to Saul; ``Repent, and be baptized
for the Remission of Sins,'' said St. Peter to the Jews; and ``He hath
saved us by the Laver of Regeneration;'' \emph{i.e}. the Water of
Baptism, and the renewing of the HOLY
GHOST.
(Tit. iii. 5.) So that whatever pollutions men had upon them, if they
come to Baptism with true faith and repentance, they are thereby made
clean again.

\textbf{HICKES,
BISHOP AND
CONFESSOR.—Christian
Priesthood}.

It belonged to the Apostles and Presbyters, by
virtue of their sacerdotal office and ministry, to be advocates and
intercessors with GOD
... I need not insist upon their power of baptizing for the remission of
sins with fasting and prayer, which was a most solemn act of
expiation for washing away all the past sins of the baptized.

\textbf{JOHNSON,
PRESBYTER.—Unbloody
Sacrifice.—Vol. i. ch. ii}.

sect. 1. p. 165.

I think the only immediate effect of the SPIRIT
in Baptism, is the remission of all sin, and removing our natural
disability to the worship and service of GOD,
and the sentence of condemnation under which we were all born: (Rom. v.
16.) and that other graces are wrought in us by that HOLY SPIRIT, which by Baptism receives us under its
protection, gradually, and according to the capacity of the recipient:
and this doctrine I learned from those words of St. Barnabas, in his
Epistle, cap. vi.: ``After, therefore, that CHRIST
had renewed us by the remission of our sins, he made us [in] another
shape, so as to have an infant-like soul, even as he himself reformed
us:'' where he plainly makes renovation to consist in forgiving sins;
and makes the new moulding, or reformation of our minds, to be not
performed at the same time with the other, or all at once, but to be
consequent upon the former renovation: and CHRIST
is always thus reforming us, from our Baptism to our death. And I look
on these words of St. Barnabas to be a better explication of the
renovation, or regeneration of Christians by Baptism, than whole volumes
of modern writers upon the same subject. And I may here very reasonably
observe, that as the HOLY
SPIRIT
is present in our Baptism, to seal the remission of sins, and to infuse
the beginnings of Christian Life; for He is present in confirmation, to
shed further influences on them that receive it, for the further
suscitation of the gift of GOD
bestowed in Baptism and in the Eucharist, as will hereafter appear at
large for our further progress and increase in grace.

\textbf{LESLIE,
PRESBYTER
AND CONFESSOR.—On
Water Baptism, § 5}.

The end of CHRIST'S Baptism was, to instate us into all the inconceivable
glories, and high eternal prerogatives which belong to the members of
His body, of His flesh, and of His bones, (Eph. v. 30.) that we might
receive the adoption of sons. (Gal. iv. 5.) Henceforth no more
servants, but sons of GOD,
and heirs of heaven! These are ends so far transcendent above the ends
of all former Baptisms, that, in comparison, other Baptisms are not only
less, but none at all; like the glory of the stars in presence of the
sun, they not only are a lesser light, but when he appears they become
altogether invisible.

And as a pledge or foretaste of these future and
boundless joys, the gift of the HOLY
GHOST is
given upon earth, and is promised as an effect of the Baptism of CHRIST:
as Peter preached (Acts ii. 38), ``Repent, and be baptized,'' \&c.
And (Gal. iii. 27.) ``As many of you as have been baptized into CHRIST, have put on CHRIST.''

This of the gift of the HOLY GHOST was not added to any Baptism before CHRIST'S,
and does remarkably distinguish it from all others.

\textbf{WILSON,
BISHOP,
CONFESSOR,
AND DOCTOR.—Maxims
of Piety, Vol. i}.

p. 310.

The HOLY
SPIRIT
at Baptism takes possession of us, and keeps possession, till men grieve
Him; then He forsakes us, and an evil spirit succeeds.

By Baptism we contract and oblige ourselves, all
our life long, to complete and perfect the image of JESUS CHRIST in ourselves. The blessings and
excellencies of Baptism:—It separates us from Adam, and engrafts us
into CHRIST.—It
is a resurrection from sin to grace.—It discharges us from the debt
owing to the justice of GOD,
by our sins, now fully satisfied by faith in the suffering and death of
CHRIST.—It
cancels the law of death and malediction which was against us.—In
Baptism our sins did indeed die, and were buried; but the seed and root
remain in us. These we are to mortify all our lives long.

\textbf{BINGHAM,
PRESBYTER.—On
Lay Baptism, part ii. ch. vi}.

…
What it [indelible character] was taken to signify in Baptism? For an
indelible character was always supposed to be imprinted as much in
Baptism as in ordination; though I do not remember that any ancient
Council expressly used that term about either, but only say
something that may be reckoned equivalent to it; and that is this, as it
relates to baptism: that a man, who is once truly baptized, can never do
any thing that can so far erase or cancel his Baptism, as that he shall
need, upon any occasion, to be re-baptized with a second Baptism. Thus
far the ancients believed an indelible character in Baptism. Though a
man turns his back on Christianity, and totally apostatize and fall away
from the profession of it; though he turn heretic or schismatic; though
he excommunicate himself, or be excommunicated by the Church; though he
embrace Paganism or Judaism, or any other opposite Religion; though he
curse and blaspheme CHRIST
in a synagogue or in a temple, as many of the old apostates did; though
he become a \emph{Julian} or an \emph{Ecebolius}, and ``trample under
foot the SON
of GOD,
and count the Blood of the Covenant an unholy thing, and do despite to
the Spirit of Grace;'' yet, after all, if this man turn again to
Christianity, he was not to be received by a second Baptism. His
repentance, and the Church's absolution, was sufficient in that case
to re-instate him in his ancient profession, and he was not to be
re-baptized to be made again a Christian. The Church had but one Baptism
for the remission of sins, and the virtue of that was so far indelible,
that it would always qualify the man that had received it, to be
admitted to communion again after the greatest apostasy, only by a true
repentance and reconciliatory imposition of hands, without re-baptizing.
This was what the ancients understood by what we now call the indelible
character of Baptism. But they were far from thinking, that a man who
was such an apostate had any right or authority, whilst he was an
apostate, to challenge any of the common privileges of a Christian. They
did not think, whilst he was a Pagan or a Jew, that he was properly a
member of the Christian Church still, because of his Baptism; or that he
had any right to be called Christian, or to be admitted to the prayers
of the Church, and much less to the communion with other faithful
laymen: and yet, after all, there was so much of a Christian in him, by
virtue of his Baptism, that he needed not to be baptized again as a mere
\emph{Jew} or \emph{Pagan}. His Baptism was such as nothing could
obliterate; it would remain with him when he was an apostate, and
either go to hell with him to his condemnation, or bring him back to
heaven and the Church by way of repentance, not re-baptization. Now, if
any one should ask whether such an apostate, while he continued an
apostate, was a Christian? the answer must be in the negative; but yet
there is something of a Christian in this apostate, that is, his
Baptism; in respect of which he is not so perfectly a no-Christian, as
one that never was baptized.

\textbf{SKELTON,
PRESBYTER.—Vol.
ii. Disc. xxi}.

Our blessed SAVIOUR and MEDIATOR, who hath procured the benefit of this covenant for us by
the ``sacrifice of His Blood,'' hath appointed the Sacrament of
Baptism as the means whereby the contracting parties, GOD and the new Christian, solemnly plight
their promises to each other; and hath likewise made the other
Sacrament, that of His Last Supper, the seal which renews and confirms
the covenant with every penitent transgressor. In both He communicates
the assistance of the HOLY
SPIRIT,
which ``helps our infirmities,'' and enables us, if we are not
shamefully wanting to ourselves, to observe and perform the conditions
promised on our part.

We have already seen, in general, what we are to
expect as the fruits of peace with GOD; namely, eternal life, eternal happiness and glory. Our
present assurance of this is represented in various lights by the
Scriptures. We are made one with CHRIST,
as He is one with the FATHER.
We are united into one Church, or Spiritual Body, whereof ``He is the
head.'' All together ``we are the body of CHRIST, and members in particular.'' Thus joined to Him, who is
by nature the SON
of GOD,
we also become, by a ``new birth in Baptism,'' the adopted sons or
children of GOD.
``We I have received the Spirit of adoption, whereby we cry, Abba,
Father:'' and being taken into the family of GOD,
are made His children by ``faith in JESUS CHRIST.''
The provision made for us is suitable to the grandeur of our new
relation; no less than an eternal kingdom, ``which it is the FATHER'S good pleasure to give us,'' as His
beloved children, and consequently, ``heirs of GOD, and joint heirs with CHRIST:
insomuch much that being one with Him, ``where He is, there
shall we be also,'' partakers both of His nature, and of His inheritance
in happiness and glory. We need not say, since these are the promises of
GOD,
that they cannot possibly fail of performance, provided we do our utmost
to fulfil the promises made on our part.

\textbf{HORNE,
BISHOP AND
DOCTOR.—Vol.
ii. Disc. xviii}.

The first portion of sanctifying grace is given at
Baptism, which is the seal of justification, and the beginning of
sanctification; inasmuch as the sinner being thus sacramentally buried
with CHRIST
into His death, arises with Him in the power of His resurrection,
justified from the guilt of sin through repentance and faith in His
blood, and renewed unto holiness by the operation of His SPIRIT.
This total renewal, at first conferred by the baptismal laver, is styled
\emph{regeneration}, and answers in things natural to the birth of an
infant. But then, as an infant, though born complete in all its parts,
yet comes to its full stature and strength by slow and imperceptible
degrees; by being supplied with proper kinds of food for its nourishment
when in health, and proper medicines for its recovery when otherwise; so
is it with the regenerate spirit of a Christian: while it is (as St.
Peter calls it) a babe in CHRIST, it must be fed with the milk of the
word; when it is more grown in grace, with the strong meat of its
salutary doctrines; when it is infirm, it must be strengthened by the
comforts of its promises; and when sick, or wounded by sin, it must be
recovered and restored by godly counsel and wholesome discipline, by
penance and absolution, by the medicines of the word and sacraments as
duly and properly administered in the Church, by the lawfully and
regularly appointed delegates and representatives of the Physician of
souls. This gradual and complex work of our sanctification is carried
on, through our lives, by the SPIRIT of GOD, given, in due degree and proportion, to
every individual for that purpose. And it is marvellous to behold, as
the excellent Bishop Andrews observes, how, from the laver of
regeneration, to the administration of the \emph{viaticum}, this good SPIRIT helpeth us, and poureth His benefits
upon us, having a grace for every season. When we are troubled with
erroneous opinions, He is the SPIRIT of truth; when assaulted with temptations, He is the SPIRIT
of holiness; when dissipated with worldly vanity, He is the SPIRIT of compunction; when broken with worldly
sorrow, He is the HOLY GHOST
the Comforter. It is He who, after having regenerated us in our baptism,
confirms us by the imposition of hands; renews us to repentance, when we
fall away; teaches us, all our life long, what we know not; puts us in
mind of what we forget; stirs us up when we are dull; helps us in our
prayers; relieves us in our infirmities; consoles us in our heaviness;
gives us songs of joy in the darkest night of sorrow; seals us to the
day of our redemption; and raises us up again in the last day; when that
which was sown in grace shall be reaped in glory, and the work of
sanctification, in spirit, soul, and body, shall be completed.

\textbf{JONES,
PRESBYTER.—On
the figurative language of the Holy Scriptures. Lect. vi}.

p. 156.

As the Ark was prepared by Noah, so hath CHRIST
prepared His Church, to conduct us in safety through the waves of
trouble and the perils of the world, in which so many are lost. And as
the waters of the flood carried Noah and his family into a new world
after the old was drowned; so do the waters of baptism carry us into a
new state with JESUS CHRIST, who passed over the waves of death, and
is risen from the dead. And this practical inference is to be made in
favour of the ordinance of the Church, that as the ark could not be
saved but by water, so must all the Church of CHRIST be baptized.

\emph{Ibid}. p. 167.—We know that Satan has not
that sovereignty over baptized Christians as he hath over men in the
state of nature. After baptism a Christian is no longer the subject of
that tyrant, but the child of GOD,
who undertakes thenceforth to conduct him through all the trials and
dangers of this life to the inheritance promised to the fathers.

\textbf{H}\textbf{EBER,
BISHOP}\textbf{.—Sermons
in England, xviii}.

It (justification) is the same with that
regeneration of which baptism is the outward symbol, and which marks
out, wherever it occurs, (that it ordinarily occurs in baptism, I am for
my own part firmly persuaded,) our admission into the number of the
children of GOD, and the heirs of everlasting happiness. It
is the commencement of that state of salvation, in which, if a man
continues, death has no power over him, in as much as the grave, which
our nature so greatly fears, is to him no extinction of life, but a
passage to a life more blessed and more glorious.

\textbf{JEBB,
BISHOP.—Pastoral
Instructions. Disc. vi}.

p. 112.

But how, may it be asked, are the benefits and
blessings of spiritual regeneration conferred upon infants in their
tender years? To this inquiry we need not be careful to reply: we need
only state, that in this, as in various other instances, it hath pleased
ALMIGHTY GOD to set limits to the presumptuousness of
human curiosity; and thus, at once to try our humility and our faith. It
is enough for us, rest assured, that GOD is now, and ever, the same all-good and gracious Parent; that,
as in times past, it was ``out of the mouths of babes and sucklings he
perfected praise:'' and as ``He revealed unto babes those things which
were hidden from the wise and prudent,'' so He is, at all times,
abundantly able to pour forth the dew of His blessing upon infants who
are faithfully brought to the baptism of his SON.
It is enough for us to believe and cherish the prevalent sentiment of
the universal Church, as it has been maintained from the age of the
Apostles, that at the time of baptism, a new nature is divinely
communicated, and gracious privileges are especially vouchsafed, in such
measure and degree that, whosoever are clothed with this white garment,
may, through GOD'S
help, ``keep their baptism pure and undefiled, for the remainder of their
lives, never wilfully committing any deadly sin.''

\textbf{VAN
MILDERT,
BISHOP.—Bampton
Lectures, vi}.

Regeneration is represented, by a certain class of
interpreters, as an instantaneous, perceptible, and irresistible
operation of the HOLY
SPIRIT
upon the heart and mind; which, whether the person have been baptized or
not, affords the only certain evidence of his conversion to a saving and
justifying faith. By others, it is regarded as a continued and
progressive work of the SPIRIT,
or as a state commencing in baptism, but not completed until, by
perseverance to the end, the individual has ``finished his course, and is
about to enter upon his final reward.'' Others again, separating what the
Scriptures state to be joined together in the work of the new birth,
maintain a distinction between baptismal and spiritual regeneration; the
former taking place in the Sacrament of Baptism; the latter subsequent
to it, and, whether progressive or instantaneous in its operation,
equally necessary with baptism to a state of salvation.

But here the analogy of faith seems to be violated
throughout. For how can any of these views of regeneration consist with
the plain and simple notion of it as an entrance upon a new state, or a
sacramental initiation into the Christian covenant? Nay, how can they
consist with the terms and conditions of the covenant itself? If the
Gospel be a covenant, admission into which, on the terms of faith and
repentance, gives an immediate title to its present privileges, with an
assurance of the spiritual helps necessary for the attainment of
salvation; and if baptism be the divinely-appointed means of admission
into that covenant, and of a participation in those privileges, is not
the person so admitted actually brought into a new state? Has he not
obtained that thing which by nature he cannot have? ``And being thus
regenerated and born anew of water and of the HOLY
GHOST,''
to what subsequent part of his Christian life can a term so peculiarly
expressive of his first entrance upon it be with propriety applied?

\textbf{MANT,
BISHOP.—Bampton
Lectures, vi}.

Verily, verily, I say unto thee, Except a man be
born of water and of the SPIRIT,
\&c. … It should appear, I say, that he was here alluding by
anticipation to the Sacrament of Baptism, which he intended to ordain:
and to that supernatural grace which was thereby to be conferred
through the instrumentality of water, and by the HOLY GHOST;
adopting, not only the ceremony itself, which he meant to exalt to more
noble and spiritual purposes; but also the very term, by which the Jews
had described the change wrought in the baptized, although he
undoubtedly employed it, in a similar sense indeed, but in an infinitely
more dignified sense. To the proselyte from heathenism to the Jewish
faith, baptism had been a death to his natural incapacities, and a new
birth to the civil privileges of a Jew: to him, who should be admitted
to a profession of the Christian faith, and who should be ``born not of
blood, nor of the will of the flesh,'' \&c. … it was to be a death
unto sin, and a new birth unto those spiritual privileges, which should
accompany his deliverance ``from the bondage of corruption into the
glorious liberty of the children of GOD.'' The Jewish proselyte had been baptized with water: the
Christian was to be baptized, not with water only, but with the HOLY
GHOST.

OXFORD,

\emph{The Feast of St. Michael and All Angels}.

[FOURTH EDITION.]

———————————————————————

These Tracts are continued in
Numbers, and sold at the price of 2d. for each sheet, or 7s. for 50
copies.

LONDON: PRINTED FOR
J. G. F. \& J. RIVINGTON,

ST. PAUL'S CHURCH YARD, AND WATERLOO PLACE.

1840.

\xchapter{On Purgatory (Against Romanism.—No. 3.)}

Tract No. 79

THE
extract from Archbishop Ussher's Answer to a Jesuit, contained in
Tract 72, on the subject of the ancient Commemorations for the Dead in CHRIST,
may fitly be succeeded by an inquiry as to what degree and sort of proof
remains for the Roman tenet of Purgatory, after deducting from the
evidence those usages or statements of the early Church, which are
commonly supposed, but, as Ussher shows, improperly, to countenance it.
Ussher's explanations have had the effect, it is presumed, of cutting
away the \emph{primâ facie} evidence, on which the doctrine is usually
rested; and it now remains to see what is left when it is withdrawn.
With this view it is proposed in the following pages to draw out in
detail the evidence alleged by the Romanists in behalf of their belief,
with such remarks as may be necessary, in order to form a fair estimate
of it. A plain statement of the doctrine itself, and of its rise, shall
be also attempted, as not unseasonable at a time when the strength of
Romanism rests in no small degree in its opponents mistaking the points
in debate, and making or refuting propositions which but indirectly or
partially bear upon the errors which they desire to combat.

Before commencing, it is necessary to warn the
reader against estimating the magnitude or quality of any of those
errors by its apparent dimensions in the theory. What seems to be a
small deviation from correctness in the abstract system, becomes
considerable and serious when it assumes a substantive form. This is
especially the case with all doctrinal discussions, in which the
undeveloped germs of many diversities of practice and moral character
lie thick together and in small compass, and as if promiscuously and
without essential differences. The highest truths differ from the most
miserable delusions by what appears to be a few words or letters. The
discriminating mark of orthodoxy, the Homoousion, has before now been
ridiculed, however irrationally, as being identical, all but the letter \emph{i},
with the heretical symbol of the Homoiousion. What is acknowledged in
the Arian controversy, must be endured without surprise in the Roman, in
whatever degree it occurs. We may be taunted as differing from the
Romanists only in phrases and modes of expression; and we may be
taunted, or despised, according to the fate of our Divines for three
centuries past, as taking a middle, timid, unsatisfactory ground,
neither quite agreeing nor quite disagreeing with our opponents. We may
be charged with dwelling on trifles and niceties, in a way inconsistent
with plain, manly good sense; but in truth it is not we who are the
speculatists, and unpractical controversialists, but they who forget
that \emph{hæ nugæ seria ducunt in mala}.

But again there is another reason, peculiar to the
Roman controversy, which occasions a want of correspondence between the
appearance presented by the Roman theology in theory, and its appearance
in practice. The separate doctrines of Romanism are very different, in
position, importance, and mutual relation, in the abstract, and when
developed, applied, and practised. Anatomists tell us that the skeletons
of the most various animals are formed on the same type; yet the animals
are dissimilar and distinct, in consequence of the respective
differences of their developed proportions. No one would confuse between
a lion and a bear; yet many of us at first sight would be unable to
discriminate between their respective skeletons. Romanism in the theory
may differ little from our own creed; nay, in the abstract type, it
might even be identical, and yet in the actual framework, and still
further in the living and breathing form, it might differ essentially.
For instance, the doctrine of Indulgences is in the theory entirely
connected with the doctrine of Penance; that is, it has relation solely
to \emph{this world}, so much so that Roman apologists sometimes speak
of it without even an allusion to its bearings elsewhere: but we know
that in practice it is mainly, if not altogether, concerned with the
next world,—with the alleviation of sufferings in Purgatory.

And further still, as regards the doctrine of
Purgatorial suffering, there have been for many ages in the Roman Church
gross corruptions of its own doctrine, untenable as that doctrine is
even by itself. The decree of the Council of Trent, which will presently
be introduced, acknowledges the fact. Now we believe that those
corruptions still continue; that Rome has never really set herself in
earnest to eradicate them. The pictures of Purgatory so commonly seen in
countries in communion with Rome, the existence of Purgatorian
societies, the means of subsistence accruing to the clergy from belief
in it, afford a strange contrast to the simple wording and apparent
innocence of the decree by which it is made an article of faith. It is
the contrast between a drug in its lifeless seed, and the same
developed, thriving, and rankly luxuriant in the actual plant.

And lastly, since we are in no danger of becoming
Romanists, and may bear to be dispassionate, and (I may say)
philosophical in our treatment of their errors, some passages in the
following account of Purgatory are more calmly written than would
satisfy those who were engaged with a victorious enemy at their doors.
Yet, whoever be our opponent, Papist or Latitudinarian, it does not seem
to be wrong to be as candid and conceding as justice and charity allow
us. Nor is it unprofitable to weigh accurately how much the Romanists
have committed themselves in their formal determinations of doctrine,
and how far, by GOD'S merciful providence, they had been restrained and
overruled; and again how far they must retract, in order to make amends
to Catholic truth and unity.

§ 1. Statement of the Roman
Doctrine concerning Purgatory

§ 2. Proof of the Roman Doctrine concerning
Purgatory

§ 3. History of the Rise of the Doctrine of
Purgatory, and Opinions in the Early Church concerning it

§ 4. The Council of Florence

———————

\xsubsection{§ 1. Statement of the Roman Doctrine concerning Purgatory}

THE
Roman doctrine is thus expressed in the Creed of Pope Pius IV.

Constanter
teneo Purgatorium esse, animasque ibi detentas fidelium suffragiis
juvari.

``I hold without wavering that there is a
Purgatory, and that souls there detained are aided by the suffrages of
the faithful.''

The words of this article are taken from the decree
of the Council of Trent on the subject, (Sess. 25,) which runs as
follows:

``Whereas
the Church Catholic, fully instructed by the HOLY
GHOST, hath from the sacred
Scriptures and ancient tradition of the Fathers, in sacred Councils, and
last of all in this present Œcumenical Synod, taught that there is
Purgatory, and that souls there detained are aided by the suffrages of
the living, and above all by the acceptable sacrifice of the altar, this
holy Synod enjoins on Bishops, to make diligent efforts that the sound
doctrine concerning Purgatory, handed down from the holy Fathers and
sacred Councils, be believed, maintained, taught, and everywhere
proclaimed by the disciples of CHRIST.
At the same time, as regards the uneducated multitude, let the more
difficult and subtle questions, such as tend not to edification nor
commonly increase piety, be excluded from popular discourses. Moreover,
let them disallow the publication and discussion of whatever is
uncertain or suspicious; and prohibit whatever is of a curious or
superstitious nature, or savours of filthy lucre, as the scandals and
stumbling-blocks of believers. And let them provide, that
the
suffrages of believers living, that is, the sacrifices of masses,
prayers, alms, and other works of piety, which believers living are wont
to perform for other believers dead, be performed according to the rules
of the Church, piously and religiously; and whatever are due for them
from the endowments of testators, or in other way, be fulfilled, not in
a perfunctory way, but diligently and accurately by the Priests and
Ministers of the Church, and others who are bound to do this service.''

Such is the Roman doctrine; and taken in the mere
letter there is little in it against which we shall be able to sustain
formal objections. Purgatory is not spoken of at all as a place of pain;
it need only mean, what its name implies, a place of purification. There
is indeed much presumption in asserting definitively that there is such
a place; and assuredly there is not only presumption, but very great
daring and uncharitableness in including belief in it, as Pope Pius'
Creed goes on to do, among the conditions of salvation; but if we could
consider it as confined to the mere opinion that that good which is
begun on earth is perfected in the next world, the tenet would be
tolerable. The word ``\emph{detentas}'' indeed expresses a somewhat
stronger idea; yet after all hardly more than that the souls in
Purgatory would be happier out of it than in it, and that they cannot of
their own will leave it; which is not much to grant. Further, that the
prayers of the living benefit the dead in CHRIST,
is, to say the least, not inconsistent, as Ussher shows us, with the
primitive belief. So much as to the \emph{letter} of the decree; but it
is not safe to go by the letter: on the contrary, we are bound to take
the universal and uniform doctrine taught and received \emph{in} the
Roman Communion, as the real and true interpreter of words which are in
themselves comparatively innocent. What that doctrine is, may be
gathered from the words of the Catechism of Trent, in which the spirit
of Romanism, not being bound by they rules which shackle it in the
Council, speaks out. The account of Purgatory which that formulary
supplies, shall here be taken as our text, and Cardinal Bellarmine's
Defence shall be used as a comment upon it.

The Catechism then speaks as follows:

``Est
Purgatorius ignis, quo piorum animæ ad definitum tempus cruciatæ
expiantur,
ut eis in æternam patriam ingressus patere possit, in quam nihil
coinquinatum ingreditur.''—Part i. De Symb. 5.

``There is a Purgatorial fire, in which the souls
of the pious are tormented for a certain time, and cleansed, in order
that an entrance may lie open to them into their eternal home, into
which nothing defiled enters.''

In like manner Bellarmine says,

``Purgatory is a certain place in
which, as if in a prison, souls are purged after this life, which have
not been fully purged in it, in order, (that is,) that thus purged they
may be enabled to enter heaven, which nothing defiled shall enter.''

A painful light is at once cast by these comments
on the Synodal Decree. ``There is a Purgatory'' in the Decree, is
interpreted by Bellarmine ``there is a sort of \emph{prison};'' and by
the Catechism, ``there is a Purgatorial \emph{fire}.'' And whereas the
Decree merely declares that souls are ``detained there,'' the
Catechism says they are ``tormented and cleansed.'' Moreover, both the
Catechism and Bellarmine imply that this is the ordinary mode of
attaining heaven, inasmuch as no one scarcely can be considered, and no
one can be surely known, to leave this world ``fully purged;'' whereas
the Decree speaks vaguely of ``the souls there.'' So much at first
sight; now to consider the persons with which Purgatory is concerned,
the sins, condition of souls, place, time, punishment, and remedies;
Bellarmine likening it to a \emph{carcer}, the Catechism saying that the
``animæ \emph{piorum} ad \emph{definitum} tempus cruciatæ expiantur
purgatorio \emph{igne}.''

1. The Persons who are reserved in Purgatory

THE
Roman Church holds that Christians or believers only are tenants of
Purgatory, as for Christians only are offered their prayers, alms, and
masses. The question follows, whether all Christians? not all
Christians, but such as die in GOD'S
favour, yet with certain sins unforgiven. Some Christians die simply in
GOD'S
favour with all their sins forgiven; others die out of His favour, as
the impenitent, whether Christians or not; but others, and that the
great majority, die, according to the Romanists, in GOD'S
favour, yet more or less under the bond of their sins. And so far we may
unhesitatingly allow to them, or rather we ourselves hold the same,
if we hold that after Baptism there is no plenary pardon of sins in this
life to the sinner, however penitent, \emph{such} as in Baptism was once
vouchsafed to him. If for sins committed after Baptism we have not yet
received a simple and unconditional absolution, surely penitents from
this time up to the day of judgment may be considered in that double
state of which the Romanists speak, their persons accepted, but certain
sins uncancelled. Such a state is plainly revealed to us in Scripture as
a real one, in various passages, to which we appeal as well as the
Romanists. Let the case of David suffice. On his repentance Nathan said
to him, ``The LORD also \emph{hath put away thy sin; thou shalt not die; howbeit},
because by this deed thou hast given great occasion to the enemies of
the LORD
to blaspheme, the child also that is born unto thee shall surely die.''
2 Sam. xii. 13, 14. Here is a perspicuous instance of a penitent
restored to GOD'S favour at once, yet his sins
afterwards visited; and it needs very little experience in life to be
aware that such punishments occur continually, though no one takes them
to be an evidence that the sufferer himself is under GOD'S displeasure, but rather accounts them
punishments even when we have abundant proofs of his faith, love,
holiness, and fruitfulness in good works. So far then we cannot be said
materially to oppose the Romanists. They on the other hand agree with us
in maintaining that CHRIST'S
death \emph{might}, if GOD
so willed, be applied for the removal even of these specific punishments
of sins, which they call \emph{temporal} punishments, as fully as it
really is for the acceptance of the \emph{soul} of the person punished,
or the removal of eternal punishment. Further, both parties agree, that
in matter of fact it is not so applied; the experience of life shows it;
else every judgment might be taken as evidence of the person suffering
it being under GOD'S wrath. The death of the disobedient
prophet from Judah would, in that case, prove that he perished
eternally, which surely would be utterly presumptuous and uncharitable.
As far as this then we have no violent difference of \emph{principle}
with the Romanists; but at this point we separate from them: \emph{they}
say these temporal punishments on sin are inflicted on the faults
incurring them, in a certain fixed proportion; that every sin of a
certain kind has a definite penalty or price; in consequence, that
if it is not fully discharged in this life, it must be hereafter; and
that Purgatory is the place of discharging it.

2. The sins for which persons are confined in
Purgatory

The next question is, what are the sins which are
thus punished? not all sins of Christians, for some incur an eternal
punishment. There are sins, it is maintained, which in themselves merit
eternal damnation, are directly opposed to love or charity, quench
grace, and throw the doer of them out of GOD'S
favour. These in consequence are called mortal; such as murder,
adultery, or blasphemy. Such sins do not lead to Purgatory; hell is
their portion if unrepented of. But all but these, all but \emph{unrepented
mortal} sins are in the case of Christians punished in Purgatory. Of
these it follows there are two kinds, sins though \emph{repented} of,
and sins though not \emph{mortal}; concerning which a few words shall be
said.

1.
Mortal sins, though repented of, and though the offender cease to be
under GOD'S
displeasure, yet have visibly their own punishment in many cases, as in
the instance of David. But the Romanists consider that these sins have
their penalty assigned to them as if by weight and measure; moreover,
that we can ourselves take part in discharging it, and by our own act
anticipate and supersede GOD'S
judgment, according to the text: ``If we would judge ourselves, we
should not be judged.'' This voluntary act on our part is called
Penance, and is said to expiate the sin, that is, to wash away its
temporal effects. Should we die before the full temporal punishment, or
satisfaction, has been paid for all our mortal sins, we must pay the
rest hereafter, \emph{i.e}. in Purgatory.

2.
Sins which are not mortal, are called \emph{venial}, and are such as do
not quench grace, or run counter to love. Bellarmine thus contrasts
them:

``Mortal sins are they which
absolutely turn us from GOD,
and merit eternal punishment; Venial, those which somewhat impede our
course to Him, but do not turn it, and are with little pains blotted
out. The former are crimes, the latter sins ... Mortal sin is like a
deadly wound, which suddenly kills: Venial is a slight stroke, which
does not endanger life, and is easily healed. The former fights with
love, which is the soul's life; the latter is rather beside than
against love.''—\emph{De Amiss}. \emph{Grat}. i. 2.

Venial sin differs from Mortal in two ways, in \emph{kind}
and \emph{degree}. An idle word, excessive laughter, and the like, are
sins in kind distinct from perjury or adultery. Again, anger is a venial
sin when slight and undesigned, but when indulged interferes with love
and is mortal; a theft of a large sum may be mortal, of a small venial.

Venial sins, being such, are considered by
Romanists not to deserve so much as eternal punishment,—to be
pardonable not merely by an express and immediate act of GOD'S mercy, or again through the virtue of our state of
regeneration, but to be intrinsically venial, to offend GOD,
but not so as to alienate Him. They rest this doctrine upon such
passages as the following: ``Sin, when it is \emph{finished}, bringeth
forth \emph{death},'' James i. 15; therefore, before it is finished or
perfected, it has no such fearful power. Still they say it requires \emph{some}
punishment; which it receives in the next world, should it not receive
it in this, that is, in Purgatory.

Such then are the sins of GOD'S true servants, penitent believers, for
which, according to the Romanists, they suffer in Purgatory; mortal sins
repented of, and those sins of infirmity which befal them so continually
and so secretly, that they cannot repent of them specifically if they
would, and which do not deserve eternal punishment, though they do not.
They consider the Purgatorial punishment of \emph{venial} sins to be
meant by the Apostle, when he speaks of those who, building on the true
foundation ``wood, hay, and stubble,'' are ``saved so \emph{as by fire};''
and the punishment for \emph{mortal} sins, in our SAVIOUR'S
declaration, that certain prisoners shall not go out till they have ``paid the very last mite.'' Luke xii. 59. It may be added, that
Martyrdom is supposed to be a full expiation of whatever guilt of sin
still rests on the Christian undergoing it; and therefore to stand
instead of Purgatory. Martyrs then are at once admitted to the Beatific
vision, which is the privilege in which Purgatory terminates.

From this account of the inmates of Purgatory, and
the causes why they are there detained, we gather what has already been
hinted, that the one main or rather sole reason of the appointment, is a
satisfaction of GOD'S
justice. The persons concerned are believers destined for bliss eternal;
but before they pass on from earth to heaven, the course of their
existence is, as it were, suspended, and they are turned aside to
discharge a debt; how they effect it, or in what length of time, or with
what effect on themselves, being questions as beside the mark, as if
they were used with reference to the payment of a charge in worldly
matters. It is an appointment altogether without bearing upon their
moral character or eternal prospects; and after it is over, is wiped out
as though it had never been.

3. The moral condition of souls in Purgatory

Bellarmine well illustrates the supposed mental
state of believers while in Purgatory, by comparing them to travellers
who come up to a fortified town after nightfall, and have to wait at the
gates till the morning. Such persons have come to the end of their
journey; they are not on the way, they have attained; they are sure of
admittance, which is a matter of time only.

Accordingly the Romanists hold that souls in
Purgatory become neither better nor worse, neither sin nor add to their
good works; they are one and all perfect in love, and ready for heaven,
were it not for this debt, which hangs about them as so much rust or
dross, and cannot be purged away except for certain appointed external
remedies. They support this view of the stationary condition of the soul
in Purgatory by such texts as the following: ``The night cometh when no
man can work.'' ``Where the tree falls, there it shall lie.'' ``We
must all appear before the judgment-seat of CHRIST,
that every one may receive the things \emph{done in his body}.'' John
ix. 4. Eccles. xi. 3. 2 Cor. v. 10.

Next, with the exception of some few theologians,
they consider that souls in Purgatory are comforted with the assurance
that their eternal happiness is secured to them. Their state in
consequence is thus described by Bellarmine (ii. 4).

``You
will object that they may be in doubt whether they are in hell or
in
purgatory. Not so; for in hell GOD
is blasphemed, in Purgatory He is praised; in hell there is neither
habit of faith, nor hope, nor love of GOD,
in Purgatory all of these. A soul then which shall understand that it
hopes in GOD, praises and
loves GOD, will clearly know
it is not in hell. But perhaps it will fear it is to be sent to hell,
though not there yet; neither can this be, for the same faith remains in
it, which it had here. Here it believed according to the plain word of
Scripture, that after death none can become of good bad, or of bad good,
and none but the bad are to be sent into hell. When then it perceives
that it loves GOD, and is
therefore good, it will not fear damnation.''

4. The place and time of Purgatory

On this subject the Church has not formally
determined any thing: but the common opinion of the Schoolmen is, that
it is one of four prisons or receptacles, which are situated in the
heart of the earth, Hell for the damned, the \emph{Limbus Puerorum} for
children dying without baptism, the \emph{Limbos Patrum} for the just
who died before the passion of CHRIST, and who since that time have all
been transferred from it to heaven, and Purgatory for believers under
punishment. In other words, whereas all punishment is either for a time
or eternal, either positive (\emph{pœna sensûs}) or negative (\emph{pœna
damni}), that of good men before CHRIST'S
coming was the \emph{pœna damni}, or absence of GOD'S light and joy for a time, that of
unbaptized infants is the \emph{pœna damni} for ever, that of Purgatory
the \emph{pœna sensûs} for a time, that of Hell the \emph{pœna sensûs}
for ever. To these some Romanists have added a fifth, that is, of
faithful souls, who without being yet admitted into heaven, are yet
secured against all pain; but these, according to Bellarmine, as at
least enduring the \emph{pœna damni}, are to be considered in
Purgatory, though in the most tolerable place in it, as being but in the
condition of the old Fathers before CHRIST came.

The time of Purgatory depends of course upon the
state of the debt which is to be liquidated in each case, and varies
consequently with the individual. Martyrs, as has been above stated, are
supposed to satisfy it in the very act of martyrdom; others will not be
released till the day of judgment. Again, the period of suffering
depends upon the exertions of survivors, by prayers, alms, and
masses, which have power not only to relieve but to shorten the pain.

5. The nature of the Punishment

Here the Roman Church has defined nothing; its
catechism, as we have seen, and its theologians in accordance, consider
it to be material fire, but in the Council of Florence, the Greeks would
not do more than subscribe to the existence o Purgatory; they denied
that the punishment was fire; the question accordingly remains open,
that is, it is not determined either way \emph{de fide}. The difficulty,
how elementary fire, or any thing of a similar nature, can affect the
disembodied soul, is paralleled by St. Austin by the mystery of the
union of soul and body.

The pains of Purgatory are considered to be
horrible, and far exceeding any in this life; ``\emph{Pœnas Purgatorii
esse atrocissimas; et cum illis nullas pœnas hujus vitæ comparandas,
docent constanter Patres},'' says Bellarmine (ii. 14), and proceeds
to refer to Austin, Pope Gregory, Bede, Anselm, and Bernard. Yet on this
point theologians differ. Some consider the chief misery to consist in
the \emph{pœna damni}, or absence of GOD'S
presence, which to holy souls, understanding and desiring it, would be
as intolerable as extreme thirst or hunger to the body; and in this way
seem to put all purgatorial pain on a level, or rather assign the
greater pain to the more spiritually-minded. Others consider the \emph{pœna
damni} to be alleviated by the certainty of heaven and of the
continually lessening term of their punishment. With them then the \emph{pœna
sensûs}, or the fire, is the chief source of torment, which admits
of degrees according to the will of GOD.

6. The efficacy of the suffrages of the Church

By suffrages are meant co-operations of the living
with the dead; prayers, masses, and works, such as alms, pilgrimages,
fastings, \&c. These aids, which individuals can supply, alms,
prayers, \&c., only avail when offered by good persons; for he who is not accepted himself, cannot do acceptable service for another.
Moreover these aids may be directed either to the benefit of all souls
in Purgatory indiscriminately, or specially to the benefit of a certain
soul in particular.

There is one other means of escaping the penalties
due to sin in Purgatory, which may briefly be mentioned, viz. by the
grant of indulgences; these are dispensed on the following theory.
Granting that a certain fixed temporal penalty attached to every act of
sin, in such case, it would be conceivable that, as the multitude of
Christians did not discharge their total debt in this life, so some
extraordinary holy men might more than discharge it. Such are the
Prophets, Apostles, Martyrs, Ascetics, and the like, who have committed
few sins, and have undergone extreme labours and sufferings, voluntary
or involuntary. This being supposed, the question rises, what becomes of
the overplus; and then there seems a fitness that what is not needed for
themselves, should avail for their brethren who are still debtors. It is
accordingly stored, together with CHRIST'S merits, in a kind of treasure-house, to be dispensed
according to the occasion, and that at the discretion of the Church. The
application of this treasure is called an Indulgence, which stands
instead of a certain time of penance in this life, or for the period,
whatever it be, to which that time is commuted in Purgatory. In this
way, the supererogatory works of the Saints are supposed to go in
payment of the debts of ordinary Christians.

\xsubsection{§ 2. Proof of the Roman Doctrine concerning Purgatory}

1. Proofs from supernatural appearances

THE
argumentative ground, on which the belief in Purgatory was actually
introduced, would seem to lie in the popular stories of apparitions
witnessing to it. Not that it rose in consequence of them historically,
or that morally it was founded in them; only that when persons came to
ask themselves why they received it, this was the ultimate ground of
evidence on which the mind fell back; viz. the evidence of miracles, not
of Scripture, or of the Fathers.

Bellarmine enumerates it as one of the confirmatory
arguments. With this view he refers in particular to some relations of
Gregory of Tours, A.D.
573; of Pope Gregory, A.D.
600; of Bede, A.D.
700; of Peter Damiani, A.D.
1057; of St. Bernard, A.D.
1100; and of St. Anselm, A.D.
1100. The dates are worth noticing, if it be true, as is here assumed,
that such supernatural accounts as then were put forth, are really the
argument on which the doctrine was and is received; for it would thence
appear, first, that the doctrine was not taught as divine before the end
of the sixth century; next, that when it was put forth, it was grounded
on the authority of an (alleged) \emph{new revelation}. This indeed is
confessed in the Dialogues of Pope Gregory, in which after he has
related to his friend Peter the Deacon, an instance of a soul which was
seen in the Purgatorial fire, Peter remarks, ``How is this, that \emph{in
these last times} so many things come to light about souls, which
before were hidden?'' lib. iv. 40. But the passage will hereafter come
before us in another connection. The following miraculous narratives are
found in a Protestant Selection from Roman writers, published in 1688,
and entitled ``Purgatory proved by Miracles.''

``St. Gregory the Great writes,
that the soul of Paschasius appeared to St. Germanus, and testified to
him, that he was freed from the pains of Purgatory for his prayers.

``When the same St. Gregory was
abbot of his Monastery, a monk of his, called Justus, now dead, appeared
to another monk, called Copiosus, and advertized him, that he had been
freed from the torments of Purgatory, by thirty
Masses, which Pretiosus, Prefect of the Monastery, by the order of St.
Gregory, had said for his soul, as is recounted in his life.

``St. Gregory of Tours writes of a
holy damsel, called Vitaliana, that she appeared to St. Martin, and told
him she had been in Purgatory for a venial sin which she had committed,
and that she had been delivered by the prayers of the Saint.

``Peter Damiani writes, that St.
Severin appeared to a clergyman, and told him that he had been in
Purgatory, for not having said the Divine Service at due hours, and that
afterwards GOD had
delivered him, and carried him to the company of the blessed.

``St. Bernard writes, that St.
Malachy freed his sister from the pains of Purgatory by his prayers; and
that the same sister had appeared unto him, begging of him that relief
and favour.

``And St. Bernard himself by his
intercession freed another, who had suffered a whole year the pains of
Purgatory: as William Abbot writes in his life.''—\emph{Flowers of the
Lives of the Saints}, p. 830.

These instances among others are adduced by
Bellarmine; and he adds, ``plura similia legi possunt apud, \&c. …
\emph{sed quæ attulimus, sunt magis authentica}.''—1. 11.

2. Proofs from the Old and New Testaments

Bellarmine adduces the following texts from the Old
and New Testaments; in doing which he must not be supposed to mean that
each of them contains in itself the evidence of its relevancy and
availableness, or could be understood without some authoritative
interpretation; only, if it is asked, ``is Purgatory the doctrine of
Holy Scripture, and where?'' he would answer, that in matter of fact it
is taught in the following passages, according to the explanations of
them found in various writers of consideration.

1. 2 Mac. xii. 42-45. ``Besides that noble Judas
exhorted the people to keep themselves from sin, forsomuch as they saw
before their eyes the things that come to pass for the sins of those
that were slain. And when he had made a gathering throughout the company
to the sum of two thousand drachms of silver, he sent it to Jerusalem, \emph{to
offer a sin offering}, doing therein very well and honestly, in that
he was mindful of the Resurrection; for if he had not hoped that they
that were slain should have risen again, it had been superfluous and
vain to pray for the dead. And also, in that he perceived that
there was great favour laid up for those that died godly, \emph{it was an
holy and good thought}. Whereupon he made \emph{a reconciliation for the
dead, that they might be delivered from sin}.''

2. Tob. iv. 17. ``\emph{Pour out thy bread on the
burial of the just}, but give nothing to the wicked;'' that is, at
the burial of the just, give alms; which were given to gain for them the
prayers of the poor.

3. 1 Sam. xxxi. 13. ``And they took their
bones,'' [of Saul and his Sons,] ``and \emph{buried} them under a tree
at Jabesh, and \emph{fasted seven} days.'' Vid. also 2 Sam. i. 12; iii.
35. This fasting was an offering for their souls.

4. Ps. xxxviii. 1. ``O LORD, rebuke me not in Thy \emph{wrath};
neither chasten me in \emph{Thy hot displeasure}.'' By \emph{wrath} is
meant Hell; by \emph{hot displeasure} Purgatory.

5. Ps. lxvi. 12. ``We went through \emph{fire} and
through \emph{water}, but Thou broughtest us out into a wealthy place''
(\emph{refrigerium}). Water is Baptism; fire is Purgatory.

6. Is. iv. 4. ``When the LORD shall have washed away the filth of the
daughters of Zion, and shall have \emph{purged} the blood of Jerusalem
from the midst thereof, by the spirit of judgment and by the spirit of \emph{burning}.''

7. Is. ix. 18. ``Wickedness burneth as the fire;
it shall devour the briers and thorns.''

8. Mic. vii. 8, 9. ``Rejoice not against me, O
mine enemy; when I fall, I shall arise; when \emph{I sit in darkness},
the LORD
shall be a light unto me. I will bear the \emph{indignation} of the LORD, because I have sinned against Him,
\emph{until he plead my cause}, and execute judgment for me: \emph{He will
bring me forth to the light}, and I shall behold His
righteousness.''

9. Zech. ix. 11. ``As for Thee also, by the blood
of Thy covenant, I have \emph{sent forth Thy prisoners out of the pit},
wherein is no water.'' This text is otherwise taken to refer to the \emph{Limbus
Patrum}.

10. Mal. iii. 3. ``He shall sit as a refiner and
purifier of silver; and He shall \emph{purify} the sons of Levi, and \emph{purge}
them as gold and silver,'' \&c.

From the New Testament he adduces the following
texts;

1. Matt. xii. 32. ``Whosoever speaketh against the
Holy Ghost, it shall not be forgiven him, neither in this world, \emph{neither
in the world to come};'' that means, ``neither in Purgatory,'' for
in hell the very supposition of forgiveness is excluded.

2. 1 Cor. iii. 15. ``He himself shall be saved; \emph{yet
so as by fire}.''

3. 1 Cor. xv. 29. ``Else what shall they do, which
are \emph{baptized}'' \emph{i.e}. who undergo the baptism of tears and
humiliation, who pray, fast, give alms, \&c. ``for the dead, if the
dead rise not at all?''

4. Matt. v. 25, 26.—Luke xii. 58, 59. ``Agree
with thine adversary quickly, whilst thou art in the way with him; lest
at any time the adversary deliver thee to the judge, and the judge
deliver thee to the officer, and thou be cast into prison. Verily, I say
unto thee, \emph{thou shalt by no means come out thence}, till thou hast
paid the uttermost farthing.'' By the \emph{way}, is meant this present
life; by the \emph{adversary}, the Law; by the \emph{Judge}, our SAVIOUR;
by the \emph{officer}, or executioner, the Angels; by the prison,
Purgatory.

5. Matt. v. 22. ``Whosoever is angry with his
brother without a cause, shall be in danger of the \emph{judgment}; and
whosoever shall say to his brother, Raca, shall be in danger of the \emph{Council}:
but whosoever shall say, Thou fool, shall be in danger of hell fire.''
Here are three kinds of punishments spoken of. Hell belongs to the next
world; therefore also do the other two. Hence there are in the next
world, besides eternal punishment, punishments short of eternal.

6. Luke xvi. 9. ``Make to yourselves friends of
the mammon of unrighteousness, that, when ye fail, they may receive you
into everlasting habitations.'' To \emph{fail}, is to die: the \emph{friends}
are the Saints in glory, and they \emph{receive} us, \emph{i.e}. from
Purgatory, in consequence of their prayers.

7. Luke xxiii. 42. ``Lord, \emph{remember} me, \emph{when}
Thou comest into Thy kingdom.'' That is, there is a remembrance and a
remission of sin, not only in this life, but after it, in
Christ's future kingdom.

8. Acts ii. 24. ``Whom God hath raised up, having
loosed the pains of death (\emph{inferi}); because it was not possible
that He should be holden of it.'' CHRIST
Himself was released from no \emph{pains} on being raised, nor were the
ancient Fathers in the \emph{Limbus}: nor were lost souls released at
all. Therefore the pains which GOD
loosed, were those of souls in Purgatory.

9. Phil. ii. 10. ``That at the name of JESUS
every knee should bow, of things in heaven, and things in earth, \emph{and
things under the earth}.'' Vid. also Rev. v. 3. ``And no \emph{man}
in heaven, nor in earth, \emph{neither under the earth}, was able to
open the book, neither to look thereon.''

Now as to many of these texts, we who have not been
educated in the belief of Purgatory, may well wonder how they come to be
enlisted in support of Purgatory at all. This may be explained in some
such way as the following,—which may be of use in helping us to
understand the state of mind under which the Romanists view them. It is
obvious, as indeed has been already remarked, that they do not of
themselves \emph{prove} the doctrine, nor are they chosen by Bellarmine
himself, but given on the authority of writers of various times. Could
indeed competent evidence be brought from other quarters, that the
doctrine really was true and Apostolical, we should not unreasonably
have believed that some of them did allude to it; especially if writers
of name, who might speak from tradition, so considered. We could not
have taken upon ourselves to say at first sight that it certainly was
not contained in them, only we should have waited for evidence that it
was. Some of the texts in question are obscure, and seem to desiderate a
meaning; and so far it is a sort of gain when they have any meaning
assigned them, as though they were unappropriated territory which the
first comer might seize. Again, the coincidence of several of them in
one and the same mode of expression, implies that they have a common
drift, whatever that drift is,—that there is something about them
which seems to have reference to secrets untold to man. Amid these dim and broken lights, the text in the Apocrypha first quoted, comes as
if to combine and steady them. All this is said \emph{by way of analysing
how} it is that such a class of texts, though of so little cogency
critically, has that influence with individuals, which it certainly
sometimes has. The reason seems to be that the doctrine of Purgatory
professes to interpret texts which GOD'S
word has left in obscurity. Yet whatever be the joint force of such
arguments from Scripture, in favour of the doctrine, it vanishes surely,
at once and altogether, before one single clear text, such as the
following: ``Blessed are the dead \emph{which die in the Lord}, from
henceforth; yea, saith the Spirit, \emph{that they may rest from their
labours}.'' Or again, if any one is destined to endure Purgatory for
the temporal punishment of sins, one should think it would be persons
circumstanced as the thief on the cross,—a dying penitent; yet to him
it is expressly said, ``Verily I say unto thee, \emph{to day} shalt
thou be with me in Paradise.''

3. Proofs from Antiquity

After Scripture, Bellarmine brings the testimony of
early Churches in Council, as follows:

1. The African Church: ``Let the Altar Sacrament
be celebrated fasting; if, however, there be any Commendation of the
Dead made in the afternoon, let prayers only be used.''—Conc. Carth.
IV. c. 79.

2. The Spanish enjoins that suicides should not be
prayed for, \&c.—Conc. Bracar. I. c. 39.

3. The Gallic: ``It has seemed fit, that in all
celebrations of the Eucharist, the Lord shall be interceded with in a
suitable place in Church, for the spirits of the dead.''—Conc.
Cabilon.

4. The German defines, (Conc. Wormat. c. 10.) that
prayers and offerings should be made even for those who are executed.

5. The Italic declares (Conc. VI. under Symmachus),
that it is sacrilege to defraud the souls of the dead of prayer, \&c.

6. The Greek in like manner.

Moreover, the Liturgies of St. James, St. Basil,
\&c. all contain prayers for the dead.

Now these professed instances are here enumerated
in order to show how plainly and entirely they fall short of the point
to be proved. Not one of them implies the doctrine of Purgatory; or goes
beyond the doctrine which Archbishop Ussher (vide Tract 72.) has shown
to have existed in the early Church, that the Saints departed were not
at once in their full happiness, and that prayers benefited them. One of
these instances indeed is somewhat remarkable, the allowing prayers for malefactors
executed; but all were the subject of prayer who were not excluded from
hope, and malefactors are, even by us, admitted to Holy Communion, and
are allowed the Burial Service. To pray for them was merely the
expression of hope.

Next, Bellarmine appeals to the Fathers, of whom I
shall only cite those within the first five hundred years; viz.
Tertullian, Cyprian, Eusebius, Cyril of Jerusalem, Gregory Nazianzen,
Ambrose, Jerome, Chrysostom, Paulinus, Augustine, Theodoret, and one or
two others. Now in order to keep the point in controversy clearly in
view, let it be recollected that we are not disputing the existence in
the Ritual of the Church, of the custom of praying for the dead in CHRIST;
but \emph{why} prayer was offered was a question in dispute, a point
unsettled by any Catholic tradition, but variously treated by various
Doctors at various times. There is nothing contrary to the genius of
religion, natural and revealed, that duties should be prescribed, yet
the reasons for them not told us, as Bishop Butler has abundantly
showed; and the circumstance that the ancients do agree in the usage,
but differ as to the reasons, shows that the reasons were built upon the
usage, not the usage on the reasons. And while this variety of opinions
in the early Church, as to the meaning of the usage, forfeits for any
one of these any claim to be considered apostolical, of course it
deprives the doctrine of Purgatory of authority inclusively, even
supposing for argument's sake it was received by some early writers as
true. Purgatory is but a violent hypothesis to give meaning to a usage,
for which other hypotheses short of it and very different from it, and
equally conjectural with it, may be assigned, nay, and were assigned before it, and far more extensively. Let it be remembered then,
when the following lists of passages, professedly in behalf of
Purgatory, is read, that what we have to look for is, not evidence of a
certain usage, which we grant did exist, but of an opinion, of a
particular opinion explaining it; not of Prayer for the dead simply, nor
of the opinion that Prayer for the dead profits, but that such Prayer is
intended and tends to rescue them from a state of suffering. Further,
what we look for is not the testimony of one or two writers to the truth
of this opinion, even if one or two could be brought, but an agreement
of all in its favour. If however it be said that the usage of Prayer in
itself tends to the doctrine of Purgatory, I answer, that so far from
it, in its primitive form it included prayers for the Virgin Mary and
Apostles, which while retained were an indirect but forcible standing
witness against the doctrine.

Tertullian, in his de Coronâ, § 3. speaks of `` oblationes pro defunctis,'' offerings for the dead.

Again, ``Let her'' [the widow] ``pray for the
soul of'' [her deceased husband] ``and ask for him a place of
refreshment in the interval before the judgment, and a fellowship in the
first resurrection, and let her offer on the anniversary of his falling
to sleep.''—\emph{De Monogam}. § 10. Vid, also \emph{de Pudicit}.

Cyprian. ``The Bishops our predecessors ...
decreed that no one dying should nominate clerics as guardians or
executors, and if any one had done this, no offering should be made for
him, or sacrifice celebrated for his sleeping well.'' —\emph{Epist}.
i. 9. \emph{et infra}.

Eusebius (vid. Constant. iv.) says that Constantine
had wished to be buried in a frequented Church, in order to have the
benefit of many prayers. On his death they offered the Holy Eucharist
over his remains \footnote{Vid. also passage in Records of the Church, No xii. ``The Adversary contrived that his [Polycarp's] poor body might not
be obtained by us, though many much desired to secure it, and to
communicate over his holy remains.''}.

Cyril of Jerusalem. ``We pray for all our
community who are dead, believing that this is the greatest benefit to
those souls for whom the offering is made.''—\emph{Mystagog}. 5.

Gregory Nazianzen. ``Let us commend to GOD
our own souls, and the souls of those who, as men more advanced on the
same road, have arrived before us at their resting place.''—\emph{Orat.
in Cæsar. fin}.

Ambrose. ``Therefore she is, I think, not so much
to be lamented as to be followed with your prayers; she is not to be
mourned over with your tears, but rather her soul is to be commended to
GOD by
your oblations.''—\emph{Ep}. ii. 8. \emph{ad Faustinum}. Vid. also \emph{de
ob. Theod., \&c. \&c}.

Jerome. ``Other husbands scatter on their wives'
graves violets, roses, lilies, and purple flowers; but our Pammachius
waters her holy ashes and reverend relics with the balsams of
almsgiving; with such embellishments and perfumes he honours the
sleeping remains, knowing what is written, 'As water quenches fire, so
doth alms sin.'''—\emph{Ad Pammach}.

Chrysostom. ``The dead is aided not by tears, but
by prayers, by supplications, by alms … Let us not weary in giving aid
to the dead, offering prayers for them.''—\emph{Hom}. 41. \emph{in}
1. \emph{ad Cor}.

Again. ``Not without purpose has it been ordained
by the Apostles, that in the awful Mysteries a commemoration should be
made of the dead; for they know that thence much gain accrues to them,
much advantage.''—\emph{Hom}. 69. \emph{ad pop}. Vid. also \emph{Hom}.
32. \emph{in Matt}. \emph{In Joan}. \emph{Hom}. 84. \emph{In Philipp}.
3. \emph{in Act. Apost}. 21.

Paulinus, writing to Delphinus, Bishop of Bordeaux:
``Do thy diligence that he may be granted to thee, and that from the
least of thy sacred fingers the dews of refreshment may sprinkle his
soul.''

Augustine. ``We read in the book of Maccabees that
sacrifice was offered for the dead; but though it were not even found in
the Old Scriptures, the authority of the universal Church is not slight,
which is explicit as to this custom, viz. that in the Priests' prayers
which are offered to the LORD
GOD at
His altar, the commendation of the dead is included.''—\emph{De Cur, pro
mortuis}, c. ii. \emph{et alibi}.

Theodoret (Hist. v. 26.) mentions that Theodosius
the younger fell down at the tomb of St. John Chrysostom, and prayed for
the souls of his parents, then dead, Arcadius and Eudoxia.

Isidore. ``Unless the Church Catholic believed
that sins are remitted to the dead in Christ, she would not do alms, or
offer sacrifice to GOD
for their spirits.''—\emph{De off. div}. i. 18.

Gregory the Great. ``Much profiteth souls even
after death the sacred oblation of the lifegiving Sacrifice, so that the
souls of the dead themselves sometimes seem to ask for it.''—\emph{Dial}.
iv. 55.

Again: ``They who are not weighed down by grievous
sins, are profited after death by burial in the Church, because that
their relatives, whenever they come to the same sacred places, remember
their own kin whose tombs they behold, and pray to the LORD for them.''

It is evident that the above passages go no way to
prove the point in debate, being nothing more in fact than Ussher allows
to be found in the early Fathers. They contain the musings of serious
minds feeling a mystery, and attempting to solve it, at least by
conjecture. They state that prayers benefit the dead in CHRIST,
but \emph{how}, is either not mentioned, or vaguely, or hesitatingly, or
discordantly. Accordingly, Bellarmine begins anew, and draws out a
series of authorities for the doctrine of \emph{Purgatory} expressly;
and this certainly demands our attention more than the former. It
contains such as the following:—

For instance, Origen says that ``he who is saved,
is saved by fire, that if he has any alloy of lead, the fire may melt
and separate it, that all may become pure gold.''—\emph{Hom}. 6. \emph{in
Exod}.

Tertullian speaks of our being ``committed into
the prison beneath, which will detain us till every small offence is
expiated, during the delay of the resurrection.''—\emph{De Anim}. 17.

Cyprian contrasts the being purged by torment in
fire, and by martyrdom.—Epist. iv. 2.

Gregory Nazianzen speaks of the last Baptism being ``one of fire, not only more bitter, but longer than the first
Baptism.''—\emph{In Sancta lum. circ. fin}.

Ambrose speaks of our being ``saved through faith,
as if through fire,'' which will be a trial under which grievous
sinners will fall, while others will pass safe through it.—In \emph{Pa}.
xxxvi.

Basil speaks of the ``Purgatorial fire,'' in cap.
ix. \emph{Isa}.

Gregory Nyssen, of ``our recovering our lost happiness by prayer and religiousness in this life, or after death by the \emph{purgatorial
fire}.''—\emph{Orat. pro Mort}. Elsewhere too he speaks of the
Purgatorial fire.

Eusebius Emissenus uses such determined words, as to require quoting. ``This punishment under the earth will await those, who, having lost
instead of preserving their Baptism, will perish for ever; whereas those
who have done deeds calling for temporal punishments, shall pass over
the fiery river and that fearful water the drops of which are fire.''

Hilary declares that we have to undergo ``that ever-living fire, which
is a punishment of the soul in cleansing of sin.''—In \emph{Ps}.
cxviii. Lactantius speaks to the same effect.—\emph{Div. Inst}. vii.
21.

Jerome contrasts the eternal torments of the devil, and of atheists and
infidels, with ``the judgment tempered with mercy, of sinners and
ungodly men, yet Christian, whose works are to be tried, purified in the
fire.''—\emph{In fin. Comment. in Is}. In another place in a like
contrast he speaks of Christians, if overtaken in a fault, being saved
after punishment.—Lib. i. \emph{in Pelag}.

Augustine has various passages in point, such as \emph{Civ. Dei}, xxi.
24, where, speaking of believers who die with lighter sins, he says, ``It is certain that these being purified before the day of judgment by
means of temporal punishment, which their souls suffer, are not to be
given over to eternal fire.'' Pope Gregory the first expresses the same
doctrine, as do some others.

These instances are at first sight to the point, and demand serious
consideration. Yet there is nothing in them really to alarm the inquirer
whither he is being carried. I say this, that no one may be surprised at
the deliberateness and over-patience with which I may seem to loiter
over the explanation of them. First, then, let it be observed, were they
ever so strong in favour of something more than we believe, it does not
therefore follow that they take that very view which the Romanists take,
nay, it does not necessarily follow that they take any one view at all,
or agree with each other. Now it so happens neither the one nor the
other of these suppositions is true, as regards those passages, though
they ought both to hold, if the Roman doctrine is to be satisfactorily maintained. These Fathers, whatever they teach, do not
teach Purgatory, they do not teach any one view at all on the subject.
Romanists consider Purgatory to be an article of faith, necessary to be
believed in order to salvation; or, in Bellarmine's words, ``Purgatory is an article of faith, so that he who disbelieves its
existence, will never have experience of it, but will be tormented in
hell with everlasting fire.'' Now it can only be an article of faith,
supposing it is held by Antiquity, and that unanimously. For such things
only are we allowed to maintain, as come to us from the Apostles; and
that only (ordinarily speaking) has evidence of so originating, which is
witnessed by a number of independent witnesses in the early Church. We
must have the unanimous ``consent of Doctors,'' as an assurance that
the Apostles have spoken; and much less can we tolerate their actual
disagreement, in a case where unanimity was promised us. Now as regards
Purgatory, not only are early writers silent as to the modern view of
Rome, but they do not agree with each other; which proves they knew
little more about the matter than ourselves, whatever they might
conjecture; that they possessed no Apostolic Tradition, only at most
entertained floating opinions on the subject. Nay, it is obvious, if we
wished to believe them, we could not; for \emph{what} is it we are to
believe? if, as I shall show, various writers speak various things,
which of their statements is to be taken? If this or that, it is but the
language of an individual: if all of them at once, a doctrine results
discordant in its details, and in general outline, if it have any, vague
and imperfect at the best.

Now as to the passages quoted by Bellarmine, it will be observed that in
the number are extracts from the works of Origen, St. Ambrose, St.
Hilary, St. Jerome, and Lactantius. He introduces the list with these
words, ``\emph{Sunt apertissima loca in Patribus, ubi asserunt Purgatorium,
quorum pauca quædam afferam},'' i. 10. ``There are most
perspicuous passages in the Fathers, in which they assert Purgatory, of
which I will adduce some few.'' Will it be believed that in his second
book these Fathers, nay, for the most part in the very extracts, which
he has before adduced in proof of the doctrine, are enumerated as at
variance with it, and mistaken in their notion of it? He quotes a
passage of Origen, (not the same) the very same two passages from St.
Ambrose, the very same passage from St. Hilary, the very same from
Lactantius, and a passage (not the same) from St. Jerome. Then he says, ``\emph{Hæc sententia, accepta ut sonat, manifestum errorem continet};
for'' (he proceeds) ``it is defined in the Council of Florence,
\&c'' ii. 1. Next he observes, ``\emph{Adde, quod Patres adducti,
Origene excepto … videntur sano modo intelligi posse}.'' At
length, after he has given the two most favourable explanations
assignable to their words, he adds of one of the two, ``\emph{Sane hanc
sententiam} [\emph{quæ docet omnes transituros per ignem, licet non
omnes lædendi sint ab igne}] \emph{nec auderem pro vera asserere, nec
ut errorem improbare}.'' ``The only alleviation of this strange
inconsistency,'' says a work which has recently appeared, ``is that he
quotes not the very same sentences both for and against his Church, but
adjoining ones.'' The work referred to, thus comments on Bellarmine's
conduct, as throwing light upon the state of feeling under which
Romanists engage in controversy. ``A Romanist,'' the writer says,''
cannot really argue in defence of the Roman doctrines. He has too firm a
confidence in their truth, if he is sincere in his profession, to enable
him critically to adjust the due weight to be given to this or that
evidence. He assumes his Church's conclusion as true; and the facts or
witnesses he adduces, are rather brought to receive an interpretation
than to furnish a proof. His highest aim is to show the mere consistency
of his theory, its possible adjustment, with the records of antiquity. I
am not here inquiring how much of high but misdirected moral feeling is
implied in this state of mind; certainly as we advance in perception of
the truth, we all of us become less fitted to be controversialists. If
this, however, be the true explanation of Bellarmine's strange error,
the more it tends to exculpate him, the more deeply it criminates his
system. He ceases to be chargeable with unfairness, only in proportion
as the notion of the infallibility of Rome is admitted to be the
sovereign and engrossing tenet of his communion, the foundation stone,
or (as it may be called) the fulcrum of its theology. I consider then,
that when he first adduces the aforementioned Fathers in proof of
Purgatory, he was really but interpreting them; he was teaching what
they ought to mean, what in charity they must be supposed to mean, what
they might mean as far as the very words went, \emph{probably} meant
considering the Church so meant, and might be taken to mean, even if
their authors did not so mean, from the notion that they spoke vaguely,
and, as children, really meant something besides what they formally
said, and that after all they were but the spokesmen of the then
existing Church, which, though in silence, held, as being the Church,
the same doctrine which Rome has since defined and published. This is to
treat Bellarmine with the same charity with which he has on this
supposition treated the Fathers, and it is to be hoped, with a nearer
approach to the matter of fact. So much as to his first use of them: but
afterwards, in noticing what he considers erroneous opinions on this
subject, he treats them, not as organs of the Church infallible, but as
individuals, and interprets their language by its literal sense or by
the context, and in consequence condemns it ... How hopeless then is it
to contend with Romanists, as if they practically agreed to the
foundation of faith, however much they pretend to it! Ours is antiquity:
theirs the existing Church. Its infallibility is their first principle;
belief in it is a deep prejudice, quite beyond the reach of any thing
external. It is quite clear that the combined testimonies of all the
Fathers, supposing such a case, would not have a feather's weight
against a decision of the Pope in Council, nor would matter at all,
except for the Fathers' sakes who had by anticipation opposed it. They
consider that the fathers ought to mean what Rome has since decreed, and
that Rome knows their meaning better than they themselves did. That
venturesome Church has usurped their place, and thinks it merciful, only
not to banish outright the rivals she has dethroned. By an act, as it
were, of grace she has determined, that when they contradict her, though
not available as witnesses against her, yet, as living in times of
ignorance, they are only heterodox, and not heretical; and she keeps
them around her, to ask their advice when it happens to agree with her
own.

``Let us then understand the position of the Romanists towards us;
they do not really argue from the Fathers, though they seem to do so.
They may affect to do so on our behalf, happy if by an innocent
stratagem they are able to convert us; but all the while in their own
feelings, they are taking a far higher position. They are teaching, not
disputing or proving. They are interpreting what is obscure in
antiquity, purifying what is alloyed, correcting what is amiss,
perfecting what is incomplete, harmonizing what is various. They claim
and use all its documents as ministers and organs of that one infallible
Church, which once forsooth kept silence, but since has spoken, which by
a divine gift must ever be consistent with itself, and which bears with
her her own evidence of divinity.''

Leaving Bellarmine then, let us proceed to inquire
what the opinion of the Fathers in the foregoing passages really is.

\xsubsection{§ 3. History of the Rise of the Doctrine of Purgatory and Opinions in the Early Church concerning it}

THE
argumentative ground of the doctrine of Purgatory as far as the
Infallibility of the Church has not superseded any, has ever been, I
conceive, the report of miracles and visions attesting it; but the
historical origin is to be sought elsewhere, viz. in the anxious
conjectures of the human mind about its future destinies, and the
apparent coincidences of these with certain obscure texts of Scripture.

These may be supposed to have operated as follows;
as described in the work already cited. ``How ALMIGHTY GOD
will deal with the mass of Christians, who are neither very good nor
very bad, is a problem with which we are not concerned, and which it is
our wisdom, and may be our duty, to put from our thoughts. But, when it
has once forced itself upon the mind, we are led in self-defence, with a
view of keeping ourselves from dwelling unhealthily on particular cases,
which come under our experience and perplex us, to imagine modes, not by
which GOD does, (for that would be presumptuous to
conjecture,) but by which He may solve the difficulty. Most men, to our
apprehensions, are too unformed in religious habits either for heaven or
for hell, yet there is no middle state when CHRIST
comes in judgment. In consequence it is obvious to have recourse to the
interval before His coming, as a time during which this incompleteness
might be remedied; a season, not of changing the spiritual bent and
character of the soul departed, whatever that be, for probation ends
with mortal life, but of developing it into a more determinate form,
whether of good or of evil. Again, when the mind once allows itself to
speculate, it will discern in such a provision, a means whereby those,
who not without true faith at bottom yet have committed great crimes, or
those who have been carried off in youth while still undecided, or who
die after a barren, though not an immoral or scandalous life, may
receive such chastisement as may prepare them for heaven, and render it
consistent with GOD'S
justice to admit them thither. Again, the inequality of the sufferings
of Christians in this life, compared one with another, would lead the
unguarded mind to the same speculations, the intense suffering, \emph{e.g}.
which some men undergo on their death-bed, seeming as if but an
anticipation in their case of what comes after death upon others, who
without greater claims on GOD'S
forbearance, have lived without chastisement and die easily. I say, the
mind will inevitably dwell upon such thoughts, unless it has been taught
to subdue them by education or by the experience of their dangerousness.

``Various suppositions have, accordingly, been made,
as pure suppositions, as mere specimens of the capabilities, (if one may
so speak,) of the Divine Dispensation, as efforts of the mind reaching
forward and venturing beyond its depth into the abyss of the divine
counsels. If one supposition could be produced to solve the problem, ten
thousand others were conceivable, unless indeed the resources of GOD'S Providence are exactly commensurate with man's discernment
of them. Religious men, amid these searchings of heart, have naturally
gone to Scripture for relief, to see if the inspired word anywhere gave
them any clue for their inquiries. And from what was there found, and
from the speculations of reason upon it, various notions have been
hazarded at different times; for instance, that there is a certain
momentary ordeal to be undergone by all men after this life, more or
less severe according to their spiritual state; or that certain gross
sins in good men will be thus visited, or their lighter failings and
habitual imperfections; or that the very sight of divine perfection in
the invisible world will be in itself a pain, while it constitutes the
purification of the imperfect but believing soul; or that, happiness
admitting of various degrees of intensity, penitents late in life may
sink for ever into a state, blissful as far as it goes, but more or less
approaching to unconsciousness; infants dying after baptism may be as
gems paving the courts of heaven, or as the living wheels in the Prophet's
vision; while matured Saints may excel in capacity of bliss, as well as
in dignity, the highest Archangels. Such speculations are dangerous
when indulged; the event proves it; from some of these in fact seems to
have resulted the doctrine of Purgatory.

``Now the texts to which the minds of the early
Christians seem to have been principally drawn, and from which they
ventured to argue in behalf of these vague notions, were these two: 'The
fire shall try every man's work,' \&c.; and 'He shall baptize you
with the Holy Ghost and with fire.' These texts, with which many more
were found to accord, directed their thoughts one way, as making mention
of '\emph{fire},' whatever was meant by the word, as the instrument of
trial and purification; and that, at some time between the present time
and the judgment, or at the judgment. And accordingly, without perhaps
having any definite or consistent meaning in what they said, or being
able to say whether they spoke literally or figuratively, and with an
indefinite reference to this life, as well as to the intermediate state,
they sometimes named fire as the instrument of recovering those who had
sinned after their baptism. That this is the origin of the notion of a
Purgatorial fire, I gather from these circumstances, first, that they do
frequently insist on the texts in question, next, that they do not agree
in the particular sense they put upon them. That they quote them shows
they rest upon them; that they vary in explaining them, that they had no
Catholic sense to guide them. Nothing can be clearer, if these facts be
so, than that the doctrine of the Purgatorial fire in all its senses, as
far as it was more than a surmise, and was rested on argument, was the
result of private judgment exerted in defect of Tradition, upon the text
of Scripture …

``As this doctrine, thus suggested by certain
striking texts, grew in popularity and definiteness, and verged towards
its present Roman form, it seemed a key to many others. Great portions
of the books of Psalms, Job, and the Lamentations, which express the
feelings of religious men under suffering, would powerfully recommend it
by the forcible and most affecting and awful meaning which they received
from it. When this was once suggested, all other meanings would seem
tame and inadequate.

``To these must be added various passages from the
Prophets, as that in the beginning of the 3rd chapter of Malachi,
which speaks of fire as the instrument of judgment and purification when
CHRIST
comes to visit His Church.

``Moreover there were other texts of obscure and
indeterminate bearing, which seemed on this hypothesis to receive a
profitable meaning; such as our LORD'S
words in the Sermon on the Mount, 'Verily, I say unto thee, thou shalt
by no means come out thence till thou hast paid the uttermost farthing;'
and St. John's expression in the Apocalypse, that 'no man in heaven, nor
in earth, \emph{neither under the earth}, was able to open the book.'

``Further, the very circumstance that no second
instrument of a plenary and entire cleansing from sin was given after
Baptism, such as Baptism, led Christians to expect that the unknown
means, when accorded, would be of a more painful nature than that which
they had received so freely and instantaneously in infancy, and
confirmed, not only the text already cited, 'He shall baptize you with
the Holy Ghost and with fire,' but also St. Paul's announcement of the 'judgment
and fiery indignation' which awaits those who sin after having been once
'enlightened,' and by CHRIST'S warning to the impotent man to sin no
more lest a worse thing come unto him.

``Lastly, the universal and apparently apostolical
custom of praying for the dead in CHRIST,
called for some explanation, the reason for it not having come down to
posterity with it. Various reasons may be supposed quite clear of this
distressing doctrine, but it supplied an adequate and a most
constraining motive for its observance to those who were not content to
practise it in ignorance.''

Should any one for a moment be startled by any
thing that is here said, as if investing the doctrine with some approach
to plausibility, I would have him give GOD
thanks for the safeguard of Catholic Tradition, which keeps us from
immoderate speculation upon Scripture or a vain indulgence of the
imagination, by authoritatively declaring the contents and the limits of
the Creed necessary to salvation and profitable to ourselves.

There seem, on the whole, to be two chief opinions
on the subject embraced in the early Church. One of these is Origen's,
which I shall first exhibit in the language of St. Ambrose, being
the very passage referred to by Bellarmine. The notion is this, that the
fire at the day of judgment will burn or scorch every one in proportion
to his remaining imperfections. St. Ambrose then thus comments on Psalm
xxxvii. (38) 14.

``'Thou hast proved us by fire,' says
David; therefore we shall all be proved by fire, and Ezekiel (Malachi)
says, 'Behold the LORD
ALMIGHTY
cometh, and who may abide the day of his coming? \&c. … for he is
like a refiner's fire, and like fullers' soap; and He shall sit as a
refiner and purifier of silver; and He shall purify the sons of Levi,
and purge them as gold and silver, \&c.' Therefore the sons of Levi
will be purged by fire; \emph{by fire Ezekiel, by fire Daniel}. But
these, though proved by fire, yet shall say, 'We \emph{passed} through
fire and water,' (Ps. lxvi. 12.) Others shall remain in the fire: and
the fire shall be as dew to them, (Song of Three Children, 27.) as to
the Hebrew Children who were exposed to the fire of the burning furnace.
But the Ministers of impiety shall be consumed in the avenging flame.
Woe is me should my work be burned, and I suffer this worsting of my
labour! Although the LORD
will save His servants, we shall be saved by faith, but so saved as by
fire. Although we shall not be consumed, yet we shall be burned. But how
some remain in the fire, others escape through it, learn from another
Scripture. The Egyptians were drowned in the Red Sea, the Israelites
passed over; Moses escaped to land, Pharaoh sank, for his heavy sins
drowned him. In like manner the irreligious will sink in the lake of
burning fire.''

It is plain that St. Ambrose, so far from imagining
a Roman Purgatory, definite in period, place, and subjects, speaks of an
ordeal by fire which \emph{all} Christians must undergo at the last day,
and grounds it on the solemn text already referred to, 1 Cor. iii.
12-15. which whether rightly so interpreted or not, a point we cannot
determine, since it is an [\emph{apax legomenon}] in Scripture, yet at
least may be so understood without violence to the wording. ``If any man
build upon this foundation, gold, silver, precious stones, wood, hay,
stubble, every man's work shall be made manifest; for the Day shall
declare it, because it (the Day) shall be revealed in fire; and the fire
shall try every man's work, of what sort it is. If any man's work abide
which he hath built thereupon, he shall receive a reward. If any man's
work shall be burned, he shall suffer loss; but he himself shall be
saved, yet so as by fire.'' Now it would seem plain that in this passage
the searching process of final Judgment, essaying our works of
righteousness, is described by the word \emph{fire}. Not that we may
presume to \emph{limit} the word fire to that meaning, or on the other
hand to say it is a merely \emph{figurative} expression denoting
judgment; which seems a stretching somewhat beyond our measure.
Doubtless there is a mystery in the word \emph{fire}, as there is a
mystery in the word \emph{day of judgment}. Yet it any how has reference
to the \emph{instrument} or \emph{process} of judgment. And in this way
the Fathers seem to have understood the passage; referring it to the
last Judgment, as Scripture does, but at the same time religiously
retaining the use of the word \emph{fire}, as not affecting to interpret
and dispense with what seems some mysterious economy, lest they should
be wiser than what is written.

Next let us turn to the same Father's 20th Sermon
on Ps. cxix. which is also referred to by Bellarmine.

``As long as the Israelites were in
Egypt, they were in the iron furnace, that is, in the furnace of
temptation, in the furnace of affliction, where they were afflicted by
cruel tyranny. Whence also it is written, 'I brought them forth out of
the land of Egypt, from the iron furnace.' The furnace was iron,
because, while the people was yet in Egypt, no one's works were
illuminated by holiness. no one's gold had been there assayed, no one's
lead of iniquity burned away. It was a cruel furnace, a furnace of
perpetual death, which none could escape, which consumed every one, in
which pain and sorrow dwell only. But the furnace, in which Ananias,
Azarias, and Misael sang their hymn to the Lord was a golden furnace,
not an iron; by means of which wisdom hath shown forth in the faith of
true obedience all over the world. It was indeed in Babylon, where
spiritual gold was not, unless perchance in captivity, for 'the LORD
led captivity captive.' This is the gold in GOD'S
saints who were captives among the Babylonians in body, but in spirit
were freemen with GOD,
delivered from the chains of human captivity, and bearing the yoke of
spiritual grace. And perchance the same furnace would be iron to the
unstable, and gold to those who persevere.

``\emph{All must be proved through fire},
as many as desire to return to Paradise; for it is not said for nothing,
that when Adam and Eve were expelled from Paradise, GOD
placed at the outlet a fiery sword which turned every way. \emph{All must
pass through the flames, whether he be John the Evangelist}, whom the
LORD
so loved as to say to Peter of him, 'If I wish him to tarry, what is
that to thee? Follow thou me.' Some have doubted of his death: of his
passage through the fire we cannot doubt, for he is in Paradise, not
separated from CHRIST.
\emph{Or whether he be Peter}; he who received the keys of the kingdom
of Heaven, who walked upon the sea, must
still say, 'We passed through fire and water, and Thou broughtest us out
into a place of refreshment.' But the fiery sword will soon be turned by
St. John, for iniquity is not found in him, whom righteousness itself
loved. Whatever human defect was in him, Divine Love melted it away; for
her wings are as the wings of fire, (Cant. viii. 6.)

``He who possesses this fire of love,
will have no cause to fear there the fiery sword. To Peter, who so often
exposed his life for CHRIST,
He will say, 'Go and sit down to meat.' But he shall say, Thou hast
tried us with fire, as silver is tried; for, when many waters do not
drown love, how can fire consume then?' But he shall be tried as silver,
I as lead; \emph{I shall burn till the lead melts away}. If no silver be
found in me, ah me! I shall be plunged down into the lowest pit, or
consume entire as the stubble. Should ought of gold or silver be found
in me, not for my works, but through the mercy and grace of CHRIST,
by the ministry of the priesthood, I shall peradventure say, 'They that
hope in Thee, shall not be ashamed.'

``The fiery sword then shall consume
iniquity, which is placed on the leaden scale. One only could not feel
that fire, CHRIST
the Righteousness of GOD,
who did no sin; for the fire found nought in Him which it might consume.''

It is now sufficiently clear what St. Ambrose's
belief was. The only point of approximation between it and the doctrine
of Purgatory is this; that he conceived that for all but the highest
saints, in whom love dissolved all remaining dross whatever, some
transient suffering, more or less in duration, was in store in the day
of judgment. And hence the force of the ordinary prayers of the early
Church, as based on Scripture, (and described at length by Archbishop
Ussher, in Tract, No. 72,) that departed believers might have ``a
merciful trial at the last day.''

St. Hilary is another witness, whom Bellarmine, in
his former book quotes, in his latter surrenders. He, too, will be found
to hold this same view of the purgatorial nature of the fire of the last
judgment.

``The prophet [the Psalmist]
observes, that it is difficult and most perilous to human nature, to
desire GOD'S
judgments: For, since no one is clean in His sight, how can His judgment
be desirable? Considering we shall have to give account for every idle
word, shall we long for the day of judgment, in which we must undergo
that ever-living fire, and those heavy penalties for cleansing the soul
from its sins? Then will a sword pierce through the soul of Mary that
the thoughts of many hearts may be revealed. If that Virgin which could
compass GOD
is to come into the severity of the judgment, who shall dare desire
to be judged of GOD?
Job, when he had finished his warfare with all calamities of man and had
triumphed, who, when tempted, said, 'The LORD
gave,' and confessed himself but [dust and] ashes when he heard GOD'S
voice from the cloud, and determined that he ought not to speak another
word. And who shall venture to desire GOD'S
judgments, whose voice from heaven neither so great a Prophet endured,
nor the Apostles again, when they were with the Lord in the Mount?''—Tract,
in Ps. cxviii. (cxix.) lit. 3. § 12. vid. also § 5.

Again,

``He [John the Baptist] marks the
season of our salvation and judgment in the LORD,
saying, 'He shall baptize you with the HOLY
GHOST and with fire;' for to
those who are baptized in the HOLY
GHOST it remains to be
perfected in the fire of judgment.''—Comm. in Matt. ii. § 4.

Let us now proceed to Origen, who is historically
the first who has put forward the theory under review. Even Origen, be
it remembered, is at first alleged by Bellarmine, though afterwards
absolutely relinquished. His words, as quoted by that author himself,
are as follows:

``I consider that even after the
resurrection from the dead, we need a sacrament to wash us throughly and
cleanse us; for no one will rise without dross upon him, nor can the
soul be found which at once is free from all defects.''—Hom. 14. in
Luc.

Again,

``We must all come to that fire, be
we Paul or Peter,'' in Ps. xxxvii.

Lactantius expresses the same, or almost the same
doctrine in the following passage, as referred to by Bellarmine.

``Moreover, when He shall have judged
the just, He will also try them in the fire. Then they whose sins
prevail in weight or number, will be tortured in the fire and partially
burned; but they, who are mature in righteousness and ripeness of
virtue, shall not feel that flame; for they have somewhat of GOD
within them, to repel and throw off the force of the flame. Such is the
force of innocence, that from it that fire recoils without mischief, as
having received this property from GOD
to burn the irreligious, to recede from the righteous.''—Div. Inst.
vii. 21.

Two more writers may be mentioned, as holding the
same view, both of whom are quoted by Bellarmine in his favour. St.
Jerome, as referred to by him, speaks as follows:

``The fire,'' he says, commenting on
Amos vii. 4, ``being called for judgment, devours first the deep; that
is, all kinds of sins, wood, hay, stubble, and afterwards consumes also
a part, that is, reaches to his saints, who are accounted the LORD'S
portion.''

St. Paulinus of Nola is the other, who thus writes
to Severus:

``If we attain by these works to be
citizens with the saints, our works shall not be burned; and that
sagacious fire will, on our passing its ordeal, surround us with no
severe heat of punishment; but as if we were commended to its care it
will play around us with a kind caress, so that we may say, 'We have
passed through fire and water,' \&c.''—Ep. 28. (9.)

To these passages, others similar might be added
from St. Basil and St. Gregory Nazianzen.

So much on this speculation or foreboding
concerning the fire of the last judgment. Before proceeding to consider
the second notion of a Purgatory, which existed in the early Church, I
stop to make a remark. What has been said will illustrate what is meant
by Catholic Tradition, and how it may be received without binding us to
accept every thing which the Fathers say. It must be \emph{Catholic} to
be of authority; that is, all the writers who mention the subject, must
agree together in their view of it, or the exceptions, if there be any,
must be such as \emph{probare regulam}. And again, they must profess it \emph{is}
Traditionary teaching. For instance, supposing all the Fathers agreed
together in their interpretation of a certain text, I consider that
agreement would invest that interpretation with such a degree of
authority, as to make it at first sight most rash (to say the very
least) to differ from them; yet it is conceivable that on some points,
as the interpretation of unfulfilled prophecy, they might be mistaken.
It is abstractedly conceivable, that a modern commentator \emph{might}
on certain occasions plausibly justify his dissent from them:—this is
conceivable, I say, \emph{unless} they were explaining a doctrine of the
Creed, which is otherwise known to come from the Apostles,—or
professed, (which would be equivalent) that such an interpretation had
ever been received in their respective Churches as coming from the
Apostles. Catholic Tradition is something more than Catholic teaching.
Great as is the authority of the latter, (and we cannot well put it too
high,) Tradition is something beyond it. This remark is in point
here, for it might be objected that so many Fathers agree together in
the notion of a last-day Purgatory, that, were it not for the accident
of others speaking differently, we should certainly have received it as
Catholic Tradition. I answer, no; whatever the worth of so many
witnesses would have been, and it certainly for safety's sake ought to
have been taken for very much,—still, Origen, Hilary, Ambrose, and the
rest, do not approximate in their remarks to the authoritative language
in which they would speak of the Trinity or the benefits of Baptism.
They do not profess to be delivering an article of the Faith once
delivered to the saints.—Now, to consider the second theory in the
early Church on the subject of Purgatory.

While the Greek Churches, and thence the Italian
held the doctrine of a judgment Purgatory, a doctrine far more like the
Roman is found from an early age in the African Church; at the same
time, it was so far from being considered as a necessary article of
faith, that even St. Austin, who brings it out most fully, expresses his
doubt about its truth. It was in fact only an opinion or conjecture.

Tertullian speaks thus, when discussing the
question, whether souls suffer in the intermediate state, or wait till
the resurrection of the body:

``In short, considering we understand
that prison, which the Gospel discloses, to be the places under the
earth (inferos), and, explain the very last farthing to mean, that every
slightest fault is then to be washed away in the interval before the
resurrection, no one will doubt that the soul pays something in those
nether places without intrenching on the fulness of the resurrection
also through the flesh.''—De Anim. fin.

Next comes St. Cyprian. Cyprian is arguing in
favour of readmitting the lapsed, when penitent, and his argument seems
to be, that it does not follow we absolve them simply, by restoring them
to the Church; we do but admit them to present privileges, the judgment
being reserved in GOD'S
hands. He thus writes to Antonianus.

``Neither suppose, dearest brother,
that the virtue of the brethren will be impaired, or martyrdoms fail,
though penitence be indulged to the lapsed, and hope
of reconciliation set before the penitent. Strength unmoveable abides
with those who have true faith; and to those who fear and love GOD
with their whole heart, integrity endures in firmness and in courage.
Even to adulterers a period of penitence is granted by us, and
reconciliation allowed; yet not on that account does virginity decline
in the Church, or the glorious resolve of continence languish through
the sins of others. The Church is still embellished by the crown of so
many virgins, and chastity and purity are as glorious as before; nor,
though the adulterer is indulged with penitence and pardon, is the
vigour of continence relaxed. It is one thing to stand for pardon,
another to arrive safe at glory; one to be sent to prison, there to
remain till the last farthing be paid; another to receive at once the
reward of faith and virtue; one thing to be tormented for sin in long
pain, and so to be cleansed, and to be purged a long while in the fire,
another to have washed away all sin in martyrdom; one thing in short, to
wait for the LORD'S sentence
in the day of judgment, another at once to be crowned by Him.''—Ep. 55.
ad Antonian.

Rigaltius, Fell, and some others, understand this
passage to refer to the penitential discipline of the Church which was
imposed on the penitent; and, as far as the context goes, certainly no
sense could be more apposite. Yet, if, I may venture on an opinion apart
from such high authorities, the words in themselves seem to go beyond
any mere ecclesiastical, though virtually divine censure, especially ``missum
in carcerem,'' and ``purgari diu igne.''

Further, the passage in Tertullian, weak in itself,
for it was perhaps written after he was a Montanist, fixes a sense,
though it rests for authority on Cyprian's language. Tertullian explains
Cyprian; Cyprian sanctions Tertullian. It should be recollected,
moreover, that Cyprian used to call Tertullian his Master; and the
inference deducible from all this is greatly strengthened; when we come
to consider the views of St. Austin, another African. At the same time
it is worth noticing, the \emph{occasion} and \emph{manner} of St.
Cyprian's statement, whatever it means. He will be found to speak
conjecturally, and as if in disputation. He is \emph{accounting} for a
difficulty; as if he said,—``You suppose that, should the lapsed be
received, this makes it all one as if they had never fallen. Far from
it; they do not receive an \emph{absolute} pardon; they are reserved to
the judgment of the great day. Had they endured and suffered martyrdom,
they would have had their pardon sealed at once; as it is, it is
uncertain, and \emph{who} \emph{knows but} in GOD'S judgments such a recompense is in store
for them as will allow the Church to be merciful to them without GOD'S
ceasing to be just?''

St. Austin is lastly to be mentioned; who speaks
neither in one uniform way, nor with one and the same degree of
certainty. Sometimes he seems to hold the Greek opinion of the final
purgatorial conflagration. In the following passage, after alluding to
Abraham's sacrifice, (Gen. xv.) in which the beasts were divided, but
not the birds, and ``when the sun went down,'' ``a smoking furnace and a
burning lamp passed between those pieces,'' and interpreting the birds of
the spiritual members of the Church, and the beasts of carnal men, some
of whom are within, some outside the Church, he says,

``The smoking furnace will come; for
Abraham sat there till the evening, and then comes the great terror of
the day of judgment. For the evening is the end of the world, and the
furnace is the coming day of judgment. It went between those things
which were already divided, separating them to the right and left. Thus
there are certain carnal men who are yet in the Church's bosom, living
according to their own way, who are in danger of seduction from
heretics. While they remain carnal, they are divisible; He did not
divide the birds, but the carnal are divided. 'I could not speak unto
you as unto spiritual, but as unto carnal.' … Whoso shall remain such,
and in a way of life suitable to the carnal, and yet has not receded
from the bosom of the Church, not been seduced by heretics, so as to be
divided off the other way, the furnace will come, nor will he be able to
stand on the right without undergoing it. If then he would escape that
furnace, let him be changed now into the turtle-dove and pigeon. Let him
receive it, who can. But if not, but he shall have built on the
foundation, wood, hay, stubble; that is, if he has heaped over the
foundation of his faith worldly likings,—yet if CHRIST
be there, so as to have the first place in his heart, above all other
objects, such are endured, are suffered. The furnace shall come, and
shall burn the wood, hay, and stubble; and 'he shall be saved, yet so as
by fire.' This will the furnace do; separating off some to the
left,—others it will in a manner strain off onto the right: but it did
not divide the birds.''—In Ps. civ. Serm. iii. and de Civ. Dei, xvi.
24. vid. also, in Ps. vi. de Civ. Dei, xx. 25; xxi. 16, and in Gen.
contr. Man. ii. 20 fin.

This is one notion St. Austin had of Purgatory;
another was, that it would be of a certain duration, in proportion to
the sins of each individual. Without asserting that this view is plainly
inconsistent with the former, it fairly may be called a distinct one.
The following passage will be found to contain it:

``Some suppose that those who do not
renounce the name of CHRIST,
and are baptized in his font in the Church, nor are cut off therefrom by
any schism or heresy, whatever be their crimes, though neither washed
away by penitence nor ransomed by alms, but persevered in obstinately to
the last day of life, will yet be saved by fire, punished indeed
according to the greatest of their excesses and wickednesses, but not
with eternal fire … But since those clear and positive apostolical
testimonies to the contrary (James ii. 14. 17. 1 Cor. vi. 9, \&c.)
cannot be false, the former obscure text concerning those who build on
this foundation, which is CHRIST,
not gold, silver, precious stones, but wood, hay, stubble, … must be
so explained as not to contradict passages which are clear. Wood, then,
hay, stubble, may naturally mean such desires of lawful things of this
world as cannot be foregone without some pain of mind. But when that
pain burns, if CHRIST abides
in the heart as a foundation, so that nothing is preferred to him, and
the man who feels the fire of that pain, had rather lose the things
which he so loves than CHRIST,
he is saved through fire … The trial of tribulation is a certain fire,
of which Scripture speaks plainly in another place. 'Earthen vessels are
proved by the furnace, and righteous men by the trial of tribulation.'
That fire fulfils the Apostle's words in this life; for instance, should
it befall two Christians, one caring for the things of GOD,
how he may please GOD; that
is, building upon the foundation of CHRIST,
gold, silver, precious stones; the other caring for the things of the
world, how he may please his wife, that is, building on the same
foundation wood, hay, stubble; the work of the former is not burned
away, for he has not loved things the loss of which would distress him;
but the other's work is burned away, since those things are not lost
without suffering which are possessed with enjoyment. But since, when an
alternative comes, he had rather lose them than CHRIST,
nor from apprehension of losing such things renounces CHRIST,
though he may feel a pain during the loss, yet he is 'saved so as by
fire;' for though the loss of what he loved is a burning pain, yet it
does not subvert or consume one who is secured by the firmness and
indestructibility of his foundation. \emph{Such a suffering too, it is not
impossible may happen after this life}; and it is a fair question,
whether it can be settled or not; viz. that some Christians, according
to their love of the perishing goods of this world, attain salvation
more slowly or speedily through a certain purgatorial fire; not such,
however, of whom it is said, 'that they shall not inherit the kingdom of
GOD,' unless they repent
suitably, and gain remission of their crimes.''—Enchirid. 68. 89. vid.
also ad Dulcitium, § 6-13. de Fide et Operib. § 16.

In his \emph{de Civitate Dei}, after speaking (as
above noticed) of the fire at the judgment, he goes on to change its
position in the course of the Divine Economy, and places it between
death and the resurrection; yet still he observes his hesitating and
conjectural tone.

``After the death of the body, until
the arrival of that last day of condemnation and reward after the
resurrection [of the body], should it be said that in this interval the
spirits of the dead suffer a fire, such as they do not feel who had not
habits and likings in the life of this body, which requires their wood,
hay, and stubble to be burned up, but they feel who have not carried
with them the like worldly tabernacles, whether these only, or how and
then, or not then because here, though they experience the fire of
transitory tribulation rescuing venial offences from damnation by
consuming them, \emph{I do not oppose, for perchance it is true}.''

He then proceeds to speak as before, of the other
senses of the word \emph{fire}, as used in the text, which affords
matter for his inquiry.

And now the reader has before him the whole extent
of Augustine's much-talked-of admissions in behalf of Purgatory; and he
may see how hesitating and incomplete they are. It is remarkable that
the passages on which Bellarmine chiefly relies, are rejected by the
Benedictines as not Augustine's; so that Romanists, if they would use
this celebrated Father in the controversy, must betake themselves to
such as the two extracts last quoted, in which Augustine speaks but
doubtfully, and which (it is remarkable) Bellartnine introduces, not in
his own favour, but on an opponent's challenge, to explain, as if from
their conjectural tone rather making against him. It really would
appear, as if in the African Church, there had been no advance in
definiteness of doctrine in this matter since the days of Cyprian; but
that what was a speculation then, remained as little insisted on or
settled when St. Austin wrote.

If it were necessary to add any other evidence, how
little the Fathers knew on this mysterious subject, I might mention,
that in one place St. Austin implies that the impenitent are in
Purgatory; and that St. Jerome seems to say, all baptized persons,
however they suffer in Purgatory, are eventually saved \footnote{Faber's Difficulties of Romanism, i. 12. This reference to Trevern is
made on the authority of Mr. Faber.}.

I have now finished my account of what the early
Fathers said about Purgatory; but very imperfect justice is done to the
subject, till the reader is put into possession of those decisive
testimonies of the Fathers the other way, (that is, in favour of the
peace and rest of the intermediate state to true believers,) which will
reduce the opinions already described to a mere conjecture, pious
indeed and solemnly made, yet received one moment, and abandoned the
next. Without determining whether the strict wording of the following
passages be such as necessarily to exclude the doctrine of Purgatory,
which is a poor way of seeking after what the fact really was, simply
consider whether persons who \emph{practically held} that doctrine, who
kept it simply before them as the whole truth and acted upon it, could
possibly have written them.

Cyprian, on occasion of the famous plague of A.D.
252,

``Let him fear death, who has never
been born anew of water and the Spirit, and is sold over to the flames
of hell; him, who has not been given an interest in the cross and
passion of CHRIST; who is to
pass from temporal to the second death; whose departure from the world
will be followed by the torments of eternal flame of punishment; who by
a longer delay gains but a longer respite from pangs and groans. Many of
our people are dying in this pestilence, that is, are delivered from the
world; and what is truly a plague to Jews, heathen, and enemies of CHRIST,
is to GOD'S servants an end
bringing salvation. That you witness righteous and wicked dying together
without any distinction of man from man, is no reason for your supposing
that destruction is common to good and evil; \emph{the righteous ore called
to a place of refreshment, the wicked are hurried to punishment},
shelter is promptly afforded to the believing, punishment to infidels.
We are undiscerning and ungrateful, well-beloved brethren, in return for
GOD'S benefits, nor do we
recognise the mercy vouchsafed to us. Lo! the virgins depart in peace
safe, and with their glory secured, without the dread of the threats,
the seductions, and the impurities of approaching Anti-Christ; youths
escape the perils of their anxious age, and happily receive the prize of
continence and chastity; the delicate matron no more fears the tortures,
the fury of persecution, the violent hands and the cruelties of the
executioner, receiving the gain of a speedy death. By fear of the
pestilence the lukewarm are kindled, the languid are braced, the
slothful are roused, deserters are driven back, the heathen are
constrained to believe; \emph{the multitude of those who are already
believers is called to peace}; recruits are collected in abundance
and with increased strength, prepared to fight without fear of death,
when the action comes on, as having joined in a season when death was
busy.''—De Mortal. 9.

``Our brethren should not cause us
sorrow, whom the LORD'S call
has delivered from the world, knowing as we do that they are not lost to
us but sent before us, they do not recede, but precede: we should behave
as towards men going a journey or a voyage, regret but not deplore them,
nor go into mourning \emph{for those who have already put on white raiment},''
\&c.—Ibid. 14.

``It is not an exit, but a passage, a
travelling to things eternal, when time has been journeyed through. \emph{Who
would not hasten to what is better}?''—Ibid. 15.

That in this last passage St. Cyprian is speaking
of heavenly felicity after the resurrection, is certain from the
context; but it is as plain that he looks upon the intermediate state as
the beginning of it, or the out-post, which he could not do, unless he
thought that at least, on the whole, and to the generality, it was a
state of rest and peace.

St. Ambrose:

``Death is in every way a good;
because it puts away those principles in us which war against each
other, and because it is \emph{a sort of harbour for those} who after
tossing on the wide sea of this life, \emph{seek for an anchorage of secure
peace}; and because it puts an end to the chance of deterioration,
but, as it finds a man in that condition it consigns him to the future
judgment, and comforts him with the rest itself, and withdraws him from
such present goods as raise envy, and quiets him with the expectation of
the future.''—De Bono Mortis, 4.

``Unwise persons fear death as the
greatest of ills; but the wise desire, \emph{as if a rest after toil, and
the end of ills}.''—Ibid. 8.

``Relying on these considerations,
let us betake ourselves courageously to our Redeemer JESUS;
courageously to the council of Patriarchs, to our father Abraham, when
our day shall arrive; courageously to that holy assembly and
congregation of the just. We shall go to our fathers, to our preceptors
in the faith, so that, though our works fail us, our faith may succour
us, our birthright plead for us. We shall go where holy Abraham opens
his arms to receive the poor, as he received Lazarus; where they rest
who in this life have endured heavy and sharp inflictions … We shall
go to those, who sit down in the kingdom of GOD
with Abraham, Isaac, and Jacob, because when asked to supper they did
not excuse themselves. We shall go thither, where there is a paradise of
delight, where Adam, who fell among thieves, has forgotten to lament his
wounds, where too the thief himself rejoices in the fellowship of the
kingdom of heaven; \emph{where are no clouds}, where no thunder, no
lightning, no storm of wind, no darkness, no evening, no summer, no
winter will vary the seasons. There will be no cold, hail, rain, nor the
presence of this sun, moon, or stars; but the brightness of Light will
alone shine forth.''—Ibid. 12.

St. Hilary:

``The vengeance of hell overtakes us
at once: and immediately we depart from the body, if we have so lived,
we 'perish from the right way.' The rich and poor man in the Gospel show
us this: the one placed by angels in the abode of the blessed and in
Abraham's bosom, the other at once received into the place of
punishment. So quickly did punishment come upon the dead, that even his
brothers were still alive. There is no deferring or delaying there. For,
as the day of judgment is the eternal award either of bliss or
punishment, so the time of death orders the
interval for every man by its own laws, committing every one to Abraham
or to punishment till the judgment.''—In Psalm ii. § 48.

Nazianzen thus speaks on the death of his
father:—

``There is but one life, to look
forward towards life; and one death, even sin, which is the destruction
of the soul. Whatever else men exult in, is but a vision in sleep in
mockery of realities, and a phantom seducing the soul. If these be our
feelings, O my Mother, we shall neither exult in life, nor be much
distressed at death. What heavy misfortune has befallen us, if we have
passed hence to the true life, released from meat and drink, from
dizzinesses, from surfeiting, from base money-getting, and placed amid
stable not transitory possessions, as lesser lights, circling in festive
dance round the Great Luminary?''—Orat. 19. fin.

Macarius, in answer to the question what shall
become of those who have two principles, of sin and grace, within them,
answers that they will go to that place on which their heart is stayed:
for

``The LORD,
beholding thy mind, that thou lightest and lovest Him with thy whole
soul, separateth death from thy soul in one hour, (for it is not for him
to do so,) and receiveth thee unto His bosom and to light. For He
snatcheth thee in an hour's turn from the mouth of darkness, and
forthwith translates thee into His kingdom. For to GOD
all things are easy to do in an hour's turn, so that thou hast the love
of Him.''—Hom. 26.

The hour's space spoken of seems to imply that the
hour of death would supply the necessary purification of the soul from
sin \footnote{Faber's Difficulties of Romanism, i. 12. This reference to Trevern is
made on the authority of Mr. Faber.}; but, whatever it
means, the passage is quite irreconcilable with the Roman tenet, for the
\emph{state} of the dead is made one of bliss, and that ``forthwith'' upon
death. The following passage is to the same effect; after saying that
the guilty soul is upon death carried away by the devil, he proceeds,

``When they'' (the righteous) ``depart
from the body, the choirs of angels receive their souls to their own
place, to the pure world, and so bring them to the LORD.''—Hom.
22.

St. Jerome:

``Let the dead be bewailed, but it
must be he whom hell receives, whom the pit swallows up, for whose
punishment the everlasting fire is in motion. We, whose departure a crowd of angels accompanies, whom CHRIST
goes out to meet, let us rather feel distress, if we have longer to
dwell in this tabernacle of death, for as long as we delay here, we are
pilgrims from the LORD.''—Ep.
25.

So much on the theology of the first five hundred
years. But it may be shown that not even Pope Gregory at the end of that
period, held the doctrine of Purgatory in the modern Roman form of it.
He seems to have gone little further than maintain the Greek notion of
the fire of judgment, as above explained, but, from the circumstance of
his considering the end of the world close at hand, he so expressed
himself as to give it a different character. Nothing has been more
common in every age than to think the day of judgment approaching; and
perhaps it was intended that the Church should ever so suppose. Perhaps
so to suppose is even a mark of a Christian mind; which at least will
ever be on its watch-tower, to see whether it be coming or no, from
desire of its SAVIOUR'S
return. But any how, as at other times, so in St. Gregory's case, this
expectation prevailed; and, as thinking that the end was all but
arrived, he seems to have fancied that ``fire upon earth'' was almost ``kindled,''
that last judicial and purgatorial trial, which the Greeks and some of
the Latins had made attendant upon it. If then he speaks of Purgatory in
language since adopted by Romanism, it was not as intending thereby to
sanction the idea to which it is appropriated in that theology, viz.
that of a regular and ordinary \emph{system} of fiery cleansing in the
intermediate state; but, because he imagined the world was on the eve
and under the incipient symptoms of an extraordinary crisis, when the
sun was to be darkened, and the earth dissolved, and the graves opened,
and all souls to be judged which were in earth and under the earth. He
says,

``As, when night is ending and day
beginning, before the sun rises there is a sort of twilight, while the
remains of the departing darkness are changing perfectly into the
radiance of the day which succeeds, so the end of this world is already
mingling with the commencement of the next, and the very gloom of what
remains has begun to be illuminated with the incoming of things
spiritual.''—Dial. iv. 41.

To the same effect, he puts into the mouth of Peter
the Deacon, the following words: —

``Why is it, I ask, that in these
last times so many things \emph{begin to be clear} about
souls which before were hidden; so that by open revelations and
disclosures the age to come seems forcing itself on us and to be
dawning?''—Ibid. 40.

Conformably with this view, he considered the pains
of Purgatory to be diverse and various in their modes and circumstances,
in this earth as well as under the earth, and consisting in other
torments as well as those of fire, being but the pangs and shudderings
of intellectual natures, when their Judge was approaching, and
disclosing themselves in a supernatural agony parallel to that trembling
of the earth or the failing of the sun, which will precede the
dissolution of the physical world. Occasion has already been taken to
speak of the belief in visions and miracles, as occurring in attestation
of the doctrine, and of the predispositions of the popular mind to
receive it. The state of the evidence, of the popular feeling, and of
the doctrine itself, is strikingly set before the reader in the
following passage of Bishop Jeremy Taylor, though perhaps with somewhat
less of considerateness in the wording of it, than such a subject might
bear.

``The people of the Roman Communion have been
principally led into belief of Purgatory by their fear, and by their
credulity; they have been softened and enticed into this belief, by
perpetual tales and legends, by which they loved to be abused. To this
purpose, their priests and friars have made great use of the apparition
of St. Jerome, after death, to Eusebius, commanding him to lay his sack
upon the corpse of three dead men, that they, arising from death, might
confess Purgatory, which formerly they, had denied. The story is written
in an epistle imputed to St. Cyril; but the ill luck of it was, that St.
Jerome outlived St. Cyril, and wrote his life, and so confuted that
story; but all is one for that, they believe it nevertheless; but these
are enough to help it out and if they be not firmly true, yet if they be
firmly believed, all is well enough. In the \emph{Speculum Exemplorum}
it is said, that a certain priest, in an ecstacy, saw the soul of
Constantinus Turritanus in the eaves of his house, tormented with frosts
and cold rains, and afterwards climbing up to heaven upon a shining
pillar. And a certain monk saw some souls roasted upon spits, like pigs,
and some devils basting them with scalding lard; but a while after,
they were carried to a cool place, and so proved Purgatory. But Bishop
Theobald standing upon a piece of ice to cool his feet, was nearer
Purgatory than he was aware, and was convinced of it, when he heard a
poor soul telling him, that under that ice he was tormented; and that he
should be delivered if for thirty days continual he would say for him
thirty masses. And some such thing was seen by Conrade and Udelric in a
pool of water: for the place of Purgatory was not yet resolved on, till
St. Patrick had the key of it delivered to him; which, when one Nicholas
borrowed of him, he saw as strange and true things there, as ever Virgil
dreamed of in his Purgatory, or Cicero in his dream of Scipio, or Plato
in his Gorgias, or Phædo, who indeed are the surest authors to prove
Purgatory. But, because to preach false stories was forbidden by the
Council of Trent, there are yet remaining more certain arguments, even
revelations made by angels, and the testimony of St. Odilio himself, who
heard the devil complain ... that the souls of dead men were daily
snatched out of his hands, by the alms and prayers of the living; and
the sister of St. Damianus being too much pleased with hearing of a
piper, told her brother, that she was to be tormented for fifteen days
in Purgatory.

``We do not think that the wise men in the Church of
Rome believe these narratives; for if they did, they were not wise; but
this we know, that by such stories the people were brought into a belief
of it, and having served their turn of them, the master builders used
them as false arches and centres, taking them away when the parts of the
building were made firm and stable by authority. But even the better
sort of them do believe them; or else they do worse, for they urge and
cite the Dialogues of St. Gregory, \&c.''—\emph{Dissuasive from Popery},
part. i. ch. i. § 4.

Yet not even after Pope Gregory's times was the
doctrine unhesitatingly received. Ussher (Answer ch. vi.) quotes the
words of the Council of Aix la Chapelle in Charlemagne's time, near 250
years after Gregory, to the effect that there are ``three ways in which
sins are punished; two in this life, and the third in the life to come;
that of the former one is the punishment with which the sinner, GOD inspiring, by penitence, takes vengeance on
himself, the other the punishment which ALMIGHTY GOD inflicts; and that tile third is that of
everlasting fire. He also quotes the author of the tracts \emph{de Vanitate
Sæculi}, and \emph{de Rectitudine Catholicæ Conversationis},
wrongly ascribed to St. Austin; the former of which says, ``Know that
when the soul is separated from the body, presently it is either placed
in paradise for its good works, or plunged into the bottom of hell for
its sins:'' and the latter, ``The departing soul, which is invisible to
eyes of flesh, is received by the angels, and placed either in Abraham's
bosom, if it be faithful, or, if a sinner, in the keeping of the prison
beneath, till the appointed day arrive for it to receive its own body
again and give account of its works before the judgment seat of CHRIST,
the true Judge.'' Even in the days of Otto Frisingensis, A.D. 1146, the doctrine of Purgatory was
considered but a private opinion, not an article of faith universally
received; for he writes, ``\emph{Some affirm} there is in the unseen
state a place of Purgatory, in which those who are to be saved are
either troubled with darkness only, or are refined by the fire of
expiation.''

However, without entering further into the history
of the gradual reception of the doctrine, which, if the circumstances of
its rise be clear, is unnecessary, even could it be given, I conclude
this head of the subject with one or two avowals on the part of
Romanists confirmatory of what has been said.

As to the text of Scripture, we have the candid
admission of the celebrated M. Trevern, present Bishop of Strasburgh,
that it is silent as regards this doctrine, at least so Mr. Faber
understands him.

``Instead of vainly labouring to establish the
doctrine on some one or two misinterpreted texts of the New Testament,
he fairly and honestly confesses, that we have received no revelation
concerning it from JESUS
CHRIST.
Hence he judiciously wastes not his time in adducing passages of Holy
Writ, which are altogether irrelevant. 'Had it been necessary for us,'
says he, 'to be instructed in such questions, JESUS
would doubtless have revealed the knowledge of them. He has not done so.
We can, therefore only form conjectures on the subject more or less
probable.'' \footnote{Faber's Difficulties of Romanism, i. 12. This reference to Trevern is
made on the authority of Mr. Faber.}

It seems then the doctrine is not taught in \emph{Scripture}.
The silence of \emph{Antiquity} concerning it is avowed by Fisher,
Bishop of Rochester, Alphonsus a de Castro, and Polydore Virgil.

Of these the celebrated Cardinal Fisher speaks as
follows:

``It weighs perhaps with many, that
we lay such stress upon indulgences, which are apparently of but recent
usage in the Church, not being found among Christians till a very late
date. I answer, that it is not clear from whom the tradition of them
originated. They are said not to be without precedent among the Romans
from the most ancient times; as may be understood from the numerous
stations in that city. Moreover Gregory the First is said to have
granted some in his own time. We all indeed are aware, that by means of
the acumen of later times many things both from the Gospels and the
other Scriptures are now more clearly developed and more exactly
understood than they once were; whether it was that the ice was not yet
broken by the ancients, and their times were unequal to the task of
accurately sounding the open sea of Scripture, or that it will ever be
possible in so extensive a field, let the reapers be ever so skilful, to
glean somewhat after them. For there are even now a great number of
obscure passages in the Gospel, which I doubt not posterity will
understand much better. Why should we despair of it when the Gospel is
given for this very purpose, to be understood thoroughly and exactly?
Seeing then that the love of CHRIST
towards His Church continues not less strong now than before, nor His
power less, and that the HOLY
GHOST is her perpetual
guardian and restorer, whose gifts flow into her as unceasingly and
abundantly as from the beginning, who can question that the minds of
posterity will be enlightened unto the clear knowledge of those things
which remain still unknown in the Gospel?''

After a sentence or two, he adds:

``Whoever reads the commentaries of
the ancient Greeks, will find no mention, as far as I see, or the
slightest possible concerning Purgatory. Nay, even the Latins did not
all at once, but only gradually, enter into the truth of this matter …
For a while it was unknown, at a late date it was known, to the Church
Universal. Then it was believed by some, by little and little, partly
from Scripture, partly from revelations.''—Assert. Luther. Confutat.
18.

It will be observed how accurately Bishop Fisher's
words bear out, as far as they go, our foregoing account. First, he candidly gives up the Greek Church, and almost gives up the Latin. He
says it was gradually introduced, that at length it became universal.
What can we desire more in disproof of the Roman doctrine? He implies
too, that the doctrine, though not suggested by the plain text of
Scripture, was recommended by it, when once suggested in whatever way;
as if what it did, was just what has been above supposed, viz. bring out
in a touching way a certain possible deep sense which the sacred text
could not be said to teach but might contain; else why should it be
understood only after a long delay? Further, he illustrates and confirms
what has above been observed, that the Church of Rome, relying on its
supposed gift of enunciating the truth, cares not to \emph{prove} its
doctrines ancient, and rather interprets the Fathers by its present
teaching than thinks it necessary to depend upon them. And lastly, he is
a witness that, as far as Rome has cared to argue in this matter, she
has rested the doctrine on \emph{revelations};—a true and honest
account of the matter of fact, but decidedly opposed to the more
accurate, though inapplicable, theory established after his death at
Trent, which is this, that the revelation was concluded once for all in
the Apostles, that all that the Church does is to discriminate and
define their doctrine, and that he is Anathema, though an angel from
heaven, who adds to it. ``That alone is matter of faith,'' says Bellarmine,
``which is revealed by GOD
either mediately or immediately; but divine revelations are \emph{partly
written partly unwritten}. The decrees of Councils, and Popes, and
the consent of Doctors, … \emph{then only} make a doctrine an article
of faith, when they \emph{explain the Word of God or deduce} any thing
from it.'' \footnote{Bellarm. de Purg. i. 15.}

Polydore Virgil appeals to Fisher's statement as
above given, and adds, ``Moreover by the Greeks, even to this day, the
doctrine is not believed.'' Alphonsus de Castro says, ``Concerning
Purgatory there is scarcely any mention, especially among the Greek
writers; for which reason, even to this day, it is not believed by the
Greeks.'' \footnote{These three passages are from Taylor's Dissuasive, part 2. ii. 2.}

Lastly, the following is the avowal of the
Benedictine Editor of St. Ambrose's Works in his preface to the \emph{de
Bono Mortis}, on certain passages concerning the state of the dead,
some of which have been above extracted in the course of these remarks.

``If we interpret the words of our author strictly
and literally, we must plainly confess that in his judgment souls are
kept shut up in certain dwellings till the general resurrection, and
there wait the award due to their deeds, which will not however be paid
them before the last day; meanwhile that they are visited with some good
or punishment, according as each of them has deserved. Lastly, the joy
of the righteous is dispensed according to certain ranks.

``It is not surprising that Ambrose should have
written in this way concerning the state of souls; but what might seem
almost incredible, is, the uncertainty and inconsistency of the Holy
Fathers on the subject from the very times of the Apostles down to the
Pontificate of Gregory XI. and the Council of Florence, that is, for
nearly the whole of fourteen centuries. For, not only do they differ one
from the other, as commonly happens in such questions not yet defined by
the Church, but they are not even consistent with themselves, sometimes
appearing to grant that those souls enjoy the clear sight of the divine
nature, of which at other times they deprive them.''

\xsubsection{§ 4. The Council of Florence}

It remains to give a brief notice of the Council of
Florence, by which the doctrine of Purgatory was first made an article
of faith. With it I shall bring this paper to an end.

The Council of Constance, which had been summoned
principally with a view to the reformation of the clergy, terminated in
April 1418, without having taken any effectual measures for their
object. Five years afterwards the remonstrance which the existing state
of things occasioned, obliged the then Pope Martin V. to summon another,
which, in consequence of his sudden death, eventually opened at Basle,
23d of July, 1431, in the pontificate of Eugenius, under the presidency
of Cardinal Julian Cæsarini. Basle, as being across the Alps, was
removed from the influence of the Roman see: and the Fathers assembled
at once applied themselves to determine a question, which had already
been agitated at Constance, the superiority, viz. of a General Council
to the Pope. They passed a decree that the jurisdiction of the
representatives of the Church Catholic in Council Assembled was supreme
and universal, and that they could not be dissolved, prorogued, or
transferred without their own consent. They proceeded to summon,
threaten, and censure Eugenius; and at length when he resisted their
proceedings, they suspended him from all his powers unless he submitted
to them within 60 days. In these acts they were supported by the Emperor
and other chief powers of Europe, as well as by the clergy; and the Pope
was forced to submit.

They next attempted to reconcile the Greeks to the
Latin Church. At this time Constantinople was much pressed by the
Turkish arms; and the Emperor John Palæologus, the second of that name,
after the example of his father, hoped by holding out the prospect of a
union of the Churches to gain succours from the West. The Fathers of
Basle invited him to attend their meeting with the Patriarch and other
chief ecclesiastics of his division of Christendom; but, on his
objecting to a journey across the Alps, an opening was afforded to
Eugenius, who was not slow to avail himself of it, to propose to the
Greeks to transfer the seat of the Council from the Rhine to Italy. In
spite of the opposition of the Fathers at Basle, Eugenius was successful
in his overtures. The Greek Emperor and ecclesiastics accepted the place
of meeting which he proposed, which was Ferrara, and proceeded thither,
that is, besides Palæologus himself, the Patriarch, and twenty chief
bishops, among whom were the metropolitans of Heraclea, Cyzicus, Nice,
Nicomedia, Ephesus, and Trebizond; representatives also attended from
Alexandria, Antioch, and Jerusalem; and the Primate of Russia. Such were
the members of the Greek Church present at this Council, who, however,
high in station as they were, evidently were too few to express the
voice of the East. It is well known that on the ancient principle of
Councils, decisions were made not by authority, but by the independent
and concordant testimony of all the Bishops of Christendom, or what was
virtually all, to the doctrines declared. On the side of the Latins
there were but five archbishops, eighteen bishops, and ten abbots, the
greater part of whom were subjects or countrymen of the Pope. This
scanty representation however of the Latin Church received, as it
happened, a considerable reinforcement from Basle; for a reaction taking
place there in the Pope's favour, some chief members of the rival
Council coming over to him, the whole number of subscribers which he at
last obtained to the synodical decree, amounted to eight cardinals, two
patriarchs, eight archbishops, fifty-two bishops, and forty-five abbots.
After all, however, these are at first sight scarcely to be considered
representatives of the whole of Christendom; yet such was the
composition of the assembly, known in history as the Council of
Florence, (whither a plague had driven it from Ferrara) which
established the doctrine of Purgatory.

This is a sketch of its external history: but the
point to be considered is the part taken by the Greeks in its
proceedings. At the first glance here is this circumstance, almost in
itself decisive against its authority, that the Greeks were actuated by
motives of interest, and at least by the influence and the presence of a
Sovereign. Were they in number fifty times as many, they would not
have appeared in Italy at all, had not the Ottomans been at the gates of
Constantinople. Next they were unprotected in a strange country,
depending even for their daily food on the bounty of those who were bent
upon the reconciliation of the Churches; and they were detained by
delays which, whether necessary or not, were sufficient to alarm them,
and to make them impatient to bring their dispute to a termination.
After the first session of the Council at Ferrara, the public
proceedings were adjourned about six months. The Greek ecclesiastics
were allowed each three or four gold florins a month; at one time there
was an arrear of four months in the payment, at another of three, and at
the time of their agreeing to unite with the Latins, of five and a half.
Besides, even had they the means, their withdrawal from the Council was
absolutely forbidden: passports were required at the gates of Ferrara,
the Venetian Government had engaged to intercept all fugitives, and
civil punishment awaited them at Constantinople. Their condition is
vividly described by Syropulus or Sguropulus, the ecclesiarch or
preacher, who was present at the Council as one of the Patriarch's five
attendants, and whose history of its proceedings is extant. Some
extracts shall be introduced from his work; which, besides proving what
I have said about the position of the Greeks, will introduce us in
particular to the course taken in their discussions on the subject of
Purgatory. There were four points of difference between the Churches:
the use of leaven in the Eucharistic bread, the supremacy of the Pope,
the nature of Purgatory, and the double procession of the Holy Ghost.
Concerning the subject which alone here concerns us, Syropulus says,

``At our fourth meeting the bishop of
Ephesus said, 'In our last meeting, venerable Fathers, you laid before
us four heads for discussion, out of which we might take our choice …
Julian (the legate of Eugenius at Basle) said … it seems to us best,
to treat first of the purgatorial fire, that our own minds may be
cleared by the discussion. Let us then now dispute upon this subject.
The Bishop of Ephesus answered, Be it so as you have decided; but tell
us first, whence has your Church her traditions about it, and when did
she receive and profess it, and what is her exact doctrine on the
subject. These inquiries will help us forward. This was agreed to, and
we separated.

``\emph{Meanwhile our allowance of
provisions was demanded, but not given us}. \emph{Though we
made frequent demands on account of our need, it was not given until we
came into the proposed conditions. When we had come round, we received
the second monthly allowance on the} 12\emph{th} \emph{of May}.

``While we were so circumstanced,
serious news kept coming that Amurath was preparing an attack upon
Constantinople. The Venetians sent the despatches to our Emperor and the
Patriarch; afterwards came letters from the city itself, intimating the
same, and begging them to do their utmost to gain succours. On hearing
this, we were sadly afflicted, were sick of life, prayed to GOD
for help, took it to heart, and with groans and tears begged for some
escape from so great a calamity ... The Emperor had much talk with the
Cardinals on this subject, and made representations through them to the
Pope. We, indignant at their unbecoming conduct, betook ourselves to
such private friends as we might have among them. When some of us had
intreated in this way brother Ambrose, he said to them, 'Be not out of
heart, \emph{but do your utmost to bring about an union}, and then we
shall make great preparations, and will send a formidable force to
Constantinople.'

``Meanwhile some of our company said,
that if a subscription for raising forces was proposed to our
Archbishops, they would be ready according to ther power. The Emperor
catching at this, immediately went to the Patriarch, and called us all
together, and made us a speech concerning contribution, saying that he
himself had set the pattern by borrowing money to fit ot a vessel of his
own, that he felt confident the Pope would send some also, and that it
was a duty in the case of those who had the means to be liberal in the
service of their country. To this the principal Archbishops made answer,
that were they in Constantinople, they would contribute even more than
they could well afford; but, being at present in a foreign land, and not
knowing what was coming upon them, they felt it necessary to keep what
they had, even supposing some among them had any thing left; …
however, under the necessity, they would each give something.
Accordingly four of them promised 50 aspers apiece.

``The Bishop of Nicea (the celebrated
Bessarion) said, 'I have no ducats, but I have three urns, of which I
will contribute two.' The Bishop also who came next said, 'I have no
ducats, but I have two woollen cloaks, and I give one of them.' The
Emperor on hearing as far as this, gave up the attempt as vain, for he
had reckoned that the Archbishops together might have almost fitted out
one vessel …

``In the fifth meeting, Julian began
to discuss the subject of Purgatory, and said that the Roman Church,
even from the very first, had received and held this doctrine, from the
time of the Holy Apostles, receiving it from St. Peter and St. Paul, …
and then from the Doctors of the Church who succeeded them.''

To complete the imbecility of the Greek party, they
were at variance with each other, Bessarion of Nicea inclining to the
Latins, Gregory the Penitentiary taking either side as it happened, and both opposing Mark of Ephesus, the resolute defender of the Greek
doctrines. The Latins having put their argument on paper, the Greeks had
to do the same, and the Emperor commanded Mark to draw it up, who
declined the office, unless it was understood that what he should
present would be accepted. The following childish scene ensued, which is
here introduced merely to show that the Greek cause was not fairly
represented in that Council, since it was in the hands, as will be seen,
of two rival Bishops and an Emperor as umpire, and not as if to imply
that a Council must be composed of none but superior men in order to
come to a right conclusion.

``It appeared proper that some among
ourselves should stay with the Bishop of Ephesus, and that the paper
should be drawn in our presence and hearing, and with our assistance, if
it happened to be needed. Accordingly the Bishop of Nicea, the great
Ecclesiarch'' (the writer), ``Gregory the Penitentiary, the Secretary of
the Holy Consistory, met him. The Bishop of Nicea began to converse
carelessly, and to digress into a variety of subjects. The Penitentiary
followed, and rivalled him in the irrelevancy of his discourse. They
took up each other, and emulated each other in wasting time on trifles
and impertinences. I at intervals begged them to spare words and attend
to the writing, but they persisted; when good part of the day was thus
wasted, the Bishop of Ephesus said, 'At this rate I shall not be able to
write a word: leave me with the Secretary of the Consistory, and I will
draw up something. Afterwards you shall look over it, and correct any
thing that is amiss?' On this we left the room. Then the Bishop of
Ephesus began to write; but the Bishop of Nicea did the same, at the
suggestion of the Penitentiary, who praised what he drew up to the
Emperor, and wished him to send it to the Latins, as more striking in
style, and more eloquent. At his command both compositions were brought
to him, and read in the presence of select judges. Then the Emperor said
to the Bishop of Ephesus, 'Your composition is good: it has many strong
points. But it has some things too which will give advantage to the
Latins, such as the story of St. Macarius asking the skull (of an
idolater) and receiving an answer; for you can bring no unexceptionable
testimony to this, and they will at once put it aside, and some other
arguments also. Better let alone what can be easily met, and urge a
little and strong than a parade of arguments, some of which may be
easily overset, for your opponent will fix on your weak points, and if
he masters you on one or two, he will appear to the many, or rather he
will be heralded forth, as having defeated you altogether. Therefore put
out these passages.' ... Then turning to the Bishop of Nicea, he
remarked, 'You too have your own faults, you begin by saying, 'O men of
Latium;' this is unsuitable. It is more becoming to say, 'Venerable
Fathers,' or something of the same respectful
and acceptable nature; you have other mistakes too.' He ended by saying
that the proem and previous statements of the Bishop of Nicea were the
better, but the course of the argument, the proofs, and collateral
remarks, stronger in the paper of the Bishop of Ephesus; and that it
seemed advisable to take the commencement of the former, and any other
serviceable passages, and the body of the latter.'' …

The reply thus compounded by two men of discordant
sentiments was submitted to the Latins, and an answer drawn up to it in
due form. A reply followed, and the discussion became animated.

``Meanwhile in private conversations
the Latins begged the Bishop of Ephesus to propound plainly the doctrine
which our Church holds concerning souls departed hence. \emph{But he did
not state it, being hindered by the Emperor}. And in proportion as
they perceived him resisting, and not wishing to set forth our Church
doctrine on the matter, so much the more did they press him, and entreat
him, and remonstrate with him, and asked what he meant by his reserve,
saying, that every regular member of any Church was bound, when asked
what was the Church's view on any question, at once to give it, without
hesitation or ambiguity. But the Bishop had his mouth stopped by the
royal command.''

John, a Spanish Bishop, then entered into a
discussion with the Bishop of Ephesus with great dialectic skill, and
Bessarion deserted to the Latins; at length, however, the Emperor
consented to Mark's speaking out, and he put the Latins into full
possession of the Greek notions on the subject of Purgatory. The next
sentences run as follows:—

``Our allowance was expended, and
nothing more was given us in spite of our frequent demands: but, when we
yielded to their demand, and told them our Church's opinion on the
question in discussion, then they gave us three months' allowance on the
30th of June, 689 florins.'' 5. § 18.

This was all that passed on the subject of
Purgatory, before the final decree, which, as in other points, so in
this, was over-ruled by the determination of the Latins and the need of
the Emperor. But here let me instance another hardship inflicted on the
Greeks, for which I have already prepared the reader.

``We sat down in sorrow, not only
because of existing and expected perils, but for the loss of liberty,
for we were shut up as slaves. And when three months and more were
passed, and all were indignant at our dependence upon strangers,
the straits we were in, and our want of provision, … three clerics,
under the spur of necessity, found an escape … But the Patriarch
learning it, and being indignant at it, wrote at once to the Doge of
Venice, who found out the men, and sent them to him.''

After many months' discomfort from the causes that
have been enumerated, the Greeks came to an understanding with the
Latins: indeed, from the first, they had very little trust or attachment
to their view. Their doctrine is said to have been, that the souls of
imperfect Christians went to a place of darkness and sadness, where they
were for some time in affliction and deprived of the light of GOD'S countenance, in which state they were benefited by
Eucharistic offerings and by alms; to this the Latins wished to add,
that souls without stain enter at once into heavenly glory, while those
who have repented of sin, but have not had time to complete the
necessary penance, are consigned for a longer or shorter time to
purgatorial fire. This was the difference between the Churches, and they
compromised the matter thus: the Latins did not press the doctrine of \emph{fire},
and the Greeks gave up—not a word, but a truth,—they allowed,
contrary to the belief with which they had come to the Council, that
those who are not in Purgatory are immediately beatified, and enjoy the
sight of GOD.

It may be objected, and readily admitted, that the
narrative of which the above are extracts, is drawn up by a writer
unfriendly and unfair to the Latins. But it would seem to prove as much
as this, viz. what was the popular view in Greece on the subject of
these discussions and their termination, immediately upon it.

A high ecclesiastic, as Syropulus was, would hardly
have ventured to have set himself against a recent and solemn act of his
own Church, sanctioned by the Court, unless he had had a strong feeling
with him. The very fact of his opposition proves that the conduct of the
Greeks at Florence was but the act of a party at most in the Church;
while the line of the history, their sufferings and compelled decision,
is too clearly guaranteed to us as true by the known circumstances of
the case. But we need not thus painfully deduce the real dissatisfaction
of the Greek Church with the articles imposed upon its delegates at
Florence. On their return home, they had to encounter so general an
indignation and

resentment at their conduct, that they were obliged
at once to recant and confess their weakness, and throw themselves on
the mercy of their brethren. Mark of Ephesus had not signed the decree,
and became a rallying point for all who held by the popular religion;
while the successor of the Patriarch was deserted even by his
cross-bearers, and presided in an empty Cathedral.

The feeling spread north and south; the patriarchs
of Alexandria, Antioch, and Jerusalem assembled a numerous Council, and
disowned the acts of their representatives in Italy; and Isidore, the
Primate of Russia, on returning to his country, was synodically
condemned and imprisoned in a monastery.

Again, it may be objected that the great article of
difference between Greeks and Latins was the question of the Procession,
not that of Purgatory, and after all, that the real point of repulsion
between them lay in national jealousies; whereas they agreed together,
as the Council shows, or at least with the slightest difference, on the
question in which we are concerned, while the subsequent resentment of
the Greeks at home had little or no reference to it; and that their
agreement under such circumstances was only the more remarkable. It may
be replied, that the object of the foregoing account has been to show
that the Greeks at Florence were not trustworthy, that they had neither
the ease of circumstances, the learning, or the composure of mind to be
witnesses of the traditionary and universal doctrine of their Churches.
If this is proved by after circumstances, by the popular indignation as
regards one doctrine, it takes all credit from their testimony as
regards another. Moreover as regards the doctrine of Purgatory, they did
not agree with the Latins in an important point, yet that point they
gave up to them; most unfaithfully, considering them as stewards of
Gospel truth; and, had they discerned the bearings of the Latin
doctrine, which doubtless they did not, most treacherously. They
admitted, against the national belief, the beatification of souls under
specific circumstances, before the judgment, and in so doing they
admitted practically almost as much, as if they had subscribed to the
doctrine of purgatorial fire. For, as the mention of fire on the one
hand is definite, and ascertains Purgatory to be strictly a place
of punishment, which the general expressions of the Greeks did not
strictly imply, so in like manner to separate off from it all the
perfected saints, and transfer them to a better and heavenly state, does
in effect sink it, by the contrast, to a place of privation and
suffering. The presence of the souls of all saints, (to speak in general
terms, that is, not to include the Martyrs whom the early Church has
excepted) in Hades, Paradise, or Abraham's bosom, or by whatever other
name we designate the Intermediate State, is our guarantee for the
substantial blessedness of that State. We cannot spare the higher Saints
from Paradise, in that they are our pledges for its heavenly character
in the case of all believers. Thus as regards their own doctrine, the
Greeks made most important admissions to the Latins, for making which
they had no warrant, and therefore cannot be considered of authority in
witnessing a Purgatory at all, any more than in the account they gave of
it.

And with these remarks shall terminate a
discussion, which has extended far beyond the limits which were
originally proposed by the writer.

OXFORD,

\emph{The Feast of the Annunciation}.

[THIRD EDITION.]

———————————————————————

These Tracts are continued in
Numbers, and sold at the price of 2d. for each sheet, or 7s. for 50
copies.

LONDON:
PRINTED FOR J. G. F. \& J. RIVINGTON,

ST. PAUL'S CHURCH YARD AND WATERLOO PLACE.

1840.

\xchapter{Daily Prayers}

\xsection{Preparation}

\xsubsection{1. Times of Prayer}

Always. \emph{Luke}
xviii. 1.

Without ceasing. 1
\emph{Thes}. v. 17.

At all times. \emph{Eph}.
vi. 18.

Samuel among such
as call upon His name \footnote{Transferred from p. 4 of Edition of 1675.}.
\emph{Ps}. xcix. 6.

GOD
forbid that I should sin against the Lord in ceasing to pray

for you and showing you the good and the right way. 1 \emph{Sam}.
xii. 23.

We will give
ourselves continually to prayer and to the ministry

of the word. \emph{Acts} vi. 4.

He kneeled upon his knees three times a day,
and prayed and

gave thanks before his GOD,
as he did aforetime. \emph{Dan}. vi. 10.

In the evening,
and morning, and at noon day will I pray, and that

instantly, and He shall hear my voice. \emph{Ps}. lv. 18.

Seven times a day do I praise Thee. \emph{Ps}.
cxix. 164.

1. In the morning, a great while before day. \emph{Mark}
i. 35.

2. In the morning
watch. \emph{Ps}. lxiii. 6. [vid. also \emph{Ps}. cxxx. 6.]

3. The third hour
of the day. \emph{Acts} ii. 15.

4. About the sixth
hour. \emph{Acts} x. 9.

5. The hour of
prayer, the ninth. \emph{Acts} iii. 1.

6. The eventide. \emph{Gen}.
xxiv. 63. 7. By night. \emph{Ps}.
cxxxiv. 2.

At midnight. \emph{Ps}. cxix. 62.

\xsubsection{2. Places of Prayer}

In all places
where I record My Name, I will come to thee,

and I will bless thee. \emph{Exod}. xx. 24.

Let \footnote{Transferred from pp. 5, 6,
and 9, of Edition 1675.} Thine eyes be open towards this house night and

day, even toward the place of which Thou hast said, My Name

shall be there; that Thou mayest hearken unto the prayer which

Thy servant shall make towards this place. 1 \emph{Kings} viii.
29.

Thou that hearest the prayer

unto Thee shall all flesh come.

The fierceness of man shall turn to Thy praise,

and the fierceness of them shalt Thou refrain.

As for me, I will come into Thy house

even upon the multitude of Thy mercy,

and in Thy fear will I worship

toward Thy Holy Temple.

Hear the voice of my humble petitions,

when I cry unto Thee;

when I hold up my hands

towards the mercy-seat of Thy Holy Temple.

We wait for Thy loving-kindness, O GOD,

in the midst of Thy Temple.

1. Among the faithful and in the
congregation. \emph{Ps}. cxi. 1.

2. Enter into thy
closet and, when thou hast shut thy door,

pray to thy Father which is in secret. \emph{Matt}. vi. 6.

3. They went up
into an upper room. \emph{Acts} i. 13.

4. He went up upon
the housetop to pray. \emph{Acts} x. 9.

5. They went up
together into the Temple. \emph{Acts} iii. 1.

6. We kneeled down
on the shore, and prayed. \emph{Acts} xxi. 5.

7. He went forth
over the brook Cedron, where was a garden.

\emph{John} xviii. 1. 8. Let them
rejoice in their beds. \emph{Ps}. cxlix. 5.

9. He departed
into a desert place and there prayed. \emph{Mark} i. 35.

10. In every place
lifting up holy hands without wrath and doubting.

1 \emph{Tim}. ii. 8.

\xsubsection{3. Circumstances of Prayer}

1. Kneeling,
\emph{humiliation}.

He kneeled down and prayed. \emph{Luke} xxii. 41.

He went a little further, and fell on His face, and prayed.

\emph{Matt}. xxvi. 39.

My soul is brought low, even unto the dust,

my belly cleaveth unto the ground.

2. Sinking the head,
\emph{shame}.

Drooping the face. [\emph{Ezr}. ix. 6.]

3. Smiting the breast, [\emph{Luke}
xviii. 13.]
\emph{indignation}. *

4. Shuddering, [\emph{Acts}
xvi. 29.]
\emph{fear}. *

5. Groaning, [\emph{Isai}. lix.
11.]
\emph{sorrow}. *

Clasping of hands.

6. Raising of eyes and hands,
[\emph{Ps}. xxv. 15. cxliii. 6.]
\emph{vehement desire}. *

7. Blows, [\emph{Ps}. lxxiii.
14.]
\emph{revenge}. *

* 2 Cor. vii. 11.

\xchapter{Order of Matin Prayer}

\emph{Litany}.

Glory be to Thee, O LORD, glory to Thee.

Glory to Thee who givest me sleep

to recruit my weakness,

and to remit the toils

of this fretful flesh.

To this day and all days,

a perfect, holy, peaceful, healthy, sinless course,

Vouchsafe O LORD.

The Angel of peace, a faithful guide,

guardian of souls and bodies,

to encamp around me,

and ever to prompt what is salutary,

Vouchsafe O LORD.

Pardon and remission

of all sins and of all offences

Vouchsafe O LORD.

To our souls what is good and convenient,

and peace to the world

Vouchsafe O LORD.

Repentance and strictness

for the residue of our life,

and health and peace to the end,

Vouchsafe O LORD.

Whatever is true, whatever is honest,

whatever just, whatever pure,

whatever lovely, whatever of good report, 

if there be any virtue, if any praise,

such thoughts, such deeds,

Vouchsafe O LORD.

A Christian close,

without sin, without shame,

and, should it please Thee, without pain,

and a good answer

at the dreadful and fearful judgment-seat

of JESUS
CHRIST
our LORD,

Vouchsafe O LORD.

\emph{Confession}.

Essence beyond essence, Nature increate,

Framer of the world,

I set Thee, L\textsc{ord}, before my face,

and I lift up my soul unto Thee.

I worship Thee on my knees,

and humble myself under Thy mighty hand.

I stretch forth my hands unto Thee,

my soul gaspeth unto Thee as a thirsty land.

I smite on my breast

and say with the Publican,

GOD
be merciful to me a sinner,

the chief of sinners;

to the sinner above the Publican,

be merciful as to the Publican.

Father of mercies,

I beseech Thy fatherly affection,

despise me not

an unclean worm, a dead dog, a putrid corpse,

despise not Thou the work of Thine own hands,

despise not Thine own image

though branded by sin.

LORD,
if Thou wilt, Thou canst make me clean,

LORD,
only say the word, and I shall be cleansed.

And Thou, my SAVIOUR CHRIST,

CHRIST
my SAVIOUR,

SAVIOUR
of sinners, of whom I am chief,

despise me not, despise me not, O LORD,

despise not the cost of Thy blood,

who am called by Thy Name;

but look on me with those eyes

with which Thou didst look upon

Magdalen at the feast,

Peter in the hall,

the thief on the wood;—

that with the thief I may call on Thee humbly,

Remember me, LORD,
in Thy kingdom;

that with Peter I may bitterly weep and say,

O that mine eyes were a fountain of tears

that I might weep day and night;

that with Magdalen I may hear Thee say,

Thy sins be forgiven thee,

and with her to love much,

for many sins yea manifold

have been forgiven me.

And Thou, All-holy, Good, and Life-giving Spirit,

despise me not, Thy breath,

despise not Thine own holy things;

but turn Thee again, O LORD,
at the last,

and be gracious unto Thy servant.

\emph{Commendation}.

Blessed art Thou, O LORD,

Our GOD,

the GOD
of our Fathers;

who turnest the shadow of death into the morning;

and lightenest the face of the earth;

who separatest darkness from the face of the light

and banishest night and bringest back the day;

who lightenest mine eyes,

that I sleep not in death;

who deliverest me from the terror by night,

from the pestilence that walketh in darkness; 

who drivest sleep from mine eyes,

and slumber from mine eyelids;

who makest the outgoings of the morning and evening

to praise Thee;

because I laid me down and slept and rose up again,

for the LORD
sustained me;

because I waked and beheld,

and my sleep was sweet unto me.

Blot out as a thick cloud my transgressions,

and as a cloud my sins;

grant me to be a child of light, a child of the day,

to walk soberly, holily, honestly, as in the day,

vouchsafe to keep me this day without sin.

Thou who upholdest the falling and liftest the fallen,

let me not harden my heart in provocation,

or temptation or deceitfulness of any sin.

Moreover, deliver me today

from the snare of the hunter

and from the noisome pestilence;

from the arrow that flieth by day,

from the sickness that destroyeth in the noon day.

Defend this day against my evil,

against the evil of this day defend Thou me.

Let not my days be spent in vanity,

nor my years in sorrow.

One day telleth another,

and one night certifieth another.

O let me hear Thy loving-kindness betimes in the morning,

for in Thee is my trust;

show Thou me the way that I should walk in,

for I lift up my soul unto Thee.

Deliver me O LORD,
from mine enemies,

for I flee unto Thee.

Teach me to do the thing that pleaseth Thee,

for Thou art my GOD;

let Thy loving S\textsc{pirit} lead me forth

into the land of righteousness.

Quicken me, O LORD, for Thy Name's sake,

and for Thy righteousness' sake

bring my soul out of trouble;

remove from me foolish imaginations,

inspire those which are good

and pleasing in Thy sight.

Turn away mine eyes

lest they behold vanity;

let mine eyes look right on,

and let mine eyelids look straight before me.

Hedge up mine ears with thorns

lest they incline to undisciplined words.

Give me early the ear to hear,

and open mine ears to the instruction of Thy oracles.

Set a watch, O LORD,
before my mouth,

and keep the door of my lips.

Let my word be seasoned with salt,

that it may minister grace to the hearers.

Let no deed be grief unto me

nor offence of heart.

Let me do some work

for which Thou wilt remember me, LORD, for good,

and spare me according to the greatness of Thy mercy.

Into Thine hands I commend

my spirit, soul, and body,

which Thou hast created, redeemed, regenerated,

O LORD,
Thou GOD
of truth;

and together with me

all mine and all that belongs to me.

Thou hast vouchsafed them to me,

LORD,
in Thy goodness.

Guard us from all evil,

guard our souls,

I beseech Thee, O LORD.

Guard us without falling,

and place us immaculate

in the presence of Thy glory

in that day. 

Guard my going
out and my coming in

henceforth and for ever.

Prosper, I pray Thee, Thy servant this day,

and grant him mercy

in the sight of those who meet him.

O GOD,
make speed to save me,

O LORD,
make haste to help me.

O turn Thee then unto me,

and have mercy upon me;

give Thy strength unto Thy servant,

and help the son of Thine handmaid.

Show some token upon me for good,

that they who hate me may see it and be ashamed,

because Thou, LORD,
hast holpen me

and comforted me.

\xchapter{Order of Evening Prayer }

\footnote{Page 196, edit. 1675.}

\emph{Meditation}.

The day is gone,

and I give Thee thanks, O LORD.

Evening is at hand,

furnish it with brightness.

As day has its evening

so also has life;

the even of life is age,

age has overtaken me,

furnish it with brightness.

Cast me not away in the time of age;

forsake me not when my strength faileth me.

Even to my old age be Thou H\textsc{e},

and even to hoar hairs carry me;

do Thou make, do Thou bear,

do Thou carry and deliver me.

Abide with me, LORD,

for it is toward evening,

and the day is far spent

of this fretful life.

Let Thy strength be made perfect

in my weakness.

Day is fled and gone,

life too is going,

this lifeless life.

Night cometh,

and cometh death,

the deathless death.

Near as is the end of day,

so too the end of life.

We then, also remembering it,

beseech of Thee

for the close of our life, 

that Thou wouldest direct it in peace,

Christian, acceptable,

sinless, shameless,

and, if it please Thee, painless,

LORD,
O LORD,

gathering us together

under the feet of Thine Elect,

when Thou wilt, and as Thou wilt

only without shame and sins.

Remember we the days of darkness,

for they shall be many,

lest we be cast into outer darkness.

Remember we to outstrip the night

doing some good thing.

Near is judgment;—

a good and acceptable answer

at the dreadful and fearful judgment-seat

of JESUS
CHRIST

vouchsafe to us, O LORD.

By night I lift up my hands in the sanctuary,

and praise the LORD.

The LORD
hath granted His loving-kindness

in the day time;

and in the night season did I sing of Him,

and made my prayer unto the GOD of my life.

As long as I live will I magnify Thee on this manner,

and lift up my hands in Thy Name.

Let my prayer be set forth in Thy sight

as the incense,

and let the lifting up of my hands

be an evening sacrifice.

Blessed art Thou, O LORD,
our GOD,

the GOD
of our fathers,

who hast created the changes of days and nights,

who givest songs in the night,

who hast delivered us from the evil of this day,

who hast not cut off like a weaver my life,

nor from day even to night made an end of me.

\emph{Confession}.

LORD,

as we add day to day,

so sin to sin.

The just falleth seven times a day;

and I, an exceeding sinner,

seventy times seven,

a wonderful, a horrible thing, O LORD.

But I turn with groans

from my evil ways,

and I return into my heart,

and with all my heart I turn to Thee,

O GOD
of penitents and Saviour of sinners;

and evening by evening I will return

in the innermost marrow of my soul;

and my soul out of the deep

crieth unto Thee.

I have sinned, O LORD,
against Thee,

heavily against Thee;

alas, alas, woe is me! for my misery.

I repent, O me! I repent, spare me, O LORD,

I repent, O me, I repent,

help Thou my impenitence.

Be appeased, spare me, O LORD;

be appeased, have mercy on me:

I said, L\textsc{ord}, have mercy upon me,

heal my soul, for I have sinned against Thee.

Have mercy upon me, O LORD,
after Thy great goodness,

according to the multitude of Thy mercies

do away mine offences.

Remit the guilt,

heal the wound,

blot out the stains,

clear away the shame,

rescue from the tyranny,

and make me not a public example. 

O bring Thou me out of my trouble,

cleanse Thou me from secret faults,

keep back Thy servant also from presumptuous sins.

My wanderings of mind

and idle talking

lay not to my charge.

Remove the dark and muddy flood

of foul and wicked thoughts.

O LORD,

I have destroyed myself;

whatever I have done amiss, pardon mercifully.

Deal not with us after our sins,

neither reward us after our iniquities.

Look mercifully upon our infirmities;

and for the glory of Thy All-holy Name,

turn from us all those ills and miseries,

which by our sins, and by us through them,

are most righteously and worthily deserved.

\emph{Commendation}.

To my weariness, O LORD,

vouchsafe Thou rest,

to my exhaustion

renew Thou strength.

Lighten mine eyes that I sleep not in death.

Deliver me from the terror by night,

the pestilence that walketh in darkness.

Supply me with healthy sleep,

and to pass through this life without fear.

O keeper of Israel,

who neither slumberest nor sleepest,

guard me this night from all evil,

guard my soul, O LORD.

Visit me with the visitation of Thine own,

reveal to me wisdom in the visions of the night.

If not, for I am not worthy, not worthy,

at least, O loving LORD,

let sleep be to me a breathing time

as from toil, so from sin. 

Yea, O LORD,

nor let me in my dreams imagine

what may anger Thee,

what may defile me.

Let not my loins be filled with illusions,

yea, let my reins chasten me in the night season,

yet without grievous terror.

Preserve me from the black sleep of sin;

all earthly and evil thoughts

put to sleep within me.

Grant to me light sleep,

rid of all imaginations

fleshly and satanical.

LORD,
Thou knowest

how sleepless are mine unseen foes,

and how feeble my wretched flesh,

who madest me;

shelter me with the wing of Thy pity;

awaken me at the fitting time,

the time of prayer;

and give me to seek Thee early,

for Thy glory, and for Thy service.

Into Thy hands, O LORD,
I commend myself,

my spirit, soul, and body:

Thou didst make, and didst redeem them;

and together with me, all my friends

and all that belongs to me.

Thou hast vouchsafed them to me, LORD,

in Thy goodness.

Guard my lying down and my rising up,

from henceforth and for ever.

Let me remember Thee on my bed,

and search out my spirit;

let me wake up and be present with Thee;

let me lay me down in peace, and take my rest:

for it is Thou, LORD,
only

that makest me dwell in safety.

\xchapter{Course of Prayers for the Week}

\xsection{The First Day}

\emph{Introduction}.

Through the tender mercies of our GOD

the day-spring from on high hath visited us.

Glory be to Thee, O LORD
glory to Thee.

Creator of the light,

and Enlightener of the world,—

of the visible light,

the Sun's ray, a flame of fire,

day and night,

evening and morning,—

of the light invisible,

the revelation of GOD,

writings of the Law,

oracles of Prophets,

music of Psalms,

instruction of Proverbs,

experience of Histories,—

light which never sets.

GOD
is the LORD
who hath showed us light;

bind the sacrifice with cords,

yea even unto the horns of the altar.

O by Thy resurrection raise us up

unto newness of life,

supplying to us frames of repentance.

The GOD
of peace,

who did bring again from the dead

the great Shepherd of the sheep,

through the blood of the everlasting covenant,

our LORD
JESUS
CHRIST,

perfect us in every good work,

to do His will,

working in us what is acceptable before Him,

through JESUS
CHRIST,

to whom be glory for ever.

Thou who didst send down on Thy disciples

on this day

Thy Thrice-Holy SPIRIT,

withdraw not Thou the gift, O LORD, from us,

but renew it in us, day by day,

who ask Thee for it.

(1) \emph{Confession}.

Merciful and pitiful LORD,

Long-suffering and full of pity,

I have sinned, LORD,
I have sinned against Thee;

O me, wretched that I am,

I have sinned, Lord, against Thee,

much and grievously,

in attending on vanities and lies.

I conceal nothing:

I make no excuses.

I give Thee glory, O LORD,
this day,

I denounce against myself my sins;

Truly I have sinned before the LORD,

and thus and thus have I done.

I have sinned and perverted

that which was right,

and it profited me not.

And what shall I now say?

or with what shall I open my mouth?

What shall I answer, seeing I have done it?

Without plea, without defence, self-condemned, am I.

I have destroyed myself.

Unto Thee, O LORD,
belongeth righteousness,

but unto me confusion of face,

because Thou art just in all that is come upon me;

for Thou hast done right,

but I have done wickedly.

And now, LORD,
what is my hope?

Truly my hope is even in Thee,

if hope of salvation remain to me,

if Thy loving-kindness cover

the multitude of my iniquities.

O remember, what my substance is,

the work of Thine hands,

the likeness of Thy countenance,

the cost of Thy blood,

a name from Thy Name,

a sheep of Thy pasture,

a son of the covenant,

Despise not Thou the work of Thine own hands.

Hast Thou made for nought

Thine own image and likeness?

for nought, if Thou destroy it.

And what profit is there in my blood?

Mine enemies will rejoice;

May they never rejoice, O LORD!

Grant not to them my destruction.

Look upon the face of Thine Anointed,

and in the Blood of Thy covenant,

the propitiation for the sins of the whole world,

LORD,
be propitious unto me, a sinner;

even unto me, O LORD,
of sinners

chief, chiefest and greatest,

For Thy Name's sake be merciful unto my sin,

for it is great: it exceeds.

For Thy Name's sake, that Name,

beside which, none other under heaven

is given among men,

whereby we must be saved,

the SPIRIT
Himself helping our infirmities,

and making intercession for us,

with plaints unutterable, 

for the tender yearnings of the FATHER,

the bloody wounds of the SON,

the unutterable plaints of the SPIRIT,

give ear, O LORD,
have mercy, O LORD,

O LORD,
hearken and do;

defer not, for Thine own sake,

O my GOD.

For me, I forget not my sins,

they are ever before me;

I remember them in the bitterness of my soul;

I am anxious about them;

I turn away and groan,

I have indignation and revenge

and wrath against myself.

I despise and bruise my own self,

that my penitence, LORD,
O LORD,

is not deeper, is not fuller;

help Thou mine impenitence.

And more, and still more,

pierce Thou, rend, crush my heart;

and remit, forgive, pardon

what things are grief to me,

and offence of heart.

Cleanse Thou me from secret faults,

and keep Thy servant also from presumptuous sins.

Magnify Thy mercies towards the wretched sinner;

and in season, LORD,
say to me,

Be of good cheer; thy sins are forgiven thee;

My grace is sufficient for thee.

Say unto my soul, I am thy salvation.

Why art thou so heavy, O my soul?

and why art thou so disquieted within thee?

Return unto thy rest, O my soul,

for the LORD
hath rewarded thee.

O LORD,
rebuke me not in Thine indignation,

neither chasten me in Thy displeasure. 

I said, I will confess my sins unto the LORD,

and so Thou forgavest the wickedness of my sin.

LORD,
Thou knowest all my desire,

and my groaning is not hid from Thee.

Have mercy upon me, O GOD,

after Thy great goodness,

according to the multitude of Thy mercies

do away mine offences.

Thou shalt arise, and have mercy on me, O LORD,

for it is time that Thou have mercy upon me,

yea, the time is come.

If Thou, O LORD,
shouldest mark iniquities,

O LORD,
who shall stand?

Enter not into judgment with Thy servant, O LORD,

for in Thy sight shall no man living be justified.

(2) \emph{Prayer for grace}.

My hands will I lift up

unto Thy commandments which I have loved.

Open Thou mine eyes that I may see,

incline my heart that I may desire,

order my steps that I may follow,

the way of Thy commandments.

O LORD
GOD,
be Thou to me a GOD,

and beside Thee none else,

none else, nought else with Thee.

Vouchsafe to me, to worship Thee and serve Thee

in truth of spirit,

in reverence of body,

in blessing of lips

in private and in public;

to pay honour to them that have the rule over me,

by obedience and submission

to show affection to my own,

by carefulness and providence;

to overcome evil with good;

to possess my vessel in sanctification and honour;

to have my converse without covetousness,

content with what I have; 

to speak the truth in love;

to be desirous not to lust,

not to lust passionately,

not to go after lusts.

(\emph{The Hedge of the Law}.

\emph{i.e}. \emph{precautions}.)

1. To bruise the serpent's head. \emph{Gen}. iii. 15.

2. To remember my latter end. \emph{Deut}. xxvii. 29.

3. To cut off opportunities. 2 \emph{Cor}. xi. 12.

4. To be sober. 1 \emph{Pet}. v. 8.

5. Not to sit idle. \emph{Matt}. xx. 6.

6. To shun the wicked. \emph{Ps}. xxvi. 5.

7. To cleave to the good. \emph{Rom}. xii. 9.

8. To make a covenant with the eyes. \emph{Job} xxxi. 1.

9. To bring my body into subjection. 1 \emph{Cor}. ix. 27.

10. To give myself unto prayer. 1 \emph{Cor}. vii. 5.

11.To betake myself to penitence. 2 \emph{Pet}. iii. 9.

Hedge up my way with thorns,

that I find not the path

for following vanity.

Hold Thou me in with bit and bridle,

lest I fall from Thee.

O LORD
compel me to come in to Thee.

(3) \emph{Profession}.

I believe, O LORD,

in Thee, Father, Word, Spirit, One GOD;

that by Thy fatherly love and power

all things were created;—

that by Thy goodness and love to man

all things have been begun anew

in Thy Word,—

Who for us men and for our salvation,

was made flesh,

was conceived and born,

suffered and was crucified,

died and was buried,

descended and rose again,

ascended and sat down,

will return and will repay;— 

that by the shining-forth and working

of Thy Holy Spirit,

hath been called out of the whole world

a peculiar people into a polity,

in belief of the truth

and sanctity of living:—

that in it we are partakers

of the communion of saints

and forgiveness of sins

in this world,—

that in it we are waiting

for resurrection of the flesh

and life everlasting

in the world to come.—

This most holy faith

which was once delivered to the saints

I believe, O LORD;

help Thou mine unbelief,

and vouchsafe to me

to love the FATHER
for His fatherly love,

to reverence the ALMIGHTY
for His power,

as a faithful CREATOR,
to commit my soul to Him

in well doing;

vouchsafe to me to partake

from JESUS
of salvation,

from CHRIST
of anointing,

from the ONLY-BEGOTTEN
of adoption;

to worship the LORD

for His conception in faith,

for His birth in humility,

for His sufferings in patience and hatred of sin;

for His cross to crucify beginnings,

for His death to mortify the flesh,

for His burial to bury evil thoughts in good works,

for His descent to meditate upon hell,

for His resurrection upon newness of life, 

for His ascension, to mind things above,

for His sitting on high, to mind the good things on His right,

for His return, to fear His second appearance,

for judgment, to judge myself ere I be judged.

From the SPIRIT

vouchsafe me the breath of salutary grace.

In the Holy Catholic Church

to have my own calling, and holiness, and portion.

and a fellowship \footnote{This passage serves to illustrate St. Leo's
language as it

occurs in Tract 75. Beatum Apostolum Petrum, cujus

suffragantibus meritis, quæ noscimus, impetrare possimus

per Dominum nostrum Jesum Christum.}

of her sacred rites, and prayers,

fastings and groans,

vigils, tears, and sufferings,

for assurance of remission of sins,

for hope of resurrection and translation

to eternal life.

(4) \emph{Intercession}.

O Hope of all the ends of the earth,

and of them that remain in the broad sea;

O Thou on whom our fathers hoped,

and Thou didst deliver them;

on whom they waited,

and were not confounded;

O my Hope from my youth,

from my mother's breasts;

on whom I have been cast from the womb,

be Thou my hope

now and evermore,

and my portion in the land of the living:

In Thy nature,

in Thy names, in Thy types,

in word and in deed,

My Hope,

let me not be disappointed of my hope.

O the Hope of all the ends of the earth,

remember Thy whole creation for good, 

visit the world in Thy compassion;

O guardian of men,

O loving LORD,

remember all our race.

Thou who hast shut up all in unbelief,

on all have pity, O LORD.

O Thou who didst die and rise again,

to be LORD
both of the dead and living,

live we or die we,

Thou art our LORD;

LORD,
have pity on living and dead.

O helper of the helpless,

seasonable aid in affliction,

remember all who are in necessity,

and need Thy succour.

O GOD
of grace and truth,

establish all who stand in truth and grace,

restore all who are sick with heresies and sins.

O wholesome Defence of Thine anointed,

remember Thy Congregation

which Thou hast purchased and redeemed of old.

O grant to all believers

one heart and one soul.

Thou that walkest amid the golden candlesticks,

remove not our candlestick

out of its place.

Amend what are wanting,

establish what remain,

which Thou art ready to cast away,

which are ready to die.

O LORD
of the harvest

send forth labourers, made sufficient by Thee,

into Thy harvest.

O portion of those

who wait in Thy temple,

grant to our clergy. 

rightly
to divide the word of truth,

rightly to walk in it;

grant to Thy Christian people

to obey and submit to them.

O King of nations, unto the ends of the earth;

strengthen all the states

of the inhabited world,

as being Thy ordinance,

though a creation of man.

Scatter the nations that delight in war,

make wars to cease in all the earth.

O expectation of the isles and their hope,

L\textsc{ord}, save this island,

and all the country in which we sojourn,

from all affliction, peril, and need.

LORD of lords, Ruler of rulers,

remember all rulers

to whom Thou hast given rule in the earth,

and O remember specially

our divinely-guarded king,

and work with him more and more,

and prosper his way in all things.

Speak good things unto his heart,

for Thy Church, and all Thy people,

grant to him profound and perpetual peace,

that in his tranquility

we may lead a quiet and peaceable life

in all godliness and honesty.

O Thou by whom are ordained the powers that be,

grant to those who are chief in court,

to be chief in virtue and Thy fear;

grant to the Parliament Thy holy wisdom;

to our great men, to do nothing against

but for the truth;

to the courts of law, Thy judgments,

to judge in all things concerning all

without preference, without partiality. 

O G\textsc{od} of armies,

give a prosperous course and strength

to all the Christian army,

against the enemies of our most holy faith.

Grant to our population

to be subject unto the higher powers,

not only for wrath, but also for conscience-sake.

Grant to farmers and graziers good seasons;

to the fleet and fishers fair weather;

to tradesmen, not to overreach one another;

to mechanics, to pursue their business lawfully,

down to the meanest workman,

down to the poor.

O G\textsc{od}, not of us only but of our seed,

bless our children among us,

to advance in wisdom as in stature,

and in favour with Thee and with men.

Thou who wouldest have us provide for our own,

and hatest the unnatural,

remember, L\textsc{ord}, my relations according to
the flesh;

grant me to speak peace concerning them,

and to seek their good.

Thou who willest us to make return

to our benefactors,

remember, L\textsc{ord}, for good,

all from whom I have received good;

keep them alive that they may be blessed upon earth,

and deliver them not

into the will of their enemies.

Thou who hast noted

the man who neglects his own, as worse than an infidel,

remember in Thy good pleasure

all those in my household.

Peace be to my house,

the S\textsc{on} of peace upon all in it.

Thou who wouldest that our righteousness exceed

the righteousness of sinners, 

grant me, LORD, to love those who love me,

my own friend, and my father's friend, and my friend's children,

never to forsake.

Thou who wouldest that we overcome

evil with good,

and pray for those who persecute us,

have pity on mine enemies, LORD,

as on myself;

and lead them together with me to Thy heavenly kingdom.

Thou who grantest the prayers of Thy servants

one for another,

remember, LORD,
for good,

and pity all those

who remember me in their prayers,

or whom I have promised to remember in mine.

Thou who acceptest diligence in every good work,

remember, LORD,
as if they prayed to Thee,

those who for any good reason

give not time to prayer.

Arise, and have mercy

on those who are in the last necessity,

for it is time that Thou hast mercy upon them,

yea the time is come.

Have mercy on them, O LORD,

as on me also, when in extremities.

Remember, LORD,

infants, children, the grown, the young, the middle aged, the
old,

hungry, thirsty, naked, sick,

prisoners, foreigners, friendless, unburied,

all in extreme age and weakness,

possessed with devils, and tempted to suicide,

troubled by unclean spirits,

the hopeless, the sick in soul or body, the weak-hearted,

all in prison and chains, all under sentence of death;

orphans, widows, foreigners, travellers, voyagers,

women with child, women who give suck,

all in bitter servitude, or mines, or galleys, or in loneliness.

Thou, LORD,
shalt save both man and beast,

how excellent is thy mercy, O GOD!

And the children of men shall put their trust

under the shadow of Thy wings.

The LORD
bless us, and keep us,

and show the light of His countenance upon us,

And be merciful unto us,

The LORD
lift up His countenance upon us,

And give us peace!

I commend to Thee, O LORD,

my soul, and my body,

my mind, and my thoughts,

my prayers, and my vows,

my senses, and my limbs,

my words, and my works \footnote{Page 172, edit. 1675.},

my life, and my death;

my brothers, and my sisters,

and their children;

my friends, my benefactors, my well wishers,

those who have a claim on me;

my kindred, and my neighbours,

my country, and all Christendom.

I commend to Thee, LORD,

my impulses, and my startings,

my intentions, and my attempts,

my going out, and my coming in,

my sitting down, and my rising up.

(5) \emph{Praise}.

Up with our hearts;

we lift them to the LORD.

O how very meet, and right, and fitting, and due,

in all, and for all,

at all times, places, manners,

in every season, every spot,

everywhere, always, altogether, 

to remember Thee, to worship Thee,

to confess to Thee, to praise Thee,

to bless Thee, to hymn Thee,

to give thanks to Thee,

Maker, nourisher, guardian, governor,

preserver, worker, perfecter of all,

LORD
and Father,

King and God,

fountain of life and immortality,

treasure of everlasting goods.

Whom the heavens hymn,

and the heaven of heavens,

the Angels and all the heavenly powers,

one to other crying continually,—

and we the while, weak and unworthy,

under their feet,

HOLY,
HOLY,
HOLY

LORD
the GOD
of Hosts;

full is the whole heaven,

and the whole earth,

of the majesty of Thy glory.

Blessed be the glory of the LORD

out of His place,

for His Godhead, His mysteriousness,

His height, His sovereignty, His almightiness,

His eternity. His providence.

The LORD
is my strength, my stony rock, and my defence,

my deliverer, my succour, my buckler,

the horn also of my salvation and my refuge.

\xchapter{The Second Day}

\emph{Introduction}.

My voice shalt Thou hear betimes, O LORD,

early in the morning will I direct my prayer unto Thee,

and will look up.

Blessed art Thou, O LORD,

who didst create the firmament of heaven,

the heavens and the heaven of heavens,

the heavenly powers,

Angels, Archangels,

Cherubim, Seraphim,

waters above the heavens,

mists and exhalations,

for showers, dew, hail, snow as wool,

hoar frost as ashes, ice as morsels,

clouds from the ends of the earth,

lightnings, thunders, winds out of Thy treasures, storms;

waters beneath the heavens,

for drinking and for bathing.

\emph{Confession}.

I will confess my sins,

and the sins of my fathers,

for I have transgressed and neglected Thee, O LORD,

and walked perversely before Thee.

Set not, O LORD,
set not my misdeeds before Thee,

nor my life in the light of Thy countenance,

But pardon the iniquity of Thy servant,

according to Thy great mercy;

as Thou hast been merciful to him from a child,

even so now.

I have sinned, what shall I do unto Thee,

O Thou preserver of men? 

Why hast Thou set me as a mark against Thee,

so that I am a burden to myself?

O pardon my transgression,

and take away mine iniquity.

Deliver me from going down to the pit,

for Thou hast found a ransom.

Have mercy on me, SON
of DAVID,

LORD,
help me.

Yea, LORD,
even the dogs eat of the crumbs

which fall from their masters' table.

Have patience with me, LORD,

yet I have not wherewith to pay,

I confess to Thee;

forgive me the whole debt, I beseech Thee.

How long wilt Thou forget me, O LORD? for ever?

How long wilt Thou hide Thy face from me?

How long shall I seek counsel in my soul,

and he vexed in my heart day and night?

How long shall mine enemies triumph over me?

Consider and hear me, O LORD
my GOD,

lighten mine eyes that I sleep not in death,

lest mine enemy say I have prevailed against him,

for if I be cast down, they that trouble me will rejoice at it;

but my trust is in thy mercy.

(2) \emph{Prayer for grace}.

Remove from me

(\emph{The Ten Com- mandments})
1. all iniquity and profaneness,
superstition,

and hypocrisy.

2. worship of idols, of individuals.

3. rash oath and curse.

4. neglect or indecency of worship.

5. haughtiness and recklessness.

6. strife and wrath.

7. passion and corruption.

8. indolence and fraud.

9. lying and injuriousness. 

10.every evil notion, every impure thought,

every base
desire, every unseemly

thought.

Grant to me,

1. to be religious and pious.

2. to worship and serve.

3. to bless and swear truly.

4. to confess meetly in the congregation.

5. affection and obedience.

6. patience and good temper.

7. purity and soberness.

8. contentedness and goodness.

9. truth and incorruptness.

10. good thoughts, perseverance to the end.

(3) \emph{Profession}.

I believe in GOD,

1. the FATHER,
Almighty, Maker of heaven and earth.

2. And in JESUS
CHRIST,
His Only-begotten Son, our LORD.

(1.) conceived of the HOLY
GHOST,

(2.) born of Mary, ever-virgin,

(3.) suffered under Pontius Pilate,

(4.) crucified,

(5.) dead,

(6.) buried.—

(1.) descended into hell,

(2.) risen from the dead,

(3.) ascended into heaven,

(4.) set down on the right hand.

(5.) to return thence,

(6.) to judge both quick and dead.

3. And in the HOLY GHOST,

(1.) The Holy Church,

(2.) Catholic,

(3.) communion of saints,

(4.) remission of sins,

(5.) resurrection of flesh,

(6.) life everlasting.

And now, LORD, what is my hope?

Truly my hope is even in Thee;

in Thee, O LORD,
have I trusted,

let me never be confounded.

(4) \emph{Intercession}.

Let us pray GOD,

for the whole creation;

for the supply of seasons,

healthy, fruitful, peaceful;

for the whole race of mankind;

for those who are not Christians;

for the conversion of Atheists, the ungodly;

Gentiles, Turks, and Jews;

for all Christians;

for restoration of all

who languish in faults and sins;

for confirmation of all

who have been granted truth and grace;

for succour and comfort of all

who are dispirited, infirm, distressed, unsettled,

men and women;

for thankfulness and sobriety in all

who are hearty, healthy, prosperous, quiet,

men and women;

For the Catholic Church,

its establishment and increase;

for the Eastern,

its deliverance and union;

for the Western,

its adjustment and peace;

for the British,

the supply of what is wanting in it,

the strengthening of what remains in it \footnote{An allusion apparently to the Church in Sardis.—Rev.
iii. 2.};

for the episcopate, presbytery, christian
people; 

for the states of the inhabited world;

for christian states,

far off, near at hand;

for our own;

for all in rule;

for our divinely-guarded king,

the queen and the prince;

for those who have place in the court;

for parliament and judicature,

army and police,

commons and their leaders,

farmers, graziers, fishers, merchants,

traders, and mechanics,

down to mean workmen, and the poor;

for the rising generation;

for the good nurture of all the royal family,

of the young ones of the nobility;

for all in universities, in law colleges,

in schools in town or country,

in apprenticeships;

for those who have a claim
on me from relationship,

for brothers and sisters,

that GOD'S
blessing may be on them,

and on their children;

or from benefits conferred,

that Thy recompence may be on all

who have benefitted me,

who have ministered to me in carnal things;

or from trust placed in me,

for all whom I have educated,

all whom I have ordained:

for my college, my parish,

Southwell, St. Paul's, Westminster,

Dioceses of Chichester, Ely, and my present,

clergy, people, helps, governments,

the deanery in the chapel royal,

the almonry, 

the colleges committed to me \footnote{As
Visitor.};

or from natural kindness,

for all who love me,

though I know them not;

or from Christian love;

for those who hate me without cause,

some too, even on account of truth and righteousness;

or from neighbourhood,

for all who dwell near me

peaceably and harmlessly;

or from promise,

for all whom I have promised to remember

in my prayers;

or from mutual offices,

for all who remember me in their prayers,

and ask of me the same;

or from stress of engagements,

for all who on sufficient reasons fail to call upon Thee;

for all who have no intercessor

in their own behalf;

for all who at present are in agony

of extreme necessity or deep affliction;

for all who are attempting any good work

which will bring glory to the Name of GOD

or some great good to the Church;

for all who act nobly

either towards things sacred or towards the poor;

for all who have ever been offended by me

either in word or in deed.

GOD
have mercy on me and bless me;

GOD
show the light of His countenance upon me

and pity me.

GOD
bless me, even our GOD,

G\textsc{od} bless me and receive my prayer; 

O direct my life towards Thy commandments,

hallow my soul,

purify my body,

correct my thoughts,

cleanse my desires,

soul and body, mind and spirit,

heart and reins.

Renew me thoroughly, O GOD,

for, if Thou wilt, Thou canst.

(5) \emph{Praise}.

The LORD,
the LORD
GOD,

merciful and pitiful,

long-suffering and full of pity, and true,

keeping pity for thousands,

taking away iniquities and unrighteousnesses and sins;

not clearing the guilty one,

bringing sins of fathers upon children.

I will bless the LORD
at all times,

His praise shall ever be in my mouth.

Glory to GOD
in the highest,

and on earth peace,

goodwill towards men.

The Angels,

Archangels,

Powers,

Thrones,

Dominions,

Principalities,

Authorities,

Cherubim,

Seraphim,

guardianship;

glory;

marvels:

judgment;

beneficence;

government;

against devils;

knowledge;

love.

\xchapter{The Third Day}

\emph{Introduction}.

O GOD,
Thou art my GOD,

early will I seek Thee.

Blessed art Thou, O LORD,

who gatheredst the water into the sea,

and broughtest to sight the earth,

and madest to sprout

herb and fruit tree.

There are the depths and the sea as on an heap,

lakes, rivers, springs;

earth, continent, and isles,

mountains, hills, and valleys;

glebe, meadows, glades,

green pasture, corn, and hay;

herbs and flowers

for food, enjoyment, medicine;

fruit trees bearing

wine, oil and spices,

and trees for wood;

and things beneath the earth,

stones, metals, minerals, coal,

blood and fire, and vapour of smoke.

(1) \emph{Confession}.

Who can understand his errors?

Cleanse Thou me from secret faults.

Keep back Thy servant also from presumptuous sins,

lest they have the dominion over me.

For Thy Name's sake,

be merciful unto my sin,

for it is great. 

My iniquities have taken such hold upon

me that I am not able to look up,

yea, they are more in number than the hairs of my head,

and my heart hath failed me.

Be pleased, O LORD,
to deliver me,

Make haste, O LORD,
to help me.

Magnify Thy mercies upon me,

O Thou who savest them that trust in Thee.

I said, LORD,
have mercy upon me,

heal my soul, for I have sinned against Thee;

I have sinned but I am confounded,

and I turn from my evil ways,

and I turn unto mine own heart,

and with my whole heart I turn unto Thee;

and I seek Thy face,

and I beseech Thee, saying,

I have sinned, I have committed iniquity,

I have done unjustly.

I know, O LORD,
the plague of my heart,

and lo, I turn to Thee with all my heart,

and with all my strength.

And Thou, O LORD,
now from Thy dwelling-place,

and from the glorious throne of Thy kingdom in heaven

O hear the prayer

and the supplication of Thy servant;

and be propitious towards Thy servant

and heal his soul.

O GOD,
be merciful to me a sinner,

be merciful to me the chief of sinners.

Father, I have sinned against heaven, and against Thee,

and am no more worthy to be called Thy son,

make me one of Thy hired servants;

Make me one, or even the last,

or the least among all.

What profit is there in my blood,

when I go down to the pit?

shall the dust give thanks unto Thee?

or shall it declare Thy truth? 

Hear, O LORD, and have mercy upon me;

LORD,
be Thou my helper;

Turn my heaviness into joy,

my dreamings into earnestness,

my falls into clearings of myself,

my guilt, my offence into indignation,

my sin into fear,

my transgression into vehement desire,

my unrighteousness into strictness,

my pollution into revenge.

(2) \emph{Prayer for grace}.

Hosanna in the highest \footnote{Vide p. 186, edit. 1675.}.

Remember, me, O LORD,

with the favour that Thou bearest unto Thy people,

O visit me with Thy salvation;

that I may see the felicity of Thy chosen,

and rejoice in the gladness of Thy people,

and give thanks with Thine inheritance.

There is glory which shall be revealed;

for when the Judge cometh

some shall see Thy face cheerful,

and shall be placed on the right,

and shall hear those most welcome words,

``Come, ye blessed.''

They shall be caught up in clouds

to meet the LORD;

they shall enter into gladness,

they shall enjoy the sight of Him.

they shall be ever with Him.

These alone, only these are blessed

among the sons of men.

O to me the meanest grant the meanest place

there under their feet;

under the feet of Thine elect,

the meanest among them. 

And that this may be,

let me find grace in Thy sight

to have grace,

so as to serve Thee acceptably

with reverence and godly fear. (\emph{Heb}.
xii. 28.)

Let me find that second grace,

not to receive in vain (2 \emph{Cor}. vi. 1.)

the first grace,

not to come short of it; (\emph{Heb}. xii. 15.)

yea, not to neglect it, (1 \emph{Tim}. iv. 14.)

so as to fall from it, (\emph{Gal}. v. 4.)

but to stir it up, (2 \emph{Tim}. i. 6.)

so as to increase in it, (2 \emph{Pet}. iii. 18.)

yea, to abide in it

till the end of my life.

And O, perfect for me what is lacking

of Thy gifts,

of faith,

of hope,

of love,
help Thou mine unbelief,

establish my trembling hope,

kindle its smoking flax.

Shed abroad Thy love in my heart,

so that I may love Thee,

my friend in Thee, my enemy for Thee.

O Thou who givest grace to the humble-minded,

also give me grace to be humble-minded.

O Thou who never failest those who fear Thee,

my Fear and my Hope,

let me fear one thing only,

the fearing ought more than Thee.

As I would that men should do to me

so may I do to them;

not to have thoughts beyond what I should think,

but to have thoughts unto sobriety.

Shine on those who sit in darkness,

and the shadow of death;

guide our feet into the way of peace,

that we may have the same thoughts

one with another, 

rightly to divide, rightly to walk,

to edify,

with one accord, with one mouth,

to glorify GOD;

and if ought otherwise,

to walk in the same rule

as far as we have attained;

to maintain order,

decency and stedfastness.

(3) \emph{Profession}.

Godhead, paternal love, power,

providence:

salvation, anointing, adoption,

lordship;

conception, birth, passion,

cross, death, burial,

descent, resurrection, ascent,

sitting, return, judgment;

Breath and Holiness,

calling from the Universal,

hallowing in the Universal,

communion of saints, and of saintly things \footnote{Vide p. 172, edit. 1675.},

resurrection,

life eternal.

(4) \emph{Intercession}.

Hosanna on the earth \footnote{Vide p. 172, edit. 1675.}.

Remember, O LORD,

to crown the year with Thy goodness;

for the eyes of all look towards Thee,

and Thou givest their food in due season.

Thou openest Thine hand,

and fillest all things living with plenteousness,

And on us, O LORD,
vouchsafe

the blessings of heaven and the dew above, 

blessings of fountains and the deep beneath,

courses of sun, conjunctions of moons,

summits of eastern mountains, of the everlasting hills,

fulness of the earth and of produce thereof,

good seasons, wholesome weather,

full crops, plenteous fruits,

health of body, peaceful times,

mild government, kind laws,

wise councils, equal judgments,

loyal obedience, vigorous justice,

fertility in resources, fruitfulness in begetting,

ease in bearing, happiness in offspring,

careful nurture, sound training,

That our sons may grow up as the young plants,

our daughters as the polished corners of the temple,

that our garners may be full and plenteous

with all manner of store,

that our sheep may bring forth thousands

and ten thousands in our streets:

that there be no decay,

no leading into captivity

and no complaining in our streets.

(5) \emph{Praise}.

\footnote{Vide p. 172, edit. 1675.}
Thou, O LORD, art praised in Sion,

and unto Thee shall the vow be performed in Jerusalem.

Thou art worthy, O LORD
our God,

the HOLY
ONE

to receive glory, and honour, and power.

Thou that hearest the prayer,

unto Thee shall all flesh come,

my flesh shall come.

My misdeeds prevail against me,

O be Thou merciful unto our sins;

that I may come and give thanks

with all Thy works,

and bless Thee with Thy holy ones. 

O LORD,
open Thou my lips,

and my mouth shall show forth Thy praise.

My soul doth praise the LORD,

for the goodness He hath done

to the whole creation,

and to the whole race of man;

for Thy mercies towards myself,

soul, body, and estate,

gifts of grace, nature, and fortune;

for all benefits received,

for all successes, now or heretofore,

for any good thing done;

for health, credit, competency,

safety, gentle estate, quiet.

Thou hast not cut off as a weaver my life,

nor from day even to night made an end of me.

He hath vouchsafed me life and breath

until this hour,

from childhood, youth, and hitherto

even unto age.

He holdeth our soul in life

and suffereth not our feet to slip;

rescuing me from perils, sicknesses,

poverty, bondage,

public shame, evil chances;

keeping me from perishing in my sins,

fully waiting my conversion,

leaving in me return into my heart,

remembrance of my latter end,

shame, horror, grief,

for my past sins;

fuller and larger, larger and fuller,

more and still more, O my LORD,

storing me with good hope

of their remission,

through repentance and its works,

in the power of the thrice-holy Keys,

and the mysteries in Thy Church. 

Wherefore day by day

for these Thy benefits towards me,

which I remember,—

wherefore also for others very many

which I have let slip

from their number, from my forgetfulness—

for those which I wished, knew and asked,

and those I asked not, knew not, wished not,—

I confess and give thanks to Thee,

I bless and praise Thee, as is fit, and every day

I pray with my whole soul,

and with my whole mind I pray.

Glory be to Thee, O LORD,
glory to Thee;

glory to Thee, and glory to Thine All-holy Name,

for all Thy Divine perfections in them;

for Thine incomprehensible and unimaginable goodness,

and thy pity towards sinners

and unworthy men,

and towards me of all sinners

far the most unworthy.

Yea, O LORD,

through this, and through the rest,

glory to Thee, and praise, and blessing, and thanksgiving,

with the voices and concert of voices

of Angels and of men,

of all Thy saints in heaven,

and all Thy creatures in heaven or earth,

and of me, beneath their feet,

unworthy and wretched sinner,

Thy abject creature,

now, in this day and hour,

and every day till my last breath,

and till the end of the world,

and for ages upon ages.

\xchapter{The Fourth Day}

\emph{Introduction}.

I have thought upon Thee, O LORD,

when I was waking,

for Thou hast been my helper.

Blessed art Thou, O LORD,

who madest the two Lights, Sun and Moon,

greater and lesser,

and the stars

for light, for signs, for seasons,

spring, summer, autumn, winter,

days, weeks, months, years,

to rule over day night.

(1) \emph{Confession}.

Behold, Thou art angry, for we have sinned.

We are all as an unclean thing

and all our righteousnesses

as filthy rags.

We all do fade as a leaf,

and our iniquities, like the wind,

have taken us away.

But now, O LORD,
Thou art our Father,

we are clay, Thy handiwork all.

Be not wroth very sore,

nor remember iniquity for ever,

behold, see, we beseech Thee,

we are all Thy people.

O LORD,
though our iniquities testify against us,

do Thou it for Thy Name's sake;

for our backslidings are many,

we have sinned against Thee. 

Yet thou, O LORD,
art in the midst of us,

and we are called by Thy Name,

leave us not.

O Hope of Israel,

The SAVIOUR
thereof in time of trouble,

why shouldest Thou be as a stranger in the land,

and as a wayfaring man that turneth aside

to tarry for a night?

why shouldest Thou be as a man astonished,

as a mighty man that cannot save?

Be merciful to our unrighteousnesses,

and our iniquities remember no more.

L\textsc{ord}, I am carnal,

sold under sin;

there dwelleth in me, that is, in my flesh,

no good thing;

for the good that I would, I do not,

but the evil which I would not, that I do.

I consent unto the law that it is good,

I delight in it after the inner man;

But I see another law in my members,

warring against the law of my mind,

and enslaving me to the law of sin.

Wretched man that I am,

who shall deliver me from the body of this death?

I thank GOD
through JESUS
CHRIST,

that where sin abounded,

grace hath much more abounded.

O LORD,
Thy goodness leadeth me to repentance:

O give me sometime repentance

to recover me from the snare of the devil,

who am taken captive by him

at his will.

Sufficient for me the past time of my life

to have done the will of lusts,

walking in lasciviousness, revelling, drunkenness, 

and in other excess of profligacy.

O LAMB
without blemish and without spot,

who hast redeemed me with Thy precious Blood,

in that very Blood pity me and save me;

in that Blood,

and in that very Name,

besides which is none other under heaven

given among men,

by which we must be saved.

O GOD,
Thou knowest my foolishness,

and my sins are not hid from Thee.

LORD,
Thou knowest all my desire,

and my groaning is not hid from Thee.

Let not them that trust in Thee,

O LORD
GOD
of hosts,

be ashamed for my cause;

let not those that seek Thee be confounded through me,

O LORD
GOD
of Israel.

Take me out of the mire that I sink not;

O let me be delivered from them that hate me

and out of the deep waters;

Let not the water flood drown me,

neither let the deep swallow me up,

and let not the pit shut her mouth upon me.

(2) \emph{Prayer for grace}:

(\emph{against deadly sins}.)

(\emph{covetousness}.)

(\emph{sloth}.)

[Defend me from]

Pride

envy 

wrath 

gluttony 

lechery 

the cares of life

lukewarm indifference

Amorite.

Hittite.

Perizzite.

Girgashite.

Hivite.

Canaanite.

Jebusite.

[Give me]

Humility, pitifulness, patience,

sobriety, purity, contentment, ready zeal. 

One thing have I desired of the LORD, which I will require \footnote{Vide p. 194, edit. 1675.},

that I may dwell in the house of the LORD

all the days of my life,

to behold the fair beauty of the LORD,

and to visit His temple.

Two things have I required of Thee, O LORD,

deny Thou me not before I die;

remove far from me vanity and lies;

give me neither poverty nor riches,

feed me with food convenient for me;

lest I be full and deny Thee

and say, who is the LORD?

or lest I be poor and steal,

and take the Name of my GOD
in vain.

Let me learn to abound,

let me learn to suffer need,

in whatsoever state I am,

therewith to be content.

For nothing earthly, temporal, mortal,

to long nor to wait.

Grant me a happy life

in piety, gravity, purity,

in all things good and fair,

in cheerfulness, in health, in credit,

in competency, in safety, in gentle estate, in quiet;

a happy death,

a deathless happiness.

(3) \emph{Profession}.

I believe

in the FATHER,
benevolent affection;

in the Almighty, saving power;

in the Creator, providence

for guarding, ruling, perfecting the universe. 

In JESUS,
salvation,

in CHRIST,
anointing;

in the Only-begotten SON,
sonship,

in the LORD,
a master's treatment,

in His conception and birth

the cleansing of our unclean conception and birth;

in His sufferings, which we owed, that we might not pay;

in His cross the curse of the law removed;

in His death the sting of death;

in His burial eternal destruction in the tomb;

in His descent, whither we ought, that we might not go;

in His resurrection, as the first fruits of them that sleep;

in His ascent, to prepare a place for us;

in His sitting, to appear and intercede;

in His return, to take unto Him His own;

in His judgment, to render to each according to his works.

In the HOLY
GHOST,
power from on high,

transforming unto sanctity

from without and invisibly,

yet inwardly and evidently.

In the Church, a body mystical

of the called out of the whole world,

unto intercourse in faith and holiness.

In the communion of Saints, members of this body,

a mutual participation in holy things \footnote{Vide above, p. 24, 42.},

for confidence of remission of sins,

for hope of resurrection, of translation,

to life everlasting.

(4) \emph{Intercession}.

And I have hoped in Thy mercy

from everlasting to everlasting.

How excellent is Thy mercy, O LORD;

If I have hope, it is in Thy mercy,

O let me not be disappointed of my hope.

Moreover we beseech Thee,

remember all, LORD,
for good; 

have pity upon all, O Sovereign LORD,

be reconciled with us all.

Give peace to the multitudes of Thy people;

scatter offences;

abolish wars;

stop the uprisings of heresies.

Thy peace and love

vouchsafe to us, O GOD
our SAVIOUR,

the Hope of all the ends of the earth.

Remember to crown the year

with Thy goodness;

for the eyes of all wait upon Thee,

and thou givest them their meat in due season.

Thou openest Thy hand,

and fillest all things living with plenteousness.

Remember Thy Holy Church,

from one end of the earth to the other;

and give her peace,

whom Thou hast redeemed with Thy precious blood;

and establish her

unto the end of the world.

Remember those who bear fruit, and act nobly,

in Thy holy Churches,

and who remember the poor and needy;

recompense to them

Thy rich and heavenly gifts;

vouchsafe to them,

for things earthly, heavenly,

for corruptible, incorruptible,

for temporal, eternal.

Remember those who are in virginity,

and purity and ascetic life;

also those who live in honourable marriage,

in Thy reverence and fear.

Remember every Christian soul

in affliction, distress, and trial,

and in need of Thy pity and succour; 

also
our brethren in captivity, prison, chains,

and bitter bondage;

supplying return to the wandering,

health to the sick,

deliverance to the captives.

Remember religious and faithful kings,

whom Thou hast given to rule upon the earth;

and especially remember, L\textsc{ord},

our divinely-guarded king;

strengthen his kingdom,

subdue to him all adversaries,

speak good thing to his heart,

for Thy Church, and all Thy people.

Vouchsafe to him deep and undisturbed peace,

that in his serenity

we may lead a quiet and peaceable life

with all godliness and honesty.

Remember, L\textsc{ord}, all power

and authority,

our brethren in the court,

those who are chief in council and judgment,

and all by land and sea

waging Thy wars for us.

Moreover, L\textsc{ord}, remember graciously

our holy Fathers,

the honourable Presbytery, and all the clergy,

rightly dividing the Word of Truth,

and rightly walking in it.

Remember, L\textsc{ord}, our brethren around us,

and praying with us in this holy hour,

for their zeal and earnestness-sake.

Remember also those who on fair reasons are away,

and pity them and us

in the multitude of Thy pity.

Fill our garners with all manner of store,

preserve our marriages in peace and concord,

nourish our infants, 

lead forward our youth,

sustain our aged,

comfort the weak-hearted,

gather together the scattered,

restore the wanderers,

and knit them to Thy Holy Catholic Apostolic Church.

Set free the troubled

with unclean spirits,

voyage with the voyagers,

travel with the travellers,

stand forth for the widow,

shield the orphan,

rescue the captive,

heal the sick.

Those who are on trial, in mines, in exile, in galleys,

in whatever affliction, necessity, and emergence,

remember, O G\textsc{od};

and all who need Thy great mercy;

and those who love us,

and those who hate;

and those who have desired us unworthy

to make mention of them in our prayers;

and all Thy people remember, O L\textsc{ord}, our G\textsc{od},

and upon all pour out Thy rich pity,

to all performing their requests for salvation;

and those of whom we have not made mention,

through ignorance, forgetfulness, or number of names,

do Thou Thyself remember, O G\textsc{od},

who knowest the stature and appellation of each,

who knowest every one from his mother's womb.

For thou art, O L\textsc{ord}, the Succour of the
succourless,

the Hope of the hopeless,

the Saviour of the tempest-tost,

the Harbour of the voyager,

the Physician of the sick,

do Thou Thyself become all things to all men. 

O Thou who knowest each man and his petition,

each house, and its need,

deliver, O LORD,
this city,

and all the country in which we sojourn,

from plague, famine, earthquake, flood,

fire, sword, hostile invasion,

and civil war.

End the schisms of the Churches,

quench the haughty cries of the nations,

and receive us all into Thy kingdom,

acknowledging us as sons of light;

and Thy peace and love

vouchsafe to us, O LORD,
our GOD.

Remember O LORD,
our GOD,

all spirits and all flesh

which we have remembered, and which we have not.

And the close of our life,

LORD,
LORD,
direct in peace,

Christianly, acceptably, and, should it please Thee,

painlessly,

gathering us together under the feet of Thine elect,

when Thou wilt and how Thou wilt,

only without shame and sins.

The brightness of the LORD
our GOD
be upon us,

prosper Thou the work of our hands upon us,

O prosper Thou our handy work.

Be, LORD,

within me to strengthen me,

without me to guard me,

over sue to shelter me,

beneath me to stablish me,

before me to guide me,

after me to forward me,

round about me to secure me. 

(5) \emph{Praise}.

Blessed art Thou, LORD, GOD
of Israel,

our FATHER,

from everlasting to everlasting.

Thine, O LORD,

is the greatness and the power,

the triumph and the victory,

the praise and the strength,

for Thou rulest over all

in heaven and on earth.

At Thy face every king is troubled,

and every nation.

Thine, O LORD,
is the kingdom

and the supremacy over all,

and over all rule.

With Thee is wealth, and glory is from Thy countenance;

Thou rulest over all, O LORD,

the Ruler of all rule;

and in Thine hand is strength and power,

and in Thine hand to give to all things

greatness and strength.

And now, LORD,
we confess to Thee,

and we praise Thy glorious Name.

\xchapter{The Fifth Day}

\emph{Introduction}.

We are satisfied with Thy
mercy, O LORD,

in the morning.

Blessed art Thou, O LORD,

who broughtest forth from the water

creeping things of life,

and whales,

and winged fowl.

Be Thou exalted, O GOD, above the heavens,

and Thy glory above all the earth.

By Thy Ascension, O LORD,

draw us too after Thee,

that we savour of what is above,

not of things on the earth.

By the marvellous mystery

of the Holy Body and precious Blood,

on the evening of this day,

LORD,
have mercy.

(1) \emph{Confession}.

Thou who hast said,

``As I live, saith the LORD,

I will not the death of a sinner,

but that the ungodly return from his way

and live;

turn ye, turn ye from your wicked way,

for why will ye die, O house of Israel?''

turn us, O LORD,
to Thee,

and so shall we be turned. 

Turn us from all our ungodlinesses,

and let them not be to us for punishments.

I have sinned, I have committed iniquity, I have done wickedly,

from Thy precepts, and Thy judgments.

To Thee, O LORD,
righteousness,

and to me confusion of face,

as at this day,

in our despicableness, wherewith Thou hast despised us.

LORD,
to us confusion of face,

and to our rulers

who have sinned against Thee.

LORD,
in all things is Thy righteousness,

unto all Thy righteousness;

let then Thine anger and Thy fury be turned away,

and cause Thy face to shine

upon Thy servant.

O my God, incline Thine ear and hear,

open Thine eyes and see my desolation.

O LORD
hear, O LORD
forgive, O LORD
hearken and do;

defer not for Thine own sake, O my GOD,

for Thy servant is called by Thy Name.

In many things we offend all;

LORD,
let Thy mercy rejoice against Thy judgment

in my sins.

If I say I have no sin, I deceive myself,

and the truth is not in me;

but I confess my sins many and grievous,

and Thou, O LORD,
art faithful and just,

to forgive me my sins when I confess them.

Yea, for this too

I have an Advocate with Thee to Thee,

Thy Only-begotten SON,
the Righteous.

May He be the propitiation for my sins,

who is also for the whole world.

Will the LORD
cast off for ever?

and will He be no more intreated? 

Is His mercy clean gone for ever?

and is His promise come utterly to an end for evermore?

Hath GOD
forgotten to be gracious?

and will He shut up His loving kindness in displeasure?

And I said, It is mine own infirmity;

but I will remember the years of the right hand

of the most Highest.

(2) \emph{Prayer for grace}.

[Give me grace]

to put aside every weight,

and the sin that doth so easily beset us;

all filthiness

and superfluity of naughtiness,

lust of the flesh, of the eyes,

pride of life,

every motion of flesh and spirit

alienated from the will of Thy sanctity:

to be poor in spirit,

that I have a portion in the kingdom of heaven;

to mourn, that I be comforted;

to be meek, that I inherit the earth;

to hunger and thirst for righteousness, that I be filled;

to be pitiful, that I be pitied;

to be pure in heart, that I see GOD;

to be a peace-maker that I be called the son of GOD;

to be prepared for persecutions and revilings

for righteousness' sake,

that my reward be in heaven,—

all this, grant to me, O LORD.

(3) \emph{Profession}.

I, coming to G\textsc{od},

believe that He is,

and that He is a rewarder of them

that diligently seek Him.

I know that my Redeemer liveth,

that He is CHRIST,
the SON
of the Living GOD,

that He is truly the SAVIOUR
of the world, 

that He came into the world to save sinners,

of whom I am chief.

Through the grace of JESUS
CHRIST

we believe that we shall be saved

like as our fathers.

I know that my skin shall rise up upon the earth,

which undergoes these things.

I believe to see the goodness of the LORD

in the land of the living.

Our heart shall rejoice in Him,

because we have hoped in His holy Name,

in the Name of the Father,

of the Saviour, Mediator, Intercessor, Redeemer,

of the two-fold Comforter,

under the figures of the Lamb and the Dove.

Let Thy merciful kindness, O LORD, be upon us,

like as we do put our trust in Thee.

(4) \emph{Intercession}.

Let us beseech the LORD in peace,

for the heavenly peace,

and the salvation of our souls;—

for the peace of the whole world;

for the stability of GOD'S
holy Churches,

and the union of them all;—

for this holy house,

and those who enter it with faith and reverence;

for our holy Fathers,

the honourable Presbytery, the Diaconate in CHRIST,

and all, both clergy and people;—

for this holy retreat, and all the city and country,

and all the faithful who dwell therein;—

for salubrious weather, fruitfulness of earth,

and peaceful times;—

for voyagers, travellers,

those who are in sickness, toil, and captivity,

and for their salvation.

Aid, save, pity, and preserve them,

O GOD,
in Thy grace. 

Making mention

of the all-holy, undefiled, and wore than blessed

Mary, Mother of GOD
and Ever-Virgin,

with all saints,

let us commend ourselves, and each other, and all our life,

to CHRIST
our GOD.

To Thee, O LORD,
for it is fitting,

be glory, honour, and worship.

The grace of our LORD,
JESUS
CHRIST,

and the love of GOD,

and the communion of the HOLY GHOST,

be with me, and with all of us. Amen.

I commend me and mine, and all that belongs to me,

to Him who is able to keep me without falling,

and to place me immaculate

before the presence of His glory,

to the only wise GOD
and our SAVIOUR;

to whom be glory and greatness,

strength and authority,

both now and for all ages. Amen.

(5) \emph{Praise}.

O LORD,
my LORD,

for my being, life, reason,

for nurture, protection, guidance,

for education, civil rights, religion,

for Thy gifts of grace, nature, fortune,

for redemption, regeneration, catechising,

for my call, recall, yea, many calls besides;

for Thy forbearance, long-suffering, long long-suffering

to me-ward,

many seasons, many years, up to this time;

for all good things received, successes granted me, good things
done;

for the use of things present,

for Thy promise, and my hope

of the enjoyment of good things to come;

for my parents honest and good,

teachers kind,

benefactors never to be forgotten, 

religious intimates congenial,

hearers thoughtful,

friends sincere,

domestics faithful,

for all who have advantaged me

by writings, homilies, converse,

prayers, patterns, rebukes, injuries;

for all these, and all others

which I know, which I know not,

open, hidden,

remembered, forgotten,

done when I wished, when I wished not,

I confess to Thee and will confess, I bless Thee and will bless,

I give thanks to Thee and will give thanks,

all the days of my life.

Who am I, or what is my father's house,

that Thou shouldst look upon a dead dog,

the like of me?

What reward shall I give unto the LORD

for all the benefits which he hath done unto me?

What thanks can I recompense unto GOD,

for all He has spared and borne with me until now?

Holy, Holy, Holy,

worthy art Thou,

O LORD
and our GOD,
the Holy One,

to receive the glory, and the honour, and the power,

for Thou hast made all things,

and for thy pleasure they are,

and were created.

\xchapter{The Sixth Day}

\emph{Introduction}.

Early shall my prayer come before Thee.

Blessed art Thou, O LORD,

who broughtest forth of the earth, wild beasts, cattle,

and all the reptiles,

for food, clothing, help;

and madest man after Thine image, to rule the earth,

and blessedst him.

The fore-counsel, fashioning hand,

breath of life, image of GOD,

appointment over the works,

charge to the Angels concerning him,

paradise.—

Heart, reins, eyes, ears, tongue, hands, feet,

life, sense, reason, spirit, free will, memory, conscience,

the revelation of GOD, writing of the law,

oracles of prophets, music of psalms,

instruction of proverbs, experience of histories,

worship of sacrifices.

Blessed art Thou, O LORD,

for Thy great and precious promise

on this day

concerning the Life-giving Seed,

and for its fulfilment in fulness of the times

at this day.

Blessed art Thou, O LORD,

for the holy passion

of this day. 

O by thy salutary sufferings

on this day,

save us, O LORD.

(1) \emph{Confession}.

I have withstood Thee, LORD,

but I return to Thee; (\emph{Hosea}.)

for I have fallen by mine iniquity.

But I take with me words,

and I return unto Thee and say,

take away all iniquity and receive us graciously,

so will we render the calves of our lips.

Spare us, LORD,
spare,
(\emph{Joel}.)

and give not Thine heritage to reproach,

to Thine enemies.

LORD,
LORD,
be propitious,

cease, I beseech Thee,
(\emph{Amos}.)

by whom shall Jacob arise?

for he is small.

Repent, O LORD,
for this,

and this shall not be.

While observing lying vanities
(\emph{Jonah}.)

I forsook my own mercy,

and am cast out of Thy sight.

When my soul fainted within me,

I remembered the LORD;

yet will I look again toward Thy Holy Temple;

Thou hast brought up my life from corruption.

Who is a G\textsc{od} like unto Thee,
(\emph{Micah}.)

that pardoneth iniquity

to the remnant of His heritage?

He retaineth not His anger for ever,

because He delighteth in mercy.

Turn again and have compassion upon us, O LORD,

subdue our iniquities,

and cast all our sins into the depths of the sea,

according to Thy truth, and according to Thy mercy. 

O LORD,
I have heard thy speech and was afraid,
(\emph{Habakkuk}.)

in wrath remember mercy.

Behold me, LORD,
clothed in filthy garments;
(\emph{Zechariah}.)

behold Satan standing at my right hand;

yet, O LORD,
by the blood of Thy covenant,

by the fountain opened for sin and for uncleanness,

take away my iniquity,

and cleanse me from my sin.

Save me as a brand

plucked out of the fire.

Father, forgive me, for I knew not,

truly I knew not, what I did

in sinning against Thee.

LORD,
remember me

when Thou comest in Thy kingdom.

LORD,
lay not mine enemies' sins to their charge,

LORD,
lay not my own to mine.

By Thy sweat bloody and heavy,

Thy soul in agony,

Thy head crowned with thorns, bruised with staves,

Thine eyes swimming with tears,

Thine ears full of insults,

Thy mouth moistened with vinegar and gall,

Thy face dishonourably stained with spitting,

Thy neck weighed down with the burden of the cross,

Thy back ploughed with the wheals and gashes of the scourge,

Thy hands and feet stabbed through,

Thy strong cry, Eli, Eli,

Thy heart pierced with the spear,

the water and blood thence flowing,

Thy body broken,

Thy blood poured out,

LORD,
forgive the offence of Thy servant,

and cover all his sins.

Turn away all Thy displeasure,

and turn Thyself from Thy wrathful indignation. 

Turn me then, O GOD
our SAVIOUR,

and let Thine anger cease from us.

Wilt Thou be displeased at us for ever,

and stretch out Thy wrath from one generation to another?

Wilt Thou not turn again and quicken us,

that Thy people may rejoice in Thee?

Show us Thy mercy, O LORD,

and grant us Thy salvation.

(2) \emph{Prayer for grace}.

\textbf{.
. . .}

the works of the flesh,

adultery, fornication, uncleanness, lasciviousness,

idolatry, witchcraft,

enmities, strifes,

emulations, heats,

quarrels, parties,

heresies, envyings, murders,

drunkennesses, revellings, and such like.

\textbf{.
. . .}

the fruits of the SPIRIT,

love, joy, peace,

long-suffering, gentleness, goodness,

faith, meekness, temperance;

the spirit of wisdom, of understanding.

of counsel, of might,

of knowledge, of godliness,

of fear of the LORD:—

and the gifts of the SPIRIT,

the word of wisdom, of knowledge,

faith, gifts of healing, working of miracles,

prophecy, discerning of spirits,

kinds of tongues, interpretation of tongues.

May Thy strong hand, O LORD
\footnote{Vide p. 146, edit. 1675.},

become my defence;

Thy mercy in CHRIST

my salvation; 

Thy all-veritable word,

my instructor;

the grace of Thy life-bringing Spirit,

my consolation,

all along, and at last.

The Soul of CHRIST
hallow me,

and the Body strengthen me,

and the Blood ransom me,

and the Water wash me,

and the Bruises heal me,

and the Sweat refresh me,

and the Wound hide me.

The peace of GOD

which passeth all understanding,

keep my heart and thoughts

in the knowledge and the love of GOD.

(3) \emph{Profession}.

I believe

that Thou hast created me;

despise not the work of Thine own hands;—

that Thou madest me after Thine image and likeness,

suffer not Thy likeness to be blotted out;—

that Thou hast redeemed me in Thy blood,

suffer not the cost of that redemption to perish;

that Thou hast called me Christian after Thy name,

disdain not Thine own title;

that Thou hast hallowed me in regeneration,

destroy not Thy holy work;—

that Thou hast grafted me into the good olive tree,

the member of a mystical body;

the member of Thy mystical body

cut not off.

O think upon Thy servant as concerning Thy word,

wherein Thou hast caused me to put my trust.

My soul hath longed for Thy salvation,

and I have good hope because of Thy word. 

(4) \emph{Intercession}.

[I pray]

for the prosperous advance and good condition

of all the Christian army,

against the enemies of our most holy faith;

for our holy fathers,

and all our brotherhood in CHRIST;

for those who hate and those who love us,

for those who pity and those who minister to us;

for those whom we have promised

to remember in prayer;

for the liberation of captives;

for our fathers and brethren absent;

for those who voyage by sea;

for those who lie in sickness.

Let us pray also for fruitfulness of the earth;

and for every soul of orthodox Christians.

Let us bless pious kings,

orthodox high-priests,

the founders of this holy retreat,

our parents,

and all our forefathers

and our brethren departed.

(5) \emph{Praise}.

Thou who, on man's
trangressing Thy command,

and falling,

didst not pass him by, nor leave him, GOD of goodness;

but didst visit in ways manifold,

as a tender Father,

supplying him with Thy great and precious promise,

concerning the Life-giving Seed,

opening to him the door of faith,

and of repentance unto life,

and in fulness of the times,

sending Thy CHRIST
Himself

to take on Him the seed of Abraham;

and, in the oblation of His life,

to fulfil the Law's obedience; 

and, in the sacrifice of His death,

to take off the Law's curse;

and, in His death,

to redeem the world;

and, in His resurrection,

to quicken it:—

O Thou, who doest all things,

whereby to bring again our race to Thee,

that it may be partaker

of Thy divine nature and eternal glory;

who hast borne witness

to the truth of Thy gospel

by many and various wonders,

in the ever-memorable converse of Thy saints,

in their supernatural endurance of torments,

in the overwhelming conversion of all lands

to the obedience of faith,

without might, or persuasion, or compulsion:—

Blessed be Thy Name,

and praised and celebrated,

and magnified, and high exalted,

and glorified, and hallowed;

its record, and its memory,

and every memorial of it,

both now and for evermore.

Worthy art Thou to take the book,

and to open the seals thereof,

for Thou wast slain, and hast redeemed us to GOD

by Thy blood,

out of every kindred and tongue,

and people, and nation.

Worthy is the Lamb that was slain

to receive the power, and riches, and wisdom,

and strength, and honour, and glory, and blessing.

To Him that sitteth upon the Throne,

and to the Lamb,

be the blessing, and the honour, and the glory, and the might,

for ever and ever. Amen.

Salvation to our GOD,
which sitteth upon the throne,

and to the Lamb.

Amen : the blessing and the glory and the wisdom,

and the thanksgiving and the honour,

and the power and the strength,

be unto our GOD,

for ever and ever,

Amen.

\xchapter{The Seventh Day}

\emph{Introduction}.

O LORD,
be gracious unto us,

we have waited for Thee;

be Thou our arm every morning,

our salvation also in the time of trouble.

Blessed art Thou, O LORD,

who restedst on the seventh day

from all Thy works,

and blessedst and sanctifiedst it:

[concerning the Sabbath,

concerning the Christian rest instead of it,

concerning the funeral rites of CHRIST,

and the resting from sin,

concerning those who are already gone to rest.]

(1) \emph{Confession}.

I am ashamed, and blush, O my GOD,

to lift up my face to Thee,

for mine iniquities are increased

over my head,

and my trespass is grown up unto the heavens;

since the days of youth

have I been in a great trespass

unto this day;

I cannot stand before Thee because of this.

My sins are more in number than the sand of the sea,

my iniquities are multiplied,

and I not worthy to look up

and see the height of heaven,

from the number of my unrighteousnesses; 

and I have no relief,

because I have provoked Thine anger,

and done evil in Thy sight;

not doing Thy will,

not keeping Thy commandments.

And now my heart kneels to Thee,

beseeching Thy goodness.

I have sinned, O LORD,
I have sinned,

and I know mine iniquities;

and I ask and beseech,

remit to me, O LORD,
remit to me,

and destroy me not in mine iniquities;

nor be Thou angry for ever,

nor reserve evil for me;

nor condemn me

in the lowest parts of the earth.

Because Thou art GOD,
the GOD
of penitents,

and Thou shalt show in me all Thy loving-kindness;

for Thou shalt save me unworthy,

according to Thy much pity,

and I will praise Thee alway.

LORD,
if Thou wilt, Thou canst cleanse me;

LORD,
only say the word, and I shall be healed.

LORD,
save me;

Carest Thou not that we perish?

Say to me, Be of good cheer, thy sins are remitted to thee.

JESU,
Master, have mercy on me;

Thou Son of David, JESU,
have mercy on me;

JESU,
Son of David, Son of David.

LORD,
say to me, Ephphatha.

LORD,
I have no man \footnote{John v. 7.};

LORD,
say to me, be loosed from thine infirmity.

Say unto my soul, I am thy salvation.

Say unto me, My grace is sufficient for thee.

LORD,
how long wilt Thou be angry?

shall Thy jealousy burn like fire for ever? 

O, remember not our old sins;

but have mercy on us and that soon,

for we are come to great misery;

Help us, O GOD
of our salvation;

for the glory of Thy Name.

O deliver us and be merciful unto our sins,

for Thy Name's sake.

(2) \emph{Prayer for grace}.

[O LORD,
remit]

all my failings, shortcomings, falls,

offences, trespasses, scandals,

transgressions, debts, sins,

faults, ignorances, iniquities,

impieties, unrighteousnesses, pollutions.

The guilt of them,

be gracious unto,

remit,

be propitious unto,

impute not,
pardon;

forgive;

spare;

charge not, remember not.

The stain,

pass by,

disregard,

hide,

blot out,
pass over;

overlook;

wash away;

cleanse.

The hurt,

remit,

take off,

abolish,
heal,

remove, 

annul,
remedy;

away with;

disperse, annihilate;

that they be not found, that they exist not.

Supply

to faith,

to virtue,

to knowledge,

to continence,

to patience,

to godliness,

to brotherly love,
virtue;

knowledge;

continence;

patience;

godliness;

brotherly love;

charity.

That I forget not my cleansing from my former
sins,

but give diligence to make my calling and election sure

through good works.

(3) \emph{Profession}.

I believe in Thee the FATHER;

Behold then, if Thou a Father and we sons,

as a father pitieth sons,

be Thou of tender mercy towards us, O LORD.

I believe in Thee, the LORD;

behold Then, if Thou art LORD and we servants,

our eyes are upon Thee our LORD,

until Thou have mercy upon us.

I believe, that though we be neither sons nor servants,

but dogs only,

yet we have leave to eat of the crumbs

that fall from Thy Table.

I believe that CHRIST
is the Lamb of GOD;

O Lamb of GOD
that takest away the sins of the world,

take Thou away mine.

I believe that JESUS
CHRIST
came into the world

to save sinners;

Thou who earnest to save sinners

save Thou me, of sinners

chief and greatest.

I believe that CHRIST
came to save what was lost;

Thou who earnest to save the lost,

never suffer, O LORD,
that to be lost which Thou hast saved.

I believe that the SPIRIT
is the Lord and Giver of life;

Thou who gavest me a living soul,

give me that I receive not my soul in vain.

I believe that the SPIRIT
gives grace

in His sacred things;

give me that I receive not His grace in vain,

nor hope of His sacred things.

I believe that the SPIRIT
intercedes for us

with plaints unutterable;

grant me of His intercession and those plaints

to partake, O LORD.

Our fathers hoped in Thee,

they trusted in Thee, and Thou didst deliver them.

They called upon Thee and were holpen,

they put their trust in Thee, and were not confounded.

As Thou didst our fathers

in the generations of old,

so also deliver us, O LORD,

who trust in Thee.

(4) \emph{Intercession}.

O Heavenly KING,

confirm our faithful kings,

stablish the faith,

soften the nations,

pacify the world,

guard well this holy retreat,

and receive us in orthodox faith and repentance,

as a kind and loving Lord.

The power of the FATHER
guide me,

the wisdom of the SON
enlighten me,

the working of the SPIRIT
quicken me.

Guard Thou my soul,

stablish my body,

elevate my senses,

direct my converse,

form my habits,

bless my actions,

fulfil my prayers,

inspire holy thoughts,

pardon the past,

correct the present,

prevent the future.

(5) \emph{Praise}.

Now unto Him that is able to do

exceeding abundantly

above all that we ask or think,

according to the power that worketh in us. 

to Him be glory

in the Church in CHRIST

unto all generations

world without end. Amen.

Blessed, and praised, and celebrated,

and magnified, and exalted, and glorified, and hallowed,

be Thy Name, O LORD,

its record, and its memory,

and every memorial of it;

for the all-honourable senate of the Patriarchs,

the ever-venerable band of the Prophets,

the all-glorious college of the Apostles,

the Evangelists,

the all-illustrious army of the Martyrs,

the Confessors,

the assembly of Doctors,

the Ascetics,

the beauty of Virgins,

for Infants the delight of the world,—

for their faith,

their labours,

their blood,

their diligence,

their purity,

their hope,

their truth,

their zeal,

their tears,

their beauty.

Glory to Thee, O LORD, glory to Thee,

glory to Thee who didst glorify them,

among whom we too glorify Thee.

Great and marvellous are Thy works,

LORD,
the GOD
ALMIGHTY;

just and true are Thy ways,

O King of Saints.

Who shall not fear Thee, O LORD,

and glorify Thy Name?

for Thou only art holy.

for all the nations shall come and worship before Thee,

for Thy judgments are made manifest. 

Praise our GOD, all ye His servants,

and ye that fear Him,

both small and great.

Alleluia,

for the LORD
GOD
Omnipotent reigneth;

let us be glad and rejoice, and give honour to Him.

Behold the tabernacle of GOD is with men,

and He will dwell with them;

and they shall be His people,

and GOD
Himself shall be with them,

and shall wipe away all tears from their eyes.

And there shall be no more death;

neither crying, neither pain any more,

for the former things are passed away.

\xchapter{Deprecations}

\xsection{1.}

O LORD
\footnote{Vide p. 92, edit. 1675.}, Thou knowest, and canst,
and willest

the good of my soul.

Miserable man am I;

I neither know, nor can, nor, as I ought,

will it.

Thou, O LORD,
I beseech Thee,

in Thine ineffable affection,

so order concerning me,

and so dispose,

as Thou knowest to be most pleasing to Thee,

and most good for me.

[Thine is]

goodness, grace;

love, kindness;

benignity, gentleness, consideration;

forbearance, long suffering;

much pity, great pity;

mercies, multitude of mercies, yearnings of mercies;

kind yearnings, deep yearnings;

in passing over,

in overlooking, in disregarding;

many seasons, many years;

[punishing] unwillingly, not willingly;

not to the full,

not correspondently,

in wrath remembering mercy,

repenting of the evil, 

compensating doubly,

ready to pardon,

to be reconciled,

to be appeased.

———————

\xsection{2. Litany}

\footnote{Page 180, edit. 1675.}

FATHER,
the Creator,

SON, the
Redeemer,

SPIRIT,
the Regenerator,

destroy me not,

whom Thou hast created, redeemed, regenerated.

Remember not, LORD,
my sins,

nor the sins of my forefathers;

neither take vengeance for our sins, theirs, nor mine.

Spare us, LORD,
them and me \footnote{Thus in St. Gregory's Sacramentary.—``Præsta, quæsumus,

omnipotens Deus, ut animam famuli, \&c. … in congregatione

justorum æternæ beatitudinis jubeas esse consortem. Per

Dominum, \&c.'' `` Præsta, \&c. ut animam, \&c. ab angelis
lucis

susceptam in præparata habitacula deduci facias beatorum.

Per Dominum, \&c'' ``Suscipe, Domine, preces nostras, pro

anima famuli tui ... ut si quæ ei maculæ de terrenis contagiis

adhæserunt, remissionis tuæ remedio deleantur.'' For similar

language in the second century vide Ussher as quoted in

Tract 72.},

spare Thy people,

and, among Thy people, Thy servant,

who is redeemed with Thy precious blood;

and be not angry with us for ever.

Be merciful, be merciful; spare us, LORD,

and be not angry with us for ever.

Be merciful, be merciful; have pity on us, LORD,

and be not angry with us to the full. 

Deal not, O LORD,

deal not with me after mine iniquities,

neither recompense me according to my sins;

but after Thy great pity,

deal with me,

and according to the multitude of Thy mercies,

recompense me;

after that so great pity,

and that multitude of mercies,

as Thou didst to our fathers

in the times of old:—

by all that is dear unto Thee.

From all evil and adversity,

in all time of need;

from this evil and this adversity,

in this time;

raise me, rescue me, save me, O LORD.

Deliver me, O LORD,

and destroy me not.

On the bed of sickness;

in the hour of death;

in the day of judgment,

in that dreadful and fearful day,

rescue me, LORD,
and save me;—

from seeing the Judge's face overcast,

from being placed on the left,

from hearing the dreadful word, Depart from Me,

from being bound in chains of darkness,

from being cast into the outer darkness,

from being tormented in the pit of fire and brimstone,

where the smoke of the torments

ascendeth for ever.

Be merciful, be merciful,

spare us, pity us,

O LORD:

and destroy us not for ever,

deliver and save us.

Let it not be, O LORD;
and that it be not,

take away from me, O LORD,

hardness of heart,

desperateness after sinning,

blindness of heart,

contempt of Thy threats,

a cauterized conscience,

a reprobate mind,

the sin against the HOLY
GHOST,

the sin unto death,

the four crying sins \footnote{Wilful murder, the sin of Sodom, oppressing the poor,

defrauding workmen of their wages.};

the six which forerun \footnote{Despair of salvation, presumption of God's mercy,

impugning known truth, envy at another's grace, obstinacy

in sin, and impenitence.}

the sin against the HOLY
GHOST.

Deliver me

from all ills and abominations of this world,

from plague, famine, and war;

earthquake, flood, and fire,

the stroke of immoderate rain and drought,

blast and blight;

thunder, lightning and tempest;

epidemic sickness, acute and malignant,

unexpected death;

from ills and difficulties in the Church,

from private interpretation,

from innovation in things sacred,

from heterodox teaching;

from unhealthy inquiries and interminable disputes,

from heresies, schisms, scandals,

public and private,

from making gods \footnote{[\emph{T\={e}s apothe\={o}se\={o}s}], vid. Acts xii. 22. Mr.
Waller ``going to

see the king at dinner, overheard a very extraordinary

conversation between his Majesty [King James] and two prelates,

the Bishop of Winchester [Andrews] and Dr. Neale, Bishop of

Durham, who were standing behind the king's chair. His majesty

asked the bishops, 'My lords, cannot I take my subjects' money

when I want it, without all this formality in parliament?' The Bishop

of Durham readily answered, 'God forbid, Sir, but you should;

you are the breath of our nostrils.' Whereupon the king turned

and said to the Bishop of Winchester, 'Well, my lord, what say

you?' 'Sir,' replied the bishop, 'I have no skill to judge of

Parliamentary cases.' The king answered, 'No put-off's, my lord,

answer me presently.' 'Then, Sir,' said he, 'I think it lawful for

you to take my brother Neale's money, for he offers it.'''—

Waller's Life, quoted in Biograph. Brit.} of kings,

from flattering of the people,

from the indifference of Saul,

from the scorn of Michal,

from the greediness of Hophni,

from the plunder of Athaliah,

from the priesthood of Micah,

from the brotherhood of Simon and Judas,

from the doctrine of men unlearned and unestablished,

from the pride of novices,

from the people resisting the priest:—

from ills and difficulties in the state,

from anarchy, many rulers, tyranny,

from Asher, Jeroboam, Rehoboam, Gallio, Haman,

the profligacy of Ahithophel,

the foolishness of Zoan \footnote{Isai. xix.},

the statutes of Omri,

the justice of Jezebel,

the overflowings of Belial \footnote{Ps. xviii. 4.},

the courage of Peor,

the valley of Achor,

pollution of blood or seed,

incursion of enemies,

civil war,

bereavement of good governors,

accession of evil and unprincipled governors; 

from an intolerable life,

in despondence, sickness, ill-fame,

distress, peril, slavery, restlessness:

from death

in sin, shame, tortures,

desperateness, defilement, violence, treachery;

from death unexpected,

from death eternal.

\xchapter{Forms of Intercession}

\xsection{1.}

\footnote{Vide p. 90, edit. 1675.}

For all creatures,

men,

persons compassed

with infirmity,

Churches

Catholic,

Eastern,

Western,

British.

The Episcopate,

Presbytery,

clergy,

Christian people.

States

of the whole earth,

Christian,

neighbouring,

our own,

Rulers,

kings,

religious kings,

our own,

councillors,

judges,

nobles,

soldiers,

sailors,

the people,

the rising generation,

schools,

those at court,

in cities,

the country.

Those who serve the soul;

those who serve the body,

in food,

clothing,

health,

necessaries.

[Those who have a claim on

my prayers,]

in nature,

by benefits,

from trust,

formerly or now,

in friendship,

in love,

in neighbourhood;

from promise,

from mutual offices,

from want of leisure,

from destitution,

from extremity.

\xsection{2.}

\footnote{Page 170, edit. 1675.}

Thy whole creation,

our whole race,

the states of the world,

the Catholic Church,

the separate Churches,

the separate states,

our Church,

our state,

the orders in each,

the persons in the orders,

the world,

the inhabited earth,

the Christian religion,

our country,

the priesthood,

the person of the King, of the Prince,

the City,

the parish in which I was baptized, All Hallows, Barking.

My two schools,

my University,

my College,

the parish committed to me, St. Giles's,

the three Churches

of Southwell,

St. Paul's,

Westminster;

the three Dioceses of Chichester,

Eli,

Winton,

my home,

my kindred

those who show me pity,

those who minister to me;

my neighbours,

my friends,

those who have a claim on me.

\xsection{3.}

\footnote{Page 206, edit. 1675.}

The
creation,

the race of man,

all in affliction and in prosperity,

in error, 

in sin,
and in truth,

and in grace;

the Church Ecumenical,

Eastern, Western, our own,

Rulers, Clergy, people.

States of the earth,

Christian, neighbouring, our own,

the King, the Queen, the Prince,

the nobles.

Parliament, Law Courts, army, police.

The Commons,

farmers, merchants, artisans,

down to mean workmen,

and poor.

Those who have a claim on me,

from kindred,

benefaction,

ministration of things temporal,

charge formerly or now,

natural kindness,

Christian love,

neighbourhood,

promise on my part,

their own desire,

their lack of leisure,

sympathy for their extreme misery;

any good work,

any noble action,

any scandal from me,

having none to pray for them.

\xsection{4.}

\footnote{Page 210, edit. 1675.}

World,

Church

throne,

Council-chamber, 

schools,

Infants,

the grown,

men,

aged,

The possessed,

sick,

orphans,

foreigners,

travellers,

with child

in bitter bondage,

overladen.

earth inhabited.

kingdom,

altar.

law courts,

work-places.

boys,

youths,

elderly,

decrepit.

weak-hearted,

prisoners,

widows,

voyagers,

who give suck,

in desolateness,

\xchapter{Meditations}

\xsection{1. On Christian Duty}

What shall I do that I may inherit eternal life?

Keep the commandments. \emph{Mark}
x. 17.

What shall we do?

Repent and be baptized every one of
you. \emph{Acts} ii. 37, 38.

What must I do to be saved?

Believe on the LORD
JESUS CHRIST.
\emph{Acts} xvi. 31.

What shall we do then?

(\emph{To the multitude}.) He who hath two coats, let him impart to

him that hath none.

He that hath meat, let him do
likewise.

(\emph{To the publicans}.) Exact no more than is appointed you.

(\emph{To soldiers}.) Do violence to no man; neither accuse any

falsely;

be content with your wages. \emph{Luke}
iii. 10-14.

The knowledge and faith

of [GOD'S]
justice
[GOD'S]
mercy 

[leads] unto

fear,

abasement,

repentance,

fasting,

prayers,

patience,

a sacrifice \footnote{\emph{E.g}. Concede nobis, Domine Deus noster, ut hæc

hostia salutaris, et nostrorum fiat purgatio delictorum,

et tuæ propitiatio Majestatis, per, \&c.—St. Gregory's

Lib. Sacrament. p. 57.},
hope,

consolation,

thanksgiving,

almsgiving,

hymns,

obedience,

an oblation.

\xsection{2. On the Day of Judgment}

FATHER
Unoriginate, Only-begotten SON,

Life-giving SPIRIT,

merciful, pitiful, long-suffering,

full of pity, full of kind yearnings,

who lovest the just and pitiest the sinful,

who passest by sins and grantest petitions,

GOD of
penitents,

SAVIOUR
of sinners,

I have sinned before Thee, O LORD,

and thus and thus have I done.

Alas, alas! woe, woe.

How was I enticed by my own lust!

How I hated instruction!

Nor felt I fear nor shame

at Thy incomprehensible glory,

Thy awful presence,

Thy fearful power,

Thy exact justice,

Thy winning goodness.

I will call if there be any that will answer me;

to which of the Saints shall I turn?

O wretched man that I am,

who shall deliver me from the body of this death?

how fearful is Thy judgment, O LORD?

when the thrones are set

and Angels stand around,

and men are brought in,

the books opened,

the works inquired into,

the thoughts examined

and the hidden things of darkness.

What judgment shall be upon me?

who shall quench my flame?

who shall lighten my darkness,

if Thou pity me not? 

LORD, as
Thou art loving,

give me tears,

give me floods, give me today.

For then will be the incorruptible Judge,

the horrible judgment-seat,

the answer without excuses,

the inevitable charges,

the shameful punishment,

the endless Gehenna,

the pitiless Angels,

the yawning hell,

the roaring stream of fire,

the unquenchable flame,

the dark prison,

the rayless darkness,

the bed of live coals,

the unwearied worm,

the indissoluble chains,

the bottomless chaos,

the impassable wall,

the inconsolable cry,

none to stand by me,

none to plead for me,

none to snatch me out.

But I repent, LORD,
O LORD,
I repent,

help Thou mine impenitence,

and more, and still more,

pierce, rend, crush my heart.

Behold, O LORD,
that I am

indignant with myself,

for my senseless, profitless,

hurtful, perilous passions;

that I loathe myself,

for these inordinate, unseemly,

deformed, insincere,

shameful, disgraceful

passions, 

that my confusion is daily before me,

and the shame of my face hath covered me.

Alas! woe, woe—

O me, how long?

Behold, LORD,
that I sentence myself

to punishment everlasting,

yea, and all miseries of this world.

Behold me, LORD,
self-condemned;

Behold, LORD,
and enter not into judgment with Thy servant.

And now, LORD,

I humble myself under Thy mighty hand,

I bend to Thee, O LORD,
my knees,

I fall on my face to the earth.

Let this cup pass from me!

I stretch forth my hands unto Thee;

I smite my breast, I smite on my thigh.

Out of the deep my soul crieth unto Thee,

as a thirsty land;

and all my bones,

and all that is within me.

LORD,
hear my voice.

\xsection{3. On Human Frailness}

Have mercy on me, LORD, for I am weak;

remember, LORD,
how short my time is;

remember that I am but flesh,

a wind that passeth away, and cometh not again.

My days are as grass, as a flower of the field;

for the wind goeth over me, and I am gone,

and my place shall know me no more.

I am dust and ashes,

earth and grass,

flesh and breath,

corruption and the worm,

a stranger upon the earth,

dwelling in a house of clay, 

few and evil my days,

today, and not tomorrow,

in the morning, yet not until night,

in a body of sin,

in a world of corruption,

of few days, and full of trouble,

coming up, and cut down like a flower,

and as a shadow, having no stay.

Remember this, O LORD,
and suffer, remit;

what profit is there in my blood,

when I go down to the pit?

By the multitude of Thy mercies,

by the richess and excessive redundance

of Thy pity;

by all that is dear to Thee,

all that we should plead,

and before and beyond all things, by Thyself,

by Thyself, O LORD,
and by Thy CHRIST.

LORD,
have mercy upon me, the chief of sinners.

O my LORD,
let Thy mercy rejoice

against Thy judgment in my sin.

O LORD,
hear, O LORD,
forgive,

O LORD,
hearken,

O LORD,
hearken and do,

do and defer not for Thine own sake,

defer not, O LORD
my GOD.
